\documentclass[reqno]{amsart}

%\documentclass[reqno]{article}
\usepackage{graphicx}
%\usepackage[backend=bibtex]{biblatex} \addbibresource{Jun.bib,band.bib}

\usepackage[margin=1in]{geometry}
%\documentclass{amsart}
\usepackage{amssymb, amsthm, mathtools}
\usepackage{amsfonts}
\usepackage{mathrsfs}
\usepackage{hyperref}
\usepackage{cleveref}
\usepackage{xcolor}
\usepackage{stmaryrd}
\SetSymbolFont{stmry}{bold}{U}{stmry}{m}{n}
\usepackage{comment}
%\usepackage{showkeys}
%\usepackage{showlabels}
\allowdisplaybreaks
\setcounter{tocdepth}{1}


\usepackage{array}
\usepackage{tikz}
\usetikzlibrary{arrows}
\usetikzlibrary{arrows.meta}
\usetikzlibrary{decorations.pathmorphing}
\usetikzlibrary{decorations.pathreplacing,decorations.markings}
\usetikzlibrary{calc}
\usetikzlibrary{shapes,petri, topaths, automata}
\usepackage{tikz-feynman} % for Feynman diagrams
\tikzfeynmanset{compat=1.1.0}
%/tikzfeynman/warn luatex=false

\tikzset{
  % style to apply some styles to each segment of a path
  on each segment/.style={
    decorate,
    decoration={
      show path construction,
      moveto code={},
      lineto code={
        \path [#1]
        (\tikzinputsegmentfirst) -- (\tikzinputsegmentlast);
      },
      curveto code={
        \path [#1] (\tikzinputsegmentfirst)
        .. controls
        (\tikzinputsegmentsupporta) and (\tikzinputsegmentsupportb)
        ..
        (\tikzinputsegmentlast);
      },
      closepath code={
        \path [#1]
        (\tikzinputsegmentfirst) -- (\tikzinputsegmentlast);
      },
    },
  },
  % style to add an arrow in the middle of a path
  mid arrow/.style={postaction={decorate,decoration={
        markings,
        mark=at position .5 with {\arrow[#1]{Stealth[scale=1.2]}}
      }}},
}

\DeclareMathOperator{\re}{{\mathrm{Re}}}
\DeclareMathOperator{\im}{{\mathrm{Im}}}
\renewcommand{\Dot}{\mathbf {Dot}}

\hypersetup{
     colorlinks=true
}
\pagestyle{plain}

\def\P{\mathbb{P}}
\def\E{\mathbb{E}}
\def\IE{\mathbb{IE}}
\def\Z{\mathbb{Z}}
\def\R{\mathbb{R}}
\def\N{\mathbb{N}}

\newcommand{\be}{\begin{equation}}
\newcommand{\ee}{\end{equation}}
\newcommand{\ii}{\mathrm{i}}
\newcommand{\dd}{\mathrm{d}}
\newcommand{\e}{{\varepsilon}}
\newcommand{\cK}{\mathcal K}
\newcommand{\cL}{\mathcal L}
\newcommand{\sig}{\sigma}
\newcommand{\bsig}{{\boldsymbol{\sigma}}}
\newcommand{\ba}{{\mathbf{a}}}
\newcommand{\bal}{{\boldsymbol{\alpha}}}
\newcommand{\bbeta}{{\boldsymbol{\beta}}}
\newcommand{\bfb}{{\mathbf{b}}}
\newcommand{\cut}{\mathcal{G}}
\newcommand{\cutL}{(\mathcal{G}_L)}
\newcommand{\cutR}{(\mathcal{G}_R)}
\newcommand{\lenk}{l_{\mathcal K}}
\DeclareMathOperator{\TSP}{{\mathbf T}}
% \oddsidemargin=0in
% \evensidemargin=0in
% \textwidth=6.5in
% \setlength{\unitlength}{1cm}
% \setlength{\parindent}{0.6cm}

\newcommand{\BA}{\mathrm{BA}}
\newcommand{\WO}{\mathrm{WO}}
\newcommand{\dist}{\mathrm{dist}}
\newcommand{\alt}{\mathrm{alt}}
\newcommand{\id}{\mathrm{Id}}
\newcommand{\fd}{{\mathfrak d}}
\newcommand{\bv}{{\bf{v}}}
\newcommand{\bw}{{\bf{w}}}
\newcommand{\ord}{{\mathsf{ord}}}
\newcommand{\size}{{\mathsf{size}}}
\newcommand{\Mol}{\mathcal{M}}
\newcommand{\new}{\mathrm{new}}
\newcommand{\aux}{\mathrm{aux}}
\newcommand{\ngh}{\mathrm{ng}}
\newcommand{\C}{\mathbb C}
\renewcommand{\bar}{\overline}
\newcommand{\wt}{\widetilde}
\newcommand{\wh}{\widehat}
\newcommand{\al}{\alpha}
\newcommand{\qqq}[1]{\llbracket{#1}\rrbracket}
\newcommand{\br}[1]{\left[{#1}\right]}
\newcommand{\p}[1]{\left({#1}\right)}
\newcommand{\abs}[1]{\left|{#1}\right|}
\newcommand{\avg}[1]{\left\langle{#1}\right\rangle}
\newcommand{\sB}{{\mathsf{B}}}
\newcommand{\fP}{\mathfrak{P}}
\newcommand{\bpi}{{\boldsymbol{\pi}}}
%\newcommand{\Zn}{\widetilde{\mathbb Z}_L^d}
\newcommand{\Zn}{{\mathbb Z}_L^d}
\newcommand{\ZL}{{\mathbb Z}_{WL}^d}
\newcommand{\Gc}{{\mathring G}}
\newcommand{\bD}{{\mathbf{D}}}
\newcommand{\cT}{\mathcal T}
\newcommand{\cE}{\mathcal E}
\newcommand{\wT}{\widetilde{\mathcal T}}
\newcommand{\zTheta}{\mathring \Theta}
\newcommand{\dashed}{\text{diffusive}}
% \newcommand{\pself}{\text{pseudo-self-energy}}
% \newcommand{\pselfs}{\text{pseudo-self-energies}}
% \newcommand{\Pself}{\text{Pseudo-self-energy}}
% \newcommand{\Pselfs}{\text{Pseudo-self-energies}}
% \newcommand{\pTexp}{\text{pseudo-$T$-expansion}}
% \newcommand{\PTexp}{\text{Pseudo-$T$-expansion}}
% \newcommand{\pTeq}{\text{pseudo-$T$-equation}}
% \newcommand{\PTeq}{\text{Pseudo-$T$-equation}}
\newcommand{\dthn}{\boldsymbol{\chi}}

% \def\mathcalC{\mathcal{C}}
% \def\mathcalV{\mathcal{V}}
% \def\mathcalW{\mathcal{W}}
\newcommand{\fa}{{\mathfrak a}}
\newcommand{\fc}{{\mathfrak c}}
\newcommand{\sM}{{\mathsf M}}
\newcommand{\Err}{{\mathcal Err}}

\DeclareMathOperator{\OO}{O}
\DeclareMathOperator{\oo}{o}

% \newcommand{\Sdelta}{{\Sigma}_T}
% \newcommand{\Sdeltak}{{\Sigma}_{T,k}}
% \newcommand{\wtSdelta}{\Sigma}
% \newcommand{\wtSdeltan}{\Sigma^{(n)}}
% \newcommand{\wtSdeltak}{\Sigma^{(k)}}
% \newcommand{\Sinfty}{\Sigma^{\infty}}
% \newcommand{\wtSinfty}{\wt\Sigma^{\infty}}

% \newcommand{\Sele}{{\mathcal E}}
% \newcommand{\Smax}{{\mathcal E}^{\max}}
% %\newcommand{\Selek}{{\mathcal E}^{(\omega)}}
% \newcommand{\Selek}{{\mathcal E}}
% \newcommand{\pSele}{{\mathcal P}}
% \newcommand{\Selekl}{{\mathcal E}^{[\eta]}}
% \newcommand{\Seleinf}{{\mathcal E}^{\infty}}
%\newcommand{\wtSeleinf}{\widetilde{\mathcal E}^{(\kappa),\infty}}
% \newcommand{\wtSeleinf}{ {\mathcal E}^{\infty}}
% \newcommand{\Sint}{\mathfrak S}
% \newcommand{\Sintk}{{\mathfrak S}^{(\omega)}}
% \newcommand{\Cele}{{\mathscr E}}
\newcommand{\bpsi}{\boldsymbol{\psi}}
% \newcommand{\zeroself}{\boldsymbol{\phi}}
%\newcommand{\nwd}{4}
\newcommand{\qll}[1]{[\![{#1}]\!]}
\newcommand{\rep}[1]{({#1})}
\newcommand{\too}{\to \to}
%\newcommand{\conc}{\mathfrak g}
%\newcommand{\soe}{d_\eta}
%\newcommand{\etas}{\lambda^{-3}W^{-7/2} L^{5-d}}
% \newcommand{\etas}{\lambda^{-2}W^{-4} L^{5-d}}
% \newcommand{\etasa}{\beta(\lambda) W^{-5} L^{5-d}}
% \newcommand{\blam}{\beta(\lambda)}
%\newcommand{\etars}{\lambda^{-5} W^{-5/2+2\mathfrak d_2}L^{-(d-5)}}

\newcommand{\ilambda}{g}
\DeclareMathOperator{\diag}{diag}
\DeclareMathOperator{\tr}{Tr}
\DeclareMathOperator{\var}{Var}


\newcommand{\cor}{\color{red}}
\newcommand{\nc}{\normalcolor}

\newcommand{\cal}{\mathcal}
\newcommand{\sT}{{\mathsf T}}
%\newcommand{\ba}{{\bf{a}}}
\newcommand{\thn}{\vartheta}
\newcommand{\contract}{\text{loop-contraction}}
\newcommand{\Contract}{\text{Loop-contraction}}

\theoremstyle{plain} %plain, definition, remark
\newtheorem{theorem}{Theorem}[section]
\newtheorem*{theorem*}{Theorem}
\newtheorem{lemma}[theorem]{Lemma}
\newtheorem{assumption}[theorem]{Assumption}
\newtheorem*{lemma*}{Lemma}
\newtheorem{corollary}[theorem]{Corollary}
\newtheorem*{corollary*}{Corollary}
\newtheorem{proposition}[theorem]{Proposition}
\newtheorem*{proposition*}{Proposition}
\newtheorem{claim}[theorem]{Claim}
\newtheorem*{claim*}{Claim}
\newtheorem{notation}[theorem]{Notation}
\newtheorem{definition}[theorem]{Definition}  
\newtheorem*{definition*}{Definition}         
\theoremstyle{remark}
\newtheorem{example}[theorem]{Example}
\newtheorem*{example*}{Example}
\newtheorem{remark}[theorem]{Remark}
\newtheorem{remarks}[theorem]{Remarks}
\newtheorem*{remark*}{Remark}
\newtheorem*{remarks*}{Remarks}
\newtheorem{strategy}[theorem]{Strategy}
%\newtheorem*{notation}{Notation}
 \numberwithin{equation}{section}

\crefname{theorem}{theorem}{theorems}
\crefname{lemma}{lemma}{lemmas}
\crefname{definition}{definition}{definitions}
\crefname{proposition}{proposition}{propositions}
\crefname{corollary}{corollary}{corollaries}
\crefname{claim}{claim}{claim}
\crefname{example}{example}{Example}
\crefname{remark}{remark}{remark}


%letter labeled theorems
\newtheorem{thmx}{Theorem}
\renewcommand{\thethmx}{\Alph{thmx}}


\makeatletter
\let\old@setauthors\@setauthors
\def\@setauthors{%
  \begingroup
    \let\MakeUppercase\relax
    \old@setauthors
  \endgroup
}
\makeatother


\title{Delocalization of Non-Mean-Field Random Matrices \\ in Dimensions $d\ge 3$}

\author{Sofiia Dubova$^\star$} 
\author{Fan Yang$^\dagger$}
\author{Horng-Tzer Yau$^\ddagger$}
\author{Jun Yin$^\S$}

\thanks{$^\star$Department of Mathematics, Northwestern University, \href{mailto:sdubova@northwestern.edu}{sdubova@northwestern.edu}}
\thanks{$^\dagger$Yau Mathematical Sciences Center, Tsinghua University, \href{mailto:fyangmath@mail.tsinghua.edu.cn}{fyangmath@mail.tsinghua.edu.cn}}
\thanks{$^\ddagger$Department of Mathematics, Harvard University, \href{mailto:htyau@math.harvard.edu}{htyau@math.harvard.edu}}
\thanks{$^\S$Department of Mathematics, University of California, Los Angeles, \href{mailto:jyin@math.ucla.edu}{jyin@math.ucla.edu}}

% \author{Sofiia Dubova\thanks{Department of Mathematics, Northwestern University, \href{mailto:sdubova@northwestern.edu}{sdubova@northwestern.edu}} \and
% Fan Yang\thanks{Yau Mathematical Sciences Center, Tsinghua University, \href{mailto:fyangmath@mail.tsinghua.edu.cn}{fyangmath@mail.tsinghua.edu.cn}} \and
%     Horng-Tzer Yau\thanks{Department of Mathematics, Harvard University, \href{mailto:htyau@math.harvard.edu}{htyau@math.harvard.edu}} \and
%   Jun Yin\thanks{Department of Mathematics, University of California, Los Angeles, \href{mailto:jyin@math.ucla.edu}{jyin@math.ucla.edu}}
% }
%\date{}  % Or specify a custom date, e.g., \date{October 14, 2024}

\begin{document}
%\addtokomafont{author}{\raggedright}
 



%\vspace{-10pt}
\begin{abstract}

\small 
We study $N \times N$ random band matrices $H = (H_{xy})$ with independent mean-zero Gaussian entries, where $x,y$ lie on the discrete torus \smash{$(\mathbb{Z} / \sqrt[d]{N} \mathbb{Z})^d$} in dimensions $d \ge 3$. The variance profile satisfies $\mathbb{E}|H_{xy}|^2 = S_{xy}$, with $S_{xy} = 0$ whenever the distance between $x$ and $y$ exceeds a bandwidth parameter $W$. We prove that if $W \geq N^{\mathfrak{c}}$ for some constant $\mathfrak{c} > 0$, then in the large-$N$ limit, bulk eigenvectors are delocalized, quantum unique ergodicity (QUE) holds, and the local bulk eigenvalue statistics are universal. Our proof is based on the tree approximation of the loop hierarchy \cite{YY_25} and diagrammatic techniques developed in earlier works \cite{yang2021random, yang2021delocalization, yang2022delocalization, DY, DYYY25}.


The random band matrix model above is the classical Wegner orbital model; alongside it we also study a second classical non-mean-field model, the block Anderson model. Specifically, we consider Hermitian matrices $H = V + \ilambda \Psi$ on the same discrete torus \smash{$(\mathbb{Z} / \sqrt[d]{N} \mathbb{Z})^d$},  
where $V$ is a random block potential consisting of i.i.d.~complex Gaussian diagonal blocks of size $W^d \times W^d$, and $\Psi$ encodes the interactions between neighboring blocks---random in the Wegner orbital model and deterministic in the block Anderson model. 
The parameter $\ilambda > 0$ represents the coupling strength between blocks. Assuming again that $W \geq N^{\mathfrak{c}}$, we establish delocalization of bulk eigenvectors, QUE, and bulk universality under the condition \smash{$W^{-d/2+\varepsilon}\le \ilambda \le \varepsilon^{-1}$} for any small constant $\varepsilon>0$. 
Combined with the localization results of~\cite{PelSchShaSod} for \smash{$\ilambda \ll W^{-d/2}$}, this identifies a localization--delocalization transition at the scale $\ilambda=W^{-d/2}$ in dimensions $d \ge 3$.
\end{abstract}

\maketitle
 
%\newpage
{
\hypersetup{linkcolor=black}
\tableofcontents
}
%\newpage

 

\section{Introduction}


The tight-binding model introduced by Anderson \cite{Anderson} describes electron transport in disordered semiconductors. It is formulated as a discrete random Schr{\"o}dinger operator on $\mathbb{Z}^{d}$ of the form
\begin{align}\label{Anderson_orig}
H = - \Delta + \lambda V,
\end{align}
where $\Delta$ denotes the graph Laplacian on $\mathbb{Z}^{d}$, $V$ is a random potential with i.i.d.~entries, and $\lambda>0$ is a coupling parameter representing the strength of disorder. In his seminal work \cite{Anderson}, Anderson predicted a transition from delocalized to localized behavior as the disorder strength $\lambda$ increases.

Localization in the one-dimensional (1D) Anderson model, valid for all $\lambda>0$, is well established; see, for example, \cite{GMP, KunzSou, Carmona1982_Duke, Damanik2002}. In dimensions $d \geq 2$, Anderson localization was first rigorously proved by Fr{\"o}hlich and Spencer \cite{FroSpen_1983} using multi-scale analysis. An alternative approach was later developed by Aizenman and Molchanov \cite{Aizenman1993} based on the fractional moment method. Since then, a substantial body of work has deepened our understanding of Anderson localization; see, for example, \cite{FroSpen_1985, Carmona1987, SimonWolff, Aizenman1994, ASFH2001, Bourgain2005, Germinet2013, DingSmart2020, LiZhang2019}. 
In contrast, the existence of delocalized states in the finite-dimensional Anderson model remains unproven---both for arbitrary disorder strengths and in all finite dimensions. In infinite dimensions, however, Aizenman and Warzel \cite{aizenman2013resonant, Bethe_PRL} rigorously established the presence of a delocalized phase for the Anderson model on the Bethe lattice, an infinite regular tree. More recently, the mobility edge phenomenon has also been proved in this setting \cite{Bethe-Anderson}. The tree geometry differs fundamentally from Euclidean lattices, as tree graphs are loop-free.

To bridge the understanding between random Schr{\"o}dinger operators and random matrix theory, the random band matrix model was introduced \cite{PhysRevLett.64.1851, PhysRevLett.64.5, PhysRevLett.66.986}. A $d$-dimensional random band matrix $H = (H_{xy})$ is an $N\times N$ Hermitian random matrix defined on the discrete torus \smash{$(\mathbb{Z}/\sqrt[d]{N}\mathbb{Z})^{d}$} (assuming \smash{$\sqrt[d]{N}\in \mathbb{N}$} for simplicity). Subject to Hermitian symmetry, the entries $H_{xy}$ are independent real or complex Gaussian variables with mean zero. The variance profile $S_{xy} = \E|H_{xy}|^{2}$ decays to zero when the distance between $x$ and $y$ exceeds a bandwidth parameter $W$, and satisfies the normalization $\sum_y S_{xy} \equiv 1$. 
It was conjectured in \cite{PhysRevLett.64.1851, PhysRevLett.64.5, PhysRevLett.66.986,PhysRevLett.67.2405} that this model exhibits a localization--delocalization transition, accompanied by a change in spectral statistics from Poisson to GOE/GUE (Gaussian Orthogonal/Unitary Ensemble) statistics, at a critical bandwidth $W_c(N)$. For eigenvalues in the bulk of the spectrum, i.e., $|E| \le 2 - \kappa$ with some fixed $\kappa > 0$, localization and Poisson statistics are expected when $W \ll W_c$, while delocalization and GOE/GUE statistics are expected when $W \gg W_c$, where
\begin{equation}\label{eqn:transition exponents}
W_c =
\begin{cases}
\sqrt N, & d=1, \\
\sqrt{ \log N}, & d=2, \\
\OO(1), & d\ge 3.
\end{cases}
\end{equation}
See \cite{PB_review, Spencer2, Spencer3, Spencer1} for further discussion of these conjectures.

The density of states for random band matrices in dimensions $d \ge 1$ was established in \cite{DisPinSpe2002}.
Delocalization of 1D band matrices under the assumption $W \geq N^{a}$, for various exponents $a > 1/2$, was proved in a series of works \cite{EK_band1, ErdKno2011, erdHos2013local, BaoErd2015, HeMa2018, bourgade2017universality, bourgade2019random, bourgade2020random, yang2021random, DY}. The full conjecture was recently resolved in \cite{YY_25}, which established both delocalization and universality of bulk eigenvalue statistics under the optimal condition $W \gg \sqrt N$. These results were later extended to non-Gaussian random band matrices in \cite{erdHos2025zigzag}. 
On the other hand, localization has been proved under the condition $W \ll N^{a}$ for various exponents $a < 1/2$ in \cite{Sch2009, PelSchShaSod,Cipolloni2024, Chen2022}, and under the sharp condition $W \ll \sqrt N$ in \cite{Localization1_2}. Therefore, the localization-delocalization transition at the critical bandwidth $W = \sqrt N$ is now rigorously established for 1D random band matrices.
In two dimensions, delocalization and bulk universality under the condition $W \ge N^{\varepsilon}$ (for arbitrarily small constant $\varepsilon > 0$) was proved in \cite{DYYY25}.


In this paper, we study band matrices $H$ defined on the $d$-dimensional torus $\mathbb{Z}_{WL}^d=\mathbb{Z}^d / ((WL) \cdot \mathbb{Z}^d)$ with $N=(WL)^d$ and $d \geq 3$. For simplicity, we assume that the torus $\mathbb{Z}_{WL}^d$ is partitioned into $d$-dimensional cubes of side length $W$, indexed by the smaller torus ${\mathbb Z}_{L}^d$. For $x, y\in {\mathbb Z}_{WL}^d$ belonging to blocks labeled by $a, b \in {\mathbb Z}_L^d$, the variance profile is given by
\begin{equation}\label{eq:variance-profile}
S_{xy} = c_{a-b} W^{-d}  \mathbf{1}(\|a-b\|_1\le 1),
\end{equation}
where $(c_{a-b}\mathbf{1}(\|a-b\|_1\le 1):a, b \in {\mathbb Z}_L^d)$ is a symmetric doubly stochastic matrix. Our main results establish delocalization of bulk eigenvectors and universality of bulk eigenvalue statistics for the band matrix $H$ in dimensions $d \ge 3$, under the assumption $W \ge N^{\varepsilon}$ for any fixed $\varepsilon>0$. In higher dimensions $d \ge 7$, weaker forms of delocalization and bulk universality were previously obtained in~\cite{yang2021delocalization, yang2022delocalization, Xu:2024aa}.

Beyond standard band matrices, the block Anderson model provides an even closer analogue to the original Anderson model~\eqref{Anderson_orig}. This model, inspired by Wegner's orbital model~\cite{Wegner1, Wegner2, Wegner3}, describes quantum particles with multiple internal degrees of freedom (e.g., orbitals or spin) moving in a disordered medium. Concretely, we consider the matrix $H = \Psi + \lambda V$, where $V=\diag(V_1,\ldots, V_{L^d})$ and the blocks $V_i$ are i.i.d.~$W^d\times W^d$ Gaussian random matrices. The matrix $\Psi$ encodes the hopping between neighboring blocks, and in the block Anderson model, we take $\Psi$ to be the block Laplacian with the diagonal removed.

It was shown in~\cite{PelSchShaSod} that, in all dimensions $d \ge 1$, the block Anderson model exhibits localization with localization length of order $\OO(W)$ whenever $\lambda \gg W^{d/2}$. In this paper, we prove that all results obtained for random band matrices extend to the block Anderson model, provided $c \le \lambda \le W^{d/2-c}$ for a small constant $c > 0$. Combined with the localization results of~\cite{PelSchShaSod}, this establishes the localization--delocalization transition for the block Anderson model in all dimensions $d \ge 3$. In dimensions $d\in\{1,2\}$, delocalization and bulk universality were previously proved in~\cite{RBSO1D}, while partial delocalization results for dimensions $d \ge 7$ were obtained in~\cite{yang2024Del}. 
However, for the physically most relevant case $d=3$, no rigorous results on delocalization were available for either random band matrices or the block Anderson model prior to this work. One reason why $d=3$ is the final dimension in which these conjectures are resolved lies in certain critical difficulties that are intrinsic to three dimensions, as we now explain.

 
%The resolvent (or Green's function) of random walks on the $\mathbb{Z}_{WL}^d$, which governs the behavior of the resolvents of random band and block Anderson models, decays like $1/|x|$ in $d=3$. In dimensions $d\in\{1,2\}$, this resolvent is essentially constant (up to a logarithmic factor in $d=2$) within a cutoff distance. Starting from $d=3$, however, it exhibits polynomial decay like $|x-y|^{-d+2}$. Therefore, we need to track the precise \emph{pointwise behavior} of the resolvent of random matrices, i.e., pointwise bounds on $(H-z)^{-1}_{xy}$. 
%To explain the fundamental difficulty of our problem in dimension $d=3$, 
Denote by $G(z)=(H-z)^{-1}$ the resolvent (Green's function) of the random band and block Anderson models, for $z\in \C$. 
It is believed that the size of the resolvent entries $|G_{0x}(z)|^2$ is approximated by the resolvent of a certain random walk on \smash{$\mathbb{Z}_{WL}^d$}, see e.g., \cite{erdos2013delocalization,yang2021delocalization}. 
In dimensions $d\in\{1,2\}$, the random-walk resolvent remains roughly constant (in magnitude, up to a logarithmic correction when $d=2$) below a certain cutoff scale. Hence, an $L^\infty$-bound on $G(z)$, together with control of this cutoff scale, suffices to estimate $G(z)$. 
Starting from $d=3$, however, $|G_{0x}|^2  \asymp |x|^{-d+2}$. This forces us to track the  \emph{pointwise estimate} of $G_{0x}$. 
At the other extreme, in high dimensions, the lattice $\mathbb Z^d$ increasingly resembles a Bethe lattice. Although this analogy is difficult to exploit rigorously, we note that $(|x|^{-d+2})^\alpha$ is integrable whenever $\alpha > d/(d-2) \to 1$ as $d\to \infty$. In other words, $|x|^{-d+2}$ becomes ``almost integrable'' for large $d$. This observation was crucially used in the proofs for $d\ge 7$ in \cite{yang2021delocalization, yang2022delocalization, Xu:2024aa}.  
In dimension $d=3$, we need to combine the \emph{tree-approximation} method for $d\in\{1,2\}$ \cite{YY_25, DYYY25} with the \emph{diagrammatic expansion} approach for $d\ge 7$ \cite{yang2021delocalization, yang2022delocalization, Xu:2024aa}. 


% We first note that  
%  the $\Theta$-propagator \eqref{prop:ThfadC}, which governs the  off-diagonal resolvent entries of band matrices, are bounded by 
% \begin{align}\label{just}
% \E|(G_t)_{xy}|^2 \approx \cL^{(2)}_{t,(+,-),(a,b)} \approx W^{-d}\Theta^{(+,-)}_{t,ab}\lesssim \frac{e^{-c_d|a-b|/\ell_t } }{W^d} \cdot \begin{cases}
%  B_{t,|a-b|} ,\quad &d\ge 3
% \\
%    (\ell_t^d\eta_t)^{-1},\quad &d=1,2  
% \end{cases}, \ \ \text{for}\ \ x\in [a], \  y\in [b]. 
% \end{align}
% Here we neglected the logarithmic corrections in $d=2$. As explained in the introduction, we need to estimate $|(G_t)_{xy}|^2$ pointwisely. 

\medskip

To illustrate the above points, we briefly explain the main ideas developed in this paper.
For simplicity, we focus on the random band matrix model with variance profile \eqref{eq:variance-profile}; the analysis for the block Anderson model is analogous, up to some additional technical details. 
We begin by recalling the matrix Brownian motion used in \cite{DY, YY_25, DYYY25}:
\[
\dd (H_{t})_{xy} = \sqrt{S_{xy}} \, \dd (\boldsymbol{B}_{t})_{xy}, \quad \text{with}\quad H_0 = 0,
\]
where \( (\boldsymbol{B}_{t})_{xy} \) are standard independent complex Brownian motions for \( x, y \in  \Z_{WL}^d\), subject to the Hermitian symmetry condition \( (\boldsymbol{B}_{t})_{xy} = \bar{(\boldsymbol{B}_{t})_{yx}} \).  Following \cite{10.1214/19-ECP278, Sooster2019, DY}, we consider the Green's function of \( H_t \) with a carefully chosen time-dependent spectral parameter \( z_t \) (see \eqref{eq:zt} below), defined by
\[
G_t := (H_t - z_t)^{-1},
\]
whose dynamics are naturally renormalized at leading order. 
We then define the $n$-\( G \)-loop observable $\mathcal{L}^{(n)}$ as an $n$-tensor:
\be\label{eq:nGloops}
{\cal L}^{(n)}_{a_1,\ldots, a_n} \equiv {\cal L}^{(n)}_{a_1,\ldots, a_n} (t):= \mathrm{Tr} \prod_{i=1}^n \p{G_t \cdot E_{a_i}}, \quad \text{for}\quad a_i\in \Z_L^d, \ \  1\le i\le n,
\ee
where each $E_a$ is a block-averaging matrix, defined by $(E_a)_{xy} = W^{-d}\delta_{xy}$ when $x, y$ belong to the $a$-th block, and $(E_a)_{xy}=0$ otherwise. In the proof, the $G$-loops considered in this paper involve combinations of $G_t$ and its Hermitian conjugate $G^\ast_t$; however, we omit this detail in the following heuristic discussion for simplicity.


The \( G \)-loops ${\cal L}^{(n)}(t)$ defined above satisfy a system of evolution equations known as the \emph{loop hierarchy}; see equation \eqref{eq:mainStoflow} below for its precise formulation. The dynamics of \( n \)-loops depend on \((n+1)\)-loops and a martingale term, whose quadratic variation involves \( (2n+2) \)-loops. Schematically, the loop hierarchy takes the form
\be\label{eq:introloop_hier}
\dd {\cal L}^{(n)} = ( {\cal L}  \ast {\cal L})^{(n)}\dd t + {\cal E}^{\Gc,(n)}\dd t + \dd {\cal E}^{M,(n)}, 
\ee
where each term on the right-hand side has the following interpretation.
\begin{itemize}
	\item The \emph{light-weight} term \( {\cal E}^{\Gc,(n)} \) depends on \( (n+1) \)-loops, typically of the form
	\be\label{eq:introLWtermS}
	W^d \sum_{u,v\in \Zn} S^{(\sB)}_{uv} \cdot \tr[\Gc_t E_u]\cdot {\cal L}^{(n+1)}_{a_1, \ldots, a_n, v}, \quad \text{with}\quad \Gc_t:= G_t-mI_N.
	\ee
	Here, $\tr[\Gc_t E_u]$ is referred to as a \emph{light-weight}, \smash{$S^{(\sB)}_{uv} = c_{u-v} \mathbf{1}(\|u-v\|_1\le 1)$} denotes the block variance profile corresponding to \eqref{eq:variance-profile}, and $m \equiv m_{\rm sc}$ is the Stieltjes transform of the semicircle law (see \eqref{eq:defmzsc}).
	
	\item The \emph{martingale} term \( \dd {\cal E}^{M,(n)} \) has a quadratic variation depending on  \( (2n+2) \)-loops.
	
	\item The \emph{convolution} term \( ({\cal L} \ast {\cal L})^{(n)} \) consists of sums of the form \( \mathcal{L}^{(p)} \diamond \mathcal{L}^{(q)} \) with \( p+q = n+2\) and \(p,q \ge 2 \), where
	\be\label{eq:diamond}
	({\cal L}^{(p)} \diamond {\cal L}^{(q)})_{a_1, \ldots, a_n} := W^d \sum_{u,v} S^{(\sB)}_{uv} \cdot {\cal L}^{(p)}_{a_1, \ldots, a_{p-1}, u} \cdot {\cal L}^{(q)}_{a_p, \ldots, a_n, v}.
	\ee
\end{itemize}
Although the formulation above is heuristic, it captures the essential structure of the loop hierarchy. 


The loop equation \eqref{eq:introloop_hier} can be generalized to random band matrices without block structure. This issue has already been addressed in earlier works; see \cite{erdHos2025zigzag} for $d=1$ and \cite{Block_reduction} for dimensions $d\in\{1,2\}$. In fact, our method can be automatically extended to the case where the variance matrix takes the following form:
\begin{align*} 
 S_{xy} = \sum_{u,v \in \ZL} \chi(x - u) \wt S_{uv} \chi(v - y),
\end{align*} 	
where $\wt S$ is a symmetric doubly stochastic matrix and \smash{$\chi : \ZL \to [0,\infty)$} is a mollifier at scale $W^{1 - \e}$. In this case, let $F_u$ denote the diagonal mollifier matrices with entries $(F_{u})_{xy}= \delta_{xy} \chi(x - u)$  for $u, x, y \in \ZL.$ Then, for \smash{$(x_1, \ldots, x_n) \in (\ZL)^n$}, we define the $n$-$G$-loops \smash{$\cL^{(n)}_{x_1,\ldots, x_n}(t)$} in analogy with \eqref{eq:nGloops}, replacing the matrices $E_{a_i}$ with the mollifier matrices $F_{x_i}$. The dynamics of these loops satisfy equations of the same form as \eqref{eq:introloop_hier}, with \smash{$S^{(\sB)}$} in \eqref{eq:introLWtermS} and \eqref{eq:diamond} replaced by \smash{$\wt S$}. Consequently, all  our arguments extend to this setting with only minor notational modifications.
To streamline the presentation, we continue to work with the simpler block structure in \eqref{eq:variance-profile}, which allows us to focus on the core technical challenges in dimensions $d\ge 3$. 




A key observation from \cite{YY_25, DYYY25} is that both the light-weight and martingale terms are small errors. Consequently, \( \mathcal{L}^{(n)} \) can be approximated by \( \mathcal{K}^{(n)} \), the solution of the deterministic \emph {convolution tree  equation}: 
\be\label{eq:Knequation}
\dd \mathcal{K}^{(n)} = (\mathcal{K} \ast \mathcal{K})^{(n)}\dd t.
\ee
For any fixed \( n \), the right-hand side of equation \eqref{eq:Knequation} depends only on \( \mathcal{K}^{(k)} \) with \( k \leq n \), forming a \emph{closed system} of equations for \( (\mathcal{K}^{(2)}, \ldots, \mathcal{K}^{(n)}) \) that can be solved inductively and admits explicit solutions, as discovered in \cite{YY_25}. Since these solutions have an explicit tree representation, we refer to \( \mathcal{K}^{(n)} \) as the \emph {tree approximation} to \( \mathcal{L}^{(n)} \). 
In particular, for \( n = 2 \), equation \eqref{eq:Knequation} has the solution \( \mathcal{K}^{(2)}(t) = W^{-d} \Theta_t\), where $\Theta_t$ denotes the propagator defined in \Cref{def_Theta}. 
This propagator can carry different charges. For simplicity, in the discussion below, we use the notation $\Theta_t$ exclusively for the most relevant $(+,-)$-charged propagator, denoted \smash{$\Theta_t^{(+,-)}$}. It can be viewed as the resolvent of a Laplacian on the torus $\Zn$,  describing the classical diffusion of a random walk. 
%This propagator can have different charges. For simplicity, in the following discussion, we will always use $\Theta_t$ to denote the more crucial $(+,-)$-charged propagator \smash{$\Theta_t^{(+,-)}$}, which can be regarded as the resolvent of a Laplacian on the torus $\ZL$, describing classical diffusion of a random walk.
This propagator serves as the basis for establishing quantum diffusion, delocalization, and the universality of eigenvalue statistics---provided the approximation \( \mathcal{L}^{(2)} \approx \mathcal{K}^{(2)} \) can be justified in a sufficiently strong sense.

To rigorously justify the tree approximation, three main ingredients are required:
\begin{enumerate}
	\item[(1)] A strategy to manage the dependence of \( {\cal E}^{\Gc,(n)} \) and \( \dd {\cal E}^{M,(n)} \) on higher-order loops \( {\cal L}^{(k)} \) with \( k > n \).
	\item[(2)]  A stability theory for the perturbed equation \( \dd {\cal L}= ({\cal L} \ast {\cal L})\dd t + \text{errors} \).
	\item[(3)] Explicit quantitative bounds showing that  \( {\cal E}^{\Gc,(n)} \) and \( \dd {\cal E}^{M,(n)} \) are indeed small.
\end{enumerate} 
In the absence of additional structure in the hierarchy, issue (1) reflects a classical obstacle in many-body dynamics, reminiscent of the well-known BBGKY hierarchy. 
In our setting, however, a key observation is that higher-order loops appear only in the error terms. This feature enables the development of a general inductive and self-improving framework, first introduced in \cite{YY_25}, based on a bootstrap argument that applies across all dimensions. It should be noted that the smallness of the martingale and light-weight terms is a highly nontrivial fact. 
Without the correct perspective, these terms may seem much larger than the leading contribution.
Indeed, in the early development of the theory of random band matrices, researchers often encountered the seemingly impenetrable issue that the error terms in the so-called \emph{$T$-equation} diverge. While we do not elaborate on the historical developments here, we will demonstrate below that the light-weight term is, in fact, a genuinely small error.


For both issues (2) and (3), we need a suitable norm to control the error terms and develop a stability theory for the loop hierarchy.  
In dimensions \( d \in \{1,2\} \), the natural choice is the \( L^\infty \)-norm, justified by the fact that \( \Theta_t \) is approximately constant up to a cutoff.  
Surprisingly, the loop hierarchy remains stable in dimensions \( d \geq 3 \) under this norm. In dimensions $d\ge 3$, the propagator satisfies
\[
\Theta_{t,ab} \lesssim (|a-b|+1)^{-d+2} e^{- c|a-b|/\ell_t }
\]
for some constant \( c>0 \), where \( \ell_t  = \min(|1-t|^{-1/2},L)\). The tree approximation \( \mathcal{K}^{(n)}_{a_1, \ldots, a_n} \) is a complicated function of the propagator entries. Hence, the \( L^\infty \)-norm is a very crude norm for these functions, and one would not a priori expect stability in this norm.   
For clarity, in the following heuristic discussion, we assume \( \im z_t \ge L^{-2} \), so that \( \ell_t^{-2} = |1-t| \asymp \im z_t \). The key new ingredient enabling \( L^\infty \)-stability of the loop hierarchy in dimensions \( d \ge 3 \) is a class of \emph{$\contract$ inequalities} (\Cref{ygdhmsgq}), which extend the classical Ward's identity:
\be\label{eq:Ward_intro1}
\sum_y |G_{xy}(z)|^2=\sum_y |G_{yx}(z)|^2 =  \im G_{xx}/ \im z ,
\ee
and apply to partial sums involving the absolute values of the $n$-\(G\)-loops. 


%$ = W^{-d} \Theta_{t,ab}$ in the sense 
%However, in order to establish quantum diffusion and quantum unique ergodicity, and for the {\cor stability of the tree approximation}, we need a \emph{pointwise estimate} valid for all \( a, b \), i.e.,

However, the \( L^\infty \)-stability of the tree approximation remains incomplete without a sufficiently precise \emph{pointwise estimate} \smash{${\cal L}^{(2)}_{ab} \approx \mathcal{K}^{(2)}_{ab}$} for the 2-$G$-loops, that is,  
\be\label{eq:n=2approx}
\left|\left({\cal L}-\cK\right)^{(2)}_{ab} \right| \ll W^{-d} \Theta_{t,ab}, \quad  \forall a,b\in \Zn.
\ee
Such a strong pointwise control is essential because of the leading term $\cal L^{(2)}\diamond \cal L^{(n)}$ in the loop hierarchy \eqref{eq:introloop_hier}. Subtracting the corresponding term $\cal K^{(2)}\diamond \cal K^{(n)}$ from the convolution tree equation \eqref{eq:Knequation} produces three terms $\cal K^{(2)}\diamond (\cal L-\cal K)^{(n)}$, $(\cal L-\cal K)^{(2)}\diamond \cal K^{(n)}$, and $(\cal L-\cal K)^{(2)}\diamond (\cal L-\cal K)^{(n)}$. While it is intuitively clear that the last two terms are negligible compared to the first, a rigorous justification requires establishing the stability estimate \eqref{eq:n=2approx}. 
This requires a significantly more delicate analysis of the equation \eqref{eq:introloop_hier} for $n=2$.  
One main difficulty is that this equation involves \( \mathcal{L}^{(n)} \) with \( n \geq 3 \), for which only \( L^\infty \)-bounds are available. Therefore, we need to develop a method to estimate the error terms in the 2-loop equation using only the \( L^\infty \)-bounds on the higher-order loops. 
We will discuss the details of the $L^\infty$-stability and $\contract$ inequalities in \Cref{Sec:Steps34}, and the proof of the pointwise stability estimate \eqref{eq:n=2approx} in \Cref{subsec:step2_pf}.



Issue (3) involves bounding the martingale term \( \dd {\cal E}^{M,(n)} \) and the light-weight term \( {\cal E}^{\Gc,(n)} \). In dimensions \( d \ge 3 \), estimating these terms presents a serious challenge. We will discuss the difficulties related to bounding the martingale term in \Cref{subsec:pf_lem: EMn2_N}. For now, we focus on estimating the light-weight term for 
 \( n=2 \): 
\[
{\cal E}^{\Gc,(2)}_{ab} = W^d \sum_{u, v} S^{(\sB)}_{uv} \cdot \tr[\Gc_t E_{u}] \cdot \mathcal{L}^{(3)}_{a,b,v}, \quad\text{with}\quad 
{\cal L}^{(3)}_{a,b,v} = W^{-3d} \sum_{x \in [a], y \in [b], z \in [v]} (G_t)_{xy} (G_t)_{yz} (G_t)_{zx},
\]
where \( x \in [a] \) means that the lattice point \( x \) belongs to the block labeled by \( a \).  
We expect an averaged local law to hold in the form \smash{\( |\tr[\Gc_t E_{u}]| \leq W^{-c} \)} for some constant \( c > 0 \).  
Next, assuming the \emph{sharp entrywise local law} for $G_t$: 
\be\label{eq:aprioriG2} 
|(G_t)_{xy}- m \delta_{xy}| \lesssim W^{-d/2}
\big({\Theta_{t,ab}}\big)^{1/2}, \quad \forall x\in [a], \ y \in[b],
\ee 
we obtain the estimate  
\begin{align} 
	|{\cal E}^{\Gc,(2)}_{ab}|  %& \lesssim W^{-c + d} \sum_{|z-x|\vee|z-y| \lesssim W\ell_t} |(G_t)_{xy} (G_t)_{xz} (G_t)_{zy}| \\
    &\lesssim W^{-c-d/2} \sum_{|v-a|\vee |v-b| \lesssim\ell_t} \p{\Theta_{t,ab} \Theta_{t,bv} \Theta_{t,va}}^{1/2} \nonumber  \\
	& \lesssim W^{-c-d/2} \big(\Theta_{t,ab}\big)^{1/2} \ell_t^2 = W^{-c-d/2} \big(\Theta_{t,ab}\big)^{1/2}\cdot (1 - t)^{-1}. 
    \label{eq:introGc}
\end{align}
Upon integrating in time, the singularity \( (1 - t)^{-1} \) produces only a harmless logarithmic factor. However, the resulting bound contains only \smash{\( (\Theta_{t,ab})^{1/2} \)} rather than the optimal \smash{\( \Theta_{t,ab} \)}. This gap cannot be compensated by the \( W^{-c} \) factor from the averaged local law, as \( c \) is at best \( c=d \).  
Without the extra factor \smash{\( (\Theta_{t,ab})^{1/2} \)}, the light-weight term would dominate the leading term and violate the approximation in \eqref{eq:n=2approx}. In the actual proof, the situation is even more delicate, since the estimate \eqref{eq:aprioriG2} is not available a priori.

To overcome this difficulty, we employ the \emph{diagrammatic methods} developed in \cite{yang2021delocalization, yang2022delocalization, yang2021random}.  
While this method is technically involved and was previously applied only in high dimensions \( d \geq 7 \), our objective here is more modest: we aim to generate an additional ``\emph{long edge}''---that is, a resolvent entry \( G_{x'y'} \) with \( |x'-y'| \gtrsim W|a-b| \)---in order to recover the missing factor \smash{\( (\Theta_{t,ab})^{1/2} \)}.  
For this purpose, the diagrammatic expansions can be extended to all dimensions \( d \geq 3 \). 

Specifically, we apply Gaussian integration by parts to express the expectation of high moments \smash{\( \E |{\cal E}_{ab}^{\Gc,(2)}|^{2p}\)} as a sum of graphs, where each graph either contains \emph{\( 2p \) additional long edges} or contributes to a sufficiently small error.   
Roughly speaking, each application of Gaussian integration by parts produces a \( W^{-d/2} \) factor, unless it generates a long edge.  
To ensure that these \( W^{-d/2} \) factors compensate for the missing \smash{\( (\Theta_{t,ab})^{1/2} \)} when \( |a-b| \asymp L \), we must expand the graphs to a sufficiently high order $k$, chosen so that \( W^{-kd/2} \le (\Theta_{t,ab})^{p} \). 
As a result, the contributing graphs are extremely complicated. To bound them, we need to sum over their internal vertices in such a way that the final contribution retains at least the \( 2p \) additional long edges. 
We will use Ward's identity \eqref{eq:Ward_intro1} as a key tool to control the summation over internal vertices, provided the order of the summations follows a carefully designed ``nested order''. This nested ordering is a crucial component in our estimates of the graphs arising from the complicated diagrammatic expansions. See \Cref{sec:graphs_ideas} for a detailed discussion. 

In addition to the above considerations, two key ingredients continue to play a central role in our analysis: the \emph{sum-zero property} \cite{yang2021delocalization, yang2022delocalization,YY_25}, and the CLT mechanism for the fluctuation cancellation of resolvents, which was previously employed in the \( d = 2 \) analysis \cite{DYYY25}. Our final objective is to incorporate all these components---diagrammatic tools, the sum-zero property, fluctuation cancellation, and the stability of the tree approximation---into the flow-based framework.  

In summary, our proof relies on a flow-based analysis of the tree approximation to the $G$-loop hierarchy, combining \emph{sharp max-norm estimates} for higher-order $(\cL-\cK)^{(n)}$-loops with $n\ge 3$ and a \emph{precise pointwise estimate} for the $(\cL-\cK)^{(2)}$-loops.
It is perhaps surprising that this seemingly ``simple'' strategy succeeds without tracking the precise pointwise decay of higher-order $G$-loops. This is made possible by several key technical inputs, including \emph{$\contract$ inequalities} and \emph{diagrammatic techniques}.
The main advantage of our approach is that it closes the analysis using only a minimal set of technical ingredients, thereby avoiding an unnecessarily complicated treatment in dimensions $d\ge 3$.
%The main advantage of our approach is that it closes the analysis using only a minimal set of technical inputs, thereby avoiding unnecessary complications in dimensions $d\ge3$. 
A complete implementation of this strategy will be presented in the subsequent sections after we state the main results in \Cref{sec_model-results}.
%Several technical arguments, which are similar to those in earlier works \cite{yang2021delocalization, YY_25, DYYY25, RBSO1D, yang2024Del}, together with the proofs of some standard deterministic estimates, are deferred to the appendix without affecting the logical flow of the proof. 







\medskip
{\noindent \bf Notations.}
To facilitate the presentation, we introduce some necessary notations that will be used throughout this paper. We will use the set of natural numbers $\N=\{1,2,3,\ldots\}$ and the upper half complex plane $\C_+:=\{z\in \C:\im z>0\}$.  
In this paper, we are interested in the asymptotic regime with $N\to \infty$. When we refer to a constant, it will not depend on $N$ or $W$. Unless otherwise noted, we will use $C$, $D$ etc.~to denote large positive constants, whose values may change from line to line. Similarly, we will use $\e$, $\delta$, $\tau$, $c$, $\fc$, $\fd$ etc.~to denote small positive constants. 
For any two (possibly complex) sequences $\xi_N$ and $\zeta_N$ depending on $N$, $\xi_N = \OO(\zeta_N)$, $\zeta_N=\Omega(\xi_N)$, or $\xi_N \lesssim \zeta_N$ means that $|\xi_N| \le C|\zeta_N|$ for some constant $C>0$, whereas $\xi_N=\oo(\zeta_N)$ or $|\xi_N|\ll |\zeta_N|$ means that $|\xi_N| /|\zeta_N| \to 0$ as $N\to \infty$. We say that $\xi_N \asymp \zeta_N$ if $\xi_N = \OO(\zeta_N)$ and $\zeta_N = \OO(\xi_N)$. For any $\al,\beta\in\R$, we denote $\qqq{\al, \beta}:= [\al,\beta]\cap {\mathbb Z}$, $\qqq{\al}:=\qqq{1,\al}$, $\al\vee \beta:=\max\{\al, \beta\}$, and $\al\wedge \beta:=\min\{\al, \beta\}$. 
Given a vector $\mathbf v$, $|\mathbf v|\equiv \|\mathbf v\|_2$ denotes the Euclidean norm and $\|\mathbf v\|_p$ denotes the $L^p$-norm. 
Given a matrix $\cal A = (\cal A_{ij})$, $\|\cal A\|$, $\|\cal A\|_{p\to p}$, and $\|\cal A\|_{\infty}\equiv \|\cal A\|_{\max}:=\max_{i,j}|\cal A_{ij}|$ denote the operator (i.e., $L^2\to L^2$) norm,  $L^p\to L^p$ norm (where we allow $p=\infty$), and maximum (i.e., $L^\infty$) norm, respectively. We will use $\cal A_{ij}$ and $ \cal A(i,j)$ interchangeably in this paper. %Moreover, we introduce the following simplified notation for trace: $ \left\langle \cal A\right\rangle=\tr (\cal A).$ 
We will use $I_n$ to denote an $n\times n$ identity matrix.

Given an event $\Xi$, let $\mathbf 1_\Xi$ or $\mathbf 1(\Xi)$ denote its indicator function. We will say an event $\Xi$ holds with high probability (w.h.p.) if for any constant $D>0$, $\mathbb P(\Xi)\ge 1- N^{-D}$ for large enough $N$. More generally, we say an event $\Omega$ holds $w.h.p.$ in $\Xi$ if for any constant $D>0$, $\P( \Xi\setminus \Omega)\le N^{-D}$ for large enough $N$. 
For clarity of presentation, we will use the following notion of stochastic domination introduced in \cite{Average_fluc}. %This framework will be used throughout the analysis to control the behavior of random matrices and their spectral properties.
%\begin{definition}[Stochastic domination]\label{stoch_domination}
%	{\rm{(i)}} 
Let \[\xi=\left(\xi^{(N)}(u):N\in\mathbb N, u\in U^{(N)}\right),\hskip 10pt \zeta=\left(\zeta^{(N)}(u):N\in\mathbb N, u\in U^{(N)}\right),\]
	be two families of non-negative random variables, where $U^{(N)}$ is a possibly $N$-dependent parameter set. We say $\xi$ is stochastically dominated by $\zeta$, uniformly in $u$, if for any fixed (small) $\tau>0$ and (large) $D>0$, 
	\be\label{stoch_domination}\mathbb P\bigg(\bigcup_{u\in U^{(N)}}\left\{\xi^{(N)}(u)>N^\tau\zeta^{(N)}(u)\right\}\bigg)\le N^{-D}\ee
	for large enough $N\ge N_0(\tau, D)$, and we will use the notation $\xi\prec\zeta$. 
		If for some complex family $\xi$ we have $|\xi|\prec\zeta$, then we will also write $\xi \prec \zeta$ or $\xi=\OO_\prec(\zeta)$. 
	% \vspace{5pt}
	% \noindent {\rm{(ii)}} 
    As a convention, for two \emph{deterministic} non-negative quantities $\xi$ and $\zeta$, we will write $\xi\prec\zeta$ if and only if $\xi\le N^\tau \zeta$ for any constant $\tau>0$. 

 % \vspace{5pt}
	% \noindent {\rm{(iii)}} Let $A$ be a family of random matrices and $\zeta$ be a family of non-negative random variables. Then, we use $A=\OO_\prec(\zeta)$ to mean that $\|A\|\prec \xi$, where $\|\cdot\|$ denotes the operator norm. 
	
	%\vspace{5pt}
	%\noindent {\rm{(iii)}} 
    
	%\end{itemize}
%\end{definition}

\medskip
\noindent{\bf Acknowledgement.}   
We would like to thank Kevin Yang for fruitful discussions.


\section{The model and main results}\label{sec_model-results}


We focus in this paper on the classical block random band matrix model (i.e., the \emph{Wegner orbital model}) and the \emph{block Anderson model} with identity interactions. As noted in the introduction, these are the models studied in \cite{PelSchShaSod}, where localization of bulk eigenvectors was established under the condition $\lambda \gg W^{d/2}$. In contrast, the present work establishes a complementary result describing the delocalized regime of the Wegner orbital and block Anderson models when $\lambda \ll W^{d/2}$. 
Our approach can be readily extended to a much broader class of random band matrices and block Anderson models with general (non-identity) interactions---by combining the techniques developed here with those from \cite{RBSO1D} and the block reduction method of \cite{Block_reduction}. However, the models considered in this paper allow us to avoid inessential technical complications and to focus on the core challenges that arise in dimensions $d \ge 3$. Extensions to more general settings will be addressed in future work.
For clarity of presentation, we rescale our models by 
\be\label{def:ilambda}\ilambda=\lambda^{-1},\ee
and consider matrices of the form $H=V+\ilambda\Psi$. 
This rescaling has the advantage that, in the regime $\ilambda\ll 1$ (i.e., $\lambda\gg 1$), the density of states of $H$ may be viewed as a small perturbation of the semicircle law. 

%\subsection{The model setting}\label{subsec_setting}

Our models are defined on a $d$-dimensional discrete torus \(\ZL\subset\mathbb{Z}^d\) with $d\ge 3$, consisting of $N$ lattice points and side length $WL$:
\[\ZL:=\qqq{ -(WL)/2+1 , (WL)/2}^d\quad \text{with}\quad N=(W\cdot L)^d.\] 
We partition $\ZL$ into $L^d$ disjoint blocks, indexed by \(\Zn:=\qqq{ -L/2+1 , L/2}^d,\)
which we refer to as the \emph{block lattice}.  Each index $a=(a(1),\ldots, a(d))\in \Zn$ corresponds to a block 
\be\label{eq:blockIa}
[a] := \prod_{i=1}^d\left\llbracket (a(i)-1)W + 1, \; a(i)W \right\rrbracket, 
\ee
where $a(i)$ denotes the $i$-th coordinate of $a$ for $i\in\qqq{d}$.\footnote{In the above definitions, we implicitly assume that $L$ is even; the case of odd $L$ can be treated similarly by defining the block lattice as $ \Zn:=\qqq{ -(L-1)/2 , (L-1)/2}^d$.}
%where $a=(a(1),\ldots, a(d))$ belongs to the \emph{lattices of blocks} $\Zn$, indexed by \[ \Zn:=\qll{ -L/2+1 , L/2}^d\] with $a(i)$ denoting its $i$-th coordinate for $i\in\qqq{d}$. 
We view both $\ZL$ and $\Zn$ as discrete $d$-dimensional tori. We will denote the vertices of $\ZL$ by $x,y,\ldots,$ and denote those of $\Zn$ by $a,b,\ldots$ Given $x,y\in \ZL$ and $a,b\in \Zn$, we denote the periodic representatives of $x-y$ and $a-b$ by
$\rep{x-y}_{WL}$ and $\rep{a-b}_L$, respectively:  
\be\label{representativeL}\rep{x-y}_{WL}:= \left((x-y)+(WL)\Z^d\right)\cap \ZL,\quad \rep{a-b}_L:= \left((a-b)+L\Z^d\right)\cap \Z_L^d.\ee
For definiteness, we use the $L^\infty$-metric to define (periodic) distances on $\ZL$ and $\Zn$: 
\[ |x-y|\equiv \|\rep{x-y}_{WL}\|_{\infty},\quad \forall x,y\in \ZL,\quad \text{and}\quad |a-b|\equiv \|\rep{a-b}_L\|_{\infty},\quad \forall a,b\in \Zn.\]
%\[\dist(x,y):=\|\rep{x-y}_{WL}\|_1,\quad \forall x,y\in \ZL,\quad \text{and}\quad \dist(a,b):=\|\rep{a-b}_L\|_1,\quad \forall a,b\in \Zn,\] 
%which correspond to the (periodic) distances on $\ZL$ and $\Zn$, respectively. The normalization by $1/d$ ensures that \smash{$\max_{a,b\in \Zn}|a-b|\le L/2$}. 
We write $x\sim y$ if $x$ and $y$ are neighbors in $\ZL$, and similarly $a\sim b$ if $a$ and $b$ are neighbors in $\Zn$. 
% Finally, given any $\ZL\times \ZL$ matrix $A$, we denote \textcolor{blue}{remove it, this notation is only used once in the whole paper}its $(a,b)$-th block by 
% \be\label{eq:submatrixnotation}A|_{[a][b]}:=(A_{xy}:x\in [a], y\in [b]), 
% \ee
% which is a $W^d\times W^d$ submatrix of $A$, corresponding to the interaction between blocks $a$ and $b$.  
%denote the $(a,b)$-th block of $A$---the block $(A_{xy}:x\in \cal I_a, y\in \cal I_b)$---which is a $W^d\times W^d$ matrix. 
% \begin{align}\label{Japanesebracket} |x-y|\equiv \|x-y\|_L,\quad &\langle x-y \rangle \equiv \|x-y\|_L + W,\quad \text{for} \ \ x,y \in \ZL, \\
% \label{Japanesebracket2} |[x]-[y]|\equiv \|[x]-[y]\|_n,\quad &\langle [x]-[y] \rangle \equiv \|[x]-[y]\|_n + 1, \quad \text{for} \ \  x,y \in \Zn.
%\end{align}
%We use $x\sim y$ to mean that $x$ and $y$ are neighbors on $\ZL$, i.e., $|x-y|=1$. Similarly, $a\sim b$ means that $a$ and $b$ are neighbors on $\Zn$. 



\subsection{Main results for the random band matrix model}\label{subsec:main}
%\paragraph{Wegner orbital ($\WO$) model.} 
Our \emph{random band matrix} (or referred to as the \emph{Wegner orbital model} following \cite{PelSchShaSod}) is defined by a complex Hermitian random block Hamiltonian \( H=(H_{xy}:x,y\in \ZL) \), where the entries $H_{xy}$ are independent (up to the Hermitian symmetry $H_{xy}=\overline{ H}_{yx}$) Gaussian random variables. More precisely, given a symmetric doubly stochastic variance matrix $S=(S_{xy}:x,y\in \ZL)$, the diagonal entries of $H$ are real Gaussian random variables, and the off-diagonal entries are complex Gaussian random variables, distributed as follows: %\textcolor{blue}{use ${\bf 1}_{x=y}$, ${\bf 1}_{x\ne y}$ to write it into 1 line }
%the diagonal and off-diagonal entries of $H$ are {\bf real} and {\bf complex} Gaussian random variables, respectively, distributed as:
\be\label{bandcw0}
H_{xy}\sim \mathcal{N}_{\R}(0, S_{xy}) \cdot \mathbf 1_{x=y} + \mathcal{N}_{\C}(0, S_{xy}) \cdot \mathbf 1_{x\ne y}.
% \begin{cases} 
% ,\quad &\text{for} \ \  x = y \\
% ,\quad &\text{for} \ \  x\ne y
% \end{cases}. 
\ee
For $x\in [a]$ and $y\in [b]$, the variance matrix $S\equiv S(\ilambda)$ is given by 
\be\label{eq:variancematrix}
S_{xy}\equiv\var (H_{xy}) := W^{-d}S^{(\sB)}_{ab}(\ilambda),\quad \text{with}\quad 
S^{(\sB)}_{ab}(\ilambda):=\frac{1}{1+2d\ilambda^2} \mathbf 1_{a=b} + \frac{\ilambda^2}{1+2d\ilambda^2} \mathbf 1_{a\sim b},
\ee
where $\ilambda>0$ (recall \eqref{def:ilambda}) is a coupling parameter that quantifies the interaction strength between neighboring blocks, and \( S^{(\sB)}(\ilambda) \) denotes an \( L ^d\times L^d \) matrix. 
Informally, the matrix $H$ consists of i.i.d.~GUE blocks on the diagonal and i.i.d.~Ginibre blocks (up to Hermitian symmetry and scaling) on the off-diagonal. 
Note that when $\ilambda=1$, the present model reduces to the standard block random band matrices studied in \cite{YY_25, DYYY25}.


%The parameter $\ilambda^2$ plays the role of a coupling parameter, quantifying the interaction strength between neighboring blocks. This perspective aligns with that of the Wegner orbital model (see e.g., \cite{PelSchShaSod}). By combining the techniques developed in this work with the tools from \cite{RBSO1D}, our results can be extended to the sharp delocalization threshold $\ilambda^2\gg W^{-d}$ for the Wegner orbital model. 
%Simply speaking, in our model $H$, the diagonal blocks are i.i.d.~GUE matrices, while the off-diagonal blocks are i.i.d.~Ginibre matrices up to Hermitian symmetry and some rescaling parameters. Here, $\ilambda^2$ is a key \emph{coupling parameter}, as it describes the coupling strength between neighboring blocks. This is the viewpoint taken by the Wegner orbital model \cite{PelSchShaSod}. By combining the techniques developed here with the tools developed in \cite{RBSO1D}, we can push our result to the sharp threshold $\ilambda^2\gg W^{-d}$ for the Wegner orbital model. 
% such that
% \[
% H_{xy} \sim \mathcal{CN}(0, S_{ij}), \quad  x,y \in (\mathbb{Z}_{WL})^d,
% \quad W, L\in \mathbb N\] 
% Clearly, \( S = S^T \), and it satisfies
% \[
% \sum_x S_{xy} = 1.
% \]
% We will follow the convention that indices \( a, b, c, \dots \) are elements of \( (\mathbb{Z}_L)^d \), while \( x, y, z, \dots \) represent indices in \( (\mathbb{Z}_{WL})^d \). 
 

Let \( \lambda_1 \leq \lambda_2 \leq \dots \leq \lambda_N \) denote the eigenvalues of \( H \). The corresponding normalized eigenvectors of $H$ are denoted by \( \left( \boldsymbol{\psi}_k \right)_{k=1}^N \). 
% such that 
% \[
% H \boldsymbol{\psi}_k = \lambda_k \boldsymbol{\psi}_k,\quad  k\in \qqq{N}. 
% \]
It is well-known that the empirical spectral measure $N^{-1}\sum_{k=1}^N \delta_{\lambda_k}$ converges almost surely to the Wigner semicircle law \cite{Wigner} with density 
\(
\rho_{\mathrm{sc}}(x) = \sqrt{(4 - x^2)_+}/{2\pi}.
\)
Define the Green's function (or resolvent) of the Hamiltonian $H$ as
\be\label{def_Green}
G(z):=(H-z)^{-1} ,  \quad z\in \C_+.
\ee
It is also known that as $N\to \infty$, $G(z)$ converges to the scalar matrix $m(z)I_N$ entrywise, where $m(z)$ denotes the Stieltjes transform of \( \rho_{\mathrm{sc}} \), defined by
\be\label{eq:defmzsc}
m(z)\equiv m_{\mathrm{sc}}(z):=\frac{-z+\sqrt{z^2-4}}{2}= \int_{\mathbb{R}} \frac{\rho_{\mathrm{sc}}(x)}{x - z} \dd x.
\ee
Moreover, we define the matrix\footnote{We introduce this notation to maintain consistency with the block Anderson model, where $M(z)$ is no longer a scalar matrix proportional to $m(z)$; see \eqref{self_m} and \eqref{def_G0}.}
\be\label{eq:defMzsc} 
M(z)\equiv M_N(z):=m(z) I_N. 
\ee


 





We now state the main results for the random band matrix model. %The model setup and corresponding main results for the block Anderson model will be presented separately in \Cref{sec:main_BA}. 
Our first main result establishes the delocalization of the bulk eigenvectors of $H$ in dimensions $d\ge 3$.  

\begin{theorem}[Delocalization]\label{MR:decol}
Fix any dimension $d\ge 3$, and consider the random band matrix model defined above. Assume there exist constants $\fc,\fd>0$ such that 
\begin{equation}\label{Main_DEL_COND}
    W \ge N^{ \mathfrak c} ,
\end{equation}
and that $\ilambda$ satisfies the condition 
\be\label{eq:WO} W^{-d/2+\fd}\le \ilambda \le \fd^{-1}.\ee
% and for the block Anderson model, assume 
% \be\label{eq:BA} W^{-d/2+\fd}\le \ilambda \le W^{-\fd}.\ee
Then, for any small constants $\kappa,\tau>0$ and large constant $D>0$, the following delocalization estimate holds, provided $N$ is sufficiently large:
\be\label{eq:psikLinfty}
\mathbb{P}\left( \max_{k:|\lambda_k|\le 2-\kappa}\left\|\boldsymbol{\psi}_k\right\|_{\infty}^2 \le N^{-1+\tau}  \right) \ge 1-N^{-D}.
\ee
\end{theorem}

Since the eigenvectors $\boldsymbol{\psi}_k$ are $L^2$-normalized, the $L^\infty$-bound in \eqref{eq:psikLinfty} implies that bulk eigenvectors have localization length at least $\Omega((WL)^{1-\e})$, for any small constant $\e>0$. The upper bound in condition \eqref{eq:WO} is not essential; we include it only for clarity of presentation, as our main interest is the small $\ilambda$ regime.\footnote{All results below remain valid for $\ilambda \gg 1$, provided $\ilambda$ in the relevant equations is replaced by $\ilambda \wedge 1$.} 
The key aspect of condition \eqref{eq:WO} is the lower bound $\ilambda\ge W^{-d/2+\fd}$. Indeed, by the fractional moment method \cite{Aizenman1993}, it was shown in \cite{PelSchShaSod} that when $\ilambda\ll W^{-d/2}$, the eigenvectors of both the random band matrix model and the block Anderson model (defined in \Cref{sec:main_BA} below) are localized with localization length of order $\OO(W)$. Combined with this result, \Cref{MR:decol} demonstrates a localization--delocalization transition for the random band matrix model in the bulk spectrum as $\ilambda$ crosses the critical threshold $W^{-d/2}$.

%The crucial aspect of condition \eqref{eq:WO} is the lower bound $\ilambda\ge W^{-d/2+\fd}$. In fact, utilizing the fractional moment method \cite{Aizenman1993}, it was shown in \cite{PelSchShaSod} that in the regime \(\ilambda\ll W^{-d/2}\), the eigenvectors of both the random band matrix model and the block Anderson model (defined in \Cref{sec:main_BA} below) are localized with localization length of order $\OO(W)$. Combined with this result, \Cref{MR:decol} gives a localization-delocalization transition of the random band matrix model inside the bulk of the spectrum as $\ilambda$ crosses the critical threshold $W^{-d/2}$.   



\Cref{MR:decol} follows immediately from the following sharp local law for the Green's function of $H$, defined as in \eqref{def_Green}. For simplicity of notation, given constants $\kappa,\e>0$, define the spectral domain 
\be\label{eq:spectral_domain}
\mathbf D_{\kappa,\e}:=\{z=\hat{E}+\ii\eta \in \C_+: |\hat{E}|\le 2-\kappa, N^{-1+\e}\le \eta\le 1\},
\ee
and the simplified notation 
\be\label{eq:calBetaK}\cal B_{\eta,K}:= \frac{(\ilambda^{2}+\eta)^{-1}}{W^2(K+W)^{d-2}} + \frac{1}{N\eta},\quad \forall K\ge 0.\ee

\begin{theorem}[Local semicircle law]\label{MR:locSC} 
In the setting of \Cref{MR:decol}, for any (small) constants $\kappa,\e, \tau>0$ and (large) constant $D>0$, the following events hold with probability $\ge 1-N^{-D}$ for large enough $N$: 
\begin{align}\label{G_bound}
&\bigcap_{z=\hat{E}+\ii\eta\in \mathbf D_{\kappa,\e}}\bigcap_{x,y\in \ZL} \left\{ |G_{xy}(z) - M_{xy}(z)|^2  \le W^\tau  \mathcal B_{\eta,|x-y|}\right\}\, , 
\\
\label{G_bound_ave}
&\bigcap_{z=\hat{E}+\ii\eta\in \mathbf D_{\kappa,\e}}\bigg\{\max_{a\in\Zn} \Big|W^{-d}\sum_{x\in [a]}G_{xx}(z) - m(z) \Big| \le W^\tau  \mathcal B_{\eta,0}\bigg\}\, , 
\end{align}
where $G$ is defined in \eqref{def_Green}, $m(z)$ and $M(z)$ are defined in \eqref{eq:defmzsc} and \eqref{eq:defMzsc}, and $[a]$ is defined in \eqref{eq:blockIa}. % for the Wegner orbital model, and $m(z)$ and $M(z)$ are defined in \eqref{self_m} and \eqref{def_G0} for the block Anderson model. 
\end{theorem}

\begin{proof}[\bf Proof of \Cref{MR:decol}]
Theorem \ref{MR:decol} follows directly from the entrywise local law \eqref{G_bound} via the bound 
\be\label{eq:ukx}|\boldsymbol{\psi}_k(x)|^2 \le \eta\im G_{xx}(\lambda_k + \ii \eta),\quad \forall \eta>0\, .\ee 
Applying the local law \eqref{G_bound} to $G_{xx}(\lambda_k+\ii\eta)$ with $\eta = N^{-1+\tau}$, we deduce that $\im G_{xx}(\lambda_k+\ii\eta)=\OO(1)$ with high probability. Combined with \eqref{eq:ukx}, this completes the proof.
\end{proof}

Under the assumptions of \Cref{MR:decol}, we can further establish a stronger \emph{quantum unique ergodicity} (QUE) estimate for the bulk eigenvectors of $H$, albeit at the cost of a slightly weaker probability bound. 
Roughly speaking, the QUE estimates \eqref{Meq:QUE} and \eqref{Meq:QUE2} below indicate that every bulk eigenvector of $H$ is asymptotically uniformly distributed (in the sense of $L^2$-mass) across all scales larger than $W$. In particular, this implies that the localization length of every bulk eigenvector is of order $\Omega(WL)$ with probability $1-\oo(1)$.

\begin{theorem}[Quantum unique ergodicity]\label{MR:QUE} 
In the setting of \Cref{MR:decol}, given $E\in [-2+\kappa,2-\kappa]$ and a constant $\e_0\in (0,\fd/2)$, define the interval  
\be\label{eq:defIE} {\cal I}_E\equiv {\cal I}_E(\e_0):=\left\{x: |x-E|\le W^{-\e_0}\big({\ilambda W^{d/2}}/N\big)\right\}. %\quad \text{where}\quad  \cor \eta_{0}:={W\ilambda}/{N}.
\ee
For each $d\ge 3$ and constant $0<c<\e_0\wedge(\fd/5)$, the following estimate holds for any small constant $\tau>0$:
\begin{equation}\label{Meq:QUE}
\sup_{E: |E|\le 2-\kappa}  
\max_{ a\in \Zn} \P\bigg(\max_{i,j:\lambda_i, \lambda_j \in {\cal I}_E} 
 \bigg| \sum_{x\in[a]}\overline \bpsi_i(x)\bpsi_j(x)-\frac{W^{d}}{N}\delta_{ij}  \bigg| \ge \frac{W^{d-c}}{N} \bigg) \le  W^{-(2\e_0)\wedge(2\fd/5)+2c+\tau} \, 
\end{equation}
provided $N$ is large enough. More generally, for any subset $A\subset \Zn$, we have
\begin{equation}\label{Meq:QUE2}
\sup_{E: |E|\le 2-\kappa} \mathbb{P}\bigg(\max_{k: \lambda_k\in {\cal I}_E} 
\bigg|\sum_{a\in A}\sum_{x\in [a]}\left|\bpsi_k(x)\right |^2 -\frac{W^d}{N}|A|\bigg| \ge  \frac{W^{d-c}}{N}|A| \bigg) \le  W^{-(2\e_0)\wedge(2\fd/5)+2c+\tau} .
\end{equation}  

% \begin{align}\label{tau} 
% 0< \tau <  \frac {  \mathfrak c } 2, 
% \end{align} 
% where $ \mathfrak c$ was defined in \eqref{Main_DEL_COND}.
% Then generalized quantum unique holds in high probability, i.e.,  for any small constant $\kappa>0$  and large $D>0$,   there exists  a large  $N_0$ such that for all $N \geqslant N_0$, 
% \begin{equation}\label{Meq:QUE}
% \max_{E: \,|E|<2-\kappa}\; 
% \max_{ a\,\in\, \mathbb Z_L}\;\mathbb{P}\left(\max_{i, j \in {\cal J}_E} 
%  \left|N (\psi_i^*\left(E_a-N^{-1}\right) \psi_j)\right|^2 \ge N^{-\tau /6} \right) \le  N^{-\tau  /6}.
% \end{equation}
% Here 
% $$
% \mathcal{J}_E:=\left\{k:\left|\ilambda_k-E\right| \leq N^{-1-\tau} W^{2/3}\right\} .
% $$
% Furthermore, for  any  subset $A\subset \mathbb Z_L$,
% \begin{equation}\label{Meq:QUE2}
% \max_{E: \,|E|<2-\kappa}\; 
% \max_{ A\,\subset\, \mathbb Z_L}\;\mathbb{P}\left(\max_{k\in {\cal J}_E} 
% \left|\,\sum_{x\in {\cal I}_a}\sum_{a\,\in\, A}\left|\boldsymbol{\psi}_k(x)\right |^2 -\frac{|A|\cdot W}{N}\right| \geqslant \frac{|A|\cdot W}{N^{1+\tau /6}}  \right) \le N^{-\tau /6}
% \end{equation}
\end{theorem}

As an important consequence of the above QUE estimates, and by employing the Green's function comparison argument developed in \cite{Xu:2024aa}, we can derive the following \emph{universality of the local bulk eigenvalue statistics} for our random band matrices. 
Let $p_{H}(\lambda_1,\ldots, \lambda_{N})$ denote the joint symmetrized probability density of the (unordered) eigenvalues of $H$. For any $1\le n \le N$, define the $n$-point correlation function as
$$
p_{H}^{(n)}\left(\lambda_1, \cdots, \lambda_n\right)
:=\int_{\R^{N-n}} p_{H}\left(\lambda_1, \cdots, \lambda_{N}\right) \mathrm{d} \lambda_{n+1} \cdots \mathrm{d} \lambda_{N}.
$$
Moreover, let $p_{\rm{GUE}}^{(n)}$ denote the corresponding $n$-point correlation function for an $N\times N$ GUE matrix.

\begin{theorem}[Bulk universality]\label{Thm: B_Univ}
In the setting of \Cref{MR:decol}, let $O \in C_c^{\infty}\left(\mathbb{R}^n\right)$ be an arbitrary smooth, compactly supported function. Then, for any $|E|\le 2-\kappa$ and fixed $n\in \N$, we have
\begin{equation}\label{eq:universality}
 \lim_{N\to \infty}	\int_{\mathbb{R}^n} \mathrm{~d} \boldsymbol{\alpha}\; O(\boldsymbol{\alpha}) \left[p_{H}^{(n)}-p_{\rm{GUE}}^{(n)}\right]\left(E+\frac{\alpha_1}{N  }, \ldots, E+\frac{\alpha_n}{N}\right)   =0, 
 \end{equation} 
where $\bal$ denotes $\boldsymbol{\alpha}=\left(\alpha_1, \ldots, \alpha_n\right)$. 
%Suppose that  the assumptions of Theorem \ref{MR:decol} holds. For any fixed $k\in \mathbb N$, the $k$ point correlation function of $H$ converges to that of GUE in the following sense.  
 \end{theorem}


This bulk universality result shows that, for our random band matrices, the local eigenvalue gap statistics near any fixed bulk energy level $E$ asymptotically coincide with those of Wigner matrices. However, we note that the distribution of an individual bulk eigenvalue $\lambda_k$ of $H$ may differ significantly from that of Wigner matrices, as $\lambda_k$ can exhibit fluctuations that are much larger than $N^{-1}$. 

%As suggested by Theorem 1 of \cite{Sch2015}, the fluctuations of bulk eigenvalues in band matrices are of order at least $(WN)^{-1/2}$, whereas the fluctuations for the eigenvalues of Wigner matrices are of order $ N^{-1}$ \cite{erdHos2012rigidity}.

 



Similar to the one- and two-dimensional cases \cite{YY_25,DYYY25,erdHos2025zigzag,RBSO1D}, our random band matrix model in dimensions $d\ge 3$ also satisfies the \emph{quantum diffusion conjecture}, which serves as a key input for establishing the QUE estimates in \Cref{MR:QUE}.
To state it, we define the following $\Theta$-matrices:
\be \label{def:Theta}
\Theta^{(+,-)}(z)=\Theta^{(-,+)}(z):= \frac{1}{1-|m(z)|^2 S^{(\sB)}} ,\quad \Theta^{(+,+)}(z)=(\Theta^{(-,-)}(z))^*:= \frac{1}{1-m(z)^2 S^{(\sB)}} , 
\ee
where we recall the matrix $S^{(\sB)}$ defined in \eqref{eq:variancematrix} and $m(z)$ defined in \eqref{eq:defmzsc}. 


% {To define the quantum diffusion, we need to define some notations. It is believed that $G(z)$ exhibits \emph{quantum diffusion}, in the sense that $\E|G_{xy}(z)|^2$ is well-approximated by the entries of the deterministic matrix $|m|^2(1 - |m|^2 S)^{-1}$. Using \eqref{eq;sproduct}, we can write that 
% \[
% \frac{1}{1 - |m|^2 S} = I + \frac{|m|^2S^{(\sB)}(\ilambda)}{1 - |m|^2 S^{(\sB)}(\ilambda)} \otimes S_W.
% \]
% In Lemma \ref{lem_propTH}, we will show that the matrix $(1 - |m|^2 S^{(\sB)}(\ilambda))^{-1}$ exhibits exponential decay at the block level on the length scale 
% \be\label{def_ellz}
%  \ell(z) := \min \left( \ilambda\eta^{-1/2}+1, L \right), \quad \text{with}\quad \eta = \mathrm{Im}(z),
% \ee
% which is referred to as the \emph{diffusion length} for $G(z)$.}

\begin{theorem}[Quantum diffusion]\label{MR:QuDiff}
    In the setting of \Cref{MR:decol}, for any (small) constants $\kappa,\e, \tau>0$ and (large) constant $D>0$, the following events hold with probability $\ge 1-N^{-D}$ for all $a,b\in \Zn$ and for large enough $N$: 
    \begin{align}
    &\bigcap_{z=\hat{E}+\ii\eta\in \mathbf D_{\kappa,\e}}
    	\bigg\{\bigg|\frac{1}{W^{2d}}
    	\sum_{x\in[a],y\in[b]}|G_{xy}(z)|^2 - \frac{|m|^2\Theta_{ab}^{(+,-)}(z)}{W^d}\bigg| \le W^\tau \br{\p{(\mathcal B_{\eta,0})^{\frac 1 5} \mathcal B_{\eta,W|a-b|}}\wedge (\mathcal B_{\eta,0})^{2}}\bigg\}\, , %\frac{1}{W^{d}}\p{\frac{|m|^2}{1-|m|^2S^{(\sB)}}}_{ab}
    	\label{eq:diffu1}\\
    	&\bigcap_{z=\hat{E}+\ii\eta\in \mathbf D_{\kappa,\e}}
    	\bigg\{\bigg|\frac{1}{W^{2d}}
    \sum_{x\in[a],y\in[b]}\p{G_{xy}G_{yx}}(z) - \frac{m^2\Theta^{(+,+)}_{ab}(z)}{W^d} \bigg| \le W^\tau \br{\p{(\mathcal B_{\eta,0})^{\frac 1 5} \mathcal B_{\eta,W|a-b|}}\wedge (\mathcal B_{\eta,0})^{2}}\bigg\}\, %\frac{1}{W^{d}}\p{\frac{m^2}{1-m^2S^{(\sB)}}}_{ab}
    	\label{eq:diffu2}.
    \end{align}
%where $G$ and $m$ are defined in \eqref{def_Green} and \eqref{eq:defmzsc}, respectively. 
% \end{theorem} 
% \begin{theorem}[Quantum diffusion with expectation for small $\eta$]\label{MR:QDiff_EE}
Moreover, for each $z=\hat{E}+\ii\eta\in \mathbf D_{\kappa,\e}$, the expectations $\E|G_{xy}(z)|^2$ and $\E\p{G_{xy}(z)G_{yx}(z)}$ satisfy the following bounds for any small constant $\tau>0$ and large enough $N$: 
\begin{align}\label{Meq:QdS1}
 \max_{a,b }  \bigg|\frac{1}{W^{2d}}
    	\sum_{x\in[a],y\in[b]}\E|G_{xy}(z)|^2 - \frac{|m|^2\Theta_{ab}^{(+,-)}(z)}{W^d} \bigg| \le W^\tau (\cal B_{\eta,0})^{2} \left(\p{\ilambda^2W^{d}}^{-1/ 5}+\cal B_{\eta,0}\right),   \\
 \label{Meq:QdS2}
 \max_{a,b }  \bigg|\frac{1}{W^{2d}}
    	\sum_{x\in[a],y\in[b]}\E(G_{xy}G_{yx})(z) - \frac{m^2\Theta_{ab}^{(+,+)}(z)}{W^d}\bigg|\le W^\tau (\cal B_{\eta,0})^{2} \left(\p{\ilambda^2W^{d}}^{-1/ 5}+\cal B_{\eta,0}\right) .
\end{align}
\end{theorem}

Note that when $\eta \le \ilambda^{2}/L^{d}$, we have 
\be\label{eq:BetaK}
\cal B_{\eta,K}\asymp (N\eta)^{-1} \ge (\ilambda^2W^d)^{-1}, \quad \forall K\ge 0 .
\ee
In this case, the expected estimates in \eqref{Meq:QdS1} and \eqref{Meq:QdS2} provide improvements over those in \eqref{eq:diffu1} and \eqref{eq:diffu2}, which in turn lead to the QUE estimates in \Cref{MR:QUE}. 
 

\begin{proof}[\bf Proof of Theorem \ref{MR:QUE}]  
With the quantum diffusion estimates \eqref{Meq:QdS1} and \eqref{Meq:QdS2}, the proof of \Cref{MR:QUE} is similar to those for \cite[Theorem 2.5]{YY_25}, \cite[Theorem 2.4]{DYYY25}, and \cite[Theorem 2.2]{RBSO1D}. We now briefly outline the proof without giving all details. 

Take $z=E+\ii\eta$ with $|E|\le 2-\kappa$ and $\eta=W^{-\e_0}({\ilambda W^{d/2}}/N)\ge W^{\fd-\e_0}/N$. With the spectral decomposition of $G(z)$ and the definition of $ {\cal I}_E$, we find that 
\begin{align}\label{ssfa2}
\E \sum_{i,j: \lambda_i,\lambda_j\in {\cal I}_E} \left|\bpsi_i^*\left(E_{a}-N^{-1}\right) \bpsi_j\right|^2 &\lesssim \eta^2 \E \tr \left[\operatorname{Im} G(z)\left(E_{a}-N^{-1}\right) \operatorname{Im} G(z)\left(E_{a}-N^{-1}\right)\right] \nonumber \\
&=\frac{\eta^2}{L^{2d}}\sum_{b,b'\in \Zn}
\E  
 \tr \left[\operatorname{Im} G(z)\left(E_{a}-E_{b}\right) \operatorname{Im} G(z)\left(E_{a}-E_{b'}\right)\right],
\end{align}
where we recall that $E_a$ is the block-averaging matrix restricted to the $a$-th block: $(E_a)_{xy} = W^{-d}\mathbf 1(x=y\in [a])$. Expanding $\im G$ as $\im G=(G-G^*)/(2\ii)$ and applying the QUE estimates \eqref{Meq:QdS1} and \eqref{Meq:QdS2}, we can bound the right-hand side of \eqref{ssfa2} by   
\begin{align}\label{ssfa2_deter}
    W^\tau \eta^2 \p{\frac{1}{N\eta}}^2\left[\frac{1}{N\eta} + \frac{1}{(\ilambda^2W^d)^{1/5}}\right]+  \frac{C\eta^2}{W^d} \max_{\sig,\sig'\in \{+,-\}}\max_{a,b,b'}\left|\Theta^{(\sig,\sig')}_{ab}-\Theta^{(\sig,\sig')}_{ab'}\right|
\end{align}
for any small constant $\tau>0$ and a large constant $C>0$, where the first term comes from the application of \eqref{Meq:QdS1}--\eqref{eq:BetaK}. 
Then, with the estimate \eqref{prop:ThfadC0} below for the $\Theta$-matrices (note that $\Theta^{(\sig,\sig')}_{ab}-\Theta^{(\sig,\sig')}_{ab'}=\zTheta^{(\sig,\sig')}_{ab}-\zTheta^{(\sig,\sig')}_{ab'}$ by \eqref{def_Thxi0}), we get 
$$
\max_{\sig,\sig'\in \{+,-\}}\max_{a,b,b'}\left|\Theta^{(\sig,\sig')}_{ab}-\Theta^{(\sig,\sig')}_{ab'}\right|\prec \ilambda^{-2} .
$$
Plugging it into \eqref{ssfa2_deter} and further into \eqref{ssfa2} yields that 
\begin{align*}
\E \sum_{i,j: \lambda_i,\lambda_j\in {\cal I}_E} \left|\bpsi_i^*\left(E_{a}-N^{-1}\right) \bpsi_j\right|^2 &\lesssim {W^\tau}N^{-2}\cdot \p{W^{-\fd+\e_0}+W^{-2\fd/5}+W^{-2\e_0}}. 
\end{align*}
Finally, applying Markov's inequality concludes \eqref{Meq:QUE}. The proof of \eqref{Meq:QUE2} follows a similar argument, where we simply replace $E_{a}$ in \eqref{ssfa2} with $|A|^{-1}\sum_{a\in A} E_{a}$, after which all subsequent arguments remain valid.  
% \textcolor{blue}{Short version. The proof is basically the same as the one for $d=2$\cite{DYYY25}, which comes from the $d=1$ in \cite{YY_25}. With the same argument and \eqref{Meq:QdS1} and \eqref{Meq:QdS2}, we have 
%  \begin{align}
%  \;& \eta^2 \E  \tr \left(\operatorname{Im} G(z)\left(E_a-N^{-1}\right) \operatorname{Im} G(z)\left(E_a-N^{-1}\right)\right)
% \nonumber \\
% \le \;& N^{-2}(W^{-d/5}+\frac{1}{N\eta})\cdot W^{\delta} + \eta^2 W^{-d}\max_{a,b,a',b'}\;\max_{ \xi=|m|^2 \text{ or } \xi=m^2}
% \left| \left(\frac{1}{1-\xi\cdot S^{(B)}}\right)_{ab}-\left(\frac{1}{1-\xi\cdot S^{(B)}}\right)_{a'b'}\right|
% \nonumber 
%  \end{align}
% for any $\delta>0$ small.
%  Furthermore, with $\Theta$ property, we have 
% $$
% \max _{a, b, a^{\prime}, b^{\prime}} \max _{\xi=|m|^2 \text { or } \xi=m^2}\left|\left(\frac{1}{1-\xi \cdot S^{(B)}}\right)_{a b}-\left(\frac{1}{1-\xi \cdot S^{(B)}}\right)_{a^{\prime} b^{\prime}}\right| \prec 1  .
% $$
% Combine them, we obtain the key estimate 
% $$
% \eta^2 \E  \tr \left(\operatorname{Im} G(z)\left(E_a-N^{-1}\right) \operatorname{Im} G(z)\left(E_a-N^{-1}\right)\right)\le N^{-2-\tau/3}
% $$
% which completes the proof as in \cite{DYYY25} for $d=2$} 
\end{proof}
  
%Finally, we prove the bulk universality stated in Theorem \ref{Thm: B_Univ}.
 
\begin{proof}[\bf Proof of Theorem \ref{Thm: B_Univ}]
%Using the delocalization estimate from \Cref{MR:decol}, the local laws from \Cref{MR:locSC}, 
Using \Cref{MR:decol}, \Cref{MR:locSC}, and the QUE estimate \eqref{Meq:QUE} as inputs, the proof of Theorem \ref{Thm: B_Univ} follows from the Green's function comparison method developed in \cite{Xu:2024aa}. 
The argument is essentially identical to that used for one-dimensional \cite[Theorem~2.6]{YY_25} and two-dimensional \cite[Theorem~2.6]{DYYY25} random band matrices. For instance, adapting the proof of \cite[Theorem~2.6]{DYYY25}, we only need to adjust certain parameters in the proofs of equations (2.22) and (2.23) therein.
In that setting, given $y\in \ZL$, the ``bad'' event $\mathcal{B}\equiv \mathcal{B}(y)$ was defined by the existence of an index $\alpha$ such that $\lvert \lambda_\alpha - E\rvert \le N^{-1+\mathfrak{c}/6}$ and $\lvert M_{y,\alpha}\rvert \ge N^{-\mathfrak{c}/18}$ ($\fc$ is the constant in the assumption $W \ge N^{\mathfrak{c}}$ and $M_{y,\al}$ is defined below equation (2.24) of \cite{DYYY25}).\footnote{In dimension $d\ge 3$, $M_{y,\al}$ is defined with the weights of the variance matrix $S$: $M_{y,\al}:=N\sum_{x\in\ZL}S_{xy}|\bpsi_\al(x)|^2-1=\sum_{b\in\Zn}S^{(\sB)}_{a_yb}\big(NW^{-d}\sum_{x\in[b]}|\bpsi_\al(x)|^2-1\big)$, where $a_y\in\Zn$ denotes the block containing $y$, i.e., $y\in[a_y]$. This is an average over at most $2d+1$ blocks (for the block Anderson model, where $S$ is given by \eqref{bandcwV}, only the block $[a_y]$ appears). Hence, if $|M_{y,\al}|\ge W^{-\fd/6}$, then the event in \eqref{Meq:QUE} with $i=j=\al$ and $c=\fd/6$ occurs for at least one of these blocks, and $\P(\mathcal B(y))$ is bounded by $2d+1$ times the right-hand side of \eqref{Meq:QUE}.} In our setting, we modify the definition of $\cal B$ to
\(
\mathcal{B}(y)
:= \bigl\{ \exists \alpha: \lvert \lambda_\alpha - E\rvert \le N^{-1} W^{\mathfrak{d}/3}, \lvert M_{y,\alpha}\rvert \ge W^{-\mathfrak{d}/6} \bigr\}.
\)
Applying \eqref{Meq:QUE} with $\e_0=\mathfrak{d}/3$ and $c=\mathfrak{d}/6$, we can deduce that
\(
\mathbb{P}(\mathcal{B}) \le W^{-\mathfrak{d}/15+\tau}.
\)
With these modified parameters and the new definition of $\cal B$, the remainder of the proof of \eqref{eq:universality} proceeds exactly as in the proof of \cite[Theorem~2.6]{DYYY25}. %since only the signs of the exponents matter, not their precise values. 
Hence, we omit further details.
\end{proof}

% \textcolor{blue}{Following the proof for the two-dimensional random band matrices in \cite[Theorem~2.6]{DYYY25}, we only need to adjust certain parameters in the proofs of (2.22) and (2.23) at there. There, the ``bad'' event $\mathcal{B}$ is defined by the existence of an index $\alpha$ such that $\lvert \lambda_\alpha - E\rvert \le N^{-1+\mathfrak{c}/6}$ and $\lvert M_{y,\alpha}\rvert \ge N^{-\mathfrak{c}/18}$. Recall that in that setting $W \ge N^{\mathfrak{c}}$.
% For $d \ge 3$, we modify the definition (for fixed $y$) to
% \[
% \mathcal{B}
% := \bigl\{ \exists \alpha:\ \lvert \lambda_\alpha - E\rvert \le N^{-1} W^{\mathfrak{d}/3},\quad \lvert M_{y,\alpha}\rvert \ge W^{-\mathfrak{d}/6} \bigr\}.
% \]
% Applying Theorem~\ref{MR:QUE} with $\epsilon_0=\mathfrak{d}/3$ and $c=\mathfrak{d}/6$, we obtain
% \[
% \mathbb{P}(\mathcal{B}) \le W^{-\mathfrak{d}/15+\tau}.
% \]
% With these parameter revisions, universality then follows as in the $d \ge 2$ case, since the argument requires only that the exponents have the correct signs, not their precise values.}
% Therefore, we omit the details of the proof.
% \textcolor{blue}{Short version. The proof would be exact the same as  $d=2$ (which origins from $d=1$ \cite{YY_25}). The key parameters are still 
% $\tau=\mathfrak{c} / 3$ and $t \in\left[0, N^{-1+\tau_U}\right]$ and $\eta=N^{-1+2 \tau_U}$ for small enough $\tau_U$. Note for this region of $\eta$, the right hand side of \eqref{Meq:QdS1} and \eqref{Meq:QdS2} become to $(N\eta)^{-3}$, which match the form of estimates used in the proof for $d=2$ paper.}    


\subsection{Main results for the block Anderson model}
\label{sec:main_BA}

%\paragraph{Block Anderson ($\BA$) model.} 

We next define the block Anderson model. We begin by introducing the \emph{random block potential} $V$, an $N\times N$ complex Hermitian random block matrix whose diagonal blocks are i.i.d.~GUE matrices. In other words, $V$ can be regarded as a random band matrix whose entries are distributed according to \eqref{bandcw0}, with variance matrix $S=(S_{xy})$ given by: 
\be\label{bandcwV}
S_{xy}\equiv \var (V_{xy}):=W^{-d}S^{(\sB)}_{ab}(0)= W^{-d} {\bf 1}\left( a =b \right), \quad \text{for} \ \  x\in [a], \ y\in [b],
\ee
where \( S^{(\sB)}(0)=I_{L^d} \) denotes the matrix defined in \eqref{eq:variancematrix} with $\ilambda=0$. 
Next, we define the \emph{block Anderson model} as a random block Schr{\"o}dinger operator of the form
\be\label{eq:H_blocka}
H \equiv H(\ilambda)=V+\ilambda \Psi ,\ee 
where $\ilambda>0$ (recall \eqref{def:ilambda}) is a coupling parameter, and $\Psi$ represents the interaction Hamiltonian that introduces hopping between neighboring blocks. 
For definiteness, we focus on the classical block Anderson model in which each block is the identity matrix: % (recall the notation in \eqref{eq:submatrixnotation}):  
\be\label{eq:Psi3D}
\Psi= \Psi^{(\sB)}\otimes
    I_{W^d},\quad \text{where}\quad  \Psi^{(\sB)}_{ab}=\mathbf 1(a\sim b).\ee 
For the block Anderson model, we again denote its eigenvalues by \( \lambda_1 \leq \lambda_2 \leq \dots \leq \lambda_N \), and the corresponding normalized eigenvectors by \( \left( \boldsymbol{\psi}_k \right)_{k=1}^N \). 
The Green's function $G(z)$ is defined analogously to \eqref{def_Green}. 
\begin{remark}
For simplicity, we have slightly abused notation by using the same symbols (such as $H$, $S$, $\lambda_i$, $\boldsymbol{\psi}_i$, and $G(z)$ defined above, as well as $m(z)$ and $M(z)$ defined below) for both the random band matrix model and the block Anderson model. When a distinction is needed, we will add a superscript $\BA$ to denote quantities associated with the block Anderson model.
%$\WO$ for the Wegner orbital model or $\BA$ for the block Anderson model setting. 
\end{remark}    

It is well-known that as $N\to \infty$, the empirical spectral measure $N^{-1}\sum_{k=1}^N \delta_{\lambda_k}$ is well approximated by a deterministic probability measure $\mu_{N}$, known as the free convolution of the semicircle law and the empirical measure of $\ilambda\Psi$. This measure has a continuous probability density $\rho_N(x)$ on $\R$ \cite{Biane}, whose support is contained in $[-2-2d\ilambda,\, 2+2d\ilambda]$. (For finite $L$, this support need not be symmetric about $0$, e.g., when $L$ is odd, and it may have gaps.)
The Stieltjes transform $m(z)\equiv m(z,\ilambda)$ of the measure $\mu_{N}$ is defined as the unique solution to the following self-consistent equation
\be\label{self_m}
\frac{1}{N}\tr \frac{1}{\ilambda \Psi -z- m(z)} = m(z)
\ee
such that $\im m(z)>0$ for $z\in \C_+$. 
In addition, we define the $N\times N$ matrix $M(z)\equiv M_N(z,\ilambda)$ and the $L^d\times L^d$ matrix \smash{$M^{(\sB)}(z)\equiv M^{(\sB)}_L(z,g)$} as
\be\label{def_G0}
M(z):= \frac{1}{\ilambda \Psi -z- m(z)} = M^{(\sB)}(z)\otimes I_{W^d},\quad M^{(\sB)}(z):=\frac{1}{\ilambda \Psi^{(\sB)} -z- m(z)}.
 \ee
Note $M$ and $M^{(\sB)}$ are both (complex) symmetric matrices. In the random matrix theory literature (see e.g., \cite{LeeSchSteYau2015,knowles2017anisotropic,He2018,AEK_PTRF,EKS_Forum} for various settings of deformed Wigner-type matrices), the Green's function $G(z)$ is known to converge to $M(z)$ in the sense of local laws. %provided $\im z\gg W^{-d}$. 
(However, existing convergence estimates in the literature are generally non-optimal in the non-mean-field setting $W^d\ll N$.) 
% By the block translation symmetry of $\Psi$ in the block Anderson model, the matrix $M(z)$ satisfies that 
% \be\label{eq:averm}
% W^{-d}\sum_{x\in[a]}M_{xx}=m, \quad \forall [a]\in \Zn.\ee 
%It is known that $m(z)$ is the Stieltjes transform of a probability measure $\mu_{N}$, referred to as the \emph{free convolution of the semicircle law and the empirical measure of $\ilambda\Psi$}. 




\begin{theorem}[Main results for the block Anderson model]\label{MR:decol_BA}
Fix any dimension $d\ge 3$, and assume there exist constants $\fc,\fd>0$ such that \eqref{Main_DEL_COND} and \eqref{eq:WO} hold. 
% together with the condition
% \be\label{eq:BA} W^{-d/2+\fd}\le \ilambda \ll 1. %\le W^{-\fd}.
% \ee
Then, for the block Anderson model defined above, the following results remain valid if, throughout their statements, we replace the energy condition $|E|\le 2-\kappa$ (for $E=\lambda_k$ in \eqref{eq:psikLinfty}, for $E=\hat E$ in \eqref{eq:spectral_domain}, and for the energy $E$ in \Cref{MR:QUE} and \Cref{Thm: B_Univ}) with the bulk condition
\be\label{eq:bulkBA}
E\in \mathbf B_\kappa:=\left\{E\in\R:\ \rho_N(E)\ge \kappa\right\},
\ee
where, by the Stieltjes inversion formula, $\rho_N(E)=\pi^{-1}\im m(E+\ii 0_+,\ilambda)$:
\begin{itemize}
    \item The delocalization estimate \eqref{eq:psikLinfty} holds.

    \item The local laws \eqref{G_bound} and \eqref{G_bound_ave} hold for large enough $N$, where $G$ is defined in \eqref{def_Green}, and $m(z)$ and $M(z)$ are defined in \eqref{self_m} and \eqref{def_G0}, respectively. 

    \item The QUE estimates \eqref{Meq:QUE} and \eqref{Meq:QUE2} hold, and the bulk universality holds in the following form, normalized by the local densities: for any fixed $n\in\N$, $O\in C_c^\infty(\R^n)$, $E\in \mathbf B_\kappa$ and $E'\in(-2,2)$,
\[
 \lim_{N\to \infty}	\int_{\mathbb{R}^n} \mathrm{~d} \boldsymbol{\alpha}\; O(\boldsymbol{\alpha}) \left[\frac{p_{H}^{(n)}\big(E+\frac{\alpha_1}{N\rho_N(E)}, \ldots, E+\frac{\alpha_n}{N\rho_N(E)}\big)}{\rho_N(E)^{n}}-\frac{p_{\rm{GUE}}^{(n)}\big(E'+\frac{\alpha_1}{N\rho_{\mathrm{sc}}(E')}, \ldots, E'+\frac{\alpha_n}{N\rho_{\mathrm{sc}}(E')}\big)}{\rho_{\mathrm{sc}}(E')^{n}}\right]   =0 .
\]
(For the random band matrix model, $\rho_N=\rho_{\mathrm{sc}}$, and the choice $E'=E$ gives back \eqref{eq:universality} after the change of variables $\boldsymbol\alpha\mapsto\rho_{\mathrm{sc}}(E)\boldsymbol\alpha$. For the block Anderson model $\rho_N(E)\neq\rho_{\mathrm{sc}}(E)$ in general, so the comparison with the GUE has to be made at the local density.)
    
    \item Recall the variance matrix \(S\) from \eqref{bandcwV} and the matrix \(M\) from \eqref{def_G0}. Define the \(\Theta\)-matrices as  
\be \label{def:Theta_BA}
\qquad  \quad \Theta^{(+,-)}(z) := \big(1 - M^{(+,-)}(z) S^{(\sB)}\big)^{-1} , \ \ 
   \Theta^{(+,+)}(z) = (\Theta^{(-,-)}(z))^* := \big(1 - M^{(+,+)}(z) S^{(\sB)}\big)^{-1} , 
\ee
where \( S^{(\sB)}=I_{L^d} \), and the matrices $M^{(+,-)}=M^{(-,+)}$ and $M^{(+,+)}=(M^{(-,-)})^*$ are defined by\footnote{Note the definition \eqref{def:Theta_BA} is consistent with \eqref{def:Theta}, where \(M^{(+,-)} = |m|^2 I\) and \(M^{(+,+)} = m^2 I\).}   
\[ M^{(+,-)}_{ab}=\frac{1}{W^d}\sum_{x\in[a],y\in[b]}|M_{xy}|^2=|M^{(\sB)}_{ab}|^2,\quad M^{(+,+)}_{ab}=\frac{1}{W^d}\sum_{x\in[a],y\in[b]}M_{xy}M_{yx}=(M^{(\sB)}_{ab})^2. \] 
% \(M^0\) and \(M^+\) are given by  
% \({M}^0_{xy} := |M_{xy}|^2\) and \(M^+_{xy} := M_{xy} M_{yx}\) for \smash{$x,y \in \ZL$}. 
Then, the quantum diffusion estimates \eqref{eq:diffu1}--\eqref{Meq:QdS2} hold if we replace $|m|^2\Theta_{ab}^{(+,-)}$ and $m^2\Theta_{ab}^{(+,+)}$ with $ (\Theta^{(+,-)}M^{(+,-)})_{ab}$ and $(\Theta^{(+,+)}M^{(+,+)})_{ab}$, respectively.
\end{itemize}

\end{theorem}

%In \eqref{eq:BA}, we have assumed $\ilambda\ll 1$ following our earlier discussion. In \eqref{eq:BA}, 
The lower bound $\ilambda\gg W^{-d/2}$ in \eqref{eq:WO} is sharp, as explained below \Cref{MR:decol}: when $\ilambda\ll W^{-d/2}$, the block Anderson model is localized, as proven in \cite{PelSchShaSod}. To establish \Cref{MR:decol_BA}, it is sufficient to prove the local laws \eqref{G_bound} and \eqref{G_bound_ave}, together with the quantum diffusion estimates \eqref{eq:diffu1}--\eqref{Meq:QdS2}, while the delocalization, QUE, and bulk universality can be derived as corollaries, as shown in \Cref{subsec:main}.



%\section{Main tools and proof of main results}




\subsection{Stochastic flow and loop hierarchy}\label{sec:tools}

The remainder of this paper is devoted to proving the local laws in \Cref{MR:locSC} and the quantum diffusion estimates in \Cref{MR:QuDiff} for both models. Before proceeding with the proofs, we first introduce some key tools and convenient notations for the remainder of this section, which will be utilized throughout the subsequent arguments. 
First, we introduce the flow framework, which is the same as that employed for 1D and 2D random band matrices \cite{YY_25,DYYY25} except for a dimension-dependent scaling $W^{-d}$. Consider the following matrix Brownian motion: 
\begin{align}\label{MBM}
	\dd (H_{t})_{xy}=\sqrt{S_{xy}}\dd (\boldsymbol{B}_{t})_{xy}, \ \ \forall x,y\in \ZL,\quad \text{where}\quad  H_{0}=\begin{cases}
	    0, \ &\text{for band matrix}\\
        \ilambda\Psi, \ &\text{for block Anderson}
	\end{cases}.
\end{align}
Here, $(\boldsymbol{B}_{t})_{xy}$ are independent complex Brownian motions up to the Hermitian symmetry \smash{$(\boldsymbol{B}_{t})_{xy}=\overline {(\boldsymbol{B}_{t})_{yx}}$}, i.e., \smash{$t^{-1/2}\boldsymbol{B}_t$} is an $N\times N$ GUE whose entries have zero mean and unit variance; $S=(S_{xy})$ is the variance matrix defined in \eqref{eq:variancematrix} or \eqref{bandcwV}. %Along this flow, we consider a time-dependent spectral parameter $z_{t}$ as defined below.
Following \cite{10.1214/19-ECP278, Sooster2019, DY}, we consider the Green's function of \( H_t \) with a carefully chosen time-dependent spectral parameter \( z_t \), whose dynamics are naturally renormalized at leading order.

%We start with the matrix Brownian motion
%\begin{align*}
%dH_{t,xy}=\sqrt{S_{xy}}dB_{t,xy}, \quad H_{0}=0, 
%\end{align*}
%where $B_{t,xy}$ are independent standard complex Brownian motions for all $x,y\in \Z_{WL}^{d}$ up to the constraint  $B_{yx} = \bar  B_{xy}$. 
 

%\begin{definition}\label{def_flow}
%For any $E \in \mathbb R$,  define $
%m^{(E)}:=\lim_{\epsilon\to 0+}m_{sc}(E+i \epsilon)$, and consider the following objects:
%\begin{align*}
%  z^{(E)}_t&=E+(1-t) m^{(E)}  ,\quad 0\le t\le 1 \\
%  G_{t}^{(E)}&:=(H_{t}-z^{(E)}_{t})^{-1}.
%\end{align*}
%Note that the imaginary part of $z_{t}^{(E)}$ is given by the following formula:
%\begin{align}\label{eta}
%\eta_t := \im  z^{(E)}_t =  (1-t)  \im m^{(E)}.
%\end{align}
%\end{definition}

\begin{definition}[Flow framework]\label{def_flow}
For any $E \in \mathbb R$ and $\ilambda>0$, we denote $m(E,\ilambda)\equiv m(E+\ii 0_+,\ilambda)$, as defined in \eqref{eq:defmzsc} and \eqref{self_m} for the random band matrix and block Anderson model, respectively. Based on this, we define the \emph{spectral parameter flow} $z_t$ by 
\be\label{eq:zt}
z_t(E,\ilambda) = E + (1-t) m(E,\ilambda),\quad \text{for}\ \ t\in [0, 1]. 
\ee
We refer to $E$ and $\ilambda$ as {\bf flow parameters}, which remain fixed throughout the flow. Let $z_t=E_t + \ii \eta_t$ with 
	\begin{align}\label{eta}
		E_t\equiv E_t(E,\ilambda)=E+(1-t)\re m(E,\ilambda),\quad \eta_t\equiv \eta_t(E,\ilambda) = (1-t)  \im m(E,\ilambda).
	\end{align}
Furthermore, we define the matrix $M(E,\ilambda)$, as in \eqref{eq:defMzsc} and \eqref{def_G0} for the random band matrix and block Anderson model, respectively, with $z$ and $m(z)$ replaced by $E$ and $m(E,\ilambda)$.  
Then, we denote the Green's function of $H_{t}$ as $G_t(z,\ilambda):=(H_t(\ilambda)-z)^{-1}$, and define the resolvent flow as 
	\begin{align}\label{self_Gt}
		G_{t;E,\ilambda}\equiv G_t(z_t(E,\ilambda),\ilambda):=(H_{t}(\ilambda)-z_t{(E,\ilambda)})^{-1}.
	\end{align}
% for given flow parameters $E$ and $\ilambda$. 
 % By It{\^o}'s formula, $G_{t;E,\ilambda}$ satisfies the following SDE:
	% \begin{align}\label{eq:SDE_Gt}
	% 	\dd G_{t;E,\ilambda}=-G_{t;E,\ilambda}(\dd H_{t})G_{t;E,\ilambda}+G_{t;E,\ilambda}\{\mathcal{S}[G_{t;E,\ilambda}]-m(E,\ilambda)\}G_{t;E,\ilambda}
	% 	\dd t.
	% \end{align}
\end{definition}

\begin{remark}
% \textcolor{blue}{Write a simple version  With $H_t$ as in \eqref{MBM} and $z_t$ as in \eqref{eq:zt}, define $M_t$ to be the solution of the matrix Dyson equation associated with $G_t=(H_t-z_t)^{-1}$:
% \[
% M_t := \bigl(g\,\Psi - z_t - t\,\mathcal S[M_t]\bigr)^{-1}.
% \]
% It is standard that $M_t$ provides the deterministic (self-consistent) limit of $G_t$. Moreover, since $\mathcal S[M(E,g)]=m(E,g)$, we have $M_t \equiv M(E,g)$ for all $t\in[0,1]$. Hence, along the flow $H_t$ with spectral parameter $z_t$, the deterministic limit of $G_t$ remains invariant.
% (the $\cal S$ is in defined in the next Def. but I suggest to move to the definition of $M$ and point out it the solution to Matrix Dyson equation}

To explain the choice of the flow $z_t$ in \eqref{eq:zt}, let $M_t(z,\ilambda)\in \cal M_N(\C)$ (where $\cal M_N(\C)$ denotes the set of $N\times N$ complex matrices) be the unique solution to the matrix Dyson equation 
\be\label{def_G0t}
M_t(z,\ilambda):= \p{H_0 -z- t \cal S[M_t(z,\ilambda)]}^{-1},
 \ee
such that $\im M_t$ is positive definite whenever \(z\in \C_+\). Here, the linear operator $\mathcal{S}$ is defined by
	\begin{align}\label{eq:opS}
		\mathcal{S}[X]_{xy}:=\delta_{xy}\sum_{w=1}^{N}S_{xw}X_{ww},\quad \text{for}\quad X\in \cal M_{N}(\C),
	\end{align}
where $S_{xy}$ is given by \eqref{eq:variancematrix} and \eqref{bandcwV} for the respective models. 
It is known that $M_t$ describes the deterministic limit of the Green's function $(H_t-z)^{-1}$. 
Moreover, since $\cal S[M(E,\ilambda)]=m(E,\ilambda)I_N$, one can check that \(M_t(z_t(E,\ilambda),\ilambda)\equiv M(E,\ilambda)\).  %solves \eqref{def_G0t}. 
In other words, the deterministic limit of $G_t$ remains invariant under the evolution.  

%In the proof, we will use a slightly different $\ilambda_0$ from $\ilambda$; see \Cref{zztE} below.    
% In other words, $m_t(z)$ is the Stieltjes transform of the semicircle law (with an extra scaling $\sqrt{t}$):
% \[m_t(z)=\int\frac{\rho_t(x)\dd x}{x-z}=\frac{-z+\sqrt{z^2-4t}}{2t},\quad \text{with}\quad \rho_t(x)=\frac{\sqrt{4t-x^2}}{2\pi t}\mathbf 1_{x\in [-2\sqrt{t},2\sqrt{t}]} \, .\]
% Then, the flow \eqref{eq:zt} is defined in a way such that
% \be\label{eq:mt_stay} m_t(z_t(E))\equiv m(E),\quad \forall t\in [0,1].\ee 
\end{remark}

% Next, we define the resolvent flow for the Green's function with $H_t$ given by \eqref{MBM} and $z_t$ given by \eqref{eq:zt}.

% \begin{definition}[Resolvent flow]\label{Def:stoch_flow}
% 	\textcolor{blue}{This whole definition was not serious used at all, we can remove it}Consider the matrix dynamics $H_t\equiv H_t(\ilambda)$ evolving according to \eqref{MBM}. Then, we denote Green's function of $H_{t}$ as $G_t(z,\ilambda):=(H_t(\ilambda)-z)^{-1}$, and define the resolvent flow as 
% 	\begin{align}\label{self_Gt}
% 		G_{t;E,\ilambda}\equiv G_t(z_t(E,\ilambda),\ilambda):=(H_{t}(\ilambda)-z_t{(E,\ilambda)})^{-1}
% 	\end{align}
%  for given flow parameters $E$ and $\ilambda$. By It{\^o}'s formula, $G_{t;E,\ilambda}$ satisfies the following SDE:
% 	\begin{align}\label{eq:SDE_Gt}
% 		\dd G_{t;E,\ilambda}=-G_{t;E,\ilambda}(\dd H_{t})G_{t;E,\ilambda}+G_{t;E,\ilambda}\{\mathcal{S}[G_{t;E,\ilambda}]-m(E,\ilambda)\}G_{t;E,\ilambda}
% 		\dd t.
% 	\end{align}
% Here, $\mathcal{S}:M_{N}(\C)\to M_{N}(\C)$ is a linear operator defined as
% 	\begin{align}\label{eq:opS}
% 		\mathcal{S}[X]_{xy}:=\delta_{xy}\sum_{y=1}^{N}S_{xy}X_{yy},\quad \text{for}\quad X\in M_{N}(\C).
% 	\end{align}
% where $S_{xy}$ is given by \eqref{eq:variancematrix} and \eqref{bandcwV} for the two models.
% \end{definition}


Given any target spectral parameter $z$, we are interested in the original resolvent $G(z)=(H-z)^{-1}$. For random band matrices, this can be achieved through the stochastic flow by carefully choosing the spectral parameter $E$. The corresponding flow for the block Anderson model will be presented later in \Cref{sec:ext-to-BA}. 


\begin{lemma}[Lemma 2.8 of \cite{YY_25}]\label{zztE}
Fix any $z\in \mathbb C_+$ with $\im z\in (0, 1]$ and $|\re z|\le 2-\kappa$. We choose 
	\be\label{eq:t0E0}t_0\equiv t_0(z)=|m(z)|^2=\frac{\im m(z)}{\im m(z)+ \im z},\quad E\equiv E(z)=-2\frac{\re m(z)}{|m(z)|}\, .\ee
	Then, for the random band matrix model, we have  
	\begin{equation}\label{eq:zztE}
		\sqrt{t_0}m(E)=m(z), \quad     z_{t_0}(E) =\sqrt{t_0} z  , \quad G(z) \stackrel{d}{=} \sqrt{t_0} G_{t_0;E} ,
	\end{equation} 
    where ``$\stackrel{d}{=}$'' means equality in distribution.
	% and the following equality in distribution (denoted by ``$\stackrel{d}{=}$''):
	% \begin{equation}\label{GtEGz}
	% 	G(z) \stackrel{d}{=} \sqrt{t_0} G_{t_0,E} .
	% \end{equation}
\end{lemma}



%and $|E_t|\le (2-\kappa)\sqrt{t}+\oo(1)$ for $|1-t|=\oo(1)$. 
%In other words, during the flow from $t=0$ to $t_0$, the imaginary part $\eta_t$ decreases from $\eta_0\asymp 1$ to $\eta_{t_0}\asymp 1-t_0 \asymp \im z $, while the spectral parameter $z_t$ stays within the bulk of the limiting spectrum of $H_t$, provided $t\ge 1 - \e$ for a small enough constant $\e>0$ depending on $\kappa$. 

 
In the main proofs for random band matrices, we will fix a target spectral parameter $z=\hat{E}+\ii \eta\in \mathbf D_{\kappa,\e}$ for an arbitrarily small constant $\e>0$ (recall \eqref{eq:spectral_domain}). Accordingly, we choose the parameters $t_0$ and $E$ as specified in \eqref{eq:t0E0}. The second identity in \eqref{eta} implies that $1-t\asymp \eta_t$ uniformly in $t\in [0,t_0]$, i.e., during the flow from $t=0$ to $t_0$, the imaginary part $\eta_t$ decreases from $\eta_0\asymp 1$ to $\eta_{t_0}\asymp 1-t_0 \asymp \eta \ge N^{-1+\e}$.  
For clarity, unless we want to emphasize their dependence on $E$ (or $g$), we will often omit this variable from various notations, such as $z_t(E,\ilambda)$, $E_t(E,\ilambda)$, $\eta_t(E,\ilambda)$, $m(E,\ilambda)$, $M(E,\ilambda)$, and most importantly, $G_{t;E,\ilambda}\equiv G_t$. Our focus will be on the dynamics of $G_{t}$ and the corresponding $G$-loops defined below. 


\begin{definition}[$G$-loop]\label{Def:G_loop}
For $\sigma\in \{+,-\}$, we denote 
 \begin{equation}\nonumber
  G_{t}(\sigma):=\begin{cases}
       (H_t-z_t)^{-1}, \ \ \text{if} \ \  \sigma=+,\\
        (H_t-\bar z_t)^{-1}, \ \ \text{if} \ \ \sigma=-.
   \end{cases}
 \end{equation}
In other words, we let $G_{t}(+)\equiv G_{t}$ and $G_{t}(-)\equiv G_{t}^*$. %Recall that $E_{a}$ ($a\in \Zn$) denotes the rescaled block identity matrix: \((E_{a})_{xy}= W^{-d}\delta_{xy}\cdot \mathbf 1_{x\in [a]} \).
% Denote by $\id_{a}$ and $E_{a}$ ($a\in \Zn$) the block identity and rescaled block identity matrices, respectively:
% \be\label{def:Ia} (\id_{a})_{xy}= \delta_{xy}\cdot \mathbf 1_{x\in [a]} , \quad E_{a}=W^{-d}\id_{a}.\ee
For any $n\in \N$, fix indices $\bsig=(\sigma_1, \ldots, \sigma_n)\in \{+,-\}^n$ and $\ba=(a_1, \ldots, a_n)\in (\Zn)^n$. We define the corresponding {\bf $n$-$G$-loop} by 
\begin{equation}\label{Eq:defGLoop}
    {\cal L}^{(n)}_{t, \bsig, \ba}= \tr \p{\prod_{i=1}^n \left(G_{t}(\sigma_i) E_{a_i}\right) }, \quad \text{where} \quad (E_{a_i})_{xy} = W^{-d}\mathbf 1(x=y\in [a_i]).
\end{equation}
%where we recall that $\avg{\cdot}$ denotes the trace operator. 
Sometimes, we will also refer to $G$-loops as $\cL$-loops. Furthermore, we denote 
\begin{equation}\label{def_mtzk}
m (\sigma ):= \begin{cases}
    m(E,\ilambda),  &\text{if} \ \ \sigma  =+ \\
    \bar m(E,\ilambda),  &\text{if} \ \ \sigma = -
\end{cases}, \quad M (\sigma ):= \begin{cases}
     M(E,\ilambda) ,  &\text{if} \ \ \sigma  =+ \\
     M(E,\ilambda)^*,  &\text{if} \ \ \sigma = -
 \end{cases}.
\end{equation}
Finally, we define the \emph{centered resolvent} $\Gc$ as
\begin{equation}\label{Eq:defwtG}
 \Gc_t(\sigma) := G_t(\sigma) - M(\sigma),\quad \forall t\in[0,1], \ \sig\in \{+,-\} .
\end{equation}
(Note that $\Gc_t(\sigma)$ was denoted by $\widetilde{G}_t(\sigma)$ in the preceding papers \cite{YY_25, DYYY25}; here we adopt a different notation for clarity.) 
For any $a\in \Zn$, we refer to $\tr[\Gc_t(\sig) E_a]$ as a \emph{light-weight}.
\end{definition}





%For the rest of the paper, we will fix $E$ such that $|E| \le 2-\kappa$ and consider only  $G_t^{(E)}$.  We now define the aforementioned $G$-loops. 
% \begin{definition}\label{Def:G_loop}
%  For any fixed $E\in\R$, consider the following resolvent as a function of $+,-$ signs:
%  \begin{equation}\nonumber
%   G^{(E)}_{t }(\sigma):=  G^{(E)}_{t,\alpha} =
%    \begin{cases}
%        (H_t-z_t)^{-1}, \quad \sigma=+\\
%         (H_t-\bar z_t)^{-1}, \quad \sigma=-
%    \end{cases}.
%  \end{equation}
% By definition, $G^{(E)}_{t,+} =\left(G^{(E)}_{t,-}\right)^\dagger$. 
% For simplicity,  we sometimes write 
% $G_t=G^{(E)}_{t,+}$. Define the block-identity matrix $E_a$ as follows.
% \[
% \left(E_{a}\right)_{xy} = \delta_{xy} \cdot W^{-d} \cdot \mathbf{1}\left(x\in \mathcal{I}_a\right), \quad x,y\in \Z_{WL}^d,\quad a\in \Z_L^d.
% \]

% \medskip

% Now, fix indices
% $$
% \bsig=(\sigma_1, \sigma_2, \cdots \sigma_n),\quad 
% \boldsymbol{a}=(a_1, a_2, \cdots a_n)
% ,\quad \sigma_i \in \{+,-\}, \quad
% a_i\in \mathbb Z_L^{d}, \quad 1\le i\le n.
% $$
% We define the corresponding ``$n$-$G$ loop'' to be the following complex number: 
% \begin{equation}\label{Eq:defGLoop}
%     {\cal L}_{t, \bsig, \mathbf{a}}=\left\langle \prod_{i=1}^n G_t(\sigma_i) E_{a_i}\right\rangle,\quad 
% \text{where} \quad \left\langle A\right\rangle=\tr A.
% \end{equation}
% We also define 
% \begin{equation}\label{def_mtzk}
% m (\sigma ):=m^{(E)}(\sigma ) = \begin{cases}
%      m^{(E)}, & \sigma  =+ \\
%      \overline{m^{(E)}} , & \sigma = -
% \end{cases},
% \end{equation}
% %and $\widetilde{G}_t $ as follows, 
% and
% \begin{equation}\label{Eq:defwtG}
%  \Gc_t(\sigma_k) := G_t(\sigma_k) - m (\sigma_k) .
% \end{equation}
% \end{definition}

%Note that the function $\Gc_t(\sigma)$ is the same function that was denoted $\widetilde{G}_t(\sigma)$ in the preceding $d=1,2$ papers \cite{YY_25, DYYY25}.

%With these notations, we can rewrite the quantity 
%$\tr G E_a G^\dagger  E_b$ in Theorem \ref{MR:QDiff} as follows:
%$$
%\tr G E_a G^\dagger  E_b=t \cdot \tr (G_{t,+})\cdot E_a\cdot (G_{t,-})\cdot  E_b=t\cdot {\cal L}_{t, (+,-),(a,b)}.
%$$

To describe the loop hierarchy for the $G$-loops, we introduce the following operations from \cite{YY_25}. (For a graphical illustration of these operations, see also the diagrams in \cite[Definition 2.10]{YY_25}.)


\begin{definition}\label{Def:oper_loop}
For any fixed $n\in \N$, take an $n$-loop of the form \eqref{Eq:defGLoop}.
% \begin{equation}\label{taahCal}
%    {\cal L}^{(n)}_{t, \bsig, \mathbf{a}} = \left\langle \prod_{i=1}^n \p{G_t(\sigma_i) E_{a_i}} \right\rangle, \quad \forall \sigma_i \in \{+, -\}, \ a_i\in \Zn, \ 1 \le i \le n .
% \end{equation}

\medskip
 
 \noindent 
 \emph{1.} For $k \in \qqq{n}$ and $a\in \Zn$, we define a ``cut-and-glue'' operator ${\cut}^{(a)}_{k}$ as follows: ${\cut}^{(a)}_{k} \circ {\cal L}^{(n)}_{t, \bsig, \mathbf{a}}$ is defined to be the loop obtained by replacing $G_t(\sigma_k)$ with $G_{t}(\sigma_k) E_a G_{t}(\sigma_k)$. In other words, it cuts the $k$-th edge $G_t(\sigma_k)$ and glues the two new ends with $E_{a}$ to get a new loop that is one unit longer. This operator can also be considered as an operator on $(\bsig,\ba)$, that is, 
\begin{align*}{\cut}^{(a)}_{k} (\bsig, \ba) =  \big( & (\sigma_1,\ldots, \sigma_{k-1}, \sigma_k,\sigma_k ,\sigma_{k+1},\ldots, \sigma_n ),\  ( a_1 ,\ldots,  a_{k-1} ,  a , a_k , a_{k+1} ,\ldots,  a_n  )\big).\end{align*}
Hence, we will sometimes write ${\cut}^{(a)}_{k} \circ {\cal L}^{(n)}_{t, \bsig, \ba}\equiv {\cal L}^{(n+1)}_{t, \;  {\cut}^{(a)}_{k} (\bsig, \ba)}.$ 

% \be\label{eq:cut1}
%     {\cut}^{[a]}_{k} \circ {\cal L}^{(n)}_{t, \bsig, \ba}:= \left \langle \prod_{i<k}  \left(G_{t}(\sigma_i) E_{[a_i]}\right)\left( G_{t}(\sigma_k) E_{[a]} G_{t}(\sigma_k)E_{[a_k]}\right)\prod_{i>k}  \left(G_{t}(\sigma_i) E_{[a_i]}\right)\right\rangle.
% \ee
%For $1 \le k \le n$, we define a ``cut-and-glue'' operator ${\cal G}^{(a)}_{k}$ as follows. Its action on a loop ${\cal L}_{t,\bsig,\mathbf{a}}$ is denoted by ${\cal G}^{(a)}_{k} \circ {\cal L}_{t, \bsig, \mathbf{a}}$, and it is defined to be the loop obtained by replacing $G_t(\sigma_k)$ by $G_{t}(\sigma_k) E_a G_{t}(\sigma_k)$.
 

 \noindent 
 \emph{2.}  For $k < l \in \qqq{n}$, we define two other types of ``cut-and-glue'' operators---${\cutL}^{(a)}_{k,l}$ from the left (``L'') of $k$, and \smash{${\cutR}^{(a)}_{k,l}$} from the right (``R'') of $k$---as follows. 
 In other words, these operators cut the $k$-th and $l$-th edges $G_t(\sigma_k)$ and $G_t(\sigma_l)$, and create two chains: the left chain to the vertex $a_k$ is of length $(n+k-l+1)$ and contains the vertex $a_n$, while the right chain to the vertex $a_k$ is of length $(l-k+1)$ and does not contain the vertex $a_n$.  
Then, \smash{${\cutL}^{(a)}_{k,l}$ (resp.~${\cutR}^{(a)}_{k,l}$)} gives an $(n+k-l+1)$-loop (resp.~$(l-k+1)$-loop) obtained by gluing the left chain (resp.~right chain) at the new vertex $a$.
Again, we can also consider the two operators to be defined on the indices $(\bsig,\ba)$: 
\begin{align*}
    &{\cutL}^{(a)}_{k,l} (\bsig, \ba) = \left((\sigma_1,\ldots, \sigma_k,\sigma_l ,\ldots, \sigma_n ),(a_1,\ldots, a_{k-1}, a,a_l,\ldots, a_n )\right),\\
    &{\cutR}^{(a)}_{k,l} (\bsig, \ba) = \left((\sigma_k,\ldots, \sigma_l),(a_k,\ldots, a_{l-1}, a)\right).
\end{align*}
Hence, we will sometimes write 
    ${\cutL}^{(a)}_{k,l} \circ {\cal L}^{(n)}_{t, \bsig, \ba}\equiv 
     {\cal L}^{(n+k-l+1)}_{t,   {\cutL}^{(a)}_{k,l} (\bsig, \ba)}$ and $ {\cutR}^{(a)}_{k,l} \circ {\cal L}^{(n)}_{t, \bsig, \ba}\equiv 
     {\cal L}^{(l-k+1)}_{t,  {\cutR}^{(a)}_{k,l} (\bsig, \ba)}.$
\end{definition}

    



 
%The following result gives the evolution of $G$-loops along the flow, which  It is derived via the It\^o formula. Throughout this paper, we let $\mathbf{a}=(a_{1},\ldots,a_{n})$ be a vector of indices in $\Z_{L}^{d}$, and we let $a\in\Z_{L}^{d}$ be a single index that is separate from $\mathbf{a}$. Finally, for convenience, we will use the notation $\partial_{(x,y)}:=\partial_{H_t(x,y)}$.

For $x,y\in \ZL$, we abbreviate $\partial_{xy}:=\partial_{(H_t)_{xy}}$. By It\^o's formula, we can derive the following SDE satisfied by the $G$-loops, called \emph{loop hierarchy}; see Lemma 2.11 of \cite{YY_25}.  

 
\begin{lemma}[Loop hierarchy] \label{lem:SE_basic}
An $n$-$G$-loop satisfies the following SDE, called the \emph{loop hierarchy}:
\begin{align}\label{eq:mainStoflow}
\dd\mathcal{L}^{(n)}_{t, \bsig, \ba} = \dd\mathcal{E}^{M,(n)}_{t, \bsig, \mathbf{a}} + \mathcal{E}^{\Gc,(n)}_{t, \bsig, \ba}\dd t +
   W^{d} \sum_{1 \le k < l \le n} \sum_{a, b} \left( {\cutL}^{(a)}_{k, l} \circ \mathcal{L}^{(n)}_{t, \bsig, \ba} \right) S^{(\sB)}_{ab} \left( {\cutR}^{(b)}_{k, l} \circ \mathcal{L}^{(n)}_{t, \bsig, \ba} \right) \dd t, 
\end{align}
where $S^{(\sB)}$ denotes $S^{(\sB)}(\ilambda)$ in \eqref{eq:variancematrix} (for the random band matrix model) or $S^{(\sB)}(0)=I_{L^d}$ (for the block Anderson model). Moreover, the martingale term \smash{$\dd\mathcal{E}^{M,(n)}_{t, \bsig, \ba}$} and the light-weight term %(as defined below \eqref{eq:lweightnota}) term 
\smash{$\mathcal{E}^{\Gc,(n)}_{t, \bsig, \ba}$} are defined by 
 \begin{align} \label{def_Edif}
\dd\mathcal{E}^{M,(n)}_{t, \bsig, \ba} :  = & \sum_{x,y\in \ZL} 
  \left( \partial_{xy}  {\cal L}^{(n)}_{t, \bsig, \ba}  \right)
 \cdot \sqrt{S _{xy}}
  \left(\dd \boldsymbol{B}_t\right)_{xy}, \\\label{def_EwtG}
\mathcal{E}^{\Gc,(n)}_{t, \bsig, \ba}:= &  {W}^d \sum_{k=1}^n \sum_{a, b\in \Zn} \; 
 \tr\p{ \Gc_t(\sigma_k) E_{a} }
   S^{(\sB)}_{ab} 
  \left( {\cut}^{(b)}_{k} \circ {\cal L}^{(n)}_{t, \bsig, \ba} \right) . 
\end{align}
We emphasize that the superscript $(n)$ indicates the length of the $G$-loop on the left-hand side of the equation. For clarity and conciseness, we may omit this superscript when its value is clear from the context. %and we will sometimes omit it from our notations when the value of $n$ is clear from the context. 
%  Above, we used the notation 
%   \begin{align} \label{def_Edif}
% \mathcal{E}^{(M)}_{t, \bsig, \mathbf{a}} :  = & \sum_{\alpha=(x,y)} 
%   \left( \partial_{  \alpha} \; {\cal L}_{t, \bsig, \mathbf{a}}  \right)
%  \cdot \left(S _\alpha\right)^{1/2} \cdot
%   \mathrm{d} \left({B}_t\right)_{\alpha} \\\label{def_EwtG}
%   \mathcal{E}^{(\Gc)}_{t, \bsig, \mathbf{a}}: = &   W^{d}\cdot \sum_{1 \le k \le n} \sum_{a, b} \; 
%  \left \langle \Gc_t(\sigma_k) E_a \right\rangle
%   \cdot S^{(B)}_{ab} \cdot
%   \left( {\cal G}^{(b)}_{k} \circ {\cal L}_{t, \bsig, \mathbf{a}} \right) dt ,\quad \quad \Gc_t = G_t - m.
% \end{align}
  \end{lemma}
  
%The factor of $W^{d}$ in ${\cal E}^{(\Gc)}$ comes from the normalization that $E_{a}$ has a factor of $W^{-d}$. 

Note that the right-hand side (RHS) of equation \eqref{eq:mainStoflow} involves $G$-loops of length larger than $n$, and hence represents a ``hierarchy'' rather than a ``self-consistent equation'' for the $G$-loops. This loop hierarchy is well-approximated by the $\cK$-loops, defined as follows. 


\begin{definition}[Tree approximation]\label{Def_Ktza}
We define the $\cK$-loop of length 1 as: 
$${\cal K}^{(1)}_{t,\sigma,a}=m(\sigma),\quad \forall t\in [0,1),\; \;\sigma\in \{+,-\},\;\; a\in \Zn.$$ 
For $n\ge 2$, we define the function \smash{${\cal K}^{(n)}_{t, \bsig, \ba}$} (of $t\in[0,1)$, $\bsig\in \{+,-\}^n$, and \smash{$\ba\in (\Zn)^n$}) to be the unique solution to the following system of equations, referred to as the \emph{\bf convolution tree equations}: 
 \begin{align}\label{pro_dyncalK}
       \partial_t{\cal K}^{(n)}_{t, \bsig, \ba} 
       =  
       W^d \sum_{1\le k < l \le n} \sum_{a, b} \left( \cutL^{(a)}_{k, l} \circ \mathcal{K}^{(n)}_{t, \bsig, \ba} \right) S^{(\sB)}_{ab} \left( \cutR^{(b)}_{k, l} \circ \mathcal{K}^{(n)}_{t, \bsig, \ba} \right) ,
    \end{align}
where the operators $\cutL$ and $\cutR$ act on ${\cal K}^{(n)}_{t, \bsig, \ba}$ through the actions on indices: 
  \be\label{calGonIND}
     {\cutL}^{(a)}_{k,l}  \circ {\cal K}^{(n)}_{t, \bsig, \ba} := {\cal K}^{(n+k-l+1)}_{t,  {\cutL}^{(a)}_{k,l}  (\bsig, \ba)} , \ \  {\cutR}^{(b)}_{k,l}  \circ {\cal K}^{(n)}_{t, \bsig, \ba} := {\cal K}^{(l-k+1)}_{t,  {\cutR}^{(b)}_{k,l}  (\bsig, \ba)}.  
       \ee
We impose the following initial condition at $t=0$: 
\be\label{eq:initial_K}
{\cal K}^{(k)}_{0, \bsig, \ba} =  {\cal M}^{(k)}_{\bsig, \ba} ,\quad \forall k\in \N, \ \bsig\in \{+,-\}^k,\ \ba\in (\Zn)^k,\ee
where ${\cal M}^{(k)}_{\bsig, \ba}$ is a $k$-$M$-loop defined as
\be\label{eq:KMloop} {\cal M}^{(k)}_{\bsig, \ba}:=\tr\p{ \prod_{i=1}^k \left(M(\sigma_i) E_{a_i}\right) } .\ee
(Note that this $M$-loop can be simplified as \smash{${\cal M}^{(k)}_{\bsig, \ba}= W^{-(k-1)d}\prod_{i=1}^k m(\sig_i) \mathbf 1(a_1=\cdots=a_k)$} for the random band matrix model.) 
% \be\label{eq:initial_K}
% {\cal K}^{(k)}_{0, \bsig, \ba} = W^{-(k-1)d}\prod_{i=1}^k m(\sig_i)\cdot \mathbf 1(a_1=\cdots=a_k) ,\quad \forall k\in \N, \ \bsig\in \{+,-\}^k,\ \ba\in (\Zn)^k.\ee
% where ${\cal M}^{(k)}_{\bsig, \ba}$ is defined as
% \be\label{eq:KMloop} {\cal M}^{(k)}_{\bsig, \ba}:=\left \langle \prod_{i=1}^k \left(M(\sigma_i) E_{[a_i]}\right)\right\rangle = W^{-(k-1)d}\prod_{i=1}^k m(\sig_i) \mathbf 1([a_1]=\cdots=[a_k]).\ee
%In \eqref{pro_dyncalK}, \textcolor{blue}{this sentense should be right below \eqref{pro_dyncalK}}
We call \smash{${\cal K}^{(n)}_{t, \bsig, \ba}$} an $n$-$\cK$-loop. Moreover, we refer to \smash{${\cal K}^{(n)}_{t, \bsig, \ba}$} as providing a \emph{\bf tree approximation} of the $G$-loop \smash{$\mathcal{L}^{(n)}_{t, \bsig, \ba}$}, since it admits an explicit tree representation, as shown in \cite{YY_25,RBSO1D} (see also \Cref{Sec:CalK} for the detailed construction).
%cite \cite{DYYY25_d3} here
\end{definition}    



%\begin{lemma}[Ward's identity]\label{lem-Ward}
Given any Hermitian matrix $\cal A$, define its resolvent as $R(z):=(\cal A-z)^{-1}$ for $z= E+ \ii \eta\in \C_+$. Then, with the algebraic identity $R-R^*=2\ii \eta RR^*=2\ii \eta R^*R$, we get the well-known Ward's identity:
    \be\label{eq_Ward0}
    \begin{split}
\sum_x \overline {R_{xy'}}  R_{xy} = \frac{1}{2\ii \eta}\p{R_{y'y}-\overline{R_{yy'}}},\quad
\sum_x \overline {R_{y'x}}  R_{yx} = \frac{1}{2\ii \eta}\p{R_{yy'}-\overline{R_{y'y}}}.
\end{split}
\ee
As a special case, if $y=y'$, we have
\be\label{eq_Ward}
\sum_x |R_{xy}( z)|^2 =\sum_x |R_{yx}( z)|^2 = {\im R_{yy}(z) }/{ \eta}.
\ee
% \end{lemma}
% \begin{proof}
% These identities follow directly from the algebraic identity $R-R^*=2\ii \eta RR^*=2\ii \eta R^*R$.
% \end{proof}
Applying \eqref{eq_Ward0} to $G_t$, we can show that the $G$-loops satisfy the following identity \eqref{WI_calL}, which we also refer to as a ``Ward's identity''. In \cite{YY_25,RBSO1D}, it is shown that a similar Ward's identity \eqref{WI_calK} holds for the $\cal K$-loops. 

\begin{lemma}[Ward's identity for $\cL$-loops and $\cK$-loops]\label{lem_WI_K} 
Given $\bsig\in\{+,-\}^n$ with $n\ge 2$ and $\sigma_1=-\sig_{n}$, we have the following identities, which are called Ward's identities at the vertex $a_n$: 
\begin{align}\label{WI_calL}
&\sum_{a_n}{\cal L}^{(n)}_{t, \bsig, \ba}=
\frac{1}{2\ii W^d\eta_t}\left( {\cal L}^{(n-1)}_{t,  \wh\bsig^{(+,n)}, \wh\ba^{(n)}}- {\cal L}^{(n-1)}_{t,  \wh\bsig^{(-,n)} , \wh\ba^{(n)}}\right),  \\
\label{WI_calK}
& \sum_{a_n}{\cal K}^{(n)}_{t, \bsig, \ba}=
\frac{1}{2\ii W^d\eta_t}\left( {\cal K}^{(n-1)}_{t,  \wh\bsig^{(+,n)}, \wh\ba^{(n)}}- {\cal K}^{(n-1)}_{t,  \wh\bsig^{(-,n)}, \wh\ba^{(n)}}\right) , 
\end{align}
where $\eta_t$ is defined in \eqref{eta}, $\wh\bsig^{(\pm,n)}$ is obtained by removing $\sigma_n$ from $\bsig$ and replacing $\sigma_1$ with $\pm$, i.e., \smash{$\wh\bsig^{(\pm,n)}:=(\pm, \sigma_2, \cdots \sigma_{n-1})$}, and \smash{$\wh\ba^{(n)}$} is obtained by removing $a_n$ from $\ba$, i.e., \smash{$\wh\ba^{(n)}:=(a_1, a_2,\cdots, a_{n-1})$}. 
\end{lemma}
\begin{proof}
For the random band matrix, this corresponds to Lemma 3.6 in \cite{YY_25}, while for the block Anderson model, it corresponds to Lemma 3.17 in \cite{RBSO1D}. 
\end{proof}


%cite \cite{DYYY25_d3} here
In \Cref{Sec:CalK}, we will present a tree representation formula for the $\cK$-loops, originally discovered in \cite{YY_25} for the random band matrix model and in \cite{RBSO1D} for the block Anderson model. Using this tree representation, we can establish the following upper bound \eqref{eq:bcal_k} for $\cK$-loops. The proof, being similar to those in \cite{YY_25,RBSO1D}, is deferred to \Cref{Sec:CalK}.


\begin{lemma}[Upper bounds on $\cK$-loops]\label{ML:Kbound}
Fix any dimension $d\ge 3$.	Then, for every $n\in \N$ and $t\in [0,1)$, the $\cK$-loops satisfy the upper bound 
	\begin{equation}\label{eq:bcal_k}
		\max_{\bsig\in\{+,-\}^n}\max_{\ba\in(\Zn)^n} \big| {\cal K}^{(n)}_{t, \bsig, \ba} \big| \prec \left(W^{-d}B_{t,0}\right)^{n-1},      
	\end{equation}
where $B_{t,0}=(\ilambda^2+|1-t|)^{-1}+(L^d|1-t|)^{-1}$ is defined precisely in \Cref{defi:ofB} below. 
\end{lemma}


\subsection{Propagators} %\texorpdfstring{$\Theta_\xi^{(\sB)}$}{Theta} and an estimate on \texorpdfstring{${\cal K}$}{ K}}

The quantum diffusion behavior of the resolvents is governed by the $\Theta$-propagators, defined as follows.
 
  
\begin{definition}\label{def_Theta}
Given $\sigma_1,\sigma_2\in \{+,-\}$, define an $L^d\times L^d$ matrix $M^{(\sigma_1,\sigma_2)}$ as a rescaled 2-$M$-loop:
\be\label{eq:Msig} M^{(\sigma_1,\sigma_2)}_{ab}:= W^d\tr\p{ M(\sigma_1) E_{a}M(\sigma_2) E_{b}},\quad \forall a,b\in \Zn.\ee
For the random band matrix model, we have $M^{(\sigma_1,\sigma_2)}=m(\sigma_1)m(\sigma_2)I_{L^d}$; for the block Anderson model, we have \smash{$M^{(\sigma_1,\sigma_2)}_{ab}=M^{(\sB)}_{ba}(\sigma_1)M^{(\sB)}_{ab}(\sigma_2)$}, where $M^{(\sB)}$ is defined in \eqref{def_G0}, with $M^{(\sB)}(+)\equiv M^{(\sB)}$ and $M^{(\sB)}(-)\equiv (M^{(\sB)})^*$. 
For $t\in[0,1)$ and $\sigma_1,\sigma_2\in \{+,-\}$, the $\Theta$-propagator \smash{$\Theta_t^{(\sigma_1,\sigma_2)}$} is an $L^d\times L^d$ matrix defined as:
\begin{equation}\label{def_Thxi}
    \Theta_{t}^{(\sigma_1,\sigma_2)} := \p{1 - t M^{(\sigma_1,\sigma_2)}S^{(\sB)}}^{-1}, 
\end{equation}
where $S^{(\sB)}$ is given by $S^{(\sB)}(\ilambda)$ in \eqref{eq:variancematrix} or $S^{(\sB)}(0)=I_{L^d}$, depending on whether we consider the random band matrix model or the block Anderson model.  
%\eqref{eq;smatrixB} or \eqref{eq;sproductBA}. 
%it can be simplified by using $S^{(\sB)}(0)=I_{L^d}$; . 
% allowing us to simplify \eqref{def_Thxi} as:
% \be\label{eq:Theta_WO}\Theta_{t}^{(\sigma_1,\sigma_2)}= \p{1-t m(\sigma_1)m(\sigma_2)S^{(\sB)}}^{-1}.\ee 
% Given $\xi \in \mathbb{C}$ with $|\xi| < 1$, we define the {\bf propagator} $\Theta^{(\sB)}_{\xi}$ to be a $L^{d}\times L^{d}$ matrix with entries {\cor here}
% \begin{equation}\label{def_Thxi}
%     \Theta^{(\sB)}_\xi := \p{1 - \xi \cdot S^{(\sB)}}^{-1} .
% \end{equation}
% Here, the superscript $(\sB)$ indicates a block-level matrix, whose indices are parameterized by $\Zn$. 
We denote the entries of \smash{$\Theta_{t}^{(\sigma_1,\sigma_2)}$} as \smash{$\Theta_{t}^{(\sig_1,\sig_2)}(a,b)$ or $\Theta_{t,ab}^{(\sig_1,\sig_2)}$}. We further define the zero-mode-removed propagator \smash{$\zTheta^{(\sig_1,\sig_2)}_{t}$} by 
\begin{equation}\label{def_Thxi0}
\zTheta^{(\sig_1,\sig_2)}_{t}(a,b):=\Theta^{(\sig_1,\sig_2)}_{t}(a,b) -{L^{-2d}}\sum_{a',b'}\Theta^{(\sig_1,\sig_2)}_{t}(a',b').
\end{equation}
\end{definition}


%For convenience, we often omit the superscript $(B)$. We remind the readers that $\Theta$  was often  used to denote the \(N \times N\) matrix  \((1 - \xi S)^{-1}\) in the literature.   In this paper,  $\Theta_\xi= \Theta_\xi^{(B)}$.  We note that
% \begin{equation}\label{deri_Thxi}
%     \partial_\xi \Theta^{(B)}_\xi = \Theta^{(B)}_\xi \cdot S^{(B)} \cdot \Theta^{(B)}_\xi.
% \end{equation}


Note that for the random band matrix model, the $\Theta$-propagators $\Theta^{(\sig,\sig')}_{t_0}$ for $\sig,\sig'\in\{+,-\}$ reduce to the $\Theta$-matrices defined in \eqref{def:Theta} under the relation \eqref{eq:zztE}. 
Similarly, for the block Anderson model, there is a corresponding reduction to the definition in \eqref{def:Theta_BA} at $t=t_0$ under the relation \eqref{eq:zztE_BA} below.
  

To capture the decay profile of the $\Theta$-propagators, we introduce the following control parameter $B_{t,K}$ in \eqref{eq_B_param}. It is an order parameter that will appear in many estimates throughout the proof. 

\begin{definition}[Definition of $B$]\label{defi:ofB}
For any $t\in [0,1)$, define 
\begin{equation}\label{eq_B_param}
B_{t, K}:=\frac{(\ilambda^{2}+|1-t|)^{-1}}{(K+1)^{d-2}}+\frac{1}{L^d|1-t|},\quad \forall K\ge 0.
\end{equation}
Since $\eta_t\asymp 1-t$ by \eqref{eta}, the parameter $\cal B_{\eta_t,K}$ defined in \eqref{eq:calBetaK} satisfies 
\be\label{eq:BtBt}
\cal B_{\eta_t,K} \asymp W^{-d}B_{t,(K/W)}.
\ee
\end{definition}

We now summarize some fundamental properties of the $\Theta$-propagators that will be used extensively in the main proof. These properties have essentially been established in previous works \cite{DYYY25, yang2024Del, RBSO1D}. For the reader's convenience, we briefly outline the proof of \Cref{lem_propTH} in \Cref{sec:pf_propTH}.
%cite \cite{DYYY25_d3} here
 
\begin{lemma}\label{lem_propTH}
For any $t\in [0,1)$, define 
\be\label{eq:ellt}
 \ell_t:= \min\p{ \max\big(\ilambda |1-t|^{-\frac 1 2},1\big)\; ,\; L} \, .
\ee
% Suppose that  $\xi\in \mathbb C$ satisfies and $|\xi|< 1$.  Define
% $$
%  \hat{\ell}(\xi) := \min\left((|1-\xi|)^{-1/2},\; L\right),\quad \xi\in \mathbb C 
% $$
%(Since $\eta_t\asymp 1-t$ by \eqref{eta} and \eqref{eq:zztE2}, the parameter $\ell(\eta_t)$ defined in \eqref{def_ellz} is of the same order as $\ell_t$.)
For any $\sig_1,\sig_2\in\{+,-\}$, $\Theta^{(\sig_1,\sig_2)}_t$ defined in \Cref{def_Theta} satisfies the following properties:
\begin{enumerate}
\item {\bf Symmetry}: $\Theta^{(\sig_2,\sig_1)}_t=(\Theta^{(\sig_1,\sig_2)}_t)^\top= \Theta^{(\sig_1,\sig_2)}_t$.
    
\item {\bf Translation invariance}: For any $a,b,r\in \Zn$, we have $\Theta^{(\sig_1,\sig_2)}_t(a+r,b+r)= \Theta^{(\sig_1,\sig_2)}_t(a,b).$
% \be\label{symmetry}
% \Theta^{(\sB)}_{\xi}([x]+[a],[y]+[a])= \Theta^{(\sB)}_{\xi}([x],[y]).
% \ee

 \item {\bf Commutativity}: We have $\big[S^{(\sB)}, \Theta^{(\sig_1,\sig_2)}_t\big]=\big[\Theta^{(\sig_1,\sig_2)}_t,\Theta^{(\sig_1,\sig_2)}_{t'}\big]=0$ for all $t\ne t'$.

 \item {\bf $(\infty\to\infty)$-norm}: For any $a,b\in \Zn$, we have $|\Theta_{t,ab}^{(\sig_1,\sig_2)}|\le \Theta_{t,ab}^{(+,-)}$. Moreover, 
 \be\label{eq:THETAinftinf} \|\Theta_{t}^{(\sig_1,\sig_2)}\|_{\infty\to\infty} = \max_a \sum_{b}\big|\Theta_{t,ab}^{(\sig_1,\sig_2)} \big|\le \max_a \sum_{b}\Theta_{t,ab}^{(+,-)}=\frac{1}{1-t}.\ee

\item {\bf Polynomial and exponential decay}: For all $\sig_1,\sig_2\in \{+,-\}$, there exist constants $c_d,C_d>0$ (depending on $d$) such that the following bound holds: 
\begin{equation}\label{prop:ThfadC}   
\big|\Theta^{(\sig_1,\sig_2)}_{t}(0,a)\big|\le C_d B_{t, |a|}\cdot {e^{-c_d |a|/ {\ell}_t}},\quad \forall a\in \Zn .  
\end{equation}
Furthermore, when $\sig_1=\sig_2$, we have a much stronger exponential decay: there exist constants $c_\kappa,C_\kappa>0$ (depending on $d$ and $\kappa$) such that 
\begin{equation}\label{prop:ThfadC_short}  \big|\Theta^{(\sig_1,\sig_2)}_{t}(0,a)\big|\le C_\kappa \left(1_{a=0} + \ilambda^2 e^{-c_{\kappa}|a|}\right),\quad \forall a\in \Zn.
\end{equation}
In the setting of the block Anderson model, these constants, including $c_d$ and $C_d$ in \eqref{prop:ThfadC}, may also depend on $\kappa$ and on $\fd$ in \eqref{eq:WO}.

\item {\bf First-order difference}: Fix any constant $c_0\in(0,1)$. The following estimate holds for all $a, r\in \Zn$ satisfying $|r|\le c_0 |a|$ (the implicit constants in $\prec$ may depend on $c_0$):  
\begin{equation}\label{prop:BD1}     
    \left| \Theta^{(\sig_1,\sig_2)}_{t}(0, a+r)-\Theta^{(\sig_1,\sig_2)}_{t}(0, a)\right|\prec \frac{1}{\ilambda^2+|1-t|}\frac{|r|}{(|a|+1)^{d-1}}.
     \end{equation}


\item {\bf Second-order difference}: Fix any constant $c_0\in(0,1)$. The following estimate holds for all $a,r\in \Zn$ satisfying $|r|\le c_0|a|$ (the implicit constants in $\prec$ may depend on $c_0$): 
\begin{equation}\label{prop:BD2}
\left|\Theta^{(\sig_1,\sig_2)}_{t} (0,a+r) + \Theta^{(\sig_1,\sig_2)}_{t} (0,a-r)-  2\Theta^{(\sig_1,\sig_2)}_{t} (0,a) \right|
\prec \frac{1}{\ilambda^2+|1-t|}\frac{|r|^2}{(|a|+1)^{d}} .
 \end{equation}

\item {\bf Propagator without zero mode.} The following estimate holds for all $a \in \Zn$: 
\be\label{prop:ThfadC0} 
\big|\zTheta^{(\sig_1,\sig_2)}_{t} (0,a)\big|\prec \frac{1}{\ilambda^2+|1-t|}\frac{1}{(|a|+1)^{d-2}}. 
\ee
\end{enumerate}

\end{lemma}


\begin{example} 
As shown in \cite{YY_25,RBSO1D}, the $2$-$\cK$ and $3$-$\cK$ loops are given by %(recall \eqref{eq:KMloop} and \eqref{eq:Msig}) 
\begin{align}\label{Kn2sol}\cK^{(2)}_{t,\bsig,\ba}&= \sum_{b}\Theta_{t}^{(\sig_1,\sig_2)}(a_1,b)\cal M^{(2)}_{\bsig,(b,a_2)}=W^{-d} \left(\Theta_t^{(\sigma_1,\sigma_2)}  M^{(\sig_1,\sig_2)}\right)_{a_1a_2},\\
\label{Kn3sol} \mathcal{K}^{(3)}_{t, \bsig, \ba} &= \sum_{b_1,b_2,b_3}\Theta_{t}^{(\sig_1,\sig_2)}(a_1,b_1)\Theta_{t}^{(\sig_2,\sig_3)}(a_2,b_2)\Theta_{t}^{(\sig_3,\sig_1)}(a_3,b_3) \cal M^{(3)}_{\bsig,\mathbf b},
\end{align}
where $\bsig=(\sig_1,\ldots,\sig_n)$ and $\ba=(a_1,\ldots,a_n)$ for $n\in \{2,3\}$, and the $M$-loops are defined in \eqref{eq:KMloop}. 

% \begin{align}    \label{Kn2sol}\cK^{(2)}_{t,\bsig, \ba}&=W^{-d}m_1m_2 \Theta^{(\sB)}_{tm_1m_2} (a_1,a_2)  ,  \\
%     \label{Kn3sol} \mathcal{K}^{(3)}_{t, \bsig, \ba} %&= \sum_{[b_1],[b_2],[b_3]}\Theta_{t}^{(\sig_1,\sig_2)}([a_1],[b_1])\Theta_{t}^{(\sig_2,\sig_3)}([a_2],[b_2])\Theta_{t}^{(\sig_3,\sig_1)}([a_3],[b_3]) \cM^{(3)}_{\bsig,\mathbf b}\\
%     &= W^{-2d}m_1m_2m_3\sum_{b}\Theta_{tm_1m_2}^{(\sB)}(a_1,b) \Theta_{tm_2m_3}^{(\sB)}(a_2,b)\Theta_{tm_3m_1}^{(\sB)}(a_3,b) ,
% \end{align}
% where $\bsig=(\sig_1,\ldots,\sig_n)$ and $\ba=(a_1,\ldots,a_n)$ for $n\in \{2,3\}$, and we abbreviate $m_i\equiv m(\sig_i)$ for $i\in\{1,2,3\}$.  
\end{example}

 
\subsection{Proof of the main results}
 
The main results, \Cref{MR:locSC,MR:QuDiff}, for random band matrices follow directly from the following key lemmas on the $G$-loop estimates, while the proof of \Cref{MR:decol_BA} for the block Anderson model will be given separately in \Cref{sec:ext-to-BA}. Recall the notions of stochastic domination, $B_{t, K}$, and $\ell_t$ defined in \eqref{stoch_domination}, \eqref{eq_B_param}, and \eqref{eq:ellt}, respectively. 


\begin{lemma}[$G$-loop estimates]\label{ML:GLoop}
In the setting of \Cref{MR:locSC}, fix any $z=\hat{E}+\ii \eta\in \mathbf D_{\kappa,\e}$ and consider the flow framework in \Cref{zztE}. For each fixed $n\in \N$, the following estimate holds uniformly in $t\in [0,t_0]$: 
\begin{align}\label{Eq:L-KGt}
 \max_{\bsig, \mathbf{a}}\left|{\cal L}^{(n)}_{t, \bsig, \mathbf{a}}-{\cal K}^{(n)}_{t, \bsig, \mathbf{a}}\right|&\prec  (W^{-d}B_{t,0})^{ n}, \\
 \label{Eq:L-KGt2}
\max_{\bsig, \mathbf{a}}\left|{\cal L}^{(n)}_{t, \bsig, \mathbf{a}} \right|&\prec (W^{-d}B_{t,0})^{ n-1}.
\end{align}
%where we recall the stochastic domination and the parameter $B_{t, K}$ defined in Definition \ref{stoch_domination} and \eqref{eq_B_param}. 
\end{lemma}

\begin{lemma}[$2$-loop estimates]\label{ML:GLoop_expec}
In the setting of \Cref{ML:GLoop}, for $\bsig\in\{+,-\}^2$ and $ \ba=(a_1, a_2)\in (\Zn)^2$, the following pointwise estimate holds uniformly in $t\in [0,t_0]$ for any large constant $D>0$: 
\begin{equation}\label{Eq:Gdecay}
 \left| {\cal L}^{(2)}_{t, \bsig, \mathbf{a}}-{\cal K}^{(2)}_{t, \bsig, \mathbf{a}}\right|\prec 
 \p{W^{-d}B_{t, 0}}^{1/5}\cdot \p{W^{-d}B_{t,  |a_1-a_2| }}e^{-\left({| a_1-a_2| }/{\ell_t}\right)^{1/2}} + W^{-D}. 
\end{equation}
Moreover, the expectation of a $2$-$G$-loop satisfies the following $L^\infty$-estimate uniformly in $t\in [0,t_0]$: 
\begin{equation}\label{Eq:Gtlp_exp}
 \max_{\bsig, \ba}\left|\E{\cal L}^{(2)}_{t, \bsig, \ba}-{\cal K}^{(2)}_{t, \bsig, \ba}\right|\prec \p{W^{-d}B_{t,0}}^2 \p{(\ilambda^2W^d)^{-1/5}+W^{-d}B_{t,0}}
 . 
\end{equation}
In the above estimates, the exponent $1/5$ can be replaced by any other positive constant less than $1/4$.
\end{lemma}

\begin{lemma}[Local law for $G_t$]\label{ML:GtLocal}
In the setting of \Cref{ML:GLoop}, the following entrywise local law holds uniformly in $t\in [0,t_0]$: 
%With the notations and assumptions of the previous lemma, we have 
\begin{equation}\label{Gt_bound}
 |(G_{t}-M)_{xy}|^2\prec W^{-d}B_{t,(|x-y|/W)} .
   \end{equation} 
\end{lemma}
 
With these lemmas, we readily conclude the proofs of our main results, \Cref{MR:locSC,MR:QuDiff}.
\begin{proof}[\bf Proof of \Cref{MR:locSC,MR:QuDiff}]
For a fixed $z=\hat{E}+\ii \eta\in \mathbf D_{\kappa,\e}$, we can choose the flow framework as in \Cref{zztE}. 
Then, with \eqref{eq:zztE}, \eqref{eq:BtBt}, and \eqref{Kn2sol}, we observe that \Cref{ML:GtLocal}, \Cref{ML:GLoop} (with $n=1$), and \Cref{ML:GLoop_expec} respectively give the entrywise local law \eqref{G_bound}, the averaged local law \eqref{G_bound_ave}, and the quantum diffusion estimates in \Cref{MR:QuDiff} at each \emph{fixed $z$}. To extend these estimates uniformly to all $z$, we can use a standard $N^{-C}$-net and perturbation argument, whose details we omit. 
\end{proof}

%  \begin{proof}[Proof of Theorem \ref{MR:locSC}]
% This proof is  the same  as that  of Theorems 2.3 in \cite{YY_25}.  More precisely \eqref{Gt_bound} in lemma \ref{ML:GtLocal} implies \eqref{Gt_bound}. The tracial local law \eqref{G_bound_ave} follows the \eqref{Eq:L-KGt} in lemma \ref{ML:GLoop}. 
% \end{proof}
 
 


%\subsection{Strategy for the proof of the loop estimates}

Before concluding this section, we now outline the proofs of \Cref{ML:GLoop,ML:GLoop_expec,ML:GtLocal}. At $t=0$, we have $G_{0}(\sigma)=M(\sigma)$ for $\sig\in\{+,-\}$. Together with \Cref{Def:G_loop,Def_Ktza}, it implies that for any fixed $n\in \N$:   
$${\cL}^{(n)}_{0, \bsig,\ba}= {\cK}^{(n)}_{0, \bsig,\ba},\quad \forall \; \bsig\in \{+,-\}^n,\ \ \ba\in (\Zn)^n. 
$$
Now, for $t>0$, we will establish the following theorem, which provides an induction result that extends the $G$-loop estimate progressively along the stochastic flow.

% By Definitions \ref{Def:G_loop} and \ref{Def_Ktza},   
% $$
% G_{0}(+)=m\cdot I_{N\times N}, \quad {\cal L}_{0, \bsig,\mathbf{a}}= {\cal K}_{0, \bsig,\mathbf{a}},\quad \forall \; \bsig,\mathbf{a}, 
% $$
% where $I_{N\times N}$ is the identity matrix. It is easy to check that  
% \begin{align}\label{ini)bigasya}
%     \hbox{    Lemmas \ref{ML:GLoop}, \ref{ML:GLoop_expec} and \ref{ML:GtLocal} hold at  $t=0$ with no error.}
% \end{align}
% For $t>0$, we will prove the following theorem.


\begin{theorem}\label{lem:main_ind} 
In the setting of \Cref{MR:locSC}, fix any $z=\hat{E}+\ii \eta\in \mathbf D_{\kappa,\e}$ and consider the flow framework in \Cref{zztE}. 
%For the block Anderson model in \Cref{def: BM}, suppose the assumptions of \Cref{zztE} hold with $t_0\le 1 - W^{\fd}\eta_*$.
Suppose the following estimates hold at a fixed $s\in [0,t_0]$: 
\begin{itemize}
\item[(a)] {\bf $G$-loop estimate}: For each fixed $n\ge 1$, we have
\be\label{Eq:L-KGt+IND}
 \max_{\bsig, \ba}\left|{\cal L}^{(n)}_{s, \bsig, \ba}-{\cal K}^{(n)}_{s, \bsig, \ba}\right|\prec (W^{-d}B_{s,0})^{n}.
\ee
\item[(b)] {\bf 2-loop estimate}: 
For $\bsig\in\{+,-\}^2$, $ \ba=(a_1,a_2)\in (\Zn)^2,$ and any large constant $D>0$, we have 
\be\label{Eq:Gdecay+IND}
\left| {\cal L}^{(2)}_{s, \bsig, \ba}-{\cal K}^{(2)}_{s, \bsig, \ba}\right|\prec \p{W^{-d}B_{s, 0}}^{1/5}\cdot \p{W^{-d}B_{s, |a_1-a_2|}} e^{-\left({| a_1-a_2| }/{\ell_s}\right)^{1/2}} +W^{-D}.
\ee
Furthermore, if $1-s\ge \ilambda^2$ (where $\ell_s= 1$ by \eqref{eq:ellt}), we assume a stronger estimate:
\be\label{Eq:Gdecay+IND_s<g}
\left| {\cal L}^{(2)}_{s, \bsig, \ba}-{\cal K}^{(2)}_{s, \bsig, \ba}\right|\prec \p{W^{-d}B_{s, 0}}^2 e^{-| a_1-a_2|^{1/2}} +W^{-D}.
\ee

\item[(c)] {\bf Local law}: We have the (maximum) entrywise local law 
\be \label{Gt_bound+IND}
 \|G_{s}-M\|_{\max} \prec (W^{-d}B_{s,0})^{1/2}.
\ee

\item[(d)] {\bf Expected $2$-loop estimate}: For all $\bsig\in\{+,-\}^2$ and $ \ba=(a_1,a_2)\in (\Zn)^2,$ we have 
 \be \label{Eq:Gtlp_exp+IND}
 \max_{\bsig, \ba}\left|\E{\cal L}^{(2)}_{s, \bsig, \ba}-{\cal K}^{(2)}_{s, \bsig, \ba}\right|\prec \p{W^{-d}B_{s,0}}^2 \p{(\ilambda^2W^d)^{-1/5}+W^{-d}B_{s,0}}. 
\ee
\end{itemize} 
%Assume for some fixed $E$: $|E|\le 2-\kappa$ and $s\in [0,1]$ that Lemmas \ref{ML:GLoop}, \ref{ML:GLoop_expec}  and \ref{ML:GtLocal} hold at time $s$, namely, %with $\eta_s:=\im z_s,\quad \ell_s:=\ell(z_s)$,   
% for $1 \le n\in \mathbb N$ and large $D>0$, 
% \begin{align}\label{Eq:L-KGt+IND}
%  \max_{\bsig, \mathbf{a}}\left|{\cal L}_{s, \bsig, \mathbf{a}}-{\cal K}_{s, \bsig, \mathbf{a}}\right|&\prec (W^{-d} B_{s,0})^{n}  
% \\\label{Eq:Gdecay+IND}
%  \left| {\cal L}_{s, \bsig, \mathbf{a}}-{\cal K}_{s, \bsig, \mathbf{a}}\right|&\prec (W^{-d}B_{s, 0})^{1/5}\cdot (W^{-d}B_{s, |a_1-a_2|})\cdot \exp \left(-\left|\frac{| a_1-a_2| }{\ell_s}\right|^{1/2}\right) +W^{-D} , \quad \bsig=(+,-)
% \\\label{Gt_bound+IND}
%  \|G^{(E)}_{s ,+}-m^{(E)}\|_{\max} & \prec 
%  (W^{-d} B_{s,0})^{1/2}.
% \end{align}
Then, there exists a constant $0<\fc_d \le 10^{-2}$ (depending on $d$, $\kappa$, $\e$, and $\fd$ in \eqref{eq:WO}) such that for any $s< t < 1$ satisfying 
\begin{equation}\label{con_st_ind}
  \p{W^{-d}B_{t,0}}^{\fc_d}\le \frac{1-t}{1-s} < 1, 
\end{equation}
the estimates \eqref{Eq:L-KGt}--\eqref{Gt_bound} hold.  
In the proof, we do not track the exact value of $\fc_d$ (although it can be done by keeping track of the constants carefully in our proof).
In addition, if $1-t\ge \ilambda^2$, we have
\be\label{Eq:Gdecay+s<g}
\left| {\cal L}^{(2)}_{t, \bsig, \ba}-{\cal K}^{(2)}_{t, \bsig, \ba}\right|\prec \p{W^{-d}B_{t, 0}}^2 e^{-| a_1-a_2|^{1/2}} +W^{-D},\quad \forall \bsig\in\{+,-\}^2,\  \ba \in (\Zn)^2.
\ee
\end{theorem}


With \Cref{lem:main_ind} in hand, we can establish \Cref{ML:GLoop,ML:GLoop_expec,ML:GtLocal} via induction on $t$. 
\begin{proof}[\bf Proof of \Cref{ML:GLoop,ML:GLoop_expec,ML:GtLocal}]
We first perform induction from $t=0$ up to $t=1-\ilambda^2$ (when $\ilambda^2\le 1/2$) or $t=1/2$ (when $\ilambda^2>1/2$), establishing \eqref{Eq:L-KGt}--\eqref{Gt_bound} at $t_1:=(1-\ilambda^2)\vee (1/2)$. In the case $\ilambda^2\le 1/2$, we additionally maintain the stronger 2-loop estimate \eqref{Eq:Gdecay+s<g} throughout this step. Next, starting from $t_1$, we continue the induction using \Cref{lem:main_ind} up to $t=t_0$, thereby completing the proof of the estimates stated in \Cref{ML:GLoop,ML:GLoop_expec,ML:GtLocal}.    
\end{proof}

The proof of \Cref{lem:main_ind} is divided into six steps, which comprise the remainder of this paper. Throughout these steps, we assume the hypotheses of \Cref{lem:main_ind} hold. Moreover, each step builds upon the results obtained in the preceding ones.



\medskip 
\noindent
\textbf{Step 1} (A priori $G$-loop bound):  
The $n$-$G$-loops satisfy the a priori bound: 
 \begin{equation}\label{lRB1}
   {\cal L}^{(n)}_{u,\bsig, \mathbf{a}}\prec \p{\frac{1-s}{1-u}}^{n-1}  
   (W^{-d}B_{s,0})^{n-1},\quad  \forall  u\in [s, t].
\end{equation}
%The right-hand side involves $B_s$, though in most applications we replace it with $B_u$ due to monotonicity.  
Furthermore, a weak local law holds in the sense 
\begin{equation}\label{Gtmwc}
    \|G_u-M\|_{\max}\prec  (W^{-d}B_{u,0})^{1/4},\quad \forall  u\in [s, t] .
\end{equation} 
%Here the exponent  is $1/4$ instead of $1/2$  in \eqref{Gt_bound+IND}. 
  
%\textbf{Step 1.5}: \textcolor{blue}{Not needed any more}and  for any   $\bsig=(+,-)$,   $s\le u\le t$, and $D>0$,   
% \begin{equation}\label{Eq:Gdecay_w2a}
%\left| {\cal L}_{u, \bsig, \mathbf{a}}-{\cal K}_{u, \bsig, \mathbf{a}}\right| \prec    (W^{-d}B_{t, 0})^{1/4}\cdot  W^{-d}\left(B_{t,  |a_1-a_2|} +B_{t,  \ell_t}\right)  . \quad 
%\end{equation}  

 
 %\medskip
 \noindent 
\textbf{Step 2}  (A priori $2$-$G$-loop decay): 
The following sharp local laws hold uniformly for all $u\in [s,t]$: 
 \begin{align}\label{Gt_bound_flow}
     |(G_u-M)_{xy}|^2&\prec W^{-d}B_{u,(|x-y|/W)},\quad \forall x,y\in \ZL, \\
     \label{Gt_avgbound_flow} 
    \max_{a\in \Zn}\left|\tr\p{(G_u-M)E_a}\right| &\prec   W^{-d}B_{u,0}.
\end{align}
In particular, the entrywise local law  \eqref{Gt_bound} and the $1$-loop estimate in \eqref{Eq:L-KGt} hold at time $t$. 
In addition, there exists a constant $C_d>0$ (independent of $\fc_d$ in \eqref{con_st_ind}) such that, for any $\bsig\in\{+,-\}^2$, $\ba=(a_1,a_2)\in(\Zn)^2$, and $u\in[s,t]$, the following estimate holds for arbitrarily large constant $D>0$: 
 \begin{equation}\label{Eq:Gdecay_w}
\left| {\cal L}^{(2)}_{u, \bsig, \mathbf{a}}-{\cal K}^{(2)}_{u, \bsig, \mathbf{a}}\right| \prec \p{\frac{1-s}{1-u}}^{C_d}\p{W^{-d}B_{u, 0}}^{1/5}\cdot \p{W^{-d}B_{u,  |a_1-a_2| }} e^{-\left({| a_1-a_2| }/{\ell_u}\right)^{1/2}} +W^{-D} . 
\end{equation}
Hence, at this step, the estimate \eqref{Eq:Gdecay+IND} deteriorates at time $t$ by a factor of $(|1-s|/|1-t|)^{C_d}$. 
We note that the exponent $1/5$---as an arbitrarily chosen positive constant less than $1/4$---in \eqref{Eq:Gdecay_w} is not optimal. Achieving the optimal decay rate for a 2-loop estimate would require this exponent to be 1, which is beyond the scope of the current paper (see also \Cref{rmk_bottleneck} below).



\medskip
\noindent 
\textbf{Step 3}  (Sharp $n$-loop bound): 
The following bound on $n$-$G$-loops holds for any fixed $n\in \N$: 
\begin{equation}\label{Eq:LGxb}
\max_{\bsig, \mathbf{a}}\left| {\cal L}^{(n)}_{u, \bsig, \mathbf{a}} \right|
\prec 
  (W^{-d}B_{u, 0})^{n-1} ,\quad \forall u\in[s,t] .
\end{equation}
In particular, this gives the estimate \eqref{Eq:L-KGt2} at time $t$. 

 \medskip
 \noindent 
\textbf{Step 4} (Sharp  $(\cal L - \cal K)$-loop estimate): The following  estimate on $(\cL-\cK)$-loops holds for any fixed $n\in \N$: 
\begin{equation}\label{Eq:L-KGt-flow}
 \max_{\bsig, \mathbf{a}}\left|{\cal L}^{(n)}_{u, \bsig, \mathbf{a}}-{\cal K}^{(n)}_{u, \bsig, \mathbf{a}}\right|\prec (W^{-d}B_{u, 0})^{n},\quad \forall u\in[s,t] .
\end{equation}
In particular, this gives the estimate \eqref{Eq:L-KGt} at time $t$. 


 \medskip
 \noindent 
\textbf{Step 5} (Pointwise estimate for $2$-loops): For any $\bsig\in\{+,-\}^2$, $ \ba=(a_1,a_2)\in(\Zn)^2,$ and $u\in[s,t]$, the following estimate holds for any large constant $D>0$:
 \begin{equation}\label{Eq:Gdecay_flow}
\left| {\cal L}^{(2)}_{u, \bsig, \ba}-{\cal K}^{(2)}_{u, \bsig, \ba}\right| \prec  \p{W^{-d}B_{u, 0}}^{1/5}\cdot \p{W^{-d}B_{u,  |a_1-a_2| }} e^{-\left({| a_1-a_2| }/{\ell_u}\right)^{1/2}} +W^{-D} . 
\end{equation}
Hence, the pointwise decay estimate \eqref{Eq:Gdecay} holds at time $t$. Furthermore, if $1-t\ge \ilambda^2$, then we have 
\be\label{Eq:Gdecay+s<g_flow}
\left| {\cal L}^{(2)}_{u, \bsig, \ba}-{\cal K}^{(2)}_{u, \bsig, \ba}\right|\prec \p{W^{-d}B_{u, 0}}^2 e^{-| a_1-a_2|^{1/2}} +W^{-D},\quad \forall u\in[s,t].
\ee

%\medskip
\noindent 
\textbf{Step 6} (Expected 2-$G$-loop estimate): 
The following estimate holds: 
\begin{equation}\label{Eq:Gtlp_exp_flow}
 \max_{\bsig, \ba}\left|\E{\cal L}^{(2)}_{u, \bsig, \ba}-{\cal K}^{(2)}_{u, \bsig, \ba}\right|\prec \p{W^{-d}B_{u,0}}^2 \p{(\ilambda^2W^d)^{-1/5}+W^{-d}B_{u,0}},\quad  
 \forall u \in [s, t] .
\end{equation}
Hence, the estimate \eqref{Eq:Gtlp_exp} holds at time $t$. 

\medskip

We remark that the estimates established in each step hold uniformly in $u\in[s,t]$ (recall \eqref{stoch_domination}) due to a standard $N^{-C}$-net and perturbation argument. For simplicity of presentation, we will not emphasize this uniformity at every step of the proof. %Moreover, in Steps 1--5, the induction hypothesis \eqref{Eq:Gtlp_exp+IND} will not be used, i.e., the estimates \eqref{Eq:L-KGt}--\eqref{Gt_bound} and \eqref{Eq:Gdecay+s<g} (in the case $1-t\ge \ilambda^2$) still hold without assuming \eqref{Eq:Gtlp_exp+IND}.



\section{Steps 1 and 2: A priori \texorpdfstring{$G$}{G}-loop estimates}\label{Sec:Steps12}


The remainder of the paper is devoted to proving \Cref{lem:main_ind}, according to the six steps outlined following its statement.
Most of the arguments extend directly to the block Anderson model, with technical differences arising in Step 1 and in the treatment of the light-weight term \smash{$\mathcal{E}^{\Gc,(2)}$} (defined in \eqref{def_EwtG}) when controlling the 2-$G$-loops---specifically, in the proofs of \Cref{lem:LWterm,lem: EWGn2_N} below. %and \Cref{lem:LWterm_EXP} (Step 6). 
We will describe the necessary modifications to adapt the argument to the block Anderson model in \Cref{sec:ext-to-BA}.


\subsection{Proof of Step 1} 

Our proof depends crucially on the following lemma, which provides estimates on the resolvent entries via bounds on 2-$G$-loops. Specifically, \eqref{GiiGEX} and \eqref{GavLGEX} establish bounds on both the entrywise and averaged differences between $G$ and $M$ in the $\max$-norm sense, while \eqref{GijGEX} provides a finer estimate on the off-diagonal resolvent entries, which allows us to derive the decay of the resolvent entries from the decay of the 2-$G$-loops.  


\begin{lemma}\label{lem_GbEXP}
Consider the setting of the random band matrix model. For any $t\in [0,t_0]$, define the event 
\begin{equation}\label{def_asGMc}
    \Omega(t, \e_0) := \big\{\|G_t - M\|_{\max} \leq W^{-\e_0} \big\}
\end{equation}
for a small constant $\e_0>0$. Then, the following \emph{entrywise local law} holds for any constants $\tau,D>0$: 
\begin{equation}
 \P\Big(\mathbf 1(\Omega(t, \e_0))\cdot \|G_{t}-M\|_{\max}^2 \le W^\tau \max_{a,b\in \Zn} \cL^{(2)}_{t,(-,+),(a,b)} \Big)\ge 1-W^{-D} \, .\label{GiiGEX}
\end{equation}
Furthermore, we have the following pointwise decay estimate of $G_t$ for all $a,b\in \Zn$:
\begin{align}\label{GijGEX}
\mathbf 1(\Omega(t, \e_0))\cdot \max_{\substack{x\in [a], y \in [b], x\ne y}} |(G_t)_{xy}|^{2}  \prec & \sum_{\substack{|a'-a|\le 1,|b'-b|\le 1,\\ \bsig\in \{(+,-),(-,+)\}}}  \cL^{(2)}_{t,\bsig,(a',b')}+ W^{-d}\mathbf 1_{|a-b|\le 1}  \, .
\end{align}
Finally, suppose the following estimates hold for a deterministic control parameter $W^{-d/2}\le \Psi_t \le W^{-\e_0}$:\footnote{This parameter $\Psi_t$ should be distinguished from the matrix $\Psi$ appearing in the block Anderson Hamiltonian \eqref{eq:H_blocka}.}
\be\label{initialGT2} 
\|G_{t} - M \|_{\max}\prec W^{-\e_0} ,\quad \max_{a,b\in \Zn} \cL^{(2)}_{t,(-,+),(a,b)} \prec \Psi_t^2 \, . 
\ee
Then, we have the \emph{averaged local law}: 
\begin{align}
\max_{a} \left|\tr\p{\left(G_{t}-M\right)E_{a}} \right| &\prec \Psi_t^2  \, . \label{GavLGEX}
\end{align}
\end{lemma}
\begin{proof}
These estimates have been proven as Lemma 4.1 in \cite{YY_25} for 1D random band matrices. However, their proofs are dimension-independent and use standard arguments based on resolvent identities and large deviation estimates as in \cite{erdHos2012rigidity,erdos2013delocalization}.     
\end{proof}

Step 1 of the proof of \Cref{lem:main_ind} is similar to that in \cite[Section 5.1]{YY_25}, where the core is to establish the following continuity estimate for the $G$-loops. 

\begin{lemma}\label{lem_ConArg}
Fix any \(\e \le s \leq t \leq 1\) for a constant \(\e > 0\). In the flow setting given by \Cref{def_flow} and \Cref{zztE}, assume that the following bound holds at time \(s\) for any fixed $n\in \N$: 
\begin{equation}\label{eq:loopbound_s}
    \max_{\bsig, {\ba}} \abs{{\cal L}^{(n)}_{s, \bsig, {\ba}}} \prec(W^{-d}B_{s,0})^{n-1}.
\end{equation}
% Next, we let $\Omega$ denote the following event:
% \begin{equation}\label{usuayzoo}
%     \Omega := \left\{\|G_{t_2}\|_{\max} \leq 2\right\}.
% \end{equation}
Then, on the event \(\Omega_t= \left\{\|G_{t}\|_{\max} \leq C_0\right\}\) for a constant $C_0>0$, the following estimate holds for any fixed \(n \in \N\) with $n\ge 2$:
\begin{equation}\label{res_lo_bo_eta}
    {\bf 1}(\Omega_t) \cdot \max_{\bsig, {\ba}} \big|{\cal L}^{(n)}_{t, \bsig, {\ba}} \big|\prec \left(
    (W^{-d}B_{s,0})\cdot\frac{\eta_{s}}{\eta_{t}}\right)^{n-1}
    \le  \left(
    (W^{-d}B_{t,0})\cdot\frac{\eta_{s}}{\eta_{t}}\right)^{n-1}.
\end{equation}
\end{lemma} 
\begin{proof}%[Proof of Lemma \ref{lem_ConArg}]%\label{sec:Inh_LE}
The proof of this lemma is exactly the same as that for Lemma 5.1 in \cite{YY_25}, except for some minor changes in notation. Hence, we omit the details.\end{proof}

%Now, combining \Cref{lem_GbEXP,lem_ConArg}, we can readily conclude the proof of Step 1 with a standard continuity argument. 

%\begin{proof}[\bf Step 1: Proof of \eqref{lRB1} and \eqref{Gtmwc}]
With \Cref{lem_GbEXP,lem_ConArg}, Step 1 of the proof of \Cref{lem:main_ind} for the random band matrix model (i.e., the proof of \eqref{lRB1} and \eqref{Gtmwc}) is the same as that in \cite[Section 5.1]{YY_25}. Hence, we omit the details. 
%\end{proof}


\subsection{Dynamics of  \texorpdfstring{$(\cal L-\cal K)$-loops}{L-K}}\label{subsec:DyLK}

We begin by presenting some representations of the $G$-loop dynamics, formulated using the loop hierarchy \eqref{eq:mainStoflow}, Duhamel's principle, and certain evolution kernels, which we introduce below. This dynamics has already been derived for random band matrices in \cite{YY_25} and for the block Anderson model in \cite{RBSO1D}, and will be the basis for Step 2 and subsequent steps of the proof of \Cref{lem:main_ind}. 
For any fixed $n\in \N$, combining the loop hierarchy \eqref{eq:mainStoflow} and the convolution tree equation \eqref{pro_dyncalK}, we obtain the following SDE for the $(\cL-\cK)$-loops as shown in equation (5.12) of \cite{YY_25}: 
\begin{align}\label{eq_L-Keee}
    \dd(\mathcal{L} - \mathcal{K})^{(n)}_{t, \bsig, \ba} = &\left[\cK^{(2)} \sim (\mathcal{L} - \mathcal{K})\right]^{(n)}_{t, \bsig, \ba} \, \dd t+\sum_{\lenk=3}^n \left[\cK^{(\lenk)}\sim (\mathcal{L} - \mathcal{K})\right]^{(n)}_{t, \bsig, \ba}\, \dd t \nonumber\\
    &+ \mathcal{E}^{(\cL-\cK)\times (\cL-\cK),(n)}_{t, \bsig, \ba}\dd t+\mathcal{E}^{\Gc,(n)}_{t, \bsig, \ba}\dd t + \dd\mathcal{E}^{M,(n)}_{t, \bsig, \ba} .
\end{align} 
% \begin{align}\label{eq_L-K-1}
%        \dd(\mathcal{L} - \mathcal{K})^{(n)}_{t, \bsig, \ba} 
%     =&~ W^d \sum_{1 \leq k < l \leq n} \sum_{a,b} 
%    (\mathcal{L} - \mathcal{K})^{(n+k-l+1)}_{t, \cutL^{(a)}_{k, l}\left(\bsig,\, \ba\right)} S^{(\sB)}_{ab}
%    \mathcal{K}^{(l-k+1)}_{t,\cutR^{(b)}_{k,l}\left(\bsig ,\ba\right)}\, \dd t \nonumber\\
%     +&~ W^d \sum_{1 \leq k < l \leq n} \sum_{a,b} \mathcal{K}^{(n+k-l+1)}_{t, \cutL^{(a)}_{k, \,l}\left(\bsig,\ba\right)} 
%    S^{(\sB)}_{ab}\left(\cL-\cK\right)^{(l-k+1)}_{t,\cutR^{(b)}_{k, l}\left(\bsig,\ba\right)} \, \dd t \nonumber\\
%    +&~ \mathcal{E}^{(\cL-\cK)\times (\cL-\cK),(n)}_{t, \bsig, \ba}\dd t+\mathcal{E}^{\Gc,(n)}_{t, \bsig, \ba}
%     \dd t + 
%     \dd\mathcal{E}^{M,(n)}_{t, \bsig, \ba}  ,
% \end{align}
Here, every \smash{$\cK^{(\lenk)}$}, $2\le \lenk\le n$, is regarded as a linear operator acting on the $(\cL-\cK)$-loops, defined as:
\begin{align}\label{DefKsimLK}
&\left[\cK^{(\lenk)}\sim (\mathcal{L} - \mathcal{K})\right]_{t, \bsig, \ba}^{(n)}:= W^d \sum_{1\le k < l \leq n : l-k=\lenk-1} \sum_{a,b} 
   (\mathcal{L} - \mathcal{K})^{(n-\lenk+2)}_{t, \cutL^{(a)}_{k, l}\left(\bsig,\, \ba\right)} 
S^{(\sB)}_{ab}\mathcal{K}^{(\lenk)}_{t,\cutR^{(b)}_{k,l}\left(\bsig ,\ba\right)} \nonumber\\
&\quad +  W^d \sum_{1 \leq k < l \leq n:l-k=n-\lenk+1} \sum_{a,b} \mathcal{K}^{(\lenk)}_{t, \cutL^{(a)}_{k, \,l}\left(\bsig,\ba\right)} S^{(\sB)}_{ab}
   \left(\cL-\cK\right)^{(n-\lenk+2)}_{t,\cutR^{(b)}_{k, l}\left(\bsig,\ba\right)} ,
\end{align}
$\mathcal{E}^{\Gc}$ and $\dd\mathcal{E}^{M}$ are defined in \eqref{def_EwtG} and \eqref{def_Edif}, respectively, and the term $\mathcal{E}^{(\cL-\cK)\times (\cL-\cK)}$ is defined by 
\begin{equation}\label{def_ELKLK}
\mathcal{E}^{(\cL-\cK)\times (\cL-\cK),(n)}_{t, \bsig, \ba} :=
    W^d \sum_{1 \leq k < l \leq n} \sum_{a,b} 
   (\mathcal{L} - \mathcal{K})^{(n+k-l+1)}_{t, \cutL^{(a)}_{k, l}\left(\bsig,\, \ba\right)} 
  S^{(\sB)}_{ab} (\cL-\mathcal{K})^{(l-k+1)}_{t,\cutR^{(b)}_{k,l}\left(\bsig ,\ba\right)}\,  .
\end{equation}
In \eqref{eq_L-Keee}, the superscript ``$(n)$'' indicates the length of the $(\cL-\cK)$-loop on the LHS of the equation, and we will sometimes omit it from our notations when the value of $n$ is clear from the context. 


We now rewrite \eqref{eq_L-Keee} into an integral equation using Duhamel's principle. First, we introduce the evolution kernel associated with this integral equation.


\begin{definition}[Evolution kernel]\label{DefTHUST}
For each $t\in [0,1)$, fixed $n\ge 2$ and $\bsig=(\sigma_1,\ldots,\sig_n)\in \{+,-\}^n$, we define the linear operator ${\varTheta}^{(n)}_{t, \bsig}$ acting on $n$-dimensional tensors ${\cal A}: (\Zn)^{n}\to \mathbb C$ as follows (recall \eqref{eq:Msig}):
\begin{align}\label{def:op_thn}
    \left({\varTheta}^{(n)}_{t, \bsig} \circ \mathcal{A}\right)_{\ba} & = \sum_{i=1}^n \sum_{b_i\in\Zn} \left(\frac{M^{(\sig_i,\sig_{i+1})}S^{(\sB)}}{1 - t M^{(\sig_i,\sig_{i+1})}S^{(\sB)}}\right)_{a_ib_i}  \mathcal{A}_{\ba^{(i)}(b_i)},
\end{align} 
where $\ba=(a_1,\ldots, a_n)\in (\Zn)^n$, \(\ba^{(i)}(b_i):= (a_1, \ldots, a_{i-1}, b_i, a_{i+1}, \ldots, a_n),\) and we adopt the cyclic convention that $\sigma_{n+1}=\sig_1$.  
The evolution kernel corresponding to ${\varTheta}^{(n)}_{t, \bsig}$ is given by
    \begin{align}\label{def_Ustz}
        \left(\mathcal{U}_{s, t, \bsig}^{(n)} \circ \mathcal{A}\right)_{\ba} = \sum_{\mathbf{b} = (b_1, \ldots, b_{n})} \prod_{i=1}^n \left(\frac{1 - s \cdot M^{(\sig_i,\sig_{i+1})} S^{(\sB)} }{1 - t \cdot M^{(\sig_i,\sig_{i+1})} S^{(\sB)}}\right)_{a_i b_i} \cdot \mathcal{A}_{\mathbf{b}} . 
    \end{align}
% Define the linear map ${\varTheta}_{t, \bsig}$ via the following action on a tensor ${\cal A}: \mathbb Z^n_L\to \mathbb C$: 
% \begin{align}\label{def:op_thn}
%     \left({{\varTheta}}_{t, \bsig} \circ \mathcal{A}\right)_{{\ba}} & = \sum_{i=1}^n \sum_{b_i} \left(\frac{m_i m_{i+1}  }{1 - t m_i m_{i+1} S^{(B)}}\right)_{a_i b_i} \cdot \mathcal{A}_{{\ba}^{(i)}},\quad m_i=m(\sigma_i) \nonumber \\
%     \quad {\ba}^{(i)}&  = (a_1, \ldots, a_{i-1}, b_i, a_{i+1}, \ldots, a_n).
% \end{align}  
% We also define the semigroup operator $\mathcal{U}_{s,t,\bsig}$ associated to $\Theta_{t,\bsig}$ as follows
%     \begin{align}\label{def_Ustz}
%         \left(\mathcal{U}_{s, t, \bsig} \circ \mathcal{A}\right)_{{\ba}} = \sum_{b_1, \ldots, b_n} \prod_{i=1}^n \left(\frac{1 - s \cdot m_i m_{i+1} S^{(B)}}{1 - t \cdot m_i m_{i+1} S^{(B)}}\right)_{a_i b_i} \cdot \mathcal{A}_{\textbf{b}}, \quad \textbf{b} = (b_1, \ldots, b_n), \quad m_i=m(\sigma_i)
%     \end{align}
\end{definition}


By the definition of the 2-$\cK$-loop in \eqref{Kn2sol}, we observe that the first term on the RHS of \eqref{eq_L-Keee} can be rewritten as ${\varTheta}^{(n)}_{t, \bsig} \circ (\mathcal{L} - \mathcal{K})^{(n)}_{t,\bsig}$. With this fact and using Duhamel's principle, we obtain the following lemma.


\begin{lemma}[Integrated loop hierarchy, Lemma 5.3 of \cite{YY_25}] \label{Sol_CalL}
First, \eqref{eq_L-Keee} is equivalent to the integral equation: for any stopping time $\tau\ge s$ with respect to the matrix Brownian motion $\{H_t\}$,
\begin{align}\label{int_K-L_ST}
({\cal L}-{\cal K})^{(n)}_{\tau,\bsig,{\ba}}&=({\cal L}-{\cal K})^{(n)}_{s,\bsig,{\ba}}+\int_{s}^{\tau}\p{\varTheta^{(n)}_{u,\bsig}\circ({\cal L}-{\cal K})^{(n)}_{u,\bsig}}_{{\ba}}\dd u+\sum_{l_\mathcal{K}=3}^n \int_s^\tau \Big[\mathcal{K}^{(\lenk)} \sim (\mathcal{L} - \mathcal{K})\Big]^{(n)}_{u, \bsig, {\ba}}\dd u \nonumber\\
&+ \int_s^\tau \mathcal{E}^{(\mathcal{L} - \mathcal{K}) \times (\mathcal{L} - \mathcal{K}),(n)}_{u, \bsig, {\ba}}\dd u + \int_{s}^{\tau}\mathcal{E}^{\Gc,(n)}_{u, \bsig, {\ba}}\dd u + \int_s^\tau \dd \mathcal{E}^{M,(n)}_{u, \bsig, {\ba}} .
\end{align}
%    
%Then we have the (simple version) of loop hierarchy (\textcolor{blue}{new})  \begin{align}\label{int_K-LcalE_simple}    (\mathcal{L} - \mathcal{K})_{t, \bsig, {\ba}} \;=\; \; &     (\mathcal{L} - \mathcal{K})_{s, \bsig,{\ba}}     \; + \;\sum_{l_\mathcal{K}\ge 2} \int_{s}^t    \Big[\mathcal{K} \sim (\mathcal{L} - \mathcal{K})\Big]^{l_\mathcal{K}}_{u, \bsig, {\ba}} du \nonumber \\   &+ \int_{s}^t   \mathcal{E}^{((\mathcal{L} - \mathcal{K}) \times (\mathcal{L} - \mathcal{K}))}_{u, \bsig, {\ba}} du \;+\; \int_{s}^t  \mathcal{E}^{(\widetilde{G})}_{u, \bsig, {\ba}} du + \int_{s}^t  \mathcal{E}^{(M)}_{u, \bsig, {\ba}}  
%\end{align}
Second, applying Duhamel's principle to \eqref{eq_L-Keee}, we obtain the following \emph{integrated loop hierarchy}:
\begin{align}\label{int_K-LcalE}
   & (\mathcal{L} - \mathcal{K})^{(n)}_{t, \bsig, \ba}  =
    \left(\mathcal{U}^{(n)}_{s, t, \bsig} \circ (\mathcal{L} - \mathcal{K})^{(n)}_{s, \bsig}\right)_{\ba} + \sum_{l_\mathcal{K} =3}^n \int_{s}^t \left(\mathcal{U}^{(n)}_{u, t, \bsig} \circ \Big[\cK^{(\lenk)}\sim (\mathcal{L} - \mathcal{K})\Big]^{(n)}_{u, \bsig}\right)_{\ba} \dd u \nonumber \\
    &+ \int_{s}^t \left(\mathcal{U}^{(n)}_{u, t, \bsig} \circ \mathcal{E}^{(\cL-\cK)\times(\cL-\cK),(n)}_{u, \bsig}\right)_{\ba} \dd u  + \int_{s}^t \left(\mathcal{U}^{(n)}_{u, t, \bsig} \circ \mathcal E^{\Gc,(n)}_{u, \bsig}\right)_{\ba} \dd u + \int_{s}^t \left(\mathcal{U}^{(n)}_{u, t, \bsig} \circ \dd \mathcal E^{M,(n)}_{u, \bsig}\right)_{\ba} .
\end{align}
% \begin{align}\label{int_K-LcalE}
%     (\mathcal{L} - \mathcal{K})_{t, \bsig, {\ba}} \;=\; \; &
%     \left(\mathcal{U}_{s, t, \bsig} \circ (\mathcal{L} - \mathcal{K})_{s, \bsig}\right)_{{\ba}} \nonumber \\
%     &+ \sum_{l_\mathcal{K} > 2} \int_{s}^t \left(\mathcal{U}_{u, t, \bsig} \circ \Big[\mathcal{K} \sim (\mathcal{L} - \mathcal{K})\Big]^{l_\mathcal{K}}_{u, \bsig}\right)_{{\ba}} du \nonumber \\
%     &+ \int_{s}^t \left(\mathcal{U}_{u, t, \bsig} \circ \mathcal{E}^{((\mathcal{L} - \mathcal{K}) \times (\mathcal{L} - \mathcal{K}))}_{u, \bsig}\right)_{{\ba}} du  \nonumber \\
%     &+ \int_{s}^t \left(\mathcal{U}_{u, t, \bsig} \circ \mathcal{E}^{(\widetilde{G})}_{u, \bsig}\right)_{{\ba}} du + \int_{s}^t \left(\mathcal{U}_{u, t, \bsig} \circ \mathcal{E}^{(M)}_{u, \bsig}\right)_{{\ba}}  
% \end{align}
% \begin{align}\label{int_K-L_ST}
%    (\mathcal{L} - \mathcal{K})_{\tau , \bsig, {\ba}} \;=\; \; &
%     \left(\mathcal{U}_{s, \tau, \bsig} \circ (\mathcal{L} - \mathcal{K})_{s, \bsig}\right)_{{\ba}}  \nonumber \\
%     &+ \sum_{l_\mathcal{K} > 2} \int_{s}^\tau \left(\mathcal{U}_{u, \tau, \bsig} \circ \Big[\mathcal{K} \sim (\mathcal{L} - \mathcal{K})\Big]^{l_\mathcal{K}}_{u, \bsig}\right)_{{\ba}} du  \nonumber \\
%     &+ \int_{s}^\tau \left(\mathcal{U}_{u, \tau, \bsig} \circ \mathcal{E}^{((\mathcal{L} - \mathcal{K}) \times (\mathcal{L} - \mathcal{K}))}_{u, \bsig}\right)_{{\ba}} du  \nonumber \\
%     &+ \int_{s}^\tau \left(\mathcal{U}_{u, \tau, \bsig} \circ \mathcal{E}^{(\widetilde{G})}_{u, \bsig}\right)_{{\ba}} du + \int_{s}^\tau \left(\mathcal{U}_{u, \tau, \bsig} \circ \mathcal{E}^{(M)}_{u, \bsig}\right)_{{\ba}}.
% \end{align}
\end{lemma}


Before returning to the proof of \Cref{lem:main_ind}, we introduce the following notation, which controls the quadratic variation of the martingale term in \eqref{def_Edif}. More precisely, $\dd\cal E^{M,(n)}$ is a sum over the $n$ edges of the loop at which the loop is cut, and the tensor defined below keeps only the terms with the same cut on both sides. Hence, by the Cauchy--Schwarz inequality over the $n$ cuts, the quadratic variation of any linear combination $\sum_{\ba}\varphi_{\ba}\,\dd\cal E^{M,(n)}_{u,\bsig,\ba}$ is bounded by $n\sum_{\ba,\ba'}\varphi_{\ba}\overline{\varphi_{\ba'}}\,(\cal E\otimes\cal E)^{M,(n)}_{u,\bsig,\ba,\ba'}\,\dd u$; in general, the two are not equal because of the cross terms between different cuts. The factor $n$ is absorbed into the constant $C_{n,p}$ in \Cref{lem:DIfREP} below.

 
\begin{definition}[Quadratic variation loop]\label{def:CALE} 
For $t\in [0,1]$ and $\bsig=(\sigma_1,\ldots,\sig_n)\in \{+,-\}^n$, we introduce the $(2n)$-dimensional \emph{quadratic variation tensor} for any \(\ba=(a_1,\ldots, a_n) \) and \(\ba'=(a_1',\ldots, a_n')\):
 %  \begin{align*}
 % \left( \cal E \otimes \cal E  \right)^{M,(n)}_{t, \bsig, \ba, \ba'} \; := \; &
 % \sum_{k=1}^n \left( \cal E \otimes \cal E \right)^{M,(n;k)}_{t,  \bsig, \ba, \ba'},\quad %\forall \ba=(a_1,\ldots, a_n) ,\ \ba'=(a_1',\ldots, a_n'), %\in (\Zn)^n,
 % \end{align*}
 % where each $\left( \cal E \otimes \cal E \right)^{M,(n;k)}_{t,  \bsig, \ba, \ba'}$ is defined as 
\be\label{defEOTE}
 \left( \cal E \otimes \cal E  \right)^{M,(n)}_{t, \bsig, \ba, \ba'} := 
 \sum_{k=1}^n \left( \cal E \otimes \cal E \right)^{M,(n;k)}_{t,  \bsig, \ba, \ba'},\ \ \text{where}\ \ \left( \cal E \otimes \cal E \right)^{M,(n;k)}_{t,  \bsig, \ba, \ba'} :
 =   W^d \sum_{b,b'} S^{(\sB)}_{bb'}{\cal L}^{(2n+2)}_{t, (\bsig\otimes\overline\bsig)^{(k) },(\ba\otimes \ba')^{(k)}(b,b')}.
\ee
Here, ${\cal L}^{(2n+2)}$ denotes a $(2n+2)$-loop obtained by cutting the $k$-th edge of ${\cal L}^{(n)}_{t,\bsig,\ba}$ and then gluing it (with indices $\ba$) with its conjugate loop (with indices $\ba'$) along the newly introduced vertices $b$ and $b'$. Formally:
\begin{align*}
   & {\cal L}^{(2n+2)}_{t, (\bsig\otimes\overline\bsig)^{(k) },(\ba\otimes \ba')^{(k )}(b,b')}:=  \tr \bigg\{ \prod_{i=k}^n \left(G_{t}(\sigma_i) E_{a_i}\right) \cdot \prod_{i=1}^{k-1} \left(G_{t}(\sigma_i) E_{a_i}\right) \cdot G_t(\sigma_k) E_{b} G_t(-\sigma_k) \\
    &\qquad \times \prod_{i={1}}^{k-1} \left(E_{a_{k-i}'} G_{t}(-\sigma_{k-i}) \right)\cdot \prod_{i=k}^{n} \left(E_{a_{n+k-i}'} G_{t}(-\sigma_{n+k-i}) \right) E_{b'}\bigg\} ,
\end{align*}
where the notations $(\ba\otimes \ba')^{(k )}$ and $(\bsig\otimes\overline \bsig)^{(k) }$ (with $\overline\bsig$ denoting $(-\sig_1,\ldots,-\sig_n)$) represent respectively
\begin{align}\label{def_diffakn_k}
 (\ba\otimes \ba')^{(k )}(b,b')&=( a_k,\ldots, a_n, a_1,\ldots , a_{k-1}, b, a'_{k-1},\ldots, a_1', a_n',\ldots, a'_{k},b'),\nonumber \\
  (\bsig\otimes \overline\bsig)^{(k )}&=(  \sigma_k, \ldots, \sigma_n,   \sigma_1,\ldots ,\sigma_{k}, -\sigma_{k}, \ldots, -\sigma_1 ,   -\sigma_n, \ldots, -\sigma_{k }).
\end{align}
%We remark that the symbol ``$\otimes$'' in the above notations is used to emphasize the symmetric structure and does not represent any kind of ``product''. 

% Denote 
%   \begin{align}\label{defEOTE}
%  \left( \mathcal{E}\otimes  \mathcal{E} \right)_{t, \,\bsig, \,{\ba}, \,{\ba}'} \; := \; &
%  \sum_{k=1}^n\left( \mathcal{E}\otimes  \mathcal{E} \right)^{(k)}_{t, \,\bsig, \,{\ba}, \,{\ba}'}
%  \nonumber \\
%  \left( \mathcal{E}\otimes  \mathcal{E} \right)^{(k)}_{t, \,\bsig, \,{\ba}, \,{\ba}'} :
% \; = \; & W^{d} \sum_{b,b'} S^{(B)}_{b,b'}{\cal L}_{t, \bsig^{(k ) },{\ba}^{(k )}},\quad \bsig^{(k ) }\in \{+,-\}^{2n+2}, 
% \end{align}
% where ${\cal L}_{t, \bsig^{(k ) },{\ba}^{(k )}}$ is obtained by cutting the $k$-th edge of ${\cal L}_{t,\bsig,{\ba}}$ and then attaching itself (with indices ${\ba}$) to its complex conjugate loop (with indices ${\ba}'$) into a bigger loop that is labeled by indices $b$ and $b'$, so that
% %where ${\ba}^{(k )}(m)$ is the $m$-th component of ${\ba}^{(k )}$. More precisely
% \begin{align}\label{def_diffakn_k}
%    & {\ba}^{(k )}=( a_k,a_{k+1},\cdots a_n, a_1,\cdots a_{k-1}, b', a'_{k-1}\cdots a_1', a_n'\cdots a'_{k},b)
%     \nonumber \\
%   &  \bsig^{(k )}=( \, \sigma_k, \, \sigma_{k+1},\cdots \sigma_n,  \, \sigma_1,\cdots \sigma_{k}, \,  \overline\sigma_{k},  \cdots \overline\sigma_1 , \,  \overline\sigma_n \cdots \overline\sigma_{k })
% \end{align}
%The symbol $\otimes$ is used to emphasize the symmetric structure of $\mathcal{E} \otimes \mathcal{E}$; it is not a tensor product.  


\end{definition}

 


Under the above notations, applying the Burkholder-Davis-Gundy inequality, we obtain the following lemma that provides high-moment bounds on the martingale term. %is an immediate consequence of the Burkholder-Davis-Gundy inequality.  We now estimate the martingale term in \eqref{int_K-L_ST} via Lemma 5.5 in \cite{YY_25}, which we record below.

\begin{lemma}[Lemma 5.5 of \cite{YY_25}]\label{lem:DIfREP}
Let $\tau$ be a stopping time with respect to the matrix Brownian motion $\{H_t\}$. Then, for any fixed $p\in \N$, there exists a constant $C_{n,p}$ such that
\begin{align}\label{aaswtghh}
   \E \left [    \int_{s}^\tau \dd  \mathcal{E}^{M,(n)}_{u, \bsig, {{\ba}}}  \right ]^{2p} 
   &\le C_{n,p}  
   \E \left(\int_{s}^\tau 
   \left( 
   \left( \mathcal{E} \otimes  \mathcal{E} \right)
  ^{M,(n)}_{u, \bsig, \ba,\ba}  
   \right)\dd u 
   \right)^{p},\\
   \label{alu9_STime}
   \E \left [    \int_{s}^t\left(\mathcal{U}^{(n)}_{u, t, \bsig} \circ \dd\cal E^{M,(n)}_{u, \bsig}\right)_{\ba} \right ]^{2p} 
   &\le C_{n,p}
   \E \left(\int_{s}^t 
   \left(\left(
   \mathcal{U}^{(n)}_{u, t,  \bsig}
   \otimes 
   \mathcal{U}^{(n)}_{u, t,  \overline{\bsig}}
   \right) \;\circ  \;
   \left( \cal E \otimes \cal E\right)^{M,(n)}
   _{u, \bsig  }
   \right)_{\ba, \ba}\dd u 
   \right)^{p}.
  \end{align} 
Here, $\mathcal{U}^{(n)}_{u, t,  \bsig} \otimes    \mathcal{U}^{(n)}_{u, t,  \overline{\bsig}}$ denotes the tensor product of the evolution kernel defined in \eqref{def_Ustz}: 
  $$
  \left[\left(
   \mathcal{U}^{(n)}_{u, t,  \bsig}
   \otimes 
   \mathcal{U}^{(n)}_{u, t,  \overline{\bsig}}
   \right)\circ \cal A\right]_{\ba,\ba'}=
   \sum_{\bfb,\bfb'\in(\Zn)^{n}}  \prod_{i=1}^n \left(\frac{1 - u\cdot M^{(\sig_i,\sig_{i+1})} S^{(\sB)}}{1 - t\cdot  M^{(\sig_i,\sig_{i+1})} S^{(\sB)}}\right)_{a_i b_i} \prod_{i=1}^n \left(\frac{1 - u\cdot M^{(-\sig_i,-\sig_{i+1})} S^{(\sB)} }{1 - t\cdot M^{(-\sig_i,-\sig_{i+1})} S^{(\sB)} }\right)_{a'_i b'_i} \mathcal{A}_{\bfb,\bfb'}
  $$
  for any $(2n)$-dimensional tensor ${\cal A}: (\Zn)^{2n}\to \mathbb C$ and $\bfb=(b_1,\ldots, b_n)$, $\bfb'=(b_1',\ldots, b_n')$.
  
% For any stopping time $T$ with respect to the filtration generated by $H_t$ we have
%  \begin{equation}\label{aaswtghh}
%    \E \left [    \int_{s}^\tau\left(  \mathcal{E}^{(M)}_{u, \bsig, {{\ba}}}\right) \right ]^{2p} 
%    \le C_{n,p} \; 
%    \E \left(\int_{s}^\tau 
%    \left( 
%    \left( \mathcal{E} \otimes  \mathcal{E} \right)
%    _{u, \,\bsig, {{\textbf a}, {\textbf a}}  }
%    \right)du 
%    \right)^{p}
%   \end{equation} 
  
  
%   \begin{equation}\label{alu9_STime}
%    \E \left [    \int_{s}^\tau\left(\mathcal{U}_{u, t, \bsig} \circ \mathcal{E}^{(M)}_{u, \bsig}\right)_{{\ba}} \right ]^{2p} 
%    \le C_{n,p} \; 
%    \E \left(\int_{s}^\tau 
%    \left(\left(
%    \mathcal{U}_{u, t,  \bsig}
%    \otimes 
%    \mathcal{U}_{u, t,  \overline{\bsig}}
%    \right) \;\circ  \;
%    \left( \mathcal{E} \otimes  \mathcal{E} \right)
%    _{u, \,\bsig  }
%    \right)_{{\textbf a}, {\textbf a}}du 
%    \right)^{p}
%   \end{equation} 
%   where $\tau:=t\wedge T$, where  $\overline{\bsig}$ is the conjugate of ${\bsig}$, and where
%   $$
%   \left[\left(
%    \mathcal{U}_{u, t,  \bsig}
%    \otimes 
%    \mathcal{U}_{u, t,  \overline{\bsig}}
%    \right)\circ \cal A\right]_{{\ba},{\ba}'}=
%    \sum_{\textbf{b},\;\textbf{b}'} \; \prod_{i=1}^n \left(\frac{1 - u \cdot m_i m_{i+1} S^{(B)}}{1 - t \cdot m_i m_{i+1} S^{(B)}}\right)_{a_i b_i}
%    \;\cdot \;\prod_{i=1}^n \left(\frac{1 - u \cdot \overline{m_i m_{i+1}} S^{(B)}}{1 - t \cdot \overline{m_i m_{i+1}} S^{(B)}}\right)_{a'_i b'_i} \cdot \mathcal{A}_{\textbf{b},\textbf{b}'}.
%   $$
\end{lemma}



 
  





\subsection{Proof of Step 2}\label{subsec:step2_pf}

In this subsection, we focus on the proof of \eqref{Eq:Gdecay_w}. The local laws \eqref{Gt_bound_flow} and \eqref{Gt_avgbound_flow} then follow directly from \eqref{Eq:Gdecay_w} together with Lemma \ref{lem_GbEXP}. 
Because of the hierarchical structure of \eqref{eq:mainStoflow}, establishing the pointwise stability estimate \eqref{Eq:Gdecay_w} would, in principle, require \emph{pointwise control} of all higher-order $G$-loops with length $n \ge 3$. Truncating this ``pointwise'' loop hierarchy (even if possible) would be substantially more involved than handling the ``maximum'' loop hierarchy used in our current approach.
Fortunately, by combining two key technical ingredients---the diagrammatic techniques developed for the light-weight estimate and the $\contract$ inequality (see \Cref{ygdhmsgq0}) used in the martingale estimate---we can overcome this difficulty and effectively truncate the pointwise loop hierarchy at order $n=2$. This is one of the pivotal components of our proof: without it, both Step 2 and the $L^\infty$-stability of the tree approximation of the loop hierarchy, established in Steps 3 and 4 below, would break down. 

We begin by defining the following tail functions to quantify the pointwise decay.

 
\begin{definition}[Tail functions]\label{def: TTfunc}
For any $t\in [0,1)$, we define a tail function $\cal T_{t}:[0,\infty)\to [0,\infty)$ as
\begin{align}\label{defTUL}
  {\cal T}_{t}(r):=  B_{t,r}\cdot e^{- \left( 
  r/\ell_t \right) ^{1/2} },\quad \forall r\ge 0,
  \end{align}
which corresponds to the $\Theta$-propagator bound in \eqref{prop:ThfadC}. 
Given any $0\le \ell \le L$ and large constant $D>0$, we will also use a truncated tail function \smash{$\widetilde{{\cal T}}$} defined as: 
\begin{align}\label{defWTTlD}
   \wT^{\ell}_{t,D}(r):= \max\left( {\cal T}_{t}(r\wedge \ell), \;  W^{-D}\right).
   %,\quad a\in (\mathbb Z_L)^d
\end{align}
 \end{definition}

Note that ${\cal T}_{t}$ is a decreasing function of $r\ge 0$. Moreover, for $0\le r\le L$, when $1-t\ge \ilambda^2/L^{2}$, the term $(L^{d}|1-t|)^{-1}$ is always dominated by the term $(\ilambda^2+|1-t|)^{-1}/(r+1)^{d-2}$ in $B_{t,r}$, while when $1-t\le \ilambda^2/L^{2}$, $\ell_t=L$ and the exponential factor \smash{$\exp (- (r/\ell_t ) ^{1/2})$} in \eqref{defTUL} is of constant order.  
The tail function $\cT_t$ satisfies the following elementary estimate, whose proof is provided in \Cref{sec:pfpropT}. %cite \cite{DYYY25_d3} here

\begin{lemma}[Property of $\cal T$]\label{lem:propT}
There exists a constant $C_d>0$ (depending only on $d$) such that the following bound holds 
% for all $t\in [0,1)$ and $a,b\in \Zn$:
%\begin{align}\label{TTT}
%    \sum_{c\in \Zn}{\cal T}_{t}(|a-c|)\cdot {\cal T}_{t}(| c -b |) \le \frac{C_d}{1-t}\cdot {\cal T}_{t}(|a-b|) .
%    % \label{TTT2}
%    % \sum_{c\in \Zn}{\cal T}_{t}(|a-c|)\cdot {\cal T}_{t}(| c-b|) \le \frac{C_d}{1-t}\cdot  {\cal T}_{t} (|a-b|). 
%\end{align}
%More generally, 
for any $0\le u\le t<1$ satisfying (i) $1-u\ge 1- t\ge \ilambda^2/L^{2}$, or (ii) $1-t\le 1-u \le \ilambda^2/L^{2}$:
\begin{align}\label{TTT2}
 \sum_{c\in \Zn}{\cal T}_{u}(|a-c|)\cdot {\cal T}_{t}(| c -b |) \le \frac{C_d}{1-u}\cdot {\cal T}_{t}(|a-b|) ,\quad \forall a,b\in \Zn.
\end{align}
\end{lemma}



Our core strategy for the proof of Step 2 involves iterating the following self-improving estimates: for any large constant $D>0$, we have that  %\textcolor{blue}{(for some $\ell$ and $\ell'$ roughly satisfying \eqref{ll'})}
\begin{align}\label{awi2iks}
\sup_{u\in [s,t]}\max_{\ba=(a,b)} \frac{\big|{\cal L}^{(2)}_{u,(+,-),\ba}\big|}{W^{-d} \widetilde{{\cal T}}^{\ell}_{u,D}\p{ |a-b|}}\prec 1  &\implies \sup_{u\in [s,t]}\max_{\ba=(a,b)}\max_{\bsig} \frac{\big|{(\cal L-\cal K)}^{(2)}_{u,\bsig,\ba}\big|}{W^{-d} \widetilde{{\cal T}}^{\ell} _{u,D}\p{|a-b|}}\prec \p{\frac{1-s}{1-t}}^{C_d}(W^{-d}B_{t,0})^{1/5} \nonumber
\\
&\implies \sup_{u\in [s,t]}\max_{\ba=(a,b)} \frac{\big|{\cal L}^{(2)}_{u,(+,-),\ba}\big|}{W^{-d} \widetilde{{\cal T}}^{\ell'}_{u,D}\p{ |a-b|}}\prec 1 ,
\end{align}
where the two scales $\ell'>\ell$ are roughly connected through the relation \({\cal T}_{t} \p{\ell'}\asymp  (W^{-d}B_{t,0})^{1/6}\cdot {\cal T}_{t}\p{\ell}.\)
% \begin{equation}\label{ll'}
%     {\cal T}_{t} \p{\ell'}\sim  (W^{-d}B_{t,0})^{1/6}\cdot {\cal T}_{t}\p{\ell}.
% \end{equation} 
Here, the decrease from $1/5$ to $1/6$ accounts for the prefactor $(|1-s|/|1-t|)^{C_d}$. In other words, \eqref{awi2iks} shows that once we have obtained a sharp 2-$G$-loop bound up to the scale $\ell$, then after one iteration, we can push this bound to a slightly larger scale $\ell'$. After $\OO(1)$ iterations, we can improve the 2-$G$-loop bound up to a scale \smash{$\wh\ell$} with \smash{$\cal T_t(\wh\ell)=\OO(W^{-D})$}. At this scale, we have 
\[{(\cal L-\cal K)}^{(2)}_{u,\bsig,\ba}\prec W^{-d}\widetilde{{\cal T}}^{\wh\ell}_{u,D}\p{|a-b|}\asymp W^{-d}\widetilde{{\cal T}}^{L}_{u,D}\p{|a-b|}.\]
Together with the middle step of \eqref{awi2iks}, it implies the desired estimate \eqref{Eq:Gdecay_w}. 

To establish \eqref{awi2iks}, we use equation \eqref{int_K-L_ST} with $n=2$:
\begin{align}\label{LK_simple}
({\cal L}-{\cal K})^{(2)}_{\tau,\bsig,{\ba}}&=({\cal L}-{\cal K})^{(2)}_{s,\bsig,{\ba}}+\int_{s}^{\tau}\p{\varTheta^{(2)}_{u,\bsig}\circ({\cal L}-{\cal K})^{(2)}_{u,\bsig} + \mathcal{E}^{(\mathcal{L} - \mathcal{K}) \times (\mathcal{L} - \mathcal{K})}_{u, \bsig} + \mathcal{E}^{\Gc}_{u, \bsig} }_{{\ba}}\dd u + \int_s^\tau \dd \mathcal{E}^{M}_{u, \bsig, {\ba}} ,
\end{align} 
where we omit ``$(2)$'' from some superscripts. %By equation \eqref{int_K-L_ST}, a similar equation as \eqref{LK_simple} with $t$ replaced by $\tau$ also holds. 
Define the following (random) control parameter:
\begin{align}\label{defCALJ}
\wh{\cal J}^{\ell}_{u,D}:=\max_{\ba=(a,b)} \max_{\bsig}{\big|\left({\cal L-\cal K}\right)^{(2)}_{u,\bsig, \ba}\big|}\big/\big[ W^{-d} \widetilde{{\cal T}}^{\ell}_{u,D}\p{|a-b|}\big].
\end{align}
The first two terms in the integrand on the RHS of \eqref{LK_simple} can be bounded as in the following lemma, whose proof is postponed to  
\Cref{subsec:pf_lem:newKLK}.

\begin{lemma}\label{lem:newKLK} 
In the setting of \Cref{lem:main_ind}, suppose the weak local law \eqref{Gtmwc} holds. Then, there exists a constant $C>0$ (depending only on $c_d$ and $C_d$ in \eqref{prop:ThfadC}) such that the following estimates hold with high probability for all $0\le \ell \le L$:
\begin{align}\label{juwo2=klk}      
\max_{\ba=(a,b)} \max_{\bsig}\abs{\p{\varTheta^{(2)}_{u,\bsig}\circ({\cal L}-{\cal K})^{(2)}_{u,\bsig}}_{{\ba}}}\Big/\br{W^{-d}\widetilde{{\cal T}}^{\ell}_{u,D}\p{|a-b|}} & \le  \frac{C}{1-u}\cdot  \wh{\cal J}^{\ell}_{u, D }  ,\\
\label{juwo=Lklk}      
\max_{\ba=(a,b)}\max_{\bsig} \left|\mathcal{E}^{(\mathcal{L} - \mathcal{K}) \times (\mathcal{L} - \mathcal{K})}_{u, \bsig, {\ba}} \right|\Big/\br{W^{-d}\widetilde{{\cal T}}^{\ell}_{u,D}\p{ |a-b|}} &\le  \frac{C}{1-u} \cdot \left( \wh{\cal J}^{\ell}_{u,D} + \left(\wh{\cal J}^{\ell}_{u,D} \right)^2\cdot {\bf 1}_{\ell\ge 1}\right)  .
\end{align}  
\end{lemma}


Next, we bound the light-weight term $\cal E^{\Gc}$ using the following two lemmas. \Cref{lem:LWterm} is applied to recover the polynomial decay $B_{t,\cdot}$ in \eqref{eq_B_param}, while \Cref{lem: EWGn2_N} is used to establish the tail behavior described by the tail function in \eqref{defWTTlD}.

%\textcolor{blue}{In the following three lemmas, you use "polynomial and exponential decay". The readers cannot understand you as we do, maybe just "with B bounds", with $T$ bound. and you need to recall $B$ in \eqref{eq_B_param}}
 
 \begin{lemma}[Light-weight estimate: $B$-bound]\label{lem:LWterm}
In the setting of \Cref{lem:main_ind}, assume that the bounds in \eqref{initialGT2} hold for some constant $\e_0>0$ and deterministic control parameter $ W^{-d/2}\le \Psi_t \le W^{-\e_0}$. Suppose we have the estimate 
\be\label{eq:LW_assm}
{\cal L}^{(2)}_{t,\bsig,(a,b)}\prec \Psi_t^2(|a-b|),\quad \forall \ \bsig\in\{(+,-),(-,+)\}, \ a,b\in \Zn, \ee
for a class of deterministic parameters $0<\Psi_t(|a-b|)\le W^{-\e_0}$. Without loss of generality, assume that $\Psi_t(\ell)$ is monotonically decreasing for $\ell\ge 0$ (otherwise, one may replace it by the function $\sup_{\ell'\ge\ell}\Psi_t(\ell')$). 
Suppose there exist constants $C_1,C_2>1$ such that the following relations hold for any constant $C>1$: 
\be\label{eq:Psi}
W^{-d/2}\lesssim \Psi_t(0)\asymp\Psi_t(\ell) \ \  \forall \ 0\le \ell\le C,\quad \text{and}\quad \Psi_t(\ell_1)/\Psi_t(\ell_2)\le C_1(\ell_2/\ell_1)^{C_2} \ \ \forall \ell_2\ge \ell_1\ge 1.
\ee
Then, the following estimate holds: 
%there exists a constant $c>0$ such that the following estimate holds:    
\be\label{eq:LW_conclusion}
\mathcal{E}^{\Gc,(2)}_{t, \bsig, (a,b)}\prec \frac{1}{\eta_t}\Psi_t(0) \cdot \Psi^2_t\p{|a-b|}, 
\quad \forall \ \bsig\in \{+,-\}^2, \ a,b\in \Zn.
\ee
In particular, if we take \(\Psi_t(|a-b|)=  (W^{-c_0}B_{t, |a-b|\wedge K})^{1/2}\) for a constant $c_0>0$ and some $0\le K\le L$, then 
\be\label{eq:LW_conclusion2}
\mathcal{E}^{\Gc,(2)}_{t, \bsig, (a,b)}\prec \frac{1}{\eta_t}\p{W^{-c_0}B_{t,0}}^{1/2}\cdot W^{-c_0} B_{t,|a-b|\wedge K}  .
\ee
\end{lemma}

\begin{lemma}[Light-weight estimate: $\cal T$-bound]\label{lem: EWGn2_N} %\label{lem:LWterm_exp}
Given any $t \in [0,1)$, assume that the bounds in \eqref{initialGT2} hold for some constant $\e_0>0$ and deterministic control parameter $W^{-d/2}\le \Psi_t\le W^{-\e_0}$. Suppose for some $0\le \ell \le (\log W)^{10}\ell_t$, the following estimate holds for any large constant $D>0$:
\be\label{eq:LW_assm_exp}
{\cal L}^{(2)}_{t,\bsig,(a,b)}\prec W^{-d}\wT_{t,D}^{\ell}(|a-b|),\quad \forall \ \bsig\in\{(+,-),(-,+)\}, \ a,b\in \Zn . \ee
Then, the following estimate holds for any large constant $D>0$: 
\be\label{eq:LW_conclusion_exp}
\mathcal{E}^{\Gc,(2)}_{t, \bsig, (a,b)} \prec \frac{1}{\eta_t}(W^{-d}B_{t,0})^{1/2} \cdot W^{-d} \wT_{t,D}^{\ell}(|a-b|)  , 
\quad \forall \ \bsig\in \{+,-\}^2, \ a,b\in \Zn.
\ee
\end{lemma}
\begin{remark}\label{rmk:poly_exp}
The second condition in \eqref{eq:Psi} implies that the parameter $\Psi_t(r)$ decays polynomially in $r$. In particular, the parameter \smash{$W^{-d}\wT_{t,D}^{\ell}(r)$} in \eqref{eq:LW_assm_exp} does not satisfy this condition because of the exponential factor \smash{$\exp\p{-(r/\ell_t)^{1/2}}$} in \eqref{defTUL}. Consequently, \Cref{lem: EWGn2_N} cannot be deduced directly from \Cref{lem:LWterm}. The corresponding proof is technically more delicate, since factors of the form $\Psi_t(c|a-b|)$, for some constant $c\in(0,1)$, can no longer be replaced by $\Psi_t(|a-b|)$ as in the polynomial-decay setting.
In \Cref{lem: EWGn2_N}, the assumption $\ell \le (\log W)^{10}\ell_t$ is made without loss of generality, since $\cal T_{t}(\ell) \le W^{-D}$ whenever $\ell > (\log W)^{10}\ell_t$. 
%Moreover, if the assumption \eqref{eq:LW_assm_exp} contains a prefactor $J\gg 1$, i.e., ${\cal L}^{(2)}_{t,\bsig,(a,b)}\prec J\cdot W^{-d}\wT_{t,D}^{\ell}(|a-b|)$, then we cannot derive a version of \eqref{eq:LW_conclusion_exp} with an additional prefactor $J^{3/2}$, as one may expect from the estimate \eqref{eq:LW_conclusion2}. 
%     Note that under the assumption \eqref{eq:LW_assm_exp}, we can take $\Phi_t=(W^{-d}B_{t,0})^{1/2}$. We keep the factor in \eqref{eq:LW_conclusion_exp} solely for easy comparison with the result \eqref{eq:LW_conclusion}. %Moreover, in our proof below, we will only use \Cref{lem:LW_moment_exp}, while \Cref{lem:LW_moment} will be used as an intermediate step in establishing \Cref{lem:LW_moment_exp}. We present \Cref{lem:LW_moment} separately because its proof illustrates the key idea concerning the graphical tools and nested properties.
\end{remark}

The proofs of \Cref{lem:LWterm,lem: EWGn2_N} are based on bounding the high moments using Gaussian integration by parts and some diagrammatic tools developed in \cite{yang2021delocalization,yang2021random}. We will present the details in \Cref{Sec:graph}.  
Finally, the martingale term is bounded as in the following lemma, whose proof will be postponed to 
\Cref{subsec:pf_lem: EMn2_N}. 

\begin{lemma}[Martingale estimate]\label{lem: EMn2_N} 
Given any $t\in[0,1)$, suppose the setting of \Cref{lem:LWterm} holds. Consider the term \smash{$\left( \cal E\otimes  \cal E \right)^{M}_{t, \bsig, \ba,\ba}\equiv \left( \cal E\otimes  \cal E \right)^{M,(2;k)}_{t, \bsig, \ba,\ba}$} defined in \eqref{defEOTE} for any \smash{$\ba=(a,b)\in (\Zn)^2$}, $\bsig=(\sig_1,\sig_2)\in \{+,-\}^2$, and $k\in\{1,2\}$.
First, the following estimate holds:
\be\label{eq:MG_conclusion}
\left( {\cal E}\otimes  {\cal E} \right)^{M}_{t, \bsig, \ba, \ba} \prec \frac{1}{\eta_t}\Psi_t(0) \cdot  \Psi_t^4\p{ |a-b|} .
\ee
In particular, if we take \(\Psi_t(|a-b|)= (W^{-c_0}B_{t, |a-b|\wedge K})^{1/2}\) for a constant $c_0>0$ and some $0\le K\le L$, then 
\be\label{eq:MG_conclusion2}
\left( \cal E\otimes  \cal E \right)^{M}_{t, \bsig, \ba, \ba} \prec \frac{1}{\eta_t}\p{W^{-c_0}B_{t,0}}^{1/2}\cdot  \p{W^{-c_0} B_{t,|a-b|\wedge K}}^2.
\ee
Second, suppose for some $0\le \ell\le (\log W)^{10}\ell_t$, the estimate \eqref{eq:LW_assm_exp} holds for any large constant $D>0$. Then, the following estimate holds for any large constant $D>0$:
\be\label{eq:MG_conclusion3}
(\mathcal{E} \otimes \mathcal{E})^{M}_{t,\bsig, \ba,\ba} \prec \frac{1}{\eta_t} \left[\left(W^{-d}B_{t,0}\right)^{1/2} +  \left(\wh{\mathcal{J}}_{t,D}^{\ell}\right)^3\right]\cdot \left(W^{-d} \wT^{\ell}_{t,D}(|a-b|)\right)^2.
\ee
% \be\label{eq:MG_conclusion3}
% (\mathcal{E} \otimes \mathcal{E})^{M}_{t,\bsig, \ba,\ba} \prec \frac{1}{\eta_t} \left[\left(W^{-d}B_{t,0}\right)^{1/2}+ \left(W^{-d}B_{t,0}\right)^{1/2}\left(\mathcal{J}_{t,D}^{\ell}\right)^2 +  \left(\mathcal{J}_{t,D}^{\ell}\right)^3\right]\cdot \left(W^{-d} \wT^{\ell}_{t,D}(|a-b|)\right)^2,
% \ee
%where the deterministic parameter $\mathcal{J}_{t,D}^{\ell}$ is defined in \eqref{defCALJ_deter}.
\end{lemma}
\begin{remark}\label{rmk_bottleneck}
We remark that the martingale estimate is the main bottleneck of our proof, preventing us from improving the exponent $1/5$ in \eqref{Eq:Gdecay_flow} to the optimal value $1$. This might be overcome by incorporating certain (not necessarily sharp) spatial decay properties of higher-order $G$-loops into the loop hierarchy analysis. However, such a refinement is far from a straightforward extension of the current argument. Even if achievable, it would entail substantial additional technical complexity while yielding only a ``marginal'' improvement in the quantum diffusion estimates. Since this refinement does not affect the main results---delocalization (\Cref{MR:decol}), local laws (\Cref{MR:locSC}), QUE (\Cref{MR:QUE}), and bulk universality (\Cref{Thm: B_Univ})---we leave this direction for future investigation.

%We remark that this martingale estimate is the bottleneck of the proof, preventing us from improving the exponent $1/5$ in \eqref{Eq:Gdecay_flow} to the optimal value $1$. This might be overcome by incorporating certain (not necessarily sharp) spatial decay properties of higher-order $\cL$-loops into the $G$-loop analysis. Yet such a refinement is far from being trivial extension of our current proof. Even if this is possible, it would introduce considerable additional technical complexity while yielding only a ``marginal'' improvement in the quantum diffusion estimates. Since this does not affect the main results---delocalization (\Cref{MR:decol}), local laws (\Cref{MR:locSC}), QUE (\Cref{MR:QUE}), and bulk universality (\Cref{Thm: B_Univ})---we leave this direction for future work.
\end{remark}

With the above Lemmas \ref{lem:newKLK}--\ref{lem: EMn2_N} as inputs, we are ready to control the terms in equation \eqref{LK_simple} and complete Step 2 of the proof of \Cref{lem:main_ind}. 

\begin{proof}[\bf Step 2: Proof of \eqref{Gt_bound_flow}--\eqref{Eq:Gdecay_w}] 
Suppose the estimate \eqref{Eq:Gdecay_w} has been established. Then, using \eqref{Kn2sol} and \eqref{prop:ThfadC}, we obtain 
\be\label{eq:kn2sol_decay}
{\cal K}^{(2)}_{u, (-,+), (a_1,a_2)}\prec W^{-d}B_{u,|a_1-a_2|}, \quad \forall a_1,a_2\in \Zn.
\ee
Together with \eqref{Eq:Gdecay_w}, this implies---provided $\fc_d$ in \eqref{con_st_ind} is chosen sufficiently small---the 2-$G$-loop bound
\be\label{eq:L2_decay}
{\cal L}^{(2)}_{u, (-,+), (a_1,a_2)}\prec W^{-d}B_{u,|a_1-a_2|}, \quad \forall a_1,a_2\in \Zn.
\ee
Applying the weak local law \eqref{Gtmwc} from Step 1 to verify the first condition in \eqref{initialGT2}, and then invoking \Cref{lem_GbEXP} with the control parameter $\Psi_u:=\max\{W^{-d/2},(W^{-d}B_{u,0})^{1/2}\}$, we obtain the entrywise local law \eqref{Gt_bound_flow} and the averaged local law \eqref{Gt_avgbound_flow}. Here, the maximum in the definition of $\Psi_u$ only ensures the lower bound $\Psi_u\ge W^{-d/2}$ required in \Cref{lem_GbEXP}: since $B_{u,0}\ge(\ilambda^2+1)^{-1}\ge(\fd^{-2}+1)^{-1}$ by \eqref{eq:WO}, $\Psi_u$ is comparable to $(W^{-d}B_{u,0})^{1/2}$, so the second condition in \eqref{initialGT2} follows from \eqref{eq:L2_decay}, and the resulting bounds are the same as with $\Psi_u=(W^{-d}B_{u,0})^{1/2}$.

It remains to establish the estimate \eqref{Eq:Gdecay_w}. From the bound \eqref{lRB1} proved in Step 1, we know 
\begin{equation} \label{lokis2}
\max_{\bsig,{\ba}}\abs{{\cal L}^{(2)}_{u,\bsig,{\ba}}}\prec \frac{1-s}{1-t}(W^{-d}B_{t,0}) =: W^{-c_0},\quad \forall u\in [s,t],
\end{equation}
where $c_0\equiv c_0(W)\gtrsim 1$ under the condition \eqref{con_st_ind}. We first prove the following maximum bound:
\be\label{eq:opt_L2}
\max_{\bsig,{\ba}}\abs{\left({\cal L}-{\cal K}\right)^{(2)}_{u,\bsig,{\ba}}}\prec  W^{-d}B_{u,0},\quad \forall u\in [s,t].
\ee
Using \Cref{lem:newKLK} (for the case $\ell=0$), \Cref{lem:LWterm}, and the inductive hypothesis \eqref{Eq:L-KGt+IND} with $n=2$, we obtain from the flow \eqref{LK_simple} (with $\tau=t$) that 
\begin{align}\label{l>0EQ}
\max_{\bsig,\ba}\left|\left({\cal L}-{\cal K}\right)^{(2)}_{t,\bsig, {\ba}}\right|
\le \int_{s}^t \frac{C_0}{1-u}\max_{\bsig,\ba}\left|\left({\cal L}-{\cal K}\right)^{(2)}_{u,\bsig,\ba}\right| \dd u  +\OO_\prec \p{(W^{-d}B_{s,0})^2+W^{-\frac32c_0}} + \abs{\int_s^t \dd \mathcal{E}^{M}_{u, \bsig, {\ba}}}   
\end{align}
for a constant $C_0>0$ depending only on $c_d$ and $C_d$ in \eqref{prop:ThfadC}. Combining the estimate \eqref{eq:MG_conclusion} in \Cref{lem: EMn2_N} with \Cref{lem:DIfREP}, and applying Markov's inequality, we obtain that 
\[\abs{\int_s^t \dd \mathcal{E}^{M}_{u, \bsig, {\ba}}} \prec \p{\int_{s}^{t}\frac{W^{-5c_0/2}}{\eta_u}\dd u}^{1/2}\prec W^{-\frac54c_0}.\]
Together with \eqref{l>0EQ}, it gives
\be\label{eq:Gronwall_2L_max}\max_{\bsig,{\ba}}\left|\left({\cal L}-{\cal K}\right)^{(2)}_{t, \bsig, {\ba} }\right|\le \int_{s}^t \frac{C_0}{1-u}\max_{\bsig,\ba}\left|\left({\cal L}-{\cal K}\right)^{(2)}_{u,\bsig,\ba}\right| \dd u  +\OO_\prec \p{(W^{-d}B_{s,0})^2+W^{-\frac54c_0 }} .
\ee
% Choose $\ell=0$ in lemma \ref{lem:newKLK}, then 
% \begin{align}\label{l>0EQ}
%   \Big|\,\Big[\mathcal{K} \sim (\mathcal{L} - \mathcal{K})\Big]^{l_\mathcal{K}=2}_{u, \bsig, {\ba}} \,\Big| \;+ \; 
%    \Big|\mathcal{E}^{((\mathcal{L} - \mathcal{K}) \times (\mathcal{L} - \mathcal{K}))}_{u, \bsig, {\ba}}  \Big| \le \frac{C}{\eta_u}\max_{{\ba}}\left|\left({\cal L}-{\cal K}\right)^{(2)}_{u, {\ba} }\right|
%     \end{align}
% It implies the following self bounded integral form estimate
% $$\max_{{\ba}}\left|\left({\cal L}-{\cal K}\right)^{(2)}_{t, {\ba} }\right|\le 
% O (W^{-d}B_{t,0})^2 +O_\prec (W^{-5c_0/4})+\int_{s}^t \frac{C}{\eta_u}\max_{{\ba}}\left|\left({\cal L}-{\cal K}\right)^{(2)}_{u, {\ba} }\right| du 
% $$
Recall the following classical Gr{\"o}nwall's inequality, where $\beta$ is non-negative and $\al$ is non-decreasing:
\be\label{Gronwall_inequality}
\begin{aligned}
&f(t) \leq \alpha(t)+\int_s^t \beta(u) f(u) \mathrm{d} u \implies f(t) \leq \alpha(t) \exp \left(\int_s^t \beta(u) \mathrm{d} u\right).
% u(t) \leq \alpha(t)+\int_a^t \alpha(s) \beta(s) \exp \left(\int_s^t \beta(r) \mathrm{d} r\right) \mathrm{d} s 
\end{aligned}
\ee
%  If, in addition, the function $\alpha$ is non-decreasing, then
% $$
% u(t) \leq \alpha(t) \exp \left(\int_a^t \beta(s) \mathrm{d} s\right)
% $$
Applying it to \eqref{eq:Gronwall_2L_max} yields that  
\be\label{eq:L-K2max} \max_{\bsig,\ba}\left|\left({\cal L}-{\cal K}\right)^{(2)}_{t,\bsig,\ba}\right|\prec 
 \p{\frac{1-s}{1-t}}^{C_0}\left((W^{-d}B_{t,0})^2 + W^{-\frac54c_0}\right)\le W^{-d}B_{t,0},
\ee
where, in the last step, we use that 
\[ \p{\frac{1-s}{1-t}}^{C_0}W^{-\frac54c_0} \le W^{-d}B_{t,0} \cdot \p{\frac{1-s}{1-t}}^{C_0+5/4}\p{W^{-d}B_{t,0}}^{1/4} \ll W^{-d}B_{t,0} \]
as long as we choose $\fc_d$ in \eqref{con_st_ind} sufficiently small depending on $C_0$. The last bound in \eqref{eq:L-K2max} is precisely \eqref{eq:opt_L2} at $u=t$. The same argument applies at each fixed $u\in [s,t]$, and a standard $N^{-C}$-net and perturbation argument extends it uniformly to all $u\in[s,t]$, which concludes \eqref{eq:opt_L2}.


We can obtain from \eqref{eq:opt_L2} that for $\ell^{(0)}=0$ and any large constant $D>0$,
\begin{align}\label{ksjjuw}
 %\cal J_{u,D}^{\ell^{(0)}}\prec 1 \implies   
 \max_{\ba=(a,b)}\max_{\bsig}\big|\cal L^{(2)}_{u, \bsig, {\ba} }\big|\Big/\p{W^{-d} \widetilde{\cal {T}}^{\ell^{(0)}}_{u,D}\p{|a-b|}} \prec 1, \quad \forall u\in[s,t],%, \quad \ell^{(0)} =0,\quad {\ba}=(a_1,a_2)
\end{align}
where we also use that 
\be\label{eq:simpleboundK}\cK^{(2)}_{u,\bsig,\ba}\prec W^{-d}\wT^{L}_{u,D}\p{|a-b|}\ee 
by \eqref{prop:ThfadC} and the definition of 2-$\cK$-loop in \eqref{Kn2sol}. 
Now, fix a large constant $D>0$. Suppose we have the following bound for a collection of length scales $\{K_u\ge 0:u\in [s,t]\}$:
\begin{align}\label{kwr3juw}
\max_{\ba=(a,b)}\max_{\bsig}\big|{\cal L}^{(2)}_{u,\bsig,{\ba} }\big|\Big/\p{W^{-d}\widetilde{\cal {T}}^{K_u}_{u,D}\p{|a-b|}}\prec 1,\quad \forall u\in[s,t]. 
\end{align}
Moreover, assume that $\cal T_u(K_u)$ is non-decreasing in $u$, i.e., 
\be\label{eq:monotone_Ku}
\cal T_{u}(K_u) \le \cal T_v(K_v), \quad  \wT_{u,D}^{K_u}(r) \le \wT_{v,D}^{K_v}(r),\quad \forall s\le u\le v \le t, \ r\ge 0.
\ee
Our goal is to establish the following estimate for a constant $C_0>0$ depending only on $c_d$ and $C_d$ in \eqref{prop:ThfadC}:
\be\label{eq:Gronwall_dervJuD}
\wh{\mathcal{J}}_{u,D}^{K_u} \prec\left(|1-s|/|1-u|\right)^{C_0}(W^{-d} {B} _{t,0})^{1/5}  ,\quad \forall u\in [s,\tau],  
\ee
where $\tau$ is a stopping time defined as
\be\label{eq:def2_stopping} 
\tau :=t\wedge T,\quad \text{with}\quad T:=\inf\left\{u\ge s: \wh{\cal J}_{u,D}^{K_u} \ge (W^{-d}B_{u,0})^{1/6}\right\}.\ee


For the proof of \eqref{eq:Gronwall_dervJuD}, using \Cref{lem: EWGn2_N}, we obtain that 
\be\label{eq:boundGcterm}
\int_{s}^{\tau}\mathcal{E}^{\Gc}_{u, \bsig, {\ba}}\dd u \prec \int_{s}^\tau \frac{1}{\eta_u} (W^{-d}B_{u,0})^{1/2} \cdot \p{W^{-d} \wT_{u,D}^{K_u}(|a-b|)}\dd u  \prec (W^{-d}B_{\tau,0})^{1/2} \cdot \p{W^{-d} \wT_{\tau,D}^{K_\tau}(|a-b|)},
\ee
where, in the second step, we also use the fact that $B_{u,0}$ and $\wT_{u,D}^{K_u}(r)$ are monotonically increasing in $u$. Similarly, using \Cref{lem: EMn2_N} along with \Cref{lem:DIfREP} and the definition of the stopping time $\tau$, we obtain that 
\be\label{eq:boundmgterm}
\begin{split}
\int_s^\tau \dd \mathcal{E}^{M}_{u, \bsig, {\ba}} &\prec  \Big[\left(W^{-d}B_{\tau,0}\right)^{1/4}  + \sup_{u\in[s,\tau]}\p{\wh{\cal J}_{u,D}^{K_u}}^{3/2}\Big]\cdot \left(W^{-d} \wT^{K_\tau}_{\tau,D}(|a-b|)\right) \\
&\lesssim \left(W^{-d}B_{\tau,0}\right)^{1/4} \cdot \left(W^{-d} \wT^{K_\tau}_{\tau,D}(|a-b|)\right).\end{split}
\ee
%where to get the second inequality, we choose $\cal J_{u,D}^{K_u}\le (W^{-d}B_{u,0})^{1/6}$ by the definition of $T$ in \eqref{eq:def2_stopping}. 
Applying \Cref{lem:newKLK}, \eqref{eq:boundGcterm}, and \eqref{eq:boundmgterm} to equation \eqref{LK_simple} yields that 
\begin{align}\label{LKu2p2mmvj}
    \left|\left({\cal L}-{\cal K}\right)^{(2)}_{\tau,\bsig,\ba}\right| \le &~
\left|\left({\cal L}-{\cal K}\right)^{(2)}_{s, \bsig,\ba}\right|+ \int_s^\tau  \frac{C_0}{1-u} \wh{\mathcal{J}}^{K_u}_{u,D} \cdot W^{-d} \widetilde{{\cal T}}^{K_u} _{u,D}\p{|a-b|} \dd u  
\\\nonumber
&~+\OO_\prec\p{\p{W^{-d} {B}_{\tau,0}}^{1/4}\cdot W^{-d} \widetilde{{\cal T}}^{K_\tau}_{\tau,D}\p{|a-b|}}.
\end{align}
%for a constant $C_0>0$ depending only on $d$. 
From this equation, using the induction hypothesis \eqref{Eq:Gdecay+IND} at time $s$ and the definition \eqref{defCALJ}, we obtain that 
  $$
  \wh{\mathcal{J}}^{K_\tau}_{\tau,D}\le  \int_s^\tau \frac{C_0}{1-u} \wh{\mathcal{J}}^{K_u}_{u,D}  \dd u + \OO_\prec \p{(W^{-d} {B} _{\tau,0})^{1/5}}.
  $$
Applying Gr{\"o}nwall's inequality \eqref{Gronwall_inequality} again to this equation, we derive \eqref{eq:Gronwall_dervJuD}.
% that 
% \be\label{eq:Gronwall_dervJuD}
% \wh{\mathcal{J}}_{u,D}^{K_u} \prec\left(\frac{1-s}{1-u}\right)^{C_0}(W^{-d} {B} _{t,0})^{1/5} \ll (W^{-d} {B} _{u,0})^{1/6},\quad \forall u\in [s,\tau],  
% \ee


 
By the induction hypothesis \eqref{Eq:Gdecay+IND} at time $s$, the stopping time $T$ defined in \eqref{eq:def2_stopping} satisfies $T>s$ with high probability. In \eqref{eq:Gronwall_dervJuD}, using again \eqref{con_st_ind} with $\fc_d$ chosen sufficiently small depending on $C_0$, we can ensure that with high probability, \smash{\(\wh{\mathcal{J}}_{u,D}^{K_u} \ll (W^{-d} {B} _{u,0})^{1/6}\)} for all \(u\in [s,\tau]\). This implies that $T\ge t$ with high probability, so that \eqref{eq:Gronwall_dervJuD} holds for all $u\in[s,t]$. 
%Then, the estimate \eqref{eq:Gronwall_dervJuD} shows that we must have $T\ge t$ with high probability, so \eqref{eq:Gronwall_dervJuD} holds for all $u\in[s,t]$. 
Next, we define a new collection of parameters $\{K_u'\ge 0:u\in [s,t]\}$ as follows: for each $u\in[s,t]$, let $K_u'$ be the unique positive solution to  
\be\label{eq:def_ell1} 
\cal T_u(K_u')=\cal T_u(K_u)\cdot (W^{-d} {B}_{u,0})^{1/6} . \ee
Under this definition, the monotonicity relation \eqref{eq:monotone_Ku} continues to hold for the family $\{K_u'\}_{u\in[s,t]}$. Furthermore, from \eqref{eq:simpleboundK} and \eqref{eq:Gronwall_dervJuD}, we obtain
\[ {\cal L}^{(2)}_{u,\bsig, \ba}\prec W^{-d} \p{\widetilde{{\cal T}}^{K_u}_{u,D}(|a-b|)\cdot (W^{-d} {B} _{u,0})^{1/6} + \wT_{u,D}^{L}(|a-b|)} \lesssim W^{-d}\widetilde{{\cal T}}^{K_u'}_{u,D}(|a-b|),\]
where, in the second step, we use \eqref{eq:def_ell1} along with the relation $\widetilde{{\cal T}}^{K_u}_{u,D}(|a-b|)\asymp \widetilde{{\cal T}}^{L}_{u,D}(|a-b|)$ for $K_u\ge L$. We now use this as the input (with $K_u$ replaced by $K_u'$) in \eqref{kwr3juw} for the next iteration of the above argument, and show that \eqref{kwr3juw} continues to hold at an even larger scale $K_u''$. From \eqref{eq:def_ell1}, it follows that, for any fixed $D>0$, after performing at most $\OO(1)$ iterations (where the number of iterations depends on $D$), we reach scales $\{K_u:u\in[s,t]\}$ such that $\cal T_u(K_u)\le W^{-D}$ for all $u\in[s,t]$. Consequently, we have %\textcolor{blue}{s.t $A\implies B$  is not misleading. s.t $A$, which implies $B$.}
\[\wT_{u,D}^{K_u}(|a-b|)\asymp \wT_{u,D}^{L}(|a-b|), \quad \forall u\in [s,t], \ a,b\in \Zn.\]
Finally, applying \eqref{eq:Gronwall_dervJuD} once more yields \eqref{Eq:Gdecay_w}.
%Now, applying \eqref{eq:Gronwall_dervJuD} again concludes \eqref{Eq:Gdecay_w}. 
\end{proof}

\begin{remark}
We briefly explain why the stability estimate \eqref{Eq:Gdecay_w} \emph{cannot} be established continuously along the flow using Duhamel's formula, as was done in the $d \in \{1,2\}$ case \cite{YY_25,DYYY25,erdHos2025zigzag}.
For a large constant $D>0$, these works consider the random control parameter \smash{$J_{u,D}\equiv \wh{\cal J}^{L}_{u,D}$} (defined in \eqref{defCALJ}) at the full-system scale $\ell=L$.   
% for $(\cL-\cK)^{(2)}$-loops:
% \be\label{eq:defJuD}
% J_{u,D}:=\max_{\ba=(a,b)}\max_{\bsig} \big|(\cL-\cK)^{(2)}_{u,\bsig,\ba}\big|\big/\big[W^{-d}\cal T_{u}(|a-b|) +W^{-D}\big],
% \ee
% which is of the same order as $\wh{\cal J}^{L}_{u,D}$ defined in \eqref{defCALJ}. 
Roughly speaking, the previous approach begins with an initial high-probability bound $J_{s,D}\le W^{-c}$ (for some constant $c>0$) at $u=s$ and then propagates this bound along a $W^{-C}$-net of $[s,t]$. The essential idea is that if $J_{u,D}\le W^{-c}$ holds, a perturbative argument yields a slightly weaker bound $J_{u',D}\le W^{-c+\e}$ at $u'=u+W^{-C}$. This can then be bootstrapped back to $J_{u',D}\le W^{-c}$ using the loop hierarchy \eqref{int_K-LcalE}, at the cost of a probability loss $W^{-D}$, which remains acceptable since the total accumulation over the net is only $W^{-D+C}$.
However, this mechanism breaks down once we attempt to incorporate the light-weight estimates from \Cref{lem:LWterm,lem: EWGn2_N}. They rely on a moment method that degrades the probability bound much more severely---from $W^{-D}$ in the original assumptions \eqref{eq:LW_assm} (resp.~\eqref{eq:LW_assm_exp}) to $W^{-\delta D}$ in \eqref{eq:LW_conclusion} (resp.~\eqref{eq:LW_conclusion_exp}) for some constant $\delta <1$. As a result, \Cref{lem:LWterm,lem: EWGn2_N} can each be applied only $\OO(1)$ times, and thus the flow-based propagation strategy is no longer applicable. 
Instead, we must prove the light-weight estimate simultaneously for all $u\in[s,t]$, starting from the weak local law \eqref{Gtmwc} established in Step 1. This necessitates the self-improving argument developed in \eqref{awi2iks}. 

%This necessity leads to the development of a self-improving argument as in \eqref{awi2iks}. 
%: starting from the weak a priori bound \eqref{eq:condon2loops}, we iteratively upgrade it to the sharper bound \eqref{eq:condon2loops_improve}, improving the control parameters $\Psi_t(\cdot)$ (inherited from Step 1, cf.~\eqref{lRB1}) until they match the sharp one (i.e., the denominator in \eqref{eq:defJuD}) after only $\OO(1)$ iterations. 

Moreover, since our initial control of $J_{u,D}$ is too weak (it can be as large as $W^D$), the quadratic term \smash{$\cal E^{(\cL-\cK)\times (\cL-\cK)}$} cannot be effectively bounded by $J_{u,D}^2$, unlike in \cite{YY_25,DYYY25,erdHos2025zigzag}. We must therefore rely on the bound \eqref{juwo=Lklk}. However, this estimate invalidates the Duhamel-based argument used in dimensions $d\in\{1,2\}$: when inserted into Duhamel's formula, it contributes a term of size \smash{$(|1-s|/|1-t|)\cdot W^{-d}\widetilde{{\cal T}}^{\ell}_{u,D}\p{|a-b|}$}, which already breaks the first step of \eqref{awi2iks}. As a result, we must treat \smash{$\cal E^{(\cL-\cK)\times (\cL-\cK)}$} as a main term rather than an error term, and reorganize the proof around a new Gr{\"o}nwall-type argument.

\end{remark}

% Starting from the scale $\ell_1$ in \eqref{eq:endofiteration1}, we assume a slightly different bound from that in \eqref{kwr3juw}: for some parameter $\al \ge 0$,
% \begin{align}\label{kwr3juw_type2}
% \max_{\bsig,\ba=(a,b)}\big|{\cal L}^{(2)}_{u,\bsig,{\ba} }\big|\Big/\p{W^{-d}\widetilde{\cal {T}}^{\ell_1+\al(\log W)^2\ell_u}_{u,D}\p{|a-b|}}\prec 1,\quad \forall u\in[s,t]. 
% \end{align}
% in which case we have  
% \[\widetilde{{\cal T}}^{\ell}_{u,D}(|a-b|) \asymp \widetilde{{\cal T}}^{L}_{u,D}(|a-b|) \quad \forall a,b\in \Zn. \]




\subsection{Proof of Lemma \ref{lem:newKLK}}\label{subsec:pf_lem:newKLK}

Without loss of generality, it suffices to consider the case where 
\begin{align}\label{uwu3po2[]}
     {\cal T}_{u}\p{\ell}\ge  W^{-D}. 
\end{align} 
Otherwise, we can find $K\le \ell$ such that ${\cal T}_{u}\p{K} = W^{-D}$, for which we always have \(\widetilde{{\cal T}}^{\ell}_{u,D}\p{|a-b|}=\widetilde{{\cal T}}^{K}_{u,D}\p{|a-b|}\) by definition \eqref{defWTTlD}. 
For the proof of \eqref{juwo2=klk}, by the definitions \eqref{def:op_thn}, \eqref{defWTTlD}, and \eqref{defCALJ}, we have that 
\begin{align} \label{uu2j2ois} W^d\abs{\p{\varTheta^{(2)}_{u,\bsig}\circ({\cal L}-{\cal K})^{(2)}_{u,\bsig}}_{{\ba}}}&\; \le \;C \wh{\cal J}^{\ell}_{u,D}\sum_c \big|\big(\Theta_{u}^{(\sig_1,\sig_2)}M^{(\sig_1,\sig_2)}\big)_{ac}\big|\cdot \widetilde{{\cal T}}^{\ell}_{u,D}\p{|c-b|} \nonumber \\
&\; \le \;  C \wh{\cal J}^{\ell}_{u,D}\sum_{c:|c-b|\ge \ell} \big|\Theta_{u,ac}^{(\sig_1,\sig_2)}\big|\cdot {\cal T}_{u}\p{\ell} + C \wh{\cal J}^{\ell}_{u,D}\sum_{c:|c-b|< \ell} \cal T_{u}(|a-c|){\cal T}_{u}\p{|c-b|} \nonumber\\
&\; \le \;  \frac{C}{1-u} \wh{\cal J}^{\ell}_{u,D}\br{{\cal T}_{u}\p{\ell} +\cal T_{u}(|a-b|)} \le \frac{C}{1-u} \wh{\cal J}^{\ell}_{u,D}\wT_{u,D}^{\ell}(|a-b|),
   \end{align} 
where, in the second step, we use the estimate \eqref{prop:ThfadC}, and the facts that under the condition \eqref{uwu3po2[]}, \be\label{eq;facts}\widetilde{{\cal T}}^{\ell}_{u,D}\p{|c-b|}\le {{\cal T}} _{u}\p{\ell}\ \ \text{for}\ \ |c-b|\ge \ell,\quad \widetilde{{\cal T}}^{\ell}_{u,D}\p{|c-b|}\le {{\cal T}} _{u}\p{|c-b|}\ \ \text{for}\ \ |c-b|<\ell,
\ee
and in the third step, we use \eqref{eq:THETAinftinf} and \eqref{TTT2}. This concludes \eqref{juwo2=klk}.


The proof of \eqref{juwo=Lklk} follows a similar argument. We have that for any $\bsig_1,\bsig_2\in\{+,-\}^2$,  
\begin{align}\label{i2kk2zgg}
&~ W^d \sum_{c_1,c_2} \left|  \p{\cal L-\cal K}_{u,\bsig_1,(a,c_1)}^{(2)} S^{(\sB)}_{c_1c_2} \p{\cal L-\cal K}^{(2)}_{u,\bsig_2,(c_2,b)}\right|\\
\;\le \;&~  \sum_{c } \left(|{\cal L}^{(2)}_{u,\bsig_1,(a,c)}|+|{\cal K}^{(2)}_{u,\bsig_1,(a,c)}|  \right)\cdot 
\max_{c':|c'-b|\ge (\ell-1)_+}
\wh{\cal J}^{\ell}_{u,D} \widetilde{\mathcal{T}}_{u,D}^{\ell}\p{|c'-b|}\nonumber \\ \nonumber
\;+ \;&~ \sum_{c } \left(|{\cal L}^{(2)}_{u,\bsig_2,(c,b)}|+|{\cal K}^{(2)}_{u,\bsig_2,(c,b)}|  \right)\cdot 
\max_{c':|c'-a|\ge  (\ell - 1)_+}
\wh{\cal J}^{\ell}_{u,D} \widetilde{\mathcal{T}}_{u,D}^{\ell}\p{|a-c'|} 
\\\nonumber
\;+ \;&~ \mathbf 1_{\ell\ge 1}\cdot W^{-d}\left(\wh{\cal J}^{\ell}_{u,D}\right)^2
 \sum_{c:|c-a|\vee |c-b| < \ell} \widetilde{\mathcal{T}}_{u,D}^{\ell}\p{|a-c|} \widetilde{\mathcal{T}}_{u,D}^{\ell}\p{|c-b|} .\nonumber
\end{align}
For the first term on the RHS of \eqref{i2kk2zgg}, applying the Cauchy--Schwarz inequality and Ward's identities \eqref{WI_calL} and \eqref{WI_calK}, we get that with high probability,
\begin{align*}
    \sum_{c } \left(|{\cal L}^{(2)}_{u,\bsig_1,(a,c)}|+|{\cal K}^{(2)}_{u,\bsig_1,(a,c)}|  \right) &\le 
\sum_{c } \p{ {\cal L}^{(2)}_{u,(-,+),(a,c)}+{\cal K}^{(2)}_{u,(-,+),(a,c)}}\\
&=\frac{\im m+ \im \tr\p{G_uE_a}}{W^{d} \eta_u} =\frac{1+\oo(1)}{W^{d} (1-u)}, 
\end{align*}
where we use the weak local law \eqref{Gtmwc} and the relation \eqref{eta} in the last step. With this estimate, we can bound the first term on the RHS of \eqref{i2kk2zgg} as follows with high probability:
$$
\frac{1+\oo(1)}{W^d(1-u)} \max_{c':|c'-b|\ge (\ell-1)_+}
\wh{\cal J}^{\ell}_{u,D} \widetilde{\mathcal{T}}_{u,D}^{\ell}\p{|c'-b|} \le \frac{C}{W^d(1-u)} \wh{\cal J}^{\ell}_{u,D} {\mathcal{T}}_{u}\p{\ell} ,
$$
where we have used \eqref{eq;facts} and ${\mathcal{T}}_{u}\p{K+1}\asymp {\mathcal{T}}_{u}\p{K}$ for any $K\ge 0$. The second term on the RHS of \eqref{i2kk2zgg} can be bounded in the same way. For the last term on the RHS of \eqref{i2kk2zgg}, using \eqref{eq;facts}, we bound it by 
\be\label{eq:i2kk2zgglast}\mathbf 1_{\ell\ge 1} \cdot CW^{-d}\left(\wh{\cal J}^{\ell}_{u,D}\right)^2
 \sum_{c}{\mathcal{T}}_{u}\p{|a-c|}  {\mathcal{T}}_{u}\p{|c-b|} \le \mathbf 1_{\ell\ge 1} \cdot \frac{C}{W^d(1-u)}\left(\wh{\cal J}^{\ell}_{u,D}\right)^2 \cal T_u(|a-b|),\ee
where we use \eqref{TTT2} in the second step. Combining the above two bounds, we conclude \eqref{juwo=Lklk}.

% Then with assumption \eqref{Gtmwc}, the 2nd line of \eqref{i2kk2zgg} is bounded by
% $$
% \le \frac{C\cdot{\bf 1}_{l<d\cdot L}}{\eta_u} {\cal J}^{\ell,D}_u \cdot \widetilde{\mathcal{T}}_{u,\ell }^{\ell, D}
% \le \frac{C\cdot{\bf 1}_{l<d\cdot L}}{\eta_u} {\cal J}^{\ell,D}_u \cdot \widetilde{\mathcal{T}}_{u,|a-b| }^{\ell, D}
% $$
% (Here we add ${\bf 1}_{l<d\cdot L}$, since if $\ell\ge dL$ then $|a-b|<dL\le \ell $ in $\mathbb Z_L^d$, which implies that the middle line of the \eqref{i2kk2zgg} is zero.) For the   last line of \eqref{i2kk2zgg}, using \eqref{cbell} and lemma   \ref{lem:propT}, it is bounded by 
% $$
% \frac{C}{\eta_u}\left( {\cal J}^{\ell,D}_u \right)^2\cdot \widetilde{\mathcal{T}}_{u,|a-b| }^{\ell, D}\cdot {\bf 1}_{\ell>0}
% $$
% Insert these two bounds back to \eqref{i2kk2zgg}, we obtain \eqref{juwo=Lklk} and complete the proof of lemma \ref{lem:newKLK}. 

%\end{proof}

\subsection{Proof of Lemma \ref{lem: EMn2_N}}\label{subsec:pf_lem: EMn2_N}
%\begin{proof}[Proof of lemma \ref{lem: EMn2_N}]

We only control the term \smash{\((\mathcal{E} \otimes \mathcal{E})_{t, \bsig, \boldsymbol{a}, \boldsymbol{a}}^M\equiv (\mathcal{E} \otimes \mathcal{E})_{t, \bsig, \boldsymbol{a}, \boldsymbol{a}}^{M,(2;1)}\)} without loss of generality. By definition \eqref{defEOTE}, we can write it as follows with $\ba(c,c')=(a,b,c',b,a,c)$ and $(\bsig\otimes \overline\bsig)^{(1)} = (\sigma_1, \sigma_2, \sigma_1, -{\sigma}_1, -{\sigma}_2, -{\sigma}_1)$:
\begin{equation}\label{eq_mart-sum}
    (\mathcal{E} \otimes \mathcal{E})_{t, \bsig, \boldsymbol{a}, \boldsymbol{a}}^M = W^d \sum_{c, c'}S^{(\sB)}_{cc'} \mathcal{L}^{(6)}_{t, (\bsig\otimes \overline\bsig)^{(1)}, \ba(c,c')}.
\end{equation}
We can represent the RHS using the following graph: 
\begin{equation}\label{eq:2martingale_quad}
\parbox[c]{0.3\linewidth}{
\scalebox{0.95}{
\begin{tikzpicture}[scale=1.2, every node/.style={font=\small}]

% Left graph
\coordinate (a1) at (0, 1.2);
\fill (a1) circle (2pt) node[above=2pt]{$a$};
%\node[above left] at (a1) {$a_1$};

\coordinate (a2) at (2, 1.2);
\fill (a2) circle (2pt) node[above=2pt]{$b$};

\coordinate (b) at (0.75, 0.6);
\fill (b) circle (2pt) node[left=2pt]{$c$};

\coordinate (b') at (1.25, 0.6);
\fill (b') circle (2pt) node[right=2pt]{$c'$};

\coordinate (a1') at (0, 0);
\fill (a1') circle (2pt) node[below=2pt]{$a$};

\coordinate (a2') at (2, 0);
\fill (a2') circle (2pt) node[below=2pt]{$b$};

\node at (-1,0.6) {$W^d\sum_{c\sim c'}S_{cc'}^{(\sB)}$};

\draw[thick] (a1) -- (b) node[midway, left] {$\sig_1$};
 \draw[thick] (b') -- (a2) node[midway, right] {$\sig_1$};
\draw[thick] (b') -- (a2') node[midway, right] {$-\sig_1$};
 \draw[thick] (b) -- (a1') node[midway, left] {$-\sig_1$};
\draw[thick] (a1) -- (a2) node[midway, above] {$\sig_2$};
\draw[thick] (a1') -- (a2') node[midway, below] {$-\sig_2$};
% \draw[decorate, decoration={snake, amplitude=0.4mm, segment length=2mm}, thick] (b) -- (b')
% node[midway, above] {$S^{(\sB)}$};

% Right graph
% \node (a1r) at (6, 2) {$\bullet$};
% \node[above left] at (a1r) {$[a_1]$};

% \node (a2r) at (10, 2) {$\bullet$};
% \node[above right] at (a2r) {$[a_2]$};

% \node (br) at (7.5, 1) {$\bullet$};
% \node[below] at (br) {$[b]$};

% \node (b'r) at (8.5, 1) {$\bullet$};
% \node[below] at (b'r) {$[b]$};

% \node (a1r') at (6, 0) {$\bullet$};
% \node[below left] at (a1r') {$[a_1']$};

% \node (a2r') at (10, 0) {$\bullet$};
% \node[below right] at (a2r') {$[a_2']$};

% \draw[thick] (a1r) -- (br) node[midway, left] {$-$};
%  \draw[thick] (b'r) -- (a2r) node[midway, right] {$-$};
% \draw[thick] (b'r) -- (a2r') node[midway, right] {$+$};
%  \draw[thick] (br) -- (a1r') node[midway, left] {$+$};
% \draw[thick] (a1r) -- (a2r) node[midway, above] {$+$};
% \draw[thick] (a1r') -- (a2r') node[midway, below] {$-$};
% \draw[decorate, decoration={snake, amplitude=0.4mm, segment length=2mm}, thick] (br) -- (b'r)
%         node[midway, above] {$S^{(B)}$};
        
\end{tikzpicture}}
}
\end{equation}
Each solid edge between two block vertices (e.g., $a$ and $b$) with charge $\sig\in\{+,-\}$ denotes a resolvent entry $(G_t(\sig))_{xy}$ or $(G_t(\sig))_{yx}$ with $x\in [a]$ and $y\in [b]$. The orientation of the graph is assumed to be clockwise.
To illustrate the idea for the proof of \eqref{eq:MG_conclusion}, assume without loss of generality that the vertex $c$ is closer to $b$ than to $a$. Then, the graph in \eqref{eq:2martingale_quad} already contains four ``long legs'': in addition to the two edges connecting $a$ and $b$, the two edges connecting $a$ and $c$ contribute a factor $\Psi_t^2(|a-c|) \le \Psi_t^2(|a-b|/2) \lesssim \Psi_t^2(|a-b|)$. Summing over the two remaining edges connected to $c'$ and applying Cauchy--Schwarz together with Ward's identity yields an additional factor of $\eta_t^{-1}$.
However, compared with the target estimate \eqref{eq:MG_conclusion}, we still miss the factor $\Psi_t(0)$. To recover it, we require a sharper treatment of the 6-$G$-loop that avoids applying Cauchy--Schwarz directly to the two edges connected to $c'$. Indeed, if we can directly perform the global summation over these two edges without taking absolute values, Ward's identity \eqref{eq_Ward0} produces a solid edge, which contributes the desired small factor $\Psi_t(0)$.
This is precisely achieved via the following \emph{$\contract$ inequality} for the 6-$G$-loop in \eqref{eq:2martingale_quad}, which can be regarded as a pointwise version of the more general $\contract$ inequality \eqref{u2jzooi-2} below.
For simplicity, we abbreviate $G\equiv G_t$ in the following proof. 


\begin{lemma} [$\Contract$ inequality]\label{ygdhmsgq0}
For any subset $\cal A\subset \Zn$, we have that 
%The symmetric 6-loop in \eqref{eq_mart-sum} satisfies the bound 
\begin{equation}\label{eq_sym_loop_bound}
\sum_{c'\in \cal A}\sum_{c\sim c'} \mathcal{L}^{(6)}_{t, (\bsig\otimes \overline\bsig)^{(1)}, \ba(c,c')} \lesssim \frac{1}{W^d\eta_t}\max_{c'\in \cal A}\left(\mathcal{L}^{(4)}_{t,\bsig^{(\alt)}, (c',b,c',b)}\right)^{1/2}\cdot \max_{\sig\in\{+,-\}} \big|\mathcal{L}^{(3)}_{t, (\sigma,\sig_2,-\sig_2), (a,b,a)} \big| ,
\end{equation}
where $\bsig^{(\alt)}=(\sig_1,-\sig_1,\sig_1,-\sig_1)$ is an alternating loop. By symmetry, a similar inequality holds: % with $\wt{\bsig}' = (-{\sigma}_1, -{\sigma}_2, \sigma_2, \sigma_1)$:
\begin{equation}\label{eq_sym_loop_bound2}
 \sum_{c\in \cal A}\sum_{c'\sim c}\mathcal{L}^{(6)}_{t, (\bsig\otimes \overline\bsig)^{(1)}, \ba(c,c')} \lesssim \frac{1}{W^d\eta_t}\max_{c\in \cal A} \left(\mathcal{L}^{(4)}_{t,\bsig^{(\alt)}, (a,c,a,c)}\right)^{1/2} \cdot \max_{\sig\in \{+,-\}}\big|\mathcal{L}^{(3)}_{t, (\sig, -{\sigma}_2, \sigma_2), (b,a,b)}\big|.
\end{equation}
\end{lemma}
\begin{proof}

We only prove the inequality \eqref{eq_sym_loop_bound}, while the proof of \eqref{eq_sym_loop_bound2} is similar by switching the roles of $(a,c)$ and $(b,c')$. We view the symmetric $6$-loop as a quadratic form. For any $c,c'\in \Zn$ with $|c-c'|\le 1$ and any $z\in [c]$, we define the column vector \smash{$\psi^{z}\in \C^{W^d}$} and the $W^d\times W^d$ matrix \smash{$A^{(c')}$} as
\begin{align*}
    \psi^{z}_{y} &=  (G(\sig_1)E_aG(\sig_2))_{zy}, \ \ \forall y\in [b], \quad \text{and}\quad 
    A^{(c')}_{y_1  y_2} = (G(\sig_1)E_{c'} G(-\sig_1))_{y_1 y_2}, \ \ \forall y_1, y_2 \in [b].
%    W^{-d} \sum_{z'\in \cal [c']} G_{y_1 z'}(\sig_1) G_{z'y_2}(-\sig_1), \quad \forall y_1, y_2 \in [b].
\end{align*}
% For any $z\in \mathcal{I}_c$, define the column vector $\psi^z$ and and a Hermitian matrix $A$ as follows.
% \begin{align*}
%     \psi^{z}_{y} &= W^{-d}\sum_{x\in\mathcal{I}_a}G^\ast_{zx}G_{xy}, \quad y\in \mathcal{I}_b, \\
%     A_{y_1  y_2} & = W^{-d} \sum_{z'\in \mathcal{I}_{c'}} G_{y_1 z'} G^\ast_{z'y_2}, \quad y_1, y_2 \in\mathcal{I}_b.
% \end{align*}
Then, each summand on the left-hand side (LHS) of \eqref{eq_sym_loop_bound} can be written as
\begin{equation}\label{eq_6loop_quadform}
\mathcal{L}^{(6)}_{t, (\bsig\otimes \overline\bsig)^{(1)}, \ba(c,c')} = W^{-3d} \sum_{z\in[c]} \left(\psi^{z}\right)^\ast A^{(c')} \psi^{z} \le W^{-3d} \|A^{(c')}\|  \sum_{z\in[c]} \|\psi^{z}\|_2^2.
\end{equation}
Since the operator norm of $A^{(c')}$ is bounded above by its Hilbert-Schmidt norm, we have
\begin{equation}\label{eq:L4loop}
    \|A^{(c')}\|^2 \le \sum_{y_1,y_2\in [b]} (G(\sig_1)E_{c'} G(-\sig_1))_{y_1y_2} (G (\sig_1)E_{c'}G(-\sig_1))_{y_2 y_1} =: W^{2d} \mathcal{L}^{(4)}_{t,\bsig^{(\alt)}, (c',b,c',b)}.
\end{equation}
On the other hand, the sum of the squared $L^2$-norms of $\psi^z$ is equal to
\begin{equation*}
    \sum_{z\in[c]} \|\psi^{z}\|_2^2 =  \sum_{z\in\cal [c],y\in [b]} (G(\sig_1)E_a G(\sig_2))_{z y} (G(-\sig_2)E_a G(-\sig_1))_{y z} = W^{2d} \mathcal{L}^{(4)}_{t, \bsig', (a,b,a,c)},
\end{equation*}
where $\bsig' = (\sigma_1, \sigma_2, -{\sigma}_2, -{\sigma}_1)$. Plugging the previous two displays into \eqref{eq_6loop_quadform} yields 
\begin{align}\label{eq_sym_loop_bound_old}
 \mathcal{L}^{(6)}_{t, (\bsig\otimes \overline\bsig)^{(1)}, \ba(c,c')} \le \left(\mathcal{L}^{(4)}_{t,\bsig^{(\alt)}, (c',b,c',b)}\right)^{1/2} \cdot \mathcal{L}^{(4)}_{t, \bsig', (a,b,a,c)} .
\end{align}
Summing this inequality over $c'\in \cal A$ and $c\sim c'$, we obtain 
\begin{align*}
\sum_{c'\in \cal A}\sum_{c\sim c'} \mathcal{L}^{(6)}_{t, (\bsig\otimes \overline\bsig)^{(1)}, \ba(c,c')} &\lesssim \max_{c'\in \cal A}\left(\mathcal{L}^{(4)}_{t,\bsig^{(\alt)}, (c',b,c',b)}\right)^{1/2}\cdot \sum_{c}\mathcal{L}^{(4)}_{t, \bsig', (a,b,a,c)} ,
\end{align*}
where the summation over $c$ has been extended to all of $\Zn$, at the cost of a dimensional constant, since each $c$ has at most $3^d$ neighbours $c'$. Finally, applying Ward's identity \eqref{WI_calL} to the sum over $c$ yields the desired bound \eqref{eq_sym_loop_bound}. 
\end{proof}



Now, we are ready to complete the proof of \Cref{lem: EMn2_N}. We first prove the bound \eqref{eq:MG_conclusion}. We split the summation over $c$ into two regions according to the distances from $c$ to the fixed indices $a$ and $b$: 
\begin{equation}\label{eq:S1+S2}
    (\mathcal{E} \otimes \mathcal{E})_{t, \bsig, \boldsymbol{a}, \boldsymbol{a}}^M \lesssim  W^d \sum_{\substack{c'\sim c \\ |c-b|> |c-a|}}  \mathcal{L}^{(6)}_{t, (\bsig\otimes \overline\bsig)^{(1)}, \ba(c,c')} + W^d \sum_{\substack{c'\sim c\\ |c-b|\le |c-a|}} \mathcal{L}^{(6)}_{t, (\bsig\otimes \overline\bsig)^{(1)}, \ba(c,c')} =: \mathcal{S}_1 + \mathcal{S}_2.
\end{equation}
By symmetry, it suffices to prove the bound \eqref{eq:MG_conclusion} for the sum $\mathcal{S}_1$ using \eqref{eq_sym_loop_bound}. (To control the sum $\cal S_2$, we only need to utilize the inequality \eqref{eq_sym_loop_bound2} in place of \eqref{eq_sym_loop_bound} in the following proof.) 
%In this case, by triangle inequality, we have that $|c-b|\ge |a-b|/2$. 
% We split the summation in \eqref{eq_mart-sum} into two parts as follows. Since $c\sim c'$, by triangle inequality we either have $|a-c|\ge |a-b|/3$ or $|c'-b| \ge |a-b|/3$. Since all of the terms in \eqref{eq_mart-sum} are non-negative, we have
% \begin{equation*}
%     (\mathcal{E} \otimes \mathcal{E})_{t, \bsig, \boldsymbol{a}, \boldsymbol{a}}^M \lesssim W^d \sum_{\substack{c\sim c'\\ |c'-b| \ge |a-b|/3}} \mathcal{L}^{(6)}_{t, \bsig', abc'abc} + W^d \sum_{\substack{c\sim c'\\ |c-a| \ge |a-b|/3}} \mathcal{L}^{(6)}_{t, \bsig', abc'abc} =: \mathcal{S}_1 + \mathcal{S}_2.
% \end{equation*}
%Without loss of generality, it suffices to show the bound \eqref{eq:MG_conclusion} for the sum $\mathcal{S}_1$. 
Under the conditions \(|c-c'|\le 1\) and \(|c-b| \ge |a-b|/2\), we bound each off-diagonal $G$-entry of the $4$-loop \smash{$\mathcal{L}^{(4)}_{t,\bsig^{(\alt)}, (c',b,c',b)}$} using \eqref{GijGEX}, and each diagonal $G$-entry using \eqref{GiiGEX}. This gives
\begin{align}\label{eq:pointwise_loop}
\p{ \mathcal{L}^{(4)}_{t,\bsig^{(\alt)}, (c',b,c',b)}}^{1/2} &\prec  \max_{\substack{|b'-b|\le 1,|c''-c'|\le 1,\\ \bsig\in \{(+,-),(-,+)\}}}\mathcal{L}^{(2)}_{t,\bsig, (b',c'')}+W^{-d}\mathbf 1_{|b-c'|\le 1}\prec \Psi_t^2(|b-c'|) \lesssim \Psi_t^2(|a-b|).
\end{align}
Here, the second step uses the assumption \eqref{eq:LW_assm} on the 2-loops together with condition \eqref{eq:Psi}, while the third step follows from $|c-c'|\le 1$, $|c-b|\ge |a-b|/2$, and another application of \eqref{eq:Psi}. Plugging this bound into the $\contract$ inequality \eqref{eq_sym_loop_bound}, we get that
\begin{equation}\label{eq:reduce4to3}
    \mathcal{S}_1 \prec \frac{1}{\eta_t} \Psi_t^2\left(|a-b|\right) \max_{\sig\in\{+,-\}}\left|\mathcal{L}^{(3)}_{t, (\sig,\sig_2,-\sig_2), (a,b,a)}\right|.
\end{equation}
To bound the 3-$G$-loop on the RHS, we write
\[\mathcal{L}^{(3)}_{t, (\sig,\sig_2,-\sig_2), (a,b,a)}=W^{-3d}\sum_{x,x'\in [a]}\sum_{y\in[b]}G_{xx'}(\sig)G_{x'y}(\sig_2)G_{yx}(-\sig_2), \]
and apply an argument analogous to that in \eqref{eq:pointwise_loop}, obtaining 
\begin{align}
\big| \mathcal{L}^{(3)}_{t, (\sig,\sig_2,-\sig_2), (a,b,a)}\big|& \prec \Psi_t(0) \cdot \Psi_t^2(|a-b|).\label{eq:reduce4_bdd3}
\end{align} 
Roughly speaking, the factor $ \Psi_t^2(|a-b|)$ arises from the two long legs---namely, the entries $G_{x'y}$ and $G_{yx}$ between $[a]$ and $[b]$. The short leg $G_{xx'}$, with both indices in $[a]$, contributes a factor of $\Psi_t(0)+W^{-d}$ by \eqref{GiiGEX}. (More precisely, each off-diagonal entry of $G$ contributes a factor of $\Psi_t(0)$, while diagonal entries contribute $W^{-d}$.) Substituting \eqref{eq:reduce4_bdd3} into \eqref{eq:reduce4to3} gives the desired bound \eqref{eq:MG_conclusion}.



In order to show the bound \eqref{eq:MG_conclusion3}, we define two cutoff scales 
\be\label{eq:cutoff_scales}
\ell_t^\ast = (\log W)^{3/2}\ell_t,\quad \ell_t^{\dag} = (\log W)^{7/4}\ell_t.\ee
Then, we divide the proof into two cases. 
First, suppose $|a-b| \le \ell_t^{\dag}$. In this case, the function $\widetilde{\mathcal{T}}^{\ell}_{t,D}\p{|a-b|}$ exhibits no exponential decay. Hence, it suffices to apply the polynomial decay bound \eqref{eq:MG_conclusion2} with $c_0=d$ and $K=\ell$, i.e.\ with the \emph{truncated} profile function \smash{\(\Psi_t(r) = (W^{-d} B_{t, r\wedge \ell})^{1/2}.\)}  
The conclusion \eqref{eq:MG_conclusion3} then follows directly from \eqref{eq:MG_conclusion2} and the facts 
\[
\br{\Psi_t(r)}^2 \prec  W^{-d}\widetilde{\mathcal{T}}^{\ell}_{t,D}\p{r} , \quad W^{-d}\widetilde{\mathcal{T}}^{\ell}_{t,D}\p{r} \prec \br{\Psi_t(r)}^2, \quad \forall \ 0\le r\le \ell_t^\dag,
\]
both of which hold by the definition \eqref{defWTTlD} of \smash{$\wT^{\ell}_{t,D}$}, since ${\cal T}_{t}(\rho)\asymp B_{t,\rho}\exp(-(\rho/\ell_t)^{1/2})$ and $\exp(-(\rho/\ell_t)^{1/2})\ge W^{-\oo(1)}$ for $0\le \rho\le \ell\wedge \ell_t^\dag$. We stress that the truncation is needed here: for the untruncated profile $(W^{-d}B_{t,r})^{1/2}$ the second bound fails in the range $\ell<r\le \ell_t^\dag$, where \smash{$\wT^{\ell}_{t,D}(r)={\cal T}_t(\ell)$} while $B_{t,r}\ll B_{t,\ell}$.
It remains to deal with the case 
\be\label{eq:largea-b}|a-b| > \ell_t^{\dag} = (\log W)^{7/4}\ell_t.\ee 
In this case, we split the summation \eqref{eq_mart-sum} into three parts as follows: 
\begin{align}
    (\mathcal{E} \otimes \mathcal{E})_{t, \bsig, \boldsymbol{a}, \boldsymbol{a}}^M &\lesssim W^d \bigg(\sum^\star_{\substack{|c-a| \le \ell_t^\ast \text{ or } |c'-b| > \ell}} + \sum^\star_{\substack{|c'-b| \le \ell_t^\ast \text{ or } |c-a| > \ell}} + \sum^\star_{\substack{\ell_t^\ast < |c'-b|,|c-a| \le \ell}}\bigg) \mathcal{L}^{(6)}_{t, (\bsig\otimes \overline\bsig)^{(1)}, \ba(c,c')}\nonumber\\
    &=: \widetilde{\mathcal{S}}_1 + \widetilde{\mathcal{S}}_2 + \widetilde{\mathcal{S}}_3.\label{eq;S123}
\end{align}
where $\sum^\star$ refers to the summation subject to the constraint $|c-c'|\le 1$. 
To estimate $\widetilde{\mathcal{S}}_1$, we combine \eqref{GijGEX} with \eqref{eq:LW_assm_exp} and proceed analogously to \eqref{eq:pointwise_loop}. This yields  
\begin{align}\label{eq:pointwise_loop2}
    \p{ \mathcal{L}^{(4)}_{t,\bsig^{(\alt)}, (c',b,c',b)}}^{1/2} \prec W^{-d}\wT_{t,D}^{\ell}(|b-c'|) \prec W^{-d}\wT_{t,D}^{\ell}(|a-b|),
\end{align}
where the second inequality is justified as follows: (1) if $|c'-b|>\ell$, then \smash{\(\wT_{t,D}^{\ell}(|b-c'|) = \wT_{t,D}^{\ell}(\ell)\le \wT_{t,D}^{\ell}(|a-b|);\)} (2) if $|c-a|\le \ell_t^*$, then by \eqref{eq:largea-b}, we have $|b-c'|\ge |a-b|-(\ell_t^*+1) =(1-\oo(1))|a-b|$, which implies 
\[ \wT_{t,D}^{\ell}(|b-c'|) \le \wT_{t,D}^{\ell}\p{|a-b|-(\ell_t^*+1)} \prec \wT_{t,D}^{\ell}(|a-b|). \]
Substituting \eqref{eq:pointwise_loop2} into the $\contract$ inequality \eqref{eq_sym_loop_bound}, we obtain  
\begin{align}
    \widetilde{\mathcal{S}}_1 &\prec  \frac{1}{\eta_t} \left[W^{-d} \widetilde{\mathcal{T}}^{\ell}_{t,D}\p{|a-b|}\right] \max_{\sig\in\{+,-\}}\left|\mathcal{L}^{(3)}_{t, (\sig,\sig_2,-\sig_2), (a,b,a)}\right|.\label{eq_S1tilde}
\end{align}
Next, using \eqref{GijGEX} and \eqref{eq:LW_assm_exp}, we obtain a similar bound on the 3-$G$-loop as in \eqref{eq:reduce4_bdd3}:
\[\left|\mathcal{L}^{(3)}_{t, (\sig,\sig_2,-\sig_2), (a,b,a)}\right|\prec (W^{-d}B_{t,0})^{1/2}\cdot W^{-d} \widetilde{\mathcal{T}}^{\ell}_{t,D}\p{|a-b|}.\]
Plugging this into \eqref{eq_S1tilde} gives
\begin{equation}\label{eq:boundwtS_1}
    \widetilde{\mathcal{S}}_1 \prec \frac{1}{\eta_t} \left(W^{-d}B_{t,0}\right)^{1/2}\left[W^{-d} \widetilde{\mathcal{T}}^{\ell}_{t,D}\p{|a-b|}\right]^2 , 
\end{equation}
which is controlled by the RHS of \eqref{eq:MG_conclusion3}. By symmetry, an identical bound holds for $\wt {\cal S}_2$.


% Since, all terms on the right-hand side are bounded by large constant powers of $W$ and the control parameters $B_{t,0}$, $\widetilde{\mathcal{T}}^{\ell,D}_{t,|a-b|}$ are bounded below by constant powers of $W$, the constant $\tilde{D}$ can be chosen large enough, such that
% \begin{equation*}
%      \widetilde{\mathcal{S}}_1 \prec \frac{1}{\eta_t} \left[\left(W^{-d}B_{t,0}\right)^{1/2} + \left(W^{-d}B_{t,0}\right)^{1/2}\left(\mathcal{J}^{\ell,D}_{t}\right)^2\right] \cdot \left(W^{-d} \widetilde{\mathcal{T}}^{\ell,D}_{t,|a-b|}\right)^2.
% \end{equation*}

It remains to control the sum $\widetilde{\mathcal{S}}_3$. In this case, we bound all legs of the original $6$-loop directly using \eqref{GijGEX}. Here, all legs of this loop have lengths (i.e., $|a-b|$, $|a-c|$, and $|c'-b|$) at least $\ell_t^\ast$. Then, instead of using the assumption \eqref{eq:LW_assm_exp}, we will bound the $2$-loops using the parameter defined in \eqref{defCALJ} as follows: for any $a,b\in \Zn$ satisfying $|a-b|\gtrsim \ell_t^\ast$,  
\begin{equation}\label{eq_L2-J}
    \mathcal{L}^{(2)}_{t,(-,+),(a,b)} \prec \wh{\mathcal{J}}^{\ell}_{t,D} \cdot W^{-d} \widetilde{\mathcal{T}}^{\ell}_{t,D}\p{|a-b|} + \mathcal{K}^{(2)}_{t,(-,+),(a,b)} \lesssim \big(\wh{\mathcal{J}}^{\ell}_{t,D} +W^{-D}\big)\cdot W^{-d} \widetilde{\mathcal{T}}^{\ell}_{t,D}\p{|a-b|} ,
\end{equation}
where the second step is due to the exponential decay of the 2-$\mathcal{K}$-loop given by \eqref{prop:ThfadC}. 
Then, using \eqref{GijGEX} to bound the 6 legs of the 6-loop, we obtain that 
\begin{align}
    \widetilde{\mathcal{S}}_3 \prec &~W^{-2d}\p{\wh{\mathcal{J}}^{\ell}_{t,D}+W^{-D}}^3\sum_{\substack{\ell_t^\ast < |c'-b|,|c-a| \le \ell}}^\star \widetilde{\mathcal{T}}^{\ell}_{t,D}\p{|a-b|} \widetilde{\mathcal{T}}^{\ell}_{t,D}\p{|a-c|}  \widetilde{\mathcal{T}}^{\ell}_{t,D}\p{|c-b|} ,\label{eq:boundwtS31}
\end{align}
where we also use that 
$\widetilde{\mathcal{T}}_{t,D}^{\ell}\p{|c'-b|}\asymp \widetilde{\mathcal{T}}_{t,D}^{\ell}\p{|c-b|}$ for $|c-c'|\le 1$. We can bound the RHS of \eqref{eq:boundwtS31} by
\begin{align*}
% &~W^{-2d} \left(\mathcal{J}^{\ell}_{t,D}+W^{-D}\right)^3\cdot  \widetilde{\mathcal{T}}^{\ell}_{t,D}\p{|a-b|} \sum_{\ell_t^\ast<|c-a|,|c'-b|\le \ell}^\star \widetilde{\mathcal{T}}^{\ell}_{t,D}\p{|a-c|} \widetilde{\mathcal{T}}^{\ell}_{t,D}\p{|c-b|} \\
% \lesssim 
&~\left(\wh{\mathcal{J}}^{\ell}_{t,D}+W^{-D}\right)^3\cdot W^{-2d} \widetilde{\mathcal{T}}^{\ell}_{t,D}\p{|a-b|}\cdot \sum_c \br{\mathcal{T}_{t}\p{|a-c|}+W^{-D}}\br{\mathcal{T}_{t}\p{|c-b|}+W^{-D}} \\
\lesssim &~\frac{(\wh{\mathcal{J}}^{\ell}_{t,D}+W^{-D})^3}{1-t} \cdot W^{-2d} \widetilde{\mathcal{T}}^{\ell}_{t,D}\p{|a-b|}\cdot \br{\mathcal{T}_{t}\p{|a-b|} +W^{-D}}\lesssim \frac{(\wh{\mathcal{J}}^{\ell}_{t,D}+W^{-D})^3}{\eta_t} \cdot \left(W^{-d} \widetilde{\mathcal{T}}^{\ell}_{t,D}\p{|a-b|}\right)^2,
\end{align*}
where we use \eqref{TTT2} and the fact $\sum_{a}\cal T_t\p{|a|}\lesssim (1-t)^{-1}$. This finishes the estimation of \smash{$\wt{\cal S}_3$}, and hence completes the proof of the bound \eqref{eq:MG_conclusion3}.

%We expand the RHS of \eqref{eq:boundwtS31} and notice that by choosing $D'$ sufficiently large depending on $D$, we can guarantee that each term containing the factor $W^{-D'}$ will be controlled by the RHS of \eqref{eq:MG_conclusion3}. 
%Note that replacing $c'$ by $c$ here is harmless, because $\widetilde{\mathcal{T}}^{\ell,D}_{t,|c'-b|}\prec \widetilde{\mathcal{T}}^{\ell,D}_{t,|c-b|}$. Similarly to the bound on $\widetilde{\mathcal{S}}_1$ above, we can ignore the terms involving $W^{-\tilde{D}}$, because $\tilde{D}$ can be chosen large enough to have these terms consumed by $\eta_t^{-1}\left(W^{-d}B_{t,0}\right)^{1/2}\left(W^{-d} \widetilde{\mathcal{T}}^{\ell,D}_{t,|a-b|}\right)^2$. 
%For the main term, we bound it by  %$\widetilde{\mathcal{T}}^{\ell,D}_{t,|a-c|}$ and $\widetilde{\mathcal{T}}^{\ell,D}_{t,|c-b|}$ by $\mathcal{T}_{t,|a-c|}$ and $\mathcal{T}_{t,|c-b|}$ taking into account that $|c-a|, |c-b| \le \ell$ and then extend the summation over $c$ to the entire space $\Z_L^d$. This allows us to apply \eqref{TTT2} as follows

%\end{proof}
 

 

 
\section{Steps 3 and 4: Sharp maximum estimates for \texorpdfstring{$G$}{G}-loops}\label{Sec:Steps34}


In this section, we apply the integrated loop hierarchy \eqref{int_K-LcalE} to complete the proofs of Steps 3 and 4. The key task is to show that the second through fifth terms in \eqref{int_K-LcalE} are errors in the max-norm sense.

The integrands in the second to fourth terms on the RHS of \eqref{int_K-LcalE} each involve a global sum over an index of an $\cal L$- or $(\cL-\cK)$-loop (see \eqref{DefKsimLK}, \eqref{def_ELKLK}, and \eqref{def_EwtG}). In low dimensions $d\in \{1,2\}$, these terms can be controlled directly using their max-norms, together with a factor $\ell_t^d$ that captures the range of exponential decay.
For $d\in\{1,2\}$, the factor $\ell_t^d$ is bounded by $|1-t|^{-1}$, which incurs only a harmless logarithmic contribution upon integration in time. However, in higher dimensions $d\ge 3$, an additional factor $\ell_t^{d-2}$ arises on top of the $(1-t)^{-1}$, creating a substantial term that cannot be canceled. 
To obtain sufficiently sharp bounds without relying on the precise pointwise decay of higher-order $G$-loops, we develop a new \emph{$\contract$ inequality} \eqref{yi2oslxj2}, derived from the Cauchy--Schwarz inequality and the Ward's identities \eqref{WI_calL} and \eqref{WI_calK}. Although technically simple, this inequality provides a powerful tool: it allows us to control the absolute sum of a higher-order $G$-loop in terms of lower-order ones in a sharp max-norm sense, and crucially produces the desired $(1-t)^{-1}$ factor without introducing any excess powers of $\ell_t$.
Among all these terms, one remains that cannot be effectively controlled using the $\contract$ inequalities---namely, the term of the form $(\cal L-\cal K)^{(2)}\diamond(\cal L-\cal K)^{(n)}$ (recall the notation in \eqref{eq:diamond}) in \eqref{def_ELKLK}. This term is instead shown to be an error by applying the pointwise $(\cal L-\cal K)^{(2)}$-estimate \eqref{Eq:Gdecay_w}.


Bounding the martingale term requires controlling a $(2n+2)$-loop (recall \Cref{def:CALE}). The corresponding loop structure is symmetric and involves two summation indices along the loop (see e.g., \eqref{eq:2martingale_quad}). To handle it, we derive another $\contract$ inequality \eqref{u2jzooi-2}, obtained by extending the argument in the proof of \Cref{ygdhmsgq0}. This inequality yields a non-sharp but sufficient bound for the martingale term, introducing an additional power of $W$, i.e., the factor $(W^{-d}B_{t,0})^{-1/(2p)}$ in \eqref{eq:MG_nloop}. However, this factor remains small enough for our proof for large enough $p$.


Our overall strategy for Steps 3 and 4 can now be summarized as follows. Using the $\contract$ inequalities together with a good stability estimate \eqref{Eq:Gdecay_w} for $({\cal L}-{\cal K})^{(2)}$-loops, we can bound each term in \eqref{int_K-LcalE} either by lower-order $(\cL-\cK)$-loops or by higher-order $(\cL-\cK)$-loops with an extra small factor $W^{-\varepsilon}$. Then, with the integrated hierarchy \eqref{int_K-LcalE}, we can bootstrap from a sequence of weaker bounds to tighter ones, ultimately yielding improved control on the entire hierarchy of $(\cL -\cK)$-loops in each induction. As in \cite{YY_25, DYYY25, erdHos2025zigzag}, after $\OO(1)$ iterations, this process yields (almost) sharp max-norm bounds for all $(\cL-\cK)$-loops.



%\subsection{Ward's inequality and estimation of the loop hierarchy}


%The proofs of Steps 3 and 4 rely crucially on the following (maximum) Ward's inequalities. % which are consequences of the Ward's identities \eqref{WI_calL} and \eqref{WI_calK}. %The second Ward's inequality \eqref{u2jzooi-2} is essential to bound the quadratic variation of the martingale term $\mathcal{E}^{M,(n)}$. It is a generalized $G$-loop version of the bound used in Lemma 10 of \cite{DY} to control the term $\mathcal{E}^{M,(2)}$.  

\begin{lemma} [$\Contract$ inequality]\label{ygdhmsgq}
For each $1\le k \le n-1$, we have the bound
\begin{align}\label{yi2oslxj2}
  \max_{\bsig}\sum_{a_n} |{\cal L}^{(n)}_{t, \bsig, \ba}|\le \frac 1 {W^{d}\eta_t} \left(\max_{\bsig,\mathbf{b}}  \left|{\cal L}^{(2k-1)}_{t, \bsig, \mathbf{b}}\right|\cdot \max_{\bsig,\mathbf{b}}  \left|{\cal L}^{(2n-2k-1)}_{t, \bsig, \mathbf{b} }\right|\right)^{\frac 1 2}   .
\end{align}
Furthermore, for any $n\ge 4$, $1\le k<j<l\le n-1$, and subsets $\cal A(a_n)\subset \Zn$ (which may depend on $a_n$) of cardinality $|\cal A(a_n)|\le C$ for a constant $C>0$, the following bound holds for any fixed $p\ge 1$: 
\begin{align}\label{u2jzooi-2}
   \max_{\bsig} \sum_{a_n}\sum_{a_j\in \cal A(a_n)} |{\cal L}^{(n)}_{t,\bsig, {\ba}}|\le  &~\frac{C}{W^d\eta_t} \left(\max_{\bsig,\mathbf{b}}  \left|{\cal L}^{(2k-1)}_{t, \bsig, \mathbf{b}}\right|\cdot \max_{\bsig,\mathbf{b}}  \left|{\cal L}^{(2n-2l-1)}_{t,\bsig,\mathbf{b}}\right|\right)^{\frac 1 2}   \left(\max_{\bsig,\bfb}   \left|{\cal L}^{(2(l-k)p)}_{t,\bsig, \mathbf b}\right|\right)^{\frac1{2p}} .
\end{align}
\end{lemma}

\begin{proof}%[Proof of lemma \ref{ygdhmsgq}]
We first prove \eqref{yi2oslxj2}.  Applying the Cauchy-Schwarz inequality with respect to the averages over $[a_k]$ and $[a_n]$ in the loop \smash{${\cal L}^{(n)}_{t,\bsig, \ba}$}, we obtain that for $\bsig=(\sig_1,\ldots, \sig_n)$ and $\ba=(a_1,\ldots, a_n)$, 
\begin{align}\label{eq:CS1}
     |{\cal L}^{(n)}_{t, \bsig,\ba}| 
     \; \le \;&
     \left({\cal L}^{(2k)}_{t,\bsig_1,\ba_1 }\cdot  {\cal L}^{(2n-2k)}_{t,\bsig_2,\ba_2}\right)^{1/2},
     \end{align}
where $\ba_1$, $\ba_2$, $\bsig_1$, and $\bsig_2$ are defined as     
 \begin{align}
    {\ba}_1= (a_1,\cdots, a_{k-1}, a_k, a_{k-1},\cdots ,a_1,a_n), &\quad 
    \bsig_1 = (\sigma_1, \cdots, \sigma_{k}, -{\sigma_{k}},\cdots, -\sigma_1) ,\nonumber
     \\\nonumber
    {\ba}_2= (a_{n-1},\cdots, a_{k+1}, a_k, a_{k+1}, \cdots ,a_{n-1},a_n), &\quad 
    \bsig_2 = (-{\sigma_n}, \cdots, -\sig_{k+1}, \sigma_{k+1}, \cdots ,\sigma_n).
\end{align}  
Note that the loops ${\cal L}^{(2k)}_{t,\bsig_1,{\ba}_1}$ and ${\cal L}^{(2n-2k)}_{t, \bsig_2,{\ba}_2}$ are both non-negative since $\bsig_1$ and $\bsig_2$ are symmetric.  Then, applying the Cauchy-Schwarz inequality to \eqref{eq:CS1} again, we obtain that 
\be\label{eq;genWard0}
      \sum_{a_n} |{\cal L}^{(n)}_{t, \bsig,\ba}| \le  
     \Big(\sum_{a_n}{\cal L}^{(2k)}_{t, \bsig_1,\ba_1}\Big)^{1/2} \Big(\sum_{a_n}{\cal L}^{(2n-2k)}_{t, \bsig_2,\ba_2}\Big)^{1/2}.
\ee
Using Ward's identity \eqref{WI_calL}, we can express the $G$-loops on the RHS as  
\be\label{eq;genWard1}
\sum_{a_n}{\cal L}^{(2k)}_{t, \bsig_1,\ba_1} =\frac{{\cal L}^{(2k-1)}_{t, \bsig^+_1,\ba'_1}
-{\cal L}^{(2k-1)}_{t, \bsig^-_1,\ba'_1}}{2\ii W^d\eta_t} ,\quad
\sum_{a_n}{\cal L}^{(2n-2k)}_{t, \bsig_2,\ba_2} =\frac{{\cal L}^{(2n-2k-1)}_{t, \bsig^+_2,\ba'_2}
-{\cal L}^{(2n-2k-1)}_{t, \bsig^-_2,\ba'_2}}{2\ii W^d\eta_t},
\ee
where $\ba'_1$, $\ba'_2$, $\bsig^{\pm}_1$, and $\bsig^{\pm}_2$ are defined as 
\begin{align*}
\ba'_1=(a_1,\cdots, a_{k-1}, a_k, a_{k-1}\cdots, a_1),\quad &\bsig^{\pm}_1 = ( \pm,\sigma_2, \cdots, \sigma_{k}, -{\sigma_{k}},\cdots,-{\sigma_2}), \\
\ba'_2=( a_{n-1}\cdots,   a_{k+1}, a_{k}, a_{k+1},\cdots, a_{n-1}), \quad 
    &\bsig^{\pm}_2 = (\pm, -{\sigma_{n-1}}, \cdots, - \sigma_{k+1}, \sigma_{k+1},\cdots  \sigma_{n-1}).
\end{align*}
Plugging \eqref{eq;genWard1} into \eqref{eq;genWard0}, we  conclude \eqref{yi2oslxj2}. 

For $\bsig=(\sigma_1,  \ldots, \sigma_n)\in \{+,-\}^n$ and \smash{$\ba = (a_1, \ldots, a_{n-1})\in (\Zn)^{n-1}$}, denote a $G$-chain of length $n$ by 
\begin{equation}\label{defC=GEG}
    {\cal C}^{(n)}_{t, \bsig, \ba} = \prod_{i=1}^{n-1}\left(G_t(\sigma_i) E_{a_i}\right)\cdot  G_t(\sigma_n) .
\end{equation}
To show \eqref{u2jzooi-2}, we split the loop ${\cal L}^{(n)}_{t, \bsig, \ba}$ at $a_k$, $a_l$, and $a_n$, and write it as 
$${\cal L}^{(n)}_{t, \bsig, \ba} = W^{-3d}\sum_{x_k\in\br{a_k} }\sum_{x_l\in\br{a_l}} \sum_{x_n\in \br{a_n}} \p{{\cal C}^{(k)}_{t, \bsig_1, \ba_1}}_{x_n x_k}\p{{\cal C}^{(l-k)}_{t, \bsig_2, \ba_2}}_{x_kx_l}\p{{\cal C}^{(n-l)}_{t, \bsig_3, \ba_3}}_{x_lx_n} ,$$
where $\ba_1=(a_1,\ldots,a_{k-1})$, $\bsig_1=(\sig_1,\ldots,\sig_k)$, $\ba_2=(a_{k+1},\ldots,a_{l-1})$, $\bsig_2=(\sig_{k+1},\ldots,\sig_l)$, $\ba_3=(a_{l+1},\ldots,a_{n-1})$, and $\bsig_3=(\sig_{l+1},\ldots,\sig_n)$. We can bound the RHS by using the operator norm of the $W^d\times W^d$ matrix $A=(A_{xy}:x\in  \br{a_k}, y\in \br{a_l})$ with \smash{$A_{xy}= ({\cal C}^{(l-k)}_{t, \bsig_2, \ba_2})_{xy}$}:
\begin{align}
    \left|{\cal L}^{(n)}_{t, \bsig, \ba}\right|\le 
\frac{1}{W^{2d}}\sum_{x_n\in \br{a_n}}
\bigg(\sum_{x_k\in \br{a_k}}\Big|\p{{\cal C}^{(k)}_{t, \bsig_1, \ba_1}}_{x_n x_k}\Big|^2\bigg)^{1/2}
\bigg(\sum_{x_l\in \br{a_l}}\Big|\p{{\cal C}^{(n-l)}_{t, \bsig_3, \ba_3}}_{x_l x_n}\Big|^2\bigg)^{1/2} \cdot \frac{1}{W^d}\|A\| .\label{eq:Apsipsi}
\end{align}
To control the operator norm $\|A\|$, we use the simple linear algebra fact 
\(\|A\|\le \left\{\tr[(AA^*)^p]\right\}^{\frac{1}{2p}} \)
for any $p\in \N$, which gives that  
\begin{align}\label{eq:Apsipsi1}
\frac{1}{W^d}\|A\| \le  \left\{\frac{1}{W^{2pd}}\tr \br{\left(AA^*\right)^{p}}\right\}^{\frac1{2p}}
\le \left(\max_{\bsig,\bfb}   \left|{\cal L}^{(2(l-k)p)}_{t, \bsig , \bfb}\right|\right)^{\frac1{2p}}
\end{align}
Plugging \eqref{eq:Apsipsi1} into \eqref{eq:Apsipsi} and applying the Cauchy--Schwarz inequality, we obtain that 
\begin{align*}
 \sum_{a_n}\sum_{a_j\in \cal A(a_n)} \left|{\cal L}^{(n)}_{t, \bsig, \ba}\right|&\le C  \left(\max_{\bsig,\bfb}  \left| {\cal L}^{(2(l-k)p)}_{t, \bsig, \bfb}\right|\right)^{\frac1{2p}} \bigg(\frac{1}{W^{2d}}\sum_{a_n}\sum_{x_n\in\br{a_n}}\sum_{x_k\in\br{a_k}}\Big|\p{{\cal C}^{(k)}_{t, \bsig_1, \ba_1}}_{x_n x_k}\Big|^2\bigg)^{1/2} \\
 &\qquad \qquad \times \bigg(\frac{1}{W^{2d}}\sum_{a_n}\sum_{x_n\in\br{a_n}}\sum_{x_l\in\br{a_l}}\Big|\p{{\cal C}^{(n-l)}_{t, \bsig_3, \ba_3}}_{x_l x_n}\Big|^2\bigg)^{1/2}   \\
&\le \frac{C}{W^d\eta_t}\left(\max_{\bsig,\bfb}   \left|{\cal L}^{(2(l-k)p)}_{t, \bsig, \bfb}\right|\right)^{\frac1{2p}} \bigg(\max_{\bsig,\bfb}  \left|{\cal L}^{(2k-1)}_{t, \bsig, \bfb}\right|\cdot \max_{\bsig,\bfb}  \left|{\cal L}^{(2n-2l-1)}_{t, \bsig, \bfb}\right| \bigg)^{1/2},
\end{align*}
where we applied Ward's identity \eqref{WI_calL} in the second step. This concludes \eqref{u2jzooi-2}.
\end{proof}

%cite \cite{DYYY25_d3} here
In the course of proving \Cref{ML:Kbound}, we will also establish the following bound, whose proof is deferred to \Cref{Sec:CalK}.

\begin{lemma}\label{lem_wardineq_K}
For any $n\ge 2$ and $t\in[0,1)$, we have that
\begin{align}\label{wardineq_K}
  \max_{\bsig\in\{+,-\}^n}\sum_{a_n} |{\cal K}^{(n)}_{t, \bsig, \ba}|\prec \frac 1 {W^{d}\eta_t} (W^{-d}B_{t,0})^{n-2}.  
\end{align} 
\end{lemma}
 

For any $n\ge 1$, let $\Xi^{({\cal L})}_{t, n}\ge 1$ and $\Xi^{({\cal L}-{\cal K})}_{t, n}\ge 1$ be \emph{deterministic} control parameters for $\cL$-loops and $(\cL-\cK)$-loops of length $n$ such that the following bounds hold: 
\begin{align}\label{def:XiL}
	\wh\Xi^{({\cal L})}_{t, n} &:= 1+ \max_{\bsig\in \{+, -\}^n}\max_{\ba\in (\Zn)^n} \left|{\cal L}^{(n)}_{t, \bsig, \ba} \right|\Big/ \left(W^{-d}B_{t,0}\right)^{n-1}\prec \Xi^{({\cal L})}_{t, n},\\
\wh\Xi^{({\cal L}-{\cal K})}_{t, n} &:=1+\max_{\bsig\in \{+, -\}^n}\max_{\ba\in (\Zn)^n} \left|({\cal L}-{\cal K})^{(n)}_{t, \bsig, \ba} \right|\Big/\left(W^{-d}B_{t,0}\right)^{n} \prec \Xi^{({\cal L}-{\cal K})}_{t, n}.\label{def:XIL-K}
\end{align}
Using these control parameters and the $\contract$ inequalities in \Cref{ygdhmsgq,lem_wardineq_K}, we can control the terms on the RHS of equation \eqref{int_K-LcalE} as follows. 
  
  
\begin{lemma}[Estimates of $\cal E$ terms]\label{lem:SEforLn}
In the setting of \Cref{lem:main_ind}, suppose the local laws \eqref{Gt_bound_flow} and \eqref{Gt_avgbound_flow} and the $2$-$G$-loop estimate \eqref{Eq:Gdecay_w} hold. Then, the following estimates hold for any fixed $n\ge 2$: 
 \begin{enumerate}

\item The light-weight term defined in \eqref{def_EwtG} satisfies that  
\begin{align}\label{bEwGn}
 \max_{\bsig,\ba}\abs{{\cal E}^{\Gc, (n)}_{t,\bsig,\ba}}\prec \frac{\p{W^{-d}B_{t,0}}^n}{\eta_t} \cdot\p{\wh\Xi^{({\cal L})}_{t, n_1}\cdot \wh\Xi^{({\cal L})}_{t, n_2}}^{\frac12}, 
\end{align} 
where $n_1=n-1$ and $n_2=n+1$ if $n$ is even, and $n_1=n_2=n$ if $n$ is odd.
 
\item For $3\le \lenk \le n$, the $[\cK^{(\lenk)}\sim (\cL-\cK)]^{(n)}$ term defined in \eqref{DefKsimLK} satisfies that
\begin{align}\label{eq:KsimL-K}
\max_{\bsig,\ba} \abs{[\cK^{(\lenk)}\sim (\cL-\cK)]^{(n)}_{t,\bsig,\ba}}\prec \frac{\left(W^{-d}B_{t,0}\right)^{n}}{\eta_t}\cdot \wh\Xi^{({\cal L}-{\cal K})}_{t, n-\lenk+2}. 
\end{align}


\item The ${\cal E}^{(\cL-\cK)\times (\cL-\cK),(n)}$ term defined in \eqref{def_ELKLK} satisfies that 
\begin{align}\label{eq:L-KsimL-K}
\max_{\bsig,\ba} \abs{ {\cal E}^{(\cL-\cK)\times (\cL-\cK),(n)}_{t,\bsig,\ba}}\prec &~ 
  \frac{\left(W^{-d}B_{t,0}\right)^{n}}{\eta_t}  \sum_{n'= \lceil n/2\rceil +1}^{n-1} \wh\Xi^{({\cal L}-{\cal K})}_{t,n+2-n'} \cdot \left(\wh\Xi^{({\cal L})}_{t,n_1'} \cdot  \wh\Xi^{({\cal L})}_{t,n_2'}\right)^{1/2} \nonumber\\
 &~ +\frac{\left(W^{-d}B_{t,0}\right)^{n}}{\eta_t}  \cdot (W^{-d} {B}_{t,0})^{\frac 1 6}\wh\Xi^{({\cal L}-{\cal K})}_{t,n},
\end{align}
where $n'_1=n_2'=n'-1$ if $n'$ is even, and $n'_1=n'-2$ and $n_2'=n'$ if $n'$ is odd. Note that the first term on the RHS is zero when $n\in\{2,3\}$. 
 
\item The term $\left( \cal E\otimes  \cal E \right)^{M}_{t, \bsig, \ba,\ba'}\equiv \left( \cal E\otimes  \cal E \right)^{M,(n;k)}_{t, \bsig, \ba,\ba'}$ defined in \eqref{defEOTE} satisfies the following estimate for each $k \in \qqq{n}$ and any fixed $p\in \N$:
\begin{align}\label{eq:MG_nloop}
\max_{\bsig}\max_{\ba,\ba'}\left( \cal E\otimes  \cal E \right)^{M}_{t, \bsig, \ba,\ba'}\prec \frac{(W^{-d} {B}_{t,0})^{2n-\frac{1}{2p}}}{\eta_t}  \cdot \wh\Xi^{({\cal L})}_{t,2n-1}  \left(\wh\Xi^{({\cal L})}_{t,4p}\right)^{\frac{1}{2p}}. 
\end{align}
\end{enumerate}
\end{lemma}
\begin{proof}%[Proof of lemma \ref{lem:SEforLn}]
Using the inequality \eqref{yi2oslxj2} (with $k=\lceil n/2\rceil$) and the averaged local law \eqref{Gt_avgbound_flow} established in Step 2, we can bound that 
\begin{align*}
   \abs{\mathcal{E}^{\Gc,(n)}_{t, \bsig, \ba}} &\le   {W}^d \sum_{k=1}^n \sum_{a, b\in \Zn} \; 
 \abs{\tr\p{ \Gc_t(\sigma_k) E_{a} }}
   S^{(\sB)}_{ab} 
  \abs{\left( {\cut}^{(b)}_{k} \circ {\cal L}^{(n)}_{t, \bsig, \ba} \right)} \\
 & \prec W^d\cdot (W^{-d}B_{t,0}) \cdot \frac{1}{W^d\eta_t} \left(\max_{\bsig,\bfb}  \left|{\cal L}^{(n_1)}_{t, \bsig, \bfb}\right|\cdot \max_{\bsig,\bfb}  \left|{\cal L}^{(n_2)}_{t, \bsig, \bfb}\right|\right)^{\frac 1 2},
  \end{align*}
which concludes \eqref{bEwGn} together with the definition \eqref{def:XiL}. For \eqref{eq:KsimL-K}, with the $\cK$-loop bounds \eqref{eq:bcal_k} and \eqref{wardineq_K}, we obtain that 
 \begin{align*}
\max_{\bsig,\ba}\abs{[\cK^{(\lenk)}\sim (\cL-\cK)]^{(n)}_{t,\bsig,\ba}} &\lesssim W^d \p{\max_{\bsig,\ba}\abs{(\mathcal{L} - \mathcal{K})^{(n-\lenk+2)}_{t, \bsig, \ba}}} \cdot  \max_{\bsig,\ba} \sum_{a_{\lenk}}\abs{\cK^{(\lenk)}_{t,\bsig,\ba}}    \\
&\lesssim \frac{1}{\eta_t}\p{\max_{\bsig,\ba}\abs{(\mathcal{L} - \mathcal{K})^{(n-\lenk+2)}_{t, \bsig, \ba}}} \cdot (W^{-d}B_{t,0})^{\lenk-2} ,
%\max_{\bsig,\mathbf{b}}  \left|{\cal K}^{(2\lenk-2k-1)}_{t, \mathbf{b}, \bsig}\right|\right)^{\frac 1 2} \\&\prec \frac{(W^{-d}B_{t,0})^{\lenk-2}}{\eta_t}\p{\max_{\bsig,\ba}\abs{(\mathcal{L} - \mathcal{K})^{(n-\lenk+2)}_{t, \bsig, \ba}}},  
\end{align*} 
which concludes \eqref{eq:KsimL-K} together with the definition \eqref{def:XIL-K}. 

For the estimate \eqref{eq:L-KsimL-K}, we first consider the case $n\in\{2,3\}$. In this case, we have that  
\begin{align*}
\max_{\bsig,\ba} \abs{\mathcal{E}^{(\cL-\cK)\times (\cL-\cK),(n)}_{t, \bsig, \ba}} &\lesssim  \p{\max_{\bsig,\ba}\abs{(\mathcal{L} - \mathcal{K})^{(n)}_{t, \bsig, \ba}}}\cdot \max_{\bsig,\ba}W^{d}\sum_{a_2}\abs{(\mathcal{L} - \mathcal{K})^{(2)}_{t, \bsig,\ba}} \prec \frac{(W^{-d} {B}_{t,0})^{\frac16}}{\eta_t}{\max_{\bsig,\ba}\abs{(\mathcal{L} - \mathcal{K})^{(n)}_{t, \bsig, \ba}}},
\end{align*}
which concludes \eqref{eq:L-KsimL-K} for $n\in\{2,3\}$ by using the definition \eqref{def:XIL-K}. Above, in the second step, we use the $2$-$G$-loop estimate \eqref{Eq:Gdecay_w} to get that 
\be\label{eq:sumtwoloop}W^{d}\sum_{a_2}\abs{(\mathcal{L} - \mathcal{K})^{(2)}_{t, \bsig,\ba}} \prec \p{\frac{1-s}{1-t}}^{C_d}(W^{-d} {B}_{t,0})^{1/5}\cdot \sum_{a_2} {\cal T}_{t}(|a_1-a_2|)+W^{-D} \prec \frac{(W^{-d} {B}_{t,0})^{1/6}}{\eta_t} \ee
under the condition \eqref{con_st_ind} as long as $\fc_d$ is chosen sufficiently small depending on $C_d$. 
For general $n\ge 4$, we need to bound that 
\begin{align}\nonumber
\max_{\bsig,\ba}\abs{\mathcal{E}^{(\cL-\cK)\times (\cL-\cK),(n)}_{t, \bsig, \ba}} &\lesssim \sum_{n'= \lceil n/2\rceil +1}^{n-1} \p{\max_{\bsig,\ba}\abs{({\cal L}-{\cal K})^{(n+2-n')}_{t,\bsig,\ba}}} \cdot \max_{\bsig,\ba}  
    W^d \sum_{a_{n'}}\abs{({\cal L}-{\cal K})^{(n')}_{t,\bsig,\ba}} \\
    &+ \p{\max_{\bsig,\ba}\abs{({\cal L}-{\cal K})^{(n)}_{t,\bsig,\ba}}} \cdot \max_{\bsig,\ba}  
    W^d \sum_{a_{2}}\abs{({\cal L}-{\cal K})^{(2)}_{t,\bsig,\ba}}.\label{eq:sumtwoloop2}
\end{align}
Note that the second term on the RHS can be controlled using \eqref{eq:sumtwoloop} again, while the first term on the RHS can be handled with the $\contract$ inequality \eqref{yi2oslxj2} for $\cL$-loops and the bound \eqref{wardineq_K} for $\cK$-loops:
\begin{align*}
    W^d \sum_{a_{n'}}\abs{({\cal L}-{\cal K})^{(n')}_{t,\bsig,\ba}} &\le W^d \sum_{a_{n'}}\p{\abs{{\cal L}^{(n')}_{t,\bsig,\ba}}+\abs{{\cal K}^{(n')}_{t,\bsig,\ba}}}\\
    &\le \frac{1}{\eta_t}\min_{k=1}^{n'-1}\left(\max_{\bsig,\bfb}  \left|{\cal L}^{(2k-1)}_{t, \bsig, \bfb}\right|\cdot \max_{\bsig,\bfb}  \left|{\cal L}^{(2n'-2k-1)}_{t, \bsig,\bfb}\right|\right)^{1/2}+ \frac{1}{\eta_t}(W^{-d}B_{t,0})^{n'-2}.
\end{align*}
Then, using \eqref{def:XiL} and \eqref{def:XIL-K}, and setting $k=\lceil n'/2\rceil$, we can bound the RHS of \eqref{eq:sumtwoloop2} by that of \eqref{eq:L-KsimL-K}.

Finally, to establish the estimate \eqref{eq:MG_nloop}, we apply the inequality \eqref{u2jzooi-2}---with $n$ replaced by $2n+2$, and choosing $k=n$, $j=n+1$, and $l=n+2$. This yields that for any $p\ge 1$,
\begin{align*}
\max_{\bsig}\max_{\ba,\ba'}\abs{\left( \cal E\otimes  \cal E \right)^{M}_{t, \bsig, \ba,\ba'}}&\lesssim \max_{\bsig,\ba} W^d \sum_{a_{n+1}\sim a_{2n+2}}  |{\cal L}^{(2n+2)}_{t, \bsig, \ba}|  \lesssim \frac{1}{\eta_t} \p{\max_{\bsig,\bfb}  \left|{\cal L}^{(2n-1)}_{t, \bsig, \bfb }\right|} \cdot  \left(\max_{\bsig,\bfb} \left|  {\cal L}^{(4p)}_{t, \bsig, \bfb}\right|\right)^{\frac{1}{2p}}. 
\end{align*}
Combined with the definition \eqref{def:XiL}, this completes the proof of \eqref{eq:MG_nloop}.
 % for any \smash{$\ba=(a,b),\ba'=(a',b')\in (\Zn)^2$} and $\bsig\in \{+,-\}^2$. ${\cal E}^M$: 
 % (choose $k=n+1$ and $l=n+3$ in \eqref{u2jzooi-2})
 % \[\left( \cB\times  \cB \right)^{(k)}_{t, \bsig, \ba, \ba'} :
 % =   W^d \sum_{[b],[b']} S^{(\sB)}_{[b][b']}{\cal L}^{(2n+2)}_{t, (\bsig\times\overline\bsig)^{(k) },(\ba\times \ba')^{(k )}([b],[b'])}.\]
% \begin{align}
%   \max_{\bsig,\ba} \abs{ {\cal E}^{(\cL-\cK)\times (\cL-\cK),(2)}_{t,\bsig,\ba}}\prec 
%   \sum_{n'= \lceil n/2\rceil +1}^n &\frac{1}{\eta_t}  \max_{{\ba},\bsig}\left|
%   \left({\cal L}-{\cal K}\right)_{t,\bsig,\ba}^{(n+2-n')}\right|
%   \nonumber\\ &\times\min_{k=2}^{n'}\left(\max_{\bsig,\bfb}  \left|{\cal L}^{(2k-3)}_{t, \bsig, \bfb}\right|\cdot \max_{\bsig,\textbf{b}}  \left|{\cal L}^{(2n'-2k+1)}_{t, \bsig, \bfb}\right|\right)^{1/2} .
% \end{align}
%   This one use $|(L-K)^{(m)}|\le |L^{(m)}|+|K^{(m)}|$ for $m\ge n/2+1$, and induction. 
\end{proof}

With the estimates from \Cref{lem:SEforLn}, we now proceed to analyze the integrated loop hierarchy in \eqref{int_K-LcalE}. 
Without loss of generality, the proof can be divided into two cases, depending on whether (i) $1-t \ge \ilambda^2/L^{2}$, or (ii) $1-s \le \ilambda^2/L^{2}$.\footnote{If $1-t < \ilambda^2/L^{2} < 1-s$, then we can add a middle time $u=1-\ilambda^2/L^{2}$ and perform the proofs for case (i) from $s$ to $u$, and for case (ii) from $u$ to $t$.} 
In case (i), we employ the sum-zero operator introduced in \cite{YY_25}. In contrast, case (ii) requires a new approach based on the removal of zero modes from the $\cL$-loops.

\subsection{Proof of Step 3: The case \texorpdfstring{$1-t \ge \ilambda^2/L^{2}$}{1-t >= L-2}}

Throughout this subsection, we always assume that $1-t \ge \ilambda^2/L^{2}$, in which case we have 
\be\label{eq:Bu0asymp}
B_{u,0}\asymp (\ilambda^2+|1-u|)^{-1},\quad  \ell_u\asymp \ilambda\eta_u^{-1/2}+1, \quad \forall u\in [s,t].
\ee 
In Step 2 of the proof of \Cref{lem:main_ind}, we have established an exponential decay of the 2-$G$-loops beyond the scale $\ell_u$ as shown in \eqref{Eq:Gdecay_w}. With \eqref{GijGEX}, we can easily extend this decay to general $G$-loops.

\begin{definition}[Fast decay property]\label{Def_decay}
Let ${\cal A}: (\Zn)^{n}\to \mathbb C$ be an $n$-dimensional tensor for a fixed $n\ge 2$. Given $u\in[s,t]$ and constants $\e,D>0$, we say $\cal A$ satisfies the $(u, \e, D)$-decay property if 
\begin{equation}\label{deccA}
    \max_{i,j\in \qqq{n}}|a_i-a_j|\ge W^{\e}\ell_u \ \implies \ {\cal A}_{\ba}=\OO(W^{-D})\quad \text{for}\quad \ba=(a_1,a_2,\ldots, a_n).
\end{equation}
\end{definition}


It is easy to see that the $G$-loops satisfy the $(u, \e, D)$-decay property for any constants $\e,D>0$ under the estimate \eqref{Eq:Gdecay_w} by using \Cref{lem_GbEXP}. Moreover, in \Cref{Sec:CalK}, we will present a tree representation formula for the $\cK$-loops, which is formed with the $\Theta$-propagators. Thus, the $\cK$-loops also satisfy the $(u, \e, D)$-decay property for any constants $\e,D>0$ by using \eqref{prop:ThfadC}.
%cite \cite{DYYY25_d3} here

\begin{claim}\label{lem_decayLoop} 
Suppose the estimates \eqref{Gt_bound_flow} and \eqref{Eq:Gdecay_w} hold. For any $n\ge 2$, $\bsig\in \{+,-\}^n$, $u\in[s,t]$, and constants $\e,D>0$, the loops \smash{${\cal L}^{(n)}_{u,\bsig,\ba}$} and \smash{${\cal K}^{(n)}_{u,\bsig,\ba}$} satisfy the $(u, \e, D)$-decay property with probability $1-\OO(W^{-D'})$ for any large constant $D'>0$. In other words, we have that
\begin{align}\label{res_decayLK}
\mathbb P\left( \max_{\bsig}  \Big(\left|{\cal L}^{(n)}_{u,\bsig,\ba}\right|+\left|{\cal K}^{(n)}_{u,\bsig,\ba}\right|\Big)\cdot{\bf 1}\left(\max_{  i, j\in \qqq{n} } |a_i-a_j|\ge W^{\e}\ell_u \right) \ge W^{-D}\right)\le W^{-D'}. 
\end{align}
 \end{claim}

Due to the fast decay property of the $G$-loops and $\cK$-loops, when the evolution kernels act on them, we can apply the evolution kernel estimates in \Cref{lem:sum_decay}. In particular, for non-alternating loops $\bsig$, using the estimates in \Cref{lem:SEforLn} and the estimate \eqref{sum_res_2_NAL} in \Cref{lem:sum_decay}, we readily establish the following lemma.  


\begin{lemma}[Non-alternating loops]\label{lem:STOeq_NQ}  
Under the assumptions of \Cref{lem:main_ind}, suppose $1-t\ge \ilambda^2/L^{2}$ and the estimates \eqref{Gt_bound_flow}--\eqref{Eq:Gdecay_w} hold uniformly in $u\in[s,t]$. Fix any $n\ge 2$ and $\bsig\in \{+,-\}^n$ satisfying 
\begin{equation}\label{NALsigm}
\sigma_k=\sigma_{k+1} \quad \text{for some}\quad  k\in\qqq{n}.
\end{equation} 
(Recall that $\sig_{n+1}=\sig_1$ as a convention.) 
Then, we have the following estimate for any fixed $p\ge 1$: 
\begin{align}\label{am;asoiuw}
\max_{\ba\in (\Zn)^n} \big|({\cal L}-{\cal K})^{(n)}_{t, \bsig, \ba} \big|  \big/ (W^{-d}B_{t,0})^n\prec \sup_{u\in [s,t]}
  \left( (W^{-d}B_{t,0})^{\frac 1 6}\wh\Xi^{(\cal L-\cK)}_{u, n}+ \max_{n'=2}^{n-1} \Xi^{({\cal L-\cal K})}_{u,n'}+ \max_{n'=n-1}^{n+1} \Xi^{({\cal L})}_{u,n'} \right)\\
+\sup_{u\in [s,t]}
  \left(\max_{n'=\lceil n/2\rceil + 1}^{n-1} \Xi^{({\cal L-\cal K})}_{u,n+2-n'} \p{\Xi^{({\cal L})}_{u,n_1'} \Xi^{({\cal L})}_{u,n_2'}}^{\frac 12 } +(W^{-d}B_{s,0})^{-\frac{1}{4p}}(\Xi^{(\cal L)}_{u, {2n-1}})^{\frac 1 2}(\Xi^{(\cal L)}_{u, 4p})^{\frac{1}{4p}}\right),\nonumber
\end{align} 
where, as defined below \eqref{eq:L-KsimL-K}, $n'_1=n_2'=n'-1$ if $n'$ is even, and $n'_1=n'-2$ and $n_2'=n'$ if $n'$ is odd. %More generally, this estimate still holds if we replace the LHS by $\sup_{u\in[s,t]}\max_{\ba} |({\cal L}-{\cal K})^{(n)}_{u, \bsig, \ba}| / (W^{-d}B_{u,0})^n$. 
\end{lemma}

\begin{proof}%{Proof of \Cref{lem:STOeq_NQ}}\label{sec:pf_STOeq_NQ}
By \Cref{lem_decayLoop}, the light-weight term $\cal E_{t,\bsig,\ba}^{\Gc,(n)}$ satisfies condition \eqref{deccA0} below for any constants $\e,D>0$. Then, under the assumption \eqref{NALsigm}, using the bound \eqref{bEwGn} together with the evolution kernel estimate \eqref{sum_res_2_NAL} below, we can control the fourth term on the RHS of \eqref{int_K-LcalE} as  
\begin{align*}
\int_{s}^t \left(\mathcal{U}^{(n)}_{u, t, \bsig} \circ \mathcal E^{\Gc,(n)}_{u, \bsig}\right)_{\ba} \dd u &\prec \int_{s}^t \p{\frac{\ilambda^2+|1-u|}{\ilambda^2+|1-t|}}^{n-1} \frac{\p{W^{-d}B_{u,0}}^n}{\eta_u} \p{\wh\Xi^{({\cal L})}_{u, n_1}\cdot \wh\Xi^{({\cal L})}_{u, n_2}}^{\frac12} \dd u \\
&\prec \p{W^{-d}B_{t,0}}^n\int_{s}^t\frac1{\eta_u} \max_{n'=n-1}^{n+1} \Xi^{(\cal L)}_{u,n'}\dd u ,
\end{align*}
where in the second step, we use $\frac{\ilambda^2+|1-u|}{\ilambda^2+|1-t|} B_{u,0} \le B_{t,0}$ and the control parameter in \eqref{def:XiL}.
Using the bounds \eqref{eq:KsimL-K} and \eqref{eq:L-KsimL-K}, together with the assumption \eqref{Eq:L-KGt+IND} on \smash{$(\mathcal{L} - \mathcal{K})^{(n)}_{s, \bsig, \ba}$} and the evolution kernel estimate  \eqref{sum_res_2_NAL}, the first three terms on the RHS of \eqref{int_K-LcalE} can be bounded in the same manner. This yields that 
\begin{align}\label{sahwNQ}
&(\mathcal{L} - \mathcal{K})^{(n)}_{t, \bsig, \ba}\big/ (W^{-d}B_{t,0})^n \prec 
1 + (W^{-d}B_{t,0})^{-n} \int_{s}^t \left(\mathcal{U}^{(n)}_{u, t, \bsig} \circ  \dd  \mathcal{E}^{M,(n)}_{u, \bsig}\right)_{\ba} +   \int_{s}^t \frac1{\eta_u} \max_{n'=2}^{n-1} \Xi^{({\cal L-\cal K})}_{u,n'}\dd u  
     \\
    &+ \int_{s}^t \frac1{\eta_u} \max_{n'=n-1}^{n+1} \Xi^{(\cal L)}_{u,n'}\dd u +  \int_{s}^t \frac1{\eta_u} \max_{n'=\lceil n/2\rceil + 1}^{n-1} \Xi^{({\cal L-\cal K})}_{u,n+2-n'} \p{\Xi^{({\cal L})}_{u,n_1'} \Xi^{({\cal L})}_{u,n_2'}}^{\frac 12 } \dd u + (W^{-d}B_{t,0})^{\frac16} \int_{s}^t \frac{1}{\eta_u}  \wh\Xi^{({\cal L-\cal K})}_{u,n} \dd u  .\nonumber
\end{align}
For the martingale term, by \Cref{lem:DIfREP}, and using \eqref{eq:MG_nloop} together with \eqref{sum_res_2_NAL}, we obtain that
\begin{align*} %\label{sahwNQ2}
(W^{-d}B_{t,0})^{-n}\int_{s}^t \left(\mathcal{U}^{(n)}_{u, t, \bsig} \circ  \dd  \mathcal{E}^{M,(n)}_{u, \bsig}\right)_{\ba} & \prec (W^{-d}B_{t,0})^{-n} \left\{\int_{s}^t 
   \left(\left(
   \mathcal{U}^{(n)}_{u,t,  \bsig}
   \otimes 
   \mathcal{U}^{(n)}_{u,t,  \overline{\bsig}}
   \right) \;\circ  \;
   \left( \cE \otimes  \cE \right)^{M,(n)}
   _{u, \bsig  }
   \right)_{\ba, \ba}\dd u\right\}^{1/2} \nonumber\\
   &\prec (W^{-d}B_{s,0})^{-\frac{1}{4p}}\left\{\int_{s}^t 
   \frac{1}{\eta_u}\Xi^{(\cal L)}_{u, {2n-1}}\p{\Xi^{(\cal L)}_{u, 4p}}^{\frac{1}{2p}}\dd u\right\}^{1/2} 
  \end{align*} 
for any fixed $p\ge 1$. Here, the bound \eqref{eq:MG_nloop} produces the factor $(W^{-d}B_{u,0})^{2n-\frac{1}{2p}}$ in the integrand. Together with the factors produced by \eqref{sum_res_2_NAL}, its positive powers $B_{u,0}^{2n}$ are bounded by $B_{t,0}^{2n}$ using $\frac{\ilambda^2+|1-u|}{\ilambda^2+|1-t|}B_{u,0}\le B_{t,0}$ and $B_{u,0}\le B_{t,0}$, while for the remaining negative power we can only use $B_{u,0}^{-\frac{1}{2p}}\le B_{s,0}^{-\frac{1}{2p}}$, since $B_{s,0}\le B_{u,0}\le B_{t,0}$. Plugging it into \eqref{sahwNQ} and performing the integral over $u$, we conclude \eqref{am;asoiuw}.\end{proof}

%It is clear from the above argument that the same estimate \eqref{am;asoiuw} indeed holds for each $\Xi^{({\cal L-\cal K})}_{u,n}$ with $u\in[s,t]$. Then, a standard $N^{-C}$-net argument allows us to extend it uniformly to all $u\in[s,t]$ and hence conclude \eqref{am;asoiuw} for \smash{$\sup_{u\in[s,t]}\max_{\ba} |({\cal L}-{\cal K})^{(n)}_{u, \bsig, \ba}| \cdot W^{dn}$}.


% First we need to use lemma \ref{Sol_CalL} (the one WITHOUT sum zero property) to prove Step 3 in the case that the loop is not alternative. 

% \begin{lemma}[$\cal L$ bound for non-alternative loop] \textbf{BYYY}

% Assume $\bsig$ is not alternative, i.e., 
% $$
% \exists\; i: \sigma_i=\sigma_{i-1},\quad \sigma_0=\sigma_n
% $$
% Then
% $$
% {\cal L}\prec (W^{-d}B_{0,t})^{n-1}$$
% \end{lemma}

\noindent{\textbf{Alternating cases}}: It remains to deal with the case with alternating signs, where \eqref{NALsigm} does not occur:
\begin{equation}\label{NALsig_diff}
\sigma_k=- \sigma_{k+1}, \quad \forall \ k\in\qqq{n}\, .
\end{equation}
For this purpose, we introduce another key tool---a sum-zero operator ${\cal Q}_t$.

\begin{definition}[Partial sum and sum-zero operator]\label{Def:QtPt} 
Let ${\cal A}: (\Zn)^{n}\to \mathbb C$ be an $n$-dimensional tensor for a fixed $n\in \N$ with $n\ge 2$. Define the partial sum operator ${\cal P}$ as 
$$
\left( {\cal P} \circ {\cal A}\right)_{a_1}:= \sum_{a_i: i\in\qqq{2,n}}  {\cal A}_{\ba},\quad \ba=(a_1,\ldots, a_n). 
$$
We say a tensor $\cal A$ satisfies the \emph{sum-zero property} if $ {\cal P} \circ {\cal A}\equiv 0$. For $t\in[0,1)$, we define the operator $\cal Q_t$ as
\be\label{eq:sumzero_op}
 \left({\cal Q}_{t}\circ {\cal A}\right)_{\ba}:={\cal A}_{\ba}-\left({\cal P} \circ {\cal A}\right)_{a_1}\dthn^{(n)}_{t, \ba},\ee
where the tensor $\dthn^{(n)}_{t, \ba}$ is a mollifier satisfying \be\label{eq:suma1chi}\sum_{a_2,\ldots,a_n}\dthn^{(n)}_{t, \ba}\equiv 1,\quad \forall a_1\in \Zn,
\ee
along with the following estimates for a constant $c>0$: 
\be\label{eq:derv_Theta} \dthn_{t,\ba}^{(n)} \prec (\ell_t^d)^{-(n-1)}\exp\p{-c\sum_{i=2}^n|a_i-a_1|/\ell_t} ,\quad 
\big|\partial_t\dthn_{t,\ba}^{(n)}\big| \prec |1-t|^{-1}(\ell_t^d)^{-(n-1)}\exp\p{-c\sum_{i=2}^n|a_i-a_1|/\ell_t}. \ee  
The detailed form of $\dthn^{(n)}_{t, \ba}$ is not important to us, and we will give an example in \Cref{rmk:choosechi} below. 
With equation \eqref{eq:suma1chi}, we see that \smash{$ {\cal P} \circ  \dthn^{(n)}_{t, \ba}\equiv 1$}, ${\cal P} \circ {\cal Q}_t =0$, and that for any tensor $\cal A$,
\be\label{eq:sum0PA}
{\cal P} \circ {\cal A} \equiv 0\ \ \implies\ \ {\cal P} \circ \left(\varTheta^{(n)}_{t, \bsig}  \circ {\cal A}\right) \equiv 0,
\ee
where we recall that \(\varTheta^{(n)}_{t, \bsig} \) is the operator defined in \Cref{DefTHUST}. In other words, if \({\cal A}\) satisfies the sum-zero property, then so does \smash{$\varTheta^{(n)}_{t, \bsig}  \circ {\cal A}$}. 
\end{definition}


% \begin{definition}
%     \textbf{New (important) $\vartheta$}: 
% We define $\widetilde{\Theta}^{( 1)}_{t}$
% $$
% \widetilde{\Theta}^{( 1)}_{t, ab}= \left[\left(1-tS^{( 1)} \right)^{-1}\right]_{ab},\quad a,b\in \mathbb Z_L
% $$
% where $S^{( 1)}$ is the $S$ matrix we used in $d=1$ paper. Note this function has no polynomial structure.
% Then 
% $$
% \vartheta_{a_1,a_2\cdots a_n}= \prod_{i=2}^n F_{a_1-a_i}, \quad F_{a}=(\sum_{b\in \mathbb Z_L^d}\Theta^{( 1)}_{0, [\|b\|]})^{-1}\cdot \Theta^{( 1)}_{0, [\|a\|]}  
% $$
% where $[x]$ is the largest integer which is not larger than $x$. 

% Note 
% $$
% \sum_a F_a=1,\quad F_a\prec \ell_t^{-d},\quad \ell_t\sim \eta^{-1/2}
% $$
% and 
% $$
% d F_a/dt\prec \frac{1}{\eta}\cdot \ell_t ^{-d},\quad \max_{|a|\gg \ell_t}F_a\approx 0
% $$
% \end{definition}


\begin{example}\label{rmk:choosechi}
To construct the mollifier tensor \(\dthn_{t,\ba}^{(n)}\), we first choose a compactly supported, smooth, non-negative function \smash{$f\in C_c^{\infty}(\R^d)$} that is not identically equal to zero, and rescale it as \smash{\(f_t(a):=\ell_t^{-d}f(|a|/\ell_t)\)} for $a\in \Zn$. We then define, with an appropriate normalization constant $C_{f,t}$,  
\begin{align}
 \dthn^{(n)}_{t, \ba}:=C_{f,t}\cdot \prod_{i=2}^n   f_t(a_i-a_1) ,\quad \forall \ba\in (\Zn)^n .
\end{align}
It is easy to see that this function satisfies the required properties in \eqref{eq:derv_Theta}; in particular, both $\dthn^{(n)}_{t,\ba}$ and $\partial_t\dthn^{(n)}_{t,\ba}$ vanish unless $\max_{i}|a_i-a_1|\lesssim \ell_t$. 
\end{example}



We will use the sum-zero operator to get improved estimates on the terms on the RHS of \eqref{int_K-LcalE} when \eqref{NALsig_diff} holds. 
Roughly speaking, we will decompose a tensor $\cal A$ as \smash{${\cal A}_{\ba}= ({\cal Q}_{t}\circ {\cal A})_{\ba} + \left({\cal P} \circ {\cal A}\right)_{a_1}\dthn^{(n)}_{t, \ba}$} using \eqref{eq:sumzero_op}. For the first part, we can get an improvement by using the evolution kernel estimate \eqref{sum_res_2}, while for the second part, we can apply Ward's identity to ${\cal P} \circ {\cal A}$. 

We first claim that when $\bsig$ satisfies \eqref{NALsig_diff}, the following estimate holds uniformly in $u\in[s,t]$:
\be\label{eq:Ward_typeP}
\left[\cal P\circ (\mathcal{L} - \mathcal{K})^{(n)}_{u, \bsig}\right]_{a_1} \dthn^{(n)}_{u, \ba} \prec   (W^{-d}B_{u,0})^{n}  \Xi^{(\mathcal{L}-\mathcal{K})}_{u, n-1}.
\ee
To see this, we apply Ward's identities in \Cref{lem_WI_K} at the vertex $a_n$ and get  
$$\left[\cal P\circ (\mathcal{L} - \mathcal{K})^{(n)}_{u, \bsig}\right]_{a_1}= \frac{1}{2\ii W^d \eta_u}\left[\cal P\circ (\mathcal{L} - \mathcal{K})^{(n-1)}_{u, \wh\bsig^{(+,n)}}-\cal P\circ (\mathcal{L} - \mathcal{K})^{(n-1)}_{u, \wh\bsig^{(-,n)}}\right]_{a_1}.$$
By \eqref{def:XIL-K}, the two $(\cL-\cK)^{(n-1)}$-loops on the RHS are controlled by 
$$(\mathcal{L} - \mathcal{K})^{(n-1)}_{u, \wh\bsig^{(\pm,n)},\wh \ba^{(n)}} \prec (W^{-d}B_{u,0})^{n-1}\Xi^{(\mathcal{L}-\mathcal{K})}_{u, n-1}.$$ Moreover, due to the fast decay property of the $(\cal L-\cal K)$-loops, the partial sums over the remaining $(n-2)$ vertices lead to an additional \smash{$\ell_u^{d(n-2)}$} factor up to a negligible error $\OO(W^{-D})$. This leads to
\begin{align}\label{jywiiwsoks}
 \left[\cal P\circ (\mathcal{L} - \mathcal{K})^{(n)}_{u, \bsig}\right]_{a_1}  \prec (W^d\eta_u)^{-1}(\ell_u^{d})^{n-2}\cdot (W^{-d}B_{u,0})^{n-1}   \Xi^{(\mathcal{L}-\mathcal{K})}_{u, n-1}.
\end{align}
Together with \eqref{eq:derv_Theta}, it implies that 
%we can bound \smash{$\dthn^{(n)}_{u, \ba}$ by $(\ell_u^d)^{-(n-1)}$}. Together with \eqref{jywiiwsoks}, it implies that 
\begin{align*}
\left[\cal P\circ (\mathcal{L} - \mathcal{K})^{(n)}_{u, \bsig}\right]_{a_1} \dthn^{(n)}_{u, \ba} \prec (W^d\ell_u^d\eta_u)^{-1}\cdot (W^{-d}B_{u,0})^{n-1}\Xi^{(\mathcal{L}-\mathcal{K})}_{u, n-1} \lesssim (W^{-d}B_{u,0})^{n}\Xi^{(\mathcal{L}-\mathcal{K})}_{u, n-1} 
\end{align*}
uniformly in $u\in [s,t]$, where, in the second step, we use $(\ell_u^d\eta_u)^{-1}\lesssim B_{u,0}$ by \eqref{eq:Bu0asymp}.  



It remains to control ${\cal Q}_t \circ (\mathcal{L} - \mathcal{K})^{(n)}_{t, \bsig, \ba}$. For this purpose, we need the following claim on the $(\infty\to \infty)$-norm of the sum-zero operator, which follows easily from \Cref{Def:QtPt} and the estimate \eqref{eq:derv_Theta}.  

\begin{claim}\label{lem_+Q} 
Let ${\cal A}: (\Zn)^{n}\to \mathbb C$ be an $n$-dimensional tensor for a fixed $n\ge 2$. If $\cal A$ satisfies the $(t, \e, D)$-decay property, then we have that 
\begin{align}\label{normQA}
\left\|{\cal Q}_t\circ \cal A\right\|_{\infty} \le W^{C_n \e} \|\cal A\|_\infty +W^{-D+C_n}
\end{align}
for a constant $C_n$ that does not depend on $\e$ or $D$. Furthermore, if $\|\cal A\|_\infty\le W^C$ for a constant $C>0$, then $\cal A_{\ba}-({\cal Q}_t\circ \cal A)_{\ba}=\left({\cal P} \circ {\cal A}\right)_{a_1}\dthn^{(n)}_{t, \ba}$ satisfies the $(t, \e', D')$-decay property for any constants $\e',D'>0$. 
\end{claim}
% \begin{proof}
% The $(t, \e', D')$-decay property of $\cal A_\fa-({\cal Q}_t \circ \cal A)_{\ba}$ follows easily from the exponential decay of $\Theta^{(+,-)}_{t}$ given by \eqref{prop:ThfadC}. To show \eqref{normQA}, it suffices to prove that   
% \be\label{eq:boundPA}
%  \max_{\ba}\left|\left({\cal P} \circ {\cal A}\right)_{[a_1]}\dthn^{(n)}_{t, \ba}\right| \le W^{C_n \e}\|\cal A\|_\infty + W^{-D+C_n}. 
% \ee
% Using the estimate \eqref{prop:ThfadC} and the $(t,\e,D)$-decay property of $\cal A$, we get that 
% \begin{align*}
%    \max_{\ba}\left|\left({\cal P} \circ {\cal A}\right)_{[a_1]}\dthn^{(n)}_{t, \ba}\right| &\prec  \frac{|1-t|^{n-1}}{(\ell_t^d|1-t|)^{n-1}}\cdot (W^\e\ell_t^d)^{(n-1)}\|\cal A\|_\infty + W^{-D}n^d \\
%    &= W^{\e (n-1)}\|\cal A\|_\infty + W^{-D}n^d, 
% \end{align*}
% which concludes \eqref{eq:boundPA}.
% \end{proof}



We derive from equation \eqref{eq_L-Keee} that 
\begin{align} 
\dd{\cal Q}_t \circ (\mathcal{L} - \mathcal{K})^{(n)}_{t, \bsig, \ba} 
&= {\cal Q}_t \circ \left[{\cK}^{(2)} \sim (\mathcal{L} - \mathcal{K})\right]^{(n)}_{t, \bsig, \ba}\dd t 
    + \sum_{l_\mathcal{K}=3}^n {\cal Q}_t \circ \Big[{\cK}^{(\lenk)}\sim (\mathcal{L} - \mathcal{K})\Big]^{(n)}_{t, \bsig, \ba}\dd t + {\cal Q}_t \circ \mathcal{E}^{(\cL-\cK)\times(\cL-\cK),(n)}_{t, \bsig, \ba}\dd t \nonumber
    \\
& + {\cal Q}_t \circ \mathcal{E}^{\Gc,(n)}_{t, \bsig, \ba} \dd t
    + {\cal Q}_t \circ \dd\mathcal{E}^{M,(n)}_{t, \bsig, \ba} 
    - \left[{\cal P} \circ \left(\mathcal{L} - \mathcal{K}\right)^{(n)}_{t,\bsig}\right]_{a_1} \big(\partial_t \dthn_{t,\ba}^{(n)}\big) \dd t.\label{zjuii1}
\end{align} 
Recalling \Cref{DefTHUST}, we can rewrite the first term on the RHS as 
\begin{align}\label{zjuii2}
{\cal Q}_t \circ \Big[{\cK}^{(2)}\sim (\mathcal{L} - \mathcal{K})\Big]_{t, \bsig,\ba}^{(n)}
& = \varTheta^{(n)}_{t, \bsig} \circ \left[{\cal Q}_t \circ (\mathcal{L} - \mathcal{K})_{t, \bsig}^{(n)}\right]_{\ba} 
+ \left[{\cal Q}_t , \varTheta^{(n)}_{t, \bsig} \right] \circ (\mathcal{L} - \mathcal{K})_{t, \bsig, \ba}^{(n)},
\end{align}
where  $[{\cal Q}_t , \varTheta^{(n)}_{t, \bsig} ] = {\cal Q}_t \circ \varTheta^{(n)}_{t, \bsig} - \varTheta^{(n)}_{t, \bsig}\circ \cal Q_t$ denotes the commutator
between ${\cal Q}_t $ and $\varTheta^{(n)}_{t, \bsig}$. Since \({\cal P} \circ {\cal Q}_t = 0\), we notice that the first 5 terms on the RHS of \eqref{zjuii1} satisfy the sum zero property. 
Since \smash{${\cal P}  \circ \dthn^{(n)}_{t,  \ba}\equiv 1$}, we have \smash{${\cal P} \circ \big(\partial_t \dthn_{t,\ba}^{(n)}\big)=0$}, so the last term on the RHS of \eqref{zjuii1} also satisfies the sum zero property.
% \be\label{eq:sumzero666}
%  {\cal P}  \circ \left\{\left[{\cal P} \circ \left(\mathcal{L} - \mathcal{K}\right)^{(n)}_{t,\bsig}\right]_{[a_1]} \partial_t \dthn_{t,\ba}^{(n)}\right\}
% =\left[{\cal P} \circ \left(\mathcal{L} - \mathcal{K}\right)^{(n)}_{t,\bsig}\right]_{[a_1]} {\cal P} \circ \left(\partial_t \dthn_{t,\ba}^{(n)}\right)=0.
% \ee
Next, due to \eqref{eq:sum0PA}, the first term on the RHS of \eqref{zjuii2} also satisfies the sum-zero property. 
Finally, since the LHS of \eqref{zjuii2} has the sum-zero property, the second term on the RHS of \eqref{zjuii2} also satisfies the sum-zero property.
% \begin{align}\label{pqthlk}
%    {\cal P}\circ \big[{\cal Q}_t , \thn^{(n)}_{t,\bsig} \big]\circ (\mathcal{L} - \mathcal{K}) _{t, \bsig, \ba}=0 .
% \end{align} 

With Duhamel's principle, we can derive from \eqref{zjuii1} and \eqref{zjuii2} the following counterpart of \eqref{int_K-LcalE}:
\begin{align}\label{int_K-L+Q}
  {\cal Q}_t\circ   (\mathcal{L} - \mathcal{K})^{(n)}_{t, \bsig, \ba} &= 
    \left(\mathcal{U}^{(n)}_{s, t, \bsig} \circ  {\cal Q}_s\circ  {\cal B}_0(s)\right)_{\ba} +   \int_{s}^t \left(\mathcal{U}^{(n)}_{u, t, \bsig} \circ  {\cal Q}_u\circ \sum_{k=1}^5{\cal B}_k(u) \right)_{\ba} \dd u \nonumber\\
    &+  \int_{s}^t \left(\mathcal{U}^{(n)}_{u, t, \bsig} \circ  {\cal Q}_u\circ \dd\mathcal{E}^{M,(n)}_{u, \bsig}\right)_{\ba} ,
\end{align}
where the tensors $\cal B_i$ for $i\in \qqq{0,5}$ are defined as follows:
\begin{align*}
&{\cal B}_0(s):=(\mathcal{L} - \mathcal{K})^{(n)}_{s, \bsig},\quad {\cal B}_1(u):=\sum_{\lenk=3}^n \Big[\mathcal{K}^{(\lenk)} \sim (\mathcal{L} - \mathcal{K})\Big]^{(n)}_{u, \bsig},\quad {\cal B}_2(u):=\mathcal{E}^{(\mathcal{L} - \mathcal{K}) \times (\mathcal{L} - \mathcal{K}),(n)}_{u, \bsig},\\
&{\cal B}_3(u):=\mathcal{E}^{\Gc,(n)}_{u, \bsig},\quad {\cal B}_4(u):= \big[{\cal Q}_u , \varTheta^{(n)}_{u,\bsig} \big]\circ (\mathcal{L} - \mathcal{K})^{(n)}_{u, \bsig },\quad {\cal B}_5(u):= -\br{{\cal P} \circ\left(\mathcal{L} - \mathcal{K}\right)^{(n)}_{u, \bsig}}\cdot \partial_u  \dthn^{(n)}_{u } . 
\end{align*}
%  \begin{align}\label{int_K-L+Q}
%  & {\cal Q}_t\circ  (\mathcal{L} - \mathcal{K})^{(n)}_{t, \bsig, \ba}   = 
% \left(\mathcal{U}^{(n)}_{s, t, \bsig} \circ  {\cal Q}_s\circ (\mathcal{L} - \mathcal{K})^{(n)}_{s, \bsig}\right)_{\ba}  + \sum_{l_\mathcal{K} = 3}^n \int_{s}^t \left(\mathcal{U}^{(n)}_{u, t, \bsig} \circ  {\cal Q}_u\circ \Big[\mathcal{O}_{\cK}^{(\lenk)} (\mathcal{L} - \mathcal{K})\Big]^{(n)}_{u, \bsig}\right)_{\ba} \dd u \nonumber \\
% &+ \int_{s}^t \left(\mathcal{U}^{(n)}_{u, t, \bsig} \circ  {\cal Q}_u\circ \mathcal{E}^{(n)}_{u, \bsig}\right)_{\ba} \dd u + \int_{s}^t \left(\mathcal{U}^{(n)}_{u, t, \bsig} \circ  {\cal Q}_u\circ \mathcal{W}^{(n)}_{u, \bsig}\right)_{\ba} \dd u + \int_{s}^t \left(\mathcal{U}^{(n)}_{u, t, \bsig} \circ  {\cal Q}_u\circ \dd\mathcal{B}^{(n)}_{u, \bsig}\right)_{\ba} 
%     \nonumber \\
% &+ \int_{s}^t \left(\mathcal{U}^{(n)}_{u, t, \bsig} 
% \circ \left( \left[{\cal Q}_u , \thn^{(n)}_{u,\bsig} \right]\circ (\mathcal{L} - \mathcal{K})^{(n)}_{u, \bsig }\right) \right)_{\ba} \dd u
%    - \int_{s}^t \left(\mathcal{U}^{(n)}_{u, t, \bsig} 
%      \circ \left\{\left[ {\cal P} \circ
%       \left(\mathcal{L} - \mathcal{K}\right)^{(n)}_{u, \bsig}\right] \partial_u\dthn_{u}^{(n)} \right\}\right)_{\ba} \dd u.
% \end{align}
We can control the terms on the RHS of \eqref{int_K-L+Q} by applying the improved evolution kernel estimate \eqref{sum_res_2}, which exploits both the sum-zero and fast-decay properties of the terms on the RHS of \eqref{zjuii1}. Combining this with \eqref{eq:Ward_typeP} and \Cref{lem:STOeq_NQ} for the non-alternating case, we obtain the following inductive bootstrap bounds for the $\Xi$-parameters. The proof proceeds by estimating the RHS of \eqref{int_K-L+Q} %via the evolution kernel estimate \eqref{sum_res_2}, 
in a manner analogous to the proof of \Cref{lem:STOeq_NQ}, using the bounds established in \Cref{lem:SEforLn} together with straightforward controls of the $\mathcal{B}_4$ and $\mathcal{B}_5$ terms, based on the properties in \eqref{eq:derv_Theta}. Hence, we defer the detailed proof to \Cref{subsec:support}.


\begin{lemma}[Inductive bootstrap bound for $\Xi$-parameters]
%\textcolor{blue}{$\Xi$-self bound}]
\label{lem:STOeq_Qt} 
Under the assumptions of \Cref{lem:main_ind}, suppose $1-t\ge \ilambda^2/L^{2}$ and the estimates \eqref{Gt_bound_flow}--\eqref{Eq:Gdecay_w} hold uniformly in $u\in[s,t]$. Then, for any fixed $n\ge 2$ and $p\ge 1$, the following bound holds uniformly in $u\in[s,t]$: 
\begin{align}\label{am;asoi222}
\sup_{v\in[s,u]} \wh\Xi^{({\cal L-\cal K})}_{v,n} &\prec \sup_{v\in[s,u]} \left(\p{W^{-d}B_{s,0}}^{-\frac{1}{4p}}(\Xi^{(\cal L)}_{v, {2n-1}})^{\frac 1 2}(\Xi^{(\cal L)}_{v, 4p})^{\frac{1}{4p}}\right)\\
&+ \sup_{v\in[s,u]} \left(\max_{n'=1}^{n-1} \Xi^{({\cal L-\cal K})}_{v,n'}+ \max_{n'=n-1}^{n+1} \Xi^{({\cal L})}_{v,n'} + \max_{n'=\lceil n/2\rceil + 1}^{n-1} \Xi^{({\cal L-\cal K})}_{v,n+2-n'} \p{\Xi^{({\cal L})}_{v,n_1'} \Xi^{({\cal L})}_{v,n_2'}}^{\frac 12 } \right).\nonumber
\end{align}

% we have the following bound for each fixed $p\ge 1$: 
% \begin{align}\label{am;asoi222}
% \sup_{u\in[s,t]} \wh\Xi^{({\cal L-\cal K})}_{u,n} &\prec \sup_{u\in[s,t]} \left(\p{W^{-d}B_{t,0}}^{-\frac{1}{4p}}(\Xi^{(\cal L)}_{u, {2n-1}})^{\frac 1 2}(\Xi^{(\cal L)}_{u, 4p})^{\frac{1}{4p}}\right)\\
% &+ \sup_{u\in [s,t]} \left(\max_{n'=1}^{n-1} \Xi^{({\cal L-\cal K})}_{u,n'}+ \max_{n'=n-1}^{n+1} \Xi^{({\cal L})}_{u,n'} + \max_{n'=\lceil n/2\rceil + 1}^{n-1} \Xi^{({\cal L-\cal K})}_{u,n+2-n'} \p{\Xi^{({\cal L})}_{u,n_1'} \Xi^{({\cal L})}_{u,n_2'}}^{\frac 12 } \right).\nonumber
% \end{align}
\end{lemma}

 

Now, similar to the argument in Section 5.6 of \cite{YY_25}, we will iterate the bootstrap bound \eqref{am;asoi222} to obtain the sharp $L^\infty$-bound \eqref{Eq:LGxb} on the $G$-loops in the regime $1-t\ge \ilambda^2/L^{2}$. That is, for any fixed $n\in \N$, 
\be\label{eq:Xiiter}
\sup_{u\in [s,t]} \wh\Xi^{(\cL)}_{u,n} \prec 1.
\ee
Observe that when $1-s>\ilambda^2$, we have \(B_{u,0}\asymp |1-u|^{-1}\) for all $u\in [s,1-\ilambda^2]$. Consequently, the $G$-loop bound \eqref{Eq:LGxb} follows directly from \eqref{lRB1}. 
Hence, in the remainder of the proof, it suffices to consider the case $1-t\le 1-s \le \ilambda^2$, where \(W^{-d}B_{u,0}\asymp (\ilambda^2W^d)^{-1}\) for all \(u\in [s,t].\) %When $1-s \le \ilambda^2$, 
First, the averaged local law \eqref{Gt_avgbound_flow} gives that \smash{$\wh\Xi_{u,1}^{(\cL-\cK)}\prec 1=:\Xi_{u,1}^{(\cL-\cK)}$} uniformly for $u\in[s,t]$. Second, by the $\cK$-loop bound \eqref{eq:bcal_k}, we have %\smash{$\wh\Xi^{(\cal L)}_{u,n}$ and $\wh\Xi^{(\cal L-\cal K)}_{u,n}$} control each other via
\begin{align}\label{rela_XILXILK}
    \wh\Xi^{(\cal L)}_{u,n}\prec 1+(\ilambda^2W^{d})^{-1} \wh\Xi^{({\cal L-\cal K})}_{u, n},\quad \quad \wh\Xi^{({\cal L-\cal K})}_{u, n}\prec  (\ilambda^2W^{d})\left(\wh\Xi^{(\cal L)}_{u,n}+1\right),\quad \text{for}\quad n\ge 2 .
\end{align}
Moreover, the a priori $G$-loop bound \eqref{lRB1} provides the following initial estimates, uniform in $u\in[s,t]$:
\begin{align}\label{sef8w483r324}
  \wh\Xi^{(\cal L)}_{u,n}\prec (\eta_s/\eta_u)^{n-1}, \quad 
 \wh\Xi^{({\cal L-\cK})}_{u, n}\prec (\eta_s/\eta_u)^{n-1}\cdot (\ilambda^2W^{d}).
\end{align}
We then introduce the control parameter
\begin{equation}\label{adsyzz0s8d6}
\Psi(n,k;s,t):= (\ilambda^2W^{d})^{3/4} +\p{{\eta_s}/{\eta_t}}^{n-1}(\ilambda^2W^{d})^{1-k/8}.
\end{equation} 
% \begin{equation}\label{adsyzz0s8d6}
% \Psi(n,k;s,t):=  W^{3d/4}+\p{\frac{\eta_s}{\eta_t}}^{n-1}W^{d(1-k/8)}.
% \end{equation} 
%Note that when $k\ge 1$, $\Psi_{u}(n,k;[s,t])$ does not depend on $u$. % $\Psi_{u}(n,k;[s,t])\equiv \Psi_{t}(n,k;[s,t])$ for all $u\in[s,t]$. 
The iteration will proceed simultaneously in the indices $n$ and $k$. The outcome of each step is summarized in the following lemma, whose proof---being analogous to that of (5.109) in \cite{YY_25}---is deferred to \Cref{subsec:support}.

%The iterations will be performed in both $n$ and $k$ at the same time. We summarize the result of each iteration in the next lemma, whose proof will be deferred to \Cref{sec:pf_iterations} because it is similar to that for equation (5.109) in Section 5.6 of \cite{YY_25}.  

\begin{lemma}\label{lem:iterations}
In the setting of \Cref{lem:main_ind}, suppose \eqref{lRB1}--\eqref{Eq:Gdecay_w} and \eqref{am;asoi222} hold uniformly for all $u\in[s,t]$ with $ 1-s \le \ilambda^2$. Fix any $(n,k)\in \N^2$ with $n\ge 2$ and $k\ge 1$. %Suppose the following estimate holds 
Assume that, uniformly for $u\in[s,t]$, 
    \be\label{eq:iteration_induc}
    \sup_{v\in[s,u]}\wh \Xi^{({\cal L-\cK})}_{v,r}\prec   \Psi(r,l;s,u)\ee
    %for all $(r,l)\in \mathbb N^2$ in the set
holds for all index pairs $(r,l)\in \{(r, k): 2\le r\le n-1\}\cup \{(r, k-1): 2\le r\le n+2\}.$ %satisfying one of the following conditions: (1) $l=k$ and $2\le r \le n-1$; (2) $l=k-1$ and $2\le r\le n+2$. 
Then, \eqref{eq:iteration_induc} also holds for $(r,l)=(n, k)$. %we have the following estimate uniformly in $u\in[s,t]$:
% \be\label{eq:iteration_improve}
    % \sup_{v\in[s,u]}\wh \Xi^{({\cal L-K})}_{v,n}\prec   \Psi(n,k;s,u).\ee
\end{lemma}

Roughly speaking, this lemma states that if we have already established a ``good'' bound for all shorter $G$-loops of length $r\le n-1$, along with a ``weaker'' bound for $G$-loops of length $r\le n+2$, then we can derive the ``good'' bound for all $G$-loops of length $n$. 
Using \Cref{lem:iterations}, we apply a simple iterative argument to complete the proof of \eqref{Eq:LGxb} in Step 3 for the case $1-t\ge \ilambda^2/L^{2}$. 

\begin{proof}[\bf Proof of \eqref{Eq:LGxb} when $1-t\ge \ilambda^2/L^{2}$]
As discussed above, it remains to deal with the case $1-s \le \ilambda^2$. By \eqref{sef8w483r324}, we initially have a weak bound for $G$-loops of arbitrarily large lengths, meaning that \eqref{eq:iteration_induc} holds with $l=0$ for every fixed $r\in \N$. Applying \Cref{lem:iterations} once, we obtain a slightly improved bound \eqref{eq:iteration_induc} for $r=2$ and $l=1$. Then, continuing the iteration in $r$ while keeping $l=1$ fixed, we establish the bound \eqref{eq:iteration_induc} for each fixed $r\in \N$ with $l=1$. 
Next, applying the iteration in \Cref{lem:iterations} again yields an even stronger bound \eqref{eq:iteration_induc} with $r=2$ and $l=2$. Repeating the iteration in $r$ with $l=2$ fixed, we further establish the bound \eqref{eq:iteration_induc} for every fixed $r\in \N$ with $l=2$. This process continues, progressively improving the bound to \eqref{eq:iteration_induc} for each fixed $r\in \N$ with $l=3$, and so forth.

For any given $(n,k)\in \N^2$, by repeating the above procedure for $\OO(1)$ times, we conclude that the estimate \eqref{eq:iteration_induc} holds for $(r,l)=(n,k)$.  
In particular, if we choose $\fc_d$ in condition \eqref{con_st_ind} sufficiently small and take $k$ large enough so that $(\eta_s/\eta_t)^{n-1}\cdot  (\ilambda^2W^{d})^{1-k/8} \ll (\ilambda^2W^{d})^{3/4}$, then we get from \eqref{eq:iteration_induc} that 
\begin{equation}\label{eq:iterative_pf_largeeta}\sup_{u\in[s,t]}\wh\Xi^{({\cal L-\cK})}_{u,n}\prec \Psi(n,k;s,t) \lesssim (\ilambda^2W^{d})^{3/4}.\end{equation}
Together with \eqref{rela_XILXILK}, this implies the estimate \eqref{eq:Xiiter}, % under the condition \eqref{con_st_ind}, provided $\fc_d$ is chosen sufficiently small. 
which completes the proof of the bound \eqref{Eq:LGxb}.
\end{proof}



\subsection{Proof of Step 3: The case \texorpdfstring{$1-s \le \ilambda^2/L^{2}$}{1-s <= L-2}}\label{sec:1-s<L-2}


In this setting, we have $\ell_u\equiv L$ for all $u\in[s,t]$. Here, a key difference from the $d = 2$ case in \cite{DYYY25} arises in the intermediate regime $\ilambda^2/L^{d} \leq 1 - t \leq \ilambda^2/L^{2}$, where the polynomial decay mode $\ilambda^{-2} / (|a - b| + 1)^{d - 2}$ and the zero mode $(L^d |1 - t|)^{-1}$ mix in the quantity $B_{t, |a - b|}$ (recall \eqref{eq_B_param}).
In this regime, we lose the fast decay property for $G$-loops that is required by \Cref{lem:sum_decay} below, and instead the weaker evolution kernel estimate \eqref{sum_res_Ndecay} becomes relevant. This estimate introduces a factor of $(1 - s)/(1 - t)$, which transforms the zero mode $(L^d |1 - s|)^{-1}$ at time $s$ into the zero mode $(L^d |1 - t|)^{-1}$ at time $t$. However, this transformation breaks the polynomial decay mode, potentially destabilizing the tree approximation at time $t$.
To address this difficulty, we observe that the two modes actually propagate independently along the flow. Motivated by this, we introduce a zero-mode-removing operator, which decomposes the $(\cL - \cK)$-loops into two parts: (1) the zero-mode components, which can be controlled using Ward's identities, and (2) the zero-mode-free components, whose evolution kernel satisfies sharper estimates thanks to the bound \eqref{prop:ThfadC0} (see \eqref{sum_res_Ndecay_nonzero} below).



\begin{definition}[Zero-mode-removing operators]\label{def;zero_mode_remove}
Let ${\cal A}: (\Zn)^{n}\to \mathbb C$ be an arbitrary $n$-dimensional tensor for a fixed $n\ge 1$. Define $P^{(i)}$ as the \emph{partial averaging operator} with respect to the $i$-th index of the loop: 
$$
(P^{(i)}\circ {\cal A})_{(a_1,\ldots, a_{i-1},b,a_{i+1},\ldots,a_n)}=L^{-d} \sum_{a_i}{\cal A}_{(a_1,\ldots, a_{i-1},a_i,a_{i+1},\ldots,a_n)},\quad \forall i\in \qqq{n}, \ b\in \Zn.$$
Correspondingly, we define the \emph{zero-mode-removing} operator as $Q^{(i)}:= I-P^{(i)}.$ Furthermore, for any subset $A\subset \qqq{n}$, we denote 
$$P^{(A)}:=\prod_{i\in A} P^{(i)},\quad \text{and}\quad  
Q^{(A)}:=\prod_{i\in A} Q^{(i)}.
$$
As a convention, when $A=\emptyset$, we define $P^{(\emptyset)}$ and $Q^{(\emptyset)}$ as the identity operator. 
Note that for any subsets $A$ and $A'$, the operators $P^{(A)}$, $P^{(A')}$, $Q^{(A)}$, and $Q^{(A')}$ all commute with each other. 
%Define non-zero-mode Loop $\mathring{\cal L}$
%$$
% \mathring{\cal L}_{{\ba},\bsig}:= (Q^{(A_0)}\cdot {\cal L})_{{\ba},\bsig}, \quad A_0=A_0(\bsig)
%$$
%where
%\begin{align}\label{def:A0sigma}
%     A_0(\bsig):=\{i: \sigma_{i-1}\ne \sigma_i\} ,\quad \sigma_0:=\sigma_n 
%\end{align} 

\end{definition}

By definition, the $(\infty\to \infty)$-operator norm of $Q^{(i)}$ is trivially bounded:
\begin{align}\label{normQA2}
\big\|Q^{(i)}\circ \cal A\big\|_{\infty} \le 2\|\cal A\|_\infty .
\end{align}
Given any $\bsig\in \{+,-\}^n$, we denote
\( I_{\mathrm{diff}}(\bsig):=\{i\in\qqq{n}:\sig_i\ne\sig_{i+1}\},\)
where we again adopt the cyclic convention $\sig_{n+1}=\sig_1$. Given any $\cL$- or $\cK$-loop, we can express it as a linear combination of $(Q^{(A)}\circ \cL)_{t,\bsig,\ba}$ or $(Q^{(A)}\circ \cK)_{t,\bsig,\ba}$ loops with $A\supset I_{\mathrm{diff}}(\bsig)$, by repeatedly applying \eqref{WI_calL} and \eqref{WI_calK}. 
This is summarized in \Cref{lem: newPQ}, whose proof is deferred to \Cref{subsec:support} below.

% \begin{align}
% \big(\cal Q^{(1)}\circ {\cal L}^{(3)}\big)_{t,\bsig,\ba}
% = &\big(Q^{(\{1,2\})}\circ {\cal L}^{(3)}\big)_{t,\bsig,\ba}+\frac{\big(Q^{(\{1 \})}\circ {\cal L}^{(2)}\big)_{t,(a_1,a_3),(-,+)}-
% \big(Q^{(\{1 \})}\circ {\cal L}^{(2)}\big)_{t,(a_1,a_3),(-,-)}}{2\ii N\eta_t},
% \end{align}

%The proof of \Cref{lem: newPQ} is straightforward by using Ward's identities \eqref{WI_calK} and \eqref{WI_calL}, and we will provide a proof in \Cref{Sec:pf_newPQ} below for reader's convenience. 

\begin{lemma}\label{lem: newPQ}
For any fixed $n\in \N$, $\bsig\in \{+,-\}^n$, $\ba\in(\Zn)^n$, and subset $A\subset\qqq{n}$, we have the expansion: 
\begin{align}\label{yurenAL}
\left(Q^{(A)}\circ {\cal L}^{(n)}\right)_{t,\bsig,\ba}=\left(Q^{(A_n)}\circ {\cal L}^{(n)}\right)_{t,\bsig,\ba}+\sum_\al \frac{\xi_{\al}}{(2\ii N\eta_t)^{n-k_\al}} \left( Q^{(A_\al)}\circ {\cal L}^{(k_\al)}\right)_{t,\bsig_{\al},\ba_{\al}},
 \end{align}
where $A_n:=A\cup I_{\mathrm{diff}}(\bsig)$, $\{\al\}$ denotes a collection of $\OO(1)$ many labels, $1\le k_\al \le n-1$, $\ba_{\al}\in (\Zn)^{k_\al}$ consists of a subset of indices in $\ba$, $\bsig_\al\in \{+,-\}^{k_\al}$, $\qqq{k_\al}\supset A_\al \supset I_{\mathrm{diff}}(\bsig_\al)$, and $\xi_{\al}$ denote some deterministic and integer-valued coefficients of order $\OO(1)$. The same expansion also holds for the $\cK$-loop:
\begin{align}\label{yurenAK}
\left(Q^{(A)}\circ {\cal K}^{(n)}\right)_{t,\bsig,\ba}=\left(Q^{(A_n)}\circ {\cal K}^{(n)}\right)_{t,\bsig,\ba}+ \sum_\al \frac{\xi_{\al}}{(2\ii N\eta_t)^{n-k_\al}} \left( Q^{(A_\al )}\circ {\cal K}^{(k_\al)}\right)_{t,\bsig_\al,\ba_\al }.
 \end{align}
As a special case that is of particular importance to us, the above expansions hold for $A=\emptyset$. 
    % For any fixed $n\in \mathbb N$, ${\ba}\in (\mathbb Z_L^d)^n$,  $\bsig\in \{+,-\}^n$ and $A\subset [[1,n]]$,  there exist 
    % $$
    % \left({\ba}_\alpha,\; \bsig_\alpha, \; 
    % A _\alpha, \;
    % \xi_\alpha,\; n_\alpha\right), \quad \alpha \le C_n
    % $$
    % such that
    % $$
    % {\ba} _{\;\alpha}\in  (\mathbb Z_L^d)^{n_\alpha}, \quad 
    % \bsig_\alpha\in  \{+,-\}^{n_\alpha},\quad
    % %A^p_\alpha, A^q_\alpha\subset
    % A _\alpha \subset[[1,n_\alpha]],\quad 
    %  \xi_ \alpha=\pm 1, \quad  1\le n_\alpha\le n $$
    % and 
    % $$
    %  A _\alpha \supset B_\alpha:=\left\{1\le i\le n_\alpha: (\sigma_\alpha) _i\ne (\sigma_\alpha)_{i-1} \right\}
    %  $$
    % and   
    \end{lemma}



% Our primary strategy is to apply mathematical induction with respect to the pair
% \[
% (n, |A|),
% \]
% where the induction proceeds by increasing \( n \) and decreasing \( |A| \). First it is trivial for $n=1$. For $n\ge 2$, define $B $ as
%     \begin{align}\label{def:A0sigma}
%      B :=B (\bsig):=\{i: \sigma_{i-1}\ne \sigma_i\} ,\quad \sigma_0:=\sigma_n 
%  \end{align} 
%  If $ B\subset A$, then just choose 
%  $$
%  \alpha=1, \quad \xi_\alpha =1, \quad A_1=A, \quad {\ba}_\alpha= {\ba},\quad \bsig_\alpha=\bsig
%  $$
%  Now if $B\not\subset A$, assume that there exists $i\in B\setminus A$, then we write 
%  $$
%  \left(Q^{(A )}\circ {\cal L}\right)_{{\ba},\bsig}= \left(Q^{(A\cup \{i\} )}\circ {\cal L}\right)_{{\ba},\bsig}+\left(Q^{(A )}\circ P^{(i)}\circ{\cal L}\right)_{{\ba},\bsig}
%  $$
 


Given any $n\ge 2$ and $A\subset\qqq{n}$, we derive from equation \eqref{eq_L-Keee} that 
\begin{align} 
\dd Q^{(A)} \circ (\mathcal{L} - \mathcal{K})^{(n)}_{t, \bsig, \ba} 
&= Q^{(A)} \circ \left[\varTheta_{t,\bsig}^{(n)} \circ (\mathcal{L} - \mathcal{K})^{(n)}_{t, \bsig}\right]_{\ba} \dd t
    + \sum_{l_\mathcal{K}=3}^n Q^{(A)} \circ \Big[{\cK}^{(\lenk)}\sim (\mathcal{L} - \mathcal{K})\Big]^{(n)}_{t, \bsig, \ba}\dd t \nonumber
    \\
&+ Q^{(A)} \circ \mathcal{E}^{(\cL-\cK)\times(\cL-\cK),(n)}_{t, \bsig, \ba}\dd t  + Q^{(A)} \circ \mathcal{E}^{\Gc,(n)}_{t, \bsig, \ba} \dd t
    + Q^{(A)} \circ \dd\mathcal{E}^{M,(n)}_{t, \bsig, \ba} .\label{eq_L-Keee_nonzeromode}
\end{align} 
Due to the translation invariance of $S^{(\sB)}$ and $M^{(\sig_i,\sig_{i+1})}$, they both have an eigenvector $\mathbf e$ with $\mathbf e(a)\equiv L^{-d/2}$ for all $a\in \Zn$. As a consequence, the operators $\varTheta^{(n)}$ and $\mathcal{U}^{(n)}$ defined in \eqref{def:op_thn} and \eqref{def_Ustz} commute with the operators $P^{(i)}$, and hence also commute with the operators $Q^{(A)}$ for all $A\subset \qqq{n}$. Hence, the first term on the RHS of \eqref{eq_L-Keee_nonzeromode} can be written as 
\[\left[ \varTheta_{t,\bsig}^{(n)} \circ Q^{(A)} \circ (\mathcal{L} - \mathcal{K})^{(n)}_{t, \bsig}\right]_{\ba}.\]
Then, applying Duhamel's principle to the equation \eqref{eq_L-Keee_nonzeromode}, we derive the following integral equation: 
\begin{align}\label{iisuwjyys}
   Q^{(A )}\circ (\mathcal{L} - \mathcal{K})^{(n)}_{t, \bsig, \ba} 
   &=    \left(Q^{(A)}\circ \mathcal{U}^{(n)}_{s, t, \bsig} \circ  {\cal B}_0(s)\right)_{\ba} +   \int_{s}^t \left( Q^{(A)}\circ\mathcal{U}^{(n)}_{u, t, \bsig} \circ \sum_{k=1}^3{\cal B}_k(u) \right)_{\ba} \dd u \nonumber\\
   & +  \int_{s}^t \left(Q^{(A)}\circ\mathcal{U}^{(n)}_{u, t, \bsig} \circ   \dd\mathcal{E}^{M,(n)}_{u, \bsig}\right)_{\ba} , 
   % \\
   % &
   %  \left(Q^{(A )}\circ\mathcal{U}_{s, t, \bsig} \circ (\mathcal{L} - \mathcal{K})_{s, \bsig}\right)_{{\ba}} \nonumber \\
   %  &+ \sum_{l_\mathcal{K} > 2} \int_{s}^t \left(Q^{(A )}\circ\mathcal{U}_{u, t, \bsig} \circ \Big[\mathcal{K} \sim (\mathcal{L} - \mathcal{K})\Big]^{l_\mathcal{K}}_{u, \bsig}\right)_{{\ba}} du \nonumber \\
   %  &+ \int_{s}^t \left(Q^{(A )}\circ\mathcal{U}_{u, t, \bsig} \circ \mathcal{E}^{((\mathcal{L} - \mathcal{K}) \times (\mathcal{L} - \mathcal{K}))}_{u, \bsig}\right)_{{\ba}} du  \nonumber \\
   %  &+ \int_{s}^t \left(Q^{(A )}\circ\mathcal{U}_{u, t, \bsig} \circ \mathcal{E}^{(\widetilde{G})}_{u, \bsig}\right)_{{\ba}} du + \int_{s}^t \left(Q^{(A )}\circ\mathcal{U}_{u, t, \bsig} \circ \mathcal{E}^{(M)}_{u, \bsig}\right)_{{\ba}}  
\end{align}
where we use the notations in \eqref{int_K-L+Q}. %Given any $\bsig \in \{+,-\}^{n}$, 
If $A\supset I_{\mathrm{diff}}(\bsig)$, then the new kernel \smash{$Q^{(A)}\circ \mathcal{U}^{(n)}_{s, t, \bsig}$} has a better $(\infty\to \infty)$-norm estimate (as stated in \Cref{lem:sum_decay_nonzero} below) than the original kernel \smash{$\mathcal{U}^{(n)}_{s, t, \bsig}$} when $1-s\le \ilambda^2/L^{2}$. In fact, corresponding to each $i\in I_{\mathrm{diff}}(\bsig)$, the $(\infty\to \infty)$-norm estimate of \smash{$\mathcal{U}^{(n)}_{s, t, \bsig}$} will be weaker by an $\eta_s/\eta_t$ factor.  

Combining the expansions \eqref{yurenAL} and \eqref{yurenAK} (for the case $A=\emptyset$), the equation \eqref{iisuwjyys}, the evolution kernel estimate in \Cref{lem:sum_decay_nonzero}, and the estimates in \Cref{lem:SEforLn}, we can establish a result similar to that in \Cref{lem:STOeq_Qt}. We postpone the proof of \Cref{lem:STOeq_Qt_nonzero} to \Cref{subsec:support} below.

\begin{lemma}\label{lem:STOeq_Qt_nonzero} 
Under the assumptions of \Cref{lem:main_ind}, suppose $1-s\le \ilambda^2/L^{2}$ and the estimates \eqref{Gt_bound_flow}, \eqref{Gt_avgbound_flow}, and \eqref{Eq:Gdecay_w} hold uniformly in $u\in[s,t]$. 
%Then, for any fixed $n\ge 2$, we have that: 
% \begin{align}\label{am;asoi222_smalletacase}
% \sup_{u\in[s,t]} \wh\Xi^{({\cal L-\cal K})}_{u,n} &\prec \sup_{u\in[s,t]} \left( (W^{-d}B_{t,0})^{-\frac{1}{4p}}(\Xi^{(\cal L)}_{u, {2n-1}})^{\frac 1 2}(\Xi^{(\cal L)}_{u, 4p})^{\frac{1}{4p}}\right)\\
% &+\sup_{u\in [s,t]}
%   \left(\max_{n'=1}^{n-1} \Xi^{({\cal L-\cal K})}_{u,n'}+ \max_{n'=n-1}^{n+1} \Xi^{({\cal L})}_{u,n'} + \max_{n'=\lceil n/2\rceil + 1}^{n-1} \Xi^{({\cal L-\cal K})}_{u,n+2-n'} \p{\Xi^{({\cal L})}_{u,n_1'} \Xi^{({\cal L})}_{u,n_2'}}^{\frac 12 } \right).\nonumber
% \end{align}
Then, for any fixed $n\ge 2$ and $p\ge 1$, the bound \eqref{am;asoi222} holds uniformly in $u\in[s,t]$. 
\end{lemma}

Now, we are ready to complete the proof of Step 3 for the case $1-s\le \ilambda^2/L^{2}$.


\begin{proof}[\bf Proof of \eqref{Eq:LGxb} when $1-s\le \ilambda^2/L^{2}$]
Given \Cref{lem:STOeq_Qt_nonzero}, we can prove a similar iterative result to \Cref{lem:iterations}, but with a different control parameter defined as follows: 
\begin{equation}\label{eq:psipara_smalletacase}
\Psi(n,k;s,t):=  (W^{-d}B_{s,0})^{-3/4}+(\eta_s/\eta_t)^{n-1}(W^{-d}B_{s,0})^{k/8-1}.
% \times
%     \begin{cases}
%         (W^{-d}B_{u,0})^{-1} , &  \ \text{if}\ k=0\\
%         (W^{-d}B_{s,0})^{k/8-1}, &   \ \text{if}\  k\ge 1 
%     \end{cases}.
\end{equation} 
%(Note that when $k\ge 1$, $\Psi_{u}(n,k;s,t)$ does not depend on $u$.) 
More precisely, suppose the estimate \eqref{eq:iteration_induc} holds uniformly in $u\in[s,t]$ for all index pairs $(r,l)\in \{(r, k): 2\le r\le n-1\}\cup \{(r, k-1): 2\le r\le n+2\}$. 
% \be\label{eq:iteration_induc_smalleta}
% \cor    \sup_{v\in[s,u]}\wh \Xi^{({\cal L-K})}_{v,r}\prec   \Psi(r,l;s,u).
% \ee
%for all $(r,l)$ satisfying one of the following conditions: (1) $l=k$ and $2\le r \le n-1$; (2) $l=k-1$ and $2\le r\le n+2$. 
Then, we have the following estimate uniformly in $u\in[s,t]$:
\be\label{eq:iteration_improve_smalleta}
    \sup_{v\in[s,u]}\wh \Xi^{({\cal L-\cK})}_{v,n}\prec   \Psi(n,k;s,u).\ee
Since the proof of \eqref{eq:iteration_improve_smalleta} is very similar to that for \Cref{lem:iterations} by using \Cref{lem:STOeq_Qt_nonzero}, we omit the details. 
With this result, performing exactly the same iterative argument as in the case $1-t\ge \ilambda^2/L^{2}$ (i.e., the argument around \eqref{eq:iterative_pf_largeeta}), we can conclude \eqref{Eq:LGxb} for the case $1-s\le \ilambda^2/L^{2}$.
\end{proof}





\subsection{Proof of Step 4}

We again divide the proof of \eqref{Eq:L-KGt-flow} into two cases according to whether $1-t \ge \ilambda^2/L^{2}$ or $1-s \le \ilambda^2/L^{2}$. In the former case, we apply \Cref{lem:STOeq_Qt} established in Step 3, while in the latter we apply \Cref{lem:STOeq_Qt_nonzero}. 
At this step, using the sharp $G$-loop bound \eqref{Eq:LGxb} together with the averaged local law \eqref{Gt_avgbound_flow}, we may choose the optimal parameters \smash{$\Xi^{(\cal L)}_{v,n'}=1$} for all $G$-loops and \smash{$\Xi^{(\cal L-\cK)}_{v,1}=1$}. Then, from \eqref{am;asoi222}, we get the following bound for any fixed $n\ge 2$ and $p\ge 1$: 
 \begin{align} \label{saww02}
 \sup_{v\in[s,u]}\wh\Xi_{v,n}^{(\mathcal{L}-\mathcal{K})}  \prec (W^{-d}B_{s,0})^{-\frac{1}{4p}}+ \sup_{v\in [s,u]}
  \left(\max_{n'=2}^{n-1}\Xi^{({\cal L-\cal K})}_{v,n'} \right).
\end{align}
First, taking $n=2$ in \eqref{saww02}, the second term on the RHS vanishes. Since $p$ can be chosen arbitrarily large, we conclude that \smash{\(\sup_{v\in[s,u]}\wh\Xi_{v,2}^{(\mathcal{L}-\mathcal{K})}\prec 1.\)} 
Starting from this base case, we can derive that \smash{\(\sup_{v\in[s,u]}\wh\Xi_{v,n}^{(\mathcal{L}-\mathcal{K})}\prec 1\)} 
for any fixed $n\in\N$ by applying \eqref{saww02} inductively in $n$. This completes the proof of \eqref{Eq:L-KGt-flow} in Step 4.
%The proof of \eqref{Eq:L-KGt-flow} in the case $1-s \le \ilambda^2/L^{2}$ is similar by using \Cref{lem:STOeq_Qt_nonzero} instead.

\subsection{Evolution kernel estimates}\label{ks-statements}

%cite \cite{DYYY25_d3} here
In the above proof, we have used the following estimates on the evolution kernel defined in \Cref{DefTHUST}. Their proofs are postponed to \Cref{ks}. We first have an easy bound on the $({\infty\to \infty})$-norm of \smash{$\mathcal U^{(n)}_{s,t,\bsig}$}. 

 \begin{lemma}\label{lem:sum_Ndecay}
	 	Let ${\cal A}: (\Zn)^{n}\to \mathbb C$ be an $n$-dimensional tensor for a fixed $n\ge 2$. Then, for any $0\le s \le t < 1$, we have that 
	 	\begin{align}\label{sum_res_Ndecay}
		 		\| {\cal U}^{(n)}_{s,t,\bsig}\;\circ {\cal A}\|_{\infty} \le \left(  \frac{1-s}{1-t}\right)^{n }\cdot \|{\cal A}\|_{\infty} , 
		 	\end{align}
	 	where the $L^\infty$-norm of ${\cal A}$ is defined as $\|{\cal A}\|_{\infty}=\max_{\ba\in (\Zn)^n}|\cal A_{\ba}|$.
	 \end{lemma}


If the tensor $\cal A$ exhibits faster-than-polynomial decay on scales larger than $\ell_s$, we show that the $({\infty \to \infty})$-norm of the evolution kernel \smash{${\cal U}^{(n)}_{s,t,\bsig}$} satisfies a better bound. This bound can be further improved when $\bsig$ is non-alternating or when $\cal A$ satisfies a sum-zero property.


\begin{lemma}\label{lem:sum_decay}
	Let ${\cal A}: (\Zn)^{n}\to \mathbb C$ be an $n$-dimensional tensor for a fixed $n\ge 2$. Suppose it satisfies the following fast-decay property for some small constant $\e\in(0,1)$ and large constant $D>1$:  
	\begin{equation}\label{deccA0}
	|\cal A_{\ba}|\le W^{-D}, \quad \forall\ \ba=(a_1,\ldots, a_n)\in (\Zn)^{n} \ \ \text{with}\ \	\max_{i,j\in \qqq{n}}|a_i-a_j|\ge W^{\e}\ell_s  \, .
	\end{equation}
	Fix any $0\le s \le t \le 1 -\ilambda^2/L^{2}$ such that $(1-t)/(1-s)\ge W^{-1}$. 
Then, there exists a constant $C_n>0$ that does not depend on $\e$ or $D$ such that the following bound holds:
	\begin{align}\label{sum_res_1}
		\left\|{\cal U}^{(n)}_{s,t,\bsig} \circ {\cal A}\right\|_\infty \le W^{C_n\e}\frac{\ell_t^2 }{\ell_s^2 }\p{\frac{\ilambda^2+|1-s|}{\ilambda^2+|1-t|}}^{n}\|{\cal A}\|_\infty
		+W^{-D+C_n} . %\left(\frac{\ell_s^2 |1-s|}{\ell_t^2 |1-t|}\right)^{n} .
	\end{align}
%where the $L^\infty$-norm of ${\cal A}$ is defined as $\|{\cal A}\|_{\infty}=\max_{\ba\in (\Zn)^n}|\cal A_{\ba}|$.
This estimate can be further improved in the following cases:

	\begin{itemize}
		\item[(I)] 
If $\bsig$ is non-alternating, i.e.,  $\sigma_k=\sigma_{k+1}$ for some $k\in \qqq{n}$, then we have
		\begin{align}\label{sum_res_2_NAL}
			 \left\|{\cal U}^{(n)}_{s,t,\bsig} \circ {\cal A}\right\|_\infty \le W^{C_n\e} \p{\frac{\ilambda^2+|1-s|}{\ilambda^2+|1-t|}}^{n-1}  \|{\cal A}\|_{\infty}  + W^{-D+C_n} %\left(\frac{\ell_s^2|1-s|}{\ell_t^2|1-t|}\right)^{n-1} 
			\end{align} 
		for a constant $C_n>0$ that does not depend on $\e$ or $D$.

\item[(II)] If ${\cal A}$ satisfies the following sum-zero property: 
		\begin{align}\label{sumAzero}
			\sum_{a_2,\ldots,a_n\in \Zn}{\cal A}_{\ba}=0, \quad \forall a_1\in \Zn \, ,
		\end{align}
then there exists a constant $C_n>0$ that does not depend on $\e$ or $D$ such that
		\begin{align}\label{sum_res_2}
			\left\|{\cal U}^{(n)}_{s,t,\bsig} \circ {\cal A}\right\|_\infty \le W^{C_n\e} \p{\frac{\ilambda^2+|1-s|}{\ilambda^2+|1-t|}}^{n} \|{\cal A}\|_{\infty} 
			+W^{-D+C_n}. 
		\end{align} 
\end{itemize}
\end{lemma} 
If $1-s \le \ilambda^2/L^{2}$, we obtain the following estimates for the evolution kernels with zero modes removed.
    \begin{lemma}\label{lem:sum_decay_nonzero}
Fix any $1-\ilambda^2/L^{2}\le s \le t<1$ and $\bsig\in \{+,-\}^n$. Let ${\cal A}: (\Zn)^{n}\to \mathbb C$ be an arbitrary $n$-dimensional tensor for a fixed $n\ge 2$. Then, for any subset $A\subset \qqq{n}$ satisfying $A\supset I_{\mathrm{diff}}(\bsig)$, we have  
\begin{align}\label{sum_res_Ndecay_nonzero}
	\| Q^{(A)}\circ {\cal U}^{(n)}_{s,t,\bsig}\;\circ {\cal A}\|_{\infty} \prec \|{\cal A}\|_{\infty} .
\end{align}
\end{lemma}


\subsection{Proofs of supporting lemmas}\label{subsec:support}

This subsection presents the proofs of several supporting lemmas.

\begin{proof}[\bf Proof of Lemma \ref{lem:STOeq_Qt}] %\label{sec:pf_STOeq_Qt_weak}
\Cref{lem:STOeq_NQ} already gives a good enough bound for non-alternating $(\cL-\cK)$-loops. It remains to control \smash{\(\max_{\ba} |({\cal L}-{\cal K})^{(n)}_{u, \bsig, \ba} |\big/ (W^{-d}B_{u,0})^{n}\)} for alternating $\bsig$ using \eqref{eq:Ward_typeP} and equation \eqref{int_K-L+Q}. 

First, the partial sum term \([\cal P\circ (\mathcal{L} - \mathcal{K})^{(n)}_{u, \bsig}]_{a_1} \dthn^{(n)}_{u, \ba}\) has been bounded in \eqref{eq:Ward_typeP}. Second, using  \Cref{lem_+Q}, the induction hypothesis \eqref{Eq:L-KGt+IND} at time $s$, the estimates established in \eqref{bEwGn}--\eqref{eq:MG_nloop}, and the evolution kernel estimate \eqref{sum_res_2}, and adopting an argument similar to that in the proof of \Cref{lem:STOeq_NQ}, we can control the terms involving $\cal B_{i}$ for $i\in \{0,1,2,3\}$ and the martingale term on the RHS of \eqref{int_K-L+Q} as follows:
%(at time $t=u$) by the RHS of \eqref{am;asoi222}. This leads to that 
\begin{align}
&  [\cal P\circ (\mathcal{L} - \mathcal{K})^{(n)}_{u, \bsig}]_{a_1} \dthn^{(n)}_{u, \ba}+ \left(\mathcal{U}^{(n)}_{s, u, \bsig} \circ  {\cal Q}_s\circ  {\cal B}_0(s)\right)_{\ba} +  \int_{s}^u \left(\mathcal{U}^{(n)}_{v, u, \bsig} \circ  {\cal Q}_v\circ \sum_{k=1}^3{\cal B}_k(v)  \right)_{\ba} \dd v + \int_{s}^u \left(\mathcal{U}^{(n)}_{v, u, \bsig} \circ  {\cal Q}_v\circ \dd\mathcal{E}^{M,(n)}_{v, \bsig}\right)_{\ba} \nonumber\\
\prec&~  (W^{-d}B_{u,0})^n \sup_{v\in [s,u]}
  \left((W^{-d}B_{u,0})^{\frac{1}{6}}\wh\Xi^{(\cal L-\cK)}_{v, n}+ \max_{n'=1}^{n-1} \Xi^{({\cal L-\cal K})}_{v,n'}+ \max_{n'=n-1}^{n+1} \Xi^{({\cal L})}_{v,n'} +\max_{n'=\lceil n/2\rceil + 1}^{n-1} \Xi^{({\cal L-\cal K})}_{v,n+2-n'} \p{\Xi^{({\cal L})}_{v,n_1'} \Xi^{({\cal L})}_{v,n_2'}}^{\frac 12 } \right) \nonumber\\
&~+ (W^{-d}B_{u,0})^n \sup_{v\in[s,u]} \left((W^{-d}B_{s,0})^{-\frac{1}{4p}}(\Xi^{(\cal L)}_{v, {2n-1}})^{\frac 1 2}(\Xi^{(\cal L)}_{v, 4p})^{\frac{1}{4p}}\right) .\label{eq:alternatecase1}
\end{align}
It remains to address the terms involving $\cal B_{4}$ and $\cal B_{5}$ on the RHS of \eqref{int_K-L+Q}. We claim the following bounds: %for $\cal B_4$ and $\cal B_5$:
% \begin{claim}[${\cal B}_4$ and ${\cal B}_5$ estimate]
% \label{lem: B45est} \textbf{Yin} \textcolor{blue}{done}
\begin{align}\label{y27kasdfg}
   \| {\cal B}_4(u)\|_{\infty}+\|{\cal B}_5(u)\|_{\infty}\prec \eta_u^{-1} (W^{-d}B_{u,0})^n \Xi^{(\mathcal{L}-\mathcal{K})}_{u, n-1}.
\end{align}
 % \end{claim}
 % \begin{proof}
First, combining \eqref{jywiiwsoks} with \eqref{eq:Bu0asymp} and the second estimate in \eqref{eq:derv_Theta}, we obtain that  
\be\label{A5}  
\left\|\left[ {\cal P} \circ \left(\mathcal{L} - \mathcal{K}\right)^{(n)}_{u, \bsig}\right] \partial_u\dthn_{u}^{(n)}\right\|_\infty \prec \frac{1}{\eta_u}\frac{\ilambda^2+|1-u|}{|1-u|\ell_u^{d}}\cdot (W^{-d}B_{u,0})^n  \Xi^{(\mathcal{L}-\mathcal{K})}_{u, n-1}\lesssim \eta_u^{-1} (W^{-d}B_{u,0})^n \Xi^{(\mathcal{L}-\mathcal{K})}_{u, n-1}.
\ee
Second, using the definition of $\cal Q_t$ in \eqref{eq:sumzero_op} and the definition of $\varTheta^{(n)}_{u,\bsig}$ in \eqref{def:op_thn}, we can bound that 
\begin{align}\label{A4}  
 \br{{\cal Q}_u , \varTheta^{(n)}_{u,\bsig} }\circ (\mathcal{L} - \mathcal{K})^{(n)}_{u, \bsig, \ba}
&= \varTheta^{(n)}_{u, \bsig} \circ \left[ \left({\cal P} \circ \left( \mathcal{L} - \mathcal{K} \right)^{(n)}_{u,\bsig}\right) \cdot \dthn^{(n)}_u \right]_{\ba} - \left[\left( {\cal P} \circ \varTheta^{(n)}_{u, \bsig} \circ (\mathcal{L} - \mathcal{K})^{(n)}_{u,\bsig}\right)  \cdot \dthn^{(n)}_u \right]_{\ba} \nonumber
\\ 
& \prec \eta_u^{-1}\big\|\dthn^{(n)}_{u}\big\|_\infty \cdot   \left\| {\cal P} \circ \left(  \mathcal{L} - \mathcal{K} \right)^{(n)}_{u,\bsig} \right\|_{\infty} \prec \eta_u^{-1} (W^{-d}B_{u,0})^n \Xi^{(\mathcal{L}-\mathcal{K})}_{u, n-1},
%\eta_u^{-1}\p{\ell_u^{d}\eta_u}^{-1}\cdot W^{-dn} \Xi^{(\mathcal{L}-\mathcal{K})}_{u, n-1}\lesssim \eta_u^{-1}\ell_u^{-1}\cdot W^{-dn} \Xi^{(\mathcal{L}-\mathcal{K})}_{u, n-1},
\end{align}
where in the second step, we use the simple fact that $ \|\varTheta_{u,\bsig}^{(n)}\|_{\infty\to\infty} \lesssim (1-u)^{-1}$ for any $\bsig\in \{+,-\}^n$ due to \eqref{eq:THETAinftinf}, and in the third step, we use \eqref{jywiiwsoks} and the first bound in \eqref{eq:derv_Theta}. Combining \eqref{A5} and \eqref{A4} yields \eqref{y27kasdfg}. 
%\end{proof}
To apply the evolution kernel estimate \eqref{sum_res_2}, we also need the fast-decay property \eqref{deccA0} (with $s$ replaced by $v$) for ${\cal Q}_v\circ{\cal B}_4(v)$ and ${\cal Q}_v\circ{\cal B}_5(v)$. We claim that this property holds for any constants $\e,D>0$. Note that $\|{\cal B}_4(v)\|_\infty+\|{\cal B}_5(v)\|_\infty\le W^{C}$ for a constant $C>0$, since $\|G_v\|\le \eta_v^{-1}$, the $\cK$-loops are bounded by \eqref{eq:bcal_k}, $L^d\le W^{C}$, and $(1-v)^{-1}\le L^2/\ilambda^2\le W^{C}$. Hence, by \Cref{lem_+Q}, it suffices to check the fast-decay property for ${\cal B}_4(v)$ and ${\cal B}_5(v)$ themselves. For ${\cal B}_5(v)$, it follows from the decay of $\partial_v\dthn^{(n)}_v$ given by the second estimate in \eqref{eq:derv_Theta}. For ${\cal B}_4(v)$, we use the expression in the first line of \eqref{A4}. Its second term decays by the first estimate in \eqref{eq:derv_Theta}. For its first term, the tensor ${\cal X}:=({\cal P}\circ(\cL-\cK)^{(n)}_{v,\bsig})\cdot\dthn^{(n)}_v$ satisfies the $(v,\e,D)$-decay property for any $\e,D>0$ by \Cref{lem_+Q}, and the kernel of $\varTheta^{(n)}_{v,\bsig}$ in \eqref{def:op_thn} is bounded by $C(1-v)^{-1}e^{-c|a_i-b_i|/\ell_v}$ by \eqref{prop:ThfadC} (together with \eqref{Mbound_AO} or \eqref{Mbound_AO2} in the block Anderson model). Thus, if $\max_{i,j}|a_i-a_j|\ge W^{2\e}\ell_v$, then in the $i$-th term of \eqref{def:op_thn}, either $|a_i-b_i|\ge W^{2\e}\ell_v/2$, in which case the kernel is $\OO(W^{-D})$, or $\max_{j,k}|(\ba^{(i)}(b_i))_j-(\ba^{(i)}(b_i))_k|\ge W^{\e}\ell_v$, in which case ${\cal X}_{\ba^{(i)}(b_i)}=\OO(W^{-D})$. Summing over $b_i\in\Zn$, we see that $\varTheta^{(n)}_{v,\bsig}\circ{\cal X}$ satisfies the $(v,2\e,D-C)$-decay property. Since $\e$ and $D$ are arbitrary, this proves the claim. 
Now, applying the evolution kernel estimate \eqref{sum_res_2} and the estimates in \eqref{y27kasdfg}, and performing the integral over $v$, we can obtain that 
\begin{align}
&  \int_{s}^u \left(\mathcal{U}^{(n)}_{v, u, \bsig} \circ  {\cal Q}_v\circ \sum_{k=4}^5{\cal B}_k(v)  \right)_{\ba} \dd v  \prec  (W^{-d}B_{u,0})^n \sup_{v\in[s,u]}\left(\Xi^{({\cal L-\cal K})}_{v,n-1} \right). \label{eq:alternatecase2}
\end{align}

%Finally, applying the evolution kernel estimate \eqref{sum_res_2} and the estimates \eqref{A5} and \eqref{A4}, and performing the integral over $u$, we can control the last two terms on the RHS of \eqref{int_K-L+Q} by $\p{{\ell_t^{d-1}}/{\ell_s^{d-1}}} \cdot (W^d\ell_t^d\eta_t)^{-n}\sup_{u\in[s,t]}\Xi^{({\cal L-\cal K})}_{u,n-1}$. This concludes the proof of \Cref{lem:STOeq_Qt_weak}.
% \begin{proof}[Proof of lemma \ref{lem: B45est}]We have
% \begin{align}\label{jfasiuu}
% \|{\cal B}_5\|_{\max}  \prec 
%           \max_{a} \left|\left[{\cal P} \circ
%       \left(\mathcal{L} - \mathcal{K}\right)_{u, \bsig}\right]_a\right| \cdot   \max_{{\ba}} \left| 
%           \dot \vartheta_{u,{\ba} } \right|
% \end{align}  
% Recall that $\bsig$ is alternating. In this case, we claim the following holds for any loop of length $n$:
% \begin{align}\nonumber %\label{jywiiwsoks}
%     \left[{\cal P} \circ
%     \left(\mathcal{L} - \mathcal{K}\right)_{u, \bsig}\right]_a \prec (W^d\eta_u)^{-1}
%     \cdot (W^{-d}B_{u,0})^{-n+1}\cdot  \ell_u^{d(n-2)}   \cdot \Xi^{(\mathcal{L}-\mathcal{K})}_{u, n-1}.
%  \end{align}
% To see this, since $\bsig$ is non-constant, there exists $1 < k \leq n$ such that $\sigma_k = \overline{\sigma}_{k+1}$. By Lemma \ref{lem_WI_K}, we can sum $\mathcal{L} - \mathcal{K}$ of rank $n$ over $a_k$ and write the resulting sum in terms of $\mathcal{L} - \mathcal{K}$ of rank $n-1$ with a multiplicative factor $(W^{d} \eta_u)^{-1}$. The summation over the remaining $n-2$ many indices results in a factor of $\prec\ell_{u}^{d(n-2)}$ because ${\cal L}-{\cal K}$ has the $(u,\tau,D)$ decay property for any $\tau,D>0$ by assumption \eqref{Eq:Gdecay_w}. Finally, the additional $(W^{d}B_{u,0})^{-n+1} $ comes from compensating the factor   in the definition of $\Xi^{({\cal L}-{\cal K})}_{u,n-1}$. Next, 
%     \begin{align}
%    \max_{{\ba}}\left| 
%        \dot \vartheta_{u,{\ba} } \right|\prec \ell_u^{-d(n-1)}\cdot \eta_u^{-1}  \label{eq:thetadot_bound}
%    \end{align}
% by definition of $\vartheta$ and the estimate \eqref{prop:ThfadC} for $\Theta^{(B)}$. Together with 
% $$
% \frac{1}{(\ell_u) ^d \cdot \eta_u}\le \frac{1}{L^d\cdot \eta_u}\le B_{u,0}
% $$
% Inserting these two bounds into \eqref{jfasiuu} and using \eqref{STOeq_Qt_assume_0} gives
% \begin{align}\label{kkuuwsaf}
%   \|{\cal B}_5\|_{\max}  \prec \frac{1}{\eta_u} (W^{-d}B_{u,0})^{n-1} \Xi^{(\mathcal{L}-\mathcal{K})}_{u, n-1}
%  \end{align}
% By a similar argument, for ${\cal B}_4$, we have
%  \begin{align}\label{kkuuwsaf5}
%   \|{\cal B}_4\|_{\max}
%    \;=\; & \; \left[ \left({\cal P} \circ \left(\varTheta_{u, \bsig} \circ (\mathcal{L} - \mathcal{K})\right)\right) \cdot \vartheta_u \right]_{{\ba}} 
%     - \varTheta_{u, \bsig} \circ \left|\left[ {\cal P} \circ \left( \mathcal{L} - \mathcal{K} \right) \cdot \vartheta_t \right]_{{\ba}}\right|.
% \\\nonumber
%  \;\prec \; & \;  \|\vartheta_u\|_{\max}\cdot \eta_u^{-1}\max_a\left[ {\cal P} \circ \left(  \mathcal{L} - \mathcal{K} \right)_{u,\bsig} \right]_{a} 
%   \\\nonumber
%  \;\prec \; & \;(\ell_u)^{-d\cdot (n-1)} \cdot (W^d\eta_u)^{-1}\cdot  (W^{-d}B_{u,0})^{-n+1}\cdot (\ell_u)^{ d\cdot (n-2)}     \cdot \Xi^{(\mathcal{L}-\mathcal{K})}_{u, n-1}
% \end{align} 
% which implies \eqref{y27kasdfg}. 
% \end{proof}


Finally, combining \Cref{lem:STOeq_NQ} with \eqref{eq:alternatecase1} and \eqref{eq:alternatecase2}, we obtain that 
\begin{align*}
\sup_{v\in[s,u]} \wh\Xi^{({\cal L-\cal K})}_{v,n} &\prec (W^{-d}B_{u,0})^{\frac{1}{6}}\sup_{v\in [s,u]}
   \wh\Xi^{(\cal L-\cK)}_{v, n}  + \sup_{v\in[s,u]} \left((W^{-d}B_{s,0})^{-\frac{1}{4p}}(\Xi^{(\cal L)}_{v, {2n-1}})^{\frac 1 2}(\Xi^{(\cal L)}_{v, 4p})^{\frac{1}{4p}}\right)\\
&+\sup_{v\in [s,u]}
  \left(\max_{n'=1}^{n-1} \Xi^{({\cal L-\cal K})}_{v,n'}+ \max_{n'=n-1}^{n+1} \Xi^{({\cal L})}_{v,n'} + \max_{n'=\lceil n/2\rceil + 1}^{n-1} \Xi^{({\cal L-\cal K})}_{v,n+2-n'} \p{\Xi^{({\cal L})}_{v,n_1'} \Xi^{({\cal L})}_{v,n_2'}}^{\frac 12 } \right),\nonumber
\end{align*}
which yields \eqref{am;asoi222} at each fixed $u\in[s,t]$ upon solving for \smash{$\sup_{v\in[s,u]} \wh\Xi^{({\cal L-\cal K})}_{v,n}$}. Then, applying a standard $N^{-C}$-net argument, together with the continuity estimate of \Cref{lem_ConArg} which controls the increments of the $G$-loops between two nearby times, extends it uniformly to all $u\in [s,t]$. 
%which gives \eqref{am;asoi222} with $u=t$ by solving this equation for $\sup_{u\in[s,t]} \wh\Xi^{({\cal L-\cal K})}_{u,n}$. 
\end{proof}

%\subsection{Proof of Lemma \ref{lem:iterations}}\label{sec:pf_iterations}
%In the setting of \Cref{lem:main_ind}, 


% \textcolor{blue}{You have already put the proof of the above lemma into appendix, why not state it as under the theorem 2.24 suppose $1-t\ge g^2/L^2$, then we have \eqref{am;asoi333}? The current version is quite twisted, especailly you have \eqref{am;asoi222} and \eqref{am;asoi333} }We are now ready to complete the proof of Step 3 in the case $1-t\ge \ilambda^2/L^{2}$. In this step, besides the assumptions in \Cref{lem:main_ind}, we have already established the a priori $G$-loop bound \eqref{lRB1} in Step 1, as well as the sharp local laws \eqref{Gt_bound_flow} and \eqref{Gt_avgbound_flow}, along with the weak $2$-$G$-loop decay estimate \eqref{Eq:Gdecay_w} in Step 2. In particular, 
\begin{proof}[\bf Proof of Lemma \ref{lem:iterations}]
By the averaged local law \eqref{Gt_avgbound_flow}, we have \smash{$\wh\Xi_{u,1}^{(\cL-\cK)}\prec 1=:\Xi_{u,1}^{(\cL-\cK)}$} uniformly for all $u\in[s,t]$. Substituting this into \eqref{am;asoi222}, we obtain the following bound, also uniformly in $u\in [s,t]$ under the assumption $1-s\le \ilambda^2$: 
\begin{align}\label{am;asoi333}
\sup_{v\in[s,u]} \wh \Xi^{({\cal L-\cal K})}_{v,n}&\prec \sup_{v\in[s,u]} \left(\p{W^{-d}B_{s,0}}^{-\frac{1}{4p}}(\Xi^{(\cal L)}_{v, {2n-1}})^{\frac 1 2}(\Xi^{(\cal L)}_{v, 4p})^{\frac{1}{4p}}\right) \\
&+\sup_{v\in [s,u]}
\left( \max_{n'=2}^{n-1}\Xi^{({\cal L-\cal K})}_{v,n'}+\max_{n'=n-1}^{n+1}\Xi^{({\cal L})}_{v,n'}+\max_{n'=3}^{n-1} \Xi^{({\cal L-\cal K})}_{v,n+2-n'} \p{\Xi^{({\cal L})}_{v,n_1'} \Xi^{({\cal L})}_{v,n_2'}}^{\frac 12 } \right),\nonumber
\end{align}
%for any fixed $p\ge 1$. 
where we also use that $\lceil n/2\rceil + 1\le n-1$ only when $n\ge 4$, in which case we have $\lceil n/2\rceil + 1\ge 3$. 

To show \Cref{lem:iterations} using the bootstrap bound \eqref{am;asoi333}, we first claim that under the assumption \eqref{eq:iteration_induc}, the following bound holds uniformly in $u\in[s,t]$ %for every fixed $k\ge 1$ 
as long as we choose $p\ge 4$: 
 \begin{align} \label{xiu2n+2psi}
\sup_{v\in[s,u]}\left((\ilambda^2W^d)^{\frac{1}{4p}}(\wh\Xi^{(\cal L)}_{v, {2n-1}})^{\frac 1 2}(\wh\Xi^{(\cal L)}_{v, 4p})^{\frac{1}{4p}}\right)  \prec \Psi(n,k;s,u) \ .
\end{align} 
Given the bound \eqref{xiu2n+2psi}, we obtain from \eqref{am;asoi333} that
\begin{equation} 
\sup_{v\in[s,u]}\wh\Xi_{v,n }^{(\mathcal{L}-\mathcal{K})} \prec \Psi(n,k;s,u) + 
\sup_{v\in [s,u]}
\left( \max_{n'=2}^{n-1}\Xi^{({\cal L-\cal K})}_{v,n'}+\max_{n'=n-1}^{n+1}\Xi^{({\cal L})}_{v,n'}+\max_{n'=3}^{n-1} \Xi^{({\cal L-\cal K})}_{v,n+2-n'} \p{\Xi^{({\cal L})}_{v,n_1'} \Xi^{({\cal L})}_{v,n_2'}}^{\frac 12 } \right)\label{sadui_w0}
 \end{equation}
uniformly in $u\in[s,t]$. By the induction hypothesis \eqref{eq:iteration_induc}, we may choose the $\Xi^{({\cal L-\cal K})}$-parameters as 
\be \nonumber
\Xi^{({\cal L-\cal K})}_{v,n'}= \Psi(n',k;s,u)\le \Psi(n,k;s,u) \ \  \text{for}\ \ 2\le n' \le n-1, %\ \  \text{and} \ \ \Xi^{({\cal L-\cal K})}_{v,n}=\Psi(n,k-1;s,u),
\ee
%for each $v\in[s,u]$. 
and, by \eqref{rela_XILXILK} and \eqref{eq:iteration_induc}, we can choose the $\Xi^{({\cal L})}$-parameters as
\begin{align*}
   \Xi^{({\cal L})}_{v,n'}=\begin{cases}
       1+(\ilambda^2W^{d})^{-1} \Psi(n',k;s,u), \ & \text{for} \ \  n'\le n-1,\\  
   1+(\ilambda^2W^{d})^{-1} \Psi(n',k-1;s,u)\ & \text{for} \ \ n\le n'\le n+2. 
   \end{cases} 
\end{align*}
%$$\sup_{v\in [s,u]}\wh \Xi^{({\cal L})}_{v,n'}\prec 1+(\ilambda^2W^{d})^{-1} \Psi(n',k-1;s,u):=\Xi^{({\cal L})}_{v,n'}.$$
With these choices, we readily verify that  
\[\max_{n'=n-1}^{n+1}\Xi^{({\cal L})}_{v,n'} \le  \Psi(n,k;s,u),\quad \text{and}\quad  \p{\Xi^{({\cal L})}_{v,n_1'} \Xi^{({\cal L})}_{v,n_2'}}^{\frac 12 } \lesssim 1+\p{{\eta_s}/{\eta_u}}^{n'}(\ilambda^2W^d)^{-\frac{k}{8}} \ \ \text{for} \ \ 3\le n'\le {n-1}.\]
Substituting these bounds into \eqref{sadui_w0}, we obtain  %and using that $\Psi(n',k;s,u)$ is increasing in $n'$, we obtain  
 \begin{align*}
\sup_{v\in[s,u]}\wh\Xi_{v,n }^{(\mathcal{L}-\mathcal{K})}\prec \Psi(n,k;s,u) +  \max_{n'=3}^{n-1} \Psi(n+2-n',k;s,u)\bigg[1+\p{\frac{\eta_s}{\eta_u}}^{n'}(\ilambda^2W^d)^{-\frac{k}{8}}\bigg]\lesssim \Psi(n,k;s,u), 
 \end{align*}
%where we again used the condition \eqref{con_st_ind} in the second step. 
%Inserting the above two estimates into \eqref{sadui_w0} 
which concludes the proof of \Cref{lem:iterations}.


It remains to prove \eqref{xiu2n+2psi}. For $n\in\{2,3\}$, using the a priori bound \eqref{sef8w483r324} for \smash{$\wh\Xi^{(\cal L)}_{v, {2n-1}}$ and $\wh\Xi^{(\cal L)}_{v, 4p}$}, along with the relation \eqref{rela_XILXILK}, we obtain for $p\ge 4$ that,
% we have 
% \[\sup_{v\in[s,u]} \wh\Xi^{(\cal L)}_{v, {2n-1}} \prec \Psi(2n-1,k-1;s,u).\]
% Then, using the relation \eqref{rela_XILXILK} and the a priori bound \eqref{sef8w483r324} for \smash{$\wh\Xi^{(\cal L)}_{v, 4p}$}, 
\begin{align*} 
\sup_{v\in[s,u]}\left((\ilambda^2W^d)^{\frac{1}{4p}}(\wh\Xi^{(\cal L)}_{v, {2n-1}})^{\frac 1 2}(\wh\Xi^{(\cal L)}_{v, 4p})^{\frac{1}{4p}}\right)  \prec (\ilambda^2W^d)^{\frac{1}{16}}\left( {\eta_s}/{\eta_u}\right)^{n}  \le \Psi(n,k;s,u).
\end{align*} 
We now consider the $n\ge 4$ case. Recall the $G$-chain defined in \eqref{defC=GEG}. Given an arbitrary $(2n-1)$-$G$-loop \smash{$\cL^{(2n-1)}_{v,\bsig,\ba}$} with 
$$\bsig=(\sig,\sig_1,\ldots, \sig_{n-1}, \sig',\sig_1',\ldots, \sig_{n-2}'),\quad \ba=(b,a_1,\ldots,a_{n-1},b',a_1',\ldots,a_{n-2}'),$$ 
we can write it as 
\begin{align*}
\cL^{(2n-1)}_{v,\bsig,\ba}=W^{-2d}\sum_{x\in[b],x'\in [b']} \left(\cal C^{(n)}_{v,\bsig_1,\ba_1}\right)_{xx'} \left(\cal C^{(n-1)}_{v,\bsig_2,\ba_2}\right)_{x'x} ,  
\end{align*}
where $\bsig_1=(\sig_1,\ldots, \sig_{n-1}, \sig')$, $\bsig_2=(\sig_1',\ldots, \sig_{n-2}',\sig)$, $\ba_1=([a_1],\ldots,[a_{n-1}])$, and $\ba_2=([a_1'],\ldots,[a_{n-2}'])$. Applying the Cauchy-Schwarz inequality gives
\begin{equation}\label{suauwiioo1}
\left|\cL^{(2n-1)}_{v,\bsig,\ba}\right|
\le 
\bigg( W^{-2d}\sum_{x\in[b],x'\in[b']}
\left|\left({\cal C}^{(n)} _{v,\bsig_1,\ba_1}\right)_{xx'} 
\right|^2 \bigg)^{1/2}\bigg( W^{-2d}\sum_{x\in[b],x'\in[b']}
\left|\left({\cal C}^{(n-1)} _{v,\bsig_2,\ba_2}\right)_{x'x} 
\right|^2 \bigg)^{1/2}.
\end{equation}
Following the argument for equation (5.118) in \cite{YY_25}, we obtain that 
\begin{equation}\label{eq:boundtwochains}
    \begin{split}
& W^{-2d}\sum_{x\in[b],x'\in[b']}
\left|\left({\cal C}^{(n)} _{v,\bsig_1,\ba_1}\right)_{xx'} 
\right|^2 \le (\ilambda^2W^d) \Xi^{(\cL)}_{v,2l_1}\Xi^{(\cL)}_{v,2l_2}, \\  &W^{-2d}\sum_{x\in[b],x'\in[b']}
\left|\left({\cal C}^{(n-1)} _{v,\bsig_2,\ba_2}\right)_{x'x} 
\right|^2 \le (\ilambda^2W^d) \Xi^{(\cL)}_{v,2l_3}\Xi^{(\cL)}_{v,2l_4},
\end{split}
\end{equation}
where $l_1=\lceil n/2\rceil$, $l_2=\lfloor n/2\rfloor$, $l_3=\lceil (n-1)/2\rceil$, and $l_4=\lfloor (n-1)/2\rfloor$. 
By the induction hypothesis \eqref{eq:iteration_induc} (noting $\max_i\p{2l_i}\le n+1$) and the relation \eqref{rela_XILXILK}, we have
\begin{align}\label{auskoppw2.00}
\wh\Xi_{v,2l_i}^{(\cal L)} \prec 1+ (\ilambda^2 W^{d})^{-1} \Psi(2l_i,k-1;s,u) \prec 1+ \p{{\eta_s}/{\eta_u}}^{2l_i-1} (\ilambda^2W^d)^{-(k-1)/8}.
\end{align}
Plugging \eqref{auskoppw2.00} into \eqref{eq:boundtwochains} and then applying \eqref{suauwiioo1}, we obtain for $k\ge 1$ that, 
\begin{align}\label{auskoppw2}
\sup_{v\in[s,u]}\big(\wh\Xi^{({\cal L})}_{v, 2n-1}\big)^{\frac{1}{2}}&\prec  \sup_{v\in[s,u]} (\ilambda^2W^d)^{1/2}\left[\prod_{i=1}^4\bigg(1+ \p{{\eta_s}/{\eta_u}}^{2l_i-1} (\ilambda^2W^d)^{-(k-1)/8}\bigg)\right]^{1/4} \nonumber\\
&\lesssim (\ilambda^2W^d)^{1/2}\bigg(1+ \p{{\eta_s}/{\eta_u}}^{n} (\ilambda^2W^d)^{-(k-1)/8}\bigg).
\end{align}
On the other hand, by \eqref{sef8w483r324}, we have 
\(\sup_{v\in[s,u]} (\wh\Xi^{(\cal L)}_{v, 4p})^{\frac{1}{4p}} \prec {\eta_s}/{\eta_u} .\)
Combining it with \eqref{auskoppw2}, and using the condition \eqref{con_st_ind}, we conclude \eqref{xiu2n+2psi} when $p\ge 4$.  
\end{proof}  



\begin{proof}[\bf Proof of Lemma \ref{lem: newPQ}] %\label{Sec:pf_newPQ}
We only prove the expansion \eqref{yurenAL} using Ward's identity \eqref{WI_calL}, while \eqref{yurenAK} follows from the same argument by using the identity \eqref{WI_calK} instead. 
We first record the following consequence of \eqref{WI_calL}: for any $\bsig\in \{+,-\}^n$ with $\sig_i\ne \sig_{i+1}$ and any subset $A\subset \qqq{n}\setminus\{i\}$, %there is 
\begin{align}\label{y2ussz}
\left(Q^{(A)} \circ P^{(i)}\circ {\cal L}^{(n)}\right)_{t,\bsig,\ba}
= \frac{1}{ 2\ii N\eta_t}
\left(\left(Q^{(A_{(i)})} \circ {\cal L}^{(n-1)}\right)_{t,\bsig_+,\ba_+}- \left(Q^{(A_{(i)})}\circ {\cal L}^{(n-1)}\right)_{t,\bsig_-,\ba_-}\right), 
\end{align}
where $A_{(i)}:=\{1\le j<i: j\in A\} \cup \{i\le j\le n-1: j+1\in A\}$, $\ba_\pm = (a_1,\ldots, a_{i-1}, a_{i+1},\ldots, a_n)$, and $\bsig_\pm = (\sigma_1\cdots \sigma_{i-1}, \pm, \sigma_{i+2},\cdots \sigma_n)$. 
We now prove \eqref{yurenAL} by induction in $n$ and using \eqref{y2ussz}. 

First, it is easy to see that \eqref{yurenAL} trivially holds for $n=1$ and $|A|\in\{0,1\}$. 
Next, suppose we have shown that there exists a decomposition \eqref{yurenAL} for any $n\le l-1$. We prove \eqref{yurenAL} for $n=l$. Given any $A\subset\qqq{n}$, if $A\supset I_{\mathrm{diff}}(\bsig)$, then \eqref{yurenAL} holds trivially with only one term on the RHS (i.e., the term $Q^{(A)}\circ {\cal L}^{(n)}$ itself). Otherwise, suppose there exists $i\in I_{\mathrm{diff}}(\bsig)\setminus A$, then we write
\be\label{eq:y2ussz}
 \left(Q^{(A)}\circ {\cal L}^{(n)}\right)_{t,\bsig,\ba}= \left(Q^{(A\cup \{i\} )}\circ {\cal L}^{(n)}\right)_{t,\ba,\bsig}+\left(Q^{(A )}\circ P^{(i)}\circ{\cal L}^{(n)}\right)_{t,\bsig,\ba}. 
 \ee
Applying \eqref{y2ussz} to the second term on the RHS and using the induction hypothesis for $Q^{(A_{(i)})} \circ {\cal L}^{(n-1)}$, we can expand \smash{$(Q^{(A )}\circ P^{(i)}\circ{\cal L}^{(n)})_{t,\bsig,\ba}$} into an expression that matches the form of a summand on the RHS of \eqref{yurenAL}. 
%into an expression taking the form of a summand on the RHS of \eqref{yurenAL}. 
On the other hand, in the first term of \eqref{eq:y2ussz}, the number of elements in $I_{\mathrm{diff}}(\bsig)\setminus (A\cup\{i\})$ is reduced by 1 relative to $I_{\mathrm{diff}}(\bsig)\setminus A$. If $I_{\mathrm{diff}}(\bsig)\setminus (A\cup\{i\})=\emptyset$, then the induction is completed. Otherwise, we choose the next element $j\in I_{\mathrm{diff}}(\bsig)\setminus (A\cup\{i\})$ and repeat the same argument with $A$ and $i$ replaced by $A\cup\{i\}$ and $j$, respectively. Iterating this procedure at most $|I_{\mathrm{diff}}(\bsig)\setminus A|$ times, we finally express $Q^{(A)}\circ {\cal L}^{(n)}$ in the form \eqref{yurenAL}. This completes the induction, thereby concluding the proof of \Cref{lem: newPQ}. 
\end{proof}

\begin{proof}[\bf Proof of Lemma \ref{lem:STOeq_Qt_nonzero}]
%\subsection{Proof of Lemma \ref{lem:STOeq_Qt_nonzero}}\label{sec:STOeq_Qt_nonzero}
Given any $\bsig\in\{+,-\}^n$ and $\ba\in (\Zn)^n$, we expand $\cL^{(n)}_{t,\bsig,\ba}$ and $\cK^{(n)}_{t,\bsig,\ba}$ as in \eqref{yurenAL} and \eqref{yurenAK} with $A_n=I_{\mathrm{diff}}(\bsig)$. This yields 
\begin{align}
\p{\cal L-\cal K}^{(n)}_{t,\bsig,\ba}=Q^{(A_n)}\circ \p{\cal L-\cal K}^{(n)}_{t,\bsig,\ba}+\sum_\al \frac{\xi_{\al}}{(2\ii N\eta_t)^{n-k_\al}}   Q^{(A_\al)}\circ \p{\cal L-\cK}^{(k_\al)} _{t,\bsig_{\al},\ba_{\al}}.\label{eq:expandQAempty}
\end{align}
For the $k_\al=1$ terms, using \eqref{normQA2} and the averaged local law \eqref{Gt_avgbound_flow}, we bound 
\begin{align}\label{eq:kalpha1}
    \frac{\xi_{\al}}{(2\ii N\eta_t)^{n-k_\al}}   Q^{(A_\al)}\circ (\cal L-\cK)^{(k_\al)} _{t,\bsig_{\al},\ba_{\al}} \prec (N\eta_t)^{-(n-1)}\cdot W^{-d}B_{t,0} \lesssim (W^{-d}B_{t,0})^n.
\end{align}

All other terms on the RHS of \eqref{eq:expandQAempty} involve loops of length $\ge 2$, and therefore satisfy equation \eqref{iisuwjyys}.
%Every other term on the RHS (with $(\cL-\cK)$-loops of length $\ge 2$) satisfies equation \eqref{iisuwjyys}. 
Take the first term as an example. Using the evolution kernel estimate in \Cref{lem:sum_decay_nonzero}, the bounds \eqref{bEwGn}--\eqref{eq:MG_nloop}, and the assumption \eqref{Eq:L-KGt+IND} on \smash{$(\mathcal{L} - \mathcal{K})^{(n)}_{s, \bsig, \ba}$}, we obtain that
\begin{align}\label{sahwNQ_smalleta}
&\frac{ Q^{(A_n)}\circ (\mathcal{L} - \mathcal{K})^{(n)}_{t, \bsig, \ba}}{(W^{-d}B_{t,0})^{n}} \prec 
1 + (W^{-d}B_{t,0})^{\frac{1}{6}} \int_{s}^t   \wh\Xi^{({\cal L-\cal K})}_{u,n} \frac{\dd u}{\eta_u} +   \int_{s}^t  \max_{n'=2}^{n-1} \Xi^{({\cal L-\cal K})}_{u,n'}\frac{\dd u}{\eta_u}  + \int_{s}^t \max_{n'=n-1}^{n+1} \Xi^{(\cal L)}_{u,n'}\frac{\dd u}{\eta_u}    \nonumber \\
&+  \int_{s}^t \max_{n'=\lceil n/2\rceil + 1}^{n-1} \Xi^{({\cal L-\cal K})}_{u,n+2-n'} \p{\Xi^{({\cal L})}_{u,n_1'} \Xi^{(\cal L)}_{u,n_2'}}^{\frac 12 } \frac{\dd u}{\eta_u}   + \frac{1}{(W^{-d}B_{t,0})^{n}} \int_{s}^t \left(Q^{(A_n)}\circ \mathcal{U}^{(n)}_{u, t, \bsig} \circ  \dd  \mathcal{E}^{M,(n)}_{u, \bsig}\right)_{\ba}.
\end{align}
By \Cref{lem:DIfREP}, together with \eqref{eq:MG_nloop} and \Cref{lem:sum_decay_nonzero}, we bound the martingale term as  
\begin{align} \label{sahwNQ2}
(W^{-d}B_{t,0})^{-n} \int_{s}^t \left(Q^{(A_n)}\circ \mathcal{U}^{(n)}_{u, t, \bsig} \circ  \dd  \mathcal{E}^{M,(n)}_{u, \bsig}\right)_{\ba}&\prec \p{W^{-d}B_{s,0}}^{-\frac{1}{4p}}\left\{\int_{s}^t 
   \eta_u^{-1}\Xi^{(\cal L)}_{u, {2n-1}}\p{\Xi^{(\cal L)}_{u, 4p}}^{\frac{1}{2p}}\dd u\right\}^{1/2} 
  \end{align} 
for any fixed $p\ge 1$. Plugging it into \eqref{sahwNQ_smalleta} and performing the integral over $u$, we obtain 
\begin{align}\label{am;asoiuw_smalleta}
\frac{ Q^{(A_n)}\circ (\mathcal{L} - \mathcal{K})^{(n)}_{t, \bsig, \ba}}{(W^{-d}B_{t,0})^{n}} &\prec \sup_{u\in [s,t]}
  \left((W^{-d}B_{t,0})^{\frac{1}{6}}\wh\Xi^{(\cal L-\cK)}_{u, n}+ \max_{n'=2}^{n-1} \Xi^{({\cal L-\cal K})}_{u,n'}+ \max_{n'=n-1}^{n+1} \Xi^{({\cal L})}_{u,n'} \right)\\
&+\sup_{u\in [s,t]}
  \left(\max_{n'=\lceil n/2\rceil + 1}^{n-1} \Xi^{({\cal L-\cal K})}_{u,n+2-n'} \p{\Xi^{({\cal L})}_{u,n_1'} \Xi^{({\cal L})}_{u,n_2'}}^{\frac 12 } +\p{W^{-d}B_{s,0}}^{-\frac{1}{4p}}(\Xi^{(\cal L)}_{u, {2n-1}})^{\frac 1 2}(\Xi^{(\cal L)}_{u, 4p})^{\frac{1}{4p}}\right).\nonumber
\end{align} 
This bound also applies to each term \smash{$Q^{(A_\al)}\circ \p{\cal L-\cK}^{(k_\al)} _{t,\bsig_{\al},\ba_{\al}}$} in \eqref{eq:expandQAempty}, with $n$ replaced by $k_\al$. Using the fact \smash{$(N\eta_t)^{-1}\le W^{-d}B_{t,0}$}, every summand of \eqref{eq:expandQAempty}, after being rescaled by
%it shows that all terms on the RHS of \eqref{eq:expandQAempty}, after rescaled by 
$(W^{-d}B_{t,0})^{-n}$, is controlled by the RHS of \eqref{am;asoiuw_smalleta}.
%can be controlled by the RHS of \eqref{am;asoiuw_smalleta}. 
Now, taking the maximum of both sides of \eqref{eq:expandQAempty} over $(\bsig,\ba)$ gives
\begin{align*}
\sup_{u\in[s,t]} \wh\Xi^{({\cal L-\cal K})}_{u,n} &\prec (W^{-d}B_{t,0})^{\frac{1}{6}}\sup_{u\in [s,t]}
 \wh\Xi^{(\cal L-\cK)}_{u, n}  + \sup_{u\in[s,t]} \left(\p{W^{-d}B_{s,0}}^{-\frac{1}{4p}}(\Xi^{(\cal L)}_{u, {2n-1}})^{\frac 1 2}(\Xi^{(\cal L)}_{u, 4p})^{\frac{1}{4p}}\right)\\
&+\sup_{u\in [s,t]}
  \left(\max_{n'=2}^{n-1} \Xi^{({\cal L-\cal K})}_{u,n'}+ \max_{n'=n-1}^{n+1} \Xi^{({\cal L})}_{u,n'} + \max_{n'=\lceil n/2\rceil + 1}^{n-1} \Xi^{({\cal L-\cal K})}_{u,n+2-n'} \p{\Xi^{({\cal L})}_{u,n_1'} \Xi^{({\cal L})}_{u,n_2'}}^{\frac 12 } \right),\nonumber
\end{align*}
solving which yields \eqref{am;asoi222} at $u=t$. Obviously, the same argument applies to each fixed $u\in [s,t]$. Finally, a standard $N^{-C}$-net argument, together with the continuity estimate of \Cref{lem_ConArg}, extends \eqref{am;asoi222} uniformly to the entire interval $u\in [s,t]$. 
\end{proof}

   
\section{Step 5: Pointwise estimate for \texorpdfstring{$(\cL-\cK)$-loops}{L-Kloop}}\label{Sec:Step5}

 

As in the proof for Step 3, by adding intermediate times if necessary, we can divide the proof into the following four cases: (i) $\ilambda^2/L^{2}\le 1- t\le 1-s \le \ilambda^2$, (ii) $\ilambda^2/L^{d}\le 1-t\le 1-s\le \ilambda^2/L^{2}$, (iii) $1- s\ge 1-t \ge \ilambda^2$, and (iv) $1-t\le 1-s\le \ilambda^2/L^{d}$. 
In case (iv), $\ell_t=L$ and $B_{t,0}$ is dominated by the $(L^d|1-t|)^{-1}$ term, so the desired bound already follows from the conclusion \eqref{Eq:L-KGt-flow} in Step 4. Hence, we only need to focus on the first three cases in the following proof. Our goal is to remove the $(|1-s|/|1-u|)^{C_d}$ factor in \eqref{Eq:Gdecay_w}. 
Our proofs below are all based on the integrated loop hierarchy in equation \eqref{int_K-LcalE} with $n=2$:
\begin{align} \label{int_K-LcalE_n=2}
    (\mathcal{L} - \mathcal{K})^{(2)}_{t, \bsig, \ba}= &
    \left(\mathcal{U}^{(2)}_{s, t, \bsig} \circ (\mathcal{L} - \mathcal{K})^{(2)}_{s, \bsig}\right)_{\ba}+ \int_{s}^t \left(\mathcal{U}^{(2)}_{u, t, \bsig} \circ \mathcal{E}^{(\mathcal{L} - \mathcal{K}) \times (\mathcal{L} - \mathcal{K}),(2)}_{u, \bsig}\right)_{\ba} \dd u  \nonumber \\
    &+ \int_{s}^t \left(\mathcal{U}^{(2)}_{u, t, \bsig} \circ \mathcal{E}^{\Gc,(2)}_{u, \bsig}\right)_{\ba} \dd u + \int_{s}^t \left(\mathcal{U}^{(2)}_{u, t, \bsig} \circ \dd\mathcal{E}^{M,(2)}_{u, \bsig}\right)_{\ba} . 
\end{align}
With the stability estimate \eqref{Eq:Gdecay_w} in hand, we can now show that the quadratic term $\mathcal{E}^{(\mathcal{L} - \mathcal{K}) \times (\mathcal{L} - \mathcal{K}),(2)}$ yields a negligible contribution. Using the light-weight estimate \eqref{eq:LW_conclusion_exp} and the martingale estimate \eqref{eq:MG_conclusion3} established in Step 2, we can also obtain sufficient control of the third and fourth terms on the RHS of \eqref{int_K-LcalE_n=2}.
The remaining technical challenge arises from the contribution of the initial condition, i.e., the first term on the RHS of \eqref{int_K-LcalE_n=2}.
In case (i), bounding this term requires leveraging a CLT-type cancellation mechanism, similar to that developed in \cite{DYYY25}.\footnote{Interestingly, in two dimensions \cite{DYYY25}, the CLT cancellation was used in Step~3 but not in Step~5, whereas in our case the situation is reversed. In lower dimensions, Step~5 becomes almost trivial once Steps~2 and~4 are established; however, in dimensions $d\ge 3$, the additional complexity in Step 5 arises once again due to the polynomial decay of the 2-$G$-loops.}
In case (ii), we apply the zero-mode-removing operator introduced in \Cref{def;zero_mode_remove}. Finally, case (iii) can be regarded as a special instance of the lower-dimensional proofs for $d\in \{1,2\}$, as explained in \Cref{sec:Step5_larget}. %---this is why it is presented after cases (i) and (ii).


If $|1-s|/|1-t|\le (\log W)^{10}$, then the factor $(|1-s|/|1-u|)^{C_d}\le (\log W)^{10C_d}$ can be absorbed into the stochastic domination notation ``$\prec$''. Hence, throughout the following proof, we always assume that 
\be\label{eq:assmtlarge} |1-s|/|1-t| > (\log W)^{10} \implies t\ge 1-(\log W)^{-10}.\ee 


\subsection{The case \texorpdfstring{$\ilambda^2/L^{2}\le 1- t\le 1-s \le \ilambda^2$}{1-t>L-2}}
In this case, we have \(B_{u,0}\asymp \ilambda^{-2}\) and $\ell_u\asymp \ilambda|1-u|^{-1/2}$ for all \(u\in[s,t].\)  
In Step 2, we have established the following estimates %for all $u\in[s,t]$ and 
for any large constant $D>0$:  
\be\label{eq:Step2_inputs}
 {\cal L}^{(2)}_{u,\bsig,(a_1,a_2)}\prec   W^{-d}\widetilde{\mathcal{T}}_{u,D}^{L}\p{|a_1-a_2|},\ \ ({\cal L}-{\cal K})^{(2)}_{u,\bsig,(a_1,a_2)}\prec  (\ilambda^2W^{d})^{-\frac16}\cdot   W^{-d}\widetilde{\mathcal{T}}_{u,D}^{L}(|a_1-a_2|),\ \ \forall u\in[s,t].
\ee
 % and 
 % $$
 % ({\cal L}-{\cal K})_{u, {\ba}}\prec  W^{-d/6}\cdot   \widetilde{\mathcal{T}}_{u,|a_1-a_2|}^{dL, D},\quad  {\cal J}_{u}^{dL,D}\le  W^{-d/6}
 % $$
We now apply \eqref{eq:Step2_inputs} to \eqref{i2kk2zgg} with $\ell=L$. In this case, the index sets $\{c':|c'-b|\ge L-1\}$ and $\{c':|c'-a|\ge L-1\}$ appearing in the first two terms on the RHS are empty, because $\max_{a,b\in \Zn}|a-b|\le L/2$ for the periodic $L^\infty$-distance on $\Zn$. Hence only the third term remains, and it can be bounded as in \eqref{eq:i2kk2zgglast}. Consequently, we obtain that
%, and notice that the first two terms on the RHS of \eqref{i2kk2zgg} are not necessary, \textcolor{blue}{fix the argument here, what do you mean "not necessary"} while the last term on the RHS can be bounded as in \eqref{eq:i2kk2zgglast}. 
\be \label{S5WG+M000}
\mathcal{E}_{u, \bsig, \ba}^{(\mathcal{L} - \mathcal{K}) \times (\mathcal{L} - \mathcal{K}),(2)} 
\prec   \frac{(\ilambda^2W^{d})^{-\frac13}}{\eta_u}  \br{W^{-d}\widetilde{{\cal T}}^{L} _{u,D}(|a_1-a_2|) }.
\ee
Moreover, using \Cref{lem: EWGn2_N,lem: EMn2_N}, we get that for any $\bsig\in \{+,-\}^2$,
\be \begin{aligned}\label{S5WG+M}
{\cal E}^{\Gc,(2)}_{u,\bsig, \ba}&\prec \frac{ (\ilambda^2W^{d})^{-\frac12}}{\eta_u}   \br{W^{-d} \widetilde{  \cal T} ^{L}_{u, D}(|a_1-a_2|)},\\  
% \mathcal{E}_{u, \bsig, \ba}^{(\mathcal{L} - \mathcal{K}) \times (\mathcal{L} - \mathcal{K}),(2)} 
%  &\prec   \frac{W^{-d/3}}{\eta_u} \cdot  \br{W^{-d}\widetilde{{\cal T}}^{L} _{u,D}(|a_1-a_2|) },
% \\ 
(\mathcal{E} \otimes \mathcal{E})^{M,(2)}_{u,\bsig, \ba,\ba}&\prec \frac{ (\ilambda^2W^{d})^{-\frac12} }{\eta_u} \br{W^{-d}\widetilde{\cal T}^{L}_{u,D}\p{|a_1-a_2|}}^2.
\end{aligned}\ee
We recall the definition of $\cal U_{s,t,\bsig}^{(2)}$ in \eqref{def_Ustz} and notice the following identity (recall that $t\gtrsim1$ under \eqref{eq:assmtlarge}): 
\be\label{eq:decompU}\frac{1 - s \cdot M^{(\sig_1,\sig_2)}S^{(\sB)}}{1 - t \cdot M^{(\sig_1,\sig_2)}S^{(\sB)}} = \frac{s}{t} + \frac{t-s}{t}\Theta^{(\sig_1,\sig_2)}_t.\ee
If $\sig_1=\sig_2$, then $\Theta_{t}^{(\sig_1,\sig_2)}$ decays exponentially by \eqref{prop:ThfadC_short}. Using the estimates \eqref{S5WG+M000} and \eqref{S5WG+M}, the assumption \eqref{Eq:Gdecay+IND}, the decomposition \eqref{eq:decompU}, and the estimate \eqref{prop:ThfadC_short}, we obtain from \eqref{int_K-LcalE_n=2} that 
\[ (\mathcal{L} - \mathcal{K})^{(2)}_{t, \bsig, \ba} \prec (\ilambda^2W^{d})^{-\frac15}\cdot \br{W^{-d}\wT^{L}_{t,D}(|a_1-a_2|)}.\]
%which is stronger than \eqref{Eq:Gdecay_flow}. 
Since the proof is considerably simpler than in the case $\sig_1\ne \sig_2$  discussed below, we omit the details.
 
%Since the proof is much simpler than that for the $\sig_1\ne \sig_2$ case below, we omit the details. 

We now focus on the key challenge with $\sig_1\ne \sig_2$.
When $\sig_1\ne \sig_2$, with the estimates in \eqref{S5WG+M} and using the decomposition \eqref{eq:decompU}, we can bound the second and third terms on the RHS of \eqref{int_K-LcalE_n=2} as
\begin{align}
 \int_s^t\left(\mathcal{U}^{(2)}_{u, t, \bsig} \circ \mathcal{E}_{u, \bsig}^{(\mathcal{L}-\mathcal{K}) \times(\mathcal{L}-\mathcal{K}),(2)}\right)_{\mathbf{a}} \dd u&\prec (\ilambda^2W^{d})^{-\frac13}\int_s^t \frac{{\cal A}_{u,t, \ba} }{\eta_u} \dd u \prec (\ilambda^2W^{d})^{-\frac13}\sup_{u\in[s,t]}{\cal A}_{u,t, \ba},\label{jymwons-L-K}
    \\ 
    \int_s^t\left(\mathcal{U}^{(2)}_{u, t, \bsig} \circ \mathcal{E}_{u, \bsig}^{\Gc,(2)}\right)_{\mathbf{a}} \dd u
    &\prec (\ilambda^2W^{d})^{-\frac12}\int_s^t \frac{{\cal A}_{u,t,\ba}}{\eta_u} \dd u\prec (\ilambda^2W^{d})^{-\frac12}\sup_{u\in[s,t]}{\cal A}_{u,t, \ba},\label{jymwons-Gc}
 \end{align}
where the two-dimensional tensor $\cal A$ is defined as 
\begin{align}
   {\cal A}_{u,t, \ba}:=&~W^{-d}\widetilde{\cal T}^{L}_{u,D} \p{|a_1-a_2|}+(t-u)\sum_b \br{\Theta_{t,a_1b}^{(+,-)}\cdot  W^{-d}\widetilde{  \cal T}^{L}_{u,D}\p{|b-a_2|}
+ W^{-d}\widetilde{  \cal T}^{L}_{u,D}\p{|a_1-b|}\cdot \Theta_{t, b a_2}^{(+,-)}} \nonumber\\  
&~+(t-u)^2\sum_{b_1,b_2}\Theta_{t,a_1 b_1}^{(+,-)} \cdot  W^{-d}\widetilde{ \cal T}^{L}_{u,D}\p{|b_1-b_2|} \cdot \Theta_{t,b_2a_2}^{(+,-)},\quad \forall \ba=(a_1,a_2)\in (\Zn)^2. \label{eq:def_calA5}
\end{align}
Similarly, we write the martingale term as
\begin{align}
\int_{s}^t \left(\mathcal{U}^{(2)}_{u, t, \bsig} \circ \dd \mathcal{E}^{M,(2)}_{u, \bsig}\right)_{\ba} 
% &=\int_{s}^t
%  \frac {u^2}{t^2}\dd \mathcal{E}^{M,(2)}
% _{u, \bsig, {\ba}} 
% +\int_{s}^t
% \frac{t-u}{t}\left(\Theta_{t}^{(+,-)}\cdot \dd 
% \mathcal{E}^{M,(2)}
% _{u, \bsig}\right)_{{\ba}}
% +\int_{s}^t
% \frac{t-u}{t}\left(
% \dd \mathcal{E}^{M,(2)}
% _{u, \bsig}\cdot \Theta_{t}^{(+,-)}  \right)_{{\ba}}\nonumber
% \\
% &\quad +\int_{s}^t
% \frac{(t-u)^2}{t^2}\left(\Theta_{t}^{(\sB)} \cdot
% \dd \mathcal{E}^{M,(2)}
% _{u, \bsig}\cdot\Theta_{t}^{(+,-)}  \right)_{{\ba}}
% \nonumber \\
&=\int_{s}^t \frac {u^2}{t^2}\dd \mathcal{E}^{M,(2)}
_{u, \bsig, {{\ba}}} 
+ \left( \Theta_{t}^{(+,-)}\cdot \p{\int_{s}^t\frac{t-u}{t}
\dd \mathcal{E}^{M,(2)}
_{u, \bsig}}\right)_{{\ba}}
+ \left(\left(\int_{s}^t
\frac{t-u}{t}\dd\mathcal{E}^{M,(2)}
_{u, \bsig}\right)\cdot\Theta_{t}^{(+,-)}  \right)_{{\ba}}\nonumber
\\
&\quad +
 \left(\Theta_{t}^{(+,-)}\cdot\left(\int_{s}^t
 \frac{(t-u)^2}{t^2}
\dd\mathcal{E}^{M,(2)}
_{u, \bsig}\right)\cdot\Theta_{t}^{(+,-)}  \right)_{{\ba}}.\label{jymwons}
\end{align}
Combining the second bound in \eqref{S5WG+M} with the Burkholder-Davis-Gundy inequality (recall \eqref{aaswtghh}), we get   
 \begin{align*}
\int_{s}^t (t-u)^r\cdot 
\dd \mathcal{E}^{M,(2)}_{u, \bsig, {{\ba}}} \prec (1-s)^r (\ilambda^2W^{d})^{-\frac14} \cdot W^{-d}    \widetilde{{\cal T}}^{L} _{t,D}\p{|a_1-a_2|},\quad \forall r\in\{0,1,2\}.  
 \end{align*}
Plugging it into \eqref{jymwons} and using the $(\infty\to \infty)$-norm of $\Theta_t^{(+,-)}$ given by \eqref{eq:THETAinftinf}, we obtain 
\begin{align}\label{uuwmskiow}
\int_{s}^t \left(\mathcal{U}^{(2)}_{u, t, \bsig} \circ \dd\mathcal{E}^{M,(2)}_{u, \bsig}\right)_{{\ba}}
\prec \frac{(1-s)^2}{(1-t)^2} \cdot (\ilambda^2W^{d})^{-\frac14} \max_{u\in[s,t]}{\cal A}_{u,t, \ba}.
 \end{align}
%where we also use the $(\infty\to \infty)$-norm of $\Theta_t^{(+,-)}$ given by \eqref{eq:THETAinftinf}. 
% \be\label{ThetaB_infint} \|\Theta_t^{(+,-)}\|_{\infty\to \infty} = \max_a \sum_{b}\Theta^{(+,-)}_{t,ab} = \frac{1}{1-t}. \ee
Now, inserting the bounds \eqref{jymwons-L-K}, \eqref{jymwons-Gc}, and \eqref{uuwmskiow} into \eqref{int_K-LcalE_n=2} gives that 
\begin{align}\label{iois-mtx}
(\mathcal{L} - \mathcal{K})^{(2)}_{t, \bsig, {\ba}} \prec &
    \left(\mathcal{U}^{(2)}_{s, t, \bsig} \circ (\mathcal{L} - \mathcal{K})^{(2)}_{s, \bsig}\right)_{{\ba}}+ \frac{(1-s)^2}{(1-t)^2}\cdot (\ilambda^2W^{d})^{-\frac14} \sup_{u\in [s,t]} {\cal A}_{u,t,{\ba}}  .
\end{align}
Applying the estimate \eqref{prop:ThfadC} to $\Theta_t^{(+,-)}$, along with the bound \eqref{TTT2}, we obtain that
\begin{align}\label{uwp2-92kj}
 (1-u) \sum_b \Theta_{t,a_1b}^{(+,-)} \cdot \widetilde{  \cal T}^{L}_{u,D}\p{|b-a_2|} &\prec (1-u)\sum_b \cal T_{t}\p{|a_1-b|} \br{\cal T_{u}\p{|b-a_2|} +W^{-D}} \nonumber\\
 &\prec \cal T_{t}\p{|a_1-a_2|} + \frac{1-u}{1-t}W^{-D} \lesssim \widetilde{  \cal T}^{L}_{t,D-\theta}\p{|a_1-a_2|},
\end{align}
where $\theta:=2/(d\fc)$. In the last step, we used that $(1-u)/(1-t)\le (1-s)/(1-t)\le L^2\le N^{2/d}\le W^{\theta}$ in the present regime, by \eqref{Main_DEL_COND}.
Inserting this estimate into the definition of $\cal A$ in \eqref{eq:def_calA5}, we can bound ${\cal A}_{u,t,{\ba}}$ as   
\begin{align}\label{uwftgwesj}
\max_{u\in [s,t]} {\cal A}_{u,t,{\ba}}\prec \frac{1-s}{1-t}\cdot W^{-d}\widetilde{  \cal T}^{L}_{t,D-2\theta}\p{|a_1-a_2|}.
\end{align}
Combining this with \eqref{iois-mtx}, we obtain, for any large constant $D>0$,
\begin{align}\label{iois-mtx2}
(\mathcal{L} - \mathcal{K})^{(2)}_{t, \bsig, {\ba}} \prec &
    \left(\mathcal{U}^{(2)}_{s, t, \bsig} \circ (\mathcal{L} - \mathcal{K})^{(2)}_{s, \bsig}\right)_{{\ba}}+ (\ilambda^2W^{d})^{-\frac15} \cdot W^{-d}\widetilde{  \cal T}^{L}_{t,D}\p{|a_1-a_2|}.
\end{align}
For the first term on the RHS, we claim that
\begin{align}\label{iksjuwjx0}
\left(\mathcal{U}^{(2)}_{s, t, \bsig} \circ (\mathcal{L} - \mathcal{K})^{(2)}_{s, \bsig}\right)_{{\ba}} \prec (\ilambda^2W^{d})^{-\frac15}  \cdot W^{-d} {\cal T}_{t}\p{|a_1-a_2|}+W^{-D}.
  \end{align}
Combining \eqref{iois-mtx2} and \eqref{iksjuwjx0}, we conclude \eqref{Eq:Gdecay_flow} for the case $\ilambda^2/L^{2}\le 1- t\le 1-s \le \ilambda^2$. 

\begin{proof}[\bf Proof of \eqref{iksjuwjx0}]   
%To control the first term on the RHS of \eqref{iois-mtx2}, we again 
Using the decomposition \eqref{eq:decompU}, we can expand the LHS of \eqref{iksjuwjx0} as 
\begin{align}
\left(\mathcal{U}^{(2)}_{s, t, \bsig} \circ (\mathcal{L} - \mathcal{K})^{(2)}_{s, \bsig}\right)_{{\ba}}=~&\frac{s^2}{t^2}\left(\mathcal{L} - \mathcal{K}\right)^{(2)}_{s, \bsig,{{\ba}}}
+\frac{t-s}{t} \left(\Theta_{t}^{(+,-)} \cdot \left (\mathcal{L} - \mathcal{K}\right)^{(2)}_{s, \bsig}\right)_{{\ba}}
+\frac{t-s}{t} \left(\left (\mathcal{L} - \mathcal{K}\right)^{(2)}_{s, \bsig}\cdot\Theta_{t}^{(+,-)}  \right)_{{\ba}} \nonumber\\
&~+\frac{(t-s)^2}{t^2}  \left(\Theta_{t}^{(+,-)} \cdot\left (\mathcal{L} - \mathcal{K}\right)^{(2)}_{s, \bsig}\cdot\Theta_{t}^{(+,-)}  \right)_{{\ba}} .\label{uwkxkisjwj0}
\end{align}
With the inductive hypothesis \eqref{Eq:Gdecay+IND} on $\left(\mathcal{L} - \mathcal{K}\right)_{s,\bsig,\ba}^{(2)}$ and using \eqref{uwp2-92kj} again, we get that 
\be\label{uwkxkisjwj}
(t-s) \left|\left(\Theta_{t}^{(+,-)} \cdot \left (\mathcal{L} - \mathcal{K}\right)^{(2)}_{s, \bsig}\right)_{{\ba}}\right|+(t-s) \left| \left(\left (\mathcal{L} - \mathcal{K}\right)^{(2)}_{s, \bsig}\cdot\Theta_{t}^{(+,-)}  \right)_{{\ba}}\right|\prec  (\ilambda^2W^{d})^{-\frac15} \cdot W^{-d}\widetilde{  \cal T}^{L}_{t,D}\p{|a_1-a_2 |}.
\ee
%Thus, combining \eqref{iois-mtx2} Combining \eqref{Eq:Gdecay+IND} with 
From \eqref{Eq:Gdecay+IND}, \eqref{uwkxkisjwj0}, and \eqref{uwkxkisjwj}, we obtain that  
%\textcolor{blue}{Claim the following statement earlier, claim it as a key middle step for this proof. }
\begin{align}\label{ilowkidjsw}
   \left(\mathcal{U}^{(2)}_{s, t, \bsig} \circ (\mathcal{L} - \mathcal{K})^{(2)}_{s, \bsig}\right)_{{\ba}}&\prec
(\ilambda^2W^{d})^{-\frac15} \cdot W^{-d}\widetilde{  \cal T}^{L}_{t,D}\p{|a_1-a_2|}\nonumber\\
 &+(1-s)^2 \left| \left(\Theta_{t}^{(+,-)} \cdot\left (\mathcal{L} - \mathcal{K}\right)^{(2)}_{s, \bsig}\cdot\Theta_{t}^{(+,-)}  \right)_{{\ba}}\right| .
\end{align}
%To conclude the proof with \eqref{iois-mtx2} and \eqref{ilowkidjsw}, it remains to estimate the second term on the RHS of \eqref{ilowkidjsw}.
% Now combining \eqref{iois-mtx}, \eqref{uwftgwesj} and \eqref{ilowkidjsw}, we have (for  $\bsig=(+,-)$) 
% \begin{align}\label{ilow0pkkxmwsw}
% (\mathcal{L} - \mathcal{K})_{t, \bsig, {\ba}}= O_\prec (W^{-6d/5}\cdot \widetilde{  \cal T}^{\,dL,D }_{t, |a_1-a_2 |})\;+\; \eta_s^2 \cdot \left(\Theta_t \cdot(\mathcal{L}-\mathcal{K})^{(2)}_{s,  \bsig} \cdot \Theta_t\right)_{a_1a_2} 
% \end{align}
% In the remainder of this section, we estimate the last erm in the r.h.s. of \eqref{ilow0pkkxmwsw}. Now, for the case (i) with $\ilambda^2/L^{2}\le 1- t\le 1-s$, 
To show \eqref{iksjuwjx0}, it remains to prove the following estimate for $\ba=(a_1,a_2)$ and $\bsig\in \{(+,-),(-,+)\}$:   
\begin{align}\label{iksjuwjx}
(1-s)^2 \left(\Theta_t^{(+,-)} \cdot(\mathcal{L}-\mathcal{K})^{(2)}_{s, \bsig} \cdot \Theta_t^{(+,-)}\right)_{\ba}\prec (\ilambda^2W^{d})^{-\frac15}  \cdot W^{-d} {\cal T}_{t}\p{|a_1-a_2|}+W^{-D}.
  \end{align}
%Combining \eqref{iois-mtx2}, \eqref{ilowkidjsw}, and \eqref{iksjuwjx}, we conclude \eqref{Eq:Gdecay_flow} for the case $\ilambda^2/L^{2}\le 1- t\le 1-s \le \ilambda^2$. 

%\textcolor{blue}{maybe list \eqref{iksjuwjx}and \eqref{ilowkidjsw} much earlier, then say \textbf{proof of \eqref{iksjuwjx}}, \textbf{proof of \eqref{ilowkidjsw}} for their proofs, so the structure will be clear. Right below listing the \eqref{iksjuwjx}and \eqref{ilowkidjsw}, you can add your rough explaination. like: the proof \eqref{iksjuwjx} replies on the error estimate of $\cal E$ in *** and the $\cal U$ norm in ****. *******. For the proof \eqref{iksjuwjx}, we first reduce to the $f^{far}$ defined in  ****. The other parts can be easily bounded with the behavior of $\Theta$ and ****.  For the $f^{far}$, we split into two part $\mathbb Ef^{far}$ and $f^{far}-\mathbb Ef^{far}$. For the first part, we will use the translation varanit property of $\E {\cal L}$ and for the 2nd part we will apply ****. The details will be introduce int eh proof of Lemma 6.1.  }

% {\cor Note that if $(1-s)/(1-t)\le  (\log W)^{10}$, then it follows from \eqref{uwp2-92kj}, \eqref{uwftgwesj}, and \eqref{uwkxkisjwj} that 
% $$
% (1-s)^2 \left(\Theta_t^{(\sB)} \cdot(\mathcal{L}-\mathcal{K})^{(2)}_{s, \bsig} \cdot \Theta_t^{(\sB)}\right)_{{\ba}}
%  \prec \left(\frac{1-s}{1-t}\right)^2 \cdot  W^{-6d/5}\widetilde{\mathcal{T}}^{L}_{t,D}\p{|a_1-a_2|}\prec  W^{-6d/5}\widetilde{\mathcal{T}}^{L}_{t,D}\p{|a_1-a_2|}
% $$
% for any constant $D>0$, which already implies \eqref{iksjuwjx}. }
Recall that we assume \eqref{eq:assmtlarge}, under which we have  
\be\label{eq:ells_to_ellt}
%W^{-\frac{d}{C_d}}\ge (1-s)/(1-t)\ge  (\log W)^{10} \implies 
\ell_t\ge \ell_s\cdot (\log W)^{5}.
\ee
Due to the exponential decay of $\Theta_t^{(+,-)}$ and $({\cal L-\cal K})^{(2)}_{s,\bsig}$ given by \eqref{prop:ThfadC} and \eqref{Eq:Gdecay+IND}, respectively, we only need to focus on the following regime: 
\be\label{eq:ells_to_ellt2}
|a_1-a_2|\le \frac{1}{2}(\log W)^{3/2}\cdot \ell_t+
 (\log W)^{5/2}\cdot \ell_s \le (\log W)^{3/2}\cdot \ell_t;
\ee
otherwise the LHS of \eqref{iksjuwjx} is an error of order $W^{-D}$ for any large constant $D>0$. Under this condition, we have \smash{$\exp\big(\p{|a_1-a_2|/\ell_t}^{1/2}\big)\prec 1$}, so proving \eqref{iksjuwjx} amounts to establishing the following bound:
\be\label{iksjuwjx2} (1-s)^2 \left(\Theta_t^{(+,-)} \cdot(\mathcal{L}-\mathcal{K})^{(2)}_{s, \bsig} \cdot \Theta_t^{(+,-)}\right)_{\ba}\prec \frac{(\ilambda^2W^{d})^{-\frac65}}{|a_1-a_2|^{d-2}+1}. \ee
To show this estimate, we write the LHS of \eqref{iksjuwjx2} as 
\begin{align*}
&\ \quad (1-s)^2 \sum_{b_1,b_2}\Theta_{t,a_1b_1}^{(+,-)}(\mathcal{L}-\mathcal{K})^{(2)}_{s, \bsig,(b_1,b_2)}  \p{\Theta_{t,b_2a_2}^{(+,-)}-\Theta_{t,b_1a_2}^{(+,-)}}+(1-s)^2 \sum_{b_1,b_2}\p{\Theta_{t,a_1b_1}^{(+,-)}\Theta_{t,a_2b_1}^{(+,-)}} (\mathcal{L}-\mathcal{K})^{(2)}_{s, \bsig,(b_1,b_2)} \\
&=(1-s)^2 \sum_{b_1,b_2}\Theta_{t,a_1b_1}^{(+,-)}(\mathcal{L}-\mathcal{K})^{(2)}_{s, \bsig,(b_1,b_2)}  \p{\Theta_{t,b_2a_2}^{(+,-)}-\Theta_{t,b_1a_2}^{(+,-)}}+(1-s)^2 \sum_{b_1} \Theta_{t,a_1b_1}^{(+,-)} \Theta_{t,a_2b_1}^{(+,-)}  \frac{\im \tr\p{(G_s-M)E_{b_1}}}{W^d\eta_s} \\
&=: f_{\ba} + g_{\ba}.
\end{align*}
Above, we have applied Ward's identities \eqref{WI_calL} and \eqref{WI_calK} to express $\sum_{b_2} (\mathcal{L}-\mathcal{K})^{(2)}_{s, \bsig,(b_1,b_2)}$ as 
\be\label{eq:PcalBterm}
\sum_{b_2} (\mathcal{L}-\mathcal{K})^{(2)}_{s, \bsig,(b_1,b_2)}=\frac{\im {({\cal L-\cal K})^{(1)}_{s,+, b_1}}}{W^d\eta_s}=\frac{\im \tr\p{(G_s-M)E_{b_1}}}{W^d\eta_s}\prec \frac{(\ilambda^2W^d)^{-1}}{W^{d}\eta_s}, \ee
where we use the averaged local law \eqref{Gt_avgbound_flow} for $\tr\p{(G_s-M)E_{b_1}}$ in the last step. With \eqref{eq:PcalBterm}, we can bound $g_{\ba}$ as follows for any large constant $D>0$:
\begin{align}
g_{\ba} &\prec \frac{1-s}{\ilambda^2W^{2d}} \sum_{b_1} \Theta_{t,a_1b_1}^{(+,-)} \Theta_{t,a_2b_1}^{(+,-)}  \lesssim \frac{1-s}{\ilambda^2W^{2d}} \sum_{b_1} \cT_{t}(|a_1-b_1|)\cT_{t}(|a_2-b_1|) \nonumber\\
&\lesssim \frac{1-s}{1-t} \p{\ilambda^2W^{d}}^{-1}\cdot W^{-d} \cT_{t}(|a_1-a_2|) \le \frac{(\ilambda^2W^{d})^{-\frac65} }{|a_1-a_2|^{d-2}+1},\label{eq:boundga}
% &\prec \frac{1-s}{\ilambda^2W^{2d}} \sum_{b_1:|b_1-a_1|\vee|b_1-a_2|\le (\log W)^2 \ell_t}\frac{1}{|a_1-b_1|^{d-2}+1} \frac{1}{|a_2-b_1|^{d-2}+1} +W^{-D}\nonumber\\
% & \prec \frac{1-s}{\ilambda^2 W^{2d}}\p{\frac{1}{|a_1-a_2|^{d-4}+1}+\frac{1}{\ell_t^{d-4}}} \prec \frac{(1-s)\ell_t^2}{\ilambda^2 W^{2d}}\cdot \frac{1}{|a_1-a_2|^{d-2}+1}  \le \frac{(\ilambda^2W^{d})^{-\frac15}  \cdot W^{-d} }{|a_1-a_2|^{d-2}+1},\label{eq:boundga}
\end{align}
where, in the second step, we apply the estimate \eqref{prop:ThfadC}, and in the third step, we use the bound \eqref{TTT2}.


To conclude \eqref{iksjuwjx2}, it remains to show that 
 \be\label{iksjuwjx3} 
% (1-s)^2 \left(\Theta_t^{(\sB)} \cdot {\cal B} \cdot \Theta_t^{(\sB)}\right)_{\ba}
 f_{\ba}\prec \frac{(\ilambda^{2}W^d)^{-\frac65} }{|a_1-a_2|^{d-2}+1}. \ee
To simplify notation, we abbreviate ${\cal B}:=(\mathcal{L} - \mathcal{K})^{(2)}_{s, \bsig}$. By the assumption \eqref{Eq:Gdecay+IND}, we have
\be\label{eq:propcalB} 
{\cal B}_{b_1b_2}\prec \frac{(\ilambda^{2}W^d)^{-\frac65} }{|b_1-b_2|^{d-2}+1},\quad \text{and} \quad  {\cal B}_{b_1b_2}\prec W^{-D}\quad \text{whenever}\ \  |b_1-b_2|\ge (\log W)^3\ell_s ,
\ee
for any large constant $D>0$. 
To show \eqref{iksjuwjx3}, we decompose $f_{\ba}$ into the following two parts:
\begin{align*}
    f_{\ba}^{\mathrm{near}}&:=(1-s)^2\sum_{b_1: |b_1-a_1|\wedge |b_1-a_2| \le (\log W)^4\ell_s} \sum_{b_2} \Theta_{t,a_1b_1}^{(+,-)}\cal B_{b_1b_2}\p{\Theta_{t,b_2a_2}^{(+,-)}-\Theta_{t,b_1a_2}^{(+,-)}},\\  f_{\ba}^{\mathrm{far}}&:=(1-s)^2\sum_{b_1: |b_1-a_1|\wedge |b_1-a_2| > (\log W)^4\ell_s} \sum_{b_2} \Theta_{t,a_1b_1}^{(+,-)}\cal B_{b_1b_2}\p{\Theta_{t,b_2a_2}^{(+,-)}-\Theta_{t,b_1a_2}^{(+,-)}}.
\end{align*}
For the term $f_{\ba}^{\mathrm{near}}$, using \eqref{eq:propcalB} and \eqref{prop:ThfadC}, we get that for any constant $D>0$, 
\begin{align*}
 f_{\ba}^{\mathrm{near}} &\prec  (1-s)^2\sum_{b_1}^* \sum_{b_2:|b_2-b_1|\le (\log W)^3\ell_s}\frac{\ilambda^{-2}}{|a_1-b_1|^{d-2}+1}\frac{(\ilambda^{2}W^d)^{-\frac65}}{|b_1-b_2|^{d-2}+1} \p{\frac{\ilambda^{-2}}{|b_2-a_2|^{d-2}+1}+\frac{\ilambda^{-2}}{|b_1-a_2|^{d-2}+1}} + W^{-D}
 \\
 &\prec  (1-s)^2\sum_{b_1}^* \frac{(\ilambda^{2}W^d)^{-\frac65}}{|a_1-b_1|^{d-2}+1}\cdot {\frac{\ilambda^{-4}\ell_s^2}{|b_1-a_2|^{d-2}+1}}%+(1-s)^2\sum_{b_1}^* \frac{\ilambda^{-6}W^{-6d/5}}{|a_1-b_1|^{d-2}+1}\frac{\ell_s^2}{|a_2-b_1|^{d-2}+1}
 \prec (\ilambda^{2}W^d)^{-\frac65} \cdot \frac{(1-s)^2 \ilambda^{-4}\ell_s^4}{|a_1-a_2|^{d-2}+1}  \lesssim \frac{(\ilambda^{2}W^d)^{-\frac65}}{|a_1-a_2|^{d-2}+1},
\end{align*}
where $\sum^*_{b_1}$ refers to the summation over $\{b_1:|b_1-a_1|\wedge |b_1-a_2| \le (\log W)^4\ell_s\}$. 
To control the term \smash{$f_{\ba}^{\mathrm{far}}$}, we employ a CLT-type cancellation mechanism, extending the approach developed for 2D random band matrices in \cite{DYYY25}.


\begin{lemma}\label{lem;CLT}
In the setting of \Cref{lem:main_ind}, fix any $\ilambda^2/L^{2}\le 1-t \le 1-s\le \ilambda^2$ so that \eqref{con_st_ind} holds. Then, $f_{\ba}^{\mathrm{far}}$ satisfies the following estimate under the conditions \eqref{eq:ells_to_ellt} and \eqref{eq:ells_to_ellt2}:
 \be\label{iksjuwjx3_far} f_{\ba}^{\mathrm{far}}\prec \frac{(\ilambda^{2}W^d)^{-\frac65} }{|a_1-a_2|^{d-2}+1}. \ee
\end{lemma}
With this lemma, we conclude the estimate \eqref{iksjuwjx3}, which implies \eqref{iksjuwjx}, and hence completes the proof of \eqref{iksjuwjx0} for the case $\ilambda^2/L^{2}\le 1-t \le 1-s\le \ilambda^2$ together with \eqref{ilowkidjsw}. 
%\eqref{Eq:Gdecay_flow} for the case $\ilambda^2/L^{2}\le 1-t \le 1-s\le \ilambda^2$ together with \eqref{iois-mtx2} and \eqref{ilowkidjsw}. 
\end{proof}


\begin{proof}[\bf Proof of \Cref{lem;CLT}]
We decompose $f_{\ba}^{\mathrm{far}}$ as its expectation $\E f_{\ba}^{\mathrm{far}}$ plus the fluctuation part $\IE f_{\ba}^{\mathrm{far}}$, where $\IE$ denotes $\mathbb {IE}:=1-\E$, and $\E f_{\ba}^{\mathrm{far}}$ and $\IE f_{\ba}^{\mathrm{far}}$ represent the expressions obtained by replacing $\cal B$ in $f_{\ba}^{\mathrm{far}}$ with $\E\cal B$ and $\IE \cal B$, respectively. 
Due to the translation invariance and symmetry of $\E\cal B$, the expectation part $\E f_{\ba}^{\mathrm{far}}$ involves a second-order difference of the $\Theta^{(+,-)}$-propagator, which satisfies the improved bound \eqref{prop:BD2} that is summable in $a$. 
For the fluctuation part $\IE f_{\ba}^{\mathrm{far}}$, the $\Theta^{(+,-)}$-propagator effectively transfers the estimate of $\IE \cal B$ from shorter scales of order $\ell_s$ to the larger scale $\ell_t$. Moreover, intuitively, the random variables \smash{$\IE \cal B_{b_1b_2}$ and $\IE \cal B_{b_1'b_2'}$} become asymptotically independent when $|b_1-b_1'|$ exceeds the typical decay scale $\ell_s$ of $\cal B$. 
Thus, the term $\IE f_{\ba}^{\mathrm{far}}$ can be viewed as a superposition of roughly $\ell_t^d/\ell_s^d$ asymptotically independent, centered random variables, yielding an additional improvement through a CLT-type argument. 


%\textcolor{blue}{say something about the reason it works, like $\E {\cal B}$ has the following property so the leading term in the expansion (***) will disappear. For $\mathbb {IE} B$, the sum is over the scale of $\ell_t$, but $B$ is uncorrelated beyond $\ell_s$ scale. Therefore with a CLT type moment estimate, we can obtain an extra factor $\ell_s/\ell_t$ ...  (introduce the basic logic/reason in each long tedious proof) }.
% \begin{align*}
% 	f_{\ba}^{\mathrm{far}}&=(1-s)^2\sum_{b_1}^* \sum_{b_2} \Theta_{t,a_1b_1}^{(\sB)}\cal B_{b_1b_2}\p{\Theta_{t,b_2a_2}^{(\sB)}-\Theta_{t,b_1a_2}^{(\sB)}} =\E f_{\ba}^{\mathrm{far}} + \IE f_{\ba}^{\mathrm{far}} ,
% \end{align*}


We now proceed to the formal proof. 
By the translation invariance and symmetry of our model on the block level, $\E \cal B$ is also translationally invariant and symmetric in the following sense:
\[ \E \cal B(a+c,b+c)=\E\cal B(a,b),\quad \E \cal B(a,b)=\E\cal B(b,a),\quad \forall a,b,c\in \Zn.\]
These imply that \( \E	\cal B(b_1,b_2)=\E \cal B(b_2,b_1)=\E\cal B(b_1,2b_1-b_2)\), with which we get 
\begin{align*}
    \sum_{b_2} \E\cal B_{b_1b_2}\p{\Theta_{t,b_2a_2}^{(+,-)}-\Theta_{t,b_1a_2}^{(+,-)}}&=\sum_{b_2} \E \cal B_{b_1,(2b_1-b_2)}\p{\Theta_{t,b_2a_2}^{(+,-)}-\Theta_{t,b_1a_2}^{(+,-)}} =\sum_{b_2} \E \cal B_{b_1b_2}\p{\Theta_{t,(2b_1-b_2)a_2}^{(+,-)}-\Theta_{t,b_1a_2}^{(+,-)}}.
\end{align*}
Thus, we can rewrite $\E f_{\ba}^{\mathrm{far}}$ as 
\begin{align*}
\E	f_{\ba}^{\mathrm{far}}&=(1-s)^2\sum_{b_1}^\star \sum_{b_2} \Theta_{t,a_1b_1}^{(+,-)}\E\cal B_{b_1b_2}\cdot \frac 12 \p{\Theta_{t,b_2a_2}^{(+,-)}+\Theta_{t,(2b_1-b_2)a_2}^{(+,-)} -2\Theta_{t,b_1a_2}^{(+,-)}},
\end{align*}
where $\sum^\star_{b_1}$ refers to the summation over $\{b_1:|b_1-a_1|\wedge |b_1-a_2| > (\log W)^4\ell_s\}$. Using \eqref{prop:ThfadC} and \eqref{prop:BD2}, along with the properties of $\cal B$ in \eqref{eq:propcalB}, we can bound the above expression by 
\begin{align}
	\E	f_{\ba}^{\mathrm{far}} &\prec (1-s)^2\sum_{b_1}^\star \frac{\ilambda^{-2}}{|a_1-b_1|^{d-2}+1}\frac{\ilambda^{-2}\ell_s^2}{|a_2-b_1|^{d}+1}  \sum_{b_2:|b_2-b_1|\le (\log W)^3\ell_s} \frac{(\ilambda^{2}W^d)^{-\frac65}}{|b_1-b_2|^{d-2}+1} + W^{-D} \nonumber\\
	&\prec \frac{(\ilambda^{2}W^d)^{-\frac65}\cdot (1-s)^2\ilambda^{-4}\ell_s^4}{|a_1-a_2|^{d-2}+1} \lesssim \frac{(\ilambda^{2}W^d)^{-\frac65}}{|a_1-a_2|^{d-2}+1} \label{eq:boundEfar}
\end{align}
for arbitrarily large constant $D>0$.

To deal with the term $\IE f_{\ba}^{\mathrm{far}}$, we need to explore the CLT cancellation within it as in \cite[Section 7]{DYYY25} and \cite[Appendix A.10]{RBSO1D}. By Markov's inequality and using $\ilambda^{4}(1-s)^{-2}\asymp\ell_s^4$, it suffices to prove the following $(2p)$-th moment bound for any fixed $p\in \N$:
\be\label{eq:main_challenge3}
\E\abs{\ilambda^4(1-s)^{-2}(\ilambda^{2}W^d)^{\frac65}\cdot  \IE f_{\ba}^{\mathrm{far}}}^{2p} \prec \left(\frac{\ell_s^4}{|a_1-a_2|^{d-2}+1}\right)^{2p} . 
\ee
Abbreviating \smash{$\Xi_{a_2}\p{b_1,b_2}:=\ilambda^2(\Theta_{t,b_2a_2}^{(+,-)}-\Theta_{t,b_1a_2}^{(+,-)})$} and $\mathsf B:=(\ilambda^{2}W^d)^{\frac65}\cal B$, we can write the LHS of \eqref{eq:main_challenge3} as 
\be\label{eq:2p_product}\sum_{\bfb^{(1)},\ldots, \bfb^{(2p)}}^{\star} \prod_{k=1}^{2p}\p{\ilambda^2\Theta_{t}^{(+,-)} (a_1,b_1^{(k)})\cdot \Xi_{a_2}(b_1^{(k)},b_2^{(k)})} \cdot \E\bigg\{\prod_{k=1}^{p}   \IE \mathsf B_{b_1^{(k)}b_2^{(k)}} \cdot \prod_{k=p+1}^{2p} \IE \overline{\mathsf B}_{b_1^{(k)}b_2^{(k)}}\bigg\}+\OO(W^{-D})\ee
for any large constant $D>0,$ where $\sum^\star$ refers to the summation of \smash{$\bfb^{(k)}=(b_1^{(k)},b_2^{(k)})$} over the regions 
\be\label{eq:sumregionsforb}\left\{(b_1^{(k)},b_2^{(k)}):|b_1^{(k)}-a_1|\wedge |b_1^{(k)}-a_2| > (\log W)^4\ell_s,\ |b_1^{(k)}-b_2^{(k)}| \le (\log W)^3\ell_s\right\}.\ee 

First, suppose the following pairing condition holds: 
\be\label{eq:pairingcond}
\text{for any }\ k\in \qqq{2p},\ \text{ there exists } \ l\ne k \ \text{ such that } \ |b^{(k)}_1-b^{(l)}_1|\le 10(\log W)^3\ell_s.\ee  
We can divide $\qqq{2p}$ into a disjoint union of $r$ subsets $\qqq{2p}=\sqcup_{i=1}^r A_i,$ constructed such that for every $k\in A_i$, any $l\in \qqq{2p}$ satisfying \smash{$|b^{(k)}_1-b^{(l)}_1|\le 10(\log W)^3\ell_s$} also belongs to $A_i$. According to the pairing condition, we have $r\le p$. For each $i\in\qqq{r}$, let $k_i$ denote an arbitrary index in $A_i$. Subject to the condition \eqref{eq:pairingcond} and the above decomposition $\qqq{2p}=\sqcup_{i=1}^r A_i$ (the corresponding summation being denoted by \smash{$\sum^{\star_1}$}), we can bound \eqref{eq:2p_product} as follows:
\begin{align}
    \label{eq:2p_product_pair}
&~\sum^{\star_1}_{\bfb^{(1)},\ldots, \bfb^{(2p)}} \prod_{i=1}^{r}\left(\prod_{k\in A_i}\p{\ilambda^2\Theta_{t}^{(+,-)} (a_1,b_1^{(k)})\cdot \Xi_{a_2}(b_1^{(k)},b_2^{(k)})}\cdot \big|\IE \mathsf B_{b_1^{(k)}b_2^{(k)}}\big|\right) \\
\prec &~\sum^{\star_1}_{\bfb^{(1)},\ldots, \bfb^{(2p)}} \prod_{i=1}^{r}\left( \frac{1}{\big|a_1-b_1^{(k_i)}\big|^{d-2}+1}\right)^{|A_i|}\p{\frac{\ell_s}{\big|a_2-b_1^{(k_i)}\big|^{d-1}+1}}^{|A_i|}\prod_{k\in A_i}\p{\frac{1}{\big|b_1^{(k)}-b_2^{(k)}\big|^{d-2}+1}} \nonumber\\
\prec &~\sum^{\star_2}_{b_1^{(k_1)},\ldots, b_1^{(k_r)}} \prod_{i=1}^{r}\left( \frac{1}{\big|a_1-b_1^{(k_i)}\big|^{d-2}+1}\right)^{|A_i|}\p{\frac{\ell_s}{\big|a_2-b_1^{(k_i)}\big|^{d-1}+1}}^{|A_i|} (\ell_s^{d+2})^{|A_i|-1}\ell_s^2 \nonumber \\
\lesssim &~\prod_{i=1}^{r}\left( \frac{1}{|a_1-a_2|^{d-2}+1}\right)^{|A_i|}\frac{\ell_s^d}{(\ell_s^{d-2})^{|A_i|}} \cdot (\ell_s^{d+2})^{|A_i|-1} \ell_s^2=\left(\frac{\ell_s^4}{|a_1-a_2|^{d-2}+1}\right)^{2p}.\nonumber 
\end{align}
Here, in the first step, we use the estimates \eqref{prop:ThfadC}, \eqref{prop:BD1}, and \eqref{eq:propcalB}, along with the conditions in \eqref{eq:sumregionsforb} and \eqref{eq:pairingcond}. In the second step, we take the summation over $\bfb^{(k)}$ for $k\in A_i\setminus\{k_i\}$ and over \smash{$b_2^{(k_i)}$} under the restrictions given by \eqref{eq:sumregionsforb} and \eqref{eq:pairingcond}, yielding a factor $(\ell_s^{d+2})^{|A_i|-1}\ell_s^2$ for each $i\in\qqq{r}$, and $\sum^{\star_2}$ refers to the summation under the constraints \smash{$|b_1^{(k_i)}-a_1|\wedge |b_1^{(k_i)}-a_2| > (\log W)^4\ell_s$} for $i\in\qqq{r}$. In the third step, we have used the following bound for any $k\ge 2$:
\begin{align}\label{eq:simplecalculus}
   \sum_{b:|b-a_1|\wedge |b-a_2| > (\log W)^4\ell_s} \left( \frac{1}{|a_1-b|^{d-2}+1}\right)^{k}\p{\frac{\ell_s}{|a_2-b|^{d-1}+1}}^{k}\lesssim \p{\frac{1}{|a_1-a_2|^{d-2}+1}}^k \frac{\ell_s^d}{(\ell_s^{d-2})^k} . 
\end{align}
To see why this holds, consider the case $|b-a_1|\ge |b-a_2|$---the case $|b-a_1|\le |b-a_2|$ can be handled similarly. We can bound the corresponding summation by 
\begin{align*}
  \left( \frac{1}{|a_1-a_2|^{d-2}+1}\right)^{k} \sum_{b:|b-a_1|\wedge |b-a_2| > (\log W)^4\ell_s} \p{\frac{\ell_s}{|a_2-b|^{d-1}+1}}^{k}\prec \left( \frac{1}{|a_1-a_2|^{d-2}+1}\right)^{k} \frac{\ell_s^d}{(\ell_s^{d-2})^k} , 
\end{align*}
where we use that $(|a_2-b|^{d-1}+1)^{-k}$ is summable when $d\ge 3$ and $k\ge 2$. This gives \eqref{eq:simplecalculus}.


Now, with \eqref{eq:2p_product_pair}, we obtain \eqref{eq:main_challenge3} under the pairing condition \eqref{eq:pairingcond}. It remains to control \eqref{eq:2p_product} when \eqref{eq:pairingcond} does not hold. In fact, we can prove the following result: if there exists an isolated vertex $\bfb^{(i)}$ such that \smash{\(\min_{j:j\ne i}|b_1^{(i)}-b_1^{(j)}|\ge 10(\log W)^3 \ell_s\)}, then for any large constant $D>0$,
\be\label{eq:bound_isolated}
\E\bigg\{\prod_{k=1}^{p}   \IE \mathsf B_{b_1^{(k)}b_2^{(k)}} \cdot \prod_{k=p+1}^{2p} \IE \overline{\mathsf B}_{b_1^{(k)}b_2^{(k)}}\bigg\} \le W^{-D}. 
\ee
Since the proof of this bound is exactly the same as that for \cite[equation (7.39)]{DYYY25} and \cite[equation (A.112)]{RBSO1D} (which does not depend on the dimension $d$), we omit the details here. This completes the proof of \eqref{eq:main_challenge3}, and hence concludes \Cref{lem;CLT} together with \eqref{eq:boundEfar}. 
\end{proof}

\subsection{The case \texorpdfstring{$\ilambda^2/L^{d} \leq 1-t \leq 1-s \leq \ilambda^2/L^{2}$}{1-t>L-d}}

The $\sig_1=\sig_2$ case is straightforward by analyzing equation \eqref{int_K-LcalE_n=2} with the estimates \eqref{Eq:Gdecay_w}, \eqref{eq:LW_conclusion_exp}, \eqref{eq:MG_conclusion3}, the decomposition \eqref{eq:decompU}, and the exponential bound \eqref{prop:ThfadC_short}. We therefore omit the details and focus on the case $\sig_1\ne \sig_2$ in the following.
In this case, we first remove the zero mode from the $(\cL-\cK)^{(2)}$-loop using the zero-mode-removing operator introduced in \Cref{def;zero_mode_remove}. Once the zero mode is removed, the rescaled propagator \smash{$(1-s)\zTheta_t^{(+,-)}$} becomes summable by \eqref{prop:ThfadC0}.


Applying Ward's identities \eqref{WI_calL} and \eqref{WI_calK}, we can express \smash{$(\mathcal{L} - \mathcal{K})_{t, \bsig}^{(2)}$} as: 
\begin{align}\label{zYU1}
    (\mathcal{L} - \mathcal{K})^{(2)}_{t, \bsig, {\ba}} - \left[Q^{(1)}\circ(\mathcal{L} - \mathcal{K})^{(2)}_{t, \bsig}\right]_{{\ba}}&=
\frac{\im {({\cal L-\cal K})^{(1)} _{t,+, a_2}}}{N\eta_t}  \prec {\frac{(\ilambda^{2}W^{d})^{-1}}{N\eta_t}},
\end{align}
where, in the second step, we apply the averaged local law \eqref{Gt_avgbound_flow}. 
Then, we analyze the term \smash{$Q^{(1)}\circ(\mathcal{L} - \mathcal{K})^{(2)}_{t, \bsig}$} using equation \eqref{iisuwjyys} with $n=2$. By definition, we have  \smash{$Q^{(1)}\circ\Theta_t^{(+,-)}=\zTheta_t^{(+,-)}$}. Using \eqref{prop:ThfadC0}, we obtain the following counterparts to the equations \eqref{eq:THETAinftinf} and \eqref{uwp2-92kj}: 
\begin{align}\label{ThetaBcirc_infint} 
\|\zTheta_t^{(+,-)}\|_{\infty\to \infty} = \max_a \sum_{b}\big|\zTheta^{(+,-)}_{t,ab}\big| & \prec \ilambda^{-2}L^2, \\ \label{uwp2-92kj00}
 (1-u) \sum_b \zTheta_{t,a_1b}^{(+,-)} \cdot \widetilde{  \cal T}^{L}_{u,D}\p{|b-a_2|} &\prec \widetilde{  \cal T}^{L}_{t,D}\p{|a_1-a_2|}.
\end{align}
With these estimates in hand, and employing a similar argument to the one above (i.e., the one leading to \eqref{iois-mtx2} and \eqref{ilowkidjsw}), we can derive the following estimate for  \smash{$Q^{(1)}\circ(\mathcal{L} - \mathcal{K})^{(2)}_{t, \bsig}$}: 
\begin{align}\label{zYU2}
\left[Q^{(1)}\circ(\mathcal{L} - \mathcal{K})^{(2)}_{t, \bsig}\right]_{{\ba}}
& \prec (\ilambda^{2}W^d)^{-\frac15}\cdot W^{-d}\widetilde{  \cal T}^{L }_{t,D}\p{|a_1-a_2|}+ (1-s)^2 \left(\zTheta_t^{(+,-)} \cdot(\mathcal{L}-\mathcal{K})^{(2)}_{s,  \bsig} \cdot \Theta_t^{(+,-)}\right)_{\ba} .
% \\
% & =\OO_\prec (W^{-6d/5}\cdot  \widetilde{  \cal T}^{\,dL,D }_{t, |a_1-a_2 |})\;+\; \eta_s^2 \cdot \left(Q^{(1)}\circ\Theta_t \cdot(\mathcal{L}-\mathcal{K})^{(2)}_{s,  \bsig} \cdot \Theta_t\right)_{a_1a_2} .
\end{align}
To control the second term on the RHS, we apply the bounds \eqref{prop:ThfadC}, \eqref{prop:ThfadC0}, and \eqref{Eq:Gdecay+IND} to get that 
\begin{align*}
    &~(1-s)^2 \left(\zTheta_t^{(+,-)} (\mathcal{L}-\mathcal{K})^{(2)}_{s,  \bsig}  \Theta_t^{(+,-)}\right)_{\ba}\\
    \prec &~  (1-s)^2 \frac{(\ilambda^2 W^{d})^{-\frac15}}{W^{d}} \sum_{b_1,b_2}\frac{\ilambda^{-2}}{|a_1-b_1|^{d-2}+1}\p{\frac{\ilambda^{-2}}{|b_1-b_2|^{d-2}+1}+\frac{1}{L^d(1-s)}} \p{ \frac{\ilambda^{-2}}{|b_2-a_2|^{d-2}+1} + \frac{1}{L^d(1-t)}}  \\
    \prec &~ (1-s) \frac{(\ilambda^2 W^{d})^{-\frac15}}{W^{d}}  \sum_{b_1}\frac{\ilambda^{-2}}{|a_1-b_1|^{d-2}+1} \left( \frac{\ilambda^{-2}}{|b_1-a_2|^{d-2}+1}+\frac{1}{L^d (1-t)}\right) \\
    \prec &~ \frac{(\ilambda^2 W^{d})^{-\frac15}}{W^{d}} \left(\frac{(1-s)\ilambda^{-4}L^2}{|a_1-a_2|^{d-2}+1}+\frac{(1-s)\ilambda^{-2}L^2}{L^{d}(1-t)}\right) \lesssim (\ilambda^2 W^{d})^{-\frac15} \cdot W^{-d}\widetilde{\mathcal{T}}_{t,D}^{L}\p{ |a_1-a_2|},
\end{align*}
where we use $B_{s,0}\asymp \ilambda^{-2}$ in the first step, and $1-s\le \ilambda^2/L^{2}$ in the second and fourth steps. Plugging it into \eqref{zYU2} and using \eqref{zYU1}, we complete the proof of \eqref{Eq:Gdecay_flow} for the case $\ilambda^2/L^{d} \leq 1-t \leq 1-s \leq \ilambda^2/L^{2}$. 

\subsection{The case \texorpdfstring{$1-t \ge \ilambda^2$}{1-t>lambda2}}\label{sec:Step5_larget}


In this case, the length scale $\ell_u$ remains identically equal to 1 throughout the evolution $u\in[s,t]$ (recall \eqref{eq:ellt}). With the estimate \eqref{Eq:Gdecay_w} established in Step 2, the proof of \eqref{Eq:Gdecay+s<g_flow} is relatively straightforward and may be viewed as a special case of the argument in \cite[Section 5.3]{YY_25}, where decay estimates for 2-$G$-loops were derived in the context of 1D random band matrices. 
The key reason that the argument of \cite{YY_25} applies in the regime $1-t \ge \ilambda^2$ is that the prefactor $(W^{-d}B_{u,0})^{2}$ in the target estimate \eqref{Eq:Gdecay+s<g_flow} is of order $(W^d|1-u|)^{-2}$. Unlike the regime $1-u\ll \ilambda^2$, this quantity has no polynomial dependence on $|a_1-a_2|$, precisely matching the setting in dimension 1. 
In fact, our proof here is much simpler, since we have already established exponential decay in \eqref{Eq:Gdecay_w} and obtained the optimal (maximum) $(\cal L - \cal K)$-loop estimate \eqref{Eq:L-KGt-flow}. 

Given $u\in[s,t]$ and a sufficiently large constant $D>0$, we introduce the following tail function to control the tail behavior of the $(\cL-\cK)$-loops: 
\begin{align}   %\label{def_WTu} 
%{T}^{(\cal L-K)}_u (\ell)&:= \max_{|a-b|\ge \ell} \Big| ({\cal L-\cK})^{(2)}_{u,  (+,-),  ( a,b)}\Big|, %\quad \bsig=(+,-), 
% \end{equation}
% \begin{equation}
{T} _{u, D} (r) :=  (W^d|1-u|)^{-2}
\exp(- r^{1/2})+W^{-D},\quad \forall r\ge 0 .\label{def_WTuD}
\end{align} 
%which will replace the role of the tail function in \eqref{defTUL} in our proof.  
% Note that they are both non-increasing functions in $\ell$. Furthermore, we introduce the ratio between them by  
% \begin{equation}\label{def_Ju}
%      {\cal J} _{u,D} (\ell):=   {\cal T}^{(\cal L-K)}_u (\ell) \big/{\cal T}_{u,D} (\ell) +1\, .
% \end{equation} 
% Then, to show \eqref{Eq:Gdecay_w}, it suffices to prove that for all $u\in [s,t]$ and any large constant $D>0$: 
%  \begin{align}\label{shoellJJ}
% \max_{0\le \ell \le n}   {\cal J} _{u,D} (\ell) \prec (\eta_s/\eta_u)^4 .
%   \end{align}
%Similar to \eqref{defCALJ_deter}, 
Let ${\cal J}^* _{u,D}\ge 1$ be a \emph{deterministic control parameter} such that 
\be\label{eq:def_new_J*}
\max_{\ba=(a,b)} \max_{\bsig}{\big|\left({\cal L-\cal K}\right)^{(2)}_{u,\bsig, \ba}\big|}\big/{ T}_{u,D}\p{|a-b|} \prec {\cal J}^* _{u,D}.
\ee
%We also introduce an intermediate scale parameter $\ell^{*}_u := (\log W)^{3/2}$. Note that on the scale $\ell_u^*$, the propagator \smash{$\Theta_u^{(+,-)}$} is exponentially small (i.e., faster than any polynomial decay) by \eqref{prop:ThfadC}, whereas ${T}_{u,D}$ is not (i.e., slower than any polynomial decay). 
With the tail function \eqref{def_WTuD} and the control parameter ${\cal J}^* _{u,D}$, we can obtain the following estimates on the second to fourth terms on the RHS of \eqref{int_K-LcalE_n=2}. 


\begin{lemma}\label{lem_dec_calE} In the setting of  \Cref{lem:main_ind}, suppose $1-t\ge \ilambda^2$ and the estimates \eqref{Gt_bound_flow}--\eqref{Eq:L-KGt-flow} hold. Then, for any $u\in [s,t]$, $\bsig\in\{+,-\}^2$, $\ba=(a_1,a_2)\in(\Zn)^2$, and large enough $D$ (such that $W^D\ge N^{12}$), the following estimates hold:  
\begin{align}\label{res_deccalE_lk}
& {\cal E}^{(\cL-\cK)\times(\cL-\cK),(2)}_{u, \bsig,\ba} \big/ 
  {T}_{u,D}(|a_1-a_2|) \prec  (1-u)^{-1}\cdot   \left(W^d|1-u|\right)^{-1}
 \left({\cal J}^{*}_{u,D} \right)^{2},\\
&\cal E^{\Gc,(2)}_{u, \bsig,\ba}  \big/ 
 {T}_{u,D}(|a_1-a_2|)
  \prec (1-u)^{-1}\br{{\bf 1}\left(|a_1-a_2|\le (\log W)^{\frac32}\right)  +  \left(W^d|1-u|\right)^{-\frac12}  \p{{\cal J}^{*}_{u,D}}^{\frac32}  }.
 \label{res_deccalE_wG}
\end{align} 
Furthermore, suppose $\ba'=(a_1',a_2')\in (\Zn)^2$ satisfies that \(\max_{i=1}^2 |a_i-a_i'|\le (\log W)^{3/2}.\) Then, we have that 
\begin{align}
\label{res_deccalE_dif}
\left(
   \mathcal{E}\otimes  \mathcal{E} \right)^{M,(2)}
   _{u, \bsig, \ba ,\ba'} \big/  
 {T}_{u,D}^2(|a_1-a_2|)
 & \prec (1-u)^{-1}\br{{\bf 1}\left(|a_1-a_2|\le 4(\log W)^{\frac32}\right)  +  \left(W^d|1-u|\right)^{-\frac12} \left({\cal J}^{*}_{u,D} \right)^{3}}.   
\end{align} 
\end{lemma}
\begin{proof}
The proof of this lemma is a special case of the argument for \cite[Lemma 5.7]{YY_25}, as explained above. In fact, many of the arguments can be simplified using the estimates \eqref{Gt_bound_flow}--\eqref{Eq:L-KGt-flow} established in the previous steps. We omit the details.
%is analogous to, and in fact much simpler than, the proof of \cite[Lemma 5.7]{YY_25}, by using \eqref{Gt_bound_flow}--\eqref{Eq:L-KGt-flow} and the definitions \eqref{def_WTuD} and \eqref{eq:def_new_J*}. Therefore, we omit the details.
\end{proof}

Another key ingredient in the proof is the following evolution kernel estimate for ${\cal U}^{(2)}_{s,t,\bsig}$.  

\begin{lemma}
\label{TailtoTail}
%Let \smash{${\cal A}: (\Zn)^{2}\to \mathbb C$} be a $2$-dimensional tensor. Then, for each $0\le s \le t < 1$, we have that  
%\begin{align}\label{sum_res_Ndecay}
%   \| {\cal U}^{(2)}_{s,t,\bsig}\;\circ {\cal A}\|_{\infty} \lesssim \left( |1-s|/ |1-t|\right)^{2}\cdot \|{\cal A}\|_{\infty} , 
%\end{align}
%where the $L^\infty$-norm of ${\cal A}$ is defined as $\|{\cal A}\|_{\infty}=\max_{\ba\in (\Zn)^2}|\cal A_{\ba}|$. 
Suppose $1-s\ge 1-t\ge \ilambda^2$ and ${\cal A}$ satisfies that 
$$
|{\cal A}_{\ba}|\le T_{s,D}(|a_1-a_2|),\quad \forall \ba=(a_1,a_2)\in(\Zn)^2,
$$
for some constant $D>0$. Then, we have  
\begin{align}
\label{neiwuj} 
\left({\cal U}^{(2)}_{s,t,\bsig} \circ 
{\cal A}\right)_{\ba} & \lesssim
 {T}_{t,D}( |a_1-a_2|)+(|1-s|/|1-t|)^2\cdot W^{-D} . % \quad \text{if}\ \ |a_1-a_2|\gtrsim (\log W)^{3/2} .
\end{align}
 \end{lemma}
\begin{proof}
The estimate \eqref{neiwuj} can be proved easily with \eqref{sum_res_Ndecay}, the estimate \eqref{prop:ThfadC} on $\Theta$-propagators, and the following basic calculus fact:
$\max_{a\in \mathbb R^d} \int_{x\in \mathbb R^d} \exp\big(-\sqrt{|a-x|}-\sqrt{|x|}+ \sqrt{|a|} \big)\dd x\lesssim 1.$  
We omit the details. % of the proof.
\end{proof}

Now, for an arbitrarily small constant $\e>0$, we define the stopping time
\begin{align}\label{eq:def_TTT}
T:=\inf \Big\{u\ge s:  \max_{\ba=(a,b)} \max_{\bsig}{\big|\left({\cal L-\cal K}\right)^{(2)}_{u,\bsig, \ba}\big|}\big/{ T}_{u,D}\p{|a-b|} \ge W^\e \Big\} .
\end{align}
By the definition \eqref{eq:def_TTT}, for every $u\le T$ we have \smash{$\max_{\ba,\bsig}|({\cal L}-{\cal K})^{(2)}_{u,\bsig,\ba}|\le W^{\e}\,T_{u,D}(|a-b|)$}, so all the estimates below may be applied to the stopped process $u\mapsto ({\cal L}-{\cal K})^{(2)}_{u\wedge T}$ and to the stopped hierarchy \eqref{int_K-L_ST} in place of \eqref{int_K-LcalE_n=2}. 
Using \Cref{lem_dec_calE,TailtoTail}, we can establish the following lemma by analyzing the equation \eqref{int_K-L_ST}. 

\begin{lemma}\label{lem:pf_step5} 
In the setting of \Cref{lem:main_ind}, suppose $1-t\ge \ilambda^2$ and the estimates \eqref{Gt_bound_flow}--\eqref{Eq:L-KGt-flow} hold. Then, for sufficiently small constant $\e>0$, we have that $T\ge t$ with high probability. 
\end{lemma}  
\begin{proof}
The proof of this lemma is analogous to, and in fact much simpler than, the proof of equation (2.76) in \cite[Section 5.3]{YY_25}, by using the estimates \eqref{res_deccalE_lk}--\eqref{res_deccalE_dif} and the evolution kernel estimate \eqref{neiwuj}. Hence, we omit the details of the proof. 
\end{proof}

By \Cref{lem:pf_step5}, we immediately obtain \eqref{Eq:Gdecay+s<g_flow}, since $\e>0$ can be taken arbitrarily small. Moreover, this estimate is stronger than \eqref{Eq:Gdecay_flow}. This completes Step 5 in the case $1-s\ge 1-t\ge \ilambda^2$. 


\section{Step 6: Expected 2-loop estimates}\label{Sec:Step6}

Our proof of Step 6 relies on analyzing the expectation of the loop hierarchy in \eqref{int_K-LcalE_n=2}: 
\begin{align} \label{Eexpint_K-L}
\E(\mathcal{L} - \mathcal{K})^{(2)}_{t, \bsig, \ba}=
    \left(\mathcal{U}^{(2)}_{s, t, \bsig} \circ \E(\mathcal{L} - \mathcal{K})^{(2)}_{s, \bsig}\right)_{\ba}+ \int_{s}^t \left(\mathcal{U}^{(2)}_{u, t, \bsig} \circ \E\mathcal{E}^{(\mathcal{L} - \mathcal{K}) \times (\mathcal{L} - \mathcal{K})}_{u, \bsig} + \mathcal{U}^{(2)}_{u, t, \bsig} \circ \E\mathcal{E}^{\Gc}_{u, \bsig}\right)_{\ba} \dd u ,  %+ \int_{s}^t \left(\mathcal{U}^{(2)}_{u, t, \bsig} \circ \E\mathcal{E}^{\Gc}_{u, \bsig}\right)_{\ba} \dd u  , 
\end{align}
where we omit the superscript ``$(2)$'' from some of the notation for brevity. 
In addition, we use the following expected averaged local law, which improves the (sharp) averaged local law \eqref{Gt_avgbound_flow} from Step 2 by an extra factor of $W^{-d}B_{u,0}$ arising from the expectation.
 

\begin{lemma}\label{lem:improve_exp_aver}
In the setting of \Cref{lem:main_ind}, suppose the estimates \eqref{Gt_avgbound_flow} and \eqref{Eq:L-KGt-flow} hold uniformly in $u\in[s,t]$. Then, we have 
\begin{align}\label{res_ELK_n=1} 
  \max_{a}  \left|\E\tr\p{ (G_u-M) E_{a}}\right|
  \prec (W^{-d}B_{u,0})^{2},\quad \forall\, s\le u \le t.
  \end{align} 
\end{lemma}
\begin{proof}
    The proof of this lemma is the same as that of \cite[Lemma 5.15]{YY_25} by using \eqref{Gt_avgbound_flow} and \eqref{Eq:L-KGt-flow}. %Hence, we omit the details.
\end{proof}

% The first ingredient is the following improved estimate on loops of length $1$, which is Lemma 5.15 in \cite{YY_25}.

% \begin{lemma}\label{lemma:step6-1}
% Fix any $a\in\Z_{L}^{2}$, any $\tau,\kappa>0$, any $E\in[-2+\kappa,2-\kappa]$, and $0\leq u\leq 1-N^{-1+\tau}$. We have 
% \begin{align}
% |\E{\cal L}_{u,(+),(a)}-{\cal K}_{u,(+),(a)}|=|\E\langle(G_{u}-m)E_{a}\rangle|\prec (W^{-d}B_{0,t})^{2}.
% \end{align}
% \end{lemma}
% \begin{proof}The proof is exactly the same as $d=1$ (and used for $d=2$. The key input is $$\langle(G_{u}-m)E_{b}\rangle\prec (W^{-d}B_{0,u})$$\end{proof}



%\begin{proof}[Proof of Lemma \ref{ML:exp}]
% Taking expectation of both side of  \eqref{int_K-LcalE} for the case $n=2$, we have  
% \begin{align}\label{Eexpint_K-L}
%  \E\;    (\mathcal{L} - \mathcal{K})_{t, \bsig, \textbf{a}} \;=
%  \; \; &
%     \left(\mathcal{U}_{s, t, \bsig} \circ  \E\, (\mathcal{L} - \mathcal{K})_{s, \bsig}\right)_{\textbf{a}} \; + \;  \int_{s}^t \left(\mathcal{U}_{u, t, \bsig}  \circ \E\,\mathcal{E}^{((\mathcal{L} - \mathcal{K}) \times (\mathcal{L} - \mathcal{K}))}_{u, \bsig}\right)_{\textbf{a}} du \\
%     &+   \int_{s}^t \left(\mathcal{U}_{u, t, \bsig} \circ  \E\, \mathcal{E}^{(\widetilde{G})}_{u, \bsig}\right)_{\textbf{a}} du.
% \end{align} 




 


% By assumption, we have 
% \begin{align}\label{jdasfuao01}
%   \E\;  (\mathcal{L} - \mathcal{K})_{s, \bsig,\textbf{a}}\prec \begin{cases}
%       W^{-2d-d/5},\quad & \eta_s\ge L^{-d}\\
%       (N\eta_s)^{-2}\cdot\left(W^{-d/5}+(N\eta_s)^{-1}\right),\quad & \eta_s\le L^{-d}
%   \end{cases}
% \end{align}

For the first term on the RHS of \eqref{Eexpint_K-L}, $\E(\mathcal{L} - \mathcal{K})^{(2)}_{s, \bsig,\ba}$ satisfies \eqref{Eq:Gtlp_exp+IND}. For the $\E\mathcal{E}^{(\mathcal{L} - \mathcal{K}) \times (\mathcal{L} - \mathcal{K})}$ term, combining the pointwise bound \eqref{Eq:Gdecay_flow} from Step 5 with the maximum bound \eqref{Eq:L-KGt-flow} from Step 4 yields, for $1-u \ge \ilambda^2/L^{d}$ and any $\ba=(a_1,a_2)\in (\Zn)^2$,
\begin{align}
\E{\cal E}^{(\cL-\cK)\times(\cL-\cK)}_{u,\bsig,\ba}&\prec 
(W^{-d}B_{u,0})^{\frac{11}{5}} \sum_{b} B_{u,|a_1-b|}e^{-(|a_1-b|/\ell_u)^{1/2}}\lesssim (1-u)^{-1}(W^{-d}B_{u,0})^{\frac{11}{5}},\label{eq:Exp(L-K)1}
\end{align}
while for $1-u\le \ilambda^2/L^{d}$, the maximum bound \eqref{Eq:L-KGt-flow} gives 
\begin{align}
\E{\cal E}^{(\cL-\cK)\times(\cL-\cK)}_{u,\bsig,\ba}&\prec W^d\cdot L^d(N|1-u|)^{-4} =(1-u)^{-1} (N|1-u|)^{-3}.\label{eq:Exp(L-K)2}
\end{align}
Finally, consider the light-weight term in \eqref{Eexpint_K-L}. By the definition \eqref{def_EwtG}, we may write 
\be\label{eq:LW-n=2}
\mathcal{E}^{\Gc}_{u, \bsig, \ba}= W^d \sum_{k=1}^2 \sum_{\al, \beta\in \Zn}  
 \tr\big ( \Gc_u(\sigma_k)  E_{\al} \big)
   S^{(\sB)}_{\al\beta} 
  \left( {\cut}^{(\beta)}_{k} \circ {\cal L}^{(2)}_{u, \bsig, \ba} \right) . \ee
When $1-u\le \ilambda^2/L^{d}$, its expectation can be bounded as
\begin{align}
\E\mathcal{E}^{\Gc}_{u, \bsig, \ba}&= W^d \sum_{k=1}^2 \sum_{\al, \beta}  
\br{\E \tr\big( \Gc_u(\sigma_k) E_{\al} \big) S^{(\sB)}_{\al\beta} \left( {\cut}^{(\beta)}_{k} \circ {\cal K}^{(2)}_{u, \bsig, \ba} \right) + \E \tr \big ( \Gc_u(\sigma_k)  E_{\al} \big) S^{(\sB)}_{\al\beta} \left( {\cut}^{(\beta)}_{k} \circ \p{\cal L-\cK}^{(2)}_{u, \bsig, \ba} \right)} \nonumber\\
&\prec W^d\cdot L^d(N|1-u|)^{-4} =(1-u)^{-1} (N|1-u|)^{-3}, \label{eq:ExpLWn=2_smalleta}
\end{align}
where the second step uses \eqref{res_ELK_n=1} and \eqref{eq:bcal_k} (with $n=3$) to bound the first term, and \eqref{Gt_avgbound_flow} together with \eqref{Eq:L-KGt-flow} (with $n=3$) to bound the second. On the other hand, when $1-u\ge \ilambda^2/L^{d}$, we control \smash{$\E\mathcal{E}^{\Gc}_{u, \bsig, \ba}$} as in the following lemma, which may be viewed as an analogue of \Cref{lem:LWterm} in expectation. Its proof uses the same diagrammatic techniques as in the proof of \Cref{lem:LWterm}, albeit in a considerably simpler form, and is deferred to \Cref{subsec:pf-LWterm_EXP}.



\begin{lemma}[Expected light-weight estimate]\label{lem:LWterm_EXP}
In the setting of \Cref{lem:main_ind}, suppose $1-t\ge \ilambda^2/L^{d}$ and the estimates \eqref{Gt_bound_flow}--\eqref{Eq:Gdecay_flow} hold. Then, for any $\bsig\in\{+,-\}^2$ and $\ba\in(\Zn)^2$, we have  
\begin{align}\label{eq:ExpLWn=2}
\E\mathcal{E}^{\Gc,(2)}_{t, \bsig, \ba}\prec (1-t)^{-1} (W^{-d}B_{t,0})^{\frac{5}{2}}. 
\end{align}
\end{lemma}
With the above estimates in place, we can now complete the proof of \eqref{Eq:Gtlp_exp_flow} using the tools developed in Step 3 (cf.~\Cref{Sec:Steps34}).
Here, we can obtain an additional small factor $(\ilambda^2W^d)^{-1/5}+W^{-d}B_{u,0}$ compared to \eqref{Eq:L-KGt-flow}, because the leading error in the integrated loop hierarchy \eqref{int_K-LcalE_n=2} arises from the martingale term, which vanishes after taking the expectation. 


\begin{proof}[\bf Step 6: Proof of \eqref{Eq:Gtlp_exp_flow}]
In the case $1-s\ge 1-t\ge \ilambda^2$, we have $B_{u,0}\asymp |1-u|^{-1}$, while in the case $1-t\le 1-s\le \ilambda^2/L^d$, we have $B_{u,0}\asymp (L^d|1-u|)^{-1}$. In these regimes, substituting \eqref{Eq:Gtlp_exp+IND} and \eqref{eq:Exp(L-K)1}--\eqref{eq:ExpLWn=2} into \eqref{Eexpint_K-L}, applying the evolution kernel estimate \eqref{sum_res_Ndecay} with $n=2$, and integrating over $u$, we obtain  
\be\label{eq:ExpL-Keasy} \E(\mathcal{L} - \mathcal{K})^{(2)}_{t, \bsig, \ba} \prec (W^{-d}B_{t,0})^{2}\p{(\ilambda^2W^d)^{-1/5}+W^{-d}B_{t,0}}
.\ee
It remains to deal with the regimes (i) $\ilambda^2/L^{2}\le 1- t\le 1-s \le \ilambda^2$, and (ii) $\ilambda^2/L^{d}\le 1-t\le 1-s\le \ilambda^2/L^{2}$. In these regimes, when $\sig_1=\sig_2$, substituting \eqref{Eq:Gtlp_exp+IND}, \eqref{eq:Exp(L-K)1}, and \eqref{eq:ExpLWn=2} into \eqref{Eexpint_K-L}, applying the evolution kernel estimate \eqref{sum_res_2_NAL} in regime (i) and the estimate \eqref{sum_res_Ndecay_nonzero} with $A=\emptyset$ (note that $I_{\mathrm{diff}}(\bsig)=\emptyset$) in regime (ii), and integrating over $u$, we obtain \eqref{eq:ExpL-Keasy}. (In regime (ii), \eqref{sum_res_2_NAL} is not available since \Cref{lem:sum_decay} requires $t\le 1-\ilambda^2/L^2$, and the factor $\int_s^t(1-u)^{-1}\dd u\le (d-2)\log L$ is absorbed by $\prec$.) 
% \[ \E(\mathcal{L} - \mathcal{K})^{(2)}_{t, \bsig, \ba} \prec (W^{-d}B_{t,0})^{2}\p{(\ilambda^2W^d)^{-1/5}+W^{-d}B_{t,0}}.\]
The remaining and more delicate case is when $\sig_1\ne \sig_2$. To handle this, we again employ the sum-zero operator from \Cref{Def:QtPt} in regime (i), and the zero-mode-removing operator from \Cref{def;zero_mode_remove} in regime (ii).

\medskip

\noindent
\textbf{Regime (i)}: Using Ward's identities \eqref{WI_calL} and \eqref{WI_calK}, along with \eqref{res_ELK_n=1} and \eqref{eq:derv_Theta}, we obtain 
\be\label{eq:EPL-K}
\br{\cal P\circ \E (\mathcal{L} - \mathcal{K})^{(2)}_{t, \bsig}}_{a_1}\dthn_{t,\ba}^{(2)}= \frac{\im \E \tr\p{(G_t -M)E_{a_1}}}{W^d\eta_t}\dthn_{t,\ba}^{(2)} \prec (W^{-d}B_{t,0})^{3}.
\ee
To estimate $\cal Q_t\circ \E (\mathcal{L} - \mathcal{K})^{(2)}_{t, \bsig, \ba}$, we take the expectation of equation \eqref{int_K-L+Q} with $n=2$, yielding
\begin{align}\label{int_K-L+QE}
& {\cal Q}_t\circ  \E(\mathcal{L} - \mathcal{K})^{(2)}_{t, \bsig, \ba}   = 
\left(\mathcal{U}^{(2)}_{s, t, \bsig} \circ  {\cal Q}_s\circ \E(\mathcal{L} - \mathcal{K})^{(2)}_{s, \bsig}\right)_{\ba} + \int_{s}^t \left(\mathcal{U}^{(2)}_{u, t, \bsig} \circ  {\cal Q}_u\circ \E \mathcal{E}^{(\cL-\cK)\times(\cL-\cK),(2)}_{u, \bsig}\right)_{\ba} \dd u \nonumber\\
& + \int_{s}^t \left(\mathcal{U}^{(2)}_{u, t, \bsig} \circ  {\cal Q}_u\circ \E\mathcal{E}^{\Gc,(2)}_{u, \bsig}\right)_{\ba} \dd u  + \int_{s}^t 
\left(\mathcal{U}^{(2)}_{u, t, \bsig} 
\circ \left( \left[{\cal Q}_u , \varTheta^{(2)}_{u,\bsig} \right]\circ \E(\mathcal{L} - \mathcal{K})^{(2)}_{u, \bsig }\right) \right)_{\ba} \dd u \nonumber\\
& - \int_{s}^t \left(\mathcal{U}^{(2)}_{u, t, \bsig} \circ \left\{\left[ {\cal P} \circ \E\left(\mathcal{L} - \mathcal{K}\right)^{(2)}_{u, \bsig}\right] \partial_u\dthn_{u}^{(2)} \right\}\right)_{\ba} \dd u.
\end{align}
We now estimate the RHS of \eqref{int_K-L+QE}. For the first three terms, we apply \Cref{lem_+Q}, together with the estimates \eqref{Eq:Gtlp_exp+IND}, \eqref{eq:Exp(L-K)1}, and \eqref{eq:ExpLWn=2}, and use the improved evolution kernel estimate \eqref{sum_res_2}, which exploits the sum-zero and fast-decay properties. Integrating over $u$, we can bound them by $(W^{-d}B_{t,0})^{2}[(\ilambda^2W^d)^{-1/5}+W^{-d}B_{t,0}]$.


It remains to control the last two terms in \eqref{int_K-L+QE}. Arguing as in \eqref{eq:EPL-K}, and again using Ward's identity, \eqref{res_ELK_n=1}, and \eqref{eq:derv_Theta}, we find
\be\label{eq:boundELKQ1}
\left\|\left[ {\cal P} \circ \E\left(\mathcal{L} - \mathcal{K}\right)^{(2)}_{u, \bsig}\right] \partial_u\dthn_{u}^{(2)}\right\|_\infty \prec (1-u)^{-1} (W^{-d}B_{u,0})^{3}. 
\ee 
For the commutator term, using the same argument as in \eqref{A4} (with $n=2$), we obtain
%using the definition \eqref{eq:sumzero_op} and the definition of $\varTheta^{(2)}_{u,\bsig}$ in \eqref{def:op_thn},  we obtain
\begin{align}
\big[{\cal Q}_u , \varTheta^{(2)}_{u,\bsig} \big]\circ \E(\mathcal{L} - \mathcal{K})^{(2)}_{u, \bsig, \ba} %&= \varTheta^{(2)}_{u, \bsig} \circ \left[ \left({\cal P} \circ \E\left( \mathcal{L} - \mathcal{K} \right)^{(2)}_{u,\bsig}\right) \cdot \dthn^{(2)}_u \right]_{\ba} - \left[\left( {\cal P} \circ \varTheta^{(2)}_{u, \bsig} \circ \E(\mathcal{L} - \mathcal{K})^{(2)}_{u,\bsig}\right)  \cdot \dthn^{(2)}_u \right]_{\ba} \nonumber\\ 
&\prec (1-u)^{-1}\left\|\dthn^{(2)}_{u}\right\|_\infty \cdot   \left\| {\cal P} \circ \E\left(  \mathcal{L} - \mathcal{K} \right)^{(2)}_{u,\bsig} \right\|_{\infty}\prec (1-u)^{-1} (W^{-d}B_{u,0})^{3},\label{eq:boundcommutator}
\end{align} 
%where in the second step, we use the simple fact that $ \|\varTheta_{u,\bsig}^{(2)}\|_{\infty\to\infty} \le (1-u)^{-1}$, and 
where the second step again follows from Ward's identity and the estimates \eqref{res_ELK_n=1} and \eqref{eq:derv_Theta}. 
Applying \Cref{lem_+Q}, along with the bounds \eqref{eq:boundELKQ1} and \eqref{eq:boundcommutator}, and using the evolution kernel estimate \eqref{sum_res_2} (the fast-decay property \eqref{deccA0} of these two terms follows exactly as for ${\cal B}_5$ and ${\cal B}_4$ in the proof of \Cref{lem:STOeq_Qt}, with $\cL-\cK$ replaced by $\E(\cL-\cK)$), we conclude that the contribution of the last two terms in \eqref{int_K-L+QE} is bounded by $(W^{-d}B_{t,0})^3$. This completes the proof of \eqref{Eq:Gtlp_exp_flow} in regime (i).

\medskip

\noindent
\textbf{Regime (ii)}: We use the operator $Q^{(\{1,2\})}=Q^{(1)}\circ Q^{(2)}$ in \Cref{def;zero_mode_remove} and Ward's identity to express 
\begin{align}
 (\cal L-\cK)_{t,\bsig,\ba}^{(2)}=  Q^{(\{1,2\})} \circ (\cal L-\cK)_{t,\bsig,\ba}^{(2)}  + \frac{\im \br{\tr\p{(G_t-M)(E_{a_1}+E_{a_2})} -N^{-1}\tr\p{G_t-M}}}{N\eta_t}
\end{align}
for $\sig_1\ne \sig_2$. Taking expectations and using \eqref{res_ELK_n=1}, the expectation of the second term is bounded by $(W^{-d}B_{t,0})^2\cdot (N\eta_t)^{-1}\le (W^{-d}B_{t,0})^3$. To estimate the first term, we take the expectation of equation \eqref{iisuwjyys} with $n=2$ and $A=\{1,2\}$, yielding
\begin{align}\label{iisuwjyys_exp}
Q^{(\{1,2\})}\circ \E(\mathcal{L} - \mathcal{K})^{(2)}_{t, \bsig, \ba} 
&=\left(Q^{(\{1,2\})}\circ\mathcal{U}^{(2)}_{s, t, \bsig}  \circ \E(\mathcal{L} - \mathcal{K})^{(2)}_{s, \bsig}\right)_{\ba} + \int_{s}^t \left(Q^{(\{1,2\})}\circ \mathcal{U}^{(2)}_{u, t, \bsig}  \circ \E \mathcal{E}^{(\cL-\cK)\times(\cL-\cK),(2)}_{u, \bsig}\right)_{\ba} \dd u \nonumber\\
& + \int_{s}^t \left(Q^{(\{1,2\})}\circ \mathcal{U}^{(2)}_{u, t, \bsig} \circ \E\mathcal{E}^{\Gc,(2)}_{u, \bsig}\right)_{\ba} \dd u .
\end{align}
Applying \eqref{normQA2}, together with the estimates \eqref{Eq:Gtlp_exp+IND}, \eqref{eq:Exp(L-K)1}, and \eqref{eq:ExpLWn=2}, and using the evolution kernel estimate \eqref{sum_res_Ndecay_nonzero}, we integrate over $u$ to obtain the desired bound: $(W^{-d}B_{t,0})^{2}[(\ilambda^2W^d)^{-1/5}+W^{-d}B_{t,0}]$. This completes the proof of \eqref{Eq:Gtlp_exp_flow} in regime (ii).
 % \begin{itemize}
 %     \item Sum zero operator for $\eta\ge L^{-2}$
 %     \item Zero mode operator for $L^{-d}\le \eta\le L^{-2}$
 % \end{itemize}  
\end{proof}

%  On the other hand, for $\eta_u\le L^{-d}$, use the old argument as in $d=1$ proof: 

% $\begin{aligned} \E \mathcal{E}_{u, \bsig, \mathbf{a}}^{(\widetilde{G})} & \prec N \cdot \max _a \max _{\mathbf{a}} \max _{\bsig \in\{+,-\}^3}\left|\E\left[\left\langle\left(G_u-m\right) E_a\right\rangle \cdot \mathcal{L}_{u, \mathbf{a}, \bsig}\right]\right| \\ & \prec N\cdot \max _a\left|\E\left\langle\left(G_u-m\right) E_a\right\rangle\right| \cdot \max _{\mathbf{a}} \max _{\bsig \in\{+,-\}^3}\left|\mathcal{K}_{u, \mathbf{a}, \bsig}\right| \\ & +N \cdot \max _a\left|(\mathcal{L}-\mathcal{K})_{u,(+),(a),}\right\rangle\left|\cdot \max _{\mathbf{a}} \max _{\bsig \in\{+,-\}^3}\right|(\mathcal{L}-\mathcal{K})_{u, \mathbf{a}, \bsig} \mid .
% \\\nonumber
% & \;\prec \; N (N\eta_u)^{-4}\prec \eta_u^{-1}\cdot(N\eta)^{-3}  \end{aligned}$

 








 

\section{Estimation of the light-weight term}\label{Sec:graph}

This section is devoted to proving the key estimates for the light-weight term $\cal E^{\Gc,(2)}$---namely, \Cref{lem:LWterm,lem: EWGn2_N} in Step 2. 
The proofs build on several crucial ideas and refined graphical tools developed in earlier works \cite{yang2021delocalization,yang2022delocalization,yang2021random}. 
We present the full details in the setting of the random band matrix model, while the modifications required for the block Anderson model are described separately in \Cref{subsec:LWchange-to-BA} below. 
%\subsection{Proofs of Lemma \ref{lem:LWterm} and Lemma \ref{lem: EWGn2_N}} 
Recall the light-weight term takes the form \eqref{eq:LW-n=2}. For $\bsig=(\sig',\sig)$ and $\ba=(a,b)$, we have 
\begin{align}
\mathcal{E}^{\Gc}_{t, \bsig, \ba}&= {W}^d \sum_{a_1, a_2}   S^{(\sB)}_{a_1a_2} 
 \tr \big(\Gc_t(\sig) E_{a_1}\big)\tr\big( G_t(\sig) E_{a_2} G_t(\sig) E_{b} G_t(\sig') E_{a}\big) +\p{(\sig,a)\leftrightarrow (\sig',b)} \label{eq:EGC}\\
 &= W^{-2d} \sum_{x\in [a],y\in [b]} (G_t(\sig'))_{yx} \sum_{a_1, a_2}S^{(\sB)}_{a_1a_2}\sum_{\al \in [a_2]}     
 \tr\big( \Gc_t(\sig) E_{a_1} \big) (G_t(\sig))_{x\al }(G_t(\sig))_{\al y} +\p{(\sig,a)\leftrightarrow (\sig',b)},\nonumber
\end{align} 
where $\p{(\sig,a)\leftrightarrow (\sig',b)}$ denotes the term obtained by exchanging $(\sig,a)$ and $(\sig',b)$ in the preceding term.
%Without loss of generality, we only control the term with $[d]=[d']$, while the cases with $[d]\sim [d']$ can be handled in exactly the same way. 
In the setting of \Cref{lem:LWterm}, using \eqref{GiiGEX} and \eqref{GijGEX}, we can bound $(G_t(\sig'))_{yx}$ as 
\be\label{eq:directG1}
|(G_t(\sig'))_{yx}|\prec  \delta_{xy}+\Psi_t(|a-b|).
\ee 
Note that when $1-t\le \ilambda^2/L^{2}$, we have $\ell_t=L$, so the exponential factor in the definition of $\cal T_t$ in \eqref{defTUL} is always of order 1 for $r\lesssim L$. In this case, \Cref{lem: EWGn2_N} is an immediate consequence of \Cref{lem:LWterm}. Hence, for the proof of \Cref{lem: EWGn2_N}, we only need to focus on the $1-t> \ilambda^2/L^{2}$ case, where  
\[ W^{-d}\wT_{t,D}^{\ell}(|a-b|)\asymp \br{\sT_{t}(|a-b|\wedge \ell)}^2 + W^{-d-D}.\]
Here, for simplicity of notation, we introduce the function $\sT_{t}:[0,\infty)\to [0,\infty)$ defined as: 
\[\sT_{t}(r):=\frac{(\ilambda^2+|1-t|)^{-1/2}}{W^{ d /2}(r+1)^{(d-2)/2}}\exp\p{-\frac{1}{2}\sqrt{{r}/{\ell_t}}},\quad \forall r\ge 0. \]
Then, in the setting of \Cref{lem: EWGn2_N} and for $1-t> \ilambda^2/L^{2}$, we have the following estimate by \eqref{GiiGEX} and \eqref{GijGEX}:
\be\label{eq:directG2}
|(G_t(\sig'))_{yx}|\prec \delta_{xy}+\sT_{t}(|a-b|\wedge \ell) + W^{-D} .
\ee 
%for any fixed $D>0$. 
%With \eqref{eq:directG1} and \eqref{eq:directG2} in hand, to prove \Cref{lem:LWterm,lem: EWGn2_N}, we need to bound the expression 
With \eqref{eq:directG1} and \eqref{eq:directG2} in place, the proofs of \Cref{lem:LWterm,lem: EWGn2_N} reduce to estimating
\begin{align}\label{fxyG_sum}
f_{xy}(G)&=W^{-d}\sum_{a_1, a_2}S^{(\sB)}_{a_1a_2}\sum_{{\al\in [a_2],\beta \in [a_1]}}\mathbf 1_{\al\notin\{x,y\}}
(\Gc_t)_{\beta\beta} (G_t)_{x\al }(G_t)_{\al y} \equiv \sum_{\al\notin\{x,y\}}\sum_{\beta}S_{\al\beta}
\Gc_{\beta\beta} G_{x\al}G_{\al y},
%\wt f_{xy}(G)&=\sum_{\al\in\{x,y\}}\sum_{\beta}S_{\al\beta}\Gc_{\beta\beta} G_{x\al}G_{\al y},\label{fxyG_sum2}
\end{align} 
where, without loss of generality, we take $\sig = +$ and abbreviate $G \equiv G_t$. 
We denote by \smash{$\wt f_{xy}(G)$} the corresponding contribution with $\al\in\{x,y\}$. 
%Here, we focus on the summation over $\al \notin \{x,y\}$. For the case $\al \in \{x,y\}$, 
Using \eqref{eq:directG1}, \eqref{eq:directG2}, and the averaged local law \eqref{GavLGEX}, we get 
\be\label{eq:recolterm} 
\wt f_{xy}(G) \prec \Psi^2_t(0)\cdot \p{\delta_{xy}+\Psi_t(|a-b|)}
\ee
in the setting of \Cref{lem:LWterm}, and 
\be\label{eq:recolterm2} 
\wt f_{xy}(G) \prec (W^{-d}B_{t,0}) \cdot \p{\delta_{xy} + \sT_{t}(|a-b|\wedge \ell)+ W^{-D}}
\ee
in the setting of \Cref{lem: EWGn2_N}.
%the same expression with the RHS replaced by \((W^{-d}B_{t,0}) \cdot (\delta_{xy} + \sT_{t}(|a-b|\wedge \ell)+ W^{-D}) \). 
Combining \eqref{eq:recolterm} with \eqref{eq:directG1}, and averaging over $x\in[a]$ and $y\in [b]$, yields  
\begin{align}\label{eq:recoltermwt}
W^{-2d} \sum_{x\in [a],y\in [b]} (G_t(\sig'))_{yx} \wt f_{xy}(G) \prec \Psi^2_t(0)\left[W^{-d}\delta_{ab}+\Psi_t^2(|a-b|)\right] \lesssim \Psi^2_t(0)\cdot \Psi_t^2(|a-b|), 
\end{align} 
which is bounded by the RHS of \eqref{eq:LW_conclusion}. Similarly, combining \eqref{eq:recolterm2} with \eqref{eq:directG2} gives
\begin{align}\label{eq:recoltermwt2}
W^{-2d} \sum_{x\in [a],y\in [b]} (G_t(\sig'))_{yx} \wt f_{xy}(G) \prec (W^{-d}B_{t,0}) \cdot W^{-d} \wT_{t,D}^{\ell}(|a-b|), 
\end{align}
which is bounded by the RHS of \eqref{eq:LW_conclusion_exp}. Therefore, it remains to estimate the main term $f_{xy}(G)$. 


We first claim the following $\max$-bound on $f_{xy}(G)$:
\begin{align}\label{eq:boundfxyGinf}
    |f_{xy}(G)|\prec \Psi_t^2 \sum_{\al}|G_{x\al}||G_{\al y}| \le \Psi_t^2 \p{{\im G_{xx}}/{\eta_t}}^{1/2}\p{{\im G_{yy}}/{\eta_t}}^{1/2}\prec \eta_t^{-1}\Psi_t^2,
\end{align}
where the first step uses the averaged local law \eqref{GavLGEX}, and the second applies Cauchy--Schwarz together with Ward's identity. In the setting of \Cref{lem:LWterm}, we take $\Psi_t = \Psi_t(0)$, whereas in the setting of \Cref{lem: EWGn2_N}, we take $\Psi_t = \max\{W^{-d/2},(W^{-d}B_{t,0})^{1/2}\}$. Since $B_{t,0}\ge (\ilambda^2+1)^{-1}\ge (\fd^{-2}+1)^{-1}$ by \eqref{eq_B_param} and \eqref{eq:WO}, this is comparable to $(W^{-d}B_{t,0})^{1/2}$; taking the maximum only ensures the lower bound $\Psi_t\ge W^{-d/2}$, which is required in \Cref{lem_GbEXP} and used below \eqref{eq:ordG}. For simplicity of notation, we continue to write $\Psi_t=(W^{-d}B_{t,0})^{1/2}$ in this setting.
Taking \Cref{lem:LWterm} as an example, our goal is to establish the sharper bound $|f_{xy}(G)|\prec \eta_t^{-1}\Psi_t(0)\Psi_t(|a-b|)$.
Thus, the estimate \eqref{eq:boundfxyGinf} lacks the long-edge factor $\Psi_t(|a-b|)$, which captures the spatial decay in $|a-b|$. This missing factor is a key technical obstacle---previously noted around \eqref{eq:introGc} in the introduction---because all tools developed within the stochastic flow and loop hierarchy frameworks fail to extract it. 
We have to exploit a \emph{global fluctuation averaging} mechanism in the sum over $\al$ in \eqref{fxyG_sum}, which can yield additional powers of $L^{-d/2}$ or \smash{$\ell_t^{-d/2}$} to compensate for the missing long edge.
To realize this mechanism, we estimate the high moments of $f_{xy}(G)$, adapting ideas from \cite{yang2021delocalization,yang2022delocalization,yang2021random}.
%Thus, the bound \eqref{eq:boundfxyGinf} lacks a long edge that contributes the factor $\Psi_t(|a-b|)$; this issue has also been discussed around \eqref{eq:introGc} in the introduction. Extracting this missing long edge is one of the central technical difficulties of the proof, where all tools developed within the framework of stochastic flow and loop hierarchy break down. Our only hope is to exploit a \emph{global fluctuation averaging mechanism} with respect to the global sum over the vertex $\al$ in $f_{xy}(G)$, which would yield additional factors of $L^{-1}$ (or \smash{$\ell_t^{-1}$}) that may compensate for the missing long edge.
By Markov's inequality, \Cref{lem:LWterm,lem: EWGn2_N} follow from the next estimates on the high moments of $f_{xy}(G)$.
\begin{lemma}\label{lem:LW_moment}
In the setting of \Cref{lem:LWterm}, for any fixed $p\in 2\N$, there exists a constant $c>0$ depending on $p$ such that the following estimate holds:
     \be\label{eq:LW_moment}
\E\big|f_{xy}(G)\big|^{p}\prec  \eta_t^{-p}[\Psi_t(0)]^p \cdot \br{\Psi_t\p{c|a-b|}}^p . %\quad \forall \ \bsig\in \{+,-\}^2, \ [a],[b]\in \Zn.
\ee
\end{lemma}
\begin{lemma}\label{lem:LW_moment_exp}
In the setting of \Cref{lem: EWGn2_N}, assume that $1-t>\ilambda^2/L^{2}$. Then, for each fixed $p\in 2\N$, the following estimate holds for any large constant $D>0$:
     \be\label{eq:LW_moment_exp}
\E\big|f_{xy}(G)\big|^{p}\prec \eta_t^{-p} (W^{-d}B_{t,0})^{p/2} \cdot \br{\sT_{t}(|a-b|\wedge \ell) }^{p}+W^{-D}. %\quad \forall \ \bsig\in \{+,-\}^2, \ [a],[b]\in \Zn.
\ee
\end{lemma}
\begin{proof}[\bf Proof of \Cref{lem:LWterm,lem: EWGn2_N}]
Applying Markov's inequality to \eqref{eq:LW_moment} (with arbitrarily large $p$) yields 
\[ f_{xy}(G)\prec \eta_t^{-1} \Psi_t(0) \cdot \Psi_t\p{c|a-b|} \prec \eta_t^{-1} \Psi_t(0) \cdot \Psi_t\p{ |a-b|},\]
where the second step uses the condition \eqref{eq:Psi}. Combining this bound with \eqref{eq:directG1} and \eqref{eq:recoltermwt}, we obtain \eqref{eq:LW_conclusion}.  
For the proof of \Cref{lem: EWGn2_N}, as discussed below \eqref{eq:directG1}, it suffices to assume that $1-t>\ilambda^2/L^{2}$. In this case, applying Markov's inequality to \eqref{eq:LW_moment_exp} yields, for any large constant $D>0$,
\[ f_{xy}(G)\prec \eta_t^{-1} (W^{-d}B_{t,0})^{1/2} \cdot \br{\sT_{t}(|a-b|\wedge \ell) } + W^{-D}.\]
Together with \eqref{eq:directG2} and \eqref{eq:recoltermwt2}, it gives \eqref{eq:LW_conclusion_exp}.
\end{proof}

%The rest of this section is devoted to the proofs of \Cref{lem:LW_moment,lem:LW_moment_exp}. In particular, we will first provide the proof of \Cref{lem:LW_moment}, while the proof of \Cref{lem:LW_moment_exp} is based on the argument developed during the proof of \Cref{lem:LW_moment}, along with some new arguments. Note that we have the following $\max$-bound on $f_{xy}(G)$:

The remainder of this section is devoted to proving \Cref{lem:LW_moment,lem:LW_moment_exp}. We begin with the proof of \Cref{lem:LW_moment}; the proof of \Cref{lem:LW_moment_exp} will follow the same general strategy, with several additional technical ingredients. Note that when $|a-b| \le (\log W)^{3/2}$, both \eqref{eq:LW_moment} and \eqref{eq:LW_moment_exp} follow directly from \eqref{eq:boundfxyGinf}. Thus, in the following proof, we restrict attention to the case
\be\label{eq:far_ab}
|a-b| > (\log W)^{3/2}.
\ee
%We now discuss the proof of the light-weight term estimate in \eqref{eq:smallLWMG}. As discussed around \eqref{eq:introGc}, the light-weight term lacks a long edge, even if we assume the sharp bound \eqref{eq:aprioriG2} on off-diagonal resolvent entries. Extracting this missing long edge is one of the central technical difficulties of the proof, where all tools developed within the framework of stochastic flow and loop hierarchy break down. Our only hope is to exploit a \emph{global fluctuation averaging mechanism} with respect to the global sum over the vertex $v$ in the 3-loop of \eqref{eq;LWG2}. If the strongest CLT-type cancellation were available, we would gain an additional factor of $L^{-d/2}$ (or \smash{$\ell_t^{-d/2}$} with $\ell_t$ giving the effective range), thereby compensating for the missing long edge. (Of course, in the actual proof we do not achieve such a strong cancellation.) To realize this mechanism, we estimate the high moments of the light-weight term, adapting ideas from \cite{yang2021delocalization,yang2022delocalization,yang2021random}. 
Following the approach of \cite{yang2021delocalization,yang2022delocalization,yang2021random}, we represent resolvent expressions (such as $\big|f_{xy}(G)\big|^{p}$) as graphs, expand them using Gaussian integration by parts, and then bound the resulting diagrams according to their graphical structures.
In the next subsection, we introduce the basic graphical notations, including the concepts of vertices, edges, molecules, graph values, and scaling sizes, and define the basic graph expansions that will be used in the proof.



\subsection{Graphical tools and local expansions}\label{sec:graphs}

%Following \cite{yang2021delocalization,yang2022delocalization,Xu:2024aa,yang2021random}, we organize our proofs using a graphical framework. In this subsection, we introduce the basic graphical notations, including the concepts of vertices, edges, molecules, graph values, and scaling sizes, and define the basic graph expansions. 

Our graphs are composed of matrix indices as vertices, and various types of edges representing different matrix entries. Specifically, the graphs will encode entries from the following matrices: $I, \ S, \ G, \ \Gc=G-M$, and $S^{\pm}$, where $S^{\pm}$ are defined by 
\be\label{eq:def-Spm}
S^+_{xy}(z):=W^{-d}\big(S^{(\sB)}\Theta^{(+,+)}(z)\big)_{ab} \quad \text{for}\quad x\in[a], y\in [b],\quad \text{and}\quad S^-(z) :=(S^+(z))^*, % = S\frac{1}{1-m^2S},
\ee
with the notations $\Theta^{(+,+)}(z)$ and $S^{(\sB)}$ specified earlier in \eqref{def:Theta} and \eqref{eq:variancematrix}. 
The graphical notation is defined in a general setting for an arbitrary random band matrix $H$ and its associated variance matrix $S$. As a special case, these definitions apply to our matrix flow $H_t$, where the variance matrix is given by $S_t=tS$, and the associated resolvent is denoted by $G_t\equiv (H_t-z_t)^{-1}$.  

\begin{definition}[Graphs] \label{def_graph1} 
Given a graph with vertices and edges, we assign the following structures and call it a vertex-level graph.  
\begin{itemize}
\item {\bf Vertices:} Vertices represent matrix indices in our expressions. Every graph has some external or internal vertices: external vertices represent indices whose values are fixed, while internal vertices represent summation indices that will be summed over. 
				
\item{\bf  Solid edges and weights:} We will use $(x,y)$ to denote a solid edge from vertex $x$ to vertex $y$. Every solid edge represents a resolvent entry. More precisely: 
\begin{itemize}
\item A blue (resp.~red) oriented solid edge from $x$ to $y$ represents a $G_{xy}$ (resp.~$\bar G_{xy}$) factor. 
    
\item A blue (resp.~red) oriented solid edge with a circle ($\circ$) from $x$ to $y$ represents a $(G-M)_{xy}$ (resp.~$\overline {(G-M)}_{xy} $) factor.

\item A $G_{xx}$ (resp.~$\bar G_{xx}$) factor will be represented by a blue (resp.~red) self-loop on the vertex $x$, while a $(G-M)_{xx}$ (resp.~$\overline {(G-M)}_{xx}$) factor will be represented by a blue (resp.~red) self-loop on the vertex $x$ with a circle ($\circ$). Following the convention in \cite{yang2021delocalization}, we will also call $G_{xx}$ and $\bar G_{xx}$ as blue and red {\bf weights}, and call $(G-M)_{xx}$ and $\overline {(G-M)}_{xx}$ as blue and red {\bf light-weights}.

%\item a red oriented solid edge from $x$ to $y$ represents a $\overline G_{xy}$ factor;

%\item a red oriented solid edge with a circle ($\circ$) from $x$ to $y$ represents a $\Gc^*_{xy}$ factor.
\end{itemize}
We assign a $+$ charge to each blue solid edge and a $-$ charge to each red solid edge.		
\item {\bf Waved edges:}
\begin{itemize}
\item A black waved edge between $x$ and $y$ represents an $S_{xy}$ factor. 
				
\item A blue (resp.~red) waved edge between $x$ and $y$ represents an $S^+_{xy}$ (resp.~$S^-_{xy}$) factor.
\end{itemize}

\item {\bf Dotted edges:} A black dotted edge between $x$ and $y$ represents a factor $\mathbf 1_{x=y}$ and a $\times$-dotted edge represents a factor $\mathbf 1_{x\ne y}$. There is at most one dotted or $\times$-dotted edge between every pair of vertices.
  
		
\item{\bf Coefficient:} There is a coefficient associated with every graph. Every coefficient is of order $\OO(1)$ and is a polynomial in $m$, $\overline m$, $m^{-1}$, $\overline m^{-1}$, $(1-m^2)^{-1}$, and $(1-\overline m^2)^{-1}$. 

\end{itemize}	
\end{definition}



Edges between internal vertices are called \emph{internal edges}, while edges with at least one end at an external vertex are called \emph{external edges}. The orientations of non-solid edges do not matter because they all represent entries of symmetric matrices. %Furthermore, we will say that a black edge has a neutral (zero) charge, while blue and red edges have plus ($+$) and minus ($-$) charges. 
The dotted edges are introduced solely for organizing the proof in certain steps; otherwise, we will almost always identify vertices connected by dotted edges. 
% Next, we introduce the concept of subgraphs. 
% \begin{definition}[Subgraphs]\label{def_sub}
% A graph $\cal G_1$ is said to be a subgraph of $\cal G_2$, denoted by $\cal G_1\subset \cal G_2$, if every graphical component (except the coefficient) of $\cal G_1$ is also in $\cal G_2$. Moreover, $\cal G_1 $ is a proper subgraph of $\cal G_2$ if  $\cal G_1\subset \cal G_2$ and $\cal G_1\ne \cal G_2$. Given a subset $\cal S$ of vertices in a graph $\cal G$, the subgraph $\cal G|_{\cal S}$ induced on $\cal S$ refers to the subgraph of $\cal G$ consisting of vertices in $\cal S$ and the edges between them.  
% \end{definition}
To each graph, we assign a \emph{value} as follows. For simplicity, throughout this paper, we will always abuse notation by identifying a graph (a geometric object) with its value (an analytic expression). 

\begin{definition}[Values of graphs]\label{ValG} 
Given a graph $\mathcal G$, we define its value as follows. We first take the product of all edge factors together with the coefficient associated with the graph $\cal G$. We then sum over all internal indices corresponding to internal vertices, while keeping the external indices fixed at their prescribed values.
For a linear combination of graphs $\sum_i c_i \cal G_i$, we define its value in the natural way, as the linear combination of the values of the individual graphs $\cal G_i$.

%Then, we sum over all the internal indices represented by the internal vertices, while the external indices are fixed by their given values. The value of a linear combination of graphs $\sum_i c_i \cal G_i$ is naturally defined as the linear combination of the graph values of $\cal G_i$.
\end{definition}


As noted in the previous works \cite{yang2021delocalization,yang2022delocalization,Xu:2024aa,yang2021random}, our graphs have a two-level structure: some \emph{local structures} varying on scales of order $W$, called {\bf molecules}, which are the equivalence classes of vertices connected through dotted or waved edges, along with a \emph{global structure of blocks} varying on scales up to $L$. 
Given a graph defined in \Cref{def_graph1}, if we ignore the inner structure within each molecule, we will get a {\bf molecular graph} with vertices being molecules. The molecular graph will be very useful in organizing the graphs and understanding their global structures. We now give its formal definition.


\begin{definition}[Molecules and molecular graphs]\label{def_poly}
We partition the set of all vertices into a union of disjoint subsets called molecules. Two vertices belong to the same molecule if and only if they are connected by a path of dotted and waved edges. Every molecule containing at least one external vertex is called an external molecule; otherwise, it is an internal molecule. An edge is said to be inside a molecule if both of its ending vertices belong to this molecule. 
% Molecular graphs are graphs consisting of 
% 	\begin{itemize}
% 		\item external and internal molecules;
% 		\item solid,  $\dashed$ and free edges between molecules.
% 		%\item dotted edges between external molecules or an external molecule and an internal molecule.
% 	\end{itemize}
% \medskip\noindent (ii) 
Given a graph $\cal G$, we define its molecular graph, denoted by $ \cal G_{\cal M}$, as follows:
	\begin{itemize}
		\item merge all vertices in the same molecule and represent them by a vertex;
		
		\item keep all solid edges between molecules;

		\item discard all the other components in $\cal G$ (i.e., $\times$-dotted edges, edges inside molecules, and coefficients).
	\end{itemize}
\end{definition}

By the exponential decay of the waved edges---derived from the definitions of $S$ and the estimate \eqref{prop:ThfadC_short}---we see that up to an error of order $\exp(-c(\log W)^{3/2})$ for a constant $c>0$:  
\be\label{scalemole}
 \hbox{$x$, $y$ are in the same molecule} \implies  |x-y|\le  W(\log W)^{3/2}.
\ee
We remark that molecular graphs are used solely to help with the analysis of graph structures, while all graph expansions will be applied exclusively to vertex-level graphs. 

\begin{definition}[Normal graphs]  \label{defnlvl0}  
We say a graph $\cal G$ is \emph{normal} if it satisfies the following properties:
\begin{itemize}
\item[(i)] It contains at most $\OO(1)$ many vertices and edges. % and it may have some external vertices whose values are fixed. 

\item[(ii)] There are no dotted edges between vertices.

%\item[(ii)] all its internal vertices are connected together or to external vertices through paths of dotted or waved edges.

\item[(iii)] Every pair of vertices in the graph are connected by a $\times$-dotted edge \emph{if and only if} they are connected by a solid edge.

\item[(iv)] Every weight is a light-weight.
%Every solid edge in the graph has a $\circ$ (i.e., it represents an entry of $G-M$ or $(G-M)^*$). In particular, all weights (i.e., self-loops) are light-weights. 
\end{itemize}
\end{definition}

In a normal graph, every $G$ edge is off-diagonal, while all the diagonal $G$ factors are represented by weights. Given any graph with $\OO(1)$ vertices and edges, we can rewrite it as a linear combination of normal graphs via the following \emph{dotted edge partition} operation. 

\begin{definition}[Dotted edge partition] \label{dot-def}
Given a graph $\cal G$, if there is at least one $G$-edge but no $\times$-dotted edge between a pair of vertices $\al$ and $\beta$, we write
 $1=\mathbf 1_{\al=\beta} + \mathbf 1_{\al \ne \beta} .$
If there is a $\times$-dotted edge $\mathbf{1}_{\al \ne \beta}$ but no $G$-edge between $\al$ and $\beta$, we write 
 $\mathbf 1_{\al\ne \beta} =1 - \mathbf 1_{\al = \beta}.$
Expanding the product of all such identities, we can express $\cal G$ as
 \be\label{odot}\cal G := \sum {\Dot} \cdot \cal G,\ee
 where each ${\Dot}$ is a product of dotted and $\times$-dotted edges together with a sign $\pm$. If $\Dot$ is ``inconsistent''---that is, if two vertices are connected both by a $\times$-dotted edge and by a path of dotted edges---then $\Dot \cdot \cal G = 0$.  
 For each consistent ${\Dot}$, we merge vertices connected by dotted edges in ${\Dot} \cdot \cal G$. In particular, if there is a $\times$-dotted edge between $\al$ and $\beta$, then all $G$ edges between them are off-diagonal; otherwise, the $G$-edges between them become weights once $\al$ and $\beta$ are merged. Finally, in each resulting graph, every diagonal factor $G_{xx}$ (resp.~$\overline{G}_{xx}$) is decomposed into a light-weight $(G_{xx}-m)$ plus a coefficient $m$ (resp.~$(\overline{G}_{xx}-\overline m)$ plus $\overline m$). Taking the product over all such decompositions yields a linear combination of normal graphs.
 
% we decompose each $G_{xx}$ (resp.~$\overline G_{xx}$) weight as a light-weight $(G_{xx}-m)$ edge plus an $m$ coefficient (resp.~$(\overline G_{xx}-\overline m)$ edge plus an $\overline m$ coefficient) and take the product over all these decompositions. This gives us a linear combination of normal graphs.
  \end{definition}

Given a normal graph, we can define its scaling size as in \Cref{def scaling}. 

%Given an arbitrary graph with $\OO(1)$ many vertices and edges, we can write it as a linear combination of normal graphs by expanding each $G_{xy}$ (resp.~$\overline G_{xy}$) edge as a $(G_{xy}-m\delta_{xy})$ edge plus a $m\delta_{xy}$ edge (resp.~$(\overline G_{xy}-\overline m\delta_{xy})$ edge plus a $\overline m\delta_{xy}$ edge), and by merging vertices connected via dotted edges. 


\begin{definition}[Scaling size]\label{def scaling} 
We define the scaling size of a normal graph $\Gamma$ as follows:
 % \begin{align}
	% 	\size(\Gamma): =&~ N^{n_M(\Gamma)}\cdot (\Psi_t)^{n_S(\Gamma)} \cdot W^{-d\left(n_W(\Gamma)  - n_V(\Gamma) +n_M(\Gamma)\right)}.
 %  \label{eq_defsize}
	% \end{align}
 \begin{align}
		\size(\Gamma): =&~ (L^d)^{n_M(\Gamma)}\cdot (\Psi_t)^{n_S(\Gamma)} \cdot W^{-d\left(n_W(\Gamma)  - n_V(\Gamma)\right)}.
  \label{eq_defsize}
	\end{align}
 % \begin{align}
	% 	\size(\Gamma): =&~ N^{\#\{\text{internal molecules}\}}(\Psi_t)^{\#\{\text{solid edges}\}} \nonumber\\
 %        &~\times W^{-d\left(\#\{ \text{waved edges}\}  -  \#\{\text{internal vertices}\}+\#\{\text{internal molecules}\}\right)}.
 %  \label{eq_defsize}
	% \end{align} 
where $n_S(\Gamma)$, $n_W(\Gamma)$, $n_V(\Gamma)$, and $n_M(\Gamma)$ denote the numbers of solid edges (including light-weights), waved edges, internal vertices, and internal molecules, respectively.
If a graph $\Gamma$ can be written as a sum of $\OO(1)$ normal graphs $\Gamma_k$, i.e., $\Gamma=\sum_k \Gamma_k$, then we define 
\be\label{eq_defsizemax}\size(\Gamma):=\max_k \size(\Gamma_k).\ee
As a convention, for any (possibly non-normal) graph $\Gamma$ with $\OO(1)$ vertices and edges, we define $\size(\Gamma)$ by first expanding $\Gamma$ into a sum of normal graphs via the dotted edge partition, and then applying \eqref{eq_defsizemax}.

%As a convention, given a possibly non-normal graph $\Gamma$ with $\OO(1)$ many vertices and edges, $\size(\Gamma)$ is defined by treating $\Gamma$ as a sum of normal graphs by expanding it with the dotted edge partition. 

%expanding each $G_{xy}$ (resp.~$\overline G_{xy}$) edge as a $(G_{xy}-m\delta_{xy})$ edge plus a $m\delta_{xy}$ edge (resp.~$(\overline G_{xy}-\overline m\delta_{xy})$ edge plus a $\overline m\delta_{xy}$ edge), and by merging vertices connected via dotted edges.
\end{definition}

First, by \eqref{GiiGEX}, each off-diagonal solid edge or light-weight contributes a factor of $\OO_\prec(\Psi_t)$ under the assumption \eqref{initialGT2}, which accounts for the \smash{$(\Psi_t)^{n_S(\Gamma)}$} factor in \eqref{eq_defsize}. 
Next, every molecule contains a ``free vertex'' that can range over the entire lattice \smash{$\ZL$}, giving rise to an \smash{$N^{n_M(\Gamma)}$} factor. The remaining vertices in the molecule are confined (up to a sufficiently small error, by \eqref{scalemole}) to a \smash{$\OO(W(\log W)^{3/2})$}-neighborhood of the free vertex, yielding the factor \smash{$W^{d(n_V(\Gamma)-n_M(\Gamma))}$}.
Finally, each waved edge contributes a $W^{-d}$ factor, by the definition of $S$ and the bound for $S^\pm$: there exist constants $c,C>0$ such that 
\be\label{eq:estSpm-W}S^\pm_{xy}\le CW^{-d}e^{-c|x-y|/W},\ee
which follows from the estimate \eqref{prop:ThfadC_short}. %together with the relation $S^+=m^2 (1-m^2S^{(\sB)})^{-1}\otimes S_W$. 
With these considerations, we obtain the following claim.

\begin{claim}
We have the bound \(\Gamma\prec \size(\Gamma)\) for any normal graph $\Gamma$ in the setting of \Cref{lem:LWterm} (with $\Psi_t=\Psi_t(0)$) or \Cref{lem: EWGn2_N} (with $\Psi_t=(W^{-d}B_{t,0})^{1/2}$).  
\end{claim}

In the proofs, we will only need to keep track of the number of factors of $\Psi_t$ and $W^{-d/2}$ appearing in the scaling size. To this end, we introduce a more convenient notion of \emph{scaling order}. 

\begin{definition}[Scaling order]\label{def scaling order} 
We define the \emph{scaling order} of a normal graph $\Gamma$ as
\be\label{eq:ordG} \ord(\Gamma):=n_S(\Gamma)+ 2 \left(n_W(\Gamma)  -  n_V(\Gamma)\right) . \ee
%where $n_S(\Gamma)$, $n_W(\Gamma)$, and $n_V(\Gamma)$ denote the numbers of solid edges (including light-weights), waved edges, and internal vertices, respectively.
Directly from \eqref{eq_defsize} and \eqref{eq:ordG} we have the identity
\(\size(\Gamma)= (L^d)^{n_M(\Gamma)}(\Psi_t)^{\ord(\Gamma)}\cdot \p{W^{d}\Psi_t^2}^{n_V(\Gamma)-n_W(\Gamma)},\)
so that, since $\Psi_t\ge W^{-d/2}$,
\(\size(\Gamma)\le (L^d)^{n_M(\Gamma)}(\Psi_t)^{\ord(\Gamma)}\) whenever $n_W(\Gamma)\ge n_V(\Gamma)$.
%where $n_M(\Gamma)$ denotes the number of internal molecules. 
% We also define another type of \emph{pseudo-scaling order} as 
% \be\label{eq:ordG} 
% \pord(\Gamma):=n_S(\Gamma)+ 2 \left(n_W(\Gamma)  -  n_V(\Gamma) \right) . \ee
For a general graph $\Gamma$ that can be written as a sum of $\OO(1)$ many normal graphs $\Gamma_k$, we define its scaling order as 
\(\ord(\Gamma)=\min_k \ord(\Gamma_k).\) 
\end{definition}



%\subsection{Graph expansions}\label{sec:local_exp}


We next state the graph expansion rules given in \cite{yang2021delocalization}. These expansions are stated for $G(z)=(H-z)^{-1}$, but they also apply to $G_t(z)$ with $G$ and $S$ replaced by $G_t$ and $S_t=tS$, respectively.\footnote{Here $h_{\al x}$ denotes the $(\al,x)$-th entry of $H$ (resp.~$H_t$), and, as in \cite{yang2021delocalization}, $\partial_{h_{\al x}}$ is the partial derivative with respect to $h_{\al x}$, where $h_{\al x}$ and $h_{x\al}=\overline{h_{\al x}}$ are regarded as independent variables; in particular, $\partial_{h_{\al x}}G_{ij}=-G_{i\al}G_{xj}$.}


\begin{lemma}[Weight expansion, Lemma 3.5 of \cite{yang2021delocalization}] \label{ssl} %[Self-energy expansion]
Suppose $f$ is a differentiable function of $G$. Then,
\begin{align} 
 \Gc_{xx}  f(G)=_{\E}  m \sum_{  \al} S_{x\al}  \Gc_{xx}\Gc_{\al\al} f (G)  &+m^3 \sum_{  \al,\beta}S^{+}_{x\al} S_{\al \beta}  \Gc_{\al\al}\Gc_{\beta\beta} f (G) \nonumber\\
  -  m  \sum_{ \al} S_{x \al} G_{\al x}\partial_{ h_{ \al x}} f (G) &-  m^3 \sum_{ \al,\beta} S^{+}_{x\al}S_{\al \beta} G_{\beta \al}\partial_{ h_{ \beta\al}} f(G)  \, ,\label{Owx}\end{align}
 where ``$=_{\E}$'' means ``equal in expectation''. % and we abbreviate that $\Gc\equiv G-M$.
%  where  $\cal Q_w$ is a sum of $Q$-graphs defined as
%  \begin{align} 
%   \cal Q_w &=  Q_x \left[\Gc_{xx}  f(G)\right] + \sum_{\al} Q_\al \left[  S^{+}_{x\al}\Gc_{\al\al} f(G)\right] \nonumber\\
%  & - m  Q_x\Big[ \sum_{\al} S_{x\al}  \Gc_{\al\al} G_{xx} f(G) \Big]  - m^3 \sum_\al Q_\al\Big[ \sum_{ \beta}S_{x\al}^{+}S_{\al\beta}\Gc_{\beta\beta}G_{\al\al} f(G) \Big] \nonumber\\
%   &+  m  Q_x  \Big[  \sum_\al S_{x \al}  G_{\al x}\partial_{ h_{\al x}} f(G)\Big]+ m^3 \sum_\al Q_\al \Big[  \sum_{ \beta} S^{+}_{x\al}S_{\al \beta} G_{\beta \al}\partial_{ h_{\beta \al}} f(G)\Big]. \nonumber
% \end{align} 
\end{lemma} 


\begin{lemma}[Edge expansion, Lemma 3.10 of \cite{yang2021delocalization}] \label{Oe14}
Suppose $f$ is a differentiable function of $G$. Consider %a graph
\be\label{multi setting}
\cal G := \prod_{i=1}^{k_1}G_{x y_i}  \cdot  \prod_{i=1}^{k_2}\overline G_{x y'_i} \cdot \prod_{i=1}^{k_3} G_{ w_i x} \cdot \prod_{i=1}^{k_4}\overline G_{ w'_i x} \cdot f(G).
\ee
If $k_1\ge 1$, then we have that 
\begin{align} 
  \cal G & =_{\E} m\delta_{xy_1} {\cal G}/{G_{xy_1}} + m   \sum_\al S_{x\al }\Gc_{\al \al} \cal G \label{Oe1x}\\
  &+\sum_{i=1}^{k_2} |m|^2  \left( \sum_\al S_{x\al }G_{\al y_1} \overline G_{\al y'_i}\right)\frac{\cal G}{G_{x y_1} \overline G_{xy_i'}}+ \sum_{i=1}^{k_3} m^2 \left(\sum_\al S_{x\al }G_{\al y_1} G_{w_i \al} \right)\frac{\cal G}{G_{xy_1}G_{w_i x}}      \nonumber \\
& + \sum_{i=1}^{k_2}m  \Gc^-_{xx} \left( \sum_\al S_{x\al }G_{\al y_1} \overline G_{\al y'_i}\right)\frac{\cal G}{G_{x y_1} \overline G_{xy_i'}}  + \sum_{i=1}^{k_3} m  \Gc_{xx}  \left(\sum_\al S_{x\al }G_{\al y_1} G_{w_i \al} \right)\frac{\cal G}{G_{xy_1}G_{w_i x}}   \nonumber\\
 &+(k_1-1) m  \sum_\al S_{x\al } G_{x \al} G_{\al y_1}\frac{ \cal G}{G_{xy_1}}   + k_4 m   \sum_\al S_{x\al }\overline G_{\al x} G_{\al y_1}  \frac{\cal G}{G_{xy_1}} - m    \sum_\al S_{x\al } \frac{\cal G}{G_{x y_1}f(G)}G_{\al y_1}\partial_{ h_{\al x}}f (G)   \, .\nonumber
\end{align}
% where $\cal Q_{e} $ is a sum of $Q$-graphs, 
% \begin{align*}
%   \cal Q_{e} &= Q_x \left( \cal G\right) -  \sum_{i=1}^{k_2}m Q_x \left[ \overline G_{xx}  \left( \sum_\al S_{x\al }G_{\al y_1} \overline G_{\al y'_i}\right)\frac{\cal G}{G_{x y_1} \overline G_{xy_i'}} \right] \\
% & - \sum_{i=1}^{k_3} m Q_x \left[ G_{xx} \left(\sum_\al S_{x\al }G_{\al y_1} G_{w_i \al} \right)\frac{\cal G}{G_{xy_1}G_{w_i x}}\right] -  m Q_x \left[  \sum_\al S_{x\al }\Gc_{\al \al}  \cal G\right]\\
% & -(k_1-1) m Q_x \left[  \sum_\al S_{x\al }G_{x \al} G_{\al y_1}  \frac{ \cal G}{G_{xy_1}}\right]  - k_4 m Q_x \left[  \sum_\al S_{x\al }\overline G_{\al x} G_{\al y_1}  \frac{\cal G}{G_{xy_1}}\right] \\
% & +m Q_x \left[  \sum_\al S_{x\al } \frac{\cal G}{G_{x y_1}f(G)}G_{\al y_1}\partial_{h_{\al x}}f(G)  \right].
%  \end{align*} 
Here, the fractions are used to simplify the expression. For example, the fraction ${\cal G}/({G_{x y_1} \overline G_{xy_i'}})$ is the  graph obtained by removing the factor \smash{${G_{x y_1} \overline G_{xy_i'}}$} from the product in \eqref{multi setting}.
We refer to the above expansion as the \emph{edge expansion with respect to $G_{xy_1}$}. The edge expansion with respect to other $G_{xy_i}$ can be defined in the same way. The edge expansions with respect to $\overline G_{x y'_i}$, $ G_{ w_i x}$, and $ \overline G_{ w'_i x}$ can be defined similarly by taking complex conjugates or matrix transpositions of \eqref{Oe1x}.   
 \end{lemma}


 
\begin{lemma} [$GG$ expansion, Lemma 3.14 of \cite{yang2021delocalization}]\label{T eq0}
Consider a graph $\cal G= G_{xy}   G_{y' x }  f (G)$, where $f$ is a differentiable function of $G$. We have that 
\begin{align}
 \cal G& =_{\E} m\delta_{xy}G_{y' x}f(G)+  m^3  S^{+}_{xy} G_{y' y} f(G)  + m \sum_\al  S_{x\al} \Gc_{\al \al} \cal G + m^3 \sum_{\al,\beta}  S^{+}_{x\al}  S_{\al\beta} \Gc_{\beta \beta} G_{\al y}   G_{y'\al} f(G) \nonumber \\
    &  + m\Gc_{xx }    \sum_\al S_{x\al}G_{\al y}   G_{y'\al} f(G) + m^3 \sum_{\al,\beta}  S^+_{x\al} S_{\al\beta}  \Gc_{\al\al }  G_{\beta y}   G_{y'\beta} f(G)  \nonumber\\
    & - m \sum_{ \al}  S_{x\al}G_{\al y} G_{y' x} \partial_{ h_{\al x}}f(G) - m^3 \sum_{\al,\beta} S^{+}_{x\al} S_{\al\beta}G_{\beta y} G_{y' \al} \partial_{ h_{\beta \al}}f(G)  . \label{Oe2x}
    \end{align}
The $\overline G \overline G$ expansion with respect to $\overline G_{xy}  \overline G_{y' x }$ can be derived by taking the complex conjugate.
% where $\cal Q_{GG} $ is a sum of $Q$-graphs,
%     \begin{align}
%   \cal Q_{GG}&=  Q_x \left(\cal G\right)- m\delta_{xy}Q_x\left(G_{y' x}f(G)\right)+  m^2 \sum_\al Q_\al\Big[ S^+_{x\al}  G_{\al y}   G_{y'\al} f(G) \Big]  - m^3 Q_y\Big[S^{+}_{xy}  G_{y' y} f(G)\Big]  \nonumber\\
%    &- m Q_x\Big[\sum_\al  S_{x\al} \Gc_{\al\al} \cal G\Big] - m^3 \sum_\al Q_\al\Big[\sum_{ \beta}  S^{+}_{x\al}  S_{\al\beta} \Gc_{\beta \beta} G_{\al y}   G_{y'\al} f(G)\Big] - mQ_x\Big[ G_{xx }  \sum_\al S_{x\al} G_{\al y}   G_{y'\al} f(G)\Big] \nonumber\\
% &- m^3 \sum_\al Q_\al\Big[ \sum_{ \beta}  S^{+}_{x\al}  S_{\al\beta} G_{\al\al }  G_{\beta y}   G_{y'\beta} f(G)\Big]  + m Q_x\Big[\sum_{ \al} S_{x\al} G_{\al y} G_{y' x} \partial_{ h_{ \al x}}f(G)\Big] \nonumber\\
% & + m^3\sum_\al Q_\al\Big[ \sum_{\beta} S^{+}_{x\al} S_{\al\beta}G_{\beta y} G_{y' \al} \partial_{h_{\beta \al}}f (G)\Big] .\nonumber 
%  \end{align}
\end{lemma} 

Corresponding to the three lemmas above, we can define graph operations that represent the weight, edge, and $GG$ expansions. Following \cite{yang2021delocalization}, we refer to all these operations as \emph{local expansions at a vertex $x$}, meaning that they do not create new molecules: every new vertex introduced by such an expansion is connected to $x$ by a path of dotted or waved edges. 
As discussed in \cite[Section 3]{yang2021delocalization}, these expansions decrease the scaling size (equivalently, increase the scaling order) of a graph. Moreover, relative to the original graph $\cal G$ prior to expansion, each new graph produced by an expansion either becomes ``smaller'' by a factor of $\Psi_t$ in scaling size, or moves ``closer'' to being locally standard (as defined in \Cref{deflvl1}), in one of the following ways:
\begin{itemize}
    
\item It contains one fewer weight than $\cal G$.

\item The solid-edge degree of a vertex is reduced by 2, a new vertex of degree 2 is introduced, and the degrees of all other vertices remain unchanged. Here, the degree refers only to the number of solid edges incident to the vertex.

%The solid-edge degree of a vertex is reduced by $2$, a new vertex of degree 2 is created, and the degrees of all other vertices remain unchanged. Here, the degree only concerns the solid edges attached to a vertex.  
    
\item A pair of edges of the form $G_{xy} G_{y x }$ (resp.~$\overline G_{xy} \overline G_{yx}$) is replaced by $m^4 S^{+}_{xy}$ (resp.~$\overline m^4 S^{-}_{xy}$).

\end{itemize}
Consequently, by repeatedly applying the local expansions from Lemmas \ref{ssl}--\ref{T eq0}, any normal graph can be expanded into a sum of \emph{locally standard graphs}, defined as follows. 


 \begin{definition} [Locally standard graphs] \label{deflvl1}
A graph is \emph{locally standard} if: 
\begin{itemize}
\item[(i)] It is a normal graph. % (recall Definition \ref{defnlvl0});  

\item[(ii)] It has no self-loops (i.e., weights or light-weights) on vertices.

\item[(iii)] Each internal vertex is either standard neutral, or it is not incident to any solid ($G$ or $\Gc$) edge.
\end{itemize}
Here, a vertex is called \emph{standard neutral} if it satisfies the following two properties: 
%it is connected to exactly two solid edges carrying \emph{opposite charges}.
\begin{itemize}
\item It is connected to exactly two solid edges carrying opposite charges. 

\item It has a neutral charge, where the \emph{charge} of a vertex is defined by counting the incoming and outgoing blue solid edges (with $+$ charge) and red solid edges (with $-$ charge): 
\be\label{eq:neutralcharge}
\qquad \#\{\text{incoming $+$ or outgoing $-$ solid edges}\}- \#\{\text{outgoing $+$ or incoming $-$ solid edges}\} .\ee
\end{itemize}
By definition, the two solid edges attached to a standard neutral vertex $x$ take the form $G_{xy}\overline G_{xy'}$ or $G_{yx}\overline G_{y'x}$.
\end{definition}

Simply speaking, \emph{locally standard graphs} are those in which each $G$ edge is paired with a unique $\overline G$ edge at internal vertices. As discussed in \cite[Section 3.4]{yang2021delocalization}, by repeatedly applying local expansions, any normal graph can be expanded into a sum of locally standard graphs, up to a sufficiently small error, as summarized in the following lemma. For the reader's convenience (and for later use in the proof of \Cref{lem:localregular}), we will recall the local expansion strategy from \cite[Section 3]{yang2021delocalization} in \Cref{sec:pflocalregular}. 



\begin{lemma}[Lemma 3.22 of \cite{yang2021delocalization}] \label{lvl1 lemma}
Let $\Gamma$ be an arbitrary graph with $\OO(1)$ vertices and edges. For any large constant $D>0$, we can expand $\Gamma$ into a sum of $\OO(1)$ many locally standard graphs $\Gamma_\mu$:
\begin{align}\label{expand lvl1}
\Gamma =_{\E} \sum_\mu \Gamma_\mu + \Err,
\end{align}
where every $\Gamma_\mu$ has %scaling size $\lesssim \size(\Gamma)$ and 
scaling order $\ge \ord(\Gamma)$, and $\Err$ is a sum of graphs of scaling size $\OO(W^{-D})$.
\end{lemma}


\subsection{Examples and nested property}\label{sec:graphs_ideas}

% The light-weight term \smash{\( {\cal E}^{\Gc,(2)} \)} in \eqref{eq:idea_L-Kn=2} for $n=2$ is given by 
% \be\label{eq;LWG2}
% {\cal E}^{\Gc,(2)}_{ab} = W^d \sum_{u, v} S^{(\sB)}_{uv} \cdot \tr[\Gc_t E_{u}]\cdot \mathcal{L}^{(3)}_{a,b,v}.
% \ee
% Here, the $3$-loop $\mathcal{L}^{(3)}$ takes the form
% \be\label{eq;L(3)}
% {\cal L}^{(3)}_{a,b,v} = W^{-3d} \sum_{x \in [a], y \in [b], z \in [v]} (G_t)_{xy} (G_t)_{xz} (G_t)_{zy},
% \ee
% where $x \in [a]$ means that the lattice point $x$ belongs to the block labeled by $a$. We expect a local law to hold in the form \smash{\( |\tr[\Gc_t E_{u}]| \leq W^{-c} \)} for some constant $c > 0$. Then, assuming the bound
% \be\label{eq:aprioriG2} 
% |(G_t)_{xy}| \lesssim W^{-d/2}
% \p{\Theta_{t,ab}}^{1/2} ,
% \ee 
% we obtain the estimate
% \begin{align} 
% |{\cal E}^{\Gc,(2)}_{ab}|  & \lesssim W^{-c + d} \sum_{|z| \lesssim W\ell_t} |(G_t)_{xy} (G_t)_{xz} (G_t)_{zy}| 
% \lesssim W^{-c-d/2} \sum_{|v-a|\vee |v-b| \lesssim \ell_t} \p{\Theta_{t,ab} \Theta_{t,av} \Theta_{t,vb}}^{1/2} \nonumber  \\
% & \lesssim W^{-c-d/2} (\Theta_{t,ab})^{1/2} \ell_t^2 \lesssim W^{-c-d/2} (\Theta_{t,ab})^{1/2}\cdot (1 - t)^{-1}.\label{eq:introGc}
% \end{align}
% After integrating in time, the singularity \( (1 - t)^{-1} \) produces a harmless logarithmic factor.  However, the resulting bound gives only \smash{\( (\Theta_{t,ab})^{1/2} \)} instead of the desired \smash{\( \Theta_{t,ab} \)}. The missing factor \(\Theta_{t,ab})^{1/2} \sim |a-b|^{-(d-2)/2} \) is tiny since $|a-b|$  can be as large as $L$. 
%  \normalcolor 

% It is important to note that without an extra factor \smash{\( (\Theta_{t,ab})^{1/2} \)}, the light-weight term would be much larger than the leading term and violate the approximation in \eqref{just}. Moreover, in the actual proof, the situation is even more delicate, as the bound \eqref{eq:aprioriG2} is not available a priori. \nc


%It should be noted that the smallness of the martingale and light-weight terms is a highly nontrivial fact. Without the correct perspective, these error terms may appear significantly larger than the main term. Indeed, in the early development of the theory of random band matrices, researchers often encountered the seemingly impenetrable issue that the error terms in the so-called $T$-equation diverge. While we do not delve into the historical development in this paper, we will show in the following discussion that the light-weight term is, in fact, a genuinely small error.

% To illustrate the idea, we consider the case $u=v$ in \eqref{eq;LWG2}, and disregard the light-weight term $\tr[\Gc_t E_{u}]|$, which in the end contributes to the $W^{-c}$ factor. Moreover, we fix a single term with vertices $x\neq y$ in \eqref{eq;L(3)}---the local averages over $x\in[a]$ and $y\in[b]$ will play no role in extracting the additional long edge. In \eqref{eq;L(3)}, the edge $(G_t)_{xy}$ already contributes a factor $\Psi_t(|a-b|)$ by \eqref{eq;decayGentry}. Thus, it remains to establish the following high-moment estimate for any fixed $p\in 2\N$:
% \be\label{eq:thelongedge}
% \E\Big|\sum_{z}(G_t)_{xz}(G_t)_{zy}\Big|^{p}\prec  \eta_t^{-p}[\Psi_t(|a-b|)]^p. 
% \ee
% By Markov's inequality, this yields $\sum_{z}(G_t)_{xz}(G_t)_{zy}\prec \eta_t^{-1} \Psi_t(|a-b|)$, which provides the missing long edge $\Psi_t(|a-b|)$, along with the factor $\eta_t^{-1}$. 

% We begin with the second moment, $p=2$:
% \begin{center}
% \begin{tikzpicture}
%     %First picture
%     \coordinate (a) at (0,0);
%     \fill[black] (a) circle (2pt) node[left=2pt]{$x$}; 
%     \coordinate (b) at (2,1); 
%     \fill[black] (b) circle (2pt) node[above=2pt]{$z_1$};
%     \coordinate (c) at (4,0);
%     \fill[black] (c) circle (2pt) node[right=2pt]{$y$};
%     \coordinate (d) at (2,-1);
%     \fill[black] (d) circle (2pt) node[below=2pt]{$z_2$};


% \path [draw=blue]
% (a) -- (b) -- (c);
% \path [draw=red]
% (a) -- (d) -- (c);

% % \end{tikzpicture}
% % \begin{tikzpicture}
%     % middle picture
%     \coordinate (L) at (5,0);       % incoming from left
%     \fill[black] (L) circle (2pt) node[left=2pt]{$x$};
%     \coordinate (R) at (9,0);        % outgoing to right
%     \fill[black] (R) circle (2pt) node[right=2pt]{$y$};

% \node (G1) at (6.3,0) {$\cal G_1$};

%     % Internal interaction vertices
%     \coordinate (T) at (7,1);  % top vertex
%     \fill[black] (T) circle (2pt) node[above=2pt]{$\al_{1}$};
%     \coordinate (B) at (7,-1); % bottom vertex
%     \fill[black] (B) circle (2pt) node[below=2pt]{$\al_2$};

% \path [draw=blue]
% (L) -- (T) -- (B) -- (R);
% \path [draw=red]
% (L) to [bend left] (T) to [bend left] (B) to [bend right] (R);

% %Third picture
% \coordinate (L) at (10,0);       % incoming from left
% \fill[black] (L) circle (2pt) node[left=2pt]{$x$};
% \coordinate (R) at (14,0);        % outgoing to right
% \fill[black] (R) circle (2pt) node[right=2pt]{$y$};

% \node (G2) at (11.3,0) {$\cal G_2$};

% % Internal interaction vertices
% \coordinate (T) at (12,1);  % top vertex
% \fill[black] (T) circle (2pt) node[above=2pt]{$\al_{1}$};
% \coordinate (B) at (12,-1); % bottom vertex
% \fill[black] (B) circle (2pt) node[below=2pt]{$\al_2$};

% \path [draw=blue]
% (L) -- (B) -- (T) -- (R);
% \path [draw=red]
% (L) to (T) to [bend left] (B) to (R);
% \end{tikzpicture}
% \end{center}
% The first diagram illustrates the LHS of \eqref{eq:thelongedge} in this case, where blue and red edges denote resolvent entries with charges $+$ and $-$, respectively. For simplicity, the directions of solid edges are not indicated. A key observation from \cite{yang2021delocalization,yang2022delocalization} is that unless the blue and red edges are perfectly paired at each vertex (in which case we call the graph \emph{locally regular}), the graph can be expanded according to certain rules (the \emph{local expansions} of \cite{yang2021delocalization,yang2022delocalization}), with each new graph having a smaller value than the original. 
% Hence, during the expansion, a $-$ charged edge from the path $(x\to z_2\to y)$ must be ``pulled'' to the vertex $z_1$, while a $+$ charged edge from the path $(x\to z_1\to y)$ must be pulled to $z_2$. The second and third diagrams depict two resulting graphs, $\cal G_1$ and $\cal G_2$, which represent the leading contributions. We emphasize that the vertices $\alpha_1$ and $\alpha_2$ in these graphs do not correspond to single real vertices, but rather to equivalence classes of vertices, known as \emph{molecules}. For clarity, we suppress this complication here.

%Here, the first picture depicts the LHS of \eqref{eq:thelongedge} with $p=2$, where the blue and red edges represent resolvent entries and are assigned $+$ and $-$ charges, respectively. In addition, we do not indicate the directions of the solid edges. A key observation of \cite{yang2021delocalization,yang2022delocalization} is that, unless the blue edges and red edges are perfectly paired at each vertex of a graph (in which case we say the graph is ``locally regular''), this graph can be expanded according to certain rules (which are referred to as ``local expansions'' in \cite{yang2021delocalization,yang2022delocalization}), so that every new graph from the expansion has smaller value than the original graph.  

To illustrate the ideas underlying the proof of \Cref{lem:LW_moment}, we examine the simplest case with $p=2$. 

%Then, we will bound each of them using some graphical structures of their molecule graphs. Before proceeding with the proof of the general case, we consider an example with $p=2$ and use it to illustrate the basic ideas behind the proof. 
\begin{example}\label{example:p=2}
Consider the second moment of $f_{xy}(G)$: 
\be\label{eq:p=2graph}\E|f_{xy}(G)|^{2}=\E\sum_{\al_1,\al_2\notin\{x,y\}}\sum_{\beta_1,\beta_2}S_{\al_1\beta_1}S_{\al_2\beta_2} \p{
\Gc_{\beta_1\beta_1} G_{x\al_1}G_{\al_1 y}}
 \cdot \p{\Gc^*_{\beta_2\beta_2}\overline G_{x\al_2}\overline G_{\al_2 y}},\ee
which can be represented by the following graph: %, where we omit the $\times$-dotted edges: 
\begin{center}\scalebox{0.8}{
\begin{tikzpicture}
  \begin{feynman}
    % Define vertices of the square
    \vertex (a) at (0,0);
    \fill[black] (a) circle (1pt) node[left=3pt]{$x$}; 
    \vertex (b) at (2,1); 
    \fill[black] (b) circle (1pt) node[right=2pt]{$\al_1$};
    \vertex [dot, fill] (c) at (4,0);
    \fill[black] (c) circle (1pt) node[right=3pt]{$y$};
    \vertex [dot, fill](d) at (2,-1);
    \fill[black] (d) circle (1pt) node[left=2pt]{$\al_2$};

    % External photon vertices
    \vertex [dot, fill] (top) [above=0.8cm of b];
    \fill[black] (top) circle (1pt) node[right=2pt]{$\beta_1$};
    \vertex [dot, fill] (bot) [below=0.8cm of d];
    \fill[black] (bot) circle (1pt) node[left=2pt]{$\beta_2$};
    
    % Tadpole endpoints
    \vertex (tadtop) [left=0.8cm of top];
    \vertex (tadbot) [right=0.8cm of bot];

\path [draw=blue,postaction={on each segment={mid arrow=blue}}]
(a) -- (b) -- (c);
\path [draw=red,postaction={on each segment={mid arrow=red}}]
(a) -- (d) -- (c);
% Draw the fermion square with arrows
\diagram*{      
      (b) -- [photon] (top) -- [half left, looseness=1.5, draw=blue] (tadtop) -- [half left, looseness=1.5, draw=blue] (top),
      (d) -- [photon] (bot) -- [half left, looseness=1.5, color=red] (tadbot) -- [half left, looseness=1.5, color=red] (bot),
};
  \end{feynman}
\end{tikzpicture}}
\end{center}
%By the averaged local law \eqref{GavLGEX}, the two light-weights at $\beta_1$ and $\beta_2$ contribute factors of $\Psi_t^2$ each. The main difficulty lies in handling the two summations over $\al_1$ and $\al_2$. A straightforward application of the Cauchy--Schwarz inequality and Ward's identity to these sums would yield the bound $(\eta_t^{-1}\Psi_t^2)^2$, which falls short by two long-leg factors $\Psi_t(c|a-b|)$ compared to the target estimate \eqref{eq:LW_moment}. To overcome this, assume without loss of generality that $|\al_1-x|\ge |\al_1-y|$ and $|\al_2-x|\ge |\al_2-y|$. Then, we have $|\al_1-x|\gtrsim W|a-b|$ and $|\al_2-x|\ge |x-y|/2\gtrsim W|a-b|$, so the edges $G_{x\al_1}$ and $\overline G_{x\al_2}$ yield the desired factor $[\Psi_t(c|a-b|)]^2$. However, each of the vertices $\al_1$ and $\al_2$ is now incident to only one $G$ edge. Applying the Cauchy--Schwarz inequality and Ward's identity in this setting produces a factor $N\eta_t^{-1}$, which is significantly worse than the desired $\eta_t^{-2}$.
We expand \eqref{eq:p=2graph} using the local expansions from \Cref{sec:graphs}. As shown in \Cref{lvl1 lemma}, these expansions decrease the scaling size of the graphs, while improving their structure for our purposes. 
More precisely, to obtain locally standard graphs, during the expansion process, a $-$ charged edge from the path $(x\to \al_2\to y)$ or from the light-weight at $\beta_2$ must be ``pulled'' to the molecule containing $\al_1$,\footnote{Pulling an edge $e=(\al,\beta)$ to a molecule $\Mol$ means replacing it by two edges, one connecting $\al$ to $\Mol$ and one connecting $\beta$ to $\Mol$.} whereas a $+$ charged edge from the path $(x\to \al_1\to y)$ or from the light-weight at $\beta_1$ must be pulled to the molecule containing $\al_2$. 
As an illustration, consider the left panel of \Cref{fig:p=2expansion}, which shows a graph $\cal G$ obtained by applying two weight expansions from \Cref{ssl} ($\cal G$ is not yet locally standard, but it is already sufficient for our proof). 
The two short edges $G_{\beta_1\gamma_1}$ and $G_{\beta_2\gamma_2}$ contribute a $[\Psi_t(0)]^2$ factor. For the rest of the graph, we examine its molecular graph $\cal G_1$ shown in the middle panel of \Cref{fig:p=2expansion}, where $\cal M_1$ and $\cal M_2$ denote the molecules containing $\al_1$ and $\al_2$, respectively. %By the triangle inequality, at least one of the edge pairs---between $x$ and $\cal M_1$, between $\cal M_1$ and $\cal M_2$, or between $\cal M_2$ and $y$---has length $\gtrsim W|a-b|$. 
Without loss of generality, suppose the two edges $(x,\cal M_1)$ provide the factor $[\Psi_t(c|a-b|)]^2$. For the remaining graph, we first sum over $\cal M_1$ and then over $\cal M_2$, applying the Cauchy--Schwarz inequality and Ward's identity to obtain a factor $\eta_t^{-2}$. Altogether, this gives the desired bound in \eqref{eq:LW_moment} for $p=2$. 


In the right panel of \Cref{fig:p=2expansion}, we illustrate another possible molecular graph generated from the expansion of \eqref{eq:p=2graph}. For this graph, there are choices of long edges such that, after removing them, each summation order in the resulting graphs yields a poor bound. Nevertheless, by applying an additional AM--GM inequality, the graph still leads to the correct estimate. 
%For this graph, there exist choices of long edges, after removing which, each summation order in the resulting graphs leads to a bad bound. But with an additional application of the AM--GM inequality, this graph can still produce a correct bound. 
To demonstrate this, suppose $(y,\cal M_1)$ and $(y,\cal M_2)$ are taken as long edges. Then, in the resulting graph $\cal G_2'$ (with these edges removed), both $\cal M_1$ and $\cal M_2$ have degree $3$. 
Summing over either one of them eliminates three incident edges and leaves a graph with only a single solid edge. The subsequent summation produces a poor bound of order \smash{$N^{1/2}\eta_t^{-1/2}$} via Cauchy-Schwarz and Ward's identity. 
To handle this case, we apply the AM--GM inequality $2|z_1z_2|\le |z_1|^2+|z_2|^2$ to the two edges $(x,\cal M_1)$ and $(x, \cal M_2)$, which yields two new graphs, one with a pair of edges $(x,\cal M_1)$ and one with a pair of edges $(x, \cal M_2)$. 
In each of these graphs, the summations over $\cal M_1$ and $\cal M_2$ can be carried out in a proper order, yielding the $\eta_t^{-2}$ factor by applying Ward's identity twice. 

%This produces the second and third graphs in \Cref{fig:p=2expansion2}, which can then be summed over, in a proper order, using Ward's identity to obtain a factor $\eta_t^{-2}$.

%If the pair of edges between $\cal M_1$ and $\cal M_2$ are long, then after removing them, we can sum over $\cal M_1$ and $\cal M_2$ in any order. Alternatively, suppose $\cal M_1$ and $\cal M_2$ are closer to $x$ than to $y$. In this case, the edges between $y$ and $\cal M_1$, and between $y$ and $\cal M_2$, are long edges. Removing them yields a graph that does not immediately satisfy property (2): after deleting one internal molecule together with its incident edges, the remaining molecule is incident to only one edge. 

 


\begin{figure}
\begin{center}
\begin{tabular}{m{4cm}m{4cm}m{4cm}}
\scalebox{0.9}{
\begin{tikzpicture}
  \begin{feynman}
    % External fermions
    \vertex (L) at (-2,0);        % incoming fermion
    \fill[black] (L) circle (2pt) node[left=2pt]{$x$};
    \vertex (R) at (2,0);         % outgoing fermion
\fill[black] (R) circle (2pt) node[right=2pt]{$y$};

    \node (G1) at (-0.8,0) {$\cal G$};
    
    % Internal fermion vertices    
    \vertex [dot] (a) at (0,1);   
    \fill[black] (a) circle (2pt) node[left=2pt]{$\al_1$};
    \vertex [dot] (a1) at (0,1.7);
    \fill[black] (a1) circle (2pt) node[right=2pt]{$\beta_1$};
    \vertex [dot] (a2) at (-0.7,1.7);
    \fill[black] (a2) circle (2pt) node[left=2pt]{$\gamma_1$};
    \vertex [dot] (c) at (0,-1);  
    \fill[black] (c) circle (2pt) node[right=2pt]{$\al_2$};
    \vertex [dot] (c1) at (0,-1.7); 
    \fill[black] (c1) circle (2pt) node[left=2pt]{$\beta_2$};
    \vertex [dot] (c2) at (0.7,-1.7);
    \fill[black] (c2) circle (2pt) node[right=2pt]{$\gamma_2$};

    % Tadpole attachment vertices
    \vertex (ta) at (0.8,2.5);
    \vertex (tc) at (0.8,-2.5);

\path [draw=blue,postaction={on each segment={mid arrow=blue}}](L) -- (a);
\path [draw=blue,postaction={on each segment={mid arrow=blue}}](a) to [bend right] (c1);
\path [draw=blue,postaction={on each segment={mid arrow=blue}}](c2) to [bend right] (R);
\path [draw=blue,postaction={on each segment={mid arrow=blue}}](a1) to [bend right] (a2);
\path [draw=red,postaction={on each segment={mid arrow=red}}](L) -- (a2);
\path [draw=red,postaction={on each segment={mid arrow=red}}](a1) to [bend left] (c); 
\path [draw=red,postaction={on each segment={mid arrow=red}}] (c) -- (R);
\path [draw=red,postaction={on each segment={mid arrow=red}}](c2) to [bend left] (c1);

    % Draw fermion lines
    \diagram*{
      %(L) -- [fermion] (a) -- [fermion] (b) -- [fermion] (R),
      %(L) -- [fermion] (c) -- [fermion] (d) -- [fermion] (R),

      % Photon exchange lines
      (a) -- [photon] (a1), (a1) -- [photon] (a2), 
      (c) -- [photon] (c1), (c1) -- [photon] (c2),

      % Tadpoles
      %(a) -- [photon, half left, looseness=1.5] (ta) -- [photon, half left, looseness=1.5] (a),
      %(c) -- [photon, half left, looseness=1.5] (tc) -- [photon, half left, looseness=1.5] (c),
    };
  \end{feynman}
\end{tikzpicture}}
& 
\scalebox{0.9}{
\begin{tikzpicture}
  \begin{feynman}
    % External fermion points
    \vertex (L) at (-2,0);       % incoming from left
    \fill[black] (L) circle (3pt) node[left=2pt]{$x$};
    \vertex (R) at (2,0);        % outgoing to right
    \fill[black] (R) circle (3pt) node[right=2pt]{$y$};

\node (G1) at (-0.8,0) {$\cal G_1$};

    % Internal interaction vertices
    \vertex [dot] (T) at (0,1);  % top vertex
    \fill[black] (T) circle (3pt) node[above=2pt]{$\cal M_{1}$};
    \vertex [dot] (B) at (0,-1); % bottom vertex
    \fill[black] (B) circle (3pt) node[below=2pt]{$\cal M_2$};

\path [draw=blue,postaction={on each segment={mid arrow=blue}}]
(L) -- (T) -- (B) -- (R);
\path [draw=red,postaction={on each segment={mid arrow=red}}]
(L) to [bend left] (T) to [bend left] (B) to [bend right] (R);
  \end{feynman}
\end{tikzpicture}}
&
\scalebox{0.9}{
\begin{tikzpicture}
  \begin{feynman}
    % External fermion points
    \vertex (L) at (-2,0);       % incoming from left
    \fill[black] (L) circle (3pt) node[left=2pt]{$x$};
    \vertex (R) at (2,0);        % outgoing to right
    \fill[black] (R) circle (3pt) node[right=2pt]{$y$};

    \node (G1) at (-0.8,0) {$\cal G_2$};

    % Internal interaction vertices
    \vertex [dot] (T) at (0,1);  % top vertex
    \fill[black] (T) circle (3pt) node[above=2pt]{$\cal M_{1}$};
    \vertex [dot] (B) at (0,-1); % bottom vertex
    \fill[black] (B) circle (3pt) node[below=2pt]{$\cal M_2$};

\path [draw=blue,postaction={on each segment={mid arrow=blue}}]
(L) -- (B) -- (T) -- (R);
\path [draw=red,postaction={on each segment={mid arrow=red}}]
(L) to (T) to [bend left] (B) to (R);
  \end{feynman}
\end{tikzpicture}}
\end{tabular}
\end{center}
\caption{A possible expansion of \eqref{eq:p=2graph} (left), its molecular graph (middle), and another possible molecular graph (right). The $\times$-dotted edges in the first graph have been omitted.
}
    \label{fig:p=2expansion}
\end{figure}

\end{example}




%We observe that there are sufficiently many long edges (at least two) and resolvent entries (four in total) left to carry out the summations over the vertices $\alpha_1$ and $\alpha_2$ using the Cauchy--Schwarz inequality and Ward's identity. However, even in the simplest case $p=2$, carrying out these summations in a way that yields \eqref{eq:thelongedge} is highly nontrivial: the result depends delicately on both the choice of long edges and the order of summation. %and in some instances further applications of the AM--GM inequality are necessary.
%However, even in the simplest case $p=2$, whether these summations can be performed in a way that leads to \eqref{eq:thelongedge} is highly nontrivial: the outcome depends crucially on the choice of long edges and the order of summation, and additional applications of the AM-GM inequality may be required.
%For instance, in graph $\cal G_1$, if the pair of edges between $\alpha_1$ and $\alpha_2$ are designated as long edges, then the summations over $\alpha_1$ and $\alpha_2$ may be performed in any order. By contrast, if the pair of edges between $y$ and $\alpha_2$ are chosen instead, then $\alpha_2$ must be summed first.
%The situation is more delicate for graph $\cal G_2$. There exist choices of long edges for which no summation order works directly. For example, if $(y,\alpha_1)$ and $(y,\alpha_2)$ are taken as long edges, then in the resulting graph $\cal G_2'$ (with these edges removed), both $\alpha_1$ and $\alpha_2$ have degree $3$. Summing over either one of them eliminates three incident edges and leaves a graph with only a single solid edge. The subsequent summation produces a poor bound of order \smash{$N^{1/2}\eta_t^{-1/2}$} via Cauchy-Schwarz and Ward's identity. 
%To bound $\cal G_2'$, we must first apply the AM--GM inequality $|G_{x\al_1}||G_{x\al_2}|\le |G_{x\al_1}|^2+|G_{x\al_2}|^2$, which yields two new graphs. In each of these graphs, the summations over $\alpha_1$ and $\alpha_2$ can be carried out in an appropriate order, yielding the $\eta_t^{-2}$ factor by applying Ward's identity twice.  






To prove \Cref{lem:LW_moment}, we need to extend the analysis for the $p=2$ example to arbitrarily large $p\in 2\N$. In the original graph $|f_{xy}(G)|^p$, there are $p$ edge-disjoint paths from $x$ to $y$, comprising a total of $2p$ solid edges. Heuristically, to obtain locally standard graphs with local expansions, each path from $x$ to $y$ must be pulled at least once, thereby increasing its length by 1. 
Now, the key graphical property required for our proof consists of the following two components: (i) We can identify at least $p$ ``long edges''---that is, edges of length $\gtrsim |x-y|$---in the graph. (ii) In the remaining graph, after removing these long edges, there exists a specific order of summation over the internal vertices such that, at each step, a sufficient number of edges (at least two per summation) remain to apply the Cauchy-Schwarz inequality and Ward's identity to control the summation. 
If there is no requirement of property (i), this problem was addressed in  \cite{yang2021random} using a structure called  ``\emph{nested property}'' of the graph. 
In \Cref{lem:Anp_key} below, we will extend this concept to inductively select both long solid edges and a summation order to guarantee the bound  \eqref{eq:LW_moment}.



% The local expansion rules remain exactly the same as in previous works \cite{yang2021delocalization, yang2022delocalization}. In the original graph (prior to expansion), there are $p$ edge-disjoint paths from $x$ to $y$, comprising a total of $2p$ solid edges. These $p$ paths are preserved throughout the entire expansion process. Moreover, to generate locally regular graphs---in which blue and red edges are perfectly paired at every vertex---each path must be pulled at least once, thereby increasing its length by 1. As a result, every graph produced through the expansion contains at least $3p$ edges.\footnote{We remark that, strictly speaking, this argument is not fully rigorous, as the number of edges in the graph may decrease when some vertices are identified (i.e., take the same value) during the summation.}
% By the triangle inequality, each path from $x$ to $y$ must contain at least one long edge. Our goal is to ensure that the remaining $2p$ edges are sufficient to control the vertex sums through Ward's identity. 
% However, as the $p=2$ case illustrates, merely having a sufficient number of edges (i.e., $2p$) does not automatically guarantee that the graph can be properly bounded by the RHS of \eqref{eq:LW_moment}. The ordering of the vertex summations in the graph plays a crucial role, and in some cases, additional applications of the AM--GM inequality may be required to obtain the desired bounds.


%Now, a key graphical property essential to our proof consists of the following two components: (i) We can identify at least $p$ long edges in the graph. (ii) In the remaining graph, after removing these long edges, there exists a specific order of summation over the vertices such that, at each step, a sufficient number of edges (at least two per summation) remain to apply the Cauchy-Schwarz inequality and Ward's identity to control the summation. 

%On the one hand, by the triangle inequality, it is trivial to extract one long edge for each of the $p$ paths, leaving at least $2p$ edges in the remaining graph. However, as the $p=2$ example illustrates, merely having a sufficient number of edges does not automatically guarantee that condition (ii) holds. %Even the condition (ii) holds, the ordering of the vertex summations in the graph still plays a crucial role, and in some cases, additional applications of the AM--GM inequality may be required.
%On the other hand, if we ignore condition (i) and do not extract the $p$ long edges, then the resulting graphs satisfy a structural property known as the \emph{nested property}, introduced in \cite{yang2021random}. This property guarantees the existence of a nested summation order for the vertices, ensuring that condition (ii) holds. However, this idea cannot be directly applied in our setting, because the long edges are excluded from the summation process. In particular, if the long edges are not chosen carefully, the remaining graph may violate condition (ii), as has already been seen in our earlier discussion of the graph $\cal G_2$---in some cases, additional applications of the AM--GM inequality may be required.





\subsection{Proof of Lemma \ref{lem:LW_moment}}

% In the general loop case, we have a loop (in graph) with the edges (here $\overline{ab} $ represents the edge connecting $a$ and $b$)
% $$
% {\cal L}_{\mathbf{a}}:=\overline{a_1a_2}, \;\cdots\;
% \overline{a_{n-1}a_n}, \; \overline{a_{n }a_1}
% $$
% Assume at this step, we are at $K-level$, and edge $( \epsilon_i-\epsilon_f)$ can be bounded as 
% $$
% f_K(\epsilon_i,\epsilon_f)\prec  
%     g\left(\|\epsilon_i-\epsilon_f\|_K\right),\quad  \|\epsilon_i-\epsilon_f\|_K:=\min (K, \|\epsilon_i-\epsilon_f\|)
% $$
% ($g(x)=(\Theta_{0x})^{1/2}$ in applicaiton) Therefore, we have (note we don't track $W$ factor  in this graphical expansion)
% $$
% {\cal L}_{\mathbf{a}}\prec 
% \prod_{i=1}^n f_K\left( a_i, a_{i+1} \right)
% \prec \prod_{i=1}^n g\left( \|a_i- a_{i+1}\|_K \right)$$
% The $\widetilde{G}$ term in the flow equation is like replace one edge with a path through some $z$ vertex.
% $$
% {\cal L}_{\mathbf{a}}^{(k)}:=\sum_z {\cal L}_{\mathbf{a}}\Bigg|
% \left(\overline{a_k,a_{k+1}}\quad \longrightarrow \quad\overline{a_k, z}+\overline{z,a_{k+1}}\right)
% $$
% The goal is to show that for each $k$ and $n$, we have 
% $$
% \E \left|{\cal L}_{\mathbf{a}}^{(k)}\right|^p\prec \left(\prod_{i=1}^n g\left(\|a_i-a_{i+1}\|_K\right)\right)^p\cdot \eta^{-p}
% $$



% we can get the following graphs from expanding  $\E \left|{\cal L}_{\mathbf{a}}^{(k)}\right|^p$
% \begin{itemize}
%     \item The graph has vertices
%     $$
%     a_1,\cdots,\; a_n,\;z_1,\cdots, z_p$$
%     \item There are independent $n\times p$ paths (not edges), more precisely they are $$
%    a_i\to a_{i+1} ,\quad 1\le i\le n$$
%   and there are $p$ paths for each of them. 
%   \item for each $z_k$, there are exactly two paths pass through it, i.e.,  
%     $$\deg (z_k)=4,\quad 1\le k\le p
%     $$
 %   \item  Let $A\subset\{1,2\cdots p\}$,  assume that there is not edges connecting $A$ and $A^c$, i.e., 
 %$$
%  k\in A, l\in A^c \implies  |\overline{z_k,z_l}|= 0 
% $$
% Then there are at least $|A|$ paths passing through $\{z_i: i\in A\}$. 
%\end{itemize}

% We can pick up the paths which visit through the $\{z_i: 1\le i\le p\}$. Let's say there are $m$ such paths. Let's say that the ending points of these $m$ paths are
% $$
% \left(a_{i_1}, a_{f_1}\right), \quad \left(a_{i_2}, a_{f_2}\right),\cdots \left(a_{i_m}, a_{f_m}\right),$$
% We only need to bound  the edges in these $m$ paths by 
% $$
% \prod_{s=1}^m g(\|a_{i_s}- a_{f_s}\|_K)\cdot \eta^{-p}
% $$
% Therefore, we only need to prove the following lemma

For the proof of \Cref{lem:LW_moment}, we begin by expanding $\big|f_{xy}(G)\big|^{p} = [f_{xy}(G)]^{p/2}[\overline{f_{xy}(G)}]^{p/2}$ into a sum of locally standard graphs, as in \Cref{lvl1 lemma}. By carefully tracking the local expansion process, we can show that each locally standard graph satisfies the graphical properties in the following lemma.  





% More precisely, our proof strategy can be summarized as follows. In the initial graph, there are $p$ disjoint paths from $x$ to $y$, corresponding to the $p$ factors $f_{xy}(G)$ or $\overline{f_{xy}(G)}$. Each such path passes through an internal molecule, and each internal molecule is connected by two solid edges of the same color.
% To obtain locally standard graphs via local expansions, each molecule must either (i) pull in an edge of a different color from other paths or molecules, or (ii) merge with another molecule connected by differently colored edges. In case (i), the number of solid edges increases by at least 1. In case (ii), the number of free vertices decreases by at least 1 (since each molecule contains exactly one free vertex). Consequently, the number of solid edges needed for later summations over free vertices is reduced. %This mechanism is the key source of the additional solid edges required in our proof. 
% In the locally standard graphs produced by these expansions, each path necessarily contains at least one long leg by the triangle inequality. Moreover, by mechanism (i) or (ii) above, there remain enough solid edges to carry out the summations over internal molecules: if there are $k$ internal molecules, then there are at least $2k$ solid edges connecting different molecules.
% However, as discussed in \Cref{subsec_idea}, it is not sufficient merely to count solid edges; the connection pattern and the summation order are also crucial. We must show that after removing the $p$ long edges, there still exists an ordering of the summations over internal molecules such that, at each step, the molecule being summed over is connected to at least two solid edges. This ensures that Cauchy--Schwarz and Ward's identity can be applied to control the sums. {\cor This issue was already discussed in the above example with $p=2$, but for general $p$ it becomes considerably more delicate.  We postpone the detailed discussion to \Cref{sec:pf_Anp}, where the setting is cleaner.}
% Finally, we remark that compared with the $\eta_t^{-1}$ and $\Psi_t(c|a-b|)$ factors in \eqref{eq:LW_moment}, the $\Psi_t(0)$ factors are relatively straightforward to obtain, by simply counting the scaling order of our graphs. 

%We are now ready to present the formal proof of \Cref{lem:LW_moment}, beginning with the following lemma on the local expansions of the graph $|f_{xy}(G)|^p$. {\cor We postpone its proof to \Cref{sec:pflocalregular}; explain in more details: the path properties and the scaling order condition} 

\begin{lemma}\label{lem:localregular}
Given any large constant $D>0$, we can expand $\big|f_{xy}(G)\big|^{p}$ into a linear combination of $\OO(1)$ many \emph{locally standard graphs} $\Gamma_{\mu,xy}$ in the sense of equal expectation:
\be\label{eq:local_Gs}
\big|f_{xy}(G)\big|^{p} =_{\E} \sum_{\mu}  \Gamma_{\mu,xy} + \Err,
\ee
where $\Err$ denotes a sum of $\OO(1)$ many graphs of scaling size $\OO(W^{-D})$.
Moreover, each graph $\Gamma_{\mu,xy}$ on the RHS satisfies the following properties under the assumption \eqref{eq:far_ab} (recall that, up to an exponentially small error $\exp(-c(\log W)^{3/2})$, $x$ and $y$ do not belong to the same molecule by \eqref{scalemole}):
\begin{enumerate}
    \item[(1)] $\Gamma_{\mu,xy}$ is a locally standard graph with external vertices $x$ and $y$, with the corresponding external molecules denoted by $\cal M_x$ and $\cal M_y$, respectively.

    \item[(2)] $\Gamma_{\mu,xy}$ contains $0\le q\le p$ many internal molecules, denoted by $\cal M_{i}$ with $i\in \qqq{q}$. Moreover, within each (internal or external) molecule $\Mol_i$, $i\in\{1,\ldots, q, x, y\}$, the number of waved edges (denoted by $n_W(\Mol_i)$) and the number of vertices (denoted by $n_V(\Mol_i)$) satisfy the relation 
    \be\label{eq:MolVW}
n_V(\Mol_i) \le n_W(\Mol_i) + 1.
    \ee

    \item[(3)] %If $\Mol_x$ and $\Mol_y$ are not the same molecule, then 
    There are at least $p$ many edge-disjoint paths $\fP_i$, $i\in\qqq{p}$, in the molecular graph that connect the molecules $\cal M_x$ and $\cal M_y$ via solid edges (and possibly passing through the internal molecules). 
    %If $\cal M_x$ and $\Mol_y$ are the same molecule, then there are at least $q$ many edge-disjoint paths $\fP_i$, $i\in\qqq{q}$, in the molecular graph that form loops on the molecule $\cal M_x=\cal M_y$ and pass through some internal molecules. 

    \item[(4)] For each internal molecule $\cal M_i$, $i\in \qqq{q}$, there are at least two paths $\fP_k\ne \fP_l$ passing through it on the molecular graph. As a consequence, each $\Mol_i$ has degree at least 4, i.e., 
\be\label{eq:deg_mole}
    \deg(\Mol_i)\ge 4,\quad i\in \qqq{q}.\ee

    \item[(5)] For each $A\subset \qqq{q}$, the number of paths passing through 
        $\{\Mol_i: i\in A\}$ is at least $|A|$, that is,
        $$
     \left|\left\{j: \mathrm V(\fP_j) \cap  \{\Mol_i: i\in A\}\ne \emptyset\right\}\right|\ge |A| ,$$
     where $\mathrm V(\fP_j)$ denotes the subset of vertices (i.e., molecules) in the path $\fP_j$ on the molecular graph.

     \item[(6)] The scaling order of $\Gamma_{\mu,xy}$ satisfies that
     \be\label{eq:sizeGammamu}
\ord(\Gamma_{\mu,xy}) \ge 2p .  %\size(\Gamma_{\mu,xy}) \le N^q (\Phi_t)^{2p+2q}. 
     \ee
\end{enumerate}

\end{lemma}


% Following a similar argument as above, we examine the graphs derived from the edge expansion \eqref{Oe1x} and $GG$ expansion \eqref{Oe2x} one by one. This leads to the following observation that generalizes (i) and (ii): 
% \begin{itemize}
%     \item For $1\le k \le 4$, if the number of internal molecules decreases by $k$ from $\cal G_0$ to $\cal G_1$, then the scaling order decreases at most by $(k-1)$. 
% \end{itemize}
% Specifically, the third to sixth terms on the RHS of \eqref{Oe1x} involve the $k\in\{1,2\}$ cases; the seventh and eighth terms on the RHS of \eqref{Oe1x} involves the $k=1$ case; and the last term on the RHS of \eqref{Oe1x} involves the $k\in\{1,2,3\}$ case; 
% We omit the details of the argument. 
%{\cor more details and draw some graphs}


% Now, given $\Gamma_{\mu,xy}$, let $k_0$ be the number of expansions where scenario (1) occurs for some internal molecule. For $r\in \{1,2,3,4\}$, let $k_r$ be the number of expansions where the number of internal molecules are reduced by $r$ from $\cal G_0$ to the new graph $\cal G_1$. By definition, we have 
% \[4k_4 + 3k_3 + 2k_2 + k_1 = p-q.\]
% Moreover, the number of internal molecules being reduced by $r$ means that $(r+1)$ molecules in $\cal G_0$ are merged into one single molecule in $\cal G_1$. Thus, we have that 
% \[  5k_4 + 4k_3 + 3k_2 + 2k_1 + k_0 \ge p, \]
% where the equality holds when all these merging processes only involve \emph{disjoint subset} of \emph{internal molecules}. Finally, from the above discussions, we see that an expansion where scenario (1) occurs leads to an increase of the scaling order by at least 1, while an expansion where the number of internal molecules are reduced by $r$ leads to an ``increase'' of the scaling order by at least $-(r-1)$. Thus, we obtain the following inequality:
% \begin{align*}
%     % 4k_4 + 3k_3 + 2k_2 + k_1 \le p-q \\
%     % 5k_4 + 4k_3 + 3k_2 + 2k_1 + k_0 \ge p \\
%     \ord(\cal G) \ge 3p -3k_4 -2k_3-k_2+k_0.
% \end{align*}
% From the above three inequalities, we derive that 
% \[\ord(\cal G) \ge 3p + (5k_4 + 4k_3 + 3k_2 + 2k_1 + k_0)- 2\cdot (4k_4 + 3k_3 + 2k_2 + k_1 )\ge 2p+2q. \]
% This gives \eqref{eq:sizeGammamu3}, which implies \eqref{eq:sizeGammamu2}, which further concludes \eqref{eq:sizeGammamu}.


The proof of this lemma is straightforward by applying the local expansions in \Cref{ssl,Oe14,T eq0}, and we defer the details to \Cref{sec:pflocalregular}. 
The path properties (3)--(5) in \Cref{lem:localregular} serve as the key structural ingredients for controlling the graphs $\Gamma_{\mu,xy}$ in \eqref{eq:local_Gs}.
To see them, note that in the initial graph, there are $p$ edge-disjoint paths from $x$ to $y$, corresponding to the $p$ factors $f_{xy}(G)$ or \smash{$\overline{f_{xy}(G)}$}. Each path passes through a distinct internal molecule, and each internal molecule is attached to two solid edges of the same color. From the expansion rules in \Cref{ssl,Oe14,T eq0}, one can verify that these paths are preserved throughout the expansion process, which ensures properties (3) and (5).
Furthermore, to obtain locally standard graphs via local expansions, each internal molecule must either (i) pull in an edge of a different color from another path or molecule, (ii) merge with another internal molecule connected by differently colored edges, or (iii) merge with an external molecule. In case (iii), the internal molecule disappears from $\Gamma_{\mu,xy}$, whereas in cases (i) and (ii), at least two paths will pass through it.





To streamline the estimation process, we now simplify the structure of the graphs $\Gamma_{\mu,xy}$ through its underlying molecular graph.
Specifically, we bound them by corresponding ``\emph{auxiliary graphs}'', which are obtained as quotient graphs by collapsing each molecule into a single vertex. 
%we first simplify its structure using their molecular structures. More precisely, we will bound them with the corresponding ``\emph{auxiliary graphs}'', which are actually obtained as quotient graphs with molecules reduced to vertices.
More precisely, for each internal molecule $\Mol_i$, $i\in \qqq{q}$, we select a representative vertex $\al_i$ within the molecule, referred to as the ``center'' of the molecule. Similarly, for the two external molecules, we set $x$ and $y$ as their respective centers. The auxiliary graph associated with $\Gamma_{\mu,xy}$ is then defined on the set of block-level vertices $[x]$, $[y]$, and $[\al_i]$ (for $i\in \qqq{q}$): % as described below.
%For the external molecules, let $x$ and $y$ be the corresponding centers. Then, the auxiliary graph of $\Gamma_{\mu,xy}$ is a graph defined on the vertices of blocks, $[x]$, $[y]$, and $[\al_i]$ (for $i\in \qqq{q}$) defined as follows. 

\begin{definition}[Block-level vertices] \label{def: BM2}
Recall that $\ZL$ is partitioned into $L^d$ blocks as defined in \eqref{eq:blockIa}. For any $x\in \ZL$, we use $[x]$ to denote the block that contains $x$. 
%For any $x\in \ZL$, we denote by $[x]$ the block containing $x$. 
With slight abuse of notation, we also view these blocks as vertices in the lattice $\Zn$. In other words, $[x]$ may be interpreted both as a vertex in $\Zn$ and as a subset of vertices in $\ZL$. 

%With a slight abuse of notation, we will also view these blocks as vertices of the lattice $\Zn$. In other words, we will regard $[x]$ both as a vertex of the lattice $\Zn$ and a subset of vertices on the lattice $\ZL$. 
\end{definition}


For clarity, throughout the following proof we use the notation $[\cdot]$ to indicate vertices in the auxiliary graph and the lattice $\Zn$, distinguishing them from vertices in the original graphs defined in \Cref{def_graph1} and from those in the lattice $\ZL$. 
To define the auxiliary graph formally, we still need to define its edges. 
For simplicity of notation, we will use ``$\al\sim_{\Mol} \beta$'' to mean that ``vertices $\al$ and $\beta$ belong to the same molecule''. 
For any $\beta_i\sim_{\Mol}\al_i$, it suffices to assume that  
\be\label{yixi}
|\al_i-\beta_i|\le W(\log W)^{1+\e_1}
\ee
for a small enough constant $\e_1\in (0,1/10)$; otherwise the graph is smaller than $\exp(-\Omega((\log W)^{1+\e_1}))\le W^{-D}$ for any constant $D>0$. Combining \Cref{lem_GbEXP} with \eqref{yixi}, we can bound the solid edges between different molecules as follows for any large constant $D>0$:  
\be\label{eq:Gbyxi}|(G-M)_{xy}|\prec \Psi_{t}(0)\delta_{xy}+\xi([x],[y])+W^{-D},\ee
where the variables $\xi(\cdot,\cdot)$ are defined as 
\be\label{eq:xia1a2} [\xi([a_1],[a_2])]^2:=\sum_{\bfb:\|\bfb-\ba\|_\infty\le (\log W)^{1+2\e_1}} \max_{\bsig\in\{(-,+),(+,-)\}}\cL^{(2)}_{t,\bsig,\bfb} + W^{-d}\mathbf 1(|[a_1]-[a_2]|\le (\log W)^{1+2\e_1}).\ee
With Ward's identity and the properties of $\Psi_t$ given by \eqref{eq:Psi}, we readily see the following properties for the $\xi$ variables. 
\begin{claim}
In the setting of \Cref{lem:LWterm}, %there exists an absolute constant $c_0>0$ such that 
for any $[a_1],[a_2]\in \Zn$, the following estimates hold:
\be\label{eq:Gbyxi2}
\xi([a_1],[a_2]) \prec \Psi_t(|[a_1]-[a_2]|), \quad \sum_{[a_2]}[\xi([a_1],[a_2])]^2 \prec (W^d\eta_t)^{-1}. 
\ee
\end{claim}

Now, we define the auxiliary graphs formed with edges representing the $\xi$ variables. 

\begin{definition}[Auxiliary graphs]\label{def_auxgraph}
Let $\cal G_{xy}$ be an arbitrary normal graph with $q$ internal molecules $\Mol_i$, $i\in\qqq{q}$, and two external molecules $\Mol_x$ and $\Mol_y$ containing $x$ and $y$, respectively. We fix the centers $\al_i\in\Mol_i$ for $i\in \qqq{q}$. We then define the auxiliary graph $\cal G^{\aux}_{[x][y]}$ of $\cal G_{xy}$ as follows:
\begin{itemize}
    \item The vertices are $[\al_i]$, $i\in \{1,\ldots,q,x,y\}$, where we adopt the convention $\al_x\equiv x$ and $\al_y\equiv y$.
    \item Corresponding to each solid edge between molecules $\Mol_i$ and $\Mol_j$ in $\cal G_{xy}$, we introduce a (black) solid edge between $[\al_i]$ and $[\al_j]$ in $\cal G^{\aux}_{[x][y]}$, representing the factor $\xi([\al_i],[\al_j])$. (Since $\xi$ is symmetric, these edges are not oriented.)
\end{itemize}
In this way, we obtain the auxiliary graph \smash{$\cal G^{\aux}_{[x][y]}$}. Its value is defined analogously to \Cref{ValG}, namely, as the product of all solid edges followed by summation over all internal vertices $[\al_i]$, $i\in \qqq{q}$. Finally, we define the scaling order of \smash{$\cal G^{\aux}_{[x][y]}$} in the same manner as in \eqref{eq:ordG}:
%Then, the auxiliary graph $\cal G^{\aux}_{[x][y]}$ of $\cal G_{xy}$ is a graph with vertices $[\al_i]$, $i\in \{1,\ldots, q, x, y\}$, where we adopt the convention that $\al_x\equiv x$ and $\al_y \equiv y$. There are solid edges between these vertices defined in the following way: corresponding to each solid edge between molecules $\Mol_i$ and $\Mol_j$, there is a solid edge between $[\al_i]$ and $[\al_j]$ representing the $\xi([\al_i],[\al_j])$ factor. (Note that these solid edges are not oriented since $\xi$ is symmetric.) In this way, we define the auxiliary graph \smash{$\cal G^{\aux}_{[x][y]}$}. Its value is defined as in \Cref{ValG} by taking the product of all the solid edges and summing over all internal vertices $[\al_i]$, $i\in \qqq{q}$. Finally, we define the scaling order of $\cal G^{\aux}_{[x][y]}$ in a similar way as \eqref{eq:ordG}:
\be\label{eq:ordGaux} \ord(\cal G^{\aux}_{[x][y]}):=\#\{\text{solid edges in } \cal G^{\aux}_{[x][y]}\}-2\#\{\text{internal vertices in } \cal G^{\aux}_{[x][y]}\}.\ee
\end{definition}

Now, the estimation of $\Gamma_{\mu, xy}$ can be reduced to bounding its auxiliary graph.

\begin{lemma}\label{GtoAG}
Given a locally standard graph $\Gamma_{\mu,xy}$ with external vertices $x$ and $y$ that satisfies properties (1) and (2) in \Cref{lem:localregular} (the path properties (3)--(6) are not used in this lemma, so it applies to any such graph, not only to those produced by \eqref{eq:local_Gs}), we define its auxiliary graph $\Gamma^{\aux}_{\mu, [a][b]}$ (recall that $x\in[a]$ and $y\in[b]$) and the corresponding scaling order as in \Cref{def_auxgraph}. For any large constant $D>0$, we have 
\be\label{G_by_auxG} 
\Gamma_{\mu, xy} \prec  (\Psi_t)^{\ord(\Gamma_{\mu, xy}) - \ord(\Gamma^{\aux}_{\mu, [a][b]})}  \cdot W^{qd}\Gamma^{\aux}_{\mu, [a][b]} + W^{-D}\, .
\ee
\end{lemma}

\begin{proof}
As in \Cref{def_auxgraph}, the vertices of \smash{$\Gamma^{\aux}_{\mu,[a][b]}$} are denoted by $[\al_i]$, $i\in\{1,\ldots,q,x,y\}$, corresponding to the molecule centers $\al_i$ in $\Gamma_{\mu,xy}$. 
First, by the definition of $S$ together with the estimate for $S^\pm$ in \eqref{eq:estSpm-W}, and recalling \eqref{yixi}, we may bound a waved edge between $\beta_i$ and $\gamma_i$ inside molecule $\Mol_i$ by
$$ \OO_\prec\p{W^{-d}\mathbf 1_{|\beta_i-\al_i|\vee |\gamma_i-\al_i|\le W(\log W)^{1+\e_1}} + e^{-c(\log W)^{1+\e_1}}}.$$
Second, by \Cref{lem_GbEXP}, solid edges within molecules can be bounded by $\Psi_t$, and every solid edge between vertices $\beta_i\in \Mol_i$ and $\beta_j\in \Mol_j$ can be bounded by
\[ \Big(\sum_{|[a']-[\beta_i]|\le 1,|[b']-[\beta_j]|\le 1} \max_{\bsig\in \{(+,-),(-,+)\}} \cL^{(2)}_{t,\bsig,([a'],[b'])}+ W^{-d}\mathbf 1_{|[\beta_i]-[\beta_j]|\le 1}\Big)^{1/2} \lesssim \xi([\al_i],[\al_j]), \]
where we again use \eqref{yixi} for $\beta_i\sim_{\Mol}\al_i$ and $\beta_j\sim_{\Mol}\al_j$.

With these bounds, and after summing over all \emph{non-center} vertices inside the molecules $\Mol_i$, $i\in \qqq{q}$, and inside the two external molecules $\Mol_x$ and $\Mol_y$, we obtain (cf. the notations in \Cref{def scaling})
\begin{align*}
    \Gamma_{\mu, xy} \prec (\Psi_t)^{n_S(\Gamma_{\mu, xy})-n_S(\Gamma^{\aux}_{\mu, [a][b]})}W^{-d(n_{W}(\Gamma_{\mu, xy}) - n_V(\Gamma_{\mu, xy})+q)} \cdot W^{qd}\Gamma^{\aux}_{\mu, [a][b]}.
\end{align*}
Here, we have also used \eqref{eq:MolVW} in obtaining the factor $W^{-d(n_{W}(\Gamma_{\mu, xy}) - n_V(\Gamma_{\mu, xy})+q)}$, while the compensating factor $W^{qd}$ reflects the difference between summing the molecule centers $\al_i$ (in $\ZL$ for $\Gamma_{\mu,xy}$) and summing the corresponding block representatives $[\al_i]$ (in $\Zn$ for \smash{$\Gamma^{\aux}_{\mu,[a][b]}$}). Finally, using $\Psi_t\ge W^{-d/2}$ together with the definitions \eqref{eq:ordG} and \eqref{eq:ordGaux}, we conclude \eqref{G_by_auxG}.
%Next, using $\Psi_t\ge W^{-d/2}$ and the definitions \eqref{eq:ordG} and \eqref{eq:ordGaux}, we conclude \eqref{G_by_auxG}.
\end{proof}

To conclude \Cref{lem:LW_moment}, it remains to bound the auxiliary graph $\Gamma^{\aux}_{\mu, [a][b]}$. %This is given by the following lemma.


\begin{lemma}\label{lem:Anp} 
In the setting of \Cref{GtoAG}, there exists a constant $c>0$ (depending on $p$) such that 
\be\label{eq:bddGamma_aux}
\Gamma^{\aux}_{\mu, [a][b]} \prec \left({W^d\eta_t}\right)^{-q} \cdot [\Psi_t(c|[a]-[b]|)]^p \cdot (\Psi_t)^{\ord(\Gamma^{\aux}_{\mu, [a][b]})-p}.
\ee

\end{lemma}

Before turning to the proof of \Cref{lem:Anp}, we first use it to complete the proof of \Cref{lem:LW_moment}.

\begin{proof}[\bf Proof of \Cref{lem:LW_moment}]
Using \Cref{lem:Anp,GtoAG}, we bound the locally standard graphs in \eqref{eq:local_Gs} as
    \begin{align*}
        \Gamma_{\mu, xy} \prec  \eta_t^{-q}\cdot (\Psi_t)^{\ord(\Gamma_{\mu, xy}) -p} \cdot [\Psi_t(c|[a]-[b]|)]^p  + W^{-D} \le \eta_t^{-p}\cdot \Psi_t^{p} \cdot [\Psi_t(c|[a]-[b]|)]^p  + W^{-D},
    \end{align*}
    where we use \eqref{eq:sizeGammamu} in the second step.
    Finally, we upgrade this stochastic domination bound to a bound in expectation. Since $\|G_t\|\le \eta_t^{-1}$ and $\sum_{\al,\beta}S_{\al\beta}=N$, the definition \eqref{fxyG_sum} gives the deterministic bound $|f_{xy}(G)|\lesssim N\eta_t^{-3}$. Hence, for any large constant $D'>0$, the event on which the above $\prec$-bounds fail has probability at most $N^{-D'}$ and contributes at most $N^{-D'}(CN\eta_t^{-3})^{p}$ to $\E|f_{xy}(G)|^{p}$, which is negligible once $D'$ is chosen large enough depending on $p$. This concludes \eqref{eq:LW_moment} by taking $\Psi_t=\Psi_t(0)$.
\end{proof}


\subsection{Proof of Lemma \ref{lem:Anp}}\label{sec:pf_Anp}


This subsection is devoted to establishing the key technical result, \Cref{lem:Anp}. Its proof relies crucially on the path properties (3)--(5) of $\Gamma_{\mu, xy}$ stated in \Cref{lem:localregular}, which also hold for the auxiliary graph \smash{$\Gamma^{\aux}_{\mu, [a][b]}$}. We refer to any graph satisfying these path properties as a \emph{nested graph}. 
By \Cref{def_auxgraph} of auxiliary graphs together with the bounds in \eqref{eq:Gbyxi2}, \Cref{lem:Anp} follows directly from the following combinatorial result concerning nested graphs.


\begin{lemma}[Bounding nested graphs]\label{lem:Anp_key}  
%In the setting of \Cref{lem:LWterm}, 
Let $\cal G_{\ba\bfb}$ be a graph with $2p$ external vertices, denoted by $\ba=(a_1,\ldots,a_p)\in (\Zn)^p$ and $\bfb=(b_1,\ldots,b_p)\in (\Zn)^p$, together with $0\le q\le p$ internal vertices. A solid edge between vertices $\al$ and $\beta$ represents a nonnegative random variable $\xi_{\al\beta}=\xi_{\beta\al}$ (not necessarily the specific random variables appearing in \eqref{eq:xia1a2}) that satisfies the bounds
\be\label{eq:Gbyxi3}
\xi_{\al\beta} \prec \Psi_t(|\al-\beta|), 
\quad \sum_{\beta}|\xi_{\al\beta}|^2 \prec (W^d\eta_t)^{-1}, \quad \forall \al,\beta\in \Zn.
\ee
%and ghost edges may be present.
We assume that the graph contains no self-loops. Suppose $\cal G_{\ba\bfb}$ satisfies the following path properties:
\begin{enumerate}
\item[(1)] There are $p$ edge-disjoint paths $\fP_i$, $i\in\qqq{p}$, such that each path $\fP_i$ connects the pair $a_i$ and $b_i$. % via solid and/or ghost edges. 

\item[(2)] Each internal vertex $\al_i$, for $i\in \qqq{q}$, is traversed by at least two distinct paths $\fP_k\ne \fP_l$. As a consequence, each vertex $\al_i$ has degree at least 4: 
\be\label{eq:deg_mole_aux}
\deg(\al_i)\ge 4,\quad \forall i\in \qqq{q}.\ee

\item[(3)] For any subset $A\subset \qqq{q}$, the number of paths passing through the internal vertices $\{\al_i: i\in A\}$ is at least $|A|$; that is,
$\left|\left\{j: \mathrm V(\fP_j) \cap  \{\al_i: i\in A\}\ne \emptyset\right\}\right|\ge |A| ,$
where $\mathrm V(\fP_j)$ denotes the subset of vertices in the path $\fP_j$.
\end{enumerate}
Then, there exists a constant $c>0$ such that 
\begin{equation}\label{adsuu_orig}
\cal G_{\ba\bfb} \prec (W^d\eta_t)^{-q}\cdot \Psi_t^{\ord(\cal G_{\ba\bfb})-p}\cdot \prod_{i=1}^p  \Psi_t(c|a_i-b_i|) ,
\end{equation}
where $\ord(\cal G_{\ba\bfb})$ is defined as in \eqref{eq:ordGaux}, namely \(\ord(\cal G_{\ba\bfb}):=\#\{\text{solid edges in } \cal G_{\ba\bfb}\}-2q.\)
\end{lemma}

\begin{proof}[\bf Proof of \Cref{lem:Anp}]
By \Cref{def_auxgraph} of auxiliary graphs, the bounds in \eqref{eq:Gbyxi2}, and the path properties (3)--(5) from \Cref{lem:localregular}, the auxiliary graph \smash{$\Gamma^{\aux}_{\mu, [a][b]}$} satisfies the assumptions of \Cref{lem:Anp_key}. Consequently, \eqref{eq:bddGamma_aux} becomes a special case of \eqref{adsuu_orig}, with $a_1=\cdots=a_p=[a]$ and $b_1=\cdots=b_p=[b]$. 
\end{proof}


% Before proceeding to the formal proof of \Cref{lem:Anp_key}, we briefly outline the main ideas. % and explain the motivation for introducing ghost edges. We begin by revisiting the original estimate \eqref{eq:bddGamma_aux}, which does not involve ghost edges. 

%The main challenge is then to show that, under the IPC property, a nested ordering of the internal vertices still exists---despite the removal of these long edges---such that \eqref{eq_nested} continues to hold.
% By the triangle inequality, each path contains at least one edge of length $\gtrsim |[a_i]-[b_i]|$. 
% We bound this edge by a factor $\Psi_t(c|[a_i]-[b_i]|)$, and subsequently remove it from the graph. %\footnote{This argument is not entirely rigorous, as suggested by property (4) of \Cref{lem:Anp_key}; in the actual proof, only long ending edges connected with external vertices are removed.}  

%In the proof of \Cref{lem:Anp_key}, we will construct this ordering by selecting the ``long solid edges'' and determining the nested order of the internal vertices via an inductive argument.
In the proof of \Cref{lem:Anp_key}, we will select the ``long edges'' and determine the nested summation order of internal vertices via an inductive approach. 
Specifically, we assume that \Cref{lem:Anp_key} holds for any graph with $(q-1)$ internal vertices. Then, given a graph with $q$ internal vertices that satisfies the assumptions of \Cref{lem:Anp_key}, we provide an algorithm to identify the long edges and determine the first vertex in a valid nested order, where the long edges have been removed.
We show that summing over this vertex yields the desired factors of $(W^d\eta_t)^{-1}$ and $\Psi_t(c|a_i-b_i|)$, and that the resulting graph still satisfies the assumptions of \Cref{lem:Anp_key}, now with $(q-1)$ internal vertices. (Note that in this reduced graph, the ending external vertices may change. This is precisely why we work with more general graphs allowing arbitrary external vertices, rather than restricting ourselves to the original external vertices $[a]$ and $[b]$.) Therefore, we can apply the inductive hypothesis to complete the argument.

%, in the remaining graph,
%This immediately raises one critical technical problem: how can we guarantee that throughout the summation process, 
% \be\label{eq_nested}
% \text{\emph{each internal vertex being summed over is connected to at least two solid edges throughout the summation process}}? 
% \ee
%If, at any step, we are left with only one (or zero) solid edges at a vertex, then applying \eqref{eq:Gbyxi3} together with the Cauchy-Schwarz inequality would yield a factor of either \smash{$L^{d/2}\eta_t^{-1/2}$} or $L^d$, which is too large to afford. 
%The situation becomes even more delicate because some ``long edges'' are required to contribute the $\Psi_t(c|a_i-b_i|)$ factor and therefore cannot be used in these internal vertex summations.


%In this induction argument, the selection of long edges and the choice of a nested order for internal vertices 
Our induction argument relies critically on the path properties assumed in \Cref{lem:Anp_key}. However, these properties may be violated once some long edges are removed. To clarify the graphical structure and streamline the presentation of our proof, instead of deleting the chosen long edges entirely, we replace them with a new type of edges, called \emph{ghost edges}:
\begin{itemize}
\item A {\bf ghost edge} is a black dashed edge between two vertices, representing a factor of $1$. 
\end{itemize}
% For the detailed argument, however, we need to go beyond the class of auxiliary graphs defined in \Cref{def_auxgraph} by introducing additional external vertices and a new type of edge, called a \emph{ghost edge}, defined as follows:
% \begin{itemize}
% \item A {\bf ghost edge} is a black dashed edge between two vertices, representing a factor of $1$. 
% \end{itemize}
Ghost edges do not contribute to the value of a graph (and hence do not affect its scaling order); rather, they are introduced solely to preserve the connectedness of paths in our graphs. As a consequence, they will be counted when we define the degrees of vertices, i.e., 
\[ \deg(\al)=\#\{\text{solid and ghost edges connected with }\al\}.\]
When we want to refer to the degree of solid edges for a vertex $\al$, we will use the notation $\deg_s(\al)$.  
The main advantage of ghost edges is that they allow us to maintain the essential path properties---namely properties (1)--(3) in \Cref{lem:Anp_key}---throughout the proof. In particular, they simplify the induction: without ghost edges, we would need to distinguish between paths depending on whether long edges had been removed from them. 
Now, \Cref{lem:Anp_key} is an easy corollary of the following more general result for graphs with ghost edges. 

% \begin{definition}[Ghost edges]\label{def_auxgraph_gen}
% A generalized auxiliary graph is a graph defined on multiple external and internal vertices, where every solid edge between vertices $[\al]$ and $[\beta]$ represents a $\xi([\al_i],[\al_j])$ factor, and a dashed ghost edge between vertices $[\al]$ and $[\beta]$ represents a factor of 1. Moreover, there are no self-loops inside generalized auxiliary graphs. 
% We define the scaling order of a generalized auxiliary graph $\cal G^{\aux}$ as in \eqref{eq:ordGaux}. While ghost edges will not contribute to the scaling order, they will be counted when we define the degrees of vertices, i.e., 
% \[ \deg([\al])=\#\{\text{solid and ghost edges connected with }[\al]\}.\]
% When we want to refer to the degree of solid edges for a vertex $[\al]$, we will use the notation $\deg_s([\al])$.  

% We use the notation $([\alpha], [\beta])$ to represent a solid or ghost edge between $[\alpha]$ and $[\beta]$ in our graphs. A path from $[\alpha]$ to $[\beta]$ is denoted by $([\alpha] \to [\beta])$. More generally, a path of the form
% $([\alpha_1] \to [\alpha_2] \to \cdots \to [\alpha_k])$
% indicates a sequence of edges from $[\alpha_1]$ to $[\alpha_2]$, then from $[\alpha_2]$ to $[\alpha_3]$, and so on.
% %We will use $([\al],[\beta])$ to denote a solid/ghost edge in our graphs, and use $([\al]\to[\beta])$ to denote a path from $[\al]$ to $[\beta]$. More generally, $([\al_1]\to [\al_2]\to \cdots [\al_k])$ denotes a path that goes from $[\al_1]$ to $[\al_2]$, then from $[\al_2]$ to $[\al_3]$, and so on.
% \end{definition}



\begin{lemma}\label{lem:Anp_key_gh}  
Suppose $\cal G_{\ba\bfb}$ is a nested graph that may contain ghost edges. Assume it satisfies the setting of \Cref{lem:Anp_key}, including properties (1)--(3) therein, with the convention that paths may also contain ghost edges. Furthermore, suppose that each path $\fP_i$, $i\in\qqq{p}$, contains at most one ghost edge, which---if present---must appear as an ending edge of the path. Here, the ending edges of a path $\fP_i$ are its first and last edges, necessarily incident to an external vertex $a_i$ or $b_i$. Then, there exists a constant $c>0$ such that
\begin{equation}\label{adsuu22}
\cal G_{\ba\bfb} \prec (W^d\eta_t)^{-q}\cdot \Psi_t^{\ord(\cal G_{\ba\bfb})-n_{\ngh}(\cal G_{\ba\bfb})}\cdot \prod_{i=1}^p  [\Psi_t(c|a_i-b_i|)]^{\mathbf 1(\fP_i \text{ has no ghost edge})} ,
\end{equation}
where $n_{\ngh}$ denotes the number of paths that have no ghost edge, i.e., 
\(n_{\ngh}(\cal G_{\ba\bfb}):=\sum_{i=1}^p \mathbf 1(\fP_i \text{ has no ghost edge}) .\)

%    \begin{enumerate}
        % \item There are vertices 
        % $$
        % z_1\cdots z_p, \quad a_1\cdots a_m, b_1\cdots b_m$$
    %     \item There are $m$ edge-paths, for the $k$-th path,  it starts from $a_k$,   ends at $b_k$. (Not necessary passing some $z$'s. ).   No other edges in $\Gamma$.  
    %     \item For each $A\subset [[1,p]]$, the number of the paths passing through 
    %     $\{z_i: i\in A\}$ is larger or equal to $|A|$
    %     $$
    %  \Big|j: Path(a_j, b_j)\cap  \{z_i: i\in A\}\ne \emptyset\Big|\ge |A| $$
    %    \item for each $z_k$,  (here is the difference)
    % $$\deg (z_k)\ge 4,\quad 1\le k\le p
    % $$
    %    \item Each path has at most  1 ghost edge, which  must be ending edge of some path.  The value of ghost edge  in the $k$-th path equals to 1 
    % $$
    % f^{(g)}_{ K}(\epsilon_i, \epsilon_f)=1
    % $$
    % i.e., it does not depend on the ending points.  
    % \end{enumerate}
    % For the value of this  graph is 
    % $$
    % Val (\Gamma)=\prod_{\epsilon\in E }f_K(\epsilon_i, \epsilon_f)
    % $$


% \begin{equation}\label{adsuu22}
% \sum_{z_1\cdots z_p} Val (\Gamma)\prec \eta^{-p}\cdot 
% \prod_{i=1}^m g(\|a_{i }- b_{i}\|_K)^{ P^{(i)}_{ng}},\quad P^{(i)}_{ng}:={\bf 1}(\hbox{$i$-th path has No Ghost edge}) 
% \end{equation}
\end{lemma}








% To prove this lemma, we need to extend it to more general cases. 
% \begin{definition}[Ending edges]
%     If an edge is in the path from $a $ to $b $, and it contains   $a $ and/or $b $ as the ending points, then it is an ending edge of the path. 
% \end{definition}

% \begin{lemma}\label{lem:Anp_2}   Under the assumption of previous lemma, we label some edges as special edges. We call them ghost edges. The ghost edges satisfy the following rules
% \begin{enumerate}
%     \item Each path has at most  1 ghost edge. 
%     \item The ghost edge must be ending edge of some path. 
%     \item The value of ghost edge  in the $k$-th path equals to 1 
%     $$
%     f^{(g)}_{ K}(\epsilon_i, \epsilon_f)=1
%     $$
%     i.e., it does not depend on the ending points.  
% \end{enumerate}
   
    
%   Then we claim that 
% $$\sum_{z_1\cdots z_p} Val (\Gamma)\prec \eta^{-p}\cdot 
% \prod_{i=1}^m g(\|a_{i }- b_{i}\|_K)^{ P^{(i)}_{ng}},\quad P^{(i)}_{ng}:={\bf 1}(\hbox{$i$-th path has No Ghost edge})
% $$ 
% \end{lemma}


\begin{proof}[\bf Proof of \Cref{lem:Anp_key_gh}] 
%Without loss of generality, we assume that all edges in $\cal G^{\aux}_{\ba\bfb}$ belong to the $p$ paths $\fP_i$, $i\in \qqq{p}$; otherwise, we can remove all edges that do not belong to these paths and bound them by factors of $\psi_t(0)$.
We prove the estimate \eqref{adsuu22} by induction on $q$, the number of internal vertices. Suppose we aim to establish \eqref{adsuu22} for a target graph \smash{${\cal G}_{\ba\bfb}$} satisfying the assumptions of \Cref{lem:Anp_key_gh}. Moreover, assume that \smash{${\cal G}_{\ba\bfb}$} contains $K\in \N$ many solid edges. Our induction hypothesis is the following: the estimate \eqref{adsuu22} holds for all nested graphs satisfying the assumptions of \Cref{lem:Anp_key_gh}, with $k$ internal vertices and at most $K$ solid edges. (Note that we do not assume that the number of paths is smaller than $p$.)
In the following proof, we use the notation $(\alpha, \beta)$ to denote either a solid or a ghost edge between vertices $\alpha$ and $\beta$. A path from $\alpha$ to $\beta$ is denoted by $(\alpha \too \beta)$, where the double arrow indicates that intermediate vertices may be present along the path.
More generally, a path of the form
$(\alpha_1 \too \alpha_2 \too \cdots \too \alpha_k)$ represents a sequence of edges from $\alpha_1$ to $\alpha_2$, then $\alpha_2$ to $\alpha_3$, and so on.

First, the estimate \eqref{adsuu22} is trivial in the case $q=0$, by \eqref{eq:Gbyxi3}. Now, let $1\le q\le p$, and assume the induction hypothesis holds for all $0\le k \le q-1$. We prove that \eqref{adsuu22} holds for every graph satisfying the assumptions of \Cref{lem:Anp_key_gh}, with $q$ internal vertices and at most $K$ solid edges.
Given such a graph $\cal G_{\ba\bfb}$, we partition the summation region over $\bal=(\al_1,\ldots,\al_q)$ into at most $2^{pq}$ subregions, according to whether each vertex $\al_i$ is closer to $a_j$ or to $b_j$ with $i\in\qqq{q}$ and $j\in \qqq{p}$:
\[ \bD_{\bpi}\subset (\Zn)^q:=\left\{\bal: |\al_i-a_j|\le |\al_i-b_j| \ \text{if}\ \pi_{i,j}=0, \text{ and } |\al_i-a_j| > |\al_i-b_j| \ \text{if}\ \pi_{i,j}=1\right\},\]
for $\bpi=(\pi_{1,1},\ldots, \pi_{1,p},\ldots, \pi_{q,1},\ldots, \pi_{q,p}) \in \{0,1\}^{pq}$. The union of these subregions covers the entire space \smash{$(\Zn)^q$}. Then, it suffices to prove that for each fixed $\bpi \in \{0,1\}^{pq}$,  
\begin{align}\label{kwuyayw}
        \sum_{\bal\in \bD_{\bpi}} \cal G_{\ba\bfb}(\bal) \prec (W^d\eta_t)^{-q}\cdot \Psi_t^{\ord(\cal G_{\ba\bfb})-n_{\ngh}(\cal G_{\ba\bfb})}\cdot \prod_{i=1}^p  [\Psi_t(c|a_i-b_i|)]^{\chi(\fP_i)} ,
        %\\(W^d\eta_t)^{-q}\cdot [\Psi_t(c|[a]-[b]|)]^{n_{\ngh}(\cal G^{\aux}_{[a][b]})} \cdot [\Psi_t(0)]^{\ord(\cal G^{\aux}_{[a][b]})-n_{\ngh}(\cal G^{\aux}_{[a][b]})}. 
    \end{align} 
where $\cal G_{\ba\bfb}(\bal)$ denotes the graph obtained by fixing the internal vertices in $\bal$ as external vertices, and we abbreviate $\chi(\fP_i)\equiv \mathbf{1}\left(\fP_i \text{ has no ghost edge}\right)$. 

For the proof of \eqref{kwuyayw}, we look at the ending edges of the $p$ paths. Suppose $(a_j,\al_i)$ is an ending edge of a path $\fP_j$ (note that $\fP_j$ connects $a_j$ to $b_j$ through $\al_i$, but there may be no direct solid edge between $b_j$ and $\al_i$). On the region $\bD_{\bpi}$, the ending edge $(a_j,\al_i)$ falls into one of the following four categories: 
\begin{enumerate}
    \item[{\bf A1}:] The path $\fP_j=(a_j\too b_j)$ has no ghost edge. Moreover, for $\bal\in \bD_{\bpi}$, the vertex $\al_i$ is closer to $a_j$ than to $b_j$, i.e., $|\al_i-a_j|\le  |\al_i-b_j|$. In particular, this implies $|\al_i-b_j|\ge |a_j-b_j|/2$. 
    
    \item[{\bf A2}:] The path $\fP_j=(a_j\too b_j)$ has no ghost edge. Moreover, for $\bal\in \bD_{\bpi}$, the vertex $\al_i$ is closer to $b_j$ than to $a_j$, i.e., $|\al_i-a_j| > |\al_i-b_j|$. In particular, this implies $|\al_i-a_j|\ge |a_j-b_j|/2$. 
    %The path $a_k\sim b_k$ has no ghost edge, and under the restriction of $\pi$, this $z_i$ is always closer to $b_k$ than to $a_k$,     i.e., $$\|z_i-a_k\|\ge  \|z_i-b_k\|,\quad  i.e.,\quad  \|z_i-a_k\|\ge \frac12\|a_i-b_i\|$$
    
     \item[{\bf B1}:] The path $\fP_j=( a_j \too b_j )$ contains a ghost edge, which is not $( a_j , \al_i )$. 
     
      \item[{\bf B2}:] The path $\fP_j=( a_j \too b_j )$ contains a ghost edge $( a_j , \al_i )$.
      
\end{enumerate}
By symmetry, we may also classify ending edges attached to $b_j$ by switching the roles of $a_j$ and $b_j$ in the definitions above. The only difference is in the classification of type-A1 (resp.~type-A2) edges: in this case, we require $|\al_i-b_j|< |\al_i-a_j|$ (resp.~$|\al_i-b_j|\ge |\al_i-a_j|$) instead.  


First, we notice that if an ending edge $\p{a_j,\al_i}$ is a type-A2 edge, then it contributes a factor $\Psi_t(|a_j-b_j|/2)$ by \eqref{eq:Gbyxi3}. In this case, we can bound the original graph $\cal G_{\ba\bfb}$ by the product of a new graph \smash{$\wt{\cal G}_{\ba\bfb}$}, obtained by replacing the edge $\p{a_j,\al_i}$ with a ghost edge, and the factor $\Psi_t(|a_j-b_j|/2)$. It is easy to see that if the bound \eqref{kwuyayw} holds for \smash{$\wt{\cal G}_{\ba\bfb}$}, then it also holds for the original graph. 
We perform this replacement for every A2 edge that appears in the paths $\fP_i$, $i\in\qqq{p}$.  %If a path contains two type-A2 edges, we replace only one of them with a ghost edge. 
After performing all such replacements, it suffices to prove \eqref{kwuyayw} for the resulting graph.
With a slight abuse of notation, we continue to denote this modified graph by $\cal G_{\ba\bfb}$, which now satisfies the condition: 
\be\label{eq:noA2}
\text{there is no A2 edge in }\cal G_{\ba\bfb}.
\ee
We now proceed to prove the estimate \eqref{kwuyayw}, assuming \eqref{eq:noA2} and the induction hypothesis. The proof is organized by classifying cases according to the number and type of ending edges attached to internal vertices.

\medskip
\paragraph{\bf (I) Two A1/B1 edges: the simple case}

We first consider the case where an internal vertex is connected to two or more ending edges of type A1 or B1. That is, suppose there exists an internal vertex in the graph connected to at least two such edges---either two A1 edges, two B1 edges, or one of each. Without loss of generality, let this vertex be $\al_q$. To illustrate the idea, we begin with a simple scenario where $\deg(\al_q)=4$, so exactly two paths pass through $\al_q$. We will address the more general case $\deg(\al_q)\ge 4$ in Case (II). 

By definition, the two ending edges connected to $\al_q$ must belong to distinct paths. 
Without loss of generality, assume these edges are $(a_1,\al_q)$ and $(a_2,\al_q)$, which belong to paths $\fP_1$ and $\fP_2$, respectively. We now define a new graph \smash{$\cal G^{\mathrm{new}}_{\ba'\bfb}\equiv \cal G^{\mathrm{new}}_{\ba'(\al_q),\bfb}$} as follows: it has external vertices 
\[\ba'=(a_1',a_2',a_3,\ldots,a_p,a_1)\quad \text{and}\quad \bfb=(b_1,\ldots,b_p,a_2), \quad \text{where}\ \ a_1'=\al_q, \ a_2'=\al_q,\] 
and internal vertices $\bal'=(\al_1,\ldots,\al_{q-1})$, i.e., we remove $\al_q$ from the internal vertex set. 
First, all edges not belonging to paths $\fP_1$ and $\fP_2$ remain unchanged, so the paths $\fP_i$ for $i\in\qqq{3,p}$ stay the same in $\cal G^{\mathrm{new}}_{\ba'\bfb}$. Second, we define new paths $\fP'_1=(\al_q\too b_1)$ and $\fP_2'=(\al_q\too b_2)$ obtained by removing the edges $(a_1,\al_q)$ and $(a_2,\al_q)$ from $\fP_1$ and $\fP_2$, respectively. 
Finally, we introduce an auxiliary path $(a_1\too a_2)$ consisting of a single ghost edge between $a_1$ and $a_2$. The purpose of this auxiliary path is to ensure consistency with the induction hypothesis, which requires that every external vertex is an ending point of some path. (We cannot simply remove $a_1$ and $a_2$ from the graph, since some path $\fP_i$ may pass through them.)

The new graph $\cal G^{\mathrm{new}}_{\ba'\bfb}$ produced by this construction also satisfies the assumptions of \Cref{lem:Anp_key_gh}, but it has one fewer internal vertex and two fewer solid edges than the original graph $\cal G_{\ba\bfb}$. Hence, by the induction hypothesis, we can bound it by 
\begin{align}
\sum_{\bal'} \cal G^{\mathrm{new}}_{\ba'\bfb}(\bal') \prec &~ \frac{ \Psi_t^{\ord(\cal G^{\mathrm{new}}_{\ba'\bfb})-n_{\ngh}(\cal G^{\mathrm{new}}_{\ba'\bfb})}}{(W^d\eta_t)^{q-1}}\cdot  \prod_{i=1}^2[\Psi_t(c|\al_q-b_i|)]^{\chi(\fP_i')} \cdot \prod_{i=3}^p  [\Psi_t(c|a_i-b_i|)]^{\chi(\fP_i)} \nonumber\\
\prec &~ (W^d\eta_t)^{-(q-1)}\cdot \Psi_t^{\ord(\cal G_{\ba\bfb})-n_{\ngh}(\cal G_{\ba\bfb})}\cdot  \prod_{i=1}^p  [\Psi_t(c|a_i-b_i|/2)]^{\chi(\fP_i)} \label{eq:induc_Ggraph}
\end{align}
for a constant $c>0$. In the second step, we use that 
\be\label{eq:change_of_order}\ord(\cal G^{\mathrm{new}}_{\ba'\bfb})=\ord(\cal G_{\ba\bfb}),\quad n_{\ngh}(\cal G_{\ba\bfb})=n_{\ngh}(\cal G^{\mathrm{new}}_{\ba'\bfb}),\ee 
%= \ord(\cal G^{\aux}_{\ba\bfb})-n_{\ngh}(\cal G^{\aux}_{\ba\bfb}) \ee in all cases (i.e., two $A1$ edges, two $B1$ edges, or an $A1$ and a $B1$ edge) 
since $\al_q$ is now treated as an external vertex. Moreover, we have also used that 
\be\label{eq:change_of_order2}\prod_{i=1}^2[\Psi_t(c|\al_q-b_i|)]^{\chi(\fP_i')}\le \prod_{i=1}^2[\Psi_t(c|a_i-b_i|/2)]^{\chi(\fP_i)} \ee
by the definitions of A1 and B1 edges. With \eqref{eq:induc_Ggraph}, we can bound the LHS of \eqref{kwuyayw} by 
\begin{align}\label{kwuyayw_case1}
\sum_{\bal\in \bD_{\bpi}} \cal G_{\ba\bfb}(\bal) &\prec (W^d\eta_t)^{-(q-1)}\cdot \Psi_t^{\ord(\cal G_{\ba\bfb})-n_{\ngh}(\cal G_{\ba\bfb})}\cdot \prod_{i=1}^p  [\Psi_t(c|a_i-b_i|/2)]^{\chi(\fP_i)} \sum_{\al_q}\xi_{a_1\al_q}\xi_{a_2\al_q} \nonumber\\
&\prec (W^d\eta_t)^{-q}\cdot \Psi_t^{\ord(\cal G_{\ba\bfb})-n_{\ngh}(\cal G_{\ba\bfb})}\cdot  \prod_{i=1}^p  [\Psi_t(c|a_i-b_i|/2)]^{\chi(\fP_i)} ,
\end{align}
where we use the Cauchy-Schwarz inequality and \eqref{eq:Gbyxi3} in the second step. This concludes \eqref{kwuyayw} in Case (I). 


%Key: Add in version 2 {\cor Draw a graph to illustrate the above proof.}
 
% It has the same $k$-th, $k\ge 3$,  paths as those in $\Gamma$. The $i$-th path ($i = 1, 2$) of $\Gamma'$ is the segment from $z_p$ to $b_i$ of the corresponding $i$-th path of $\Gamma$. In another word, in the new $\Gamma'$, $a_1'=a_2'=z_p$. Then 
% $$\sum_{z_1\cdots z_p} \chi_\pi \cdot Val (\Gamma)\prec \sum_{z_p} f_K(z_p,a_1)\cdot f_K(z_p,a_2)\cdot \left(\sum_{z_1\cdots z_{p-1}}Val (\Gamma')\right)
% $$
% It is easy to check that $\Gamma'$ satisfies the assumptions in lemma \ref{lem:Anp} and \ref{lem:Anp_2}.  Note, the follow quantities clearly did not change from $\Gamma$ to $\Gamma'$
% \begin{enumerate}
%     \item The degree  of $z_i$ ($1\le i\le p-1$) 
%     \item Whether a path has ghost edge or not. 
%     \item The number of paths pass through $\{z_i: i\in A\subset [[1,p-1]]\}$ for any fixed $A$
% \end{enumerate}
% By induction, and for $i=1,2$, 
% $$
% \|z_p-b_i\|\ge 1/2 \|a_i-b_i\|, 
% $$
% we have 
% $$\left(\sum_{z_1\cdots z_{p-1}}Val (\Gamma')\right)
% \prec  \eta^{-p+1}\cdot  \prod_{s=1}^m g(\|a_{i }- b_{i}\|_K) ^{P_{ng}^{(i)}}$$
% and 
% $$\sum_{z_p} f_K(z_p,a_1)\cdot f_K(z_p,a_2)\prec\eta^{-1}
% $$
% Therefore, in this case \eqref{kwuyayw} holds.

\medskip
\paragraph{\bf (II) Two A1/B1 edges: general case} 

We now consider a more general case than Case (I), where there is an internal vertex, say $\al_q$, connected to at least two ending edges of type A1 or B1, and where $\deg(\al_q)\ge 4$. Without loss of generality, assume that the two ending edges connected to $\al_q$ are $(a_1,\al_q)$ and $(a_2,\al_q)$, belonging to paths $\fP_1$ and $\fP_2$, respectively. As before, we define the new graph \smash{$\cal G^{\new}_{\ba'\bfb'}\equiv \cal G^{\new}_{\ba'(\al_q),\bfb'(\al_q)}$} obtained from \smash{$\cal G_{\ba\bfb}$} by removing the two edges $(a_1,\al_q)$ and $(a_2,\al_q)$. 

The new graph has external vertices \(\al_q, a_1,a_2, a_3, \ldots, a_p,  b_1, \ldots,  b_p\) and internal vertices $\al_1,  \ldots,  \al_{q-1}$. To apply the induction hypothesis, we need to ensure that the structure of the paths in $\cal G^{\new}_{\ba' \bfb'}$ satisfies the assumptions in \Cref{lem:Anp_key_gh}.  
For each path $\fP_j$, $j\in \qqq{3,p}$, in the original graph, suppose it takes the form 
\be\label{eq:jthpath}( a_j \too  \al_q \too \al_q \too \cdots \too  \al_q \too b_j ),\ee 
where each occurrence of $\alpha_q$ is explicitly listed. Here, $(a_j \too \alpha_q)$ denotes the initial segment of $\fP_j$ connecting $a_j$ to the first occurrence of $\alpha_q$, while $(\alpha_q \too b_j)$ denotes the final segment connecting the last occurrence of $\alpha_q$ to $b_j$. Each intermediate segment $(\alpha_q \too \alpha_q)$ represents a connection between two consecutive appearances of $\alpha_q$ along the path. Let $k_j$ denote the total number of times $\alpha_q$ appears in the path $\fP_j$. 
Correspondingly, in $\cal G^{\mathrm{new}}_{\ba'\bfb'}$, we define the following $k_j+1$ paths for each $j$: for each $1 \leq r \leq k_j + 1$, the $(j,r)$-th path $\fP'_{j,r}$ is the $r$-th segment of $\fP_j$, connecting $a_{j,r}$ to $b_{j,r}$, where  
\begin{equation}
a_{j,r}:=  \begin{cases} a_j, & \text{if} \ r=1\\
\al_q, & \text{if} \ r>1
    \end{cases},\quad  b_{j,r}:=  \begin{cases} \al_q,  &\text{if} \ r\le k_j\\
    b_j, &\text{if} \ r= k_j +1
    \end{cases}. % j=1,2\cdots m_k+1
\end{equation}
Next, for the first two paths $\fP_j$, $j \in \{1, 2\}$, assume they also take the form \eqref{eq:jthpath}, with $k_j$ appearances of $\al_q$. After removing the edges $(a_1,\al_q)$ and $(a_2,\al_q)$, we define $k_j$ paths in $\cal G^{\mathrm{new}}_{\ba'\bfb'}$ for each $j\in\{1,2\}$. Specifically, for each $1 \leq r \leq k_j$, the $(j,r)$-th path $\fP'_{j,r}$ is the $r$-th segment of $\fP_j$, connecting $a_{j,r}$ to $b_{j,r}$, where  
\begin{equation}
a_{j,r}:= \al_q,\quad  b_{j,r}:=  \begin{cases} \al_q,  &\text{if} \ r< k_j\\
    b_j, &\text{if} \ r= k_j
    \end{cases}. % j=1,2\cdots m_k+1
\end{equation}
Finally, we introduce an auxiliary path $(a_1\too a_2)$ consisting of a single ghost edge between $a_1$ and $a_2$. 

%Key: Add in version 2 {\cor Draw a graph to illustrate the above construction.}
 
Now, we let $\ba':=\{a_{j,r}:1\le j \le p, 1\le r \le k_j+\mathbf 1_{j\ge 3}\}\cup\{a_1\}$ and $\bfb':=\{b_{j,r}:1\le j \le p, 1\le r \le k_j+\mathbf 1_{j\ge 3}\}\cup\{a_2\}$. It is easy to see that the graph $\cal G^{\new}_{\ba'\bfb'}$ defined in this way satisfies the assumptions in \Cref{lem:Anp_key_gh}, but with one fewer internal vertex and two fewer solid edges than the original graph $\cal G_{\ba\bfb}$. 
Therefore, using the induction hypothesis, we can bound it in a similar way as \eqref{eq:induc_Ggraph}, that is,
\begin{align}\label{eq:bound_new_graph}
\sum_{\bal'} \cal G^{\mathrm{new}}_{\ba'\bfb'}(\bal') \prec (W^d\eta_t)^{-(q-1)}\cdot \Psi_t^{\ord(\cal G_{\ba\bfb})-n_{\ngh}(\cal G_{\ba\bfb})}\cdot  \prod_{i=1}^p  [\Psi_t(c|a_i-b_i|/2)]^{\chi(\fP_i)} ,
\end{align}
where $\bal'=(\al_1,\ldots,\al_{q-1})$, and in the derivation of this bound, we have used \eqref{eq:change_of_order} and \eqref{eq:change_of_order2}, and that 
$$
  \prod_{r=1}^{k_j+1} [\Psi_t(c|a_{j,r}- b_{j,r}|)]^{\chi(\fP'_{j,r})} \le [\Psi_t (c|a_{j}- b_{j}|/2)]^{\chi(\fP_{j})} [\Psi_t(0)]^{\sum_{r=1}^{k_j+1}\chi(\fP'_{j,r})- \chi(\fP_{j})},\quad \forall j\in\qqq{3,p}.
 $$
Finally, using \eqref{eq:bound_new_graph}, we can complete the proof of \eqref{kwuyayw} for Case (II), following a similar argument to the one used in \eqref{kwuyayw_case1}.

%With \eqref{eq:bound_new_graph}, using a similar argument as in \eqref{kwuyayw_case1}, we complete the proof of \eqref{kwuyayw} in the current case. 

\medskip
\paragraph{\bf (III) Two B2 edges}
In this case, we assume there exists an internal vertex, say $\al_q$, that is connected to two B2 ending edges and two solid edges, i.e., $\deg_s(\al_q)=2$. By definition, these B2 edges do not belong to the same path. Without loss of generality, assume that these two B2 edges are $(a_1, \al_q)$ and $(a_2, \al_q)$, belonging to paths $\fP_1$ and $\fP_2$, respectively. 
Additionally, suppose the two solid edges connected to $\al_q$ are $(\al_q, \beta_1)$ and $(\al_q, \beta_2)$, which belong to paths $\fP_1$ and $\fP_2$, respectively. Here, $\beta_1$ and $\beta_2$ may represent external vertices in $\ba, \bfb$ or internal vertices in $\bal' = (\al_1, \ldots, \al_{q-1})$.
 

We now define a new graph $\cal G^{\new}_{\ba\bfb}$ as follows: it still has the external vertices $\ba$ and $\bfb$, and the internal vertices $\bal'$. The edges that do not belong to the paths $\fP_1$ and $\fP_2$ remain unchanged in the new graph $\cal G^{\mathrm{new}}_{\ba\bfb}$. For paths $\fP_1$ and $\fP_2$, we modify them as follows:
\begin{itemize}
\item We create a new path $\fP_1'=(a_1 \too b_1)$ obtained by removing the solid edge $(\al_q, \beta_1)$ from $\fP_1$ and adding a ghost edge $(a_1, \beta_1)$ to the path.

\item Similarly, we create a new path $\fP_2'=(a_2 \too b_2)$ obtained by removing the solid edge $(\al_q, \beta_2)$ from $\fP_2$ and adding a ghost edge $(a_2, \beta_2)$ to the path.
\end{itemize}
The reader can refer to \Cref{Fig:construct2B2} for an illustration of this construction. 

  \begin{figure}[h]
    \centering
    \begin{tikzpicture}

    \coordinate (a1) at (0, 0.5);
    \coordinate (a2) at (0, -0.5);
    \coordinate (al) at (1, 0);
    \coordinate (b1) at (2, 0.5);
    \coordinate (b2) at (2, -0.5);
    \coordinate (c1) at (3, 1);
    \coordinate (c2) at (3, -1);
    \coordinate (d1) at (5, 0.5);
    \coordinate (d2) at (5, -0.5);
    \coordinate (L1) at (6, 0.5);
    \coordinate (L2) at (6, -0.5);
    
      \fill (a1) circle (1.5pt) node[above=2pt]{$a_1$};
      \fill (a2) circle (1.5pt) node[below=2pt]{$a_2$};
    \fill (b1) circle (1.5pt)  node[above=2pt]{$\beta_1$};
      \fill (b2) circle (1.5pt) node[below=2pt]{$\beta_2$};
       \fill (al) circle (1.5pt) node[above=2pt]{$\al_q$};
      \fill (c1) circle (1.5pt) node[right=16pt]{$\fP_1$};
      \fill (c2) circle (1.5pt) node[right=16pt]{$\fP_2$};
      \fill (d1) circle (1.5pt);
      \fill (d2) circle (1.5pt);
      \fill (L1) circle (1.5pt) node[above=2pt]{$b_1$};
      \fill (L2) circle (1.5pt) node[below=2pt]{$b_2$};

      
      \draw (a1) [dashed]-- (al); 
      \draw (al) -- (b1) -- (c1); 
      \draw (c1) [dotted]-- (d1); 
      \draw (d1) -- (L1); 
      \draw (a2) [dashed]-- (al);
      \draw (al) -- (b2) -- (c2); 
      \draw (c2) [dotted]-- (d2); 
      \draw (d2) -- (L2); 
    \end{tikzpicture}
    \hspace{10pt}
    \begin{tikzpicture}

    \coordinate (a1) at (0.5, 0.5);
    \coordinate (a2) at (0.5, -0.5);
    \coordinate (b1) at (2, 0.5);
    \coordinate (b2) at (2, -0.5);
    \coordinate (c1) at (3, 1);
    \coordinate (c2) at (3, -1);
    \coordinate (d1) at (5, 0.5);
    \coordinate (d2) at (5, -0.5);
    \coordinate (L1) at (6, 0.5);
    \coordinate (L2) at (6, -0.5);
    
      \fill (a1) circle (1.5pt) node[above=2pt]{$a_1$};
      \fill (a2) circle (1.5pt) node[below=2pt]{$a_2$};
    \fill (b1) circle (1.5pt)  node[above=2pt]{$\beta_1$};
      \fill (b2) circle (1.5pt) node[below=2pt]{$\beta_2$};
  
      \fill (c1) circle (1.5pt) node[right=16pt]{$\fP_1'$};
      \fill (c2) circle (1.5pt) node[right=16pt]{$\fP_2'$};
      \fill (d1) circle (1.5pt);
      \fill (d2) circle (1.5pt);
      \fill (L1) circle (1.5pt) node[above=2pt]{$b_1$};
      \fill (L2) circle (1.5pt) node[below=2pt]{$b_2$};

      \draw (a1) [dashed]-- (b1);
      \draw (b1) -- (c1); 
      \draw (c1) [dotted]-- (d1); 
      \draw (d1) -- (L1); 
      \draw (a2) [dashed]-- (b2);
      \draw (b2) -- (c2); 
      \draw (c2) [dotted]-- (d2); 
      \draw (d2) -- (L2); 
    \end{tikzpicture}
    \caption{Illustration of the construction of paths $\fP_1'$ and $\fP_2'$ (right panel) from $\fP_1$ and $\fP_2$ (left panel).}\label{Fig:construct2B2}
  \end{figure}
  
%On the other hand, we define paths $\fP'_1=([\al_q]\to [b_1])$ and $\fP_2'=([\al_q]\to [b_2])$ in the new graph as the paths obtained by removing the edges $([\al_q],[\beta_{1}])$ and $([\al_q],[\beta_{2}])$, and then adding the ghost edges $([a_1],[\beta_1])$ and $([a_2],[\beta_2])$ to the paths $\fP_1$ and $\fP_2$, respectively. 



Note that $\cal G^{\mathrm{new}}_{\ba\bfb}$ defined above is also a graph satisfying the assumptions in \Cref{lem:Anp_key_gh}, but with one fewer internal vertex and two fewer solid edges compared to the original graph $\cal G_{\ba\bfb}$. Therefore, by the induction hypothesis, we can bound it by 
\begin{align*}
\sum_{\bal'} \cal G^{\mathrm{new}}_{\ba\bfb}(\bal') \prec &~ (W^d\eta_t)^{-(q-1)}  \cdot \Psi_t^{\ord(\cal G^{\mathrm{new}}_{\ba\bfb})-n_{\ngh}(\cal G^{\mathrm{new}}_{\ba\bfb})} \cdot  \prod_{i=3}^p  [\Psi_t(c|a_i-b_i|)]^{\chi(\fP_i)} \nonumber\\
= &~ (W^d\eta_t)^{-(q-1)}\cdot \Psi_t^{\ord(\cal G_{\ba\bfb})-n_{\ngh}(\cal G_{\ba\bfb})}\cdot  \prod_{i=1}^p  [\Psi_t(c|a_i-b_i|)]^{\chi(\fP_i)}, \end{align*}
where, in the second step, we use that $\ord(\cal G^{\mathrm{new}}_{\ba\bfb})=\ord(\cal G_{\ba\bfb})$, $n_{\ngh}(\cal G^{\mathrm{new}}_{\ba\bfb})=n_{\ngh}(\cal G_{\ba\bfb})$, and $\chi(\fP_1)=\chi(\fP_2)=0$. 
Thus, we can bound the LHS of \eqref{kwuyayw} as follows: 
\begin{align}\label{kwuyayw_case3}
\sum_{\bal\in \bD_{\bpi}} \cal G_{\ba\bfb}(\bal) &\prec \sum_{\bal'} \cal G^{\mathrm{new}}_{\ba\bfb}(\bal')\sum_{\al_q}\xi_{\al_q\beta_1}\xi_{\al_q\beta_2} \prec (W^d\eta_t)^{-1}\sum_{\bal'} \cal G^{\mathrm{new}}_{\ba\bfb}(\bal') \nonumber\\
&\prec (W^d\eta_t)^{-q}\cdot \Psi_t^{\ord(\cal G_{\ba\bfb})-n_{\ngh}(\cal G_{\ba\bfb})}\cdot  \prod_{i=1}^p  [\Psi_t(c|a_i-b_i|/2)]^{\chi(\fP_i)} ,
\end{align}
where we use the Cauchy-Schwarz inequality and \eqref{eq:Gbyxi3} in the second step. This implies \eqref{kwuyayw} in Case (III). 





% For this case, we also define a new graph $\Gamma'$, whose vertices are 
% $$
% z_1\cdots z_{p-1}, \quad a_1,a_2, a_3\cdots a_m, \quad b_1,\cdots b_m
% $$
% It has the same $k$-th ($k\ge 3$) paths as those in $\Gamma$. For the $i$-th ($i=1,2$) path in $\Gamma'$, we define them as follows. First let's say $z_p$ connecting with $x_1$ and $x_2$ besides $a_1$ and $a_2$ in the 1st and 2nd path in  $\Gamma$, i.e,. 
% $$
% x_1-z_p-a_1, \quad x_2-z_p-a_2
% $$
% In $\Gamma'$, we keep the  $x_i$-to-$b_i$ part of the $i-th$ path in $\Gamma$. Then add two ghost edges  $\overline{x_ia_i}$ 


% Then   
% $$\sum_{z_1\cdots z_p} \chi_\pi \cdot Val (\Gamma)\prec \sum_{z_1\cdots z_{p-1}}\left(Val (\Gamma')\cdot\sum_{z_p} f_K(z_p,x_1)\cdot f_K(z_p,x_2) \right)
% $$
% We have: 
% $$\sum_{z_p} f_K(z_p,x_1)\cdot f_K(z_p,x_2)\prec\eta^{-1}
% $$
% Note, the follow quantities clearly did not change from $\Gamma$ to $\Gamma'$

% \begin{enumerate}
%     \item The degree  of $z_i$ ($1\le i\le p-1$) 
%     \item Whether a path has ghost edge or not. 
%     \item The number of paths pass through $\{z_i: i\in A\subset [[1,p-1]]\}$ for any fixed $A$
% \end{enumerate}
% $$\left(\sum_{z_1\cdots z_{p-1}}Val (\Gamma')\right)
% \prec  \eta^{-p+1}\cdot  \prod_{s=1}^m g(\|a_{i }- b_{i}\|_K) ^{P_{ng}^{(i)}}$$
% Therefore, in this case \eqref{kwuyayw} holds.




\paragraph{\bf (IV) The remaining case}


Finally, suppose that none of Cases (I)--(III) applies. In this case, we prove \eqref{kwuyayw} directly, without using the induction hypothesis. (Note that we cannot argue by counting the ending edges attached to internal vertices: a path may pass through external vertices, so that some of its ending edges are not incident to any internal vertex, and the number $q$ of internal vertices may be smaller than $p$.)

First, let $K$ be the number of edges of $\cal G_{\ba\bfb}$. We further partition $\bD_{\bpi}$ into at most $K^p$ subregions according to which edge of each path $\fP_i$ without a ghost edge is the longest one; it suffices to prove \eqref{kwuyayw} with $\bD_{\bpi}$ replaced by any of these subregions. On such a subregion, by the triangle inequality, the chosen edge of $\fP_i$ has length at least $|a_i-b_i|/K$, so it contributes a factor $\Psi_t(c|a_i-b_i|)$ with $c=1/K$ by \eqref{eq:Gbyxi3}. We bound the chosen edge of each path without a ghost edge in this way and remove it from the graph, and we also remove all ghost edges. Let $\cal H$ denote the resulting graph, which contains only solid edges, and let $n_s(\cal H)=\ord(\cal G_{\ba\bfb})+2q-n_{\ngh}(\cal G_{\ba\bfb})$ be the number of its solid edges. Since all graph values are nonnegative, we may then sum over all $\bal\in(\Zn)^q$, so \eqref{kwuyayw} follows from
\begin{align}\label{kwuyayw_ng}
        \sum_{\bal\in (\Zn)^q} \cal H(\bal) \prec (W^d\eta_t)^{-q}\cdot \Psi_t^{n_s(\cal H)-2q}.
    \end{align}
The key to proving \eqref{kwuyayw_ng} is to construct a \emph{nested order of summation} for the $q$ internal vertices, as in \cite{yang2021random}: when we sum over the internal vertices in this order, each vertex being summed over must be incident to at least two solid edges whose other endpoints are external vertices or internal vertices that have not yet been summed over. Summing over such a vertex yields a factor $(W^d\eta_t)^{-1}$ by the Cauchy-Schwarz inequality and \eqref{eq:Gbyxi3}, while all other solid edges incident to it are bounded by $\Psi_t$. Hence, the existence of such an order implies \eqref{kwuyayw_ng}.

We first observe that, for every nonempty subset $A$ of the internal vertices,
\be\label{eq:degali}
e_{\cal H}(A,A^c)\ge |A| \quad \text{and}\quad \#\{\text{solid edges of }\cal H\text{ incident to }A\}\ge 2|A|,
\ee
where $A^c$ is the set of all other (internal or external) vertices, and $e_{\cal H}(A,A^c)$ is the number of solid edges of $\cal H$ between $A$ and $A^c$. Indeed, removing one edge from each path $\fP_i$ splits it into two edge-disjoint segments, each of which starts at an external vertex ($a_i$ or $b_i$), and every vertex of $\fP_i$ lies on one of them. If $\fP_i$ passes through $A$, then one of its two segments enters $A$ from $A^c$, which gives a solid edge between $A$ and $A^c$; since the paths are edge-disjoint, the first bound in \eqref{eq:degali} follows from property (3) of \Cref{lem:Anp_key}. For the second bound, by property (2) of \Cref{lem:Anp_key}, every internal vertex $\al\in A$ lies on segments of at least two distinct paths, and each of these segments contains the edge through which it reaches $\al$ for the first time; these edges are incident to $A$ and are distinct for distinct pairs (segment, vertex).

We now construct the order backwards, i.e., we first determine the vertex that is summed over last. Let $\mathscr A$ denote the set consisting of the external vertices and of the internal vertices chosen so far. As long as there is an internal vertex $\al\notin\mathscr A$ with at least two solid edges to $\mathscr A$, we add it to $\mathscr A$. Suppose that this procedure stops while the set $A$ of the remaining internal vertices is nonempty. Then, each $\al\in A$ has at most one solid edge to $A^c=\mathscr A$, so by the first bound in \eqref{eq:degali} it has exactly one such edge, which we denote by $(\al,x_\al)$. Let $\cal C$ be a connected component of the subgraph of $\cal H$ induced on $A$. Applying the second bound in \eqref{eq:degali} to $\cal C$, we see that $\cal C$ contains at least $|\cal C|$ edges, and hence a cycle (possibly consisting of two parallel edges). We pick two vertices $\al,\beta\in\cal C$ that are adjacent on this cycle and use
\be\label{kwuyayw_ng_tree}
\xi_{\al x_\al}\,\xi_{\beta x_\beta}\le \frac12\left(\xi_{\al x_\al}^2+\xi_{\beta x_\beta}^2\right) .
\ee
Consider the graph in which $\xi_{\al x_\al}\xi_{\beta x_\beta}$ is replaced by $\xi_{\al x_\al}^2$ (the other term is treated in the same way, with the roles of $\al$ and $\beta$ exchanged). In this graph, $\al$ has two edges to $\mathscr A$, so it can be added to $\mathscr A$. Next, every vertex of $\cal C\setminus\{\beta\}$ that is adjacent to $\mathscr A$ has two edges to $\mathscr A$ (its edge to $x_\cdot$ and an edge to an added vertex), so all vertices connected to $\al$ in $\cal C\setminus\{\beta\}$ can be added one by one. These include both neighbors of $\beta$ on the cycle, so $\beta$ can be added as well, and then the remaining vertices of $\cal C$ can be added one by one along paths in $\cal C$ starting from $\beta$. Applying \eqref{kwuyayw_ng_tree} to one such pair in each connected component of the subgraph induced on $A$, we obtain $\OO(1)$ many graphs, for each of which the procedure exhausts all internal vertices; see \Cref{Fig:construct_trees} for an illustration. Since \eqref{kwuyayw_ng_tree} does not change the number of solid edges, summing over the internal vertices in the reverse order of their addition proves \eqref{kwuyayw_ng}, and hence \eqref{kwuyayw} in Case (IV). (In fact, this argument does not use the assumption that Cases (I)--(III) do not apply.)
\begin{figure}[h]
    \centering
    \begin{tikzpicture}[scale=2, every node/.style={circle, fill=black, inner sep=1.5pt}]
  % Nodes
 \node (A) at (0,0) {};       % root
\node[anchor=south,draw=none,fill=none] at (A) {$\al_1$};
  \node (B) at (0,-0.5) {};  
  \node[anchor=north,draw=none,fill=none] at (B) {$\al_3$};
  \node (C) at (-0.7,-0.9) {};     
  \node (D) at (-1,-1.5) {};  
  \node (E) at (-0.5,-1.5) {};    
  \node (F) at (0.7,-0.9) {};  
  \node[anchor=north,draw=none,fill=none] at (F) {$\al_2$};

\node (A1) [left of = A, distance=4pt] {};
\node[anchor=south,draw=none,fill=none] at (A1) {$a_1$};
\node (B1) [left of = B, distance=4pt] {};  
\node (C1) [left of = C, distance=4pt] {};
\node (D1) [left of = D, distance=4pt] {};
\node (E1) [right of = E, distance=4pt] {};
%\node (F1) [right of = F, distance=4pt] {};
%\node[anchor=north,draw=none,fill=none] at (F1) {$[a_2]$};

  % Black edges (tree structure)
  \draw[black, thick] (A) -- (B) -- (C) -- (D);
  \draw[black, thick] (C) -- (E);
  \draw[black, thick] (B) -- (F);
  \path [draw=blue, thick] (A) to [bend left] (A1); 
  \path [draw=blue, thick] (A) to [bend right] (A1); 
  %\draw[red, thick, bend left] (A) -- (A1);
  \draw[blue, thick] (B) -- (B1);
  \draw[blue, thick] (C) -- (C1);
  \draw[blue, thick] (D) -- (D1);
  \draw[blue, thick] (E) -- (E1);
  \draw[purple, thick] (A) -- (F);

\end{tikzpicture}
    \hspace{10pt}
    \begin{tikzpicture}[scale=2, every node/.style={circle, fill=black, inner sep=1.5pt}]
  % Nodes
  \node (A) at (0,0) {};       % root
\node[anchor=south,draw=none,fill=none] at (A) {$\al_1$};
  \node (B) at (0,-0.5) {};  
\node[anchor=north,draw=none,fill=none] at (B) {$\al_2$};
  \node (C) at (-0.7,-0.9) {};     
  \node (D) at (-1,-1.5) {};  
  \node (E) at (-0.5,-1.5) {};    
  \node (F) at (0.7,-0.9) {};
  \node[anchor=north,draw=none,fill=none] at (F) {$\al_3$};

\node (A1) [left of = A, distance=4pt] {};
\node[anchor=south,draw=none,fill=none] at (A1) {$a_1$};
%\node (B1) [right of = B, distance=4pt] {};  
\node (C1) [left of = C, distance=4pt] {};
\node (D1) [left of = D, distance=4pt] {};
\node (E1) [right of = E, distance=4pt] {};
\node (F1) [right of = F, distance=4pt] {};



  % Black edges (tree structure)
  \draw[purple, thick] (A) -- (B);
  \draw[black, thick] (A) -- (F);
  \draw[black, thick] (C) -- (D);
  \draw[black, thick] (C) -- (B);
  \draw[black, thick] (C) -- (E);
  \draw[black, thick] (B) -- (F);
  \path [draw=blue, thick] (A) to [bend left] (A1); 
  \path [draw=blue, thick] (A) to [bend right] (A1); 
  %\draw[red, thick, bend left] (A) -- (A1);
  %\draw[blue, thick] (B) -- (B1);
  \draw[blue, thick] (C) -- (C1);
  \draw[blue, thick] (D) -- (D1);
  \draw[blue, thick] (E) -- (E1);
  \draw[blue, thick] (F) -- (F1);

\end{tikzpicture}
\caption{Two examples of a connected component $\cal C$ in Case (IV), after the splitting \eqref{kwuyayw_ng_tree} with $(\al,\beta)=(\al_1,\al_2)$, in the term containing $\xi_{\al_1 a_1}^2$. Blue edges represent the solid edges to the set $\mathscr A$, here to external vertices (blue is used solely for illustration and does not indicate the $+$ charge of $G$-edges as in \Cref{def_graph1}); the purple edge is the edge between $\al_1$ and $\al_2$ on a cycle of $\cal C$; and black edges represent the other solid edges inside $\cal C$. The vertex $\al_1$ is added to $\mathscr A$ first; every other vertex except $\al_2$ is added as soon as it is adjacent to $\mathscr A$, and $\al_2$ is added once two of its neighbors belong to $\mathscr A$.}\label{Fig:construct_trees}
  \end{figure}

 \end{proof}

% Now, due to the previous two cases, we know that if a $z$ vertex connects with three or more ending edges, then it either be two short ending edges case, or two long ending edges case. Therefore, any $z$ vertex only connects at most two ending edges.  There are at least $p$ paths passing through $z$ vertices, and there two ending edges connect with these $z_1\cdots z_p$ vertices. Therefore, the only possible case (which we have not discussed) is the one with the following properties

% \begin{enumerate}
% \item There are exactly $p$ paths passing through $z_1\cdots z_p$.
%     \item each $z_i$ connects with two ending edges
%     \item for the ending edges connecting with $z_i$, one is $A1$ or $B1$ type, another one is $B2$.
% \end{enumerate}
% We notice that there are $p$ $B2$ ending edges in the $p$ paths. Therefore, all of these $p$ paths have ghost edges, i.e., for all $1\le i\le p$, 
% $$
% P_{ng}^{(i)}=0
% $$

% It implies that there is \textbf{no $A1$ types. } More precisely,  each $z_i$ connects with two ending edges, one is $B1$ type, another one is $B2$.

% Now each $z_i$ should connect with exactly two other edges, therefore, these edges form loops (with $z$ as vertices). Wlog, we can assume that there is only 1 loop connecting $z_i$'s like 
% $$
% z_1-z_2-z_3\cdots z_p-z_1$$

% Wlog, we assume that the $B1$ type ending edge connecting with $z_i$ is $\overline {z_ia_i}$. Therefore, we only need to prove that 
% $$
% \sum_{z_1\cdots z_p}\left(\prod_{i=1}^pf(z_i,a_i)\right)
% \left(\prod_{i=1}^pf(z_i,z_{i+1})\right)\prec \eta^{-p}
% $$
% The trick is applying Schwarz inequality on $\prod_{i=1}^2f(z_i,a_i)$,
% $$
% \prod_{i=1}^2f(z_i,a_i)\le \sum_i f^2(z_i,a_i)
% $$
% By symmetric, we only need to prove that 
%  $$
% \sum_{z_1\cdots z_p}f^2(z_1,a_1)\cdot \left(\prod_{i=3}^pf(z_i,a_i)\right)
% \left(\prod_{i=1}^pf(z_i,z_{i+1})\right)\prec \eta^{-p}
% $$
% To see that, one only needs to do the sum with the order $z_2$ (first), then $z_3,z_4\cdots z_p$ and $z_1$ (last).  Done. 


%\subsubsection*{\bf The mixed case: general} 
% So far, we have reduced to the case that there is no $z$'s connecting with two or more short ending edges in this graph. But we know there are at least $p$ short ending edges and there are $p$ $z$ vertices. It implies the following key result:
%  \begin{enumerate}
%      \item  Each $z_i$ connects with exactly one short ending edge, more precisely $B1$ edge. 
%      \item There are exactly $p$ paths and all of them have ghost edge. 
%  \end{enumerate}
%  Then we only need to prove 
%  $$
%  \sum_{z_1\cdots z_p} Val (\Gamma)\prec \eta^{-p}
%  $$
% More precisely, we only need to create a net without using ghost edge.  Now we reset the graph as follows: 
%  \begin{enumerate}
%      \item For each $z_i$, it connects with $a_i$ and it is a $B1$ edge. 
%      \item The other non-ghost edges are $z-z$ edges.  
%  \end{enumerate}
% Now we look at one special case: 
% $$
% \deg_{non-ghost} (z_i)=2$$
% Note, $\deg_{non-ghost} (z_i)$ is impossible to be less than 2. If it is equal to $2$, with $\deg (z_i)\ge 4$, it must connect with two ghost edges, then it can be reduced to simple case as abvoe two long ending edge case. 
% Therefore we can assume that 
% $$
% \deg_{non-ghost} (z_i)\ge 3
% $$






\subsection{Proof of Lemma \ref{lem:LW_moment_exp}}\label{subsec_pf_LW_moment_exp}

Note that when $|a-b|\le (\log W)^{3/2}\ell_t$ or $\ell\le (\log W)^{3/2}\ell_t$, $\sT_{t}(|a-b|\wedge \ell)$ does not exhibit exponential decay. In this regime, \Cref{lem:LW_moment_exp} follows directly from \Cref{lem:LW_moment}. Therefore, for the remainder of the proof, it suffices to establish \eqref{eq:LW_moment_exp} under the conditions  
\be\label{eq:far_ab_K}
|a-b|> (\log W)^{3/2}\ell_t,\quad \text{and}\quad (\log W)^{3/2}\ell_t\le \ell \le (\log W)^{10}\ell_t.
\ee
For this purpose, we decompose $f_{xy}(G)=f_{xy}^{>\ell}(G)+f_{xy}^{\le \ell}(G)$ into two parts, where
\begin{align*}
f_{xy}^{>\ell}(G)&=W^{-d}\sum_{a_1:|a_1-a|\wedge |a_1-b|>\ell}\sum_{a_2}S^{(\sB)}_{a_1a_2}\sum_{\substack{\al\in[a_2],\beta \in[a_1],\\ \al\notin\{x,y\}}}
\Gc_{\beta\beta} G_{x\al }G_{\al y},\\
f_{xy}^{\le \ell}(G)&=f_{xy}^{(a)}(G)+f_{xy}^{(b)}(G),\quad f_{xy}^{(a)}(G):=W^{-d}\sum_{a_1:|a_1-a|\le \ell}\sum_{a_2}S^{(\sB)}_{a_1a_2}\sum_{\substack{\al\in[a_2],\beta \in[a_1],\\ \al\notin\{x,y\}}}
\Gc_{\beta\beta} G_{x\al }G_{\al y},
\end{align*}
and $f_{xy}^{(b)}(G)$ is defined in the same way as $f_{xy}^{(a)}(G)$, with the summation over $a_1$ restricted to $|a_1-a|> \ell\ge |a_1-b|$.
\Cref{lem:LW_moment_exp} follows immediately from the next two estimates on \(f_{xy}^{>\ell}(G)\) and \(f_{xy}^{\le \ell}(G)\) with $\Psi_t = (W^{-d}B_{t,0})^{1/2}$:

\begin{lemma}\label{lem:LW_moment_exp_far}
In the setting of \Cref{lem:LW_moment_exp}, for any fixed $p\in 2\N$, the following estimate holds for any large constant $D>0$:
\be\label{eq:LW_moment_exp_far}
\E\big|f_{xy}^{>\ell}(G)\big|^{p}\prec \frac{1}{\eta_t^p} \Psi_t^p \cdot \br{\sT_{t}(|a-b|\wedge \ell) }^{p}+W^{-D}.
\ee
\end{lemma}
\begin{lemma}\label{lem:LW_moment_exp_near}
In the setting of \Cref{lem:LW_moment_exp}, for any fixed $p\in 2\N$, the following estimate holds for any large constant $D>0$:
\be\label{eq:LW_moment_exp_near}
\E\big|f_{xy}^{\le \ell}(G)\big|^{p}\prec \frac{1}{\eta_t^p} \Psi_t^p \cdot \br{\sT_{t}(|a-b|\wedge \ell) }^{p}+W^{-D}. 
\ee
\end{lemma}



First, the proof of \Cref{lem:LW_moment_exp_far} is similar to that for \Cref{lem:LW_moment}.
\begin{proof}[\bf Proof of \Cref{lem:LW_moment_exp_far}]
First, similar to \Cref{lem:localregular}, we expand $\big|f_{xy}^{>\ell}(G)\big|^{p}$ into a sum of locally standard graphs satisfying properties (1)--(6) listed there. Next, we bound each locally standard graph using its associated auxiliary graph, constructed as in \Cref{def_auxgraph}.
We then estimate the auxiliary graphs using a result analogous to \Cref{lem:Anp_key}. More precisely, suppose \smash{$\cal G^{\aux}_{\ba\bfb}$} is a nested graph that satisfies the assumptions of \Cref{lem:Anp_key}. In addition, all internal vertices lie in the domain  $\bD_{>\ell}:=\{c \in \Zn:|a_i-c|\wedge |b_i-c|> \ell \ \text{ for all } i\in\qqq{p}\}$. Indeed, every vertex created by the local expansions is connected to an existing vertex by a path of dotted or waved edges (see the discussion below \Cref{T eq0}), so every internal molecule contains one of the vertices $\beta\in[a_1]$ of the $p$ factors $f_{xy}^{>\ell}(G)$ or $\overline{f_{xy}^{>\ell}(G)}$, and we choose such a vertex as the center of the molecule in \Cref{def_auxgraph}. Then, we aim to show that
\begin{equation}\label{adsuu33}
\cal G^{\aux}_{\ba\bfb} \prec (W^d\eta_t)^{-q}\cdot \Psi_t^{\ord(\cal G^{\aux}_{\ba\bfb})-p}\cdot \br{\sT_{t}(|a-b|\wedge \ell) }^{p}+W^{-D} .
\end{equation}
In the proof of \Cref{lem:Anp_key_gh}, each path contributed at least one factor of $\Psi_t(c|a_i-b_i|)$ because, roughly speaking, every path contains at least one long ending edge. In the current setting, where $a_i=[a]$ and $b_i=[b]$ for all $i\in\qqq{p}$, we use the fact that the first edge of each path $\fP_i$ is ``long'': it joins $[a]$ either to $[b]$, in which case its length is $|a-b|$, or to an internal vertex, in which case its length is larger than $\ell$, since all internal vertices lie in the domain $\bD_{>\ell}$.
Since each solid edge $\xi([\al],[\beta])$ is bounded by $\sT_t(|[\al]-[\beta]|\wedge\ell)+W^{-D}$ and $\sT_t$ is decreasing, each such long edge provides a factor of $\sT_t(|a-b|\wedge\ell)$ (up to a negligible error $W^{-D}$), and we can treat it as a type-A2 edge. Replacing these A2 edges with ghost edges in all paths $\fP_i$, $i\in \qqq{p}$, yields a total factor of $\br{\sT_{t}(|a-b|\wedge\ell)}^{p}$. For the remaining graph---denoted by \(\cal G^{\new}_{\ba\bfb}\)---we need to establish the following bound for any constant $D>0$:
\begin{equation}\label{adsuu44}
\cal G^{\new}_{\ba\bfb} \prec (W^d\eta_t)^{-q}\cdot \Psi_t^{\ord(\cal G^{\new}_{\ba\bfb})} +W^{-D} .
\end{equation}
The proof of this estimate follows exactly the same argument as in \Cref{lem:Anp_key_gh}, and we omit the details.
\end{proof}

For the proof of \Cref{lem:LW_moment_exp_near}, we need a new argument that makes use of the exponential decay in $\sT_t$. 
\begin{proof}[\bf Proof of \Cref{lem:LW_moment_exp_near}]
Since $|f_{xy}^{\le \ell}(G)|^p\le 2^{p-1}(|f_{xy}^{(a)}(G)|^p+|f_{xy}^{(b)}(G)|^p)$, it suffices to bound $\E|f_{xy}^{(a)}(G)|^{p}$; the term $f_{xy}^{(b)}(G)$ is handled in the same way, with $a$ replaced by $b$ in the definition of the domain $\bD_{\le\ell}$ below.
Similar to \Cref{lem:localregular}, we expand $\big|f_{xy}^{(a)}(G)\big|^{p}$ into a sum of locally standard graphs satisfying properties (1)--(6) listed there.
Next, we bound each locally standard graph using its associated auxiliary graph, constructed as in \Cref{def_auxgraph}. This yields a class of auxiliary graphs \smash{$\cal G^{\aux}_{[a][b]}$} satisfying the assumptions of \Cref{lem:Anp_key} with $[a_i]=[a]$ and $[b_i]=[b]$ for $i\in \qqq{p}$. Moreover, choosing the centers of the internal molecules as in the proof of \Cref{lem:LW_moment_exp_far}, all internal vertices of the graph lie in the domain $\bD_{\le \ell}:=\{c \in \Zn:|a-c|\le \ell\}$. It remains to bound such graphs as follows: % with $\Psi_t=(W^{-d}B_{t,0})^{1/2}$:  
\begin{equation}\label{adsuu_exp}
\cal G^{\aux}_{[a][b]} \prec (W^d\eta_t)^{-q}\cdot \Psi_t^{\ord(\cal G^{\aux}_{[a][b]})-p}\cdot \br{\sT_{t}(|[a]-[b]|\wedge \ell)}^{p}+W^{-D} .
\end{equation}

We first control \smash{$\cal G^{\aux}_{[a][b]}$} by bounding each solid edge, say $\xi([\al],[\beta])$, by its upper bound $\sT_t(|[\al]-[\beta]|\wedge \ell)+W^{-D}$. 
Discarding the negligible error term containing $W^{-D}$ factors, we obtain a new graph, denoted by \smash{${\mathscr G}_{[a][b]}$}, which has the same graphical structure as \smash{$\cal G^{\aux}_{[a][b]}$} but with each solid edge representing a $\sT_t$ factor instead. Hence, to prove \eqref{adsuu_exp}, it suffices to establish that
\begin{equation}\label{adsuu_exp2}
\mathscr G_{[a][b]} = \sum_{\bal=([\al_1],\ldots,[\al_q])\in (\bD_{\le \ell})^q}\mathscr G_{[a][b]}(\bal) \prec (W^d\eta_t)^{-q}\cdot \Psi_t^{\ord(\mathscr G_{[a][b]})-p}\cdot \br{\sT_{t}(|[a]-[b]|\wedge \ell)}^{p} ,
\end{equation}
where, recall, \smash{$\mathscr G_{[a][b]}(\bal)$} denotes the graph obtained by fixing the external vertices to $\bal$, and $\ord(\mathscr G_{[a][b]})$ is defined in \eqref{eq:ordGaux}.
%denotes the number of solid edges in \smash{$\mathscr G_{[a][b]}$}. 
To show \eqref{adsuu_exp2}, we sum over the internal vertices in $\bal$ one by one. Unlike in the proof of \Cref{lem:Anp_key}, the order of summation here is arbitrary; for definiteness, we follow the order $[\al_1],\ldots,[\al_q]$. At each step, we apply the following key estimate.
\begin{claim}\label{claim:TTk}
Assume $1-t\ge \ilambda^2/L^{2}$. Then, for any $k \ge 2$ and any $0\le \ell \le (\log W)^{10}\ell_t$, the following bound holds with $\Psi_t=(W^{-d}B_{t,0})^{1/2}$:\footnote{This estimate can be viewed as an extension of \Cref{lem:propT} in the regime $1-t \ge \ilambda^2/L^2$, except that here we obtain a ``$\prec$'' bound, rather than the ``$\lesssim$'' bound in \eqref{TTT2}, due to the presence of additional logarithmic factors.}
\begin{align}\label{eq:key_T_reudce}
 \sum_{[\al]\in \bD_{\le \ell}} \prod_{i=1}^k \br{\sT_{t}(|[x_i]-[\al]|\wedge \ell) \cdot \sT_{t}(|[y_i]-[\al]|\wedge \ell) } \prec (W^{d}\eta_t)^{-1} \cdot {\Psi_t^{k-2}} \cdot   \prod_{i=1}^k \sT_{t}(|[x_i]-[y_i]|\wedge \ell) .
\end{align}
\end{claim}

 
We defer the proof of \Cref{claim:TTk} to \Cref{pf:claim_TTk}, where it is derived from elementary calculus estimates. (That proof only uses that $\bD_{\le\ell}$ is contained in a ball of radius $\ell$ in $\Zn$.)
%cite \cite{DYYY25_d3} here
The estimate \eqref{eq:key_T_reudce} shows that when summing over an internal vertex $[\al]$, each path of the form $([x_i]\to [\al]\to [y_i])$ consisting of two solid edges through $[\al]$ is effectively replaced by a single solid edge $[x_i]\to [y_i]$ in the resulting graph. We refer to this as the ``path-preserving phenomenon''. In addition, the summation over $[\al]\in \bD_{\le \ell}$ contributes a factor $(W^d\eta_t)^{-1}$ from each such pair of solid edges, along with the factor $\Psi_t^{k-2}$ reflecting the reduction in scaling order. Note that the case $k=2$ is critical---the estimate \eqref{eq:key_T_reudce} fails when $k=1$. \Cref{Fig:construct_keepT} illustrates the path-preserving phenomenon in the critical case $k=2$.

  \begin{figure}[h]
    \centering
    \begin{tikzpicture}
    \coordinate (L) at (-2, 0);
    \coordinate (x1) at (-1.3, 0.3);
    \coordinate (x2) at (-1.3,- 0.3);
    \coordinate (a1) at (0, 0.5);
    \coordinate (a2) at (0, -0.5);
    \coordinate (al) at (1, 0);
    \coordinate (b1) at (2, 0.5);
    \coordinate (b2) at (2, -0.5);
    %\coordinate (c1) at (3, 1);
    %\coordinate (c2) at (3, -1);
    \coordinate (d1) at (3.5, 0.4);
    \coordinate (d2) at (3.5, -0.4);
    \coordinate (R) at (4.8, 0);
   %\coordinate (R) at (6, -0.5);

    \fill (L) circle (1.5pt) node[left=2pt]{$[a]$};
    \fill (x1) circle (1.5pt) node[above=2pt]{};
      \fill (x2) circle (1.5pt) node[below=2pt]{};
      \fill (a1) circle (1.5pt) node[above=2pt]{$[x_1]$};
      \fill (a2) circle (1.5pt) node[below=2pt]{$[x_2]$};
    \fill (b1) circle (1.5pt)  node[above=2pt]{$[y_1]$};
      \fill (b2) circle (1.5pt) node[below=2pt]{$[y_2]$};
       \fill (al) circle (1.5pt) node[above=2pt]{$[\al]$};
      %\fill (c1) circle (1.5pt) node[right=16pt]{$\fP_1$};
      %\fill (c2) circle (1.5pt) node[right=16pt]{$\fP_2$};
      \fill (d1) circle (1.5pt);
      \fill (d2) circle (1.5pt);
      \fill (R) circle (1.5pt) node[right=2pt]{$[b]$};
      %\fill (L2) circle (1.5pt) node[below=2pt]{$[b_2]$};

      \draw (L) -- (x1); 
      \draw (L) -- (x2); 
      \draw (a1) [dotted]-- (x1); 
      \draw (a2) [dotted]-- (x2); 
      \draw (a1) -- (al); 
      \draw (a2) -- (al); 
      \draw (al) -- (b1); 
      \draw (b1) [dotted]-- (d1); 
      \draw (d1) -- (R); 
      \draw (al) -- (b2);
      \draw (b2) [dotted]-- (d2); 
      \draw (d2) -- (R); 
    \end{tikzpicture}
    \hspace{10pt}
    \begin{tikzpicture}
    \coordinate (L) at (-2, 0);
    \coordinate (x1) at (-1.3, 0.3);
    \coordinate (x2) at (-1.3,- 0.3);
    \coordinate (a1) at (0, 0.5);
    \coordinate (a2) at (0, -0.5);
    %\coordinate (al) at (1, 0);
    \coordinate (b1) at (1.5, 0.5);
    \coordinate (b2) at (1.5, -0.5);
    %\coordinate (c1) at (3, 1);
    %\coordinate (c2) at (3, -1);
    \coordinate (d1) at (3, 0.4);
    \coordinate (d2) at (3, -0.4);
    \coordinate (R) at (4.3, 0);
   %\coordinate (R) at (6, -0.5);

    \fill (L) circle (1.5pt) node[left=2pt]{$[a]$};
    \fill (x1) circle (1.5pt) node[above=2pt]{};
    \fill (x2) circle (1.5pt) node[below=2pt]{};
    \fill (a1) circle (1.5pt) node[above=2pt]{$[x_1]$};
    \fill (a2) circle (1.5pt) node[below=2pt]{$[x_2]$};
    \fill (b1) circle (1.5pt)  node[above=2pt]{$[y_1]$};
    \fill (b2) circle (1.5pt) node[below=2pt]{$[y_2]$};
    %\fill (al) circle (1.5pt) node[above=2pt]{$[\al_q]$};
      %\fill (c1) circle (1.5pt) node[right=16pt]{$\fP_1$};
      %\fill (c2) circle (1.5pt) node[right=16pt]{$\fP_2$};
    \fill (d1) circle (1.5pt);
    \fill (d2) circle (1.5pt);
    \fill (R) circle (1.5pt) node[right=2pt]{$[b]$};
      %\fill (L2) circle (1.5pt) node[below=2pt]{$[b_2]$};

      \draw (L) -- (x1); 
      \draw (L) -- (x2); 
      \draw (a1) [dotted]-- (x1); 
      \draw (a2) [dotted]-- (x2); 
      \draw (a1) -- (b1); 
      \draw (a2) -- (b2); 
      %\draw (al) -- (b1); 
      \draw (b1) [dotted]-- (d1); 
      \draw (d1) -- (R); 
      %\draw (a2) [dashed]-- (al);
     % \draw (al) -- (b2);
      \draw (b2) [dotted]-- (d2); 
      \draw (d2) -- (R); 
    \end{tikzpicture}
    \caption{Illustration of the path-preserving phenomenon in applying \eqref{eq:key_T_reudce}.}\label{Fig:construct_keepT}
  \end{figure}

%Key: Add in version 2 {\cor Draw a graph to illustrate the above process.}


We perform the summations of $\mathscr G_{[a][b]}$ over the internal vertices by generating a sequence of new graphs that consistently satisfy conditions (1)--(3) in \Cref{lem:Anp_key} with $[a_i]\equiv [a]$ and $[b_i]\equiv [b]$. To illustrate this procedure, consider the summation over $[\al_1]$. Suppose there are $k_1$ two-edge paths passing through $[\al_1]$ in $\mathscr G_{[a][b]}$. More precisely, assume that each path $\fP_i$ can be decomposed as
\[([a_i]\too [x_{i,1}]\to [\al_1] \to [y_{i,1}] \too [x_{i,2}]\to[\al_1]\to [y_{i,2}]\too \cdots \too [x_{i,r_{i}}]\to[\al_1]\to [y_{i,r_{i}}]\too [b_i]),\]
where we omit intermediate vertices on the path and only list those vertices $[x_{i,j}]$ and $[y_{i,j}]$ that are directly connected to $[\al_1]$, and $r_i\ge 0$ is a non-negative integer. Then, we apply \eqref{eq:key_T_reudce} to the summation 
\begin{align}\nonumber
 \sum_{[\al_1]\in \bD_{\le \ell}} \prod_{i=1}^p\prod_{j=1}^{r_i}  \br{\sT_{t}(|[x_{i,j}]-[\al_1]|\wedge \ell) \cdot \sT_{t}(|[y_{i,j}]-[\al_1]|\wedge \ell)} ,
\end{align}
which yields the bound
\[\sum_{[\al_1]\in \bD_{\le \ell}}\mathscr G_{[a][b]}(\bal) \prec (W^d\eta_t)^{-1} \Psi_t^{k_1-2} \cdot \mathscr G^{(1)}_{[a][b]}([\al_2],\ldots,[\al_q]) ,\] 
where $k_1$ is defined as $k_1=\sum_{i=1}^p r_i$, %we have used $W^{-d/2}\le \Psi_t$, 
and \smash{$\mathscr G^{(1)}_{[a][b]}$} is a new graph obtained by removing the vertex $[\al_1]$ from $\mathscr G_{[a][b]}$ and replacing each pair of solid edges $([x_{i,j}]\to[\al_1]\to [y_{i,j}])$ by $([x_{i,j}], [y_{i,j}])$.  
It is easy to see that the new graph \smash{$\mathscr G^{(1)}_{[a][b]}$} also satisfies conditions (1)--(3) in \Cref{lem:Anp_key}. Next, summing over the vertex $[\al_2]$ in \smash{$\mathscr G^{(1)}_{[a][b]}$} gives another graph \smash{$\mathscr G^{(2)}_{[a][b]}([\al_3],\ldots,[\al_q])$} that again satisfies conditions (1)--(3) in \Cref{lem:Anp_key}, along with a factor $(W^d\eta_t)^{-1} \Psi_t^{k_2-2}$ for some $k_2\ge 2$. 
Continuing this procedure, after summing over all internal vertices, we can bound the LHS of \eqref{adsuu_exp2} as  
\[\sum_{\bal\in (\bD_{\le \ell})^q}\mathscr G_{[a][b]}(\bal)\prec (W^d\eta_t)^{-q}\cdot \Psi_t^{\ord(\mathscr G_{[a][b]})-p}\cdot \mathscr G^{(q)}_{[a][b]},\]
where \smash{$\mathscr G^{(q)}_{[a][b]}$} is a graph consisting solely of $p$ solid edges between $[a]$ and $[b]$, with no internal vertices. 
%Using $\Psi_t\ge W^{-d/2}$, and 
Since each solid edge is bounded by $\sT_{t}(|[a]-[b]|\wedge \ell)$, this yields \eqref{adsuu_exp2}. Combining this with \eqref{G_by_auxG} completes the proof of \Cref{lem:LW_moment_exp_near}. 
%Combining the above three cases, we conclude \eqref{eq:key_T_reudce_pf}, which further implies \eqref{eq:key_T_reudce}. 
\end{proof}



\section{Extension to the block Anderson model}\label{sec:ext-to-BA}

The proof of \Cref{MR:decol_BA} for the block Anderson model is based on the following flow framework.

\begin{lemma}[Lemma 3.3 of \cite{RBSO1D}]\label{zztE_BA}
For the block Anderson model, fix any $\ilambda>0$ and $z\in \mathbb C_+$ with $\im z\in (0, 1]$ and $\im m(z,\ilambda)\ge \kappa$. We choose %Choosing $t_0$, $E_0$, and $\ilambda_0$ as follows: 
\be\label{eq:t0E0_BA}t_0=\frac{\im m(z,\ilambda)}{\im m(z,\ilambda)+ \im z},\quad E= \frac{t_0 \re z -(1-t_0)\re m(z,\ilambda)}{\sqrt{t_0}},\quad \ilambda_0=\sqrt{t_0}\ilambda \, .\ee
Then, we have that 
 \begin{equation}\label{eq:zztE_BA}
\sqrt{t_0}m(E,\ilambda_0)=m(z,\ilambda), \ \  z_{t_0}(E,\ilambda_0) = \sqrt{t_0}z, \ \  \sqrt{t_0}M(E,\ilambda_0)=M(z,\ilambda), \ \ G(z,\ilambda) \stackrel{d}{=} \sqrt{t_0} G_{t_0;E,\ilambda_0},
 \end{equation} 
% and the following equality in distribution:
%    \begin{equation}\label{GtEGz_BA}
%     G(z,\ilambda) \stackrel{d}{=} \sqrt{t_0} G_{t_0;E,\ilambda_0},
%    \end{equation}
where recall that $G(z,\ilambda)=(H-z)^{-1}=(V+\ilambda\Psi-z)^{-1}$.
\end{lemma}
Note that $t_0\ge \kappa/(\kappa+1)$, since $\im m(z,\ilambda)\ge\kappa$ and $\im z\le 1$, and that $\im m(E,\ilambda_0)=\im m(z,\ilambda)/\sqrt{t_0}\ge\kappa$ by \eqref{eq:zztE_BA}. In particular, the flow parameters $(E,\ilambda_0)$ satisfy the bulk condition of \Cref{lem:propM} below.



%and $|E_t|\le (2-\kappa)\sqrt{t}+\oo(1)$ for $|1-t|=\oo(1)$. 
%In other words, during the flow from $t=0$ to $t_0$, the imaginary part $\eta_t$ decreases from $\eta_0\asymp 1$ to $\eta_{t_0}\asymp 1-t_0 \asymp \im z $, while the spectral parameter $z_t$ stays within the bulk of the limiting spectrum of $H_t$, provided $t\ge 1 - \e$ for a small enough constant $\e>0$ depending on $\kappa$. 


In the following proof, we fix a target spectral parameter $z=\hat{E}+\ii \eta\in \mathbf D_{\kappa,\e}$ for an arbitrarily small constant $\e>0$, where the spectral domain $\mathbf D_{\kappa,\e}$ is now defined as %(recall \eqref{eq:spectral_domain}).
\be\label{eq:spectral_domainBA}
\mathbf D_{\kappa,\e}:=\{z=\hat{E}+\ii\eta \in \C_+: \im m(z,\ilambda)\ge \kappa, N^{-1+\e}\le \eta\le 1\}.
\ee
This domain contains the spectral parameters considered in \Cref{MR:decol_BA}: if $\hat E$ satisfies the bulk condition \eqref{eq:bulkBA}, i.e., $\rho_N(\hat E)\ge\kappa$, then $\hat E+\ii\eta\in \mathbf D_{\kappa',\e}$ for all $N^{-1+\e}\le \eta\le 1$, where $\kappa'>0$ is a constant depending only on $d$, $\kappa$, and $\fd$. To see this, let $f(\eta):=\im m(\hat E+\ii\eta,\ilambda)$ for $\eta>0$, so that $f(0^+)=\pi\rho_N(\hat E)\ge\pi\kappa$. By property (4) of \Cref{lem:propM} below, $|f'(\eta)|\le (2f(\eta)^2)^{-1}$, hence $f(\eta)^3\ge (\pi\kappa)^3-3\eta/2\ge(\pi\kappa)^3/2$ for $0<\eta\le (\pi\kappa)^3/3$. For $(\pi\kappa)^3/3\le\eta\le 1$, we use that $\mu_N$, being the free convolution of the semicircle law with a probability measure supported in $[-2d\ilambda,2d\ilambda]$, is supported in $[-R,R]$ with $R:=2+2d\ilambda\le 2+2d\fd^{-1}$, and that $|\hat E|\le R$ since $\rho_N(\hat E)>0$; thus, $f(\eta)=\int \eta\,(( x-\hat E)^2+\eta^2)^{-1}\mu_N(\dd x)\ge \eta/(4R^2+1)$. Consequently, it suffices to prove the estimates in \Cref{MR:decol_BA} for $z\in\mathbf D_{\kappa,\e}$, for any fixed constant $\kappa>0$.
Accordingly, we choose the parameters $t_0$, $E$, and $\ilambda_0$ as specified in \eqref{eq:t0E0_BA}. %The second identity in \eqref{eq:E0bulk} imply that $1-t\asymp \eta_t$ uniformly in $t\in [0,t_0]$. 
Again, for simplicity of presentation, unless we want to emphasize the dependence on the flow parameters $E$ and $g_0$, we will omit them from various notations, such as $z_t(E,\ilambda_0)$, $E_t(E,\ilambda_0)$, $\eta_t(E,\ilambda_0)$, $m(E,\ilambda_0)$, $M(E,\ilambda_0)$, and $G_{t;E,\ilambda_0}\equiv G_t$. 

Since $M$ is not a multiple of the identity matrix, the block Anderson model also involves the following \emph{twisted} quantities. For $c_1,c_2\in\Zn$, let $P_{c_1c_2}$ be the $N\times N$ matrix that maps the block $[c_2]$ to the block $[c_1]$ while preserving the position inside the block, normalized in the same way as $E_a$:
\[ (P_{c_1c_2})_{xy}:=W^{-d}\,\mathbf 1\big(x\in[c_1],\ y\in[c_2],\ x-Wc_1=y-Wc_2\big),\quad\text{i.e.,}\quad P_{c_1c_2}=W^{-d}\big(\mathbf e_{c_1}\mathbf e_{c_2}^{\top}\big)\otimes I_{W^d}, \]
where $\mathbf e_{c}$ denotes the standard basis vector of $\C^{L^d}$ associated with $c\in\Zn$. In particular, $P_{cc}=E_c$, $P_{c_1c_2}^*=P_{c_2c_1}$, and $\tr(M(\sig)P_{c_1c_2})=M^{(\sB)}_{c_2c_1}(\sig)$ by \eqref{def_G0}. We refer to $\tr\big((G_t(\sig)-M(\sig))P_{c_1c_2}\big)$ as a \emph{twisted trace}, and to
\[ \wt{\cal L}_{t,\bsig,(a;c_1,c_2)}:=\tr\big(G_t(\sig_1)E_aG_t(\sig_2)P_{c_1c_2}\big),\qquad \bsig=(\sig_1,\sig_2)\in\{+,-\}^2,\ \ a,c_1,c_2\in\Zn, \]
as a \emph{twisted $2$-loop}; when $c_1=c_2$, it is the $2$-$G$-loop ${\cal L}^{(2)}_{t,\bsig,(a,c_1)}$. The deterministic approximation of $\wt{\cal L}_{t,\bsig,(a;c_1,c_2)}$ is the \emph{twisted tree}
\[ \wt{\cal K}_{t,\bsig,(a;c_1,c_2)}:=W^{-d}\sum_{a'\in\Zn}\Theta^{(\sig_1,\sig_2)}_t(a,a')\,M^{(\sB)}_{c_2a'}(\sig_1)\,M^{(\sB)}_{a'c_1}(\sig_2), \]
which coincides with ${\cal K}^{(2)}_{t,\bsig,(a,c_1)}$ in \eqref{Kn2sol} when $c_1=c_2$, and with $\wt{\cal L}_{0,\bsig,(a;c_1,c_2)}$ at $t=0$, since $G_0(\sig)=M(\sig)$. Note also that $\tr(XE_aM(\sig)P_{c_1c_2})=W^{-d}M^{(\sB)}_{ac_1}(\sig)\tr(XP_{ac_2})$ for any $N\times N$ matrix $X$; this relates twisted $2$-loops containing $\Gc_t(\sig_2)=G_t(\sig_2)-M(\sig_2)$ in place of $G_t(\sig_2)$ to $\wt{\cal L}_{t,\bsig,(a;c_1,c_2)}$ and twisted traces.
In the flow framework of \Cref{zztE_BA}, we can establish an analogue of \Cref{lem:main_ind} for the block Anderson model.


\begin{theorem}\label{lem:main_ind_BA}
In the setting of \Cref{MR:decol_BA}, fix any $z=\hat{E}+\ii \eta\in \mathbf D_{\kappa,\e}$ and consider the flow framework in \Cref{zztE_BA}. Suppose the estimates \eqref{Eq:L-KGt+IND}--\eqref{Eq:Gtlp_exp+IND} hold at some fixed $s\in [0,t_0]$, together with the following estimate at $u=s$:
\begin{itemize}
\item[(e)] {\bf Twisted-loop estimate}: We have
\be\label{eq:TL2_BA}
\max_{\bsig,a,c_1,c_2}\left|\wt{\cal L}_{u,\bsig,(a;c_1,c_2)}-\wt{\cal K}_{u,\bsig,(a;c_1,c_2)}\right|\prec \p{W^{-d}B_{u,0}}^{3/2}.
\ee
\end{itemize}
Then, there exists a constant $0<\fc_d \le 10^{-2}$ such that for any $s< t < 1$ satisfying \eqref{con_st_ind}, the estimates \eqref{Eq:L-KGt}--\eqref{Gt_bound} hold, and \eqref{eq:TL2_BA} holds at $u=t$. In addition, if $1-t\ge \ilambda^2$, then the estimate \eqref{Eq:Gdecay+s<g} holds.
%if we do not assume \eqref{Eq:Gtlp_exp+IND}, then the estimates \eqref{Eq:L-KGt}, \eqref{Eq:L-KGt2}, \eqref{Eq:Gdecay}, and \eqref{Gt_bound} still hold at $t$.
\end{theorem}

%With \Cref{lem:main_ind}, we can prove \Cref{ML:GLoop,ML:GLoop_expec,ML:GtLocal} easily by induction in $t$. The proof of \Cref{lem:main_ind} will be divided into five steps, which constitute the remainder of this paper. Throughout these steps, we assume that the assumptions of \Cref{lem:main_ind} hold. In addition, each step will build on the conclusions established in the previous steps.

With \Cref{lem:main_ind_BA} in hand, we can establish \Cref{MR:decol_BA} by induction on $t$.  

\begin{proof}[\bf Proof of \Cref{MR:decol_BA}]
By the discussion below \eqref{eq:spectral_domainBA}, it suffices to consider $z\in\mathbf D_{\kappa,\e}$. Fix $z=\hat{E}+\ii \eta\in \mathbf D_{\kappa,\e}$, and choose the flow as in \Cref{zztE_BA}. At $t=0$, we have $G_0(\sig)=M(\sig)$, so ${\cal L}^{(n)}_{0,\bsig,\ba}={\cal K}^{(n)}_{0,\bsig,\ba}$ and $\wt{\cal L}_{0,\bsig,(a;c_1,c_2)}=\wt{\cal K}_{0,\bsig,(a;c_1,c_2)}$, and the estimates \eqref{Eq:L-KGt+IND}--\eqref{Eq:Gtlp_exp+IND} and \eqref{eq:TL2_BA} hold at $s=0$. By \Cref{lem:main_ind_BA}, applying induction on $t$ from $t=0$ to $t=t_0$ (simultaneously for all target spectral parameters, see the proof of Step 1 below) yields the estimates \eqref{Eq:L-KGt}--\eqref{Gt_bound} at $t=t_0$. Using \eqref{eq:zztE_BA}, \eqref{eq:BtBt}, and \eqref{Kn2sol}, we see that these estimates together imply the entrywise local law \eqref{G_bound}, the averaged local law \eqref{G_bound_ave}, and the quantum diffusion estimates \eqref{eq:diffu1}--\eqref{Meq:QdS2} for each fixed $z$. To extend these estimates uniformly to all $z\in \mathbf D_{\kappa,\e}$, we apply a standard $N^{-C}$-net and perturbation argument. Finally, the delocalization estimate \eqref{eq:psikLinfty}, the QUE estimates \eqref{Meq:QUE} and \eqref{Meq:QUE2}, and the bulk universality \eqref{eq:universality} then follow as consequences, as shown in \Cref{subsec:main}. For the bulk universality, two modifications are needed. First, the Green's function comparison is applied with an interpolation acting only on the random part $V=H-\ilambda\Psi$, i.e., centered at the deterministic matrix $\ilambda\Psi$, as in the flow \eqref{MBM}. Second, the argument of \cite[Theorem~2.6]{YY_25} that we follow uses the analogue of \cite[Lemma~5.15]{YY_25} (cf.~\Cref{lem:improve_exp_aver}) for the matrices obtained by adding a small GUE component. For the block Anderson model, the proof of this analogue requires, in addition, a twisted averaged law for the resolvents $G$ of these matrices, namely $\max_{c_1,c_2}|\tr((G-M)P_{c_1c_2})|\prec (N\eta)^{-1}$ at the corresponding spectral parameters $z=E+\ii\eta$, for the same reason as in the proof of \Cref{lem:improve_exp_aver} below; this bound is proved in the same way as \Cref{lem:TL1_BA}.
\end{proof}


%The proof of \Cref{lem:main_ind_BA} follows the same six-step strategy outlined below \Cref{lem:main_ind}. We now describe how the arguments in \Cref{Sec:Steps12}--\ref{Sec:graph} for the random band matrix model extend to the block Anderson model. Our proof relies on the properties of the $\Theta$-propagators stated in \Cref{lem_propTH} and properties of $m(z)$ defined in \eqref{self_m} and the matrix $\sM$ defined in \eqref{def_G0}.

The proof of \Cref{lem:main_ind_BA} follows the same six-step strategy outlined below \Cref{lem:main_ind}. We now explain how the arguments in \Cref{Sec:Steps12}--\ref{Sec:graph} for the random band matrix model extend to the block Anderson model. Our proof relies on the properties of the $\Theta$-propagators stated in \Cref{lem_propTH}, together with the following properties of $m(z)$ (defined in \eqref{self_m}) and the matrix $M^{(\sB)}$ (defined in \eqref{def_G0}).


%2-$M$-loop $M^{(\sigma_1,\sigma_2)}$ stated in \Cref{lem:propM}, which were established in \cite[Lemma~3.9]{RBSO1D}.

\begin{lemma}%[Lemma 3.9 of \cite{RBSO1D}]
\label{lem:propM}
For the block Anderson model, under the condition \eqref{eq:WO}, $m=m(E,\ilambda)$ and $M^{(\sB)}=M^{(\sB)}(E,\ilambda)$ satisfy the following properties for any $E\in\R$ with $\im m(E,\ilambda)\ge \kappa$ (recall that $E$ is the spectral parameter in the flow \eqref{eq:t0E0_BA}, and that this condition holds for the flow parameters $(E,\ilambda_0)$ of \Cref{zztE_BA} when $z\in\mathbf D_{\kappa,\e}$):

\begin{enumerate}
%\item {\bf Symmetry}: $\cal M$ is symmetric, which follows from the definition \eqref{def_G0} and the fact that $\Psi$ is real symmetric. As a result, $M^{(\sigma_1,\sigma_2)}$ is also symmetric. 

% \be\label{eq:Msym}
% (M^{(\sigma_1,\sigma_2)})^\top =M^{(\sigma_2,\sigma_1)}.
% \ee

\item[(1)] {\bf Translation invariance}: For any $a,b,r\in \Zn$, we have $M^{(\sB)}_{a+r,b+r}= M^{(\sB)}_{ab}$ and $M^{(\sB)}_{aa}\equiv m$. 

%$M^{(\sig_1,\sig_2)}(a+r,b+r)= M^{(\sig_1,\sig_2)}(a,b).$
% The following block translation invariance holds: 
% \be\label{symmetryM}
%   \qquad  M^{(\sigma_1,\sigma_2)}([x]+[a],[y]+[a])= M^{(\sigma_1,\sigma_2)}([x],[y]),\ \ \forall [x],[y],[a]\in \Zn.
% \ee
% In addition, in 2D, we have the parity symmetry: 
% \be\label{symmetryP} \qquad  M^{(\sigma_1,\sigma_2)}(0,[x])= M^{(\sigma_1,\sigma_2)}([x], 0)= M^{(\sigma_1,\sigma_2)}(0,-[x]),\ \ \forall [x] \in \Zn.\ee

\item[(2)] {\bf Ward's identity}: We have \(|m|\le 1\),
\(\im m\ge \kappa\), and %the following identity
\be\label{eq:WardM}
\sum_{b}|M^{(\sB)}_{ab}|^2=1,\quad \forall a\in \Zn \, .
\ee
% \be\label{eq:propm}
% |m|\le 1, \quad \im m\gtrsim 1.
% \ee

%\item[(3)] For all $a\in \ZL$, we have
%{\bf Cauchy-Schwarz inequality}: For any $\sig_1,\sig_2\in \{+,-\}$, we have
%\be\label{eq:WardM}
%|M^{(\sigma_1,\sigma_2)}_{ab}|\le M^{(+,-)}_{ab},\quad \text{and}\quad 
%\sum_{b}|\sM_{ab}|^2=1. %,\quad \forall a,b\in \Zn,
%\ee
%where the equalities holds if $\sigma_1\ne \sigma_2$.
% Moreover, if $\sigma_1\ne \sigma_2$, we have 
% \be\label{eq:WardM1}
% \sum_{b\in\Zn}M^{(\sigma_1,\sigma_2)}_{ab} = 1,\quad \forall a\in \Zn.
% \ee

\item[(3)] {\bf Combes--Thomas bound}: There exists a constant $C>0$ (depending only on $d$ and $\kappa$) such that the following estimate holds for $\ilambda<(2C)^{-1}$: 
\begin{align}
% \label{eq:M-msc} &\|M - m_{sc}(E)I_N\|\le C\ilambda ,\quad \|M^+ - [m_{sc}(E)]^2 I_N\|_{\infty\to\infty}\le C\ilambda ,\\ 
\label{Mbound_AO}
%\|M - m_{sc}(E)I_N\|\le C\ilambda , %\quad \|M^+ - [m_{sc}(E)]^2 I_N\|_{\infty\to\infty}\le C\ilambda ,\quad 
C^{-1}\ilambda \mathbf 1(a\sim b)\le |M^{(\sB)}_{ab}|\le (C\ilambda)^{|a-b|}. 
%\quad C^{-1}\ilambda^2 \mathbf 1(a\sim b)\le \big|M^{(\sB)}_{ab}\big|^2\le (C\ilambda)^{2|a-b|}, 
\end{align}
%where $M^{(\sig_1,\sig_2)}$ is defined in \eqref{eq:Msig}. 
%where we have used \smash{$M|_{[a][b]}$} to denote the $(a,b)$-th block of $M$ following the notation in  \eqref{eq:submatrixnotation}. Furthermore, for any $\sig_1,\sig_2\in \{+,-\}$ and $a\ne b\in \Zn$, we have 
% \be\label{Mbound_AO2}
% C^{-1}\ilambda^2 \mathbf 1(a\sim b)\le \big|\big(M^{(\sig_1,\sig_2)}S^{(\sB)}\big)_{ab}\big|\le (C\ilambda)^{2|a-b|}.   
% \ee
Furthermore, for $\ilambda\ge (2C)^{-1}$, there exists a constant $c>0$ such that the following  bound holds:
\begin{align}\label{Mbound_AO2}
|M^{(\sB)}_{ab}|\le c^{-1}\exp\p{-c|a-b|}.
\end{align}

\item[(4)] {\bf Stability}: For every $z\in\C_+$ and $\ilambda>0$, with $m=m(z,\ilambda)$ and $M^{(\sB)}=M^{(\sB)}(z)$, we have
\[ \re\Big(1-L^{-d}\tr\big(M^{(\sB)}\big)^2\Big)\ge 2(\im m)^2 . \]
Consequently, $m(z,\ilambda)$ is a smooth function of $(z,\ilambda)\in\C_+\times(0,\infty)$, and
\[ |\partial_z m(z,\ilambda)|\le \frac{1}{2(\im m(z,\ilambda))^2},\qquad |\partial_{\ilambda} m(z,\ilambda)|\le \frac{d}{(\im m(z,\ilambda))^2}. \]
\end{enumerate}
\end{lemma}
\begin{proof}
Property (1) follows directly from the translation invariance of the matrix $\Psi^{(\sB)}$ in \eqref{eq:Psi3D}. The bound $\im m \ge \kappa$ is our assumption, while Ward's identity \eqref{eq:WardM} follows by taking the imaginary part of the equation \smash{$m=M^{(\sB)}_{aa}=(\ilambda \Psi^{(\sB)} - E - m)^{-1}_{aa}$}. The identity \eqref{eq:WardM} gives directly $|m|\le 1$. 
For sufficiently small $\ilambda$, the estimates in \eqref{Mbound_AO} are obtained from the Taylor expansion 
\[ M^{(\sB)} = -\sum_{k=0}^\infty (E+m)^{-k-1}(\ilambda \Psi^{(\sB)})^k.\]
For $\ilambda$ of order 1, \eqref{Mbound_AO2} is given by the classical Combes--Thomas estimate (see, e.g., \cite[Theorem 10.5]{Aizenman_book}).
For property (4), let $\psi_1,\ldots,\psi_{L^d}$ be the eigenvalues of $\Psi^{(\sB)}$ (so that $|\psi_i|\le 2d$), and let $\mu_i:=(\ilambda\psi_i-z-m)^{-1}$ be the corresponding eigenvalues of $M^{(\sB)}$. Taking the imaginary part of \eqref{self_m}, i.e., of $m=L^{-d}\sum_i\mu_i$, gives $\im m=\im(z+m)\cdot L^{-d}\sum_i|\mu_i|^2$; hence, $L^{-d}\sum_i|\mu_i|^2\le 1$ and
\[ \re\Big(1-L^{-d}\sum_i\mu_i^2\Big)\ge L^{-d}\sum_i\big(|\mu_i|^2-\re \mu_i^2\big)=2L^{-d}\sum_i(\im\mu_i)^2\ge 2\Big(L^{-d}\sum_i\im\mu_i\Big)^2=2(\im m)^2 . \]
Differentiating \eqref{self_m} with respect to $z$ and $\ilambda$, we obtain
\[ \Big(1-L^{-d}\tr\big(M^{(\sB)}\big)^2\Big)\partial_z m=L^{-d}\tr\big(M^{(\sB)}\big)^2,\qquad \Big(1-L^{-d}\tr\big(M^{(\sB)}\big)^2\Big)\partial_{\ilambda} m=-L^{-d}\tr\big(\Psi^{(\sB)}(M^{(\sB)})^2\big), \]
where the right-hand sides are bounded in absolute value by $L^{-d}\sum_i|\mu_i|^2\le 1$ and $2d$, respectively. Together with the implicit function theorem (the derivative of \eqref{self_m} with respect to $m$ does not vanish by the above stability bound), this proves property (4).
\end{proof}

Another key ingredient in the proof of \Cref{lem:main_ind_BA} is the following analogue of \Cref{lem_GbEXP}. 

\begin{lemma}\label{lem_GbEXP_BA}
In the setting of \Cref{lem:main_ind_BA}, suppose the estimates in \eqref{initialGT2} hold for a deterministic control parameter $W^{-d/2}\le \Psi_t \le W^{-\e_0}$.
%the event in \eqref{def_asGMc}. 
% In the setting of \Cref{lem:main_ind}, fix any $t\in [0,t_0]$ and suppose we have the following estimates for any $\Psi$ satisfying the assumptions \eqref{eq:Psi1D}--\eqref{sym_cond}: there exists a small constant $\e_0>0$ such that 
% \begin{equation}\label{def_asGMc2}
% \|G_t - M\|_{\max} \prec W^{-\e_0},
% \end{equation}
Furthermore, suppose there exist deterministic control parameters $0<\Phi_t(a,b)\le W^{-\e_0}$ such that  
\be\label{eq:def_Psit}
%\max_{[a],[b]\in \Zn} \cL^{(2)}_{t,(-,+),(a,b)}+W^{-d}\prec \Psi_t^2,\quad 
{\cal L}^{(2)}_{t, (-,+), (a,b)}\prec [\Phi_t(a,b)]^2,\quad \forall a,b\in \Zn .
\ee
%Then, the following estimates hold: 
\begin{itemize}
\item[(a)] {\bf Local laws}: The following entrywise and averaged local laws hold: 
    \begin{align}\label{GiiGEX_BA} 
 \|G_t-M\|_{\max} \prec \Psi_t, \quad \max_a\left|\tr \p{ \left(G_t - M\right)E_{a}}\right| \prec \Psi_t^2  . 
\end{align}

% \item[(b)] {\bf Averaged local law}: For any $a,b\in \Zn$, we have that %(recall $I_{[a]}$ defined in \eqref{def:Ia})
% \begin{align}\label{GavLGEX_BA}
% \left\langle  \left(G_t - M\right)E_{a}\right\rangle \prec \Psi_t^2 . %, \quad  W^{-d} \left\langle  \left(G_t - M\right)I_{[a]}MI_{[b]}\right\rangle  \prec \Psi_t^2\|I_{[a]}MI_{[b]}\| \, . 
% \end{align}

\item[(b)] {\bf Entrywise decay estimate}:
There exists a constant $c_{\ilambda}>0$ such that, for any large constant $D>0$ and all \(a,b\in \Zn \), the following estimate holds:
\begin{align}\label{GijGEX_BA}
\max_{x\in [a], y \in [b]} |(G_t-M)_{xy}|  \prec & \sum_{a', b' \in \Zn}  \Phi_t(a',b') e^{-c_\ilambda(|a'-a|+|b'-b|)}  + \Psi_t e^{-c_{\ilambda}|a-b|}+W^{-D} .
\end{align}
\end{itemize}
\end{lemma}
\begin{proof}
This lemma was proved as Lemma 6.1 in \cite{RBSO1D} under the condition $\ilambda\le W^{-\e}$ for a small constant $\e>0$. For $W^{-\e}<\ilambda\le\fd^{-1}$, the argument has to be modified in two places.
First, in the proof of the entrywise bounds in (a) and (b), the expansion of $M$ in powers of $\ilambda$ is replaced by the bounds \eqref{Mbound_AO} and \eqref{Mbound_AO2}, and the self-consistent equation for the diagonal entries is inverted directly in the max norm, without any expansion in $\ilambda$. Indeed, the matrix $1-tM^{(+,+)}$, $t\in[0,1]$, is diagonally dominant: its diagonal entries are equal to $1-tm^2$, while, by \eqref{eq:WardM}, its off-diagonal row sums are equal to
\[ t\sum_{b\ne a}\big|M^{(\sB)}_{ab}\big|^2=t\big(1-|m|^2\big)=\re\big(1-tm^2\big)-(1-t)-2t(\im m)^2\le \big|1-tm^2\big|-\min\{1,2\kappa^2\}, \]
where we used $\kappa\le\im m\le 1$. Hence, $\|(1-tM^{(+,+)})^{-1}\|_{\infty\to\infty}\le (\min\{1,2\kappa^2\})^{-1}$ uniformly in $t\in[0,1]$ and $0<\ilambda\le\fd^{-1}$.
Second, the averaged local law in \eqref{GiiGEX_BA} does not follow from the argument of \cite{RBSO1D} when $\ilambda\gg W^{-\e}$. Instead, we derive it from the entrywise local law by a fluctuation averaging argument for the block averages $\tr((G_t-M)E_a)$, which uses conditional expectations with respect to the independent Gaussian blocks of $V$, the bounds $|m|\le1$ and $\im m\ge\kappa$, the above stability bound, and the row bound $\max_a\sum_b|M^{(\sB)}_{ab}|\le C$ implied by \eqref{Mbound_AO} or \eqref{Mbound_AO2}. This argument does not use the smallness of $\ilambda$ and holds uniformly in $0<\ilambda\le\fd^{-1}$.
\end{proof}



\noindent{\bf Proof of Step 1 for \Cref{lem:main_ind_BA}.}
The proof of Step 1 depends on the following continuity estimates in \Cref{lem_ConArg_BA}. To make the dependence on the spectral parameter $z$ and the coupling parameter $\ilambda$ explicit, we denote the resolvent by \(G_t(z_t, \ilambda)=(V_t+\ilambda\Psi-z_t)^{-1}\) and the corresponding $G$-loop by \smash{${\cal L}^{(n)}_{t, \bsig, \ba}(z_t, \ilambda)$}, where $z_t\equiv z_t(E,\ilambda)$ is defined in \eqref{eq:zt}. %Moreover, we define the $T$-variables $(T_t)_{x\bv}(z,\ilambda)$ and \smash{$(\wt T_t)_{x\bv}(z,\ilambda)$} in terms of the resolvent $G_t(z,\ilambda)$ as in \eqref{eq:gen_T}.

\begin{lemma}\label{lem_ConArg_BA}
Fix any \( \e\le s \leq t \leq 1\) for a constant \(\e > 0\). Given any $\ilambda$ satisfying the condition \eqref{eq:WO}, let $\ilambda_s:=\ilambda\cdot \sqrt{s/t}$, and assume that $\im m(E,\ilambda_s)\ge\kappa$. Then, in the flow setting given by \Cref{def_flow} and \Cref{zztE_BA}, we have the following continuity estimates for the $G$-loops and (generalized) resolvent entries. 
\begin{enumerate}

\item Assume that the bound \eqref{eq:loopbound_s} holds at time \(s\) for the loops ${\cal L}^{(n)}_{s, \bsig, \ba}\p{z_s,\ilambda_s}$ for each fixed \(n \in \mathbb{N}\).  
% \begin{equation}\label{55}
%  \max_{\bsig, \ba} \left|{\cal L}^{(n)}_{s, \bsig, \ba}\p{z_s,\ilambda_s}\right| \prec \left( W^d \ell_{s}^d \eta_{s} \right)^{-n+1}. 
% \end{equation}
Then, on the event \(\Omega_t:=\{\|G_t(z_t,\ilambda)\|_{\max} \leq C_0\}\) for a constant $C_0>0$, the following estimate holds for any $n\ge 2$:
\begin{align}\label{res_lo_bo_eta_BA}
{\bf 1}(\Omega_t)\cdot\max_{\bsig, \ba}
\left|{\cal L}_{t, \bsig, \ba}^{(n)}\p{z_t,\ilambda}\right| \prec %\frac{\max_{[a]}\langle \im G_t(z_t,\ilambda) E_{[a]}\rangle}{\left( W^d \ell_{s}^d \eta_{t} \right)^{n-1}} =
\left(\frac{\eta_{s}}{\eta_{t}}\cdot W^{-d}B_{s,0}\right)^{n-1} .
\end{align}
% Moreover, for an arbitrary deterministic unit vector $\bv$, we have 
% \begin{align}\label{eq:cont_T}
% (T_t)_{x\bv}(z_t,\ilambda)+(\wt T_t)_{x\bv}(z_t,\ilambda) \prec (T_s)_{x\bv}(z_t,\ilambda_s)+(\wt T_s)_{x\bv}(z_t,\ilambda_s)+\frac{\ell_{t}^d}{\ell_{s}^d}\frac{\im (G_t)_{\bv\bv}(z_t,\ilambda)}{W^d \ell_{t}^d \eta_{t} }.
% \end{align}

\item Given any deterministic unit vectors $\bv,\bw\in \C^N$, suppose 
\be\label{eq:ImGs}\im (G_s)_{\bv\bv}(z_s,\ilambda_s)\lesssim 1,\quad \im (G_s)_{\bw\bw}(z_s,\ilambda_s)\lesssim 1, \quad \left|(G_s)_{\bv\bw}(z_s,\ilambda_s)\right|\lesssim 1,
\ee
with high probability, where we adopt the simplified notation of generalized matrix entries: given a matrix $\cal A$ and any vectors $\bv,\bw$, we denote $ \cal A_{\bv\bw}:= \bv^* \cal A\bw$. Then, the following estimates hold with high probability at time $t$: 
\be\label{eq:ImGt}
\im (G_t)_{\bv\bv}(z_t,\ilambda)\lesssim {\eta_s}/{\eta_t},\quad (G_t)_{\bv\bw}(z_t,\ilambda)\lesssim {\eta_s}/{\eta_t}.
\ee
 
\end{enumerate}

\end{lemma}
\begin{proof}
The proof of this lemma follows exactly the same argument as that of Lemma 7.1 in \cite{RBSO1D}, except that the factor $W^{-d}B_{s,0}$ in \eqref{res_lo_bo_eta_BA} corresponds to $(W^d\ell_s^d\eta_s)^{-1}$ therein. The assumption $\im m(E,\ilambda_s)\ge\kappa$ ensures that \Cref{lem:propM} applies to $m(E,\ilambda_s)$ and $M(E,\ilambda_s)$, which are the deterministic limits of the resolvent at time $s$; in particular, $|m(E,\ilambda_s)|\le 1$ and $\eta_t\le \kappa^{-1}\eta_s$.
\end{proof}

With \Cref{lem_GbEXP_BA,lem_ConArg_BA}, Step 1 of the proof of \Cref{lem:main_ind_BA} for the block Anderson model (i.e., the proof of \eqref{lRB1} and \eqref{Gtmwc}) is the same as that in \cite[Section 7.1]{RBSO1D}, except that \Cref{lem_ConArg_BA} is applied only with source times close to $1$. This is needed because the deterministic part $\ilambda_0\Psi$ of $H_t$ does not scale with $t$: \Cref{lem_ConArg_BA} compares the flow with parameters $(E,\ilambda_0)$ at time $u$ with the flow with parameters $(E,g')$, $g':=\ilambda_0\sqrt{s'/u}$, at an earlier time $s'$, and the latter flow must lie in the bulk. Let $c_*:=\min\{1/2,\fd\kappa^3/(6d)\}$.
\begin{itemize}
\item For $u\le 1-c_*$, we have $\eta_u=(1-u)\im m(E,\ilambda_0)\ge c_*\kappa$, hence $\|G_u\|\le (c_*\kappa)^{-1}$ and, deterministically, $|{\cal L}^{(n)}_{u,\bsig,\ba}|\le \|G_u\|^nW^{-(n-1)d}\lesssim (W^{-d}B_{u,0})^{n-1}$, since $B_{u,0}\ge(\ilambda^2+1)^{-1}$. Thus, \eqref{lRB1} holds for $u\le 1-c_*$ without using \Cref{lem_ConArg_BA}, and \eqref{Gtmwc} then follows as in the random band matrix case.

\item For $u>1-c_*$, we apply \Cref{lem_ConArg_BA} only with the source time $s':=\max(s,1-c_*)$. Then, $0\le \ilambda_0-g'\le \ilambda_0(1-s')\le c_*/\fd$, so property (4) of \Cref{lem:propM} (applied at $E+\ii\eta$, followed by the limit $\eta\downarrow 0$) gives $\im m(E,g')^3\ge \im m(E,\ilambda_0)^3-3d(\ilambda_0-g')\ge \kappa^3/2$. Hence, the source flow satisfies the assumption of \Cref{lem_ConArg_BA} with $\kappa$ replaced by $\kappa/2$. If $s'=1-c_*$, the bound \eqref{eq:loopbound_s} for the source flow at time $s'$ holds deterministically, as in the previous item.
\end{itemize}
If $s'=s$, the bound \eqref{eq:loopbound_s} for the source flow at time $s$ is provided by the induction hypothesis, in the following sense. By the scaling relations \eqref{eq:zztE_BA}, the flow with parameters $(E,g')$ at time $s$ coincides with the flow of \Cref{zztE_BA} for the coupling $g'/\sqrt{s}$ and the target spectral parameter $z':=z_{s}(E,g')/\sqrt{s}$, stopped at its final time $t_0(z')=s$. Here, $\im m(z',g'/\sqrt{s})=\sqrt{s}\,\im m(E,g')\ge\kappa/4$, so $z'\in\mathbf D_{\kappa/4,\e/2}$ for large $N$, and the coupling $g'/\sqrt{s}$ satisfies \eqref{eq:WO} up to a constant factor. Therefore, as in \cite[Section 7.1]{RBSO1D}, the induction on $t$ in the proof of \Cref{MR:decol_BA} is carried out simultaneously for all such target spectral parameters and couplings. With these modifications, we omit the details.

\medskip
\noindent{\bf Proof of Step 2 for \Cref{lem:main_ind_BA}.} 
In Step 2, the technical lemmas---\Cref{lem:newKLK,lem:LWterm,lem: EWGn2_N,lem: EMn2_N}---remain valid in the setting of the block Anderson model, except that the estimate \eqref{eq:MG_conclusion3} is modified as follows. Let ${\cal J}^{\ell}_{t,D}\ge W^{-d}$ be a \emph{deterministic} control parameter such that \smash{\(\wh{\cal J}^{\ell}_{t,D} \prec {\cal J}^{\ell}_{t,D}.\)}
% \begin{align}\label{defCALJ_deter}
% \wh{\cal J}^{\ell}_{u,D} \prec {\cal J}^{\ell}_{u,D}.
% \end{align}
Then, \eqref{eq:MG_conclusion3} continues to hold with \smash{$\wh{\mathcal{J}}_{t,D}^{\ell}$} replaced by \smash{${\mathcal{J}}_{t,D}^{\ell}$}; that is, 
\be\label{eq:MG_conclusion3_BA}
(\mathcal{E} \otimes \mathcal{E})^{M}_{t,\bsig, \ba,\ba} \prec \frac{1}{\eta_t} \left[\left(W^{-d}B_{t,0}\right)^{1/2} +  \left({\mathcal{J}}_{t,D}^{\ell}\right)^3\right]\cdot \left(W^{-d} \wT^{\ell}_{t,D}(|a-b|)\right)^2.
\ee
Assume that \eqref{Eq:Gdecay_w} holds. We combine \eqref{Kn2sol}, \eqref{Mbound_AO} (or \eqref{Mbound_AO2}), and \eqref{prop:ThfadC} to obtain \eqref{eq:kn2sol_decay}. Together with \eqref{Eq:Gdecay_w}, this yields \eqref{eq:L2_decay}. Then, applying \Cref{lem_GbEXP_BA} gives the desired estimates \eqref{Gt_bound_flow} and \eqref{Gt_avgbound_flow}.
Finally, we prove \eqref{Eq:Gdecay_w}. Using the same argument as between \eqref{lokis2} and \eqref{eq:L-K2max}, we obtain that for all $u\in[s,t]$,
\be\label{eq:L-K2max_BA} \max_{\bsig,\ba}\left|\left({\cal L}-{\cal K}\right)^{(2)}_{u,\bsig,\ba}\right|\prec 
 \p{\frac{1-s}{1-u}}^{C_0+\frac54}(W^{-d}B_{u,0})^{\frac 5 4} \le (W^{-d}B_{u,0})^{\frac1 5}\cdot \p{W^{-d} \widetilde{\cal {T}}^{\ell^{(0)}}_{u,D}\p{|a-b|}},
\ee
where $\ell^{(0)}=0$ as chosen in \eqref{ksjjuw}. Next, we implement a similar inductive argument as that below \eqref{ksjjuw}: assume that \eqref{kwr3juw} holds for length scales $K_u$ satisfying \eqref{eq:monotone_Ku}.
Moreover, suppose we have the initial estimate
\be\label{eq:Gronwall_dervJuD_BA}
\wh{\mathcal{J}}_{u,D}^{K_u} \prec \mathcal{J}_{u,D}^{K_u}, \ \ \ \forall u\in [s,t], \quad \text{where}\quad \mathcal{J}_{u,D}^{K_u}\equiv (W^{-d}{B}_{t,0})^{\frac{1}{50}}.    
\ee
We then define the stopping time 
\be\label{eq:def2_stopping_BA} 
\tau :=t\wedge T,\quad \text{with}\quad T:=\inf\left\{u\ge s: \wh{\cal J}_{u,D}^{K_u} \ge (W^{-d}B_{u,0})^{\frac1{100}} \right\}.\ee
By \Cref{lem: EWGn2_N}, the estimate \eqref{eq:boundGcterm} remains valid.
On the other hand, using \eqref{eq:MG_conclusion3_BA} together with \Cref{lem:DIfREP}, the estimate \eqref{eq:boundmgterm} becomes
\be\label{eq:boundmgterm_BA}
\begin{split}
\int_s^\tau \dd \mathcal{E}^{M}_{u, \bsig, {\ba}} &\prec  \Big[\left(W^{-d}B_{\tau,0}\right)^{1/4}  + \sup_{u\in[s,\tau]}\p{{\cal J}_{u,D}^{K_u}}^{3/2}\Big]\cdot \left(W^{-d} \wT^{K_\tau}_{\tau,D}(|a-b|)\right) .\end{split}
\ee
Applying the same reasoning as below \eqref{eq:boundmgterm} and invoking Gr{\"o}nwall's inequality, we obtain
\be\label{eq:Gronwall_dervJuD_BA2}
\wh{\mathcal{J}}_{u,D}^{K_u} \prec\left(\frac{1-s}{1-u}\right)^{C_0}\br{(W^{-d} {B} _{t,0})^{1/5} + \p{{\cal J}_{t,D}^{K_t}}^{3/2} },\quad \forall u\in [s,\tau].  
\ee
This implies that $T\ge t$ with high probability, so \eqref{eq:Gronwall_dervJuD_BA2} holds for all $u\in [s,t]$. We then take the RHS of \eqref{eq:Gronwall_dervJuD_BA2} as the new control parameter $\mathcal{J}_{u,D}^{K_u}$ and repeat the above argument. Iterating this procedure $\OO(1)$ times yields the improved bound 
\be\label{eq:Gronwall_dervJuD_BA3}
\wh{\mathcal{J}}_{u,D}^{K_u} \prec\left(|1-s|/|1-u|\right)^{C_0}(W^{-d} {B} _{t,0})^{1/5},\quad \forall u\in [s,t].  
\ee
Next, we redefine the scale $K_u'$ as in \eqref{eq:def_ell1}. Under this choice, using \eqref{eq:Gronwall_dervJuD_BA3}, we can re-establish the initial estimate \eqref{eq:Gronwall_dervJuD_BA} with $K_u$ replaced by $K_u'$, provided the constant $\fc_d$ in \eqref{con_st_ind} is chosen sufficiently small.
Finally, repeating the above procedure $\OO(1)$ more times yields \eqref{Eq:Gdecay_w} at $u=t$. The same argument clearly applies to all $u\in [s,t]$. 

It remains to justify the validity of \Cref{lem:newKLK,lem:LWterm,lem: EWGn2_N,lem: EMn2_N} for the block Anderson model. First, the proof of \Cref{lem:newKLK} in \Cref{subsec:pf_lem:newKLK} carries over with only a minor modification: in the derivation of \eqref{uu2j2ois}, we additionally make use of \eqref{Mbound_AO} or \eqref{Mbound_AO2}. 
The proofs of \Cref{lem:LWterm,lem: EWGn2_N} require more substantial changes in the graphical tools, which will be detailed in \Cref{subsec:LWchange-to-BA} below. 
Finally, the proof of \Cref{lem: EMn2_N} (with \eqref{eq:MG_conclusion3} replaced by \eqref{eq:MG_conclusion3_BA}) follows the same argument as in \Cref{subsec:pf_lem: EMn2_N}, with the following adjustments:

\begin{proof}[\bf Proof of \Cref{lem: EMn2_N} for the block Anderson model]
Most of the arguments in \Cref{subsec:pf_lem: EMn2_N} carry over directly to the block Anderson model, provided we can establish the following resolvent estimate under the assumption \eqref{eq:LW_assm}: for any $a,b\in \Zn$ and $x\in[a]$, $y\in [b]$, 
\begin{equation}\label{eq_resolventunderpoly}
|(G_t-M)_{xy}| \prec \Psi_t(|a-b|).
\end{equation}
To prove \eqref{eq_resolventunderpoly}, we apply \eqref{GijGEX_BA} to obtain
\begin{align*}
|(G_t-M)_{xy}|   \prec \sum_{|a'-a|\vee |b'-b| \le (\log W)^{3/2}}\Psi_t(|a'-b'|)e^{-c_\ilambda(|a'-a|+|b'-b|)} + \Psi_t(0)e^{-c_{\ilambda}|a-b|} + W^{-D} \prec \Psi_t(|a-b|) ,  
\end{align*}
where $D>0$ is an arbitrarily large constant. Here, in the second step, we have used the conditions in \eqref{eq:Psi} to obtain that %for sufficiently large $D$,
\be\label{eq:Psirelations}\begin{aligned}
    &\Psi_t(|a'-b'|)\prec \Psi_t(|a-b|) ,\quad \Psi_t(|a-b|)\gtrsim C_1^{-1}\p{|a-b|+1}^{-C_2}\Psi_t(0)\ge W^{-D},\\  
    &\Psi_t(0)e^{-c_{\ilambda}|a-b|} \lesssim C_1^{-1}(|a-b|+1)^{-C_2}\Psi_t(0) \lesssim \Psi_t(|a-b|).
\end{aligned}
\ee
Under the stronger assumption \eqref{eq:LW_assm_exp}, a parallel argument yields the following analogue of \eqref{eq_resolventunderpoly}:
\begin{equation}\label{eq_resolventunderexp}
|(G_t-M)_{xy}| \prec \big(W^{-d} \wT_{t,D}^{\ell}(|a-b|) \big)^{1/2},\quad \forall x\in [a], y\in [b].
\end{equation}


We now prove \eqref{eq:MG_conclusion} using the estimate \eqref{eq_resolventunderpoly}. As in \eqref{eq:S1+S2}, we need to bound 
\[\cal S_1:= W^d \sum_{|c-b|> |c-a|}  \mathcal{L}^{(6)}_{t, (\bsig\otimes \bsig)^{(1)}, \ba(c,c)},\quad \cal S_2:= W^d \sum_{|c-b|\le |c-a|} \mathcal{L}^{(6)}_{t, (\bsig\otimes \bsig)^{(1)}, \ba(c,c)}. \]
By symmetry, it suffices to establish \eqref{eq:MG_conclusion} for $\mathcal{S}_1$. %Analogous to \eqref{eq_S1_bound}, 
In this case, we apply \eqref{eq_sym_loop_bound} with $c'=c$ to get that 
\begin{equation}\label{eq_S1_bound_BA}
    \mathcal{S}_1 \lesssim \frac{1}{\eta_t}\max_{|c-b| \ge |a-b|/2} \left(\mathcal{L}^{(4)}_{t,\bsig^{(\alt)}, (c,b,c,b)}\right)^{1/2} \cdot  \max_{\sig\in\{+,-\}} \big|\mathcal{L}^{(3)}_{t, (\sigma,\sig_2,-\sig_2), (a,b,a)} \big|. %\mathcal{L}^{(4)}_{t, \bsig', (a,b,a,c)}.
\end{equation}
We write the $(c,b,c,b)$-loop as 
\[\mathcal{L}^{(4)}_{t,\bsig^{(\alt)}, (c,b,c,b)}=W^{-4d} \sum_{z',z''\in[c]}\sum_{y_1,y_2\in [b]} G_{y_1 z'}(\sig_1) G_{z'y_2}(-\sig_1) G_{y_2 z''} (\sig_1)G_{z''y_1}(-\sig_1).\]
Expanding each resolvent entry as 
\be\label{eq:G+G-M}
G_{\al \beta}(\sig)=M_{\al \beta}(\sig)+(G-M)_{\al\beta}(\sig),\quad \forall \al,\beta\in \{y_1,y_2,z',z''\}, \ \sig\in\{+,-\},\ee 
we bound every $(G-M)_{\al\beta}$ using \eqref{eq_resolventunderpoly} and each $M_{\al \beta}$ using \eqref{Mbound_AO} or \eqref{Mbound_AO2}. This yields (with $\e>0$ a small constant): 
\begin{align*}
   \mathcal{L}^{(4)}_{t,\bsig^{(\alt)}, (c,b,c,b)} &\prec \Psi_t^4(|b-c|)+\br{W^{-d}\Psi_t^3(|b-c|)+ W^{-2d}\Psi_t^2(|b-c|) + W^{-3d}\Psi_t(|b-c|)+W^{-3d}} e^{-\e|b-c|} \\
  & \lesssim \Psi_t^4(|b-c|) \lesssim \Psi_t^4(|a-b|),
\end{align*}
where the second step uses a bound analogous to \eqref{eq:Psirelations}, and the last step follows from $|c-b|\ge |a-b|/2$ together with the condition \eqref{eq:Psi}. 
Inserting this bound into \eqref{eq_S1_bound_BA}, we obtain the estimate \eqref{eq:reduce4to3} again.
% \begin{equation}\label{eq:reduce4to3_BA}
%     \mathcal{S}_1 \prec W^d \Psi_t^2\left(|a-b|\right) \sum_{c} \mathcal{L}^{(4)}_{t, \bsig', (a,b,a,c)}\prec \frac{1}{\eta_t} \Psi_t^2\left(|a-b|\right) \max_{\sig\in\{+,-\}}\left|\mathcal{L}^{(3)}_{t, (\sig,\sig_2,-\sig_2), (a,b,a)}\right|,
% \end{equation} 
% where, in the second step, we applied Ward's identity \eqref{WI_calL}. 
To handle the 3-loop on the RHS of \eqref{eq:reduce4to3}, we again expand each resolvent entry as in \eqref{eq:G+G-M}, bound every $(G-M)_{\al\beta}$ with \eqref{eq_resolventunderpoly}, and each $M_{\al \beta}$ with \eqref{Mbound_AO} or \eqref{Mbound_AO2}. This gives the same bound as in \eqref{eq:reduce4_bdd3}, 
% \begin{align}
%  \mathcal{L}^{(3)}_{t, (\sig,\sig_2,-\sig_2), (a,b,a)}& \prec \Psi_t(0)\cdot \Psi_t^2(|a-b|)+ \Psi_t(0)\cdot W^{-d}W^{-\e|a-b|}\cdot \Psi_t(|a-b|)+ W^{-2d} W^{-\e|a-b|} \nonumber\\   
%  &\lesssim \Psi_t(0) \cdot \Psi_t^2(|a-b|),\label{eq:reduce4_bdd3_BA}
% \end{align} 
where the factor $ \Psi_t^2(|a-b|)$ comes from the entries $G_{x'y}$ and $G_{yx}$ between blocks $a$ and $b$, while the short leg $G_{xx'}$ contributes a factor $\Psi_t(0)+W^{-d}\lesssim\Psi_t(0)$ by \eqref{GiiGEX_BA}. (Specifically, $(G-M)_{xx'}$ contributes $\Psi_t(0)$, while the deterministic part $M_{xx'}$ yields a factor $W^{-d}$.) %via the $(\infty\to \infty)$-norm bound in \eqref{Mbound_AO}.} 
Substituting \eqref{eq:reduce4_bdd3} into \eqref{eq:reduce4to3} concludes  \eqref{eq:MG_conclusion}.

Next, the exponential decay bound \eqref{eq:MG_conclusion3_BA} can be established by an argument analogous to that used for \eqref{eq:MG_conclusion3} below \eqref{eq:cutoff_scales}. 
First, it follows directly from the polynomial decay bound \eqref{eq:MG_conclusion} whenever \smash{$|a-b| \le \ell_t^{\dag}$}. It therefore remains to consider the case \eqref{eq:largea-b}, where we must control the three terms \smash{\(\widetilde{\mathcal{S}}_i\)}, $i\in \{1,2,3\}$, appearing in \eqref{eq;S123} with $c'=c$. The terms \smash{$\widetilde{\mathcal{S}}_1$ and $\widetilde{\mathcal{S}}_2$} can be bounded by the same argument used between \eqref{eq:pointwise_loop2} and \eqref{eq:boundwtS_1}, except that here we apply \eqref{eq_resolventunderexp} in place of the estimates \eqref{GijGEX} and \eqref{eq:LW_assm_exp} employed there.
For the term \smash{$\wt {\cal S}_3$}, we bound all legs of the original $6$-$G$-loop directly using \eqref{GijGEX_BA} together with the control parameter \smash{$\cal J_{t,D}^{\ell}$}. In this setting, \eqref{eq_L2-J} remains valid with \smash{$\wh{\cal J}_{t,D}$} replaced by \smash{$\cal J_{t,D}^{\ell}$}, where the $2$-$\cK$-loop \smash{$\mathcal{K}^{(2)}_{t,(-,+),(a,b)}$} defined in \eqref{Kn2sol} is of order $\OO(W^{-D})$ thanks to the exponential decay of the $\Theta$-propagators in \eqref{prop:ThfadC} and the decay of the $2$-$M$-loop due to \eqref{Mbound_AO} or \eqref{Mbound_AO2}. Using \eqref{eq_L2-J} (with \smash{$\wh{\cal J}_{t,D}^{\ell}$} replaced by $\cal J_{t,D}^{\ell}$) and an argument analogous to that below \eqref{eq_resolventunderpoly}, we obtain that for any $x\in [a]$ and $y\in [b]$ with $|a-b|\gtrsim\ell_t^*$,
\begin{equation}\label{eq_resolventunderexp2}
|(G_t-M)_{xy}| \prec \big(W^{-d} \wT_{t,D}^{\ell}(|a-b|) \big)^{1/2} \implies |(G_t)_{xy}| \prec \big(W^{-d} \wT_{t,D}^{\ell}(|a-b|) \big)^{1/2},
\end{equation}
where the implication uses the bound \eqref{Mbound_AO} or \eqref{Mbound_AO2}.
Finally, combining \eqref{eq_resolventunderexp2} with the same reasoning as below \eqref{eq_L2-J}, we conclude the exponential decay bound \eqref{eq:MG_conclusion3_BA}.
\end{proof}



\noindent{\bf Proof of Steps 3--6 for \Cref{lem:main_ind_BA}.} 
The proofs in \Cref{Sec:Steps34,Sec:Step5,Sec:Step6} extend verbatim to the block Anderson model, except for the arguments in \Cref{sec:Step5_larget}, the proof of \Cref{lem:improve_exp_aver}, and the proof of \Cref{lem:LWterm_EXP}. In \Cref{sec:Step5_larget}, the random band matrix case relies on \Cref{lem_GbEXP} and follows the approach of \cite[Section 5.3]{YY_25}, whereas the corresponding arguments for the block Anderson model are based on \Cref{lem_GbEXP_BA} and parallel those in \cite[Section 7.3]{RBSO1D}. We therefore omit the details.
The proof of \Cref{lem:LWterm_EXP} is deferred to \Cref{subsec:pf-LWterm_EXP} below.
In Step 6, i.e., in the proofs of \Cref{lem:improve_exp_aver,lem:LWterm_EXP}, we use in addition the following twisted averaged law, which is a consequence of Steps 2 and 3, and the twisted-loop estimate \eqref{eq:TL2_BA} for $u\in[s,t]$, which we derive after Step 4 below.

\begin{lemma}[Twisted averaged law]\label{lem:TL1_BA}
In the setting of \Cref{lem:main_ind_BA}, suppose that the estimates \eqref{Gt_bound_flow}, \eqref{Gt_avgbound_flow}, and \eqref{Eq:LGxb} with $n\le 4$ hold at some $u\in[s,t]$. Then, we have
\[ \max_{\sig\in\{+,-\}}\max_{c_1,c_2\in\Zn}\Big|\tr\Big(\big(G_u(\sig)-M(\sig)\big)P_{c_1c_2}\Big)\Big|\prec W^{-d}B_{u,0} . \]
In particular, this bound holds uniformly in $u\in[s,t]$ after Steps 2 and 3.
\end{lemma}
\begin{proof}
Since $\tr((G_u^*-M^*)P_{c_1c_2})=\overline{\tr((G_u-M)P_{c_2c_1})}$, it suffices to consider $\sig=+$. Denote $\delta_b:=\tr((G_u-M)E_b)$ and $\Lambda:=\max_{c_1,c_2}|\tr((G_u-M)P_{c_1c_2})|$. By \eqref{def_G0} and \eqref{eq:zt}, we have $H_u-z_u=V_u+M^{-1}+u\,m$, where $V_u:=H_u-H_0$. Together with $\cal S[M]=mI_N$, this gives the identity
\be\label{eq:GminusM_BA}
G_u-M=-M\,\underline{V_uG_u}+u\,M\,\cal S[G_u-M]\,G_u,\qquad \underline{V_uG_u}:=V_uG_u+u\,\cal S[G_u]\,G_u ,
\ee
where $\cal S[G_u-M]=W^d\sum_b\delta_bE_b$ by \eqref{eq:opS} and \eqref{bandcwV}. Multiplying \eqref{eq:GminusM_BA} by $P_{c_1c_2}$ and taking the trace, we obtain from $M=M^{(\sB)}\otimes I_{W^d}$ that
\[ \tr\big((G_u-M)P_{c_1c_2}\big)=-\tr\big(M\underline{V_uG_u}P_{c_1c_2}\big)+u\sum_b\delta_bM^{(\sB)}_{c_2b}M^{(\sB)}_{bc_1}+u\sum_b\delta_bM^{(\sB)}_{c_2b}\tr\big((G_u-M)P_{c_1b}\big). \]
By \eqref{Gt_avgbound_flow}, the Cauchy--Schwarz inequality, and \eqref{eq:WardM}, the second term on the right-hand side is $\prec W^{-d}B_{u,0}$. By \eqref{Gt_avgbound_flow} and the row bound $\max_c\sum_b|M^{(\sB)}_{cb}|\le C$, which follows from \eqref{Mbound_AO} or \eqref{Mbound_AO2}, the third term is $\prec W^{-d}B_{u,0}\cdot\Lambda$. Since $W^{-d}B_{u,0}\le W^{-c}$ for a constant $c>0$ by \eqref{eq:WO}, this term can be absorbed into the left-hand side after taking the maximum over $c_1,c_2$.

It remains to show that $X:=\tr(M\underline{V_uG_u}P_{c_1c_2})\prec W^{-d}B_{u,0}$, which follows from a high-moment estimate. Applying Gaussian integration by parts to the entries of $V_u$ in one factor $X$ of $\E|X|^{2p}$, $p\in\N$, the term in which the derivative acts on the resolvent in the same factor is cancelled by the renormalization in $\underline{V_uG_u}$. In the remaining terms, the derivative acts on one of the other $2p-1$ factors, and, using \eqref{eq:GminusM_BA} again (i.e., $\underline{V_uG_u}=-M^{-1}(G_u-M)+u\,\cal S[G_u-M]G_u$), it produces sums of products of resolvent entries with deterministic weights. The leading one is $u\sum_{x,w}S_{xw}(G_uP_{c_1c_2}M)_{wx}(G_uP_{c_1c_2}M)_{xw}$ (or its analogue with one factor replaced by its complex conjugate). Since $(G_uP_{c_1c_2}M)_{wx}=W^{-d}M^{(\sB)}_{c_2[x]}(G_u)_{w\tilde x}$, where $\tilde x\in[c_1]$ has the same position inside its block as $x$, the Cauchy--Schwarz inequality and \eqref{eq:WardM} bound it by
\[ u\,W^{-d}\sum_b\big|M^{(\sB)}_{c_2b}\big|^2\,{\cal L}^{(2)}_{u,(+,-),(b,c_1)}\prec (W^{-d}B_{u,0})^2 , \]
where we used \eqref{Eq:LGxb} with $n=2$. The other terms are bounded in the same way, using \eqref{Gt_bound_flow} and \eqref{Eq:LGxb} with $n\le 4$, by $(W^{-d}B_{u,0})^{2}$ times a product of $2p-2$ factors $X$ or $\overline X$. This yields $\E|X|^{2p}\prec (W^{-d}B_{u,0})^{2}\,\E|X|^{2p-2}$, and hence $X\prec W^{-d}B_{u,0}$ by H{\"o}lder's and Markov's inequalities, since $p$ is arbitrary.
\end{proof}

\medskip
\noindent{\bf Proof of the twisted-loop estimate \eqref{eq:TL2_BA} for $u\in[s,t]$.}
Fix $\bsig=(\sig_1,\sig_2)$ and $c_1,c_2\in\Zn$, and regard $Z_u(a):=\wt{\cal L}_{u,\bsig,(a;c_1,c_2)}-\wt{\cal K}_{u,\bsig,(a;c_1,c_2)}$ as a vector indexed by $a\in\Zn$. The derivation of the loop hierarchy \eqref{eq:mainStoflow} by It{\^o}'s formula does not use the special form of the matrices $E_{a_i}$, so it also applies to $\wt{\cal L}_{u,\bsig,(a;c_1,c_2)}$, viewed as a $2$-loop in which $E_{a_2}$ is replaced by $P_{c_1c_2}$. Since $S^{(\sB)}=I_{L^d}$, the cut-and-glue term in \eqref{eq:mainStoflow} then becomes $W^d\sum_{b}{\cal L}^{(2)}_{u,\bsig,(a,b)}\wt{\cal L}_{u,\bsig,(b;c_1,c_2)}$. On the other hand, since $\partial_u\Theta^{(\sig_1,\sig_2)}_u=\Theta^{(\sig_1,\sig_2)}_uM^{(\sig_1,\sig_2)}\Theta^{(\sig_1,\sig_2)}_u$, we have $\partial_u\wt{\cal K}_{u,\bsig,(a;c_1,c_2)}=W^d\sum_{b}{\cal K}^{(2)}_{u,\bsig,(a,b)}\wt{\cal K}_{u,\bsig,(b;c_1,c_2)}$ by \eqref{Kn2sol}. Hence, with the evolution kernel $\wt{\cal U}_{v,u}:=\Theta^{(\sig_1,\sig_2)}_u(\Theta^{(\sig_1,\sig_2)}_v)^{-1}$ acting on the index $a$, we have
\be\label{eq:Duhamel_twist}
Z_u=\wt{\cal U}_{s,u}Z_s+\int_s^u\wt{\cal U}_{v,u}\Big(\wt{\cal B}_{1,v}+\wt{\cal B}_{2,v}+\wt{\cal E}^{\Gc}_v\Big)\dd v+\int_s^u\wt{\cal U}_{v,u}\,\dd\wt{\cal E}^{M}_v ,
\ee
where $\wt{\cal B}_{1,v}(a):=W^d\sum_b(\cL-\cK)^{(2)}_{v,\bsig,(a,b)}\wt{\cal K}_{v,\bsig,(b;c_1,c_2)}$ and $\wt{\cal B}_{2,v}(a):=W^d\sum_b(\cL-\cK)^{(2)}_{v,\bsig,(a,b)}Z_v(b)$, while $\wt{\cal E}^{\Gc}_v$ and $\dd\wt{\cal E}^{M}_v$ are the light-weight and martingale terms defined as in \eqref{def_EwtG} and \eqref{def_Edif}, with $E_{a_2}$ replaced by $P_{c_1c_2}$. In particular, $\wt{\cal E}^{\Gc}_v(a)$ is a sum of twisted $3$-loops $\tr(G_v(\sig_1)E_{b}G_v(\sig_1)E_aG_v(\sig_2)P_{c_1c_2})$ and $\tr(G_v(\sig_1)E_{a}G_v(\sig_2)E_bG_v(\sig_2)P_{c_1c_2})$, multiplied by $W^d\tr(\Gc_v(\sig_i)E_{b})$ and summed over $b$.

We reduce twisted loops to $G$-loops using the following two facts. First, by the Cauchy--Schwarz inequality, for any $N\times N$ matrices $X_1,X_2$,
\[ \big|\tr\big(X_1E_aX_2P_{c_1c_2}\big)\big|^2\le\tr\big(E_{c_2}X_1E_aX_1^*\big)\,\tr\big(E_aX_2E_{c_1}X_2^*\big); \]
for example, $|\wt{\cal L}_{v,\bsig,(a;c_1,c_2)}|^2\le{\cal L}^{(2)}_{v,(\sig_1,-\sig_1),(a,c_2)}\,{\cal L}^{(2)}_{v,(\sig_2,-\sig_2),(a,c_1)}$. Second, since $G-G^*=2\ii\eta GG^*=2\ii\eta G^*G$ and $\sum_aE_a=W^{-d}I_N$, we have the Ward identity
\[ \sum_a\wt{\cal L}_{v,(\sig,-\sig),(a;c_1,c_2)}=\frac{1}{2\ii W^d\eta_v}\Big(\tr\big(G_vP_{c_1c_2}\big)-\tr\big(G_v^*P_{c_1c_2}\big)\Big),\quad \sig\in\{+,-\}, \]
and the same identity holds for $\wt{\cal K}$, with $\tr(G_vP_{c_1c_2})$ and $\tr(G_v^*P_{c_1c_2})$ replaced by $M^{(\sB)}_{c_2c_1}$ and $(M^{(\sB)})^*_{c_2c_1}$, respectively. Using these two facts together with the estimates of Steps 1--4, we obtain the following bounds, uniformly in $v\in[s,u]$, $\bsig$, $a$, $c_1$, and $c_2$:
\begin{itemize}
\item $|\wt{\cal B}_{1,v}(a)|+|\wt{\cal E}^{\Gc}_v(a)|\prec \eta_v^{-1}(W^{-d}B_{v,0})^2$. For $\wt{\cal B}_{1,v}$, this follows from \eqref{Eq:L-KGt-flow} with $n=2$ and the bound $\sum_b|\wt{\cal K}_{v,\bsig,(b;c_1,c_2)}|\le (W^d(1-v))^{-1}$, which is a consequence of \eqref{eq:THETAinftinf}, the symmetry of $\Theta^{(\sig_1,\sig_2)}_v$, and \eqref{eq:WardM}. For $\wt{\cal E}^{\Gc}_v$, we use \eqref{Gt_avgbound_flow}, and we bound the sum over $b$ of the twisted $3$-loops by $(W^d\eta_v)^{-1}W^{-d}B_{v,0}$, using the above Cauchy--Schwarz inequality with the vertex $E_b$ in place of $E_a$, the Cauchy--Schwarz inequality in $b$, Ward's identity \eqref{WI_calL} for the resulting $2$- and $4$-loops at the vertex $b$, and \eqref{Eq:LGxb}.
\item The quadratic variation of the martingale term is bounded in the same way, by $\eta_v^{-1}(W^{-d}B_{v,0})^{4-\delta}$ for any constant $\delta>0$ (in the sense of $\prec$).
\item If $|Z_{v}(b)|\le \Xi\,(W^{-d}B_{v,0})^{\gamma}$ for all $v\in[s,u]$ and $b\in\Zn$, for some $\gamma\in[1,2]$ and $\Xi\ge 1$, then $|\wt{\cal B}_{2,v}(a)|\prec \Xi\,\eta_v^{-1}(W^{-d}B_{v,0})^{\gamma+1/6}$, by \eqref{Eq:Gdecay_w} and \eqref{con_st_ind} (provided $\fc_d$ is small enough).
\item All these sources, as well as $Z_v(a)$, are bounded by $W^{-D}$ for any constant $D>0$ when $|a-c_1|\ge W^{\e}\ell_v$, by the decay of the $\Theta$-propagators in \eqref{prop:ThfadC}, of $M^{(\sB)}$ in \eqref{Mbound_AO} or \eqref{Mbound_AO2}, and of the $2$-$G$-loops in \eqref{Eq:Gdecay_w}.
\end{itemize}
With these bounds, \eqref{eq:TL2_BA} at time $u$ follows from \eqref{eq:Duhamel_twist} in the same way as \eqref{Eq:L-KGt-flow} for $n=2$ in Step 4: we use the one-index analogues of the evolution kernel estimates in \Cref{lem:sum_Ndecay,lem:sum_decay,lem:sum_decay_nonzero} for $\wt{\cal U}_{v,u}$ (which are proved in the same way, with $c_1$ playing the role of the second index), the Burkholder--Davis--Gundy inequality for the martingale term, and a bootstrap in $\gamma$ starting from $\gamma=1$; the initial bound $|Z_v(a)|\le|\wt{\cal L}_{v,\bsig,(a;c_1,c_2)}|+|\wt{\cal K}_{v,\bsig,(a;c_1,c_2)}|\prec W^{-d}B_{v,0}$ follows from the above Cauchy--Schwarz inequality, \eqref{Eq:LGxb}, \eqref{prop:ThfadC}, and \eqref{eq:WardM}. Finally, in the regime $1-u\le\ilambda^2/L^2$ and for $\sig_1\ne\sig_2$, the zero mode of $Z_u$ is controlled directly by the two Ward identities above and \Cref{lem:TL1_BA}, which give $|\sum_aZ_u(a)|\prec (W^d\eta_u)^{-1}W^{-d}B_{u,0}$.

\medskip
\noindent{\bf Modification in the proof of \Cref{lem:improve_exp_aver}.}
Since $M=M^{(\sB)}\otimes I_{W^d}$ is not a multiple of the identity matrix, the argument of \cite[Lemma 5.15]{YY_25} has to be modified as follows. Taking $P_{c_1c_2}=E_a$ in the identity derived in the proof of \Cref{lem:TL1_BA} and using $\E\,\underline{V_uG_u}=0$, which holds by Gaussian integration by parts, we obtain that $\delta_b=\tr((G_u-M)E_b)$ satisfies
\be\label{eq:twisted_terms}
\E\delta_a-u\sum_{b}M^{(+,+)}_{ab}\,\E\delta_b=u\sum_{b}M^{(\sB)}_{ab}\,\E\Big[\delta_b\,\tr\big((G_u-M)P_{ab}\big)\Big],\qquad a\in\Zn,
\ee
where $M^{(+,+)}_{ab}=(M^{(\sB)}_{ab})^2$. For $b\ne a$, the right-hand side contains the twisted traces $\tr((G_u-M)P_{ab})$, which are controlled neither by \eqref{Gt_avgbound_flow} nor by \eqref{Eq:L-KGt-flow}. By \Cref{lem:TL1_BA}, \eqref{Gt_avgbound_flow}, and the row bound $\max_a\sum_b|M^{(\sB)}_{ab}|\le C$, the right-hand side of \eqref{eq:twisted_terms} is $\prec(W^{-d}B_{u,0})^2$ (to pass from $\prec$ to expectations, we use the deterministic bound $\|G_u\|\le\eta_u^{-1}$ on the exceptional event). Inverting $1-uM^{(+,+)}$, whose inverse $\Theta^{(+,+)}_u$ is bounded in the $(\infty\to\infty)$-norm by \eqref{prop:ThfadC_short}, yields \eqref{res_ELK_n=1} for the block Anderson model.

% The proofs in \Cref{Sec:Steps34,Sec:Step5,Sec:Step6} carry over verbatim to the block Anderson model, with the exception of the arguments in \Cref{sec:Step5_larget} and the proof for \Cref{lem:LWterm_EXP}. In \Cref{sec:Step5_larget}, the proofs for random band matrices rely on \Cref{lem_GbEXP} and follow similar lines to those in \cite[Section 5.3]{YY_25}. In contrast, the corresponding arguments for the block Anderson model are based on \Cref{lem_GbEXP_BA} and resemble those in \cite[Section 7.3]{RBSO1D}. For this reason, we omit the details.
% For the proof of \Cref{lem:LWterm_EXP}, the arguments are almost the same as those in \Cref{lem:LWterm_EXP}, except that we will perform the $GG$ expansion in \eqref{GGGamma} below for the block Anderson model, instead of that in \eqref{Oe2x} for random band matrices. Hence, we omit the details. 





%\newpage 

\appendix
%\addtocontents{toc}{\protect\setcounter{tocdepth}{-1}}
\section{Proofs of some deterministic estimates}

This appendix is devoted to proving several deterministic estimates used in the main proof.

\subsection{Proof of \texorpdfstring{\Cref{lem_propTH}}{DeterLemma}}\label{sec:pf_propTH}

 
%\begin{proof}[Proof of Lemma \ref{lem_propTH}]\label{sec:pf_propTH}

Property 1 follows from the underlying symmetry of $S^{(\sB)}$ in the random band matrix model, where $M^{(\sig_1,\sig_2)}$ is a scalar matrix, or from the symmetry of $M^{(\sig_1,\sig_2)}$ in the block Anderson model, where \smash{$S^{(\sB)}$} is the identity matrix. Property 2 is a consequence of the translation invariance of the matrix $M^{(\sig_1,\sig_2)}S^{(\sB)}$. Property 3 follows from the fact that \smash{$\Theta^{(\sig_1,\sig_2)}_t$} is a rational function of $S^{(\sB)}$ in the random band matrix model, and from the fact that $S^{(\sB)}=I_{L^d}$ in the block Anderson model. 
To prove property 4, we expand \smash{$\Theta^{(\sig_1,\sig_2)}_t$} using the Taylor series
\be\label{eq;Taylor}
\Theta^{(\sig_1,\sig_2)}_t= \sum_{k=0}^\infty   t^k \big(M^{(\sig_1,\sig_2)}S^{(\sB)}\big)^k.\ee
Since $|M_{ab}^{(\sig_1,\sig_2)}|\le M_{ab}^{(+,-)}$ for both models, it follows that for any $\sig_1,\sig_2\in\{+,-\}$,
\[ |\Theta_{t,ab}^{(\sig_1,\sig_2)}| \le  \sum_{k=0}^\infty t^k \big(M^{(+,-)}S^{(\sB)}\big)_{ab}^k  = \Theta_{t,ab}^{(+,-)}.\]
The estimate \eqref{eq:THETAinftinf} then follows directly from this inequality and the identity $\sum_{b}\Theta_{t,ab}^{(+,-)}\equiv (1-t)^{-1}$, which holds because \smash{$M^{(+,-)}S^{(\sB)}$} is doubly stochastic. This uses $|m|=1$ in the random band matrix model, and Ward's identity \eqref{eq:WardM} in the block Anderson model. 


To prove \eqref{prop:ThfadC_short} for the random band matrix model, we use the following shifted Taylor expansion:
\be\label{eq:shiftTaylor} \Theta_t^{(+,+)} = \frac{1}{1+\e}\sum_{k=0}^\infty \p{\frac{t m^2 S^{(\sB)} + \e}{1+\e}}^k, \ee
where $\e>0$ is a positive constant. As shown in \cite[Lemma 4.2]{bourgade2019random}, one has $\|(tm^2 S^{(\sB)} + \e)/(1+\e)\|_{\infty\to \infty} \le 1-c$ for some constant $c>0$ depending on $\e$ and $\kappa$. Applying this estimate to the expansion \eqref{eq:shiftTaylor} and noticing that the off-diagonal entries of $\p{t m^2 S^{(\sB)} + \e}/\p{1+\e}$ contain a $\ilambda^2$ factor (by the definition \eqref{eq:variancematrix}), we derive the bound \eqref{prop:ThfadC_short}. For the block Anderson model, we instead use the Taylor expansion % for $\sig_1=\sig_2=+$:
\begin{align}
\left(1-tM^{(+,+)}\right)^{-1}_{0a} &=\left[\left(1-tm^2\right)I - t M'\right]^{-1}_{0a} = \sum_{k=0}^{\infty}\left(1-tm^2\right)^{-(k+1)}
 \left(t M'\right)^{k}_{0a},\label{eq:expMLn}
\end{align}
where $M':=M^{(+,+)}- m^2I$ is obtained from $M^{(+,+)}$ by setting its diagonal entries to zero.
By property (2) of \Cref{lem:propM}, there exists a constant $\e>0$ such that $|1-tm^2|\ge \e$ and
\be\label{eq:off_diagM}
\|M'\|_{\infty\to \infty}= \sum_{a\ne 0}|M'_{0a}| = 1 - |m|^2 \le (1-\e)|1-tm^2| ,\quad \forall t\in[0,1].
\ee
Applying this bound to \eqref{eq:expMLn}, we deduce the following estimate: there exists a constant $C>0$ such that for any $\delta>0$,
\begin{align}
\Big|\big(1-tM^{(+,+)}\big)^{-1}_{0a}\Big| &\le  \sum_{0\le k\le \delta |a|}\left(1-tm^2\right)^{-(k+1)}
 \left(t M'\right)^{k}_{0a} + \e^{-1}\sum_{k> \delta |a|}(1-\e)^k \nonumber\\
&\le \sum_{0\le k\le \delta |a|} C^k \exp(-c|a|) +\e^{-2}(1-\e)^{\delta|a|} \le C' e^{-c' |a|}.\label{eq:expMLn2}
\end{align}
In the second inequality, we use the exponential decay of $M'$ by \eqref{Mbound_AO} or \eqref{Mbound_AO2}, together with the convolution bound $\sum_{b}e^{-c|a_1-b|-c|a_2-b|}\lesssim e^{-c|a_1-a_2|}$ for all $a_1,a_2\in \Zn$, and in the third inequality we choose $\delta>0$ sufficiently small so that $C^{\delta}\le e^{c/2}$. The constants $c', C'$ depend only on $C,\ c,\ \e$, and $\delta$. The bound \eqref{eq:expMLn2} yields \eqref{prop:ThfadC_short} in the regime $\ilambda\gtrsim 1$. When $\ilambda\ll 1$, we return to \eqref{eq:expMLn} and note that $\|M'\|_{\infty\to\infty}\lesssim \ilambda^2$ by \eqref{Mbound_AO}. In this case, for $a\ne 0$, the same argument as in \eqref{eq:expMLn2} applies, leading to the bound \eqref{prop:ThfadC_short}.


% When $\ilambda\ll 1$, using \eqref{eq:expMLn} and $\|M'\|_{\infty\to\infty}\lesssim \lambda^2$ by \eqref{Mbound_AO}, we can derive \eqref{prop:ThfadC_short} for the case $a\ne 0$ with a similar argument as in \eqref{eq:expMLn2}. 

%\eqref{Mbound_AO} and \eqref{Mbound_AO2}, we have that \smash{$|1-tM^{(+,+)}_{00}|\gtrsim 1$} and $\|M'\|\lesssim \lambda^2$. Applying them to \eqref{eq:expMLn} and recalling that $S^{(\sB)}=I$, we can  conclude \eqref{prop:ThfadC_short}.

%The bounds \eqref{prop:BD1}--\eqref{prop:ThfadC0} are immediate consequences of \eqref{prop:ThfadC_short} when $\sig_1=\sig_2$. It remains to handle the case $\sig_1\ne \sig_2$. In this case, the bounds \eqref{prop:ThfadC0} and \eqref{prop:BD2} were established in \cite{yang2024Del} by analyzing the Fourier series of \smash{$\zTheta_t^{(+,-)}$} via a standard summation-by-parts argument.\footnote{In \cite{yang2024Del}, the assumption $\ilambda\ll 1$ was imposed; but essentially the same (and actually slightly easier) argument applies to the case $\ilambda\gtrsim 1$.} 

The bounds \eqref{prop:ThfadC} and \eqref{prop:BD1}--\eqref{prop:ThfadC0} follow directly from \eqref{prop:ThfadC_short} in the case $\sigma_1=\sigma_2$. It therefore remains to consider the case $\sigma_1\neq \sigma_2$. In this case, the estimates \eqref{prop:ThfadC0} and \eqref{prop:BD2} were proved in \cite{yang2024Del} by analyzing the Fourier series of \smash{$\zTheta_t^{(+,-)}$} through a standard summation-by-parts argument.\footnote{In \cite{yang2024Del}, the assumption $\ilambda\ll 1$ was imposed; however, the same argument (in fact, slightly simpler) applies when $\ilambda\gtrsim 1$.} 
Specifically, \eqref{prop:ThfadC0} is proved in Lemma 3.1 of \cite{yang2024Del}, while \eqref{prop:BD2} appears as equation (E.19) therein. Although the bound \eqref{prop:BD1} is not stated explicitly in \cite{yang2024Del}, its proof proceeds analogously to those of \eqref{prop:ThfadC0} and \eqref{prop:BD2}, by applying the same summation-by-parts technique to the corresponding Fourier expansion. We therefore omit the details.
These bounds \eqref{prop:BD1}--\eqref{prop:ThfadC0} have also been derived for dimension $d=2$ in \cite[Lemma 3.10]{RBSO1D}, where the summation-by-parts argument is explained in Appendix~B. The same reasoning extends directly to dimensions $d\ge 3$.
 

%More precisely, \eqref{prop:ThfadC0} is proved in Lemma 3.1 of \cite{yang2024Del}, while \eqref{prop:BD2} is proved as (E.19) in \cite{yang2024Del}. Although the estimate \eqref{prop:BD1} does not appear explicitly in \cite{yang2024Del}, its proof is similar to those for \eqref{prop:ThfadC0} and \eqref{prop:BD2} by estimating the related Fourier series using the same summation by parts argument. Hence, we omit the details of the proof. 


Finally, it remains to prove the bound \eqref{prop:ThfadC} for the case $\sig_1\ne \sig_2$. Its proof is similar to that of \cite[Lemma 2.14]{DYYY25}, using the Taylor expansion, along with the random walk representation of \smash{$(S^{(\sB)})^k$} (for the random band matrix model), or \smash{$(M^{(+,-)})^k$} (for the block Anderson model). 
More precisely, for the random band matrix model, we consider the random walk $\{X_k:k\ge 0\}$ on $\Z^d$, with transition probabilities \smash{$\big\{p(0,a)=S_{0a}^{(\sB)}:a\in \Z^d\big\}$}.  
First, given any $a\in \Z^d\setminus\{0\}$, by applying the Bernstein inequality to $X_k\cdot \wh a$ with $\wh a$ denoting the unit vector $a/\|a\|_2$, we can derive the following large deviation estimate: there exist constants $c, C>0$ (which do not depend on $a$) such that 
\[ \P_0(X_k=a) \le C\exp\p{-c \p{\frac{|a|^2}{\ilambda^2k}\wedge |a|}},\quad \forall a\in \Z^d,\ k\ge 1 . \]
On the other hand, using the local CLT for $X_k$ (see e.g., \cite[Section 2]{Lawler_book}), we obtain that 
\[ \P_0(X_k=a) \le 1\wedge \p{C (\ilambda^2k)^{-d/2}},\quad \forall a\in \Z^d,\ k\ge 1 . \]
Combining the above two bounds and using the argument below equation (8.3) of \cite{DYYY25}, we obtain the large deviation bound
\[ \P_0(X_k=a) \le C\p{ (\ilambda^2k)^{-d/2}\wedge 1}\exp\p{-c \p{\frac{|a|^2}{\ilambda^2k}\wedge |a|}},\quad \forall a\in \Z^d,\ k\ge 1,  \]
for some constants $c,C>0$. Projecting the random walk onto the torus $\Zn$, i.e., summing this bound over the periodic images of $a$, we obtain that
\begin{equation*}\label{eq:heat_torus}\tag{$\star$}
\big(M^{(+,-)}S^{(\sB)}\big)_{0a}^k \le C\p{ (\ilambda^2k)^{-d/2}\wedge 1+L^{-d}}\exp\p{-c \p{\frac{|a|^2}{\ilambda^2k}\wedge |a|}},\quad \forall a\in \Zn,\ k\ge 1,
\end{equation*}
for some constants $c,C>0$ that do not depend on $L$. 

For the block Anderson model, the matrix $P:=M^{(+,-)}S^{(\sB)}=M^{(+,-)}$, $P_{ab}=|M^{(\sB)}_{ab}|^2$, is symmetric, translation invariant, and doubly stochastic by properties (1) and (2) of \Cref{lem:propM}; in particular, $P_{0,-a}=P_{a0}=P_{0a}$. However, $P$ depends on $L$ and $m$, and it is not the periodization of a fixed transition kernel on $\Z^d$, so the above projection argument and the local CLT do not apply. Instead, we prove \eqref{eq:heat_torus} directly on the torus, with constants $c,C>0$ depending only on $d$, $\kappa$, and $\fd$. For $x\in \Zn$ and $j\in\qqq{d}$, let $|x_j|_L$ denote the periodic distance between the $j$-th coordinate of $x$ and $0$, so that $|x|=\max_j|x_j|_L$. 

\smallskip
\noindent{\it Step 1: Fourier symbol.} The eigenvalues of $\Psi^{(\sB)}$ are $\psi_q:=2\sum_{j=1}^d\cos q_j$, $q\in (\Zn)^*:=(2\pi/L)\Zn\subset(-\pi,\pi]^d$, with eigenvectors $(L^{-d/2}e^{\ii q\cdot a})_{a\in\Zn}$. Hence, setting $f_q:=(\ilambda\psi_q-E-m)^{-1}$, we have $M^{(\sB)}_{0a}=L^{-d}\sum_{q\in(\Zn)^*}f_qe^{-\ii q\cdot a}$, and Ward's identity \eqref{eq:WardM} reads $L^{-d}\sum_q|f_q|^2=1$. In particular, $|f_q|\ge 1$ for some $q$, which gives $|E|\le 2d\ilambda+2$ (recall $|m|\le 1$), and therefore $|f_q|\ge R^{-1}$ for all $q$, where $R:=4d\ilambda+3\lesssim 1$. The Fourier transform of $P_{0\cdot}$ is 
\[\wh P(p):=\sum_{a\in\Zn}P_{0a}e^{-\ii p\cdot a}=\frac{1}{L^d}\sum_{q\in(\Zn)^*} f_q\overline{f_{q+p}},\qquad p\in(\Zn)^*,\]
which is real since $P_{0,-a}=P_{0a}$. Using $L^{-d}\sum_q|f_q|^2=1$, $f_{q+p}-f_q=\ilambda(\psi_q-\psi_{q+p})f_qf_{q+p}$, and $L\ge 3$, we obtain that
\[1-\wh P(p)=\frac{1}{2L^d}\sum_q|f_{q+p}-f_q|^2\ge \frac{\ilambda^2}{2R^4}\cdot\frac{1}{L^d}\sum_q(\psi_{q+p}-\psi_q)^2=\frac{2\ilambda^2}{R^4}\sum_{j=1}^d(1-\cos p_j).\]
On the other hand, $1+\wh P(p)=\sum_a P_{0a}(1+\cos(p\cdot a))\ge 2P_{00}=2|m|^2\ge 2(\im m)^2\gtrsim 1$ by property (2) of \Cref{lem:propM}. Since $\ilambda\le\fd^{-1}$, these two bounds yield that, for a constant $c_1>0$,
\[
|\wh P(p)|\le 1-c_1\ilambda^2|p|^2,\quad \forall p\in (\Zn)^*.
\]

\smallskip
\noindent{\it Step 2: on-diagonal bound.} By Fourier inversion and the bound in Step 1, $P^{2k}_{00}=L^{-d}\sum_{p}\wh P(p)^{2k}\le L^{-d}\sum_{p}e^{-2c_1k\ilambda^2|p|^2}\le C\big((\ilambda^2k)^{-d/2}+L^{-d}\big)$, where the term $L^{-d}$ comes from $p=0$. Moreover, by the Cauchy--Schwarz inequality, the symmetry and the translation invariance of $P$, we have $P^{2k}_{0a}=\sum_bP^k_{0b}P^k_{ba}\le (P^{2k}_{00}P^{2k}_{aa})^{1/2}=P^{2k}_{00}$ and $P^{2k+1}_{0a}=\sum_bP_{0b}P^{2k}_{ba}\le P^{2k}_{00}$. Together with $P^k_{0a}\le 1$, this gives
\[
\max_{a}P^k_{0a}\le C\p{(\ilambda^2k)^{-d/2}\wedge 1+L^{-d}},\quad \forall k\ge 1. 
\]

\smallskip
\noindent{\it Step 3: tail bound.} Fix $j\in\qqq{d}$ and $\mu\ge 0$, and let $F_\mu(x):=\cosh(\mu|x_j|_L)$. For $x,y\in\Zn$, choose integer representatives $\wt x_j,\wt y_j$ of $x_j,y_j$ with $|\wt x_j|=|x_j|_L$ and $|\wt y_j|=|y_j|_L$. Since $|x_j\pm y_j|_L\le |\wt x_j\pm \wt y_j|$, we have $F_\mu(x+y)+F_\mu(x-y)\le 2F_\mu(x)\cosh(\mu |y_j|_L)$. Together with $P_{0,-y}=P_{0y}$, this implies
\[(PF_\mu)(x)=\frac12\sum_y P_{0y}\big(F_\mu(x+y)+F_\mu(x-y)\big)\le F_\mu(x)\sum_y P_{0y}\cosh(\mu|y|)\le \p{1+C\ilambda^2\mu^2}F_\mu(x),\quad \forall\, 0\le\mu\le c,\]
where in the last step we used $\sum_{y\ne 0}P_{0y}|y|^2e^{2c|y|}\le C\ilambda^2$, which follows from \eqref{Mbound_AO} when $\ilambda<(2C)^{-1}$ and from \eqref{Mbound_AO2} when $\ilambda\ge (2C)^{-1}$. Iterating this bound and applying the exponential Chebyshev inequality, we get $\sum_{x:|x_j|_L\ge r}P^k_{0x}\le e^{C\ilambda^2\mu^2k}/\cosh(\mu r)$ for all $r\ge 0$. Optimizing over $\mu\in[0,c]$ and taking a union bound over $j\in\qqq{d}$, we obtain that
\[
\sum_{x:|x|\ge r}P^k_{0x}\le 2d\exp\p{-c\p{\frac{r^2}{\ilambda^2k}\wedge r}},\quad \forall r\ge 0,\ k\ge 1.
\]

\smallskip
\noindent{\it Step 4: conclusion.} For $k\ge 2$, write $k=k_1+k_2$ with $k_1=\lceil k/2\rceil$ and $k_2=\lfloor k/2\rfloor$. Since $|b|\ge |a|/2$ or $|a-b|\ge |a|/2$ for every $b\in\Zn$, we have
\[P^k_{0a}=\sum_b P^{k_1}_{0b}P^{k_2}_{ba}\le \max_{x}P^{k_2}_{0x}\sum_{b:|b|\ge |a|/2}P^{k_1}_{0b}+\max_xP^{k_1}_{0x}\sum_{b:|b-a|\ge|a|/2}P^{k_2}_{ab}.\]
Applying the bounds of Steps 2 and 3 to the right-hand side yields \eqref{eq:heat_torus} for $k\ge 2$, while the case $k=1$ follows directly from \eqref{Mbound_AO} and \eqref{Mbound_AO2}. 

\smallskip
Finally, for both models, applying \eqref{eq:heat_torus} to the expansion \eqref{eq;Taylor} and summing over $k$, we can derive the desired bound \eqref{prop:ThfadC}.
Since the argument closely follows that in \cite[Section 8]{DYYY25}, we omit the details. In fact, the proof here is somewhat simpler than in \cite{DYYY25}, because $d=2$ is the critical dimension, where logarithmic corrections (e.g., $\log L$) appear. In contrast, for $d\ge 3$, all relevant series (or integrals) are summable, resulting in a dimension-dependent constant $C_d$ in \eqref{prop:ThfadC}.


%\end{proof}

%With this bound, we can control $(S^{(\sB)})_{0a}^k$ by projecting this random walk onto the torus $\Zn$. Then, summing over these estimates and using \eqref{eq;Taylor}, we can derive the bound \eqref{prop:ThfadC}. Since the argument is similar to that in \cite{DYYY25}, we omit the details. In fact, the proof here is relatively easier than that in \cite{DYYY25}, because $d=2$ is the critical dimension where some $\log L$ factors would appear. On the other hand, for dimensions $d\ge 3$, all related series are summable, which leads to the constant $C_d$ in \eqref{prop:ThfadC}. 

%{The proof of Lemma \ref{lem_propTH} is based on the random walk representation and pointwise estimates for $(S^{(B)})^{k}$ that agree with the continuum Gaussian heat kernel estimates. A proof is given in Section \ref{sec_TH}; the argument is long and elementary, so we give an intuitive explanation of Lemma \ref{lem_propTH} below.

% The first four properties in Lemma \ref{lem_propTH} are standard. For the exponential decay \eqref{prop:ThfadC} for the discrete Laplacian on $\Z^{2}$ (instead of $S^{(B)}-\mathrm{Id}$) can be found in Theorem 1.3 of \cite{MS22} (see also the discussion following it therein). The difference between the full space $\Z^{2}$ and the torus $\Z_{L}^{2}$ is that if $\ell_t=L$, so that $|1-t|^{-1/2}\geq L$, then the periodic boundary conditions on $\Z_{L}^{2}$ automatically introduces a space-cutoff of scale $L$ (before the cutoff at scale $|1-t|^{-1/2}$ that is imposed by the resolvent). One can also check that \eqref{prop:ThfadC} implies 
% \begin{align*}
% \sum_{b\in\Z_{L}^{2}}(\Theta^{(B)}_{\xi})_{ab}\prec|1-t|^{-1},
% \end{align*}
% which can be checked using the random walk representation as well. 


%The gradient bounds \eqref{prop:BD1}-\eqref{prop:BD2} are standard, in that the Green's function for the Laplacian on $\R^{2}$ is $\log|x|$. Thus, the first terms on the right-hand side of \eqref{prop:BD1}-\eqref{prop:BD2} agree with its derivatives. (The support-length of the gradients is also of order $\ell_t$, but we will not need this fact.) The additional corrections come from the periodic boundary of $\Z_{L}^{2}$. The scaling of the last term in \eqref{prop:BD2} agrees with the heuristic that $\Theta^{(B)}$ is the resolvent for the diffusion of a generator, so its second-derivative should have $L^{1}$ norm of size $\mathrm{O}_{\prec}(1)$. The last term in \eqref{prop:BD1} therefore interpolates between \eqref{prop:ThfadC} and the last term in \eqref{prop:BD2}.}
\subsection{Proofs of evolution kernel estimates}\label{ks}

In this subsection, we present the proofs of \Cref{lem:sum_Ndecay,lem:sum_decay,lem:sum_decay_nonzero}. 
% During the main proof, we rely on several estimates for the evolution kernel (defined in \Cref{DefTHUST}), which are stated in this section and used to control the terms in \eqref{int_K-LcalE}. 
Parts of the proofs parallel those in \cite{YY_25,DYYY25,RBSO1D} for random band matrices and for the block Anderson model in dimensions 1 and 2. However, certain key arguments must be adapted to handle the higher-dimensional setting $d \ge 3$. We therefore outline the proofs of the evolution kernel estimates, emphasizing the modifications needed compared to the arguments in \cite{YY_25,DYYY25,RBSO1D}. 
%We first have an easy bound on the $({\infty\to \infty})$-norm of $\mathcal U^{(n)}_{t,\bsig,\ba}$. 

 % \begin{lemma}\label{lem:sum_Ndecay}
	%  	Let ${\cal A}: (\Zn)^{n}\to \mathbb C$ be an $n$-dimensional tensor for a fixed $n\in \N$ with $n\ge 2$. Then, for each $0\le s \le t < 1$, we have that 
	%  	\begin{align}\label{sum_res_Ndecay}
	% 	 		\| {\cal U}^{(n)}_{s,t,\bsig}\;\circ {\cal A}\|_{\infty} \le \left(  \frac{1-s}{1-t}\right)^{n }\cdot \|{\cal A}\|_{\infty} , 
	% 	 	\end{align}
	%  	where the $L^\infty$-norm of ${\cal A}$ is defined as $\|{\cal A}\|_{\infty}=\max_{\ba\in (\Zn)^n}|\cal A_{\ba}|$.
	%  \end{lemma}
\begin{proof}[\bf Proof of \Cref{lem:sum_Ndecay}]
	Notice the following decomposition:
	\be\label{eq:decompUalt}\frac{1 - s \cdot M^{(\sig_i,\sig_{i+1})}S^{(\sB)}}{1 - t \cdot M^{(\sig_i,\sig_{i+1})}S^{(\sB)}} =1+ \Xi^{(i)},\quad \text{where}\quad \Xi^{(i)}:=(t-s)\cdot M^{(\sig_i,\sig_{i+1})}S^{(\sB)} \Theta_{t}^{(\sig_i,\sig_{i+1})} .\ee
	Using \eqref{eq:THETAinftinf}, along with \eqref{eq:WardM} in the case of block Anderson model, we obtain that
	\be\label{Xi_infint} \|\Xi^{(i)}\|_{\infty\to \infty} = \max_{a}\sum_b |\Xi^{(i)}_{ab}| \le \p{t-s}\|\Theta^{(\sig_i,\sig_{i+1})}_{t}\|_{\infty\to \infty} \le \frac{t-s}{1-t}. \ee
	Together with \eqref{eq:decompUalt}, this implies that 
	$$\left\|\frac{1 - s \cdot M^{(\sig_i,\sig_{i+1})}S^{(\sB)}}{1 - t \cdot M^{(\sig_i,\sig_{i+1})}S^{(\sB)}}\right\|_{\infty\to \infty}\le \frac{1-s}{1-t} .$$
	With this estimate, we conclude \eqref{sum_res_Ndecay} immediately using the definition \eqref{def_Ustz}. 
	\end{proof}


\begin{proof}[\bf Proof of \Cref{lem:sum_decay}]
With the decomposition \eqref{eq:decompUalt}, we can express \smash{$\mathcal{U}^{(n)}_{s, t, \bsig} \circ \mathcal{A}$} as
\begin{align}\label{eq:decomp_U2}
      \left(\mathcal{U}^{(n)}_{s, t, \bsig} \circ \mathcal{A}\right)_{\ba}&=\sum_{\mathbf b\in (\Zn)^n} \prod_{i=1}^n \left( \delta_{a_ib_i}+\Xi^{(i)}_{a_ib_i} \right)  \cdot \mathcal{A}_{\mathbf b} =\sum_{\mathbf b\in (\Zn)^n}\sum_{A\subset \qqq{n}} \prod_{i\in A} \delta_{a_ib_i} \cdot \prod_{i\in A^c} \Xi^{(i)}_{a_ib_i}  \cdot \mathcal{A}_{\mathbf{b}}\, .
\end{align}
By the estimate \eqref{prop:ThfadC} (along with \eqref{Mbound_AO} or \eqref{Mbound_AO2} in the case of block Anderson model), we have that
\be\label{eq:decayXi}
\Xi^{(i)}_{a_ib_i}\lesssim (1-s)\frac{(\ilambda^2+|1-t|)^{-1}}{(|a_i-b_i|+1)^{d-2}} e^{-c |a_i-b_i|/ {\ell}_t} 
\ee
for a constant $c>0$. 
%Moreover, we have the following bound by \eqref{eq:THETAinftinf} (along with \eqref{eq:WardM} in the case of block Anderson model): 
%\be\label{Xi_infint} \|\Xi^{(i)}\|_{\infty\to \infty} = \max_{a}\sum_b |\Xi^{(i)}_{ab}| \le \p{1-s}\|\Theta^{(\sig_i,\sig_{i+1})}_{t}\|_{\infty\to \infty} = \frac{1-s}{1-t}. \ee
We claim that for any subset $A$ with $|A|=k\in \qqq{1,n}$, 
\be\label{sum_res_1_red0}
\sum_{\mathbf b\in (\Zn)^n} \prod_{i\in A} \delta_{a_ib_i} \cdot \prod_{i\in A^c} \Xi^{(i)}_{a_ib_i}  \cdot \mathcal{A}_{\boldsymbol{b}} \le W^{C\e}\p{\frac{\ilambda^2+|1-s|}{\ilambda^2+|1-t|}}^{n-k} \|{\cal A}\|_{\infty}+ W^{-D+C },
\ee
for a constant $C$ that does not depend on $\e$ or $D$, while for $A= \emptyset$, we claim that 
\be\label{sum_res_1_red}
\sum_{\mathbf b\in (\Zn)^n} \prod_{i=1}^n \Xi^{(i)}_{a_ib_i}  \cdot \mathcal{A}_{\boldsymbol{b}} \le W^{C\e} \frac{\ell_t^2}{\ell_s^2}\p{\frac{\ilambda^2+|1-s|}{\ilambda^2+|1-t|}}^{n} \|{\cal A}\|_{\infty}+ W^{-D+C}.
\ee
Note that combining \eqref{sum_res_1_red0} and \eqref{sum_res_1_red} concludes the proof of \eqref{sum_res_1}.   

To show the estimate \eqref{sum_res_1_red0}, we assume that $A=\qqq{1,k}$ without loss of generality. With the notations $\mathbf a'=(a_{1},\ldots,a_{k})$ and $\mathbf b'=(b_{k+1},\ldots,b_{n})$, we can bound the LHS of \eqref{sum_res_1_red0} as 
\begin{align*}
\sum_{\mathbf b'\in (\Zn)^{n-k}} \prod_{i=k+1}^n  \Xi^{(i)}_{a_ib_i}  \cdot \mathcal{A}_{\ba',\mathbf{b}'} &\lesssim  \frac{\left(1-s\right)^{n-k}}{(\ilambda^2+|1-t|)^{n-k}} \|\cal A\|_\infty \sum_{\mathbf b'} \prod_{i=k+1}^n\frac{\mathbf 1\left(|b_i-a_1|\le W^\e\ell_s \right)}{|a_i-b_i|^{d-2}+1} +  W^{-D+(n-k)}\\
    &\lesssim (W^{2\e})^{n-k}\left( \frac{\ell_s^2|1-s|}{\ilambda^2+|1-t|} \right)^{n-k} \|\cal A\|_\infty  + W^{-D+(n-k)} \\
    & \lesssim \p{W^{2\e}\frac{\ilambda^2+|1-s|}{\ilambda^2+|1-t|}}^{n-k} \|\cal A\|_\infty  + W^{-D+(n-k)} , 
\end{align*}
where we have used \eqref{eq:decayXi}, the decay property \eqref{deccA0} for $\cal A_{\mathbf b}$, and 
\be\label{eq:1-sells2}(1-s)\ell_s^2\asymp \ilambda^2+|1-s|,\quad \text{for}\quad s\le 1-\ilambda^2/L^{2},\ee 
along with the condition $(1-t)/(1-s)\ge W^{-1}$ and the bound \eqref{Xi_infint} in getting the $W^{-D+(n-k)}$ term. This concludes \eqref{sum_res_1_red0} for any constant $C>2(n-k)$.  
% Without loss of generality, we only prove \eqref{sum_res_1_red} for $k=n$; to handle the $k<n$ case, we can apply the same argument to  
% $$\sum_{\mathbf b'\in (\Zn)^k}  \prod_{i\in A} \Xi^{(i)}_{[a_i][b_i]}  \cdot \mathcal{A}'_{\boldsymbol{b}'},$$
% where $\mathcal{A}'_{\boldsymbol{b}'}$ is the $k$-dimensional tensor obtained by setting $[b_i]=[a_i]$ for $i\in A^c$. 
For \eqref{sum_res_1_red}, with the notation $\mathbf b'=(b_{2},\ldots,b_{n})$, we get that 
\begin{align}
   \sum_{\mathbf b}\prod_{i=1}^{n} \Xi^{(i)}_{a_ib_i} \cdot \mathcal{A}_{\mathbf b} &= \sum_{\mathbf b}\prod_{i=1}^{n} \Xi^{(i)}_{a_ib_i} \cdot \mathcal{A}_{\mathbf b}\cdot  \mathbf 1\left(\max_{i\ne j}|b_i-b_j|\le W^\e\ell_s \right) + \OO\left( W^{-D+n}\right) \nonumber\\
   &\lesssim \frac{(1-s)^{n-1}}{(\ilambda^2+|1-t|)^{n-1}} \|\cal A\|_\infty \sum_{b_1} \left|\Xi^{(1)}_{a_1b_1}\right| \sum_{\mathbf b'} \prod_{i=2}^n \frac{\mathbf 1\left(|b_i-b_1|\le W^\e\ell_s \right)}{|a_i-b_i|^{d-2}+1} +  W^{-D+n} \nonumber\\
    %&\lesssim W^{2(n-1) \e }\left(\frac{\ell_s^2|1-s|}{\ilambda^2+|1-t|}\right)^{n-1} \|\cal A\|_\infty \sum_{b_1} \left|\Xi^{(1)}_{a_1b_1}\right| + W^{-D+n} \nonumber\\
    &\lesssim W^{2(n-1) \e}\frac{1-s}{1-t}\left(\frac{\ell_s^2|1-s|}{\ilambda^2+|1-t|}\right)^{n-1} \|\cal A\|_\infty + W^{-D+n} \nonumber\\
    &\lesssim  W^{2(n-1) \e}\frac{\ell_t^2}{\ell_s^2}\left(\frac{\ilambda^2+|1-s|}{\ilambda^2+|1-t|}\right)^{n} \|\cal A\|_\infty + W^{-D+n}, \label{sum_res_deriv_red}
\end{align} 
where in the first step, we use the decay property \eqref{deccA0} for $\cal A_{\mathbf b}$ and the bound \eqref{Xi_infint} (in getting the $W^{-D+n}$ term); in the second step, we apply \eqref{eq:decayXi} for $i\in\qqq{2,n}$; in the third step, we take the summation over $\mathbf b'$ and apply \eqref{Xi_infint}; in the last step, we use \eqref{eq:1-sells2} again.  This gives \eqref{sum_res_1_red} for any constant $C>2(n-1)$.  


For the estimates \eqref{sum_res_2_NAL} and \eqref{sum_res_2}, %\eqref{sum_res_2_sym}, and \eqref{sum_res_2_doublezero}, 
we notice that when $|A| \ge 1$, \eqref{sum_res_1_red0} already gives a good enough bound. 
Hence, we only need to focus on the case where $A=\emptyset$. 
First, for Case I, due to \eqref{prop:ThfadC_short}, we can use the better bound \smash{$\sum_{b_1}|\Xi^{(1)}_{a_1b_1}|=\OO(1)$} in the third step of \eqref{sum_res_deriv_red}, which leads to \eqref{sum_res_2_NAL}. 

Next, for the estimate \eqref{sum_res_2} in Case II, we need to show that
\be\label{sum_res_2_red}
\sum_{\mathbf b\in (\Zn)^n} \prod_{i=1}^n \Xi^{(i)}_{a_ib_i}  \cdot \mathcal{A}_{\boldsymbol{b}} \le W^{C_n\e}\p{\frac{\ilambda^2+|1-s|}{\ilambda^2+|1-t|}}^{n}  \|{\cal A}\|_{\infty} +W^{-D+C_n}.
\ee
It suffices to assume that $\sigma_i\ne \sigma_{i+1}$ for all $i\in \qqq{n}$. 
We decompose $\Zn$ into the following two regions:
\[S_{\mathrm{far}}:=\Big\{b\in \Zn:\min_{i\in\qqq{n}}|b-a_i|>  W^{2\e}\ell_s\Big\},\quad S_{\mathrm{near}}:=\Big\{b\in \Zn:\min_{i\in\qqq{n}}|b-a_i|\le W^{2\e}\ell_s\Big\}.\]
Following a similar argument as in \eqref{sum_res_deriv_red}, and using that \smash{$\sum_{b_1\in S_{\mathrm{near}}} |\Xi^{(1)}_{a_1b_1}|\lesssim W^{4\e}(\ilambda^2+|1-s|)/(\ilambda^2+|1-t|)$} by \eqref{eq:decayXi} and \eqref{eq:1-sells2}, we obtain
\begin{align}
\sum_{b_1\in S_{\mathrm{near}}} \sum_{\bfb'}\prod_{i=1}^{n} \Xi^{(i)}_{a_ib_i} \cdot \mathcal{A}_{\mathbf b} \lesssim W^{2(n+1) \e} \left(\frac{\ilambda^2+|1-s|}{\ilambda^2+|1-t|}\right)^{n}\|\cal A\|_\infty + W^{-D+n}. \label{sum_res_deriv_red2}
\end{align} 
It remains to control the sum over $b_1\in S_{\mathrm{far}}$. In this case, we decompose $\Xi^{(i)}$ as 
\be\nonumber \Xi^{(i)}_{a_ib_i}=\Xi^{(i)}_{a_ib_1} + \Delta \Xi^{(i)}_{a_i;b_1b_i},\quad\text{with}\quad \Delta \Xi^{(i)}_{a_i;b_1b_i}:=\Xi^{(i)}_{a_ib_i}-\Xi^{(i)}_{a_ib_1}. \ee
Then, we expand the LHS of \eqref{sum_res_2_red} as 
\begin{align} \label{eq:decompXii}
\sum_{A: A\subset \qqq{2,n}}f(A),\quad \text{with}\quad 
f(A):=\sum_{b_1\in S_{\mathrm{far}}}\sum_{\bfb'\in (\Zn)^{n-1}}\prod_{i\in A^c}  \Xi^{(i)}_{a_ib_1} \cdot \prod_{i\in A} \Delta \Xi^{(i)}_{a_i;b_1b_i}  \cdot \mathcal{A}_{\mathbf{b}}. 
\end{align}
By the sum-zero property \eqref{sumAzero}, the leading term with $A=\emptyset$ vanishes. For the remaining terms, we will use  \eqref{eq:decayXi} to control the factors \smash{$\Xi^{(i)}_{a_ib_1}$}, and apply \eqref{eq:decayXi} and \eqref{prop:BD1} (along with \eqref{Mbound_AO} or \eqref{Mbound_AO2} in the case of block Anderson model) to control the factors %\smash{$\Xi^{(i)}_{a_ib_1}$} and 
\smash{$\Delta \Xi^{(i)}_{a_i;b_1b_i}$} for $|a_i-b_1|\ge W^{2\e}\ell_s$ and $|b_i-b_1|\le W^{\e}\ell_s$:
\be\label{eq:Xibb}
\Xi^{(i)}_{a_ib_1}\lesssim \frac{W^{-(d-1)\e}}{\ell_s^d}\left(\frac{\ilambda^2+|1-s|}{\ilambda^2+|1-t|}\right),\   \Delta \Xi^{(i)}_{a_i;b_1b_i} \lesssim  \frac{1-s}{\ilambda^2+|1-t|}\frac{|b_i-b_1|}{|a_i-b_1|^{d-1}}  \prec \frac{W^{-d\e}}{\ell_s^d}\left(\frac{\ilambda^2+|1-s|}{\ilambda^2+|1-t|}\right),
\ee
where we have also used $2(d-2)\ge d-1$ and \eqref{eq:1-sells2} in the derivation. 
Without loss of generality, suppose $2\in A$. Then, using the estimates in \eqref{eq:decayXi} and \eqref{eq:Xibb}, along with the decay property \eqref{deccA0} for $\cal A$, by a similar argument as in \eqref{sum_res_deriv_red}, we can bound $f(A)$ by 
\begin{align}
f(A)&\lesssim W^{\e (n-1)} \frac{(1-s)(\ilambda^2+|1-s|)^{n-2}}{(\ilambda^2+|1-t|)^{n-1}}  \|\cal A\|_\infty 
\sum_{b_1\in S_{\mathrm{far}}} \left|\Xi^{(1)}_{a_1b_1}\right| \frac{ (W^\e\ell_s)^{d+1}}{|a_2-b_1|^{d-1}} + W^{-D+n}L^{d(n-1)}\nonumber\\
&\lesssim  W^{(n+d+1) \e }\frac{(1-s)^2\ell_s^{d+1}(\ilambda^2+|1-s|)^{n-2}}{(\ilambda^2+|1-t|)^{n}}  \|\cal A\|_\infty  
\sum_{b_1\in S_{\mathrm{far}}} \frac{\exp(-c|a_1-b_1|/\ell_t)}{|a_1-b_1|^{d-2} \cdot |a_2-b_1|^{d-1}} + W^{-D+n}L^{d(n-1)} \nonumber\\
&\lesssim  W^{(n+d+1)\e } \|\cal A\|_\infty \cdot  \frac{(1-s)^2\ell_s^{d+1}}{(W^{2\e}\ell_s)^{d-3}}\log L\cdot \frac{(\ilambda^2+|1-s|)^{n-2}}{(\ilambda^2+|1-t|)^{n}} + W^{-D+n}L^{d(n-1)} \nonumber\\
&\lesssim W^{(n+5)\e }  \p{\frac{\ilambda^2+|1-s|}{\ilambda^2+|1-t|}}^{n} \|\cal A\|_\infty+ W^{-D+n}L^{d(n-1)}. 
\label{eq:bddfA}
\end{align}
In the third step, we used the lattice sum estimate
\be\label{eq:latticesum_d3}
\sum_{b\,:\,|a_1-b|\wedge|a_2-b|> R} \frac{\exp(-c|a_1-b|/\ell_t)}{|a_1-b|^{d-2} \cdot |a_2-b|^{d-1}} \lesssim \frac{\log L}{R^{d-3}},\quad \forall\, 1\le R\le L,
\ee
which holds for every $d\ge 3$; the logarithm is needed only in the borderline case $d=3$, where the summand decays exactly like $|b|^{-d}$ and the sum over $R\le |b|\lesssim \ell_t$ is of order $\log(\ell_t/R)$. In the last step we absorbed it using $\log L\le W^{\e}$, together with \eqref{eq:1-sells2} in the form $(1-s)^2\ell_s^4\asymp(\ilambda^2+|1-s|)^2$. The error term $W^{-D+n}L^{d(n-1)}$ comes from the region where $|\cal A_{\mathbf b}|\le W^{-D}$ by \eqref{deccA0}: there, since $\Xi^{(i)}_{a_ib_1}$ does not depend on $b_i$, we have $\sum_{b_i}\big(|\Xi^{(i)}_{a_ib_1}|+|\Delta \Xi^{(i)}_{a_i;b_1b_i}|\big)\le (2L^d+1)\|\Xi^{(i)}\|_{\infty\to\infty}$ for $i\in\qqq{2,n}$, so the free summations over $b_i\in\Zn$ produce factors $L^d$; the remaining factors are controlled by \eqref{Xi_infint} and the condition $(1-t)/(1-s)\ge W^{-1}$. Since $L^d\le N\le W^{1/\fc}$ by \eqref{Main_DEL_COND}, this term is bounded by $W^{-D+n+(n-1)/\fc}$. 
This concludes the estimate \eqref{sum_res_2_red} with a constant $C_n$ that depends only on $n$, $d$, and $\fc$ (and not on $\e$ or $D$), which further concludes \eqref{sum_res_2}. 
\end{proof}

%\subsection{Evolution kernel estimates without zero mode}


\begin{proof}[\bf Proof of \Cref{lem:sum_decay_nonzero}]
Let $\mathbf e\in \C^{L^d}$ denote the unit vector with $\mathbf e(a)\equiv L^{-d/2}$, and define the projection matrix ${\mathsf{Proj}}_{\mathbf e^{\perp}}$ onto the orthogonal complement of $\mathbf e$:
\( {\mathsf{Proj}}_{\mathbf e^{\perp}} =I-\mathbf e\mathbf e^\top.\) 
%For the decomposition \eqref{eq:decompUalt}, 
When $\sig_i=\sig_{i+1}$, we have
\be\label{eq:samecolor} \|\Xi^{(i)}\|_{\infty\to\infty} \lesssim t-s,\quad \text{and}\quad \|{\mathsf{Proj}}_{\mathbf e^{\perp}}\cdot  \Xi^{(i)}\|_{\infty\to\infty} \lesssim t-s,\ee
by using the definition of $\Xi^{(i)}$ in \eqref{eq:decompUalt} and the estimate \eqref{prop:ThfadC_short}. 
On the other hand, when $\sig_i\ne \sig_{i+1}$, we have  \[{\mathsf{Proj}}_{\mathbf e^{\perp}}\cdot  \Xi^{(i)}= (t-s) {\mathsf{Proj}}_{\mathbf e^{\perp}}\cdot M^{(+,-)}S^{(\sB)}\Theta_t^{(+,-)} = (t-s)M^{(+,-)}S^{(\sB)} \zTheta_t^{(+,-)}.\] 
% \begin{align*}
% {\mathsf{Proj}}_{\mathbf e^{\perp}}\cdot  \Xi^{(i)}= (t-s) {\mathsf{Proj}}_{\mathbf e^{\perp}}\cdot\Theta_t^{(\sB)} = (t-s) \zTheta_t^{(\sB)}.
% \end{align*}
Then, using \eqref{prop:ThfadC0}  (along with \eqref{eq:WardM} in the case of block Anderson model), we get 
\be\label{eq:diffcolor}
\|{\mathsf{Proj}}_{\mathbf e^{\perp}}\cdot  \Xi^{(i)}\|_{\infty\to\infty} \lesssim (t-s) \max_{a}\sum_{b}\big|\zTheta_{t,ab}^{(+,-)}\big| \prec (1-s)\cdot \ilambda^{-2}L^2 \le 1,
\ee
where we use $1-s\le \ilambda^2/L^{2}$ in the last step.
Plugging \eqref{eq:samecolor} and \eqref{eq:diffcolor} into \eqref{eq:decompUalt}, we obtain that 
%\[ \left\| \frac{1 - s \cdot m(\sig_i)m(\sig_{i+1})S^{(\sB)}}{1 - t \cdot m(\sig_i)m(\sig_{i+1})S^{(\sB)}}\right\|_{\infty\to\infty}\lesssim 1\]
%if $\sig_{i}=\sig_{i+1}$, and 
\[ \left\|{\mathsf{Proj}}_{\mathbf e^{\perp}}\cdot  \frac{1 - s \cdot M^{(\sig_i,\sig_{i+1})}S^{(\sB)}}{1 - t \cdot M^{(\sig_i,\sig_{i+1})}S^{(\sB)}}\right\|_{\infty\to\infty}\lesssim 1\]
for all $i\in \qqq{n}$.
Finally, we deduce \eqref{sum_res_Ndecay_nonzero} as follows. By \eqref{def_Ustz}, the kernel \smash{${\cal U}^{(n)}_{s,t,\bsig}$} acts on each of the $n$ indices separately, and by \Cref{def;zero_mode_remove} the operator $Q^{(A)}=\prod_{i\in A}Q^{(i)}$ acts on the indices in $A$ only, with $Q^{(i)}$ acting on the $i$-th index as \smash{${\mathsf{Proj}}_{\mathbf e^{\perp}}$}. Since \smash{${\mathsf{Proj}}_{\mathbf e^{\perp}}$} commutes with the translation-invariant matrix $M^{(\sig_i,\sig_{i+1})}S^{(\sB)}$, we may therefore write \smash{$Q^{(A)}\circ{\cal U}^{(n)}_{s,t,\bsig}$} as the tensor product over $i\in\qqq{n}$ of the one-index operators
\[{\mathsf{Proj}}_{\mathbf e^{\perp}}\cdot  \frac{1 - s \cdot M^{(\sig_i,\sig_{i+1})}S^{(\sB)}}{1 - t \cdot M^{(\sig_i,\sig_{i+1})}S^{(\sB)}}\ \ (i\in A),\qquad  \frac{1 - s \cdot M^{(\sig_i,\sig_{i+1})}S^{(\sB)}}{1 - t \cdot M^{(\sig_i,\sig_{i+1})}S^{(\sB)}}\ \ (i\notin A).\]
The operators of the first kind are bounded in $(\infty\to\infty)$-norm by the display above. For those of the second kind, the assumption $A\supset I_{\mathrm{diff}}(\bsig)$ forces $\sig_i=\sig_{i+1}$, so \eqref{eq:decompUalt} and the first bound in \eqref{eq:samecolor} give $\|1+\Xi^{(i)}\|_{\infty\to\infty}\lesssim 1$. Multiplying the $n$ one-index bounds yields \eqref{sum_res_Ndecay_nonzero}. 
\end{proof}

\subsection{Proof of \Cref{lem:propT}}\label{sec:pfpropT}
%\begin{proof}
First, we consider the case $1-t \le  1-u \le \ilambda^2/L^{2}$, where we have $\ell_t=\ell_u=L$. In this case, the exponential factor is of order 1, and we have 
\begin{align}
	\sum_{c}{\cal T}_{u}(|a-c|)\cdot {\cal T}_{t}(|c-b|) &\lesssim \sum_{c} \left(\frac{\ilambda^{-2}}{|a-c|^{d-2}+1}+\frac{1}{L^d|1-u|}\right)\left(\frac{\ilambda^{-2}}{|c-b|^{d-2}+1}+\frac{1}{L^{d}|1-t|}\right) \nonumber
	\\\nonumber
	&\lesssim  \frac{\ilambda^{-4}L^2}{|a-b|^{d-2}+1} + \frac{\ilambda^{-2}}{L^{d-2}|1-t|}+\frac{1}{1-u}\frac{1}{L^d |1-t|} 
	\\
	 &\lesssim   \frac{1}{1-u}\left(
	\frac{\ilambda^{-2}}{|a-b|^{d-2}+1}+\frac{1}{L^d |1-t|}\right)
	\lesssim   \frac{1}{1-u} {\cal T}_{t}(|a-b|),\label{eq:smalletacase}
\end{align}   
where, in the third and fourth steps, we use that $1-u\le \ilambda^2/L^2$. Next, we consider the case $1-u\ge 1-t \ge \ilambda^2/L^{2}$, where we have $\ell_u \le \ell_t=\max(\ilambda(1-t)^{-1/2},1)\le L$. In this case, the term $(L^d|1-t|)^{-1}$ can always be neglected in the function $\cT_t$. Then, we get that
\begin{align}
&~	 {{\cal T}_{u}(|a-c|)\cdot {\cal T}_{t}(|c-b|)}\big/{{\cal T}_{t}(|a-b|)} \nonumber\\
\lesssim &~ \frac{1}{\ilambda^2+|1-u|}\Bigg\{1 + \sum_{c\notin \{a,b\}}  \frac{(|a-b|+1)^{d-2}}{|a-c|^{d-2}|c-b|^{d-2}  }\exp \left(-\frac{\sqrt{|a-c|}+
	(\ell_u/\ell_t)^{1/2}(\sqrt{|c-b|}-\sqrt{|a-b|})}{\ell_u^{1/2}} \right)\Bigg\} \nonumber\\
\lesssim &~ {\ell_u^2}/\p{\ilambda^2+|1-u|} \lesssim |1-u|^{-1},\label{eq:largeetacase}
\end{align}
where, in the second step, we use the following basic calculus fact for $d\ge 3$ and any $0\le \e \le 1$:
$$
\max_{a\in \mathbb R^d} \int_{x\in \mathbb R^d}\frac{|a|^{d-2}}{|a-x|^{d-2}\cdot |x|^{d-2}}\exp\left(-\sqrt{|a-x|}-\e \p{\sqrt{|x|}-\sqrt{|a|}} \right)\dd x\le C_d$$
for a constant $C_d>0$. Combining the two cases \eqref{eq:smalletacase} and \eqref{eq:largeetacase} completes the proof of \eqref{TTT2}. 
%\end{proof}

%\begin{align}
%\frac{{\cal T}_{t}(|a-c|)\cdot {\cal T}_{t}(|c-b|)}{{\cal T}_{t}(|a-b|)} &\lesssim 1 + \sum_{c\notin \{a,b\}}  \frac{|a-b|^{d-2}}{|a-c|^{d-2}|c-b|^{d-2}  }\exp \left(-\frac{|a-c|^{1/2}+|c-b|^{1/2}-|a-b|^{1/2}}{\ell_t^{1/2}} \right) \nonumber\\
%    &\lesssim \ell_t^2 =|1-t|^{-1},\label{eq:largeetacase}
%   \end{align}
%where, in the second step, we used the following basic calculus fact for $d\ge 3$: %(It can be proved with simple changing variables) 
%$$
% \max_{a\in \mathbb R^d} \int_{x\in \mathbb R^d}\frac{|a|^{d-2}}{|a-x|^{d-2}\cdot |x|^{d-2}}\exp\left(-\sqrt{|a-x|}-\sqrt{|x|}+\sqrt{|a|} \right)dx\le C_d$$
% for a constant $C_d>0$. Combining the two cases \eqref{eq:smalletacase} and \eqref{eq:largeetacase} completes the proof of \eqref{TTT}.


\subsection{Proof of \Cref{claim:TTk}}\label{pf:claim_TTk}
To prove the estimate \eqref{eq:key_T_reudce}, we begin by partitioning the summation region over $[\al]$ into $2^{2k}$ subregions according to whether each $|[x_i]-[\al]|$ or $|[y_i]-[\al]|$ is larger than $\ell$ or not. Namely, we define
\[ \bD_{\le \ell,\bsig}:=\left\{[\al]:  |[w_{i}]-[\al]|\le \ell \ \text{if}\ \sig_i=0, \text{ and } |[w_{i}]-[\al]|>\ell \ \text{if}\ \sig_i=1, \ \forall i\in\qqq{2k} \right\},\]
where $\bsig=(\sig_1,\ldots, \sig_{2k}) \in \{0,1\}^{2k}$, and $[w_{2i-1}]=[x_i]$ and $[w_{2i}]=[y_i]$ for $i\in\qqq{k}$. It therefore suffices to show that, for each fixed $\bsig \in \{0,1\}^{2k}$,
\begin{align}\label{eq:key_T_reudce_pf}
 \sum_{[\al]\in \bD_{\le \ell,\bsig}} \prod_{i=1}^k \br{\sT_{t}(|[x_i]-[\al]|\wedge \ell) \cdot \sT_{t}(|[y_i]-[\al]|\wedge \ell) } \prec (W^d\eta_t)^{-1} \Psi_t^{k-2}  \prod_{i=1}^k \sT_{t}(|[x_i]-[y_i]|\wedge \ell) .
\end{align} 
We now analyze three cases, depending on how many paths consist entirely of ``short edges'', i.e., edges of length $\le \ell$.

%We consider the following three cases for $\bD_{\le \ell,\bsig}$, depending on the number of paths consisting only of ``short edges'', i.e., edges of length $\le \ell$. 

\medskip
\noindent{\bf 1.} Assume that there are at least two indices $i$ such that $|[x_i]-[\al]|\vee|[y_i]-[\al]|\le \ell$. Without loss of generality, suppose this condition holds for $1\le i \le r$, and $|[x_i]-[\al]|\vee|[y_i]-[\al]|> \ell$ for $r+1\le i \le k$, where $2\le r\le k$. For $1\le i\le r$, we have  
\be\label{eq:TtTt}\sT_{t}(|[x_i]-[\al]|\wedge \ell)\,\sT_{t}(|[y_i]-[\al]|\wedge \ell)\lesssim \sT_{t}(|[x_i]-[y_i]|\wedge \ell) 
\cdot 
\frac{(W^{-d}B_{t,0})^{1/2}}{(|[x_i]-[\al]| \wedge |[y_i]-[\al]|+1)^{(d-2)/{2}}}, \ee
and whenever $|[y_i]-[\al]|>\ell$ (in particular for $r+1\le i \le k$, after relabelling $[x_i]$ and $[y_i]$ if necessary), we have
\be\label{eq:KtKt}\sT_{t}(|[x_i]-[\al]|\wedge \ell)\,\sT_{t}(|[y_i]-[\al]|\wedge \ell)\lesssim  \sT_{t}(\ell) \cdot (W^{-d}B_{t,0})^{1/2}\cdot \p{|[x_i]-[\al]|\wedge \ell+1}^{-\frac{d-2}{2}} . \ee
Both bounds follow from the definition of $\sT_t$, the monotonicity of $\sT_t$, and the two elementary facts
\(\sT_t(0)\le (W^{-d}B_{t,0})^{1/2}\) and \(\sT_t(r\wedge \ell)\le \sT_t(0)\,(r\wedge \ell+1)^{-(d-2)/2}\). In \eqref{eq:KtKt} we will bound the last factor by $1$ whenever it is not needed.
Using these two estimates, we can bound the LHS of \eqref{eq:key_T_reudce_pf} by
\begin{align*}
 &~(W^{-d}B_{t,0})^{k/2} \prod_{i=1}^{k} \sT_{t}(|[x_i]-[y_i]|\wedge \ell) \cdot \sum_{[\al]\in \bD_{\le \ell,\bsig}} \prod_{i=1}^r (|[x_i]-[\al]| \wedge |[y_i]-[\al]|+1)^{-\frac{d-2}{2}} \\
 \lesssim &~\ell^2 (W^{-d}B_{t,0})^{k/2} \prod_{i=1}^{k} \sT_{t}(|[x_i]-[y_i]|\wedge \ell) \prec (W^d\eta_t)^{-1} \Psi_t^{k-2}  \prod_{i=1}^{k} \sT_{t}(|[x_i]-[y_i]|\wedge \ell) ,
\end{align*}
where in the first step, we use
\[ \sum_{[\al]\in \bD_{\le \ell,\bsig}} \prod_{i=1}^r (|[x_i]-[\al]| \wedge |[y_i]-[\al]|+1)^{-\frac{d-2}{2}} \lesssim  \sum_{i=1}^r \sum_{[\al]\in \bD_{\le \ell}}(|[x_i]-[\al]| \wedge |[y_i]-[\al]|+1)^{-(d-2)}\lesssim \ell^2,  \]
and in the second step, we use the facts that $\ell\le (\log W)^{10}\ell_t$ and $\ell_t^2B_{t,0}\lesssim |1-t|^{-1} \lesssim\eta_t^{-1}$ when $1-t\ge \ilambda^2/L^{2}$.


\medskip
\noindent{\bf 2.} Suppose there is exactly one index $i$ such that $|[x_i]-[\al]|\vee|[y_i]-[\al]|\le \ell$.  
Without loss of generality, assume that $|[x_1]-[\al]|\vee|[y_1]-[\al]|\le \ell$, and $|[y_2]-[\al]|\ge \ell$. Then, applying \eqref{eq:TtTt} for $i=1$ and \eqref{eq:KtKt} for $2\le i \le k$---keeping the last factor of \eqref{eq:KtKt} for $i=2$ and bounding it by $1$ for $i\ge 3$---we can bound the LHS of \eqref{eq:key_T_reudce_pf} as
\begin{align*}
 &~(W^{-d}B_{t,0})^{k/2} \prod_{i=1}^{k} \sT_{t}(|[x_i]-[y_i]|\wedge \ell) \cdot \sum_{[\al]\in \bD_{\le \ell,\bsig}}  (|[x_1]-[\al]| \wedge |[y_1]-[\al]|+1)^{-\frac{d-2}{2}}(|[x_2]-[\al]| \wedge \ell +1)^{-\frac{d-2}{2}} \\
 \lesssim &~\ell^2 (W^{-d}B_{t,0})^{k/2} \prod_{i=1}^{k} \sT_{t}(|[x_i]-[y_i]|\wedge \ell) \prec (W^{d}\eta_t)^{-1} \Psi_t^{k-2} \prod_{i=1}^{k} \sT_{t}(|[x_i]-[y_i]|\wedge \ell) .
\end{align*}


\medskip
\noindent{\bf 3.} Finally, assume that for all $i\in\qqq{k},$ we have $|[x_i]-[\al]|\vee|[y_i]-[\al]|> \ell$. Without loss of generality, suppose $|[y_1]-[\al]|\ge \ell$ and $|[y_2]-[\al]|\ge \ell$. Then, applying \eqref{eq:KtKt} for $1\le i \le k$---keeping the last factor for $i\in\{1,2\}$ and bounding it by $1$ for $i\ge 3$---we can bound the LHS of \eqref{eq:key_T_reudce_pf} by 
\begin{align*}
 &~(W^{-d}B_{t,0})^{k/2} \prod_{i=1}^{k} \sT_{t}(|[x_i]-[y_i]|\wedge \ell) \cdot \sum_{[\al]\in \bD_{\le \ell,\bsig}}  (|[x_1]-[\al]| \wedge \ell+1)^{-\frac{d-2}{2}}(|[x_2]-[\al]| \wedge \ell +1)^{-\frac{d-2}{2}} \\
 \lesssim &~\ell^2 (W^{-d}B_{t,0})^{k/2} \prod_{i=1}^{k} \sT_{t}(|[x_i]-[y_i]|\wedge \ell) \prec (W^{d}\eta_t)^{-1} \Psi_t^{k-2} \prod_{i=1}^{k} \sT_{t}(|[x_i]-[y_i]|\wedge \ell) .
\end{align*}

\medskip

By combining all three cases, we conclude that \eqref{eq:key_T_reudce_pf} holds, which, in turn, implies \eqref{eq:key_T_reudce}.

\subsection{Basic properties of  \texorpdfstring{$\cal K$}{K}-loops}\label{Sec:CalK}

%In this section, we record some basic properties of the ${\cal K}$-loops that will be used in the analysis of the loop hierarchy, which finally lead to the proof of the $\cK$-loop bounds stated in \Cref{ML:Kbound}. The results in this section are higher-dimensional analogues (in dimensions $d\ge 3$) of the results in \cite[Section 3]{YY_25} and \cite[Section 4]{RBSO1D}. First, we define a dimension-independent \emph{tree representation formula} for the primitive loops as discovered in \cite{YY_25} for random band matrices and in \cite{RBSO1D} for the block Anderson model. The tree representation is built upon the concept of \emph{canonical partitions of polygons}. Roughly speaking, canonical partitions of an oriented polygon $\mathcal{P}_{\ba}$ are partitions in which each edge of the polygon is in one-to-one correspondence with each region in the partition. 

In this subsection, we collect several basic properties of the ${\cal K}$-loops used in the analysis of the loop hierarchy and apply them to establish the $\cK$-loop bounds stated in \Cref{ML:Kbound,lem_wardineq_K}. The results presented here are higher-dimensional analogues (in dimensions $d \ge 3$) of those in \cite[Section 3]{YY_25} and \cite[Section 4]{RBSO1D}. We begin by introducing a dimension-independent \emph{tree representation formula} for $\cal K$-loops, first discovered in \cite{YY_25} for random band matrices and later extended to the block Anderson model in \cite{RBSO1D}. This tree representation is constructed using the notion of \emph{canonical partitions of polygons}. Roughly speaking, a canonical partition of an oriented polygon $\mathcal{P}_{\ba}$ is a partition in which each edge of the polygon is in one-to-one correspondence with each region in the partition. 

\begin{definition}[Canonical partitions]\label{def:canpnical_part}
Fix $n\ge 3$ and let $\cal P_{\ba}$ be an oriented polygon with vertices $\ba=(a_1,a_2, \ldots ,a_n)$ arranged in a (counterclockwise) cyclic order,  where we adopt the cyclic convention that $a_i=a_j$ if and only if $i=j\mod n$. Let $(a_{k-1},a_k)$ denote the $k$-th side of $\cal P_{\ba}$. A planar partition of the polygonal domain enclosed by $\cal P_{\ba}$ is called {\bf canonical} if the following properties hold:
\begin{itemize}
\item Every sub-region in the partition is also a polygonal domain. 
\item There is a one-to-one correspondence between the sides of the polygon and the sub-regions, where every side $(a_{k-1},a_k)$ belongs to exactly one sub-region, denoted by $R_k$, and each sub-region contains exactly one side of $\cal P_{\ba}$. 

\item Every vertex $a_k$ of $\cal P_{\ba}$ belongs to exactly two regions, $R_k$ and $R_{k+1}$ (with the convention $R_{n+1}=R_1$). 
  
\end{itemize}
Note that given a canonical partition, by removing the $n$ sides of the polygon ${\cal P}_{\ba}$, the remaining interior edges form a tree, with the leaves being the vertices of ${\cal P}_{\ba}$. Following the definitions in \cite{YY_25}, we define the equivalence classes of all such trees under graph isomorphism, and denote the collection of equivalence classes by $\TSP({\cal P}_{\ba})$. % (where $\TSP$ refers to ``tree structure partition''). 
We will consider each element of $\TSP({\cal P}_{\ba})$ as an abstract tree structure rather than as an equivalence class, and call it a \emph{canonical tree partition}.
\end{definition}

In a canonical tree partition $\Gamma \in \TSP({\cal P}_{\ba})$, we call an edge that contains exactly one external vertex $a_k$ an \emph{external edge}, and an edge connecting two internal vertices an \emph{internal edge}. Two regions $R_k$ and $R_l$ are said to be \emph{neighbors} if they share a common side, which may be either an external or an internal edge. In the case of an external edge, we necessarily have $k-l = \pm 1 \pmod{n}$, and we refer to $R_k$ and $R_l$ as \emph{trivial neighbors}; otherwise, they are called \emph{nontrivial neighbors}.
Given $\bsig \in \{+,-\}^n$, we assign charges to the subregions as follows: each subregion $R_k$ carries the charge of the edge $(a_{k-1},a_k)$, which is given by $\sig_k$. An illustration is provided in the left panel of \Cref{example}, which shows a canonical tree partition $\Gamma \in \TSP({\cal P}_{\ba})$ of a polygon with six vertices, where $R_4$ and $R_6$ form a pair of nontrivial neighbors. We note that the figures in this section are all reproduced from \cite{RBSO1D}.


In the context of random band matrices, we assign a value to $\Gamma$ according to the following rule.
 
\begin{definition}%\label{M-graph-value-definition}
Given any $t\in [0,1)$ and $\bsig\in\{+,-\}^n$, we define the values of the edges in $\Gamma$ as follows:
  \begin{enumerate}
    \item If $e = \p{a_k, b}$ is an external edge lying between regions $R_k$ and $R_{k+1}$, then we define 
      \begin{equation}\label{f-external}
        f_{t,\bsig}\p{e} := \Theta^{(\sig_k,\sig_{k+1})}_{t}(a_k,b).
      \end{equation}
    \item If $e = \p{b_1, b_2}$ is an internal edge lying between regions $R_k$ and $R_{l}$, then 
      \begin{equation}\label{f-internal}
      \begin{aligned}
        f_{t,\bsig}\p{e}
          &:= \big(\Theta^{(\sig_k,\sig_l)}_{t}-I\big) \p{b_1, b_2} = m(\sig_k)m(\sig_{l})\cdot \big(tS^{(\sB)} \Theta^{(\sig_k,\sig_l)}_{t}\big)\p{b_1,b_2}.
          \end{aligned}
      \end{equation}
  \end{enumerate}
%Above, we abbreviate that $m_k\equiv m(\sig_k)$ for $k\in \qqq{n}$. 
Then, we assign a value $\Gamma^{(n)}_{t,\bsig,\ba}$ to  $\Gamma$ as:
\begin{equation}\label{M-graph-value-unsummed}
\Gamma^{(n)}_{t,\bsig,\ba} := \p{\prod_{i=1}^n m(\sig_i)} \cdot \sum_{\mathbf b} \prod_{e} f_{t,\bsig}\p{e}    \, ,
\end{equation}
where $\mathbf b=(b_1,\ldots,b_{r})$ denotes the internal vertices in $\Gamma$ and $e$ denotes all the edges in $\Gamma$. 
\end{definition}




With these definitions, we recall the tree representation formula of the $\mathcal{K}$-loops in Lemma 3.4 of \cite{YY_25}. Recall that the formulas for the $\cal K$-loops of length 2 and 3 have been given in \eqref{Kn2sol} and \eqref{Kn3sol}.

\begin{lemma}[Lemma 3.4 of \cite{YY_25}]\label{tree-representation}
In the setting of random band matrices, for any $n\ge 4$, $t\in[0,1)$, $\bsig\in \{+,-\}^n$, and $\ba\in (\Zn)^n$, we have the following representation formula for $\cal K$-loops:
  \begin{equation}\label{eq_Ktree}
\cK_{t,\bsig,\ba}^{(n)}
      =W^{-d(n-1)} \sum_{\Gamma \in \TSP\p{\mathcal{P}_{\ba}}} \Gamma^{(n)}_{t,\bsig,\ba}.
  \end{equation}
\end{lemma}

When extending \eqref{eq_Ktree} to the block Anderson model, certain factors of $m$ must be replaced by entries of the matrix $M$. A canonical procedure for this replacement is described in \cite[Section 4]{RBSO1D}. To formulate it, we first extend \Cref{def:canpnical_part} to include loops that contain \emph{$M$-edges}. As the name suggests, these edges correspond to entries of $M$, while the remaining edges in our graphs are left \emph{unlabeled} (i.e., without label $M$): external edges represent entries of $\Theta_t$, and internal edges represent entries of $tS^{(\sB)}\Theta_t$.



\begin{definition}[Canonical partitions with $M$-loops]\label{m-loop-tsp}
Let $\Gamma \in \TSP\p{\mathcal{P}_{\ba}}$ be a canonical tree partition of the oriented polygon $\mathcal{P}_{\ba}$, and denote the internal vertices of $\Gamma$ by $\mathbf b = \p{b_1, \ldots, b_r}$. We define the graph {$\Gamma_{M}$} by replacing each $b_i$ with an $M$-loop in the following way.
\begin{enumerate}
\item Consider the subgraph $(V^{\p{b_i}}, E^{\p{b_i}})$ of all vertices in $\Gamma$ connected to $b_i$. More precisely, we let
      \begin{equation}\label{b_i-subgraph}
        V^{\p{b_i}} := \{b_i,c_1, \ldots, c_{k_i}\},
        \quad \text{and}\quad E^{\p{b_i}} = \{\p{c_1, b_i}, \ldots, \p{c_{k_i}, b_i}\}
      \end{equation}
      denote the subsets of all vertices (including $b_i$) and edges connected to $b_i$. Reordering the $c_k$'s if necessary, we can ensure that $\p{c_1, \ldots, c_{k_i}}$ form a loop without any crossing edges and with vertices arranged in counterclockwise order.
      
\item Next, we construct a new graph $(\widetilde{V}^{\p{b_i}}, \widetilde{E}^{\p{b_i}})$ with vertices and edges
      \begin{align}\label{eq:M-loopexample}
   \qquad      \widetilde{V}^{\p{b_i}} &= \{b_{i,1}, \ldots, b_{i,k_i}, c_1, \ldots, c_{k_i}\},\ \
           \widetilde{E}^{\p{b_i}}
     =\{\p{c_j, b_{i,j}} :  j \in \qqq{k_i}\} \cup \{\p{b_{i,j}, b_{i,j+1}; M}: j \in \qqq{k_i}\}, 
      \end{align}
   where $e = \p{e_i, e_f; M}$ refers to an edge labeled with $M$, and we adopt the cyclic convention that $b_{i,k_i+1}=b_{i,1}$. We will call $\p{e_i, e_f; M}$ an \emph{$M$-edge}. %Relabeling the $b_{i,j}$'s if needed, we ensure that $\p{b_{i,1} \to b_{i,2} \to \cdots \to b_{i,k_i}}$ are also arranged in counterclockwise order.  
      
\item Lastly, we replace the subgraph $(V^{\p{b_i}}, E^{\p{b_i}})$ with $(\widetilde{V}^{\p{b_i}}, \widetilde{E}^{\p{b_i}})$.
\end{enumerate}

  \begin{figure}[h]
    \centering
     \scalebox{0.9}{
    \begin{tikzpicture}
      \coordinate (bi) at (0, 0);
      
      \fill (bi) circle (1pt) node[left=4pt]{$b_i$};
      \foreach \i/\t in {1/72, 2/144, 3/216, 4/288, 5/360} {
        \fill (\t:2) circle (1pt);
        \draw (bi) -- (\t:2);
        \node at (\t+72:2.4) {$c_{\i}$};
      }
    \end{tikzpicture}
    $\quad$
    \begin{tikzpicture}
      \node at (-3, 0) {$\to$};
      \node at (0, 0) {$M$};
      \foreach \i/\t in {1/72, 2/144, 3/216, 4/288, 5/360} {
        \fill (\t:2) circle (1pt);
        \fill (\t:1) circle (1pt);
        \draw (\t:1) -- (\t:2);
        \draw (\t:1) -- (\t+72:1);
        \node at (\t+72:2.4) {$c_{\i}$};
        \node at (\t+92:1.3) {$b_{i,\i}$};
      }
    \end{tikzpicture}}
    \caption{Replacing $b_i$ with an $M$-loop with 5 sides.}\label{Fig:Mloop}
  \end{figure}

  In \Cref{Fig:Mloop}, we illustrate the above procedure for an example with $k_i=5$. Repeating these steps for each internal vertex $b_i$, $i\in\qqq{r}$, we get a graph \smash{$\Gamma_{M}$}, which we will refer to as the \emph{$M$-graph corresponding to $\Gamma$}. Note that the order in which we replace $b_i$'s by $M$-loops does not matter, and every boundary edge $\p{a_{k-1}, a_k}$ still belongs to exactly one polygonal region in the $M$-graph \smash{$\Gamma_{M}$}. 
With a slight abuse of notation, we still use $R_k$ to denote the sub-region containing $\p{a_{k-1}, a_k}$ in \smash{$\Gamma_{M}$}, and assign the charge $\sig_k$ of $\p{a_{k-1}, a_k}$ to $R_k$. 
% see \Cref{Fig:Mloop2} for an example.
% \begin{figure}[h]\label{m-graph-region}
%     \centering
%     \begin{tikzpicture}
%       \fill (-1, 0) circle (1pt);
%       \fill (1, 0) circle (1pt);
%       \fill (0, 2) circle (1pt);
      
%       \draw (-1, 0) node[below]{$a_{k-1}$}
%         -- (1, 0) node[below]{$a_k$}
%         -- (0, 2) node[above]{$b_i$}
%         -- cycle;
%       \node at (0, 0.75) {$R_k$};
%       \node at (0, -1) {$\Gamma_{\ba}$};
%     \end{tikzpicture}
%     \hspace{1in}
%     \begin{tikzpicture}
%       \fill (-1, 0) circle (1pt);
%       \fill (1, 0) circle (1pt);
%       \fill (0.5, 2) circle (1pt);
%       \fill (-0.5, 2) circle (1pt);
      
%       \draw (-1, 0) node[below]{$a_{k-1}$}
%         -- (1, 0) node[below]{$a_k$}
%         -- (0.5, 2)  node[above right]{$b_{i,j}$}
%         -- (-0.5, 2) node[above left]{$b_{i,j-1}$}
%                      node[pos=0.5, above]{$M$}
%         -- cycle;
%       \node at (0, 0.75) {$R_k$};
%       \node at (0, -1) {$\Gamma^{\p{M}}_{\ba}$};
%     \end{tikzpicture}
%     \caption{$R_k$ refers to the region containing $\protect\p{a_{k-1}, a_k}$ in both $\Gamma_{\ba}$ and its corresponding $M$-graph \smash{$\Gamma^{\protect\p{M}}_{\ba}$}. %Moreover, this also illustrates that the $M$-edge $\protect\p{b_{i,j-1}, b_{i,j}}$ belongs to the region shared by $\protect\p{a_{k-1}, b_i}$ and $\protect\p{a_k, b_i}$.
%     }\label{Fig:Mloop2}
%   \end{figure}
\end{definition}

We refer readers to \Cref{example} for an example of a canonical tree partition $\Gamma$ and its $M$-graph $\Gamma_{M}$. 
  \begin{figure}[h]
    \centering
     \scalebox{0.9}{
    \begin{tikzpicture}
      \coordinate (b1) at (-1, 0);
      \coordinate (b2) at (1, 0);
      
      \fill (b1) circle (1pt);
      \fill (b2) circle (1pt);

      \foreach \i/\t in {1/60, 2/120, 3/180, 4/240, 5/300, 6/360} {
        \coordinate (a\i) at (\t+90:2.5);
        \fill (a\i) circle (1pt);
        \draw (a\i) [dashed]-- (\t+150:2.5);
        \node at (\t+90:2.85) {$a_{\i}$};
        \node at (\t+60:2.5) {$\sigma_{\i}$};
      }
      
      \draw (a6) -- (b1);
      \draw (a1) -- (b1);
      \draw (a2) -- (b1);
      \draw (a3) -- (b1);
      
      \draw (a4) -- (b2);
      \draw (a5) -- (b2);
      
      \draw (b1) -- (b2);

      \node at (0, -3.5) {$\Gamma$};
    \end{tikzpicture}
    \hspace{1in}
    \begin{tikzpicture}
      \foreach \i/\t in {1/60, 2/120, 3/180, 4/240, 5/300, 6/360} {
        \coordinate (a\i) at (\t+90:2.5);
        \fill (a\i) circle (1pt);
        \draw (a\i) [dashed]-- (\t+150:2.5);
        \node at (\t+90:2.85) {$a_{\i}$};
        \node at (\t+60:2.5) {$\sigma_{\i}$};
      }
      
      \coordinate (c11) at ($(b1)!0.45!(a6)$);
      \coordinate (c12) at ($(b1)!0.45!(a1)$);
      \coordinate (c13) at ($(b1)!0.45!(a2)$);
      \coordinate (c14) at ($(b1)!0.45!(a3)$);
      \coordinate (c15) at ($(b1)!0.45!(b2)$);
      
      \coordinate (c21) at ($(b2)!0.25!(b1)$);
      \coordinate (c22) at ($(b2)!0.45!(a4)$);
      \coordinate (c23) at ($(b2)!0.45!(a5)$);
      
      \foreach \i/\nsides in {1/5, 2/3} {
        \foreach \j in {1, 2, ..., \nsides} {
          \fill (c\i\j) circle (1pt);
        }
      }
      
      \foreach \i/\nsides in {1/5, 2/3} {
        \draw (c\i1)
        \foreach \j in {1, ..., \nsides} {
           -- (c\i\j)
        }
          -- cycle;
      }
      
      \draw (c11) -- (a6);
      \draw (c12) -- (a1);
      \draw (c13) -- (a2);
      \draw (c14) -- (a3);
      
      \draw (c21) -- (c15);
      
      \draw (c22) -- (a4);
      \draw (c23) -- (a5);
      
      \node at ($0.2*(c11) + 0.2*(c12) + 0.2*(c13) + 0.2*(c14) + 0.2*(c15)$) {$M$};
      \node at ($0.333*(c21) + 0.333*(c22) + 0.333*(c23)$) {$M$};
      
      \node at (0, -3.5) {$\Gamma_{{M}}$};
    \end{tikzpicture}}
    \caption{Example of $\Gamma \in \TSP\protect\p{\mathcal{P}_{\ba}}$ and its corresponding $M$-graph $\Gamma_{M}$.}\label{example}
  \end{figure}


\begin{definition}\label{M-graph-value-definition}
Given an $M$-graph \smash{$\Gamma_{M}$}, we assign a value to $\Gamma$ according to the following rule. For any $t\in[0,1)$ and $\bsig\in\{+,-\}^n$, we define the values of the edges in $\Gamma_{M}$ as follows. 
  \begin{enumerate}
    \item If $e = \p{a_k, b}$ is an external edge lying between regions $R_k$ and $R_{k+1}$, then we define
      \begin{equation}\label{f-external2}
        g_{t,\bsig}\p{e} := \Theta^{\p{\sig_{k}, \sig_{k+1}}}_{t}(a_k,b) \, .
      \end{equation}

    \item If $e = \p{b_1, b_2}$ is an internal unlabeled edge lying between regions $R_k$ and $R_{l}$, then we define
      \begin{equation}\label{f-internal2}
        g_{t,\bsig}\p{e}
          := \big(tS^{(\sB)}\Theta^{(\sig_k,\sig_{l})}_{t}\big)\p{b_1,b_2} .
      \end{equation}

    \item Corresponding to each loop of $M$-edges as in \eqref{eq:M-loopexample}, we define its value as 
  \begin{equation}\label{eq:Mloop_defgen}
    \cal F\p{b_i} \equiv W^{(k_i-1)d} \cal M_{\bsig(b_i),\ba(b_i)}^{(k_i)}, %=W^{-d} \avg{\prod_{j=1}^{k_i} \left[M\p{\sigma_j{(b_i)}}\id_{b_{i,j}}\right]},
  \end{equation}
 where we recall the $M$-loop defined in \eqref{eq:KMloop}. %and $\id_a:=W^dE_a$ is the block identity with.  
 Here, the charges $\sig_j(b_i)$ and vertices $a_j(b_i)$ are fixed according to the following rule: for each $j \in \qqq{k_i}$, the edge $\p{b_{i,j-1} \to b_{i,j}}$ (with the cyclic convention that $b_{i,0}=b_{i,k_i}$) belongs to exactly one region $R_{k_j}$. Then, we set ${\sigma}_j{\p{b_i}} = {\sigma}_{k_j}$ and $a_j{\p{b_i}} = b_{i,j}$. 
 
%  In simpler terms, $\cal F\p{b_i}$ is an $M$-loop that shares the same orientation as  $\mathcal{P}_{\ba}$, with each edge having the same color as the region it lies in. See \Cref{Fig2:charge-example} for an example of how to assign charges and directions to the $M$-loops.
    
    % If $e = \p{[b_1], [b_2]}$ is an internal $M$-edge belonging to the region $R_k$, then %we define
    %   \begin{equation}\label{f-internalM}
    %     g_{\bsig,z}\p{e}
    %       := m_t(z,\sig_k)\delta_{[b_1][b_2]}\, .
    %   \end{equation}
      
  \end{enumerate}
Now, we assign a value $\Gamma^{(n)}_{M;t,\bsig,\ba}$ to $\Gamma_M$ as:
\begin{equation}\label{M-graph-value-unsummed2}
\Gamma^{(n)}_{M;t,\bsig,\ba} := \sum_{b_{i,j} \in \Zn,\; i\in\qqq{r},\; j\in\qqq{k_i}}  \prod_{\text{unlabeled }e} g_{t,\bsig}(e) \cdot \prod_{i=1}^r \cal F(b_i) \,  ,
\end{equation}
where $b_{i,j}$'s denote the internal vertices in $\Gamma_M$ and $e$ denotes all unlabeled (i.e., non-$M$) edges.  
\end{definition}

It is straightforward to verify that the definition in \eqref{M-graph-value-unsummed2} is consistent with \eqref{M-graph-value-unsummed} in the case of random band matrices, where $M(\sig) = m(\sig) I$. Furthermore, in the block Anderson model, the tree representation formula \eqref{eq_Ktree} extends as follows.
 

\begin{lemma}[Lemma 4.16 of \cite{RBSO1D}]\label{tree-representation_BA}
In the setting of the block Anderson model, for any $n\ge 4$, $t\in[0,1)$, $\bsig\in \{+,-\}^n$, and $\ba\in (\Zn)^n$, we have the following representation formula for $\cal K$-loops:
\begin{equation}\label{eq:tree_rep2}
\cK_{t,\bsig,\ba}^{(n)}
=W^{-d(n-1)} \sum_{\Gamma \in \TSP\p{\mathcal{P}_{\ba}}} \Gamma^{(n)}_{M;t,\bsig,\ba}.
  \end{equation}
\end{lemma}

 

Next, we recall the molecule structure of $\mathcal{K}^{(n)}_{t,\bsig, \mathbf{a}}$ given in Section 3.3 of \cite{YY_25}. We split the $M$-graphs $\Gamma_M$ with $\Gamma\in \TSP(\mathcal{P}_{\mathbf{a}})$ according to which edges of the tree are ``long''---we refer to a \smash{$\Theta_t^{(\sig_1,\sig_2)}$}-edge as a \emph{long edge} if $\sig_1\ne \sig_2$, and a \emph{short edge} otherwise. We adopt this terminology because, according to \eqref{prop:ThfadC} and \eqref{prop:ThfadC_short}, the short edge decays on a scale of order 1, which is much shorter than the decay scale $\ell_t$ for a long edge. Given any $\bsig \in \{+, -\}^{n}$, define the subset of long internal edges (i.e., the boundary between two nontrivial neighbors of different charges) as 
\[
{\cal F}_{\mathrm{long}}(\Gamma, \bsig) := \left\{\{k,l\} \in \Z^{\mathrm{off}}_n: R_k\cap R_l\neq \emptyset ,\ \sigma_k\ne \sigma_l \right\},
\]
where we define the subset 
\[
\Z^{\mathrm{off}}_n :=\left\{\{k, l\} : 1 \leq k < l \leq n,\, k - l \pmod n \notin \{1,-1\}\right\}.
\]
Given any subset \(\pi \subset \mathbb{Z}_n^{\mathrm{off}}\), we use \(\TSP(\mathcal{P}_{\ba}, \bsig, \pi) := \big\{\Gamma \in \TSP(\mathcal{P}_{\ba}) : {\cal F}_{\text{long}}(\Gamma, \bsig) = \pi\big\}\) to denote the set of all $\Gamma$ such that \(\pi\) labels the pairs of all non-trivial neighbors in $\Gamma$ (note $\pi$ can be $\emptyset$). Then, we define 
  \begin{align}\label{eq:defKpi}
    \cK^{\p{\pi}}\p{t,\bsig,\ba}
      :=\sum_{\Gamma \in \TSP\p{\mathcal{P}_{\ba}, \bsig, \pi}} \Gamma^{\p{n}}_{M;t,\bsig,\ba}  .
  \end{align}
% (meaning that they are associated to $\Theta^{(B)}_{t|m|^2}$ as opposed to $\Theta^{(B)}_{tm^2}$ or $\Theta^{(B)}_{t\overline{m}^2}$). More precisely,
% \begin{equation*}
%     TSP(\mathcal{P}_{\mathbf{a}}) = \bigsqcup_{\pi\subset\Z^{off}_n} TSP(\mathcal{P}_{\mathbf{a}}, \bsig, \pi),
% \end{equation*}
% where $\Z^{off}_n$ is the set of pares of non-consecutive edges of $\mathcal{P}_{\mathbf{a}}$. See Definition 3.8 of \cite{YY_25} for the exact details. The key observation needed in this section is that the subsets $\pi\subset\Z^{off}_n$ encode the different ``molecule structures'' of trees $\Gamma_{\mathbf{a}}\in TSP(\mathcal{P}_{\mathbf{a}})$. 
Note that $\mathcal{K}^{(n)}_{t,\bsig, \mathbf{a}}$ can be decomposed as 
\begin{equation}\label{eq_K-Kpi}
    \mathcal{K}^{(n)}_{t, \bsig, \mathbf{a}} = W^{-d(n-1)}\sum_{\pi\subset \mathbb{Z}_n^{\rm{off}}} \mathcal{K}^{(\pi)}(t, \bsig, \mathbf{a}).
\end{equation}
Moreover, we have the following molecule decomposition of $\mathcal{K}^{(\pi)}$ given in equation (3.53) of \cite{YY_25}:
\begin{equation}\label{eq:molecule-Kpi}
    \mathcal{K}^{(\pi)}(t, \bsig, \mathbf{a}) = \sum_{\mathbf{b},\mathbf{c}} \prod_{k=1}^n \Theta^{(\sig_k,\sig_{k+1})}_{t,a_k b_k} \cdot \prod_{k=1}^{r-1} \left(tS^{(\sB)} \Theta_t^{(+,-)}\right)_{c_{2k-1} c_{2k}} \cdot \prod_{k=1}^r \Sigma^{(\pi)}(t, \bsig^{(k)}, \mathbf{b}^{(k)}).
\end{equation}
Here, each term $\Sigma^{(\pi)}(t, \bsig^{(k)}, \mathbf{b}^{(k)})$ represents a \emph{molecule}, defined as a maximal subgraph consisting solely of $M$-loops and \emph{short (unlabeled) internal edges}.\footnote{
With a slight abuse of notation, we again use the term ``molecule'' here. This definition is in the same spirit as that in \Cref{def_poly}, although the graph considered here is different from the one in \Cref{def_graph1}.
} 
In other words, if each molecule is collapsed into a single vertex, the resulting quotient graph contains only long internal edges connecting different molecules, along with external edges, which may be either short or long. 
In \eqref{eq:molecule-Kpi}, we assume there are $r$ molecules in total. The terms \smash{$tS^{(\sB)} \Theta_t^{(+,-)}$} correspond to the long internal edges $(c_{2k-1},c_{2k})$ between molecules, while the terms \smash{$\Theta^{(\sig_k,\sig_{k+1})}_{t,a_k b_k}$} correspond to the external edges $(a_k,b_k)$. The vectors $\mathbf b=(b_1,\ldots, b_n)$ and $\mathbf c=(c_1, c_2,\ldots, c_{2r-3},c_{2r-2})$ denote the internal vertices that serve as endpoints of these external and long internal edges, respectively. Furthermore, $\bsig^{(k)}$ and $ \mathbf{b}^{(k)}$ specify the charges and distinguished vertices associated with the $k$-th molecule. All other internal vertices within the molecules are implicitly summed over.

%Moreover, $\bsig^{(k)}$ and $ \mathbf{b}^{(k)}$ denote the charges and vertices associated with the $k$-th molecule, while all the other internal vertices in the molecules have been summed over. 


%The terms $\Sigma^{(\emptyset)}(t, \bsig^{(k)}, \mathbf{b}^{(k)})$ are ``molecule weights'', which consist of products of $\Theta^{(B)}_{tm^2}$ and $\Theta^{(B)}_{t\overline{m}^2}$ and, thus, exhibit exponential decay when $\max_{i,j} |b^{(k)}_i - b^{(k)}_j| \gg 1$. \medskip


Using the estimate \eqref{prop:ThfadC_short}, we can derive the following exponential decay estimate on ``pure loops'' where all charges are identical.


\begin{lemma}\label{lem_pureloop} 
For $\bsig \in \{+, -\}^n$ with $\sigma_1 = \sigma_2 = \cdots = \sigma_n$, there exist constants $c_n,C_n>0$ such that 
\begin{equation}\label{res_pureKes}
\left|\mathcal{K}^{(n)}_{t, \bsig, \mathbf{a}}\right|\le C_n W^{-d(n-1)} \exp\left(-c_n\max_{i,j}|a_i-a_j|\right).
\end{equation} 
\end{lemma}
\begin{proof}
% We follow the proof of Corollary 3.5 in \cite{YY_25}. We use the tree representation in \eqref{eq_Ktree} and the decomposition of the tree terms into sums from (3.6) of \cite{YY_25}:
% \begin{equation*}
%     \Gamma_{\mathbf{a}}(t,\bsig)=\sum_{\mathbf{b}\in(\Z_{L}^{2})^{n}}\Gamma^{(\mathbf{b})}_{\mathbf{a}}(t,\bsig),
% \end{equation*}
In the tree representation \eqref{eq:tree_rep2}, each term $\Gamma^{(n)}_{M;t,\bsig,\ba}$ on the RHS contains only $M$-edges and short unlabeled edges. Using the exponential decay from \eqref{prop:ThfadC_short}, along with the bound \eqref{Mbound_AO} or \eqref{Mbound_AO2} for $M$-edges in the setting of the block Anderson model, we obtain the desired estimate. 
% \begin{equation*}
%     |\Gamma^{(n)}_{t,\bsig,\ba}|\leq C_{n}\exp\left(-c_{n}\max_{i,j}|b_{i}-b_{j}|-c_{n}\max_{i,j}|a_{i}-b_{j}|\right).
% \end{equation*}
% Plugging this into the tree representation of $\mathcal{K}$ in \eqref{eq_Ktree}, completes the proof.
\end{proof}

We are now ready to prove the key bounds on $\cK$-loops, stated in \Cref{ML:Kbound}. The proof is analogous to that of Lemma 3.11 in \cite{YY_25}, but requires additional modifications to handle the higher-dimensional setting $d \ge 3$. For the reader's convenience, we provide the proof below.



% \begin{lemma}\label{ML:Kbound+pi}
% For any $\mathbf{a} \in \left(\Z_L^d\right)^n$ and $\bsig \in \left\{+, -\right\}^n$, we have the estimate
% \begin{equation}\label{eq:bcal_k_2}
%   \left|{\cal K}^{(n)}_{t, \bsig, \mathbf{a}}\right| \prec \left(W^{-d}B_{t,0}\right)^{n-1}.    
% \end{equation}
% \end{lemma}

\begin{proof}[\bf Proof of \Cref{ML:Kbound}]
For $n=1$ we have $\cK^{(1)}_{t,\bsig,\ba}=m(\sig_1)=\OO(1)$, and for $n\in\{2,3\}$ the bound follows from the explicit formulas \eqref{Kn2sol} and \eqref{Kn3sol} together with \eqref{prop:ThfadC} and, for $n=3$, \eqref{prop:ThfadC_short} (and \eqref{Mbound_AO} or \eqref{Mbound_AO2} for the $M$-loop in the block Anderson model). Indeed, the three charge pairs $(\sig_1,\sig_2)$, $(\sig_2,\sig_3)$, and $(\sig_3,\sig_1)$ in \eqref{Kn3sol} cannot all consist of different charges, so at least one of the three $\Theta$-propagators in \eqref{Kn3sol} is short-range by \eqref{prop:ThfadC_short}, and its sum over the free index is $\OO(1)$ rather than $(1-t)^{-1}$; the other two propagators are bounded pointwise by $\OO(B_{t,0})$ using \eqref{prop:ThfadC}. This gives $|\cK^{(3)}_{t,\bsig,\ba}|\lesssim (W^{-d}B_{t,0})^2$. We may therefore assume $n\ge 4$, which is the range in which the tree representation \eqref{eq_Ktree} is available.
By \eqref{eq_K-Kpi}, it suffices to show that for any $\pi \subset \Z^{\rm{off}}_n$, we have
\begin{equation}\label{eq:K-pi-bound}
    \left|\mathcal{K}^{(\pi)}(t, \bsig, \mathbf{a})\right| \prec B_{t,0}^{n-1}.
\end{equation}
We prove the above bound by induction on the number of molecules $r$ in $\pi$. In order to carry out the induction step, we will prove \eqref{eq:K-pi-bound} for a slightly generalized quantity $\tilde{\mathcal{K}}^{(\pi)}$, defined by
\begin{equation}\label{eq:wtKpi}
    \tilde{\mathcal{K}}^{(\pi)}(t, \bsig, \mathbf{a}) = \sum_{\mathbf{b},\mathbf{c}} \prod_{k=1}^n \tilde{\Theta}_{t,a_k b_k} \cdot \prod_{k=1}^{r-1} \left(tS^{(\sB)}\Theta^{(+,-)}_{t}\right)_{c_{2k-1} c_{2k}} \cdot \prod_{k=1}^r \Sigma^{(\pi)}(t, \bsig^{(k)}, \mathbf{b}^{(k)}),
\end{equation}
where each $\tilde{\Theta}_{t}$ represents either $\Theta^{(\sig,\sig')}_{t}$ for some $\sig,\sig'\in\{+,-\}$ or $tS^{(\sB)}\Theta^{(+,-)}_{t}$.  

We start with the single molecule case with $\pi=\emptyset$. In this case, we can write 
\begin{equation*}
    \tilde{\mathcal{K}}^{(\emptyset)}(t, \bsig, \mathbf{a}) = \sum_{\mathbf{b}} \prod_{i=1}^n \tilde{\Theta}_{t,a_i b_i} \cdot \Sigma^{(\emptyset)}(t, \bsig, \mathbf{b}).
\end{equation*}
First, the $M$-edges and short unlabeled edges in $\tilde{\mathcal{K}}$, including all edges contained within the molecule,
% have weights \smash{$\Theta^{(\sB)}_{tm^2}$}, \smash{$\Theta^{(\sB)}_{tm^2} - I$}, \smash{$\Theta^{(\sB)}_{t\bar{m}^2}$}, or \smash{$\Theta^{(\sB)}_{t\bar{m}^2} - I$}, which 
have constant size and fast exponential decay outside their constant range by \eqref{prop:ThfadC_short}, together with \eqref{Mbound_AO} or \eqref{Mbound_AO2} in the block Anderson model. Thus, the molecule $\Sigma^{(\emptyset)}$ has constant size and fast exponential decay as well:
\begin{equation}\label{eq:molecule-decay}
    \left|\Sigma^{(\emptyset)}(t, \bsig, \mathbf{b})\right| \le C\exp\left\{-c\max_{i,j} |b_i-b_j|\right\}
\end{equation}
for some constants $c,C>0$. Additionally, if we have a pure loop (i.e., $\bsig$ consists entirely of the same charge), then every edge in $\tilde{\mathcal{K}}^{(\emptyset)}$ is short, so we have 
\begin{equation*}
    \left|\tilde{\mathcal{K}}^{(\emptyset)}(t, \bsig, \mathbf{a})\right| \le C\exp\left\{-c\max_{i,j} |a_i-a_j|\right\} \lesssim B_{t,0}^{n-1}.
\end{equation*}
If $\bsig$ is not a pure loop, then without loss of generality, we can assume that $\sigma_1 \neq \sigma_2$. In this case, we have 
\begin{equation*}
    \tilde{\mathcal{K}}^{(\emptyset)}(t, \bsig, \mathbf{a}) = \sum_{b_1} \tilde{\Theta}_{t,a_1b_1} \sum_{\mathbf{b}\setminus \{b_1\}} \Sigma^{(\emptyset)}(t, \bsig, \mathbf{b}) \prod_{i=2}^n  \tilde{\Theta}_{t,a_i b_i}.
\end{equation*}
We now claim that 
\begin{equation}\label{eq:ind-step-bound}
    \sum_{b_1} \bigg|\sum_{\mathbf{b}\setminus \{b_1\}} \Sigma^{(\emptyset)}(t, \bsig, \mathbf{b}) \prod_{i=2}^n \tilde{\Theta}_{t,a_i b_i}\bigg| \prec B_{t,0}^{n-2}.
\end{equation}
Combining this estimate with the bound \eqref{prop:ThfadC} for the $\Theta$-propagators, we obtain that 
\begin{align*}
    \left|\tilde{\mathcal{K}}^{(\emptyset)}(t, \bsig, \mathbf{a})\right| %&\le \sum_{b_1} \left(\tilde{\Theta}^{(sB)}_{t}\right)_{a_1b_1} \bigg|\sum_{\mathbf{b}\setminus \{b_1\}} \Sigma^{(\emptyset)}(t, \bsig, \mathbf{b}) \prod_{i=2}^n \left(\tilde{\Theta}^{(\sB)}_{tm_i m_{i+1}}\right)_{a_i b_i}\bigg| \\
    &\lesssim B_{t,0}\sum_{b_1} \bigg|\sum_{\mathbf{b}\setminus \{b_1\}} \Sigma^{(\emptyset)}(t, \bsig, \mathbf{b}) \prod_{i=2}^n \tilde{\Theta}_{t,a_i b_i}\bigg| \prec B_{t,0}^{n-1}.
\end{align*}
This concludes \eqref{eq:K-pi-bound} for the case $\pi=\emptyset$.


To prove the bound \eqref{eq:ind-step-bound}, we consider two cases depending on whether (i) $\bsig$ is non-alternating, or (ii) $\bsig$ is alternating, where $\sig_j\ne \sig_{j+1}$ for all $j\in \qqq{n}$.  

\medskip
{\noindent\bf (i)} If $\bsig$ is non-alternating, then there exists $j\in\qqq{2,n}$ such that $\sigma_j = \sigma_{j+1}$. The corresponding short edge has fast exponential decay on the constant scale by \eqref{prop:ThfadC_short}: % and \eqref{Mbound_AO} or \eqref{Mbound_AO2}:
\begin{equation}\label{eq:shortexternal}
    \left|\tilde{\Theta}_{t,a_jb_j}\right| \le C\exp\{-c|a_j - b_j|\}.
\end{equation}
For the remaining $(n-2)$ factors \smash{$\tilde{\Theta}_{t,a_i b_i}$}, we can bound them pointwise using \eqref{prop:ThfadC} as
\be\label{eq:pointwise_Theta}
    \left|\tilde{\Theta}_{t,a_ib_i}\right| \lesssim B_{t,0}, \quad \forall i \in\qqq{2,n} \setminus \{j\}.
\ee
The above two bounds, together with the exponential decay for the molecule weight in \eqref{eq:molecule-decay}, allow us to bound the summation on the LHS of \eqref{eq:ind-step-bound} by $B_{t,0}^{n-2}$. % which gives \eqref{eq:ind-step-bound} in the non-alternating case.

 

\medskip
{\noindent\bf (ii)} It remains to handle the challenging case of alternating $\bsig$. In this setting, in addition to the constant-range exponential decay, we must also use the sum-zero property of molecule weights:
\begin{equation}\label{eq:Sigma-empty-sum-zero}
    \sum_{\mathbf{b}\setminus \{b_1\}} \Sigma^{(\emptyset)}(t, \bsig, \mathbf{b}) = \OO(|1-t|),\quad \sum_{\mathbf{b}\setminus \{b_1\}} \big|\Sigma^{(\emptyset)}(t, \bsig, \mathbf{b})\big| = \OO(\ilambda^2+|1-t|).
\end{equation}
The first estimate was established in \cite[Lemma 3.10]{YY_25} for 1D random band matrices and in \cite[Lemma 4.29]{RBSO1D} for 1D and 2D block Anderson models, while the second estimate was proved in \cite[Claim 4.30]{RBSO1D}. Importantly, these proofs are dimension-independent: they rely only on the pure loop estimate (\Cref{lem_pureloop}), the short-edge bound \eqref{prop:ThfadC_short}, the $M$-edge bound \eqref{Mbound_AO} or \eqref{Mbound_AO2}, and Ward's identity \eqref{WI_calK} for $\cal K$-loops.


We now decompose the long external edges into three parts as follows:
\begin{align*}
\tilde{\Theta}^{(+,-)}_{t,a_jb_j} = f_0(a_j, s_j) + f_1(a_j, s_j) + f_2(a_j, s_j), \quad \tilde{\Theta}_{t}^{(+,-)}\in \big\{\Theta^{(+,-)}_{t},\ tS^{(\sB)}\Theta^{(+,-)}_{t}\big\}  ,
\end{align*}
where we denote $f(a_j, s_j)= \tilde{\Theta}^{(+,-)}_{t}\p{a_j,b_1+s_j}$ with $s_j = b_j - b_1$, and 
\begin{align*}
    %f(a_j, s_j) &= \big(\tilde{\Theta}^{(\sB)}_{t}\big)_{a_j,b_1+s_j},\\
    f_0(a_j, s_j) &= f(a_j, 0),\quad
    f_1(a_j, s_j) = \frac12 f(a_j, s_j) - \frac12 f(a_j, -s_j),\\
    f_2(a_j, s_j) &= \frac12 f(a_j, s_j) + \frac12 f(a_j, -s_j) - f(a_j, 0).
\end{align*}
By \Cref{lem_propTH}, we have the following bounds under the condition $|s_j|\prec 1$ for $j\in\qqq{2,n}$ (if $|a_j-b_1|\ge 4|s_j|$, the bounds on $f_1$ and $f_2$ follow from \eqref{prop:BD1} and \eqref{prop:BD2} with $c_0=1/2$, recalling that $S^{(\sB)}$ has range $1$; otherwise, $|a_j-b_1|<4|s_j|\prec 1$, and they follow from \eqref{prop:ThfadC0}, since the zero mode cancels in $f_1$ and $f_2$):
\begin{align}\label{eq:f12}
    \left|f_0(a_j, s_j)\right| &\lesssim B_{t,|a_j - b_1|} \le B_{t, 0},\quad \left|f_1(a_j, s_j)\right| \prec \frac{(\ilambda^2+|1-t|)^{-1}}{|a_j - b_1|^{d-1}+1},\quad 
    \left|f_2(a_j, s_j)\right| \prec \frac{(\ilambda^2+|1-t|)^{-1}}{|a_j - b_1|^{d}+1}.
\end{align}
We view $\Sigma^{(\emptyset)}(t, \bsig, \mathbf{b}) = g(\mathbf{s})$ as a function of the shifts $\mathbf{s} = (s_2, \ldots, s_n)$. Then, we can write 
\begin{equation*}
    \sum_{\mathbf{b}\setminus \{b_1\}} \Sigma^{(\emptyset)}(t, \bsig, \mathbf{b}) \prod_{j=2}^n \tilde{\Theta}_{t,a_j b_j} = \sum_{\mathbf{s}} g(\mathbf{s}) \prod_{j=2}^n \sum_{\xi_j\in\{0,1,2\}} f_{\xi_j}(a_j, s_j).
\end{equation*}
We split the sum above into several parts.
\begin{enumerate}

    \item Consider the terms where $\xi_j = 0$ for all $j\in\qqq{2,n}$. Then, we use the sum-zero property \eqref{eq:Sigma-empty-sum-zero} to get 
    \begin{align*}
        \sum_{b_1}\bigg|\sum_{\mathbf{s}} g(\mathbf{s}) \prod_{j=2}^n f_{0}(a_j, s_j)\bigg| &\lesssim (1-t)\sum_{b_1}\prod_{j=2}^n \left|\tilde{\Theta}_{t,a_jb_1}\right|\lesssim (1-t) B_{t, 0}^{n-2} \sum_{b_1} \left|\tilde{\Theta}_{t,a_2b_1}\right| \lesssim B_{t, 0}^{n-2},
    \end{align*}
    where in the second step, we use \eqref{eq:pointwise_Theta} for all but one $\tilde{\Theta}_{t}$ factor, and in the last step, we use \eqref{eq:THETAinftinf}. 

    \item Consider the case where exactly one $\xi_j$ equals $1$ while all other $\xi_j$'s are zero. In this situation, the sum vanishes by the skew symmetry of $f_1$ (namely, $f_1(a_j, s_j) = -f_1(a_j, -s_j)$) together with the symmetry of $\Sigma^{(\emptyset)}$ (i.e., $g(\mathbf{s}) = g(-\mathbf{s})$). The latter symmetry follows from the translation invariance and symmetry of the $\Theta$-propagators (\Cref{lem_propTH}) and of the $M$-loops (\Cref{lem:propM}).
    
    \item Suppose there is at least one $i\in\qqq{2,n}$ such that $\xi_i = 2$. Then, using \eqref{eq:f12} and \eqref{eq:Sigma-empty-sum-zero}, we get 
    \begin{equation*}
     \qquad   \sum_{b_1}\bigg|\sum_{\mathbf{s}} g(\mathbf{s}) f_2(a_i, s_i) \prod_{j\notin \{1,i\}} f_{\xi_j}(a_j, s_j)\bigg| \prec \frac{B_{t, 0 }^{n-2}}{\ilambda^2+|1-t|} \sum_{b_1} \frac{1}{|a_i-b_1|^{d}+1} \sum_{\mathbf{b}\setminus \{b_1\}} \big|\Sigma^{(\emptyset)}(t, \bsig, \mathbf{b})\big| \prec B_{t, 0 }^{n-2}.
    \end{equation*}

    \item Suppose there are at least two $i,k\in \qqq{2,n}$ such that $i\neq k$ and $\xi_i=\xi_k = 1$. In this case, using \eqref{eq:f12} and \eqref{eq:Sigma-empty-sum-zero}, we get 
    \begin{align*}
     \quad   \sum_{b_1}\bigg|\sum_{\mathbf{s}} g(\mathbf{s}) f_1(a_i, s_i) f_1(a_k, s_k)\prod_{j\notin \{1,i,k\}} f_{\xi_j}(a_j, s_j)\bigg| \prec \frac{B_{t, 0 }^{n-3}}{\ilambda^2+|1-t|} \sum_{b_1} \frac{1}{|a_i-b_1|^{d-1}+1}\frac{1}{|a_k-b_1|^{d-1}+1} \\
        \lesssim  B_{t, 0 }^{n-2}  \sum_{b_1} \left[\frac{1}{|a_i-b_1|^{d}+1}\frac{1}{|a_k-b_1|^{d-2}+1}+\frac{1}{|a_i-b_1|^{d-2}+1}\frac{1}{|a_k-b_1|^{d}+1}\right] \lesssim B_{t, 0 }^{n-2},
        % \\&
        % \lesssim B_{t, 0 }^{n-2} \sum_{b_1} \left[\frac{1}{|a_i-b_1|^{d}+1}+\frac{1}{|a_k-b_1|^{d}+1}\right] \prec B_{t, 0 }^{n-2},
    \end{align*}
    where, in the second step, we also use $(\ilambda^2+|1-t|)^{-1}\lesssim B_{t,0}$ and the elementary inequality
    \[\frac{1}{(x^{d-1}+1)(y^{d-1}+1)}\le 2\left[\frac{1}{(x^{d}+1)(y^{d-2}+1)}+\frac{1}{(x^{d-2}+1)(y^{d}+1)}\right],\qquad \forall\, x,y\ge 0 .\]
    This asymmetric splitting is the modification of the argument in \cite[Lemma 3.11]{YY_25} needed for $d\ge 3$: since $d-2\ge 1$, each of the two resulting lattice sums over $b_1$ in the second line is $\OO(1)$, uniformly in $a_i$, $a_k$, and $L$. 

\end{enumerate}
This concludes the proof of the claim \eqref{eq:ind-step-bound}. 
% Now to finish the case $\pi=\emptyset$, we use \eqref{eq:ind-step-bound} to compute
% \begin{align*}
%     \left|\tilde{\mathcal{K}}^{(\emptyset)}(t, \bsig, \mathbf{a})\right| &\le \sum_{b_1} \left(\tilde{\Theta}^{(\sB)}_{t|m|^2}\right)_{a_1b_1} \left|\sum_{\mathbf{b}\setminus b_1} \Sigma^{(\emptyset)}(t, \bsig, \mathbf{b}) \prod_{i=2}^n \left(\tilde{\Theta}^{(\sB)}_{tm_i m_{i+1}}\right)_{a_i b_i}\right| \\
%     &\prec B_{t,0}\sum_{b_1} \left|\sum_{\mathbf{b}\setminus b_1} \Sigma^{(\emptyset)}(t, \bsig, \mathbf{b}) \prod_{i=2}^n \left(\tilde{\Theta}^{(\sB)}_{tm_i m_{i+1}}\right)_{a_i b_i}\right| \prec B_{t,0}^{n-1}.
% \end{align*}


Finally, we establish \eqref{eq:K-pi-bound} by induction on the number of molecules $r$. 
Assume that \eqref{eq:K-pi-bound} holds for every $\pi$ with at most $n$ external vertices and at most $(r-1)$ molecules. 
%Suppose that \eqref{eq:K-pi-bound} holds for any $\pi$ with $\le n$ external vertices and $\le r-1$ molecules. 
Now, consider a configuration $\pi$ consisting of $r$ molecules. In the molecule-level quotient graph---where each molecule is represented as a single vertex---the molecules, together with the long internal edges, form a tree. Therefore, there exists a leaf in this tree. Without loss of generality, assume that the molecule indexed by $k=1$ in \eqref{eq:molecule-Kpi} is such a leaf, and that it is connected to $ a_1, \ldots, a_{l}$ by external edges. 
%Consider any $\pi$ with $r$ molecules. Note that in the molecule-level quotient graph, where each molecule is collapsed into a single vertex, the molecules together with the long internal edges form a tree, so there exists a leaf of this tree. Without loss of generality, assume that the molecule indexed by $k=1$ in \eqref{eq:molecule-Kpi} is a leaf and that it is connected to $a_1, \ldots, a_l$ by the external edges. 
Then, we can decompose \smash{$\tilde{\mathcal{K}}^{(\pi)}$} in \eqref{eq:wtKpi} as  
\begin{align}
    \left|\tilde{\mathcal{K}}^{(\pi)}(t, \bsig, \mathbf{a})\right| &= \bigg|\sum_{\mathbf{b}^{(1)}} \prod_{i=1}^l \tilde{\Theta}_{t,a_i b_i} \cdot \Sigma^{(\emptyset)}(t, \bsig^{(1)}, \mathbf{b}^{(1)}) \cdot\tilde{\mathcal{K}}^{(\pi')}(t, \bsig', \mathbf{a}')\bigg| \nonumber\\
    &\le \sum_{c_1}\bigg|\sum_{\mathbf{b}^{(1)}\setminus \{c_1\}} \prod_{i=1}^l \tilde{\Theta}_{t,a_i b_i} \cdot \Sigma^{(\emptyset)}(t, \bsig^{(1)}, \mathbf{b}^{(1)})\bigg| \cdot\left|\tilde{\mathcal{K}}^{(\pi')}(t, \bsig', \mathbf{a}')\right|, \label{eq:Kpipi}
\end{align}
where $\pi'$ consists of the remaining $(r-1)$ molecules, $\mathbf{b}^{(1)} = (b_1,\ldots, b_l,c_1)$ with $c_1$ denoting the vertex in the first molecule that connects to other molecules through a long internal edge, $\bsig^{(1)}=(\sig_1,\ldots,\sig_l,\sig_{l+1})$, $\mathbf{a}' = (c_1, a_{l+1}, \ldots, a_{n})$, and $\bsig' = (\sigma_{l+1}, \ldots, \sigma_n,\sig_1)$.  
Applying the induction hypothesis to bound \smash{$\tilde{\mathcal{K}}^{(\pi')}$}, and using \eqref{eq:ind-step-bound} to estimate the remaining part of the product (noting that \eqref{eq:ind-step-bound} is applicable because $\sigma_1 \neq \sigma_{l+1}$ by the definition of a molecule), we obtain that 
\begin{equation*}
    \left|\tilde{\mathcal{K}}^{(\pi)}(t, \bsig, \mathbf{a})\right| \prec B_{t, 0}^{l-1} \cdot B_{t, 0}^{n-l} = B_{t, 0}^{n-1}.
\end{equation*}
This concludes the proof of \Cref{ML:Kbound}.
\end{proof}

Finally, \Cref{lem_wardineq_K} follows from the bound \eqref{eq:ind-step-bound}, together with an induction argument analogous to the one used above in the proof of \Cref{ML:Kbound}.

\begin{proof}[\bf Proof of \Cref{lem_wardineq_K}]
By \eqref{eq_K-Kpi}, it suffices to prove that for any $\pi \subset \Z^{\rm{off}}_n$, 
\begin{equation}\label{eq:K-pi-bound_partial}
 \sum_{a_n}   \left|\mathcal{K}^{(\pi)}(t, \bsig, \mathbf{a})\right| \prec \eta_t^{-1}B_{t,0}^{n-2}.
\end{equation}
In the case $\pi=\emptyset$, using \eqref{eq:THETAinftinf} and \eqref{eq:ind-step-bound}, we can bound the LHS of \eqref{eq:K-pi-bound_partial} by 
\begin{equation}\nonumber
 \sum_{b_n} \sum_{a_n}\Theta^{(+,-)}_{t,a_n b_n} \bigg|\sum_{\mathbf{b}\setminus \{b_n\}} \Sigma^{(\pi)}(t, \bsig, \mathbf{b}) \prod_{i=1}^{n-1} \tilde{\Theta}_{t,a_i b_i}\bigg| \le \frac{1}{1-t}\sum_{b_n} \bigg|\sum_{\mathbf{b}\setminus \{b_n\}} \Sigma^{(\pi)}(t, \bsig, \mathbf{b}) \prod_{i=1}^{n-1} \tilde{\Theta}_{t,a_i b_i}\bigg| \prec \eta_t^{-1}B_{t,0}^{n-2}.
\end{equation} 
Here, in the last step, we have used that the external edge at $a_n$ is $\Theta^{(\sig_n,\sig_1)}_t$ with $\sig_n\ne\sig_1$, so that \eqref{eq:ind-step-bound} applies with $b_n$ playing the role of $b_1$. If $\sig_n=\sig_1$, this edge is short-range, and \eqref{prop:ThfadC_short} gives $\sum_{a_n}|\Theta^{(\sig_n,\sig_1)}_{t,a_nb_n}|=\OO(1)$. Then, bounding all but one of the remaining external edges pointwise as in \eqref{eq:pointwise_Theta}, summing the last one using \eqref{eq:THETAinftinf}, and summing over the other vertices of the molecule using \eqref{eq:molecule-decay}, we obtain the bound $\OO\big((1-t)^{-1}B_{t,0}^{n-2}\big)\lesssim\eta_t^{-1}B_{t,0}^{n-2}$ in this case as well. 
For $\pi\ne \emptyset$, we argue by induction as in the proof of \Cref{ML:Kbound}.
Assume that \eqref{eq:K-pi-bound_partial} holds for any $\pi$ with at most $n$ external vertices and at most $(r-1)$ molecules. Using the notation of \eqref{eq:Kpipi}, and applying \eqref{eq:ind-step-bound} together with the induction hypothesis, we obtain
\begin{align*}
\sum_{a_n} \left|\tilde{\mathcal{K}}^{(\pi)}(t, \bsig, \mathbf{a})\right| &\le \sum_{c_1}\bigg|\sum_{\mathbf{b}^{(1)}\setminus \{c_1\}} \prod_{i=1}^l \tilde{\Theta}_{t,a_i b_i} \cdot \Sigma^{(\emptyset)}(t, \bsig^{(1)}, \mathbf{b}^{(1)})\bigg| \cdot \sum_{a_n}\left|\tilde{\mathcal{K}}^{(\pi')}(t, \bsig', \mathbf{a}')\right| \\
&\prec B_{t, 0}^{l-1} \cdot \eta_t^{-1}B_{t, 0}^{n-l-1} =\eta_t^{-1} B_{t, 0}^{n-2}. 
\end{align*}
This concludes the proof of \eqref{eq:K-pi-bound_partial}.
\end{proof}


















\section{Proofs of auxiliary graphical lemmas}


\subsection{Proof of \texorpdfstring{\Cref{lem:LWterm_EXP}}{Lemma ExpLW}}\label{subsec:pf-LWterm_EXP}


Recall the light-weight term in \eqref{eq:EGC}. To prove \Cref{lem:LWterm_EXP}, it suffices to establish 
\begin{align}
 W^{-2d}\sum_{a_1, a_2}S^{(\sB)}_{a_1a_2} \sum_{x\in [a],y\in [b],\al \in [a_2]}\E       
 \tr\big( \Gc_t E_{a_1} \big)  (G_t(\sig))_{yx} (G_t)_{x\al }(G_t)_{\al y} \prec \eta_t^{-1}(W^{-d}B_{t,0})^{5/2}\label{eq:ELW_term}
\end{align} 
for any $\sig\in\{+,-\}$. Here $[a]$ and $[b]$ denote the two external blocks appearing in \eqref{eq:EGC}, whereas $a_1$ and $a_2$ are summation indices. For brevity, write $G_t\equiv G$, and assume $\sig=-$ in the following proof. The case $\sig=+$ is analogous. %---and in fact slightly simpler---since the LHS of \eqref{eq:ELW_term} contains no negatively charged resolvent entries. 
We first focus on the random band matrix model. Performing the $GG$ expansion in \eqref{Oe2x} with respect to $G_{x\al}G_{\al y}$, we can expand the LHS of \eqref{eq:ELW_term} as 
\begin{align*}
 &~mW^{-2d}\sum_{a_1}S^{(\sB)}_{a_1 b} \sum_{x\in [a],y\in [b]}\E  \tr \big(\Gc E_{a_1} \big) |G_{xy}|^2 \tag{$I_1$} \\
 +&~m^3 W^{-2d}\sum_{a_1, a_2}S^{(\sB)}_{a_1a_2} \sum_{x\in [a],y\in [b],\al \in [a_2]} S^+_{\al y} \E       
 \tr \big(\Gc E_{a_1} \big) |G_{xy}|^2 \tag{$J_1$}\\
+&~m W^{-2d}\sum_{a_1, a_2}S^{(\sB)}_{a_1a_2} \sum_{x\in [a],y\in [b],\al \in [a_2]}\sum_{\beta}\E       
 \tr \big(\Gc E_{a_1} \big) (S_{\al\beta} \Gc_{\beta\beta})  G_{x\al }G_{\al y} \overline G_{xy} \tag{$I_2$}\\
 +&~m^3 W^{-2d}\sum_{a_1, a_2}S^{(\sB)}_{a_1a_2} \sum_{x\in [a],y\in [b],\al \in [a_2]}\sum_{\beta,\gamma}S^{+}_{\al\beta}\E       
 \tr \big(\Gc E_{a_1} \big) (S_{\beta\gamma} \Gc_{\gamma\gamma})  G_{x\beta }G_{\beta y} \overline G_{xy} \tag{$J_2$}\\
 +&~m W^{-2d}\sum_{a_1, a_2}S^{(\sB)}_{a_1a_2} \sum_{x\in [a],y\in [b],\al \in [a_2]}\sum_{\beta}\E       
 \tr \big(\Gc E_{a_1} \big) (\Gc_{\al\al} S_{\al\beta})G_{x\beta }G_{\beta y} \overline G_{xy} \tag{$I_3$}\\
+&~m^3 W^{-2d}\sum_{a_1, a_2}S^{(\sB)}_{a_1a_2} \sum_{x\in [a],y\in [b],\al \in [a_2]}\sum_{\beta,\gamma}S^{+}_{\al\beta} \E       
 \tr \big(\Gc E_{a_1} \big) (\Gc_{\beta\beta} S_{\beta\gamma})G_{x\gamma }G_{\gamma y} \overline G_{xy} \tag{$J_3$}\\
-&~m W^{-2d}\sum_{a_1, a_2}S^{(\sB)}_{a_1a_2} \sum_{x\in [a],y\in [b],\al \in [a_2]}\sum_{\beta}S_{\al\beta} \E       
 G_{x\al }G_{\beta y} \partial_{\beta\al}\left(\tr \big(\Gc E_{a_1} \big) \overline G_{xy}\right) \tag{$I_4$}\\
-&~m^3 W^{-2d}\sum_{a_1, a_2}S^{(\sB)}_{a_1a_2} \sum_{x\in [a],y\in [b],\al \in [a_2]}\sum_{\beta,\gamma}S^+_{\al\beta}S_{\beta\gamma} \E      
 G_{x\beta}G_{\gamma y} \partial_{\gamma\beta}\left(\tr \big(\Gc E_{a_1} \big) \overline G_{xy}\right) \tag{$J_4$},
\end{align*} 
where we recall that $\partial_{\beta\al}\equiv \partial_{(H_t)_{\beta\al}}$. We will only estimate the terms $I_i$ for $i\in\qqq{4}$; the terms $J_i$ can be treated in exactly the same way by using \eqref{eq:def-Spm}, together with the bound \eqref{prop:ThfadC_short}. 


For the term $I_1$, we have
\begin{align}
I_1&= m \sum_{a_1}S^{(\sB)}_{a_1 b} \E \tr \big(\Gc E_{a_1} \big)\cal L^{(2)}_{t,(-,+),(a,b)} \nonumber\\
&=m \sum_{a_1}S^{(\sB)}_{a_1 b} \E \tr \big(\Gc E_{a_1} \big)\cal K^{(2)}_{t,(-,+),(a,b)} +  m \sum_{a_1}S^{(\sB)}_{a_1 b} \E \tr \big(\Gc E_{a_1} \big) (\cal L-\cK)^{(2)}_{t,(-,+),(a,b)} \prec (W^{-d}B_{t,0})^3, \label{eq:termI1}
\end{align}
where in the last step, we use \eqref{eq:bcal_k} and \eqref{res_ELK_n=1} to bound the first term, and \eqref{Gt_avgbound_flow} together with \eqref{Eq:L-KGt-flow} to bound the second.
For $I_2$, applying the averaged local law \eqref{Gt_avgbound_flow} to both $\tr( \Gc E_{a_1} )$ and $\sum_\beta S_{\al\beta} \Gc_{\beta\beta}$ yields
\begin{align}
    I_2&\prec (W^{-d}B_{t,0})^2 \cdot W^{-2d}\sum_{a_1, a_2}S^{(\sB)}_{a_1a_2} \sum_{x\in [a],y\in [b],\al \in [a_2]}\E  |G_{x\al}||G_{\al y}| |G_{xy}| \nonumber\\
    &\lesssim (W^{-d}B_{t,0})^2 \cdot W^{-2d}\sum_{x\in [a],y\in [b]} \sum_\al \E  |G_{x\al}||G_{\al y}| |G_{xy}| \nonumber\\
    &\prec \eta_t^{-1}(W^{-d}B_{t,0})^2 \cdot W^{-2d}\sum_{x\in [a],y\in [b]}\E|G_{xy}| \prec \eta_t^{-1}(W^{-d}B_{t,0})^{5/2},\label{eq:termI2}
\end{align}
where in the third step we use Cauchy--Schwarz together with Ward's identity, and in the last step the local law \eqref{Gt_bound_flow} to control the averages over $x\in[a]$ and $y\in[b]$. The term $I_3$ can be handled in exactly the same way. It remains to estimate $I_4$. Splitting according to which factor the partial derivative acts on, we obtain
\begin{align*}
I_4&=m W^{-3d}\sum_{a_1, a_2,a_3}S^{(\sB)}_{a_1a_2}S^{(\sB)}_{a_2a_3} \sum_{x\in [a],y\in [b],\al \in [a_2],\beta\in[a_3]}  \E  \tr\big( \Gc E_{a_1} \big) |G_{x\al}|^2 |G_{\beta y}|^2 \\
&\quad +m W^{-4d}\sum_{a_1, a_2, a_3}S^{(\sB)}_{a_1a_2}S^{(\sB)}_{a_2a_3} \sum_{x\in [a],y\in [b],\gamma\in[a_1],\al \in [a_2],\beta\in[a_3]} \E       
 G_{x\al }G_{\al\gamma}G_{\gamma\beta}G_{\beta y}  \overline G_{xy} \\
&= mW^d\sum_{a_1, a_2,a_3}S^{(\sB)}_{a_1a_2}S^{(\sB)}_{a_2a_3} \E  \tr\big( \Gc E_{a_1} \big) \cL^{(2)}_{t,(-,+),(a,a_2)}\cL^{(2)}_{t,(-,+),(a_3,b)} +mW^d \sum_{a_1, a_2, a_3}S^{(\sB)}_{a_1a_2}S^{(\sB)}_{a_2a_3} \E \cL^{(5)}_{t,\bsig_5,\ba_5} \\
&=:I_{41}+I_{42},
\end{align*}
where $\bsig_5=(+,+,+,+,-)$ and $\ba_5=(a_2,a_1,a_3,b,a)$. For $I_{41}$, we decompose each 2-$G$-loop as $(\cL-\cK)^{(2)}+\cK^{(2)}$, giving 
\begin{align}
    I_{41}&= mW^d\sum_{a_1, a_2,a_3}S^{(\sB)}_{a_1a_2}S^{(\sB)}_{a_2a_3} \E  \tr\big( \Gc E_{a_1} \big)  \cL^{(2)}_{t,(-,+),(a,a_2)}\cK^{(2)}_{t,(-,+),(a_3,b)} \nonumber\\
    &\quad +mW^d\sum_{a_1, a_2,a_3}S^{(\sB)}_{a_1a_2}S^{(\sB)}_{a_2a_3} \E  \tr\big( \Gc E_{a_1} \big)  \cL^{(2)}_{t,(-,+),(a,a_2)}(\cL-\cK)^{(2)}_{t,(-,+),(a_3,b)}\nonumber\\ 
    &= mW^d\sum_{a_1, a_2,a_3}S^{(\sB)}_{a_1a_2}S^{(\sB)}_{a_2a_3} \E  \tr\big( \Gc E_{a_1} \big) \cL^{(2)}_{t,(-,+),(a,a_2)}\cK^{(2)}_{t,(-,+),(a_3,b)}+ \OO_\prec\p{(W^{-d}B_{t,0})^3}\cdot W^d\sum_{ a_2}\E \cL^{(2)}_{t,(-,+),(a,a_2)} \nonumber\\
    &= mW^d\sum_{a_1, a_2,a_3}S^{(\sB)}_{a_1a_2}S^{(\sB)}_{a_2a_3} \E  \tr\big( \Gc E_{a_1} \big) \cK^{(2)}_{t,(-,+),(a,a_2)}\cK^{(2)}_{t,(-,+),(a_3,b)}\nonumber\\
    &\quad +mW^d\sum_{a_1, a_2,a_3}S^{(\sB)}_{a_1a_2}S^{(\sB)}_{a_2a_3} \E  \tr\big( \Gc E_{a_1} \big) (\cL-\cK)^{(2)}_{t,(-,+),(a,a_2)}\cK^{(2)}_{t,(-,+),(a_3,b)}+ \OO_\prec\p{\eta_t^{-1}(W^{-d}B_{t,0})^3} \nonumber\\
    &\prec (W^{-d}B_{t,0})^3\cdot W^d\sum_{a_3}\cK^{(2)}_{t,(-,+),(a_3,b)} + \eta_t^{-1}(W^{-d}B_{t,0})^3 \prec \eta_t^{-1}(W^{-d}B_{t,0})^3.\label{eq:termI41}
\end{align}
Here, we use \eqref{Eq:L-KGt-flow} and \eqref{Gt_avgbound_flow} in the second step; Ward's identity \eqref{WI_calL} for \smash{$\sum_{a_2} \cL^{(2)}_{t,(-,+),(a,a_2)}$} and the local law \eqref{Gt_bound_flow} in the third; \eqref{Eq:L-KGt-flow}, \eqref{eq:bcal_k}, and \eqref{res_ELK_n=1} in the fourth; and Ward's identity \eqref{WI_calK} for \smash{$\sum_{a_3}\cK^{(2)}_{t,(-,+),(a_3,b)}$} in the last. 
Finally, for the term $I_{42}$, we rewrite it as an average over the graphs $\cal G_{xy}$: 
\begin{align}\label{eq;I42inG}
 I_{42}=mW^{-2d} \sum_{x\in [a],y\in [b]} \E\cal G_{xy},\quad \text{where}\quad \cal G_{xy}:=\sum_{\gamma,\al,\beta}S_{\gamma\al}S_{\al\beta} \p{G_{x\al }G_{\al\gamma}G_{\gamma\beta}G_{\beta y}  \overline G_{xy}}.
\end{align}
We claim the following bound for $\E\cal G_{xy}$: 
\begin{align}\label{eq;EGxy:x=y}
\E\cal G_{xy} \prec \mathbf 1_{x = y}\cdot  \eta_t^{-1} (W^{-d}B_{t,0})^{2}+\mathbf 1_{x\ne y}\cdot \eta_t^{-1}(W^{-d}B_{t,0})^{5/2}. 
\end{align}
Substituting this estimate into \eqref{eq;I42inG} gives
\[ I_{42} \prec \eta_t^{-1} (W^{-d}B_{t,0})^{5/2}.\]
Together with \eqref{eq:termI1}--\eqref{eq:termI41}, this completes the proof of \eqref{eq:ELW_term}.


Finally, we prove the estimate \eqref{eq;EGxy:x=y}. To this end, we apply the $GG$-expansion from \Cref{T eq0} at the non-standard neutral vertices $\al,\gamma,\beta$, which expands $\cal G_{xy}$ into a collection of new graphs:
\begin{align*}
    \cal G_{xy}=_{\E} \sum_\mu \Gamma_{\mu,xy} \; .
\end{align*}
By \Cref{GtoAG}, each $\Gamma_{\mu,xy}$ can be bounded through its auxiliary graph \smash{$\Gamma^{\aux}_{\mu,[a][b]}$} as in \eqref{G_by_auxG}, with $q\in\{0,1\}$ internal molecules and $\Psi_t=(W^{-d}B_{t,0})^{1/2}$. 
If an auxiliary graph contains an internal molecule, then that molecule is attached to at least two solid edges corresponding to the $\xi$-variables defined in \eqref{eq:xia1a2}. Indeed, every vertex of $\cal G_{xy}$ is incident to an even number of solid half-edges (a light-weight counting twice), and $\cal G_{xy}$ is connected through its solid and waved edges; both properties are preserved by the local expansions and by the identification of vertices. Since different molecules are joined only by solid edges, an internal molecule of $\Gamma_{\mu,xy}$ is therefore joined to the other molecules by a positive even number of solid edges. Applying Cauchy--Schwarz together with \eqref{eq:Gbyxi2} therefore gives
\[ \Gamma^{\aux}_{\mu, [a][b]} \prec (W^d\eta_t)^{-q}(\Psi_t)^{\ord(\Gamma^{\aux}_{\mu, [a][b]})}. \]
Combining this with \eqref{G_by_auxG} implies that %, for any large constant $D>0$, 
\begin{align}\label{Gammamuxy}
    \Gamma_{\mu, xy} \prec \eta_t^{-1} (W^{-d}B_{t,0})^{\frac12\ord(\Gamma_{\mu, xy})} .% + W^{-D}.
\end{align}
With this estimate in hand, to conclude \eqref{eq;EGxy:x=y} it remains to show that the scaling order of each $\Gamma_{\mu,xy}$ satisfies
\be\label{eq:sizeGammamu_E}
\ord(\Gamma_{\mu,xy}) \ge 4\cdot \mathbf 1_{x= y}+5\cdot \mathbf 1_{x\ne y} , 
\ee
where $x$ and $y$ are regarded as external vertices.


The three remaining cases below all rest on the following fact about the graphs \smash{$\cal G^{(21)}_{xy}$} and \smash{$\cal G^{(4)}_{xy}$} defined below (and the analogue of \smash{$\cal G^{(21)}_{xy}$} in case (3)), which is verified by inspecting the terms on the right-hand side of \eqref{Oe2x}:
\be\label{eq:GGraisesord}
\text{\emph{applying the $GG$ expansion \eqref{Oe2x} at a non-standard neutral vertex strictly increases the scaling order.}}
\ee
This fact does not hold for arbitrary graphs. For instance, if the two solid edges at the expansion vertex form a $2$-cycle, i.e., $\cal G=G_{xy}G_{yx}f(G)$ in the notation of \eqref{Oe2x}, then the parts of the first two terms on the right-hand side of \eqref{Oe2x} in which the resulting diagonal entry $G_{yy}$ is replaced by $m$ have the same scaling order as $\cal G$.
To obtain \eqref{eq:sizeGammamu_E}, we decompose $\cal G_{xy}$ into four parts $\cal G^{(i)}_{xy}$, $i\in\{1,2,3,4\}$, corresponding to the cases: (1) $\al=\gamma=\beta$, (2) $\al=\gamma\ne \beta$, (3) $\al\ne \gamma = \beta$, and (4) $\al\ne \gamma$ and $\gamma\ne \beta$. It is easy to see that \smash{$\ord(\cal G^{(1)}_{xy})$} already satisfies the lower bound in \eqref{eq:sizeGammamu_E}. Next, in case (2), we write 
\[ \cal G^{(2)}_{xy}= m\sum_{\al\ne\beta}S_{\al\al}S_{\al\beta} \p{G_{x\al }G_{\al\beta}G_{\beta y}  \overline G_{xy}}+ \sum_{\al\ne\beta}S_{\al\al}S_{\al\beta} \Gc_{\al\al} \p{G_{x\al }G_{\al\beta}G_{\beta y}  \overline G_{xy}}=:\cal G^{(21)}_{xy}+\cal G^{(22)}_{xy}.\]
The second graph \smash{$\cal G^{(22)}_{xy}$} already meets the lower bound in \eqref{eq:sizeGammamu_E}, while the first graph \smash{$\cal G^{(21)}_{xy}$} has scaling order $3\cdot \mathbf 1_{x= y}+4\cdot \mathbf 1_{x\ne y} .$ To raise its scaling order, we expand \smash{$\cal G^{(21)}_{xy}$} using the $GG$ expansion \eqref{Oe2x} at the vertex $\al$. By \eqref{eq:GGraisesord}, every graph produced by this expansion has scaling order strictly larger than that of \smash{$\cal G^{(21)}_{xy}$}, and thus satisfies \eqref{eq:sizeGammamu_E}.
Case (3) is handled in exactly the same manner as case (2).
%One can check directly that all resulting graphs from the expansions have strictly higher scaling order than \smash{$\cal G^{(21)}_{xy}$}, and hence satisfy the lower bound in \eqref{eq:sizeGammamu_E}. Since checking this fact is straightforward, we omit the details. The case (3) can be handled in the same way as case (2). 
Finally, in case (4), the graph \smash{$\cal G^{(4)}_{xy}$} has scaling order $2\cdot \mathbf 1_{x= y}+3\cdot \mathbf 1_{x\ne y},$ which has to be raised by $2$. We first apply the $GG$ expansion \eqref{Oe2x} at the vertex $\al$; by \eqref{eq:GGraisesord}, every resulting graph has scaling order at least $3\cdot \mathbf 1_{x= y}+4\cdot \mathbf 1_{x\ne y}$. A second $GG$ expansion need not raise the scaling order further, because the first one may create $2$-cycles: for instance, the term in which the derivative $\partial_{h_{\al'\al}}$ in \eqref{Oe2x} acts on $G_{\gamma\beta}$ contains the factor $G_{\al'\gamma}G_{\gamma\al'}$. We therefore proceed as follows: in every graph that does not yet satisfy \eqref{eq:sizeGammamu_E}, we apply \eqref{Oe2x} at one of its remaining non-standard neutral vertices, and we repeat. A direct inspection of the resulting finite expansion tree shows that, for any choice of the expansion vertices, every branch ends after at most three $GG$ expansions in a graph satisfying \eqref{eq:sizeGammamu_E}, for some branches with equality.
This yields \eqref{eq:sizeGammamu_E} and thus completes the proof of \eqref{eq;EGxy:x=y}.

%\[\cal G^{(4)}_{xy}:=\sum_{\al,\beta}\sum_{\gamma\notin \{\al ,  \beta\}}S_{\gamma\al}S_{\al\beta} \p{G_{x\al }G_{\al\gamma}G_{\gamma\beta}G_{\beta y}  \overline G_{xy}}.\]
%One can check directly that each application of \eqref{Oe2x} strictly increases the scaling order of the resulting graphs. Thus, the two $GG$ expansions at the vertices $\al$ and $\beta$ increase the scaling order by at least 2 in total, thus yielding \eqref{eq:sizeGammamu_E} and concluding the proof of \eqref{eq;EGxy:x=y}. 
% {\cor Next, by \Cref{GtoAG}, we can bound each $\Gamma_{\mu,xy}$ via its auxiliary graph \smash{$\Gamma^{\aux}_{\mu, [a][b]}$} as in \eqref{G_by_auxG}, with $q\in \{0, 1\}$ and $\Psi_t=(W^{-d}B_{t,0})^{1/2}$. The auxiliary graphs \smash{$\Gamma^{\aux}_{\mu, [a][b]}$} are controlled by \eqref{eq:bddGamma_aux} with $p=0$ and $q\in \{0, 1\}$. Consequently,
% \begin{align*}
%     \Gamma_{\mu, xy} \prec  (\Psi_t)^{\ord(\Gamma_{\mu, xy}) - \ord(\Gamma^{\aux}_{\mu, [a][b]})}  \cdot W^{qd}\Gamma^{\aux}_{\mu, [a][b]} + W^{-D}\prec \eta_t^{-1}(\Psi_t)^{\ord(\Gamma_{\mu, xy})}.
% \end{align*}
% Combining this with \eqref{eq:sizeGammamu_E} yields the desired bound \eqref{eq;EGxy:x=y}. } 

For the block Anderson model, the proof of \Cref{lem:LWterm_EXP} follows the same outline, in the graphical notation of \Cref{subsec:LWchange-to-BA} below: we have $S^{(\sB)}=I_{L^d}$ (so that $a_2=a_1$ in \eqref{eq:ELW_term}), every solid edge is split as $G=\Gc+M$, and the $GG$ expansion \eqref{GGGamma} (and its complex conjugate) is used in place of \eqref{Oe2x}. However, the argument above does not carry over verbatim, and the following three modifications are needed.
\begin{itemize}
\item[(i)] \emph{Twisted loops.} Since $M_{uv}=0$ unless $u-W[u]=v-W[v]$ (see \eqref{eq:Psi3D} and \eqref{def_G0}), summing over the two ends of an $M$-dotted edge produces, besides $G$-loops, the twisted traces $\tr(\Gc_t(\sig)P_{c_1c_2})$ and the twisted $2$-loops $\wt{\cal L}_{t,\bsig,(a;c_1,c_2)}=\tr(G_t(\sig_1)E_aG_t(\sig_2)P_{c_1c_2})$ introduced in \Cref{sec:ext-to-BA}. In coordinates,
\[ \tr\big(\Gc_t P_{c_1c_2}\big)=W^{-d}\sum_{u\in[c_1]}(\Gc_t)_{u'u} ,\qquad \wt{\cal L}_{t,\bsig,(a;c_1,c_2)}=W^{-2d}\sum_{x\in[a]}\sum_{u\in[c_1]}\big(G_t(\sig_1)\big)_{u'x}\big(G_t(\sig_2)\big)_{xu}, \]
where, for $u\in[c_1]$, $u'$ denotes the vertex of $[c_2]$ with $u'-W[u']=u-W[u]$. For $c_1=c_2$ these are the light-weight $\tr(\Gc_t E_{c_1})$ and the $2$-loop $\cL^{(2)}_{t,\bsig,(a,c_1)}$. For $c_1\ne c_2$ they are not $G$-loops in the sense of \Cref{Def:G_loop}, and bounding them entrywise by \eqref{Gt_bound_flow} loses a factor $(W^{-d}B_{t,0})^{-1/2}$. Twisted traces replace the light-weights at the expansion vertex in the analogues of the terms $I_3$ and $J_3$, and twisted $2$-loops replace the $2$-loop $\cL^{(2)}_{t,(-,+),(a,a_2)}$ in the analogue of $I_{41}$. We bound the twisted traces by \Cref{lem:TL1_BA}, which gives the same bound $W^{-d}B_{t,0}$ as \eqref{Gt_avgbound_flow} gives for light-weights. We write each twisted $2$-loop as the twisted tree $\wt{\cal K}_{t,\bsig,(a;c_1,c_2)}$ plus a fluctuation, which is bounded by \eqref{eq:TL2_BA}; if a leg is circled, we write $\Gc_t(\sig)=G_t(\sig)-M(\sig)$, and the contribution of $M(\sig)$ is $W^{-d}$ times an entry of $M^{(\sB)}(\sig)$ times a twisted trace. By \eqref{prop:ThfadC}, \eqref{eq:THETAinftinf} and \eqref{eq:WardM}, the twisted tree satisfies $|\wt{\cal K}_{t,\bsig,(a;c_1,c_2)}|\lesssim W^{-d}B_{t,0}$ and $\sum_a|\wt{\cal K}_{t,\bsig,(a;c_1,c_2)}|\lesssim (W^d\eta_t)^{-1}$, which are the properties of the $\cK$-loops used in \eqref{eq:termI41}. With these replacements, the arguments leading to \eqref{eq:termI1}--\eqref{eq:termI41} carry over. Since \eqref{eq:TL2_BA} gives the power $3/2$ of $W^{-d}B_{t,0}$ instead of the power $2$ in \eqref{Eq:L-KGt-flow}, the analogue of \eqref{eq:termI41} is $\eta_t^{-1}(W^{-d}B_{t,0})^{5/2}$, which still suffices for \eqref{eq:ELW_term}.

\item[(ii)] \emph{Scaling orders at the level of atoms.} Scaling orders are now defined by \eqref{eq:ordG_BA}, and the analogue of \eqref{eq:GGraisesord} for \eqref{GGGamma} holds only under additional conditions. Before the new solid edges are split as $G=\Gc+M$, the first sum on the right-hand side of \eqref{GGGamma} removes two solid edges, adds one waved edge, and merges the atoms of $y$ and $y'$; hence it raises the scaling order by $2\nu$, where $\nu\in\{0,1\}$ is the number of internal atoms lost in this merger. In particular, it does not raise the scaling order if $y$ and $y'$ lie in the same atom, or if both lie in external atoms. The other terms raise the scaling order by exactly $1$. Accordingly, we apply \eqref{GGGamma} to $\Gc_{y'v}\Gc_{vy}f(G)$ at a vertex $v$ only if the following conditions hold:
\begin{itemize}
\item[(H1)] the atom of $v$ is internal, and $\Gc_{y'v}$ and $\Gc_{vy}$ are the only solid edges incident to it;
\item[(H2)] the atoms of $v$, $y$, and $y'$ are pairwise distinct;
\item[(H3)] at least one of the atoms of $y$ and $y'$ is internal.
\end{itemize}
Under these conditions, every graph produced by \eqref{GGGamma} at $v$ has scaling order at least that of the input graph plus $1$. Indeed, $\nu=1$ by (H2) and (H3). Moreover, if $j$ of the new solid edges of a term are replaced by $M$-dotted edges, the scaling order of the term changes by $2\nu'-j$, where $\nu'$ is the number of internal atoms lost in the resulting mergers. Each such $M$-dotted edge joins two different atoms, at least one of which is internal (the atom of $v$, which is internal by (H1), an atom created by the expansion, or one of the atoms of $y$ and $y'$ as in (H3)), and checking the terms of \eqref{GGGamma} one by one shows that $\nu'\ge \lceil j/2\rceil$.
The target \eqref{eq:sizeGammamu_E} is replaced by the atom-level bound
\[ \ord(\Gamma_{\mu,xy})\ge 5-2\cdot\mathbf 1\p{x \text{ and } y \text{ lie in the same atom}} . \]
The version with $\mathbf 1_{x=y}$ fails for the block Anderson model: for instance, the part of case (1) in which all solid edges are replaced by $M$-dotted edges has scaling order $4$, while $x$ and $y$ lie in the same atom. By \eqref{scaleatom}, $x\in[a]$ and $y\in[b]$ can lie in the same atom only if $x-W[x]=y-W[y]$, which holds for at most $W^d$ of the $W^{2d}$ pairs $(x,y)$. Since $B_{t,0}\ge (\fd^{-2}+1)^{-1}$ by \eqref{eq:WO}, the resulting factor $W^{-d}\lesssim W^{-d}B_{t,0}$, obtained after averaging over $x\in[a]$ and $y\in[b]$, compensates for the weaker bound. Finally, the analogue of \eqref{Gammamuxy}, which follows from \eqref{G_by_auxG_BA} and \eqref{eq:Gbyxi2_BA}, requires that each internal molecule be joined to the other molecules by at least two solid edges. This follows from two properties that hold for $\cal G_{xy}$ and are preserved by every expansion step and by the splitting $G=\Gc+M$: each atom is incident to an even number of solid half-edges (a light-weight counting twice), and the graph is connected through its solid, waved and dotted edges. Since different molecules are joined only by solid edges, an internal molecule is then joined to the other molecules by a positive even number of solid edges.

\item[(iii)] \emph{Folding.} The case analysis below \eqref{eq:GGraisesord} is replaced by the following procedure. We expand the block Anderson analogue of the graph $\cal G_{xy}$ in \eqref{eq;I42inG} by \eqref{GGGamma}, only at vertices satisfying (H1)--(H3), and we stop as soon as a graph satisfies the bound in (ii). Along each branch, at most two such expansions are needed, and every graph at which the procedure stops either satisfies the bound in (ii) and the attachment property, or is a \emph{folded graph}: a graph of scaling order $4$, with $x$ and $y$ in different atoms, whose solid edges form the cycle $\Gc_{xu_1}\Gc_{u_1'u_2}\Gc_{u_2'y}(\Gc_t(\sig))_{yx}$, where $u_i$ and $u_i'$ are joined by $M$-dotted edges, $i=1,2$. An example is the graph $\Gc_{xw}M_{w\gamma}\Gc_{\gamma v}M_{v\delta}\Gc_{\delta y}(\Gc_t(\sig))_{yx}$ (with the waved edges of $\cal G_{xy}$), obtained from $\cal G_{xy}$ by the splitting $G=\Gc+M$. In a folded graph, no vertex is incident to two solid edges of the same charge, so \eqref{GGGamma} cannot be applied, although each of the two internal atoms carries two blue solid edges. For these graphs, we do not prove the pointwise bound \eqref{eq;EGxy:x=y}, but only its average over $x\in[a]$ and $y\in[b]$, which is all that is used in \eqref{eq;I42inG}. Since $M_{uu'}\ne0$ only if $u-W[u]=u'-W[u']$, every folded graph, after summing over all vertices other than $x$, $y$, $u_1$, $u_2$ and averaging over $x\in[a]$ and $y\in[b]$, takes the form
\[ \mathfrak F_\rho:=W^{-4d}\sum_{\mathbf c\in(\Zn)^4}\rho(\mathbf c)\sum_{x\in[a],\,y\in[b]}\ \sum_{u\in[c_1],\,v\in[c_3]}(\Gc_t)_{xu}\,(\Gc_t)_{u'v}\,(\Gc_t)_{v'y}\,\big(\Gc_t(\sig)\big)_{yx}, \]
where $\mathbf c=(c_1,c_2,c_3,c_4)$, the vertices $u'\in[c_2]$ and $v'\in[c_4]$ are determined by $u'-W[u']=u-W[u]$ and $v'-W[v']=v-W[v]$, and the coefficients $\rho(\mathbf c)$ collect the entries of $M^{(\sB)}$, the waved edges and the remaining factors, which do not depend on the vertices of the cycle. In all cases, one checks using \eqref{eq:WardM} and \eqref{Mbound_AO}--\eqref{Mbound_AO2} that
\[ r_1:=\max_{c_1}\sum_{c_2,c_3,c_4}|\rho(\mathbf c)|\prec B_{t,0},\qquad r_2:=\max_{c_4}\sum_{c_1,c_2,c_3}|\rho(\mathbf c)|\prec B_{t,0}. \]
We now show that $|\mathfrak F_\rho|\prec \eta_t^{-1}(W^{-d}B_{t,0})^{5/2}$ whenever $r_1r_2\prec B_{t,0}^2$. Writing $R:=\Gc_tE_b\Gc_t(\sig)$, we have $\mathfrak F_\rho=W^{-3d}\sum_{\mathbf c,x,u,v}\rho(\mathbf c)(\Gc_t)_{xu}(\Gc_t)_{u'v}R_{v'x}$. We apply the Cauchy--Schwarz inequality to this sum, with weight $|\rho(\mathbf c)|$, putting $(\Gc_t)_{xu}(\Gc_t)_{u'v}$ into the first factor and $R_{v'x}$ into the second. Since $R_{v'x}$ does not depend on $u\in[c_1]$, we obtain
\[ |\mathfrak F_\rho|^2\le W^{-4d}\,r_1r_2\,\|\Gc_t\|_{\max}^2\,S_AS_C,\qquad S_A:=\sum_{x\in[a]}\sum_{w}|(\Gc_t)_{xw}|^2,\quad S_C:=\sum_{x\in[a]}\big\|R\mathbf e_x\big\|^2 , \]
where $\mathbf e_x$ denotes the standard basis vector. By \eqref{Gt_bound_flow}, $\|\Gc_t\|_{\max}^2\prec W^{-d}B_{t,0}$. Since $\sum_w|(G_t)_{xw}|^2=\im (G_t)_{xx}/\eta_t$ by Ward's identity, and $\sum_w|M_{xw}|^2=\sum_c|M^{(\sB)}_{[x]c}|^2=1$ by \eqref{eq:WardM}, we have $S_A\le 2\eta_t^{-1}W^d\,\im\tr(G_tE_a)+2W^d\prec\eta_t^{-1}W^d$. For $S_C$, let $\phi_x:=E_b\Gc_t(\sig)\mathbf e_x$, a vector supported in $[b]$. We have $\|\Gc_t\phi\|^2\le 2\|G_t\phi\|^2+2\|M\phi\|^2$, where $\|G_t\phi\|^2=\eta_t^{-1}\im\langle\phi,G_t\phi\rangle$ by Ward's identity, and $\|M\phi\|=\|\phi\|$ for every vector $\phi$ supported in a single block by \eqref{eq:WardM}. Splitting $\phi_x=E_bG_t(\sig)\mathbf e_x-E_bM(\sig)\mathbf e_x$ and using that $\phi\mapsto\im\langle\phi,G_t\phi\rangle$ is a nonnegative quadratic form, we obtain
\[ S_C\le 4\eta_t^{-1}W^{d}\,\big|\cL^{(3)}_{t,(-\sig,+,\sig),(b,b,a)}\big|+4\eta_t^{-1}W^{-d}\,\big|\tr(G_tE_b)\big|+2\|\Gc_t\|_{\max}^2\prec \eta_t^{-1}W^d(W^{-d}B_{t,0})^2 . \]
Here, the second term comes from $E_bM(\sig)\mathbf e_x$, which is supported on the single vertex of $[b]$ with the same offset as $x$, and in the last step we used \eqref{Eq:LGxb} with $n=3$, $|\tr(G_tE_b)|\prec 1$, $\eta_t\le 1$, and $B_{t,0}\gtrsim 1$. Combining these bounds yields $|\mathfrak F_\rho|^2\prec r_1r_2\,\eta_t^{-2}W^{-2d}(W^{-d}B_{t,0})^3\prec \eta_t^{-2}(W^{-d}B_{t,0})^5$. Since all quantities involved are deterministically bounded by $W^C$, the same bound holds for the expectations. Together with (ii), we conclude
\[ W^{-2d}\sum_{x\in[a],\,y\in[b]}\E\,\cal G_{xy}\prec \eta_t^{-1}(W^{-d}B_{t,0})^{5/2}, \]
which is the form in which \eqref{eq;EGxy:x=y} is used in \eqref{eq;I42inG}.
\end{itemize}
With these modifications, the remaining steps are the same as for the random band matrix model.



\subsection{Proof of \texorpdfstring{\Cref{lem:localregular}}{Lemma localregular}}\label{sec:pflocalregular}

Our local expansion strategy leading to the proof of \Cref{lem:localregular} proceeds as follows. Given a graph $\Gamma$ that is not locally standard, we first find all weights in it, and apply the weight expansion \eqref{Owx} to remove them one by one. 
Once the graph is free of weights, we find all vertices connected to more than two solid edges, and apply the edge expansion \eqref{Oe1x} iteratively to reduce their degrees.  
After removing weights and ensuring each vertex is connected to at most two solid edges, we identify all vertices that are not standard neutral---that is, vertices connected to two $G$ edges or two $\overline G$ edges---and apply the $GG$ expansion \eqref{Oe2x} to correct them.
It is important to note that after each expansion, the resulting graphs may require earlier steps to be re-applied. For instance, an edge expansion might introduce new weights, which must be removed via weight expansions before proceeding with further edge or $GG$ expansions. %This recursive process explains the loops shown in Figure \ref{Fig chart1}.
We now formally state the local expansion rules. Throughout the following proof, when we refer to the degree of a vertex $\al$, we mean \emph{the number of solid edges connected to it}, excluding any solid self-loops (i.e., weights).

%Finally, given a graph without weights and where all vertices are connected with at most two solid edges, we find all vertices that are not standard neutral (i.e., those connected with two $G$ edges or two $\overline G$ edges) and apply the $GG$ expansion \eqref{Oe2x} to them one by one. Note that after an expansion, we may need to perform earlier expansions to the resulting graphs. For example, after an edge expansion, we may get graphs that contain weights. Then, before performing another edge expansion, we first need to perform weight expansions to these graphs. This explains why we have loops in Figure \ref{Fig chart1}. We now formally state the local expansion rule. In the following proof, whenever we refer to the degree of a vertex $\al$, we only count the \emph{solid edges connected with it} and we do not count the self-loops of solid edges (i.e., weights).




\begin{strategy}[Local expansion strategy]\label{strat_local}
Given an arbitrarily large constant $D>0$, we apply the following local expansion strategy. %{\cor For each step, add a decomposition step $G=(G-M)+M$}
\begin{itemize}
\item[Step 1:] Given an input graph, we apply the \emph{weight expansion} $\cal O_{\mathrm{weight}}$ as follows. If the graph contains no weights (neither regular nor light), then $\cal O_{\mathrm{weight}}$ is a null operation, and the graph is passed to the next step. Otherwise, we select one of its weights, say on vertex $\al$. If this is a light-weight, we apply the expansion \eqref{Owx} directly. If it is a regular weight, i.e., $G_{\al\al}$ or \smash{$\overline G_{\al\al}$}, we decompose it as a sum of a light-weight and a factor of $m$ or $\overline m$: \smash{$(G_{\al\al}-m)+m$} or \smash{$(\overline G_{\al\al}-\overline m)+\overline m$}, and then apply the expansion \eqref{Owx} to the resulting light-weight term.  
For each new graph generated by the weight expansion, we apply the operations in \Cref{dot-def} to write it as a sum of normal graphs. For each normal graph $\cal G$, if its scaling size is sufficiently small:
\be\label{eq:smallsize}
\size(\cal G)\le W^{-D},
\ee
then we send it directly to the output; if it contains no weights, we pass it to Step 2; if it still contains weights, we restart Step 1.  


\item[Step 2:] At this stage, the input graph has no weights. We apply the edge expansion $\cal O_{\mathrm{edge}}$ as follows. If there exists a vertex $\al$ whose degree is $\notin\{0,2\}$, or whose charge is not neutral (recall the definition in \eqref{eq:neutralcharge}),\footnote{If $\al$ is attached to two $+$ solid edges and has non-neutral charge, then the two edges incident to it must be of the form $G_{\al y}G_{\al y'}$ or $G_{y\al}G_{y'\al}$. Such configurations cannot be expanded using the $GG$ expansion \eqref{Oe2x}, and therefore must first be removed via the edge expansion \eqref{Oe1x}.} then we apply the edge expansion \eqref{Oe1x} at $\al$. 
For each resulting graph, we again apply the operations in \Cref{dot-def} to express it as a sum of normal graphs. Those satisfying the small-size condition \eqref{eq:smallsize} are sent to the output, while the rest are sent back to Step 1.
If all internal vertices have degree $\in \{0,2\}$ and are neutral in charge, then $\cal O_{\mathrm{edge}}$ is a null operation, and the graph is passed to Step 3.  

\item[Step 3:] Now, the input graph has no weights, and all internal vertices have degree $\in\{0,2\}$ and neutral charge. We apply the $GG$ expansion $\cal O_{GG}$ as follows. 
If there exists a vertex $\al$ connected to two solid edges of the same charge (i.e., two $G$ edges or two $\overline G$ edges), we apply the expansion \eqref{Oe2x} at $\al$.
Each resulting graph is expanded into a sum of normal graphs using the operations in \Cref{dot-def}. Those satisfying \eqref{eq:smallsize} are sent to the output, while the rest are sent back to Step 1.
If all internal vertices are standard neutral, then $\cal O_{GG}$ is a null operation, and the graph is sent to the output. 

%\item[Stop:] Finally, if all the above operations are null, then the input graph is locally standard, and we send it to the output.  
\end{itemize}
\end{strategy}

% Finally, we collect all the output graphs of Strategy \ref{strat_local} and obtain the following lemma. 

% \begin{lemma}[Lemma 3.22 of \cite{yang2021delocalization}] \label{lvl1 lemma}
% Let $\Gamma$ be an arbitrary graph with $\OO(1)$ many vertices and edges. For any large constant $D>0$, we can expand $\Gamma$ into a sum of $\OO(1)$ many locally standard graphs $\Gamma_\mu$:
% \begin{align}\label{expand lvl1}
% \Gamma =_{\E} \sum_\mu \Gamma_\mu + \Err,
% \end{align}
% where every $\Gamma_\mu$ has scaling size $\lesssim \size(\Gamma)$ and scaling order $\ge \ord(\Gamma)$, and $\Err$ is a sum of graphs of scaling size $\OO(W^{-D})$.
% \end{lemma}


We now apply \Cref{strat_local} to complete the proof of \Cref{lem:localregular}.

\begin{proof}[\bf Proof of \Cref{lem:localregular}]
By \Cref{lvl1 lemma}, we obtain the expansion of $\big|f_{xy}(G)\big|^{p}$ as in \eqref{eq:local_Gs}, where the graphs $\Gamma_{\mu,xy}$ trivially satisfy property (1). For property (2), note that the original graph 
\be\label{eq:originGamma}\Gamma_{xy}=\big|f_{xy}(G)\big|^{p}= \sum_{\bal,\bbeta}\prod_{i=1}^p \left[S_{\al_i\beta_i} \mathbf 1_{x\ne \al_i}\mathbf 1_{y\ne \al_i} \cdot 
\Gc^{(i)}_{\beta_i\beta_i} G^{(i)}_{x\al_i}G^{(i)}_{\al_i y}\right]\ee
where $G^{(i)}:=G$ for $1\le i\le p/2$ and $G^{(i)}:=\overline{G}$ for $p/2< i\le p$ (recall $p\in 2\N$), contains $p$ internal molecules, corresponding to the vertices $\bal=(\al_1,\ldots, \al_p)$ and $\bbeta=(\beta_1,\ldots, \beta_p)$. During the expansions, the number of internal molecules never increases; it may decrease when two molecules---internal or external---are merged due to new dotted or waved edges created in the process. Consequently, $\Gamma_{\mu,xy}$ may contain strictly fewer molecules than $\Gamma_{xy}$. In particular, if all internal molecules merge with external ones, then $\Gamma_{\mu,xy}$ contains no internal molecules, in which case $q=0$. Finally, \eqref{eq:MolVW} follows directly from the definition of a molecule, since all vertices within a molecule are connected by paths of waved edges. 

For properties (3)--(5), note that the original graph $\Gamma_{xy}$ contains $p$ edge-disjoint paths between $\Mol_x$ and $\Mol_y$, each passing through an internal molecule $\Mol_i$ that contains the vertex $\al_i$, $i\in \qqq{p}$. We denote the path through $\Mol_i$ by $\fP_i$.
During the expansions, molecules may merge, but for notational convenience, we retain their original labels: if molecules $\Mol_i$ and $\Mol_j$ merge, we continue to refer to the resulting molecule as both  $\Mol_i$ and $\Mol_j$. Similarly, although the paths $\fP_i$ may change during the expansion, we keep their names and regard each $\fP_i$ as the path associated with $\Mol_i$. 
Checking \Cref{ssl,Oe14,T eq0}, we observe that these paths never disappear under local expansions. More precisely, if $e$ is a solid edge between two distinct molecules $\cal M$ and $\cal M'$, then in any local expansion either: 
\begin{itemize}
\item  the edge $e$ remains unaffected on the molecular graph; or %this edge is either unaffected (on the molecular graph) or one of the following scenarios occurs to it: 

\item $e$ is replaced by two new edges $e_1$ and $e_2$, which still form a connected path between $\Mol$ and $\Mol'$ (though this connectedness may fail at the vertex level); or %, and one edge may lie within a molecule); or
 
% The edge $e$ is replaced by two new edges $e_1$ and $e_2$, but these two edges still form a connected path between $\Mol$ and $\Mol'$ on the molecular graph. (Note that such connectedness may break on the vertex-level graph. Moreover, one of these edges may lie within a molecule.)

\item the molecules $\Mol$ and $\Mol'$ merge, in which case $e$ disappears on the molecular graph, but the connectedness between $\Mol$ and $\Mol'$ holds trivially. 

\end{itemize}
From this observation, properties (3) and (5) follow.

For property (4), let $\Mol$ be an internal molecule in $\Gamma_{\mu,xy}$. If $\Mol$ arises from merging at least two internal molecules, say $\Mol_i$ and $\Mol_j$, then both associated paths $\fP_i$ and $\fP_j$ pass through $\Mol$.
If the previous scenario does not occur, then we are in the case $\Mol=\Mol_i$, which corresponds to the following path in the original graph $\Gamma_{xy}$ (assuming, without loss of generality, that the path carries a $+$ charge):
\[ \sum_{\al_i,\beta_i}S_{\al_i\beta_i} \mathbf 1_{x\ne \al_i}\mathbf 1_{y\ne \al_i}\cdot \Gc_{\beta_i\beta_i} G_{x\al_i}G_{\al_i y}.\] 
To ensure that $\Gamma_{\mu,xy}$ is locally standard, the molecule $\Mol_i$ must pull in some $\overline G$ edges from other paths or molecules during the expansions. 
If a $\overline G$-edge from path $\fP_j$ is pulled to $\Mol_i$, then $\fP_j$ also passes through $\Mol_i$. Similarly, if a $\overline G$-edge inside $\Mol_j$ is pulled to $\Mol_i$, then the two new edges between $\Mol_i$ and $\Mol_j$ can be incorporated into path $\fP_j$, so that $\fP_j$ again passes through $\Mol_i$. In either case, property (4) holds.


Finally, we examine the scaling order of $\Gamma_{\mu,xy}$. For the original graph $\Gamma_{xy}$, the scaling order is
\be\label{eq:initial_scaling} \ord(\Gamma_{xy}) =
p, \ee
which follows from the presence of $2p$ off-diagonal solid edges, $p$ light-weights, $p$ waved edges, and $2p$ internal vertices. The scaling order itself is not monotone along the expansions; for instance, the contraction of a $2$-cycle discussed below \eqref{eq:GGraisesord} does not increase it. We therefore track a quantity that also accounts for the vertices that still have to be expanded, and that is minimized over all identifications of vertices that may occur later.

Let $\cal G$ be a normal graph. A \emph{partition} $\pi$ of $\cal G$ is a partition of its vertex set into classes; a class is \emph{external} if it contains an external vertex, and \emph{internal} otherwise. A solid edge is \emph{kept} by $\pi$ if it is a light-weight or joins two different classes, and $\cal G/\pi$ denotes the graph obtained from $\cal G$ by identifying the vertices in each class and removing the solid edges that are not kept.\footnote{Intuitively, $\pi$ describes an identification of vertices that may occur in later expansions: an off-diagonal solid edge whose two ends are identified may be replaced by its deterministic part $m$ or $\overline m$, and then loses its factor $\Psi_t$.} An internal class of $\pi$ is called \emph{elementary} if, as a vertex of $\cal G/\pi$, it either carries exactly one light-weight and no other solid edge, or carries no light-weight and exactly two solid edges, which have the same color (two $G$ edges or two $\overline G$ edges). We define the cost of $\pi$ and the local cost of $\cal G$ by
\be\label{eq:relateG1G0}
c_\pi(\cal G):=\ord(\cal G/\pi)+\#\{\text{elementary classes of }\pi\},\qquad \Phi(\cal G):=\min_{\pi}c_\pi(\cal G),
\ee
where $\ord(\cal G/\pi)=\#\{\text{solid edges kept by }\pi\}+2n_W(\cal G)-2\#\{\text{internal classes of }\pi\}$ as in \eqref{eq:ordG}, and we denote by $\Phi^{\rm far}(\cal G)$ the minimum of $c_\pi(\cal G)$ over the partitions $\pi$ in which $x$ and $y$ lie in different classes. These quantities have the following three properties.
\begin{itemize}
\item[(a)] Let $\pi_0$ be the partition of $\cal G$ into single vertices. Then $\cal G/\pi_0=\cal G$, and the elementary classes of $\pi_0$ are the internal vertices that carry either exactly one light-weight and no other solid edge, or exactly two solid edges of the same color and no light-weight. A locally standard graph has no such vertices. Hence, if $\cal G$ is locally standard and $x\ne y$, then $\ord(\cal G)=c_{\pi_0}(\cal G)\ge \Phi^{\rm far}(\cal G)$.

\item[(b)] We have $\Phi(\Gamma_{xy})=2p$ and $\Phi^{\rm far}(\Gamma_{xy})=3p$. Indeed, all internal vertices $\al_i$ and $\beta_i$ of $\Gamma_{xy}$ are elementary, so every internal class of a partition $\pi$ that consists of a single vertex is elementary. Hence $-2\#\{\text{internal classes}\}+\#\{\text{elementary classes}\}$ is at least minus the number of internal vertices lying in internal classes, and each $\beta_i$ contributes at least $-1$ to $c_\pi(\Gamma_{xy})$. Each $\al_i$ contributes at least $1$, unless its class contains both $x$ and $y$, in which case it contributes at least $0$: if $\al_i$ lies in an internal class, its two edges $G^{(i)}_{x\al_i}$ and $G^{(i)}_{\al_i y}$ are kept, and if it lies in an external class not containing both $x$ and $y$, at least one of them is kept. Together with the $p$ light-weights and the contribution $2p$ of the waved edges, we obtain
\[ c_\pi(\Gamma_{xy})\ge 2p+\#\{i\in\qqq{p}:\text{the class of $\al_i$ does not contain both $x$ and $y$}\}, \]
which gives $\Phi(\Gamma_{xy})\ge 2p$ and $\Phi^{\rm far}(\Gamma_{xy})\ge 3p$. Equality is attained by putting $x$, $y$, and all $\al_i$ in one class, and by $\pi_0$, respectively.

\item[(c)] $\Phi$ and $\Phi^{\rm far}$ do not decrease along the expansions: if $\cal G_1$ is a normal graph obtained from $\cal G_0$ by one step of \Cref{strat_local}, then $\Phi(\cal G_1)\ge\Phi(\cal G_0)$ and $\Phi^{\rm far}(\cal G_1)\ge\Phi^{\rm far}(\cal G_0)$.
\end{itemize}
To prove (c), given a partition $\pi_1$ of $\cal G_1$, we construct a partition $\pi_0$ of $\cal G_0$ with $c_{\pi_0}(\cal G_0)\le c_{\pi_1}(\cal G_1)$, such that $x$ and $y$ lie in different classes of $\pi_0$ whenever they do in $\pi_1$. Since the dotted edge partition only identifies vertices, and an identification of vertices amounts to passing to a coarser partition (with a cost that is not smaller), it suffices to consider the graphs on the right-hand sides of \eqref{Owx}, \eqref{Oe1x}, and \eqref{Oe2x} before the dotted edge partition. Each of these graphs is obtained from $\cal G_0$ by a bounded number of local operations at the expansion vertex: adding a new vertex joined to it by a waved edge, adding or removing a light-weight, moving an end of a solid edge or of a light-weight (this includes the action of the partial derivatives), and contracting a solid edge, i.e., identifying its two ends (this comes from the factors $\delta$ in \eqref{Oe1x} and \eqref{Oe2x}). Placing a new vertex in a class that contains at most one of its neighbors (counted with multiplicity) never lowers the cost, so it suffices to consider partitions $\pi_1$ in which each new vertex either forms a class of its own or lies in a class containing at least two of its neighbors. For each of the local operations, one checks directly that the partition $\pi_0$ induced by $\pi_1$ on the vertices of $\cal G_0$ (in which identified vertices are put into the same class) is a suitable choice, with one exception: when the two solid edges of a $2$-cycle $G_{uv}G_{vu}$ (or $\overline G_{uv}\overline G_{vu}$) at an internal vertex $u$ are removed without identifying $u$ and $v$, as in the second term of the example below \eqref{eq:GGraisesord}, the cost of the partition into single vertices may decrease by $1$. In this case, we put in addition $u$ into the class of $v$ in $\pi_0$. Since $u$ is internal, this never merges two external classes, so $x$ and $y$ remain in different classes. This finite case-by-case verification over the terms of \eqref{Owx}, \eqref{Oe1x}, and \eqref{Oe2x} replaces a bookkeeping of individual light-weights and vertices, which is not monotone along the expansions.

Applying (c) along the expansions in \Cref{lvl1 lemma} and using (a) and (b), we obtain that every locally standard graph $\Gamma_{\mu,xy}$ satisfies
\be\label{eq:ordGinter} \ord(\Gamma_{\mu,xy})\ge \Phi^{\rm far}(\Gamma_{\mu,xy})\ge \Phi^{\rm far}(\Gamma_{xy}) = 3p > 2p ,\ee
since $x\ne y$ under \eqref{eq:far_ab}. (Without \eqref{eq:far_ab}, the same argument with $\Phi$ in place of $\Phi^{\rm far}$ gives $\ord(\Gamma_{\mu,xy})\ge 2p$, also when $x=y$.) This concludes \eqref{eq:sizeGammamu}.
\end{proof}

\begin{remark}
The bound \eqref{eq:ordGinter} is in fact the stronger bound $\ord(\Gamma_{\mu,xy}) \ge 3p$, which in turn improves \eqref{eq:LW_moment} by replacing the factor $[\Psi_t(0)]^p$ with $[\Psi_t(0)]^{2p}$. We do not use this refinement, however, since the bottleneck of the main proof lies instead in the martingale estimate of \Cref{lem: EMn2_N} (see \Cref{rmk_bottleneck}).
\end{remark}



\subsection{Proof of \texorpdfstring{\Cref{lem:LWterm,lem: EWGn2_N}}{LemGraph1} for the block Anderson model}\label{subsec:LWchange-to-BA}

Analogously to the random band matrix model, the proofs of \Cref{lem:LWterm,lem: EWGn2_N} for the block Anderson model reduce to \Cref{lem:LW_moment,lem:LW_moment_exp}. In this reduction (see the beginning of \Cref{Sec:graph}), \Cref{lem_GbEXP_BA} gives, in place of \eqref{eq:directG1} and \eqref{eq:directG2},
\[ |(G_t(\sig'))_{yx}|\prec |(M(\sig'))_{yx}|+\Psi_t(|a-b|)\quad\text{and}\quad |(G_t(\sig'))_{yx}|\prec |(M(\sig'))_{yx}|+\sT_t(|a-b|\wedge\ell)+W^{-D}, \]
respectively, where $M(\sig')$ is defined in \eqref{def_mtzk}. Unlike $\delta_{xy}$, the term $|(M(\sig'))_{yx}|$ need not vanish for $x\ne y$; it is of order $\ilambda$ for some pairs in neighboring blocks. However, it is harmless after averaging over $x\in[a]$ and $y\in[b]$: by \eqref{eq:Psi3D} and \eqref{def_G0}, $M_{yx}\ne 0$ only if $x-W[x]=y-W[y]$, so at most $W^d$ of the $W^{2d}$ pairs $(x,y)$ contribute, each with $|M_{yx}|\le Ce^{-c|a-b|}$ by \eqref{Mbound_AO} or \eqref{Mbound_AO2}. Hence the additional contributions carry a factor $\OO(W^{-d}e^{-c|a-b|})$, which is bounded by $\Psi_t(|a-b|)$ by \eqref{eq:Psi}, and by $\sT_t(|a-b|\wedge\ell)$ since $\sT_t(0)\gtrsim W^{-d/2}$ and $\sT_t(0)e^{-cr}\lesssim\sT_t(r\wedge\ell)$. To establish \Cref{lem:LW_moment,lem:LW_moment_exp}, we introduce new graphical notations and corresponding local expansion rules.
In the setting of the block Anderson model, we continue to define the matrices $S^\pm$ as in \eqref{eq:def-Spm}, with \smash{$\Theta^{(+,+)}(z)$} defined in \eqref{def:Theta_BA} and $S^{(\sB)}=I_{L^d}$. 
In addition to the components introduced in \Cref{def_graph1}, we introduce new types of edges that represent the matrix entries of $\Psi$ and $M$, specific to this model: 
\begin{itemize}			
	
\item {\bf $\Psi$-dotted edge}: A black dotted edge labeled $\Psi$ between vertices $x$ and $y$ represents a factor $\ilambda\Psi_{xy}$.\footnote{Here $\Psi$ denotes the block-hopping matrix of the model, cf.~\eqref{eq:Psi3D}; it should not be confused with the control parameter $\Psi_t$ of \eqref{eq:Psi}.}

\item {\bf $M$-dotted edge}: A blue (resp.~red) dotted edge  labeled $M$ between vertices $x$ and $y$ represents a factor $M_{xy}$ (resp.~$\bar M_{xy}$).

\end{itemize}
With these new types of dotted edges, the definitions of molecules and molecular graphs remain the same as in \Cref{def_poly}. However, graphs in the block Anderson model carry an additional level of structure, which we refer to as \emph{atoms}. More precisely, our graphs exhibit microscopic structures, called atoms, which are equivalence classes of vertices connected via dotted edges. These atoms form mesoscopic structures within each molecule, and the global (macroscopic) structure is represented by the molecular graph.

% More precisely, our graphs have some microscopic structures, called \emph{atoms}, which are the equivalence classes of vertices connected through dotted edges, some mesoscopic structures within molecules, which are the equivalence classes of atoms connected through waved edges, and a global (macroscopic) structure given by the molecular graph. 

\begin{definition}[Atoms and atomic graphs]\label{def_atom}
We partition the set of all vertices in a graph into disjoint subsets called atoms. Two vertices belong to the same atom if and only if they are connected by a path consisting entirely of dotted edges---that is, any combination of regular dotted edges, $\Psi$-dotted edges, or $M$-dotted edges.
An atom is called external if it contains at least one external vertex; otherwise, it is called internal. 

Given a graph $\cal G$, we define its atomic graph as follows:
\begin{itemize}
\item Merge all vertices belonging to the same atom into a single vertex. 
		
\item Retain all solid and waved edges that connect different atoms. 
		
\item Discard all other components within $\cal G$, including $\times$-dotted edges, edges between vertices within the same atom, and coefficients.
\end{itemize}
\end{definition}

By the definition of (regular) dotted edges, the form of the $\Psi$-dotted edges in \eqref{eq:Psi3D}, and the bounds for the $M$-dotted edges in \eqref{Mbound_AO} or \eqref{Mbound_AO2}, we deduce that, up to an error of order $e^{-c(\log W)^{1+\e_1}}$, where $\e_1\in(0,1/10)$ is the constant fixed in \eqref{yixi}, %the following holds: %for any small constant $\e>0$
\be\label{scaleatom}
 \hbox{$x$, $y$ belong to the same atom} \implies  x-W[x] = y-W[y],\ \text{and} \ |[x]-[y]|\le (\log W)^{1+\e_1}.
\ee
Here, $[x],[y]\in \Zn$ are the block-level vertices as defined in \Cref{def: BM2}. 
The concept of atoms is introduced for two main purposes: first, to define the scaling size and scaling order of our graphs; and second, to enable a structural comparison between the block Anderson model and the random band matrix model. 
In particular, if every solid edge between two vertices of the same atom is retained as a solid self-loop of the merged vertex, then the atomic graphs in the block Anderson model have the same form as the vertex-level graphs in the random band matrix setting, where the $M$ matrices reduce to scalars. This correspondence transfers the \emph{statements} concerning molecular graphs from the random band matrix model to the block Anderson model, but not the \emph{arguments} step by step. A local expansion in the block Anderson model acts at a vertex rather than at an atom (for example, \eqref{GGGamma} below requires two solid edges at the same vertex, not merely in the same atom), and in general it is not the image of a single expansion step of the random band matrix model; the correspondence holds only for the outputs of the expansions, after grouping consecutive steps. The structural properties used below are therefore verified directly for the block Anderson expansions; see the end of \Cref{subsec:pf-LWterm_EXP} and the proof of \Cref{lem:LW_moment,lem:LW_moment_exp} below.


%One reason we introduce the concept of atoms is to define the scaling size and scaling order of our graphs. Another key reason is that our atomic graphs for the block Anderson model would have the same structure as the vertex-level graphs in the setting of the random band matrix model if we merge all vertices connected by dotted edges and discard all $\times$-dotted edges. In particular, all statements about the molecular graphs in the previous proofs for the random band matrix model will carry over verbatim to the setting of the block Anderson model.  



\begin{definition}[Normal graphs]  \label{defn_normalBA}  
We say a graph is \emph{normal} if it satisfies the following properties:
\begin{itemize}
\item[(i)] It contains at most $\OO(1)$ many vertices and edges.

\item[(ii)] There are no regular dotted edges between vertices (note that $\Psi$- or $M$-dotted edges are allowed).

\item[(iii)] Every solid edge carries a $\circ$ (i.e., it represents an entry of \smash{$\Gc$ or $\Gc^*$}). In particular, all weights are light-weights.

\end{itemize}
\end{definition}


Given an arbitrary graph with $\OO(1)$ many vertices and edges, we can decompose it into a linear combination of normal graphs by expanding each $G_{xy}$ (resp.~$G^*_{xy}$) edge into a \smash{$\Gc_{xy}$} edge plus an $M_{xy}$ edge (resp.~a \smash{$\Gc^*_{xy}$} edge plus an $M^*_{xy}$ edge), and by merging any vertices connected by regular dotted edges. For a normal graph, we define its scaling size in the same way as in \Cref{def scaling}, except that we replace the number of internal vertices with the number of internal atoms.


\begin{definition}[Scaling size and scaling order]\label{def scalingBA} 
%For a normal graph $\Gamma$, its scaling size is defined by
We define the scaling size of a normal graph $\Gamma$ as 
 \begin{align}
\size(\Gamma): = (L^d)^{n_M(\Gamma)}\cdot(\Psi_t)^{n_S(\Gamma)} \cdot W^{-d\left(n_W(\Gamma) - n_A(\Gamma) \right)},
  \label{eq_defsize_BA}
\end{align} 
where $n_S(\Gamma)$, $n_W(\Gamma)$, $n_A(\Gamma)$, and $n_M(\Gamma)$ denote the numbers of solid edges (including light-weights), waved edges, internal atoms, and internal molecules, respectively. The scaling order of $\Gamma$ is then defined as
\be\label{eq:ordG_BA} \ord(\Gamma):=n_S(\Gamma)+ 2 \left(n_W(\Gamma)  -  n_A(\Gamma)\right) . \ee
If $\Gamma$ can be expressed as a sum of $\OO(1)$ many normal graphs $\Gamma_k$, i.e., $\Gamma=\sum_k \Gamma_k$, we define its scaling size as in \eqref{eq_defsizemax}, and its scaling order by
\(\ord(\Gamma)=\min_k \ord(\Gamma_k).\)
\end{definition}
Next, we state the local expansion rules for the block Anderson model, as given in \cite{yang2024Del}.  

\begin{lemma}[Basic expansion, Lemma B.9 of \cite{yang2024Del}]\label{lanlw}
In the setting of the block Anderson model, let $f$ be a differentiable function of $G$. Then, we have the expansion
\begin{align} \label{eq:BE}
 \Gc_{xy} f(G) =_{\E}  &
   \sum_{\al,\beta}  M_{x\al} S_{\al \beta}
 \Gc_{\beta\beta} G _{\al y} f(G)
  - \sum_{\al,\beta}M_{x\al} S_{\al \beta} G _{\beta y} 
  \partial_{h_{ \beta\al}} f(G) .
 \end{align}
% where $Q$ is a sum of $Q$-graphs defined as
% \begin{align*}
% Q=-\sum_{\al,\beta}M_{x\al}S_{\al\beta}  Q_{\al}\left[\Gc_{\beta \beta}G_{\al y}\Gamma\right] +\sum_{\al,\beta}M_{x\al}S_{\al\beta} Q_{\al}\left[G_{\beta y}\partial_{h_{\beta\al}}\Gamma\right]+ \sum_{\al}M_{x\al}Q_\al\left[(M^{-1}\Gc)_{\al y} \bigGamma\right].
% \end{align*}
\end{lemma}

The purpose of the basic expansion is to expand any graph as a sum of graphs in which every vertex has a solid-edge degree $\in \{0,2\}$. If a vertex still carries self-loops (i.e., weights), we then apply the following weight expansion.

\begin{lemma} [Weight expansion, Lemma B.10 of \cite{yang2024Del}]\label{lem_lweight}
In the setting of the block Anderson model, let $f$ be a differentiable function of $G$. Then, we have the expansion
\begin{align} \label{eq:LW}
\Gc _{xx} f(G)  
=_{\E} &\sum_{y}(1+M^+S^+)_{xy}
 \Big(   \sum_{\al,\beta}  M_{y\al} S_{\al \beta} \Gc_{ \al y} 
 \Gc_{\beta\beta} f(G)
  - \sum_{\al,\beta}M_{y\al} S_{\al \beta} G _{\beta y} 
   \partial_{h_{ \beta\al}} f(G)\Big) ,
 \end{align}
where $M^+$ is the $N\times N$ matrix with entries $M^+_{xy}:=M_{xy}M_{yx}$. 
\end{lemma}



If there are vertices incident to two solid edges of the same charge, we use the following $GG$ expansion.
 
\begin{lemma}[$GG$ expansion, Lemma B.11 of \cite{yang2024Del}]
In the setting of the block Anderson model, let $f$ be a differentiable function of $G$. Then, we have the expansion
\begin{align} 
&\Gc _{y'x}\Gc _{xy} f(G) 
=_\E \sum_\beta (M^+S^+) _{x\beta } M _{\beta y}M_{y'\beta}f(G)
+ \sum_{ \beta}(M^+S^+)_{x\beta}\left(\Gc _{\beta y}M_{y'\beta}+M_{\beta y}\Gc _{y'\beta}\right)f(G)\label{GGGamma}
 \\ & +\sum_{w,\al,\beta}\left(1+M^+S^+\right)_{xw}M_{w\al} S_{\al \beta}\left(   
 \Gc_{\beta\beta} G _{\al y} \Gc _{y'w}f(G)
 + \Gc _{\al w}G _{\beta y}G_{y'\beta}
   f(G) -  G _{\beta y} \Gc _{y'w}
     \partial_{h_{ \beta\al}}   f(G)\right) .\nonumber
 \end{align}
\end{lemma}
The expansion of $\overline{\Gc_{y'x}}\,\overline{\Gc_{xy}} f(G)$ (two red solid edges) can be derived by taking the complex conjugate of \eqref{GGGamma}.

As explained in \cite[Appendix B]{yang2024Del}, repeated applications of the above local expansions to an arbitrary normal graph yield a sum of $\OO(1)$ \emph{locally standard graphs} (see \Cref{deflvl1}), subject to the additional requirement that all solid edges are \smash{$\Gc$} edges. This allows us to establish an analogue of \Cref{lvl1 lemma} for the block Anderson model. 
With the above preparations, we can now proceed to the proof of \Cref{lem:LW_moment,lem:LW_moment_exp} for the block Anderson model.
Since the argument is very similar to that in \Cref{Sec:graph}, we only sketch the proof and omit the repetitive details.

%With these preparations, we are now ready to give the proof of \Cref{lem:LW_moment,lem:LW_moment_exp} for the block Anderson model.  Due to close similarity with the arguments in \Cref{Sec:graph}, we will only outline the proof, without writing all details. 

%With the above preparations, we are ready to outline the proofs of \Cref{lem:LWterm,lem: EWGn2_N} for the block Anderson model. As for the random band matrix model, we need to establish \Cref{lem:LW_moment,lem:LW_moment_exp} for the block Anderson model. 

\begin{proof}[\bf Proof of \Cref{lem:LW_moment,lem:LW_moment_exp} for the block Anderson model]
We begin by applying the local expansions \eqref{eq:BE}, \eqref{eq:LW}, and \eqref{GGGamma} (and their complex conjugates) to the graph $|f_{xy}(G)|^p$, splitting every solid edge as $G=\Gc+M$. This yields the analogue of \Cref{lem:localregular} for the block Anderson model, with properties (1)--(6) read on the atomic graphs. In property (2), $n_V(\Mol_i)$ denotes the number of atoms in $\Mol_i$, in accordance with \Cref{def scalingBA}; counted with vertices, \eqref{eq:MolVW} fails already for the initial graph, since $M$-dotted edges join vertices of a molecule without waved edges. With atoms in place of vertices, (2) follows as in \Cref{sec:pflocalregular}, since the atoms of a molecule are connected by paths of waved edges.

These properties are not inherited from the random band matrix model step by step (see the discussion below \eqref{scaleatom}), and we verify them directly. For the path properties (3)--(5), one checks term by term that every expansion step, including the splitting $G=\Gc+M$, acts on the molecular graph in one of the three ways listed in \Cref{sec:pflocalregular}: a solid edge between two molecules is unaffected, or it is replaced by two edges forming a path between the same molecules, or it disappears because the two molecules merge. Hence the family of $p$ edge-disjoint paths from $\Mol_x$ to $\Mol_y$ (which may revisit molecules), labeled as in \Cref{sec:pflocalregular}, is preserved. This includes the terms of \eqref{GGGamma} that have no single counterpart in \eqref{Oe2x}, namely its first sum and the terms in which both expanded solid edges become edges inside one molecule.

Property (6) is used, and proved, when $x$ and $y$ lie in different atoms, which holds under \eqref{eq:far_ab} up to a negligible error by \eqref{scaleatom}. For (6), we use the local cost $\Phi^{\rm far}$ of \Cref{sec:pflocalregular}, computed on the atomic graphs, in which solid edges between different atoms are treated as off-diagonal solid edges and solid edges inside an atom as light-weights. A locally standard graph has no elementary atoms, since each of its vertices carries either no solid edge or one blue and one red solid edge. The splitting $G=\Gc+M$ of a solid edge between two atoms replaces it by an $M$-dotted edge and merges the two atoms; this does not decrease $\Phi^{\rm far}$, because every partition of the new atomic graph comes from a partition of the old one in which the two atoms lie in the same class, the removed edge is not kept by the latter, and the two quotient graphs coincide. Most of the remaining terms act on the atomic graphs as compositions of the local operations considered in \Cref{sec:pflocalregular}; for the terms that have no counterpart in \eqref{Owx}, \eqref{Oe1x}, and \eqref{Oe2x} (for instance, the first sum in \eqref{GGGamma}), the same comparison of partitions is carried out directly. Hence $\ord(\Gamma_{\mu,xy})\ge 3p>2p$ whenever $x$ and $y$ lie in different atoms.




Next, we define the auxiliary graph in a similar manner as in \Cref{def_auxgraph}, with minor modifications that we now describe. Let $\cal G_{xy}$ be a locally standard graph obtained from the local expansions, with $q$ internal molecules $\Mol_i$, $i\in\qqq{q}$, and two external molecules $\Mol_x$ and $\Mol_y$, as in \Cref{def_auxgraph}. For each $i\in \{1,\ldots, q,x,y\}$, we fix a center $\al_i\in \Mol_i$, and we also choose centers $x_{i,j}$ for the atoms $\cal A_{i,j}\subset\Mol_i$, where $j\in\qqq{k_i}$ and $k_i$ denotes the number of atoms in $\Mol_i$. Without loss of generality, we set $\al_i\equiv x_{i,1}$ and assume that $\al_i$ is connected to other molecules via solid edges.
(By the definition of locally standard graphs, the solid-edge degree of each $\al_i$ is at least 2.)
If $\p{\beta\to \beta'}$ is a solid edge connecting the atoms $\cal A_{i,j}$ and $\cal A_{i',j'}$, then by \eqref{scaleatom}, the form of the $\Psi$-dotted edges in \eqref{eq:Psi3D}, and the exponential decay estimates \eqref{Mbound_AO} or \eqref{Mbound_AO2} for the $M$-dotted edges, we obtain
\be\label{Gyiyjatom}
\big|(\Gc_t)_{\beta\beta'}\big|\prec \zeta(x_{i,j},x_{i',j'}):=\sum_{|[a]|+|[b]|\le (\log W)^{1+\e_1}}\big|(\Gc_{t})_{x_{i,j}+W[a],x_{i',j'}+W[b]}\big| + W^{-D}
\ee
for any large constant $D>0$. This quantity satisfies bounds analogous to \eqref{eq:Gbyxi2}:
\be\label{eq:Gbyxi2_BA}
\zeta(\al,\beta) \prec \Psi_t(|[\al]-[\beta]|),\quad \sum_{\al\in \ZL}(|\zeta(\al,\beta)|^2+|\zeta(\beta,\al)|^2) \prec \eta_t^{-1},
\ee
where the first estimate follows from \Cref{lem_GbEXP_BA}, and the second from Ward's identity for $G$ together with \eqref{Mbound_AO} or \eqref{Mbound_AO2} for $M$. If $\p{\beta,\beta'}$ is a waved edge connecting the atoms $\cal A_{i,j}$ and $\cal A_{i',j'}$, then by \eqref{scaleatom}, the definition of $S$ in \eqref{bandcwV}, and the estimate \eqref{prop:ThfadC_short}, we obtain
\be\label{Gyiyjatom2}
\big|S_{\beta\beta'}\big|+ \big|S^+_{\beta\beta'}\big|\prec W^{-d}\mathbf 1\p{|x_{i,j}-x_{i',j'}|\le W(\log W)^{1+\e_1}} + W^{-D}.\ee


We first define a graph $\cal G^{a}_{xy}$ as an ``atomic reduction'' of $\cal G_{xy}$:
\begin{itemize}
    \item Its vertices are the atom centers $x_{i,j}$ for $i\in\{1,\ldots,q,x,y\}$ and $j\in\qqq{k_i}$.
    
    \item For each solid edge in $\cal G_{xy}$ from the atom $\cal A_{i,j}$ to the atom $\cal A_{i',j'}$, we introduce an oriented solid edge from $x_{i,j}$ to $x_{i',j'}$ in $\cal G^{a}_{xy}$, representing the factor $\zeta(x_{i,j},x_{i',j'})$.

\item For each waved edge in $\cal G_{xy}$ connecting the atoms $\cal A_{i,j}$ and $\cal A_{i',j'}$, we introduce a waved edge between $x_{i,j}$ and $x_{i',j'}$ in $\cal G^{a}_{xy}$, representing the factor $W^{-d}\mathbf 1 \left(|x_{i,j}-x_{i',j'}|\le W(\log W)^{1+\e_1}\right)$.

\end{itemize}
Next, we further simplify the structure of $\cal G^{a}_{xy}$ and define the auxiliary graph $\cal G^{\aux}_{xy}$ as follows.
First, we remove all solid and waved edges that lie entirely within individual molecules of $\cal G^{a}_{xy}$.
Second, for each molecule $\Mol_i$, we retain only those vertices (including the center $\al_i$) that connect to other molecules through solid edges; all remaining vertices inside $\Mol_i$ are discarded. 
Third, for each retained vertex $x_{i,j}\ne \al_i$ inside $\Mol_i$, we add a new waved edge connecting $x_{i,j}$ to $\al_i$, representing the factor $W^{-d}\mathbf 1(|x_{i,j}-\al_i|\le W(\log W)^{1+2\e_1})$.
In other words, the molecular structure is reduced to the following simplified form:
\be  \nonumber%\label{simplemolecule} 
\parbox[c]{0.8\linewidth}{
\begin{center}
\scalebox{0.9}{
\begin{tikzpicture}
  \begin{feynman}
    
    % Internal fermion vertices    
    \vertex [dot] (a1) at (0,0.7);
    \fill[black] (a1) circle (2pt) node at (0.2,0.5) {$\al_i$};
    \vertex [dot] (a) at (0,0);   
    \fill[black] (a) circle (2pt) node[left=2pt]{};
      \vertex [dot] (a2) at (-0.7,0.7);
    \fill[black] (a2) circle (2pt) node[left=2pt]{$x_{i,2}$};
    \vertex [dot] (a3) at (0,1.4);   
    \fill[black] (a3) circle (2pt) node[left=2pt]{$x_{i,3}$};
     \vertex [dot] (a4) at (0.7,0.7);   
    \fill[black] (a4) circle (2pt) node[right=2pt]{$x_{i,4}$};
  \vertex (a5) at (0.5,0.2) {$\cdots$};
    
    % Draw fermion lines
    \diagram*{
      %(L) -- [fermion] (a) -- [fermion] (b) -- [fermion] (R),
      %(L) -- [fermion] (c) -- [fermion] (d) -- [fermion] (R),

      % Photon exchange lines
      (a) -- [photon] (a1), (a1) -- [photon] (a2), (a1) -- [photon] (a3), (a1) -- [photon] (a4)

      % Tadpoles
      %(a) -- [photon, half left, looseness=1.5] (ta) -- [photon, half left, looseness=1.5] (a),
      %(c) -- [photon, half left, looseness=1.5] (tc) -- [photon, half left, looseness=1.5] (c),
    };
  \end{feynman}
\end{tikzpicture}}
\end{center}}\ee
For the auxiliary graph $\cal G^{\aux}_{xy}$ constructed above, we define its scaling order by
\be\label{eq:ordGaux_BAM} \ord(\cal G^{\aux}_{xy}):=\#\{\text{solid edges in } \cal G^{\aux}_{xy}\}+2\#\{\text{waved edges in } \cal G^{\aux}_{xy}\}-2\#\{\text{internal vertices in } \cal G^{\aux}_{xy}\}.
\ee
From the construction, together with \eqref{eq:Gbyxi2_BA} and by repeating the argument used in the proof of \Cref{GtoAG}, we can control $\cal G_{xy}$ via its auxiliary graph as
\be\label{G_by_auxG_BA} 
\cal G_{xy} \prec  (\Psi_t)^{\ord(\cal G_{xy}) - \ord(\cal G_{xy}^{\aux})}  \cdot \cal G^{\aux}_{xy} + W^{-D}\, .
\ee
Finally, using \eqref{eq:Gbyxi2_BA}, the auxiliary graph $\cal G^{\aux}_{xy}$ can be bounded by an argument parallel to that in \Cref{sec:pf_Anp}.
In particular, the argument there, when applied to the molecular graph of $\cal G^{\aux}_{xy}$, provides a systematic procedure for selecting the long solid edges between molecules and for establishing a nested summation order over the internal molecules.
The only difference from the random band matrix case is that the ``summation over a molecule'' here also includes summing over certain non-center vertices within the molecule. However, only two solid edges incident to $\Mol_i$ are used in the summation step, while every other solid edge contributes either a long-edge factor or a $\Psi_t$-factor.
%when summing over vertices in $\Mol_i$ we retain only two solid edges attached to the center $\al_i$ and bound all other solid edges by $\Psi_t$-factors.
Applying the Cauchy-Schwarz inequality together with the bound \eqref{eq:Gbyxi2_BA} to these two solid edges, the summation over all vertices in the molecule $\Mol_i$ produces the desired factor of $\eta_t^{-1}$.
With this modification, the argument of \Cref{sec:pf_Anp} carries over verbatim, completing the proof of \Cref{lem:LW_moment} for the block Anderson model.

%The only difference from the random band matrix setting is that a ``summation over a molecule'' now also involves summing over some non-center vertices inside a molecule. However, recall that only two solid edges incident to $\Mol_i$ will be used in the summation. Hence, when summing over the vertices in $\Mol_i$, we only keep two solid edges incident to $\al_i$, while bounding other solid edges connected with $\Mol_i$ by $\Psi_t$ factors. Then, using \eqref{eq:Gbyxi2_BA}, the summation over $\al_i$ provides the desired $\eta_t^{-1}$ factor. With this modification in mind, the argument in \Cref{sec:pf_Anp} completes the proof of \Cref{lem:LW_moment} for the block Anderson model.

% Each waved edge now represents a factor of \smash{$W^{-d}\mathbf 1_{|\beta_k-\al_i|\le W(\log W)^{3/2}}$} up to an error of order \smash{$\exp(-c(\log W)^{3/2})$}, as implied by \eqref{scaleatom} and \eqref{scalemole}. In other words, such a contribution corresponds to the summation  
% \[W^{-kd} \sum_{\al_i}\sum_{\beta_1,\ldots, \beta_k} \mathbf 1\Big(\max_{j=1}^k|\beta_j-\al_i|\le W(\log W)^{3/2}\Big) .\] 

% For each auxiliary graph constructed as above, every solid edge $|\Gc_{\al\beta}|$ connecting different molecules satisfies an estimate analogous to \eqref{eq:Gbyxi2}:
% \be\label{eq:Gbyxi2_BA}
% |\Gc_{\al\beta}| \prec \Psi_t(|[\al]-[\beta]|)+\exp\left(-c_0(\log W)^{3/2}\right),\quad \sum_{\al}|\Gc_{\al\beta}|^2 \prec \eta_t^{-1},
% \ee
% where the first bound follows from \Cref{lem_GbEXP_BA}, and the second is obtained via Ward's identity for $G$, together with the bound \eqref{Mbound_AO} or \eqref{Mbound_AO2} for $M$. 

%except that we modify the definition of the $\xi$-variables slightly. Moreover, by the definition of $S$ in \eqref{bandcwV} and the bound \eqref{prop:ThfadC_short}, we can again assume \eqref{yixi} for $\beta_i\sim_{\Mol}\al_i$. By the definition of $S$ in \eqref{bandcwV} and the bound \eqref{prop:ThfadC_short}, we can again assume \eqref{yixi} for $\beta_i\sim_{\Mol}\al_i$. Thus, we can further bound 

% Specifically, given a graph $\cal G_{xy}$, we still select a molecular center $\al_i\in \cal M_i$ for each $i\in\{1,\ldots, q, x, y\}$. For each molecule, we bound every waved edge within it by a factor of $\OO_\prec(W^{-d})$, and every solid edge within it by a $\Psi_t$ factor. By \eqref{scaleatom}, each atom contributes a factor of $\OO_\prec(W^{d})$ in the summation. It remains to control the solid edges between molecules, say $\Mol_i$ and $\Mol_j$. 
% For any atoms $\cal A_i\subset \Mol_i$ and $\cal A_j\subset \Mol_j$, we fix their centers as $x_i\in \cal A_i$ and $x_j\in \cal A_j$. Then, by \eqref{scaleatom}, the form of $\Psi$-dotted edges in \eqref{eq:Psi3D}, and the bound \eqref{Mbound_AO} or \eqref{Mbound_AO2} for $M$-dotted edges, we see that any solid edge between vertices $y_i\in \cal A_i$ and $y_j\in \cal A_j$ can be bounded by 
% \[\zeta(x_i,x_j):=\sum_{[a],[b]\in \ZL}e^{-c(|[a]|+|[b]|)}\p{\big|\Gc_{t}\p{x_i+W[a],x_j+W[b]}\big|+\big|\Gc_{t}\p{x_j+W[b],x_i+W[a]}\big|} .\]


%Within each molecule, we bound the solid edges by $\Psi_t$ factors, while solid edges between molecules are retained and each is bounded by its absolute value, e.g., $|\Gc_{\al\beta}|$ for an edge $\Gc_{\al\beta}$. 
 



For the proof of \Cref{lem:LW_moment_exp}, the bounds in \eqref{eq:Gbyxi2_BA} remain valid with $\Psi_t(|[\al]-[\beta]|)=\sT_t(|[\al]-[\beta]|\wedge\ell)+W^{-D}$. Then, by following the same reasoning as in \Cref{subsec_pf_LW_moment_exp}, we complete the proof of \Cref{lem:LW_moment_exp} for the block Anderson model.
\end{proof}




\bibliographystyle{abbrv}
\bibliography{Jun}




\end{document}